\documentclass[11pt,openany]{book}
\setcounter{secnumdepth}{3}
\usepackage[margin=1.1in]{geometry}
\usepackage[T5]{fontenc}
\usepackage[utf8]{inputenc}
\usepackage[vietnamese]{babel}
\usepackage{amsmath,amssymb,amsthm}
\usepackage{graphicx}
\usepackage{tikz}
\usetikzlibrary{arrows.meta,decorations.markings,decorations.pathmorphing,decorations.pathreplacing,calc}
\input{figures/icons}
\usepackage[unicode,colorlinks=true,linkcolor=blue!60!black,citecolor=blue!60!black]{hyperref}
\usepackage{microtype}
\setlength{\emergencystretch}{3em} % long names (licences, references) in narrow lines
\usepackage{empheq}
\usepackage{array}
\usepackage{booktabs}
\usepackage{multirow}
\usepackage{longtable}
\usepackage{circuitikz}
\definecolor{keyyellow}{rgb}{1.0,0.97,0.78}
% key equations: \begin{empheq}[box=\keybox]{equation} ... \end{empheq} -- light-yellow box, number stays outside
\newcommand{\keybox}[1]{{\setlength{\fboxsep}{7pt}\colorbox{keyyellow}{#1}}}
\usepackage[most]{tcolorbox}
% key statements (propositions etc.): the same light-yellow background, breakable across pages
\newtcolorbox{keyblock}{enhanced,breakable,colback=keyyellow,colframe=keyyellow,boxrule=0pt,arc=0pt,left=6pt,right=6pt,top=2pt,bottom=2pt,boxsep=0pt}
\renewcommand{\chaptermark}[1]{\markboth{\small\chaptername~\thechapter.\ #1}{}}
\renewcommand{\sectionmark}[1]{\markright{\small\thesection.\ #1}}
\renewcommand{\topfraction}{0.95}
\renewcommand{\textfraction}{0.05}
\renewcommand{\floatpagefraction}{0.75}

% part pages: title at the top third, and text may follow on the same page (used for the part introductions)
\makeatletter
\renewcommand\part{\if@openright\cleardoublepage\else\clearpage\fi\thispagestyle{plain}\if@twocolumn\onecolumn\@tempswatrue\else\@tempswafalse\fi\null\vspace{0.18\textheight}\secdef\@part\@spart}
\renewcommand\@endpart{\par\vspace{3em}\if@tempswa\twocolumn\fi}
\makeatother
\newtheorem{theorem}{Định lý}[chapter]
\newtheorem{proposition}{Mệnh đề}[chapter]
\newtheorem{lemma}{Bổ đề}[chapter]
\theoremstyle{remark}
\newtheorem{remark}{Nhận xét}[chapter]
% exercises: \begin{exercise}[\lvlB] ... \end{exercise}, numbered "Exercise 4.3"; difficulty marks
\theoremstyle{definition}
\newtheorem{exercise}{Bài tập}[chapter]
\newcommand{\lvlA}{$\star$}
\newcommand{\lvlB}{$\star\star$}
\newcommand{\lvlC}{$\star\star\star$}
% answer to an exercise in the Answers appendix: \answer{ex:NN:k}{text}
\newcommand{\answer}[2]{\par\noindent\textbf{\ref{#1}.}~#2\par\smallskip}
\newcommand{\dd}{\mathrm{d}}
% neuron type code: first four letters of the primary mediator
\newcommand{\nt}[1]{\textsf{\textbf{#1}}}
\newcommand{\mV}{\,\text{mV}}
\newcommand{\ms}{\,\text{ms}}
\newcommand{\mM}{\,\text{mM}}
\newcommand{\um}{\,\text{\textmu m}}
\newcommand{\ion}[2]{\mathrm{#1}^{#2}}
\newcommand{\Na}{\ion{Na}{+}}
\newcommand{\K}{\ion{K}{+}}
\newcommand{\Cl}{\ion{Cl}{-}}
\newcommand{\Ca}{\ion{Ca}{2+}}
\newcommand{\conc}[1]{[\mathrm{#1}]}

\title{\Huge Chiếc máy tính 20 oát\\[16pt] \Large Bộ não tính toán như thế nào: cách tiếp cận định lượng dành cho kỹ sư}
\hypersetup{pdftitle={Chiếc máy tính 20 oát. Bộ não tính toán như thế nào: cách tiếp cận định lượng dành cho kỹ sư},pdfauthor={Konstantin Kladko}}
\author{\Large Konstantin~Kladko}
\date{}

\begin{document}
\frontmatter
\maketitle
\thispagestyle{empty}
\vspace*{\fill}
\noindent{\small \copyright~2026 Konstantin~Kladko.\\[4pt]
Văn bản, hình vẽ, bài tập và đáp án được phát hành theo giấy phép Creative Commons Ghi công -- Phi thương mại -- Chia sẻ tương tự 4.0 Quốc tế (CC BY-NC-SA 4.0): \url{https://creativecommons.org/licenses/by-nc-sa/4.0/deed.vi}. Có thể tự do sao chép, dùng trong giảng dạy, dịch và chỉnh sửa cuốn sách cho mục đích phi thương mại, với điều kiện ghi tên tác giả và phát hành các tác phẩm phái sinh theo cùng giấy phép.\\[4pt]
Các chương trình mô hình và biên dịch được phát hành theo giấy phép MIT.\\[4pt]
Mã nguồn, mô hình và phiên bản mới: \url{https://github.com/kladkogex/20-watt-computer}.}
\clearpage

\tableofcontents
\chapter*{Cách đọc cuốn sách này}
\addcontentsline{toc}{chapter}{Cách đọc cuốn sách này}

Không nhất thiết phải đọc sách theo thứ tự. Các chương~1--4 là nền tảng chung: các loại nơ-ron, \nt{GLUT} và \nt{GABA} tính gì và chúng nói bằng ngôn ngữ nào. Từ đó cuốn sách rẽ nhánh. Bản đồ trên hình~\ref{fig:roadmap} cho thấy chương nào dựa trên chương nào: mũi tên từ chương $A$ đến chương $B$ nghĩa là $B$ dùng các khái niệm và công thức của $A$. Bản đồ chỉ vẽ các phụ thuộc chính; những tham chiếu ngược nhỏ không cản trở việc đọc.

\begin{figure}[h]
\centering
\begin{tikzpicture}[>=Stealth,scale=0.95,
  ch/.style={draw,thick,rounded corners=3pt,align=center,font=\scriptsize,text width=1.85cm,minimum height=0.62cm,inner sep=2pt},
  base/.style={ch,fill=blue!8}, bus/.style={ch,fill=orange!14}, sum/.style={ch,fill=gray!18},
  io/.style={ch,fill=green!10}, sort/.style={ch,fill=cyan!8}, lab/.style={ch,fill=violet!10},
  dep/.style={->,thick,gray!60!black}]
% foundation
\node[base] (c1) at (0,0) {1 Các loại\\nơ-ron};
\node[base] (c2) at (2.3,0) {2 \nt{GLUT}};
\node[base] (c3) at (4.6,0) {3 \nt{GLUT} và \nt{GABA}};
\node[base] (c4) at (6.9,0) {4 Mã Morse};
% broadcast channels and the synthesis
\node[bus] (c5) at (4.6,-1.5) {5 \nt{SERO}};
\node[bus] (c6) at (7.3,-1.5) {6 \nt{DOPA}};
\node[bus] (c7) at (9.2,-0.9) {7 Các thuyết\\dopamine};
\node[bus] (c8) at (9.2,-2.1) {8 \nt{NORA}};
\node[bus] (c9) at (11.5,-2.1) {9 \nt{ACET}};
\node[sum] (c10) at (13.8,-0.9) {10 Toàn bộ\\cỗ máy};
% inputs and outputs
\node[io] (c12) at (2.3,-4.1) {12 Nhận\\dạng};
\node[io] (c11) at (4.6,-4.1) {11 Đầu vào};
\node[io] (c13) at (6.9,-4.1) {13 Cơ bắp};
\node[io] (c14) at (9.2,-3.05) {14 Đau};
\node[io] (c15) at (9.2,-4.1) {15 Học\\vận động};
\node[io] (c16) at (9.2,-5.15) {16 Hóa học\\cơ thể};
% lab, sorting, sleep
\node[lab] (c17) at (11.5,-4.1) {17 Gỡ lỗi};
\node[lab] (c21) at (13.8,-4.1) {21 Cỗ máy\\sống};
\node[lab] (c22) at (13.8,-5.15) {22 DishBrain};
\node[sort] (c18) at (11.5,-5.15) {18 Nơ-ron\\dụng cụ};
\node[sort] (c19) at (13.8,-6.2) {19 Số\\đuôi gai};
\node[bus] (c20) at (11.5,-6.2) {20 Giấc ngủ};
% dependencies
\draw[dep] (c1) -- (c2); \draw[dep] (c2) -- (c3); \draw[dep] (c3) -- (c4);
\draw[dep] (c1) -- (c5); \draw[dep] (c4) -- (c6); \draw[dep] (c5) -- (c6);
\draw[dep] (c6) -- (c7); \draw[dep] (c6) -- (c8); \draw[dep] (c8) -- (c9);
\draw[dep] (c7) -- (c10); \draw[dep] (c9) -- (c10);
\draw[dep] (c4) -- (c11); \draw[dep] (c11) -- (c12); \draw[dep] (c11) -- (c13); \draw[dep] (c5) -- (c13);
\draw[dep] (c13) -- (c14); \draw[dep] (c13) -- (c15); \draw[dep] (c13) -- (c16);
\draw[dep] (c15) -- (c17); \draw[dep] (c9) -- (c17); \draw[dep] (c17) -- (c21); \draw[dep] (c21) -- (c22);
\draw[dep] (c16) -- (c18); \draw[dep] (c18) -- (c19); \draw[dep] (c16) -- (c20);
% legend
\begin{scope}[yshift=-7.25cm]
\foreach \s/\t [count=\i] in {base/nền tảng, bus/kênh và chế độ, sum/tổng hợp, io/đầu vào và đầu ra, sort/phân loại nơ-ron, lab/phòng thí nghiệm}{
  \pgfmathtruncatemacro{\col}{mod(\i-1,3)}\pgfmathtruncatemacro{\row}{(\i-1)/3}
  \node[\s,text width=0.3cm,minimum height=0.32cm] at ({\col*4.8+0.2},{-\row*0.55}) {};
  \node[anchor=west,font=\scriptsize] at ({\col*4.8+0.55},{-\row*0.55}) {\t};}
\end{scope}
\end{tikzpicture}
\caption{Bản đồ các chương. Mũi tên đi từ một chương đến chương dựa trên nó. Các tham chiếu khác giữa các chương chỉ là chỉ dẫn, không phải phụ thuộc.}
\label{fig:roadmap}
\end{figure}

\section*{Các lộ trình đọc}

\begin{center}
\footnotesize
\setlength{\tabcolsep}{4pt}
\renewcommand{\arraystretch}{1.2}
\begin{tabular}{>{\raggedright}p{3.0cm}>{\raggedright}p{3.9cm}>{\raggedright\arraybackslash}p{6.4cm}}
\toprule
Lộ trình & Dành cho ai & Các chương theo thứ tự \\
\midrule
khóa học một quý & giảng viên, 10 tuần & tuần 1: ch.~1--2; 2: ch.~3; 3: ch.~4; 4: ch.~5--6; 5: ch.~7--8; 6: ch.~9; 7: ch.~10; 8: ch.~11--12; 9: ch.~13 và~15; 10: ch.~17 \\
AI và học máy & người biết mạng nơ-ron nhân tạo và muốn hiểu bộ não học như thế nào & 1--4, 5 (đọc lướt), 6, 7, 8, 9, 11 (đọc lướt), 12, 13 (đọc lướt), 15 \\
điện tử và phần cứng & kỹ sư thiết kế mạch, người phát triển chip neuromorphic & 1--4, 5, 11, 13, 16, 18, 19, 17 (chương 9 và 15 đọc lướt) \\
giao diện thần kinh và máy tính sống & người xây dựng giao diện não--máy tính và chip có nơ-ron sống & 1, 2, 4, 11, 13, 17 (chương 9 và 15 đọc lướt), 21, 22 (chương 12 đọc lướt) \\
tổng quan nhanh & kỹ sư có một ngày cuối tuần & 1, 2, 3, 6, 10; các tham chiếu của chương~6 đến chương~4--5 hiểu được từ ngữ cảnh \\
\bottomrule
\end{tabular}
\end{center}

``Đọc lướt'' nghĩa là: chỉ cần đọc phần đầu chương, các mệnh đề trong khung vàng và bảng tổng kết.

Cuối mỗi chương có bài tập: ước tính, suy luận, mô phỏng và thiết kế; đáp án ngắn được tập hợp trong phụ lục~\ref{sec:answers}. Mọi ký hiệu nằm trong phụ lục~\ref{sec:notation}, còn các thí nghiệm làm nền cho những mệnh đề chính của từng chương nằm trong phụ lục~\ref{sec:supplement}. Các con số trong sách được tính bằng những chương trình nhỏ trong thư mục \texttt{models/}; có thể chạy và sửa chúng.

\mainmatter

\chapter{Các loại nơ-ron}
\label{sec:neurontypes}

\section{Nơ-ron như một hộp đen}

Giả sử bạn nhận được một bao tải chứa tám mươi sáu tỷ nơ-ron~\cite{Azevedo2009} và được yêu cầu chia chúng thành từng đống. Bạn sẽ làm thế nào?

Một kỹ sư nhận được bao tải vi mạch lạ sẽ không phân loại chúng theo hình dạng vỏ. Anh ta sẽ hỏi: linh kiện có bao nhiêu đầu vào, bao nhiêu đầu ra, và đưa ra đầu ra cái gì. Hãy vẽ nơ-ron theo cách đó --- như một khối trên sơ đồ (hình~\ref{fig:blackbox}).

\begin{figure}[t]
\centering
\begin{tikzpicture}[>=Stealth,x=1cm,y=1cm]
% the box
\node[draw,very thick,minimum width=2.6cm,minimum height=2.6cm,fill=blue!6] (N) at (0,0) {};
\node at (0,0.55) {\small nơ-ron};
\node at (0,-0.15) {$\displaystyle u=\sum_i w_i x_i$};
\node at (0,-0.8) {\small $y=\Theta(u-\theta)$};
% inputs
\foreach \k/\y/\lab in {1/1.05/{x_1},2/0.55/{x_2},3/0.05/{x_3}}{
  \draw[->,thick,blue!60!black] (-4.2,\y) -- (-1.3,\y);
  \node[left] at (-4.2,\y) {$\lab$};
  \node[above,blue!60!black] at (-2.1,\y-0.06) {\scriptsize $w_{\k}$};}
\node at (-4.45,-0.4) {$\vdots$};
\node at (-2.1,-0.4) {$\vdots$};
\draw[->,thick,blue!60!black] (-4.2,-1.05) -- (-1.3,-1.05);
\node[left] at (-4.2,-1.05) {$x_n$};
\node[above,blue!60!black] at (-2.1,-1.11) {\scriptsize $w_{n}$};
\draw[decorate,decoration={brace,amplitude=5pt},gray!60!black] (-4.9,-1.2) -- (-4.9,1.2);
\node[left,align=right,gray!40!black] at (-5.1,0) {\small $n\sim10^3$--$10^5$\\[-2pt]\small đầu vào};
% single output wire, then fan-out
\draw[very thick,red!70!black] (1.3,0) -- (3.2,0);
\foreach \x in {1.6,2.2,2.34,2.9}{\draw[thick,red!70!black] (\x,0) -- (\x,0.4);}
\node[above right,font=\tiny\itshape,gray!30!black,inner sep=1pt] at (1.42,0.45) {tách\dots\ tách-tách\dots};
\node[below,red!70!black,align=center] at (2.25,-0.05) {\scriptsize một đầu ra $y$};
\fill[red!70!black] (3.2,0) circle (0.06);
\foreach \y in {1.3,0.65,0,-0.65,-1.3}{
  \draw[->,thick,red!70!black] (3.2,0) -- (5.0,\y);}
\draw[decorate,decoration={brace,amplitude=5pt},gray!60!black] (5.3,1.45) -- (5.3,-1.45);
\node[right,align=left,gray!40!black] at (5.5,0) {\small bản sao của $y$\\[-2pt]\small tới $10^3$--$10^4$\\[-2pt]\small nơ-ron khác};
\end{tikzpicture}
\caption{Nơ-ron qua con mắt kỹ sư. Nhiều đầu vào $x_i$, mỗi đầu vào có trọng số $w_i$ riêng; bên trong là tổng có trọng số $u$ và phép so sánh với ngưỡng $\theta$ ($\Theta$ là hàm bậc thang: bằng $1$ nếu đối số dương và bằng $0$ trong trường hợp còn lại); một đầu ra $y$ là chuỗi xung, được nhân bản qua các nhánh tới hàng nghìn nơi nhận.}
\label{fig:blackbox}
\end{figure}

Ở đầu ra của khối không phải là một con số, mà là chuỗi xung điện áp ngắn: ``tách'', khoảng lặng, ``tách-tách'', khoảng lặng dài. Mọi xung đều cao như nhau, nên toàn bộ thông tin nằm ở chỗ \emph{khi nào} chúng xuất hiện. Bên trong khối, ở phép gần đúng thô đầu tiên, là một bộ cộng và một bộ so sánh: các đầu vào được cộng với trọng số, và nếu tổng vượt ngưỡng, khối phát ra một xung. Ở chương sau ta sẽ thay nó bằng một mô hình thực sự làm việc trong thời gian liên tục.

Đầu ra chỉ có một, nhưng dây dẫn phân nhánh, và mỗi nhánh mang \emph{cùng một} bản sao tín hiệu. Trong điện tử, điều này gọi là hệ số phân nhánh đầu ra, fan-out. Ở nơ-ron vỏ não, con số này cỡ vài nghìn; một cổng logic trong vi mạch thường chỉ nối tới vài cổng khác.

Vậy \emph{cái gì} được truyền theo dây đầu ra tới nơi nhận? Xung chạy tới cuối mỗi nhánh và ở đó biến thành tín hiệu hóa học: nhánh nhả vào khe hẹp giữa hai tế bào một nhúm chất nhất định, \emph{chất dẫn truyền thần kinh}, làm thay đổi dòng điện vào của nơi nhận. Đây chính là tiêu chí phân loại: \emph{đầu ra của nơ-ron nhả chất dẫn truyền nào}.

\section{Một đầu ra --- một loại tín hiệu}

Việc phân loại chỉ có nghĩa nếu mỗi nơ-ron chỉ có một loại đầu ra. Có đúng vậy không? Thực nghiệm cho thấy gần như đúng.

\begin{keyblock}
\begin{proposition}[Một nơ-ron --- một loại đầu ra]
\label{prop:dale}
Hầu như mỗi nơ-ron có một chất dẫn truyền chính và nhả nó ở mọi nhánh của đầu ra. Vì vậy loại của nơ-ron được xác định bởi chất dẫn truyền chính của nó~\cite{Strata1999}.
\end{proposition}
\end{keyblock}

Ta quy ước ký hiệu loại nơ-ron bằng \emph{bốn chữ cái đầu trong tên tiếng Anh của chất dẫn truyền chính}. Chẳng hạn, nơ-ron có chất dẫn truyền chính là glutamate là nơ-ron \nt{GLUT}. Tất cả các loại được tập hợp trong bảng~\ref{tab:types}.

\begin{table}[t]
\centering
\footnotesize
\setlength{\tabcolsep}{3pt}
\renewcommand{\arraystretch}{1.15}
\begin{tabular}{>{\raggedright}p{1.0cm}>{\raggedright}p{2.05cm}>{\raggedright}p{1.75cm}>{\raggedright}p{1.6cm}>{\centering}p{0.7cm}>{\raggedright}p{1.55cm}>{\raggedright}p{1.8cm}>{\raggedright\arraybackslash}p{3.2cm}}
\toprule
Loại & Chất dẫn truyền chính & Bus & Tốc độ & Dấu & Số nơ-ron trong não & Số đầu ra mỗi nơ-ron & Vai trò trong tính toán \\
\midrule
\nt{GLUT} & glutamate & địa chỉ & nhanh và chậm & $+$ & $\sim8\cdot10^{10}$ & $\sim10^4$ & trọng số dương chính \\
\nt{GABA} & GABA & địa chỉ & nhanh và chậm & $-$ & $\sim3\cdot10^{9}$ & $10^3$--$10^4$ & trọng số âm chính \\
\nt{GLYC} & glycine & địa chỉ & nhanh & $-$ & $10^6$--$10^7$ & $10^2$--$10^3$ & trọng số âm ở tủy sống và thân não; chìa khóa thứ hai cho một loại thụ thể glutamate \\
\nt{ACET} & acetylcholine & cả hai & nhanh và chậm & $\pm$ & $\sim10^6$ & tới cơ: $10$--$10^3$; trong não: $\sim10^5$ & đầu ra tới cơ; trong não là tốc độ học (giả thuyết) \\
\nt{DOPA} & dopamine & quảng bá & chậm & $\pm$ & $\sim5\cdot10^{5}$ & $10^5$--$10^6$ & sai số dự đoán phần thưởng $\delta$ \\
\nt{SERO} & serotonin & quảng bá & chậm (một loại thụ thể nhanh) & $\pm$ & $\sim3\cdot10^{5}$ & $\sim10^5$ & tầm nhìn lập kế hoạch (giả thuyết) \\
\nt{HIST} & histamine & quảng bá & chậm & $+$ & $\sim6\cdot10^{4}$ & chưa có ước tính & tín hiệu toàn cục ``hãy thức'' \\
\nt{NORA} & noradrenaline & quảng bá & chậm & $\pm$ & $\sim5\cdot10^{4}$ & $\sim10^5$ & hệ số khuếch đại (giả thuyết) \\
\nt{ADRE} & adrenaline & quảng bá & chậm & $\pm$ & $\sim10^4$ & chưa có ước tính & ít nơ-ron; vai trò trong não chưa rõ \\
\nt{ASPA} & aspartate & địa chỉ & nhanh & $+$ & chưa có ước tính & chưa có ước tính & trọng số dương yếu; vai trò còn tranh cãi \\
\bottomrule
\end{tabular}
\caption{Các loại nơ-ron, theo số lượng nơ-ron trong não giảm dần. \emph{Bus}: địa chỉ --- chất dẫn truyền được nhả vào khe hẹp tới một nơi nhận xác định; quảng bá --- từ một nhóm nhỏ nơ-ron nguồn tới những vùng rộng lớn (mục~\ref{sec:buses}). \emph{Tốc độ}: nhanh --- qua thụ thể hướng ion, mili giây; chậm --- qua thụ thể hướng chuyển hóa, hàng trăm mili giây hoặc lâu hơn. \emph{Dấu} là dấu thường gặp; dấu cuối cùng do nơi nhận quyết định (mục~\ref{sec:receiver}), và $\pm$ nghĩa là gặp cả hai dấu. \emph{Số nơ-ron trong não} --- có bao nhiêu nơ-ron loại này trong toàn bộ não người, theo bậc độ lớn; tổng số nơ-ron khoảng $8.6\cdot10^{10}$~\cite{Azevedo2009}. Gần như toàn bộ nơ-ron \nt{GLUT} --- khoảng $7\cdot10^{10}$ --- là các nơ-ron đầu vào nhỏ của tiểu não; số \nt{GLUT} và \nt{GABA} được suy ra từ phép đếm này và tỷ lệ nơ-ron \nt{GABA} trong vỏ não~\cite{Hendry1987}. Trong các loại ít ỏi, chỉ \nt{HIST}~\cite{Airaksinen1991} và nhóm chính trong hai nhóm nơ-ron \nt{DOPA}~\cite{Pakkenberg1991} được đếm trực tiếp, nên với \nt{DOPA} đây là cận dưới; với \nt{SERO} là tổng của hai phép đếm từng phần~\cite{Baker1990,Baker1991}. Các số cho \nt{NORA}, \nt{GLYC}, \nt{ACET} và \nt{ADRE} là ước tính thô theo bậc độ lớn, không có phép đếm trực tiếp. \emph{Số đầu ra mỗi nơ-ron} --- một nơ-ron có bao nhiêu đầu tận cùng khớp thần kinh, tức là những điểm nhả chất dẫn truyền trên các nhánh của dây đầu ra (hình~\ref{fig:blackbox}), theo bậc độ lớn. Với các loại quảng bá, đầu tận cùng chưa được đếm trực tiếp: ở \nt{DOPA} con số được suy ra từ tỷ lệ vùng đích được phủ bởi sự phân nhánh của một dây đầu ra~\cite{Matsuda2009}, ở các loại còn lại ước tính còn thô hơn. \nt{ACET} có hai trường hợp: nơ-ron điều khiển cơ và nơ-ron quảng bá trong não.}
\label{tab:types}
\end{table}

% the pie is a table-type float with a figure caption, so it can never float ahead of Table~\ref{tab:types}
\begin{table}[tb]
\makeatletter\def\@captype{figure}\makeatother
\centering
\definecolor{vizglut}{HTML}{2A78D6}
\definecolor{vizgaba}{HTML}{E34948}
\definecolor{vizrest}{HTML}{8A8986}
\begin{tikzpicture}[>=Stealth]
% pie: GLUT 96.4 %, GABA 3.6 % (13.0 degrees), the rest 0.006 % (a hairline at 0 degrees)
\fill[vizglut] (0,0) -- (13:1.9) arc (13:360:1.9) -- cycle;
\fill[vizgaba] (0,0) -- (0:1.9) arc (0:13:1.9) -- cycle;
\draw[white,line width=1pt] (0,0) -- (13:1.9);
\draw[vizrest,line width=1.2pt] (0,0) -- (0:1.9);
\node[white,align=center,font=\small] at (190:0.95) {\nt{GLUT}\\$96.4\,\%$};
\draw[gray!60!black] (6.5:1.75) -- (2.05,2.1);
\node[anchor=west,font=\small] at (2.05,2.1) {\nt{GABA} $3.6\,\%$};
% zoom box for the remaining types
\draw[densely dashed,vizrest] (1.9,0) -- (3.1,1.65);
\draw[densely dashed,vizrest] (1.9,0) -- (3.1,-2.2);
\draw[vizrest,rounded corners=3pt] (3.1,-2.2) rectangle (10.1,1.65);
\node[anchor=west,font=\scriptsize] at (3.2,1.4) {tất cả các loại còn lại: $\approx0.006\,\%$; thang logarit};
\foreach \code/\lg/\y/\val/\lx in {GLYC/6.477/1.0/{$10^6$--$10^7$}/4.05,ACET/6.0/0.6/{$\sim10^6$}/3.0,DOPA/5.699/0.2/{$\sim5\cdot10^5$}/2.699,SERO/5.477/-0.2/{$\sim3\cdot10^5$}/2.477,HIST/4.778/-0.6/{$\sim6\cdot10^4$}/1.778,NORA/4.699/-1.0/{$\sim5\cdot10^4$}/1.699,ADRE/4.0/-1.4/{$\sim10^4$}/1.0}{
  \node[anchor=east,font=\scriptsize] at (4.35,\y) {\nt{\code}};
  \fill[vizrest] (4.45,\y-0.13) rectangle ({4.45+\lg-3},\y+0.13);
  \node[anchor=west,font=\scriptsize] at ({4.5+\lx},\y) {\val};}
\draw[vizrest,thick] (7.45,1.0) -- (8.45,1.0);
\draw[vizrest,thick] (7.45,0.92) -- (7.45,1.08);
\draw[vizrest,thick] (8.45,0.92) -- (8.45,1.08);
\draw[gray!60!black] (4.45,-1.72) -- (8.45,-1.72);
\foreach \e in {3,4,5,6,7}{
  \draw[gray!60!black] ({4.45+\e-3},-1.72) -- ({4.45+\e-3},-1.66);
  \node[below,font=\scriptsize] at ({4.45+\e-3},-1.72) {$10^{\e}$};}
\end{tikzpicture}
\caption{Tỷ lệ các loại nơ-ron trong não theo ước tính của bảng~\ref{tab:types}. Hình tròn gần như bị \nt{GLUT} chiếm trọn; \nt{GABA} là một hình quạt hẹp, còn tất cả các loại khác gộp lại chỉ là một sợi tóc ở mép. Bên phải là sợi tóc đó được phóng to, theo thang logarit của số nơ-ron; với \nt{GLYC}, cột kéo tới giữa khoảng $10^6$--$10^7$, còn đoạn thẳng biểu thị cả khoảng. \nt{ASPA} chưa có ước tính nên không có trên hình.}
\label{fig:typepie}
\end{table}

Các đống này chênh lệch nhau đến mức nào, có thể thấy trên hình~\ref{fig:typepie}: cứ một nghìn nơ-ron thì $964$ là \nt{GLUT}, $36$ là \nt{GABA}, còn tất cả các loại khác gộp lại chỉ khoảng sáu nơ-ron trên một trăm nghìn.

Tên GABA vốn đã gồm bốn chữ cái. Những chất hầu như không bao giờ là chất dẫn truyền chính --- neuropeptide, chất khí, lipid, purine --- không có mã riêng; về chúng xem phụ lục~\ref{sec:othermediators}. Chữ ``hầu như'' trong mệnh đề~\ref{prop:dale} là quan trọng. Nhiều nơ-ron nhả kèm chất dẫn truyền chính những chất phụ trợ~\cite{Yao2023}. Một số còn nhả chất dẫn truyền nhanh thứ hai, thậm chí trái dấu: một phần nơ-ron \nt{DOPA} nhả cả GABA~\cite{Tritsch2012}. Lại có những nơ-ron mà các chất dẫn truyền khác nhau đi ra từ các nhánh khác nhau của cùng một đầu ra~\cite{Zhang2015}. Tất cả những điều này đều đã được đo, nên ``một nơ-ron --- một chất dẫn truyền'' là quy tắc có ngoại lệ chứ không phải định luật. Nhưng với tư cách phép chia đầu tiên, nó đứng vững: trong bản đồ tế bào não chuột, nơi nơ-ron được chia theo các gen hoạt động trong chúng thành hơn năm nghìn nhóm, nơ-ron vẫn chia thành các lớp lớn trước hết theo chất dẫn truyền chính~\cite{Yao2023}.

Quy tắc này có một hệ quả giúp phân biệt ngay bộ não với mạng nơ-ron nhân tạo. Trong mạng nhân tạo, cùng một nơ-ron có thể có trọng số dương tới nơi nhận này và trọng số âm tới nơi nhận khác: các trọng số $w_{ij}$ chỉ là những con số trong ma trận. Trong não, dấu tác động chủ yếu do chất dẫn truyền quyết định, mà mỗi nơ-ron chỉ có một chất dẫn truyền. Do đó, nếu nơ-ron $j$ kích thích một nơi nhận, nó kích thích mọi nơi nhận khác. Cả \emph{cột} ma trận trọng số đi ra từ nơ-ron $j$ cùng một dấu:
\begin{empheq}[box=\keybox]{equation}
\operatorname{sign} w_{ij} \;=\; s_j \quad\text{với mọi nơi nhận } i, \qquad s_j\in\{+1,-1\}.
\label{eq:dalesign}
\end{empheq}
Nơ-ron chia thành nơ-ron kích thích ($s_j=+1$) và nơ-ron ức chế ($s_j=-1$), và đây là tính chất của chính nơ-ron chứ không phải của liên kết. Nếu mạng cần một liên kết âm từ nơ-ron kích thích, phải làm qua trung gian: nơ-ron kích thích kích thích một nơ-ron ức chế, và nơ-ron này ức chế đích: một trọng số âm trễ thêm một mắt xích.

Người ta biết hơn một trăm chất dẫn truyền, nhưng gần như toàn bộ công việc của não dựa vào khoảng sáu chất, và chúng chia thành hai lớp hoàn toàn khác nhau.

\section{Hai bus: bus địa chỉ và bus quảng bá}
\label{sec:buses}

Trong mạng máy tính có hai cách gửi thông điệp. Cách thứ nhất là \emph{gửi theo địa chỉ}, unicast: gói tin đi từ một người gửi tới một người nhận xác định, nhanh, và không ai khác thấy nó. Cách thứ hai là \emph{quảng bá}, broadcast: người gửi truyền một thông điệp tới mọi nút mạng cùng lúc. Nơ-ron dùng cả hai cách, và các chất dẫn truyền chia ra đúng theo tiêu chí này (hình~\ref{fig:phone_radio}).

\begin{figure}[t]
\centering
\begin{tikzpicture}[>=Stealth,scale=0.95]
% ---- unicast: point-to-point wiring
\node[above] at (2.0,3.3) {\small \textbf{Bus địa chỉ}: glutamate và GABA};
\foreach \i in {0,...,4}{
  \fill[blue!60!black] (0.2+0.9*\i,2.5) circle (0.1);
  \fill[blue!60!black] (0.2+0.9*\i,0.6) circle (0.1);}
\foreach \a/\b in {0/0,0/1,1/1,1/3,2/2,3/2,3/4,4/4,2/0}{
  \draw[->,blue!60!black] (0.2+0.9*\a,2.38) -- (0.2+0.9*\b,0.72);}
\fill[red!70!black] (1.1,1.55) circle (0.09);
\pic[scale=1.2] at (1.75,1.9) {letter};
\draw[->,red!70!black,thick] (1.1,1.55) -- (0.28,0.7);
\draw[->,red!70!black,thick] (1.1,1.55) -- (1.95,0.72);
\node[align=center,below] at (2.0,0.2) {\small hàng tỷ nơ-ron,\\[-2pt]\small mỗi nơ-ron có nơi nhận riêng,\\[-2pt]\small thời gian: mili giây};
% ---- broadcast: one small source, global targets
\begin{scope}[xshift=7.2cm]
\node[above] at (1.8,3.3) {\small \textbf{Bus quảng bá}: dopamine, \dots};
\foreach \i in {0,...,8}{\fill[blue!60!black] (-0.2+0.5*\i,2.75) circle (0.08);}
\fill[orange!85!black] (1.8,0.35) ellipse (0.35 and 0.18);
\pic[scale=1.5] at (2.7,0.75) {mast};
\foreach \x in {-0.2,0.3,0.8,1.3,1.8,2.3,2.8,3.3,3.8}{
  \draw[->,orange!85!black] (1.8,0.5) .. controls (1.8,1.3) and (\x,1.6) .. (\x,2.62);}
\node[align=center,below] at (1.8,0.1) {\small hàng nghìn đến hàng trăm nghìn nơ-ron,\\[-2pt]\small một thông điệp cho tất cả,\\[-2pt]\small thời gian: hàng trăm mili giây trở lên};
\end{scope}
\end{tikzpicture}
\caption{Hai cách truyền tin. Bên trái là bus địa chỉ, giống bưu điện với địa chỉ ghi trên phong bì; màu đỏ đánh dấu một nơ-ron và hai nơi nhận của nó. Bên phải là bus quảng bá, giống một đài phát thanh: một nhóm nhỏ nơ-ron nguồn, và mọi nơi đều nhận cùng một thông điệp.}
\label{fig:phone_radio}
\end{figure}

\emph{Bus địa chỉ} là glutamate và GABA. Tuyệt đại đa số nơ-ron nhả chúng. Mỗi đầu ra dẫn tới những tế bào xác định, chất dẫn truyền chỉ tác động trong khe hẹp của điểm tiếp xúc, và tác động nhanh: trong một mili giây các kênh ion của nơi nhận mở ra, sau vài mili giây mọi thứ kết thúc. Đây là bus \emph{dữ liệu}: trên nó truyền đi những gì bộ não tính toán.

\emph{Bus quảng bá} là dopamine, serotonin, noradrenaline và một phần acetylcholine. Chúng được nhả ra bởi những nhóm nơ-ron rất nhỏ, nhưng đầu ra của mỗi nơ-ron phân nhánh cực kỳ rộng. Chất dẫn truyền thường không được nhả vào khe hẹp mà vào khoảng gian bào, nơi nó loang ra và tác động lên mọi tế bào ở gần. Nó tác động chậm: không mở kênh trực tiếp mà khởi động bên trong nơi nhận một chuỗi phản ứng hóa học. Trên bus này không truyền dữ liệu mà truyền các \emph{tham số điều khiển}, chung cho những vùng lớn của mạng, làm thay đổi chế độ làm việc của nó: hệ số khuếch đại, tốc độ học, tín hiệu ``điều vừa rồi là tốt''. Vì thế các chất này được gọi là \emph{chất điều biến}. Trong một thời gian dài, người ta cho rằng mỗi kênh như vậy là một con số duy nhất cho cả bộ não. Các phép đo hiện đại đã làm rõ bức tranh: kênh không phải một dây dẫn mà là một bó nhỏ các đường song song. Các nơ-ron nguồn của cùng một chất dẫn truyền chia thành những phân hệ có nơi nhận riêng và thông điệp riêng, còn lượng chất được nhả tại chỗ còn được điều chỉnh bởi các tín hiệu cục bộ~\cite{Ren2018,Poe2020}. Số đường như vậy là vài đơn vị hoặc vài chục, chứ không phải hàng tỷ, nên kênh vẫn là kênh quảng bá.

Nếu bạn chỉ cần một hình ảnh từ chương này, thì đây là nó. Bộ não là một mạng khổng lồ truyền dữ liệu theo địa chỉ, bên trên có vài kênh quảng bá mang tín hiệu điều khiển. Lập trình viên sẽ nhận ra một cấu trúc quen thuộc: phần tính toán cộng với một tập nhỏ biến điều khiển, mỗi biến dùng chung cho một vùng lớn của mạng.

\section{\nt{GLUT}: trọng số dương}

Glutamate là chất dẫn truyền kích thích chính của não. Khi đầu ra của nơ-ron nhả glutamate, ở nơi nhận mở ra các kênh cho natri đi vào, điện áp trên màng của nơi nhận tăng lên, và khả năng nó phát ra xung của chính mình tăng lên. Theo ngôn ngữ của hình~\ref{fig:blackbox}, đây là đầu vào có trọng số \emph{dương}, $w_i>0$.

Ai nhả glutamate? Gần như tất cả những nơ-ron truyền dữ liệu đi xa: từ ba phần tư đến bốn phần năm số nơ-ron vỏ não và phần lớn các nơ-ron nối các vùng não khác nhau. Mạng của não trước hết là mạng các liên kết glutamate.

Một chi tiết kỹ thuật. Sau khi được nhả, nồng độ glutamate trong khe vọt lên hàng nghìn lần, tới khoảng một milimol, rồi giảm với hằng số thời gian khoảng một mili giây~\cite{Clements1992}: glutamate khuếch tán ra khỏi khe, còn các protein bơm chuyên dụng hút nó trở lại vào tế bào. Để làm gì? Cũng như trong đường truyền, sau mỗi ký hiệu tín hiệu được đưa về không. Nếu không dọn chất dẫn truyền, xung tiếp theo sẽ tới một đường ``đang bận'', và các xung cách nhau vài mili giây sẽ dính vào nhau.

\section{\nt{GABA}: trọng số âm}

Hãy hình dung một mạng mà mọi trọng số đều dương. Nếu trung bình mỗi xung sinh ra hơn một xung tiếp theo, hoạt động tăng theo hàm mũ và sau vài bước lan khắp mạng. Người làm điện tử sẽ nhận ra một vòng hồi tiếp dương có hệ số khuếch đại lớn hơn một: mạch đi vào bão hòa. Để tránh điều đó cần có trọng số âm. Nguồn chính của chúng là axit gamma-aminobutyric, viết tắt là GABA.

GABA mở ở nơi nhận những kênh cho ion clo đi qua. Điện thế cân bằng của clo nằm gần điện thế nghỉ hoặc thấp hơn một chút, nên các kênh clo mở giữ màng ở gần mức nghỉ và không cho điện áp tăng tới ngưỡng. Đây là đầu vào có trọng số \emph{âm}, $w_i<0$. Nơ-ron GABA chiếm từ một phần năm đến một phần tư số nơ-ron vỏ não~\cite{Hendry1987}. Đầu ra của chúng ngắn, và chúng ức chế những láng giềng gần nhất.

Ức chế trong não làm nhiều hơn hẳn một phép trừ. Nó hoạt động như một vòng hồi tiếp âm giữ toàn bộ mạng ở điểm làm việc. Người ta cho rằng trong vỏ não hoạt động bình thường, dòng kích thích và dòng ức chế đi tới một nơ-ron cân bằng nhau gần như chính xác: chế độ này được lý thuyết dự đoán~\cite{vanVreeswijk1996} và thấy được trong các phép đo dòng điện ở vỏ não sống~\cite{Haider2006}, và nơ-ron luôn nằm gần ngưỡng. Một mạch được ổn định bởi hồi tiếp âm mạnh quanh một điểm không bền phản ứng nhanh và nhạy với tín hiệu nhỏ --- một thủ thuật kỹ thuật lâu đời.

\section{\nt{DOPA}: tín hiệu sai số}

Kênh quảng bá đầu tiên là dopamine. Ở người có cỡ nửa triệu nơ-ron dopamine, khoảng một trên một trăm bảy mươi nghìn; đó là số đếm được trong nhóm chính của hai nhóm~\cite{Pakkenberg1991}, nên đây là cận dưới. Đầu ra của chúng tỏa tới những vùng chọn hành động.

Kênh này truyền gì? Câu trả lời đến từ những thí nghiệm ghi xung của nơ-ron dopamine ở con vật nhận phần thưởng~\cite{Schultz1997}. Thỉnh thoảng con vật bất ngờ được cho một giọt nước ép, và nơ-ron dopamine đáp lại bằng một chùm xung ngắn. Sau đó, trước khi cho nước ép, người ta bật một bóng đèn, luôn trước nước ép một giây. Khi con vật đã học được mối liên hệ này, nơ-ron bắt đầu đáp lại bằng chùm xung khi \emph{đèn} sáng, còn với chính giọt nước ép đến đúng hạn thì hoàn toàn không đáp lại. Còn nếu sau khi đèn sáng mà không cho nước ép, thì vào lúc lẽ ra nước ép phải đến, nơ-ron \emph{im lặng}, và tần số của chúng giảm xuống dưới mức thường.

Quy tắc giải thích cả ba trường hợp (hình~\ref{fig:rpe}): nơ-ron không báo điều tốt đến mức nào, mà báo nó \emph{tốt hơn hay tệ hơn so với dự kiến} bao nhiêu.

\begin{figure}[t]
\centering
\begin{tikzpicture}[>=Stealth,x=3.4cm,y=1cm]
\foreach \y/\lab in {2.8/{nước ép bất ngờ},1.4/{nước ép sau đèn},0/{có đèn, không nước ép}}{
  \draw[gray!70!black] (0,\y) -- (3,\y);
  \draw[densely dotted,gray!60!black] (0,\y+0.3) -- (3,\y+0.3);
  \node[left,align=right,font=\small] at (-0.03,\y+0.35) {\lab};}
% row 1: juice without a cue -> burst at the juice
\draw[very thick,orange!85!black] plot[domain=0:3,samples=120] (\x,{2.8+0.3+0.75*exp(-((\x-2.08)/0.07)^2)});
\pic[scale=0.8] at (2,2.58) {juice};
% row 2: cue then juice -> burst at the cue, nothing at the juice
\draw[very thick,orange!85!black] plot[domain=0:3,samples=120] (\x,{1.4+0.3+0.75*exp(-((\x-1.08)/0.07)^2)});
\pic[scale=0.85] at (1,1.15) {lamp}; \pic[scale=0.85] at (2,1.15) {juice};
% row 3: cue, juice omitted -> burst at the cue, dip when the juice was due
\draw[very thick,orange!85!black] plot[domain=0:3,samples=120] (\x,{0.3+0.75*exp(-((\x-1.08)/0.07)^2)-0.28*exp(-((\x-2.1)/0.12)^2)});
\pic[scale=0.85] at (1,-0.25) {lamp}; \pic[scale=0.85] at (2,-0.25) {nojuice};
\node[right,font=\scriptsize] at (2.36,0.1) {sụt};
\node[right,font=\scriptsize] at (2.13,3.75) {bất ngờ!};
\node[left,font=\scriptsize] at (0.97,2.25) {vọt lên khi có đèn};
\node[above,font=\scriptsize] at (2,1.75) {không có gì};
\draw[->,gray!70!black] (0,-0.75) -- (3.05,-0.75) node[right,black,font=\small] {$t$};
\draw[gray!60!black] (1,-0.8) -- (1,-0.7) node[below=2pt,font=\scriptsize] {đèn, $1\,\text{s}$};
\draw[gray!60!black] (2,-0.8) -- (2,-0.7) node[below=2pt,font=\scriptsize] {nước ép, $2\,\text{s}$};
\end{tikzpicture}
\caption{Tần số xung của nơ-ron \nt{DOPA} trong ba thí nghiệm (sơ đồ, theo~\cite{Schultz1997}). Đường chấm là tần số nền đều đặn. Bóng đèn là tín hiệu báo trước, giọt nước là nước ép, giọt nước bị gạch chéo là nước ép đã hứa nhưng không cho. Đáp ứng chính là sai số dự đoán~\eqref{eq:rpe}.}
\label{fig:rpe}
\end{figure}

Lập trình viên quen với học tăng cường~\cite{SuttonBarto2018} sẽ nhận ra đại lượng này ngay. Một tác tử học cách chọn hành động lưu giữ ước lượng $V(s)$ --- nó kỳ vọng nhận được bao nhiêu phần thưởng khi ở trạng thái $s$. Sau mỗi bước, nó so sánh kỳ vọng với cái nhận được:
\begin{empheq}[box=\keybox]{equation}
\delta_t \;=\; r_t \;+\; \gamma\,V(s_{t+1}) \;-\; V(s_t) ,
\label{eq:rpe}
\end{empheq}
và hiệu chỉnh ước lượng: $V(s_t)\leftarrow V(s_t)+\alpha\,\delta_t$. Ở đây $r_t$ là phần thưởng nhận được, $\gamma<1$ là hệ số chiết khấu (phần thưởng tương lai rẻ hơn phần thưởng tức thời bao nhiêu), $\alpha$ là tốc độ học. Đại lượng $\delta_t$ được gọi là \emph{sai số sai phân thời gian}.

Nó hành xử đúng hệt như nơ-ron dopamine. Nước ép bất ngờ: $V$ vẫn bằng không, $r>0$, $\delta>0$. Nước ép đã học: đèn đã dự báo phần thưởng, $V(s_t)$ trước nước ép đã bằng $r$, và $\delta=0$. Nước ép không đến: $V(s_t)=r$, còn $r_t=0$, $\delta<0$. Còn chùm xung dời sang thời điểm đèn sáng, bởi chính lúc đó kỳ vọng tăng vọt: $\gamma V(s_{t+1})-V(s_t)>0$. Các bước $t,t+1,\dots$ ở đây chỉ là quy ước của thuật toán: cơ thể không có bộ đếm nhịp chia thời gian thành bước, và ở chương~\ref{sec:dofa} ta sẽ thu được cùng đại lượng đó trong thời gian liên tục.

Vậy dopamine là sự phát quảng bá một tín hiệu sai số vô hướng $\delta$ tới mọi nơi cần học theo nó. Ở phép gần đúng đầu tiên, đó là một dây dẫn cho cả bộ não và một con số tại mỗi thời điểm; dây dẫn này thực ra là một bó đến mức nào, sẽ được phân tích ở chương~\ref{sec:dofaalt}.

\section{Các kênh quảng bá còn lại}

Với các chất điều biến còn lại, mới chỉ có những giả thuyết làm việc, chuyển vai trò của chúng sang cùng ngôn ngữ các tham số toàn cục~\cite{Doya2002}. Hãy coi chúng đúng là giả thuyết.

\emph{\nt{NORA}, noradrenaline: hệ số khuếch đại.} Nơ-ron noradrenaline còn ít hơn nơ-ron dopamine, theo ước tính thô khoảng vài chục nghìn, còn đầu ra của chúng tỏa ra gần như khắp bộ não. Chúng hoạt hóa mạnh khi xảy ra điều gì bất ngờ hoặc quan trọng. Noradrenaline thay đổi độ dốc của đặc tuyến truyền đạt của các nơ-ron nhận: đầu vào yếu trở nên yếu hơn, đầu vào mạnh trở nên mạnh hơn. Điều này được đo như sau: xung nền của nơi nhận bị ức chế mạnh hơn đáp ứng với đầu vào, và tỷ số tín hiệu trên nhiễu tăng lên~\cite{Foote1975NE}. Trong mô hình đơn giản, điều này được mô tả bằng một con số: nếu nơ-ron đáp ứng dòng vào $u$ bằng tần số $f=F\bigl(g\cdot u\bigr)$, thì noradrenaline làm tăng hệ số $g$~\cite{AstonJones2005} (chi tiết ở chương~\ref{sec:nora}). Mạng có $g$ lớn làm việc ``tương phản hơn'' và ra quyết định nhanh hơn; với $g$ nhỏ --- ``mờ'' hơn và thử nhiều phương án hơn, giống một tác tử học tăng cường ở chế độ khám phá.

\emph{\nt{ACET}, acetylcholine: tốc độ học.} Acetylcholine điều khiển việc mạng lắng nghe đầu vào hiện tại nhiều đến đâu và lắng nghe dữ liệu đã lưu của chính nó nhiều đến đâu. Khi mức acetylcholine cao, các đầu vào từ giác quan được tăng cường, còn các liên kết bên trong, ``hồi tưởng'', bị làm yếu đi, đồng thời việc thay đổi trọng số trở nên dễ hơn~\cite{Hasselmo2006}. Nói thô, đây là công tắc ``ghi/đọc'' và kèm theo nó là tốc độ học $\alpha$.

\emph{\nt{SERO}, serotonin: tầm nhìn lập kế hoạch.} Serotonin truyền gì là điều được hiểu kém nhất. Một giả thuyết gắn nó với hệ số chiết khấu $\gamma$ trong~\eqref{eq:rpe}: mức serotonin cao tương ứng với sự sẵn lòng chờ phần thưởng lớn thay vì chộp lấy phần thưởng nhỏ ngay. Nơ-ron serotonin làm việc đều và chậm, vài xung mỗi giây khi thức, và gần như im lặng trong một pha của giấc ngủ~\cite{JacobsAzmitia1992}.

Cả sáu chất dẫn truyền được tập hợp trong bảng~\ref{tab:transmitters}.

\begin{table}[t]
\centering
\small
\begin{tabular}{>{\raggedright}p{2.4cm}>{\raggedright}p{3.0cm}>{\raggedright}p{2.0cm}>{\raggedright\arraybackslash}p{5.2cm}}
\toprule
Loại nơ-ron & Tỷ lệ nơ-ron & Bus & Vai trò trong tính toán \\
\midrule
\nt{GLUT} (glutamate) & đa số & địa chỉ & trọng số dương, $w>0$ \\
\nt{GABA} (GABA) & 20--25\% trong vỏ não & địa chỉ & trọng số âm, $w<0$; ổn định hóa \\
\nt{DOPA} (dopamine) & $\sim 6\cdot10^{-6}$ & quảng bá & sai số dự đoán $\delta$ \\
\nt{NORA} (noradrenaline) & $\sim 6\cdot10^{-7}$ & quảng bá & hệ số khuếch đại $g$ (giả thuyết) \\
\nt{ACET} (acetylcholine) & nhỏ & cả hai & tốc độ học $\alpha$; ``ghi/đọc'' (giả thuyết) \\
\nt{SERO} (serotonin) & $\sim 3\cdot10^{-6}$ & quảng bá & chiết khấu $\gamma$ (giả thuyết) \\
\bottomrule
\end{tabular}
\caption{Các chất dẫn truyền chính của não theo ngôn ngữ tính toán. Cột bên phải là một sự đơn giản hóa mạnh: đáng tin với glutamate, GABA và dopamine, còn với các chất khác đó là giả thuyết.}
\label{tab:transmitters}
\end{table}

\section{Dấu do nơi nhận quyết định}
\label{sec:receiver}

Tôi đã lừa bạn một chút khi nói: glutamate là trọng số dương, GABA là trọng số âm. Không phân tử nào tự mang dấu. Phân tử chỉ là một chiếc chìa khóa. Điều gì xảy ra khi nó bị bắt giữ là do \emph{ổ khóa} quyết định --- thụ thể trên tế bào nhận.

Ví dụ tốt nhất là dopamine. Nó có thụ thể thuộc hai họ~\cite{Beaulieu2011}. Ở họ này nó làm tế bào dễ bị kích thích hơn, ở họ kia --- khó hơn. Trong vùng chọn hành động, một số nơ-ron mang thụ thể loại thứ nhất, số khác mang loại thứ hai, và cùng một tín hiệu dopamine đồng thời tăng cường nhóm này và làm yếu nhóm kia. Giống như một gói tin quảng bá được các nút mạng khác nhau diễn giải theo những cách khác nhau.

\begin{keyblock}
\begin{proposition}[Ý nghĩa của thông điệp do nơi nhận quyết định]
\label{prop:receptor}
Dấu và tốc độ tác động của một chất dẫn truyền không do bản thân chất đó quyết định, mà do thụ thể mà nó gắn vào. Thụ thể có hai loại. Thụ thể \emph{hướng ion} tự chúng là kênh ion và mở ra trong một phần nhỏ của mili giây: đó là điều khiển dòng điện trực tiếp, như cực cổng của transistor. Thụ thể \emph{hướng chuyển hóa} khởi động bên trong tế bào một chuỗi phản ứng hóa học và tác động trong hàng trăm mili giây hoặc lâu hơn: đó giống như ghi vào một thanh ghi cấu hình, làm thay đổi hoạt động tiếp theo của tế bào. Bus địa chỉ chủ yếu nói qua loại thứ nhất, bus quảng bá chủ yếu qua loại thứ hai.
\end{proposition}
\end{keyblock}

Vậy tại sao vẫn có thể nói ``glutamate kích thích''? Bởi vì gần như mọi thụ thể glutamate gặp trong thực tế đều là thụ thể kích thích, còn gần như mọi thụ thể GABA đều là thụ thể ức chế. Đây không phải định luật tự nhiên mà là một quy ước rất đáng tin cậy, và quy tắc~\eqref{eq:dalesign} dựa trên chính quy ước đó.

\section{Những điều cần nhớ}

Tuyệt đại đa số nơ-ron là bus dữ liệu địa chỉ với các trọng số cùng dấu trên mỗi nơ-ron; một thiểu số nhỏ bé là các kênh quảng bá, truyền tới những vùng não lớn vài tín hiệu điều khiển chung. Khoảng hai nơ-ron trên mỗi một trăm nghìn --- và từ chúng phụ thuộc việc toàn bộ phần còn lại của mạng học ra sao và làm việc ở chế độ nào. Với kỹ sư, điều này không có gì lạ: trong bất kỳ máy tính nào, số thanh ghi điều khiển ít hơn hẳn số ô nhớ, vậy mà thiếu chúng máy không chạy được.

Ở chương sau, ta sẽ kiểm tra xem một mạng chỉ gồm nơ-ron \nt{GLUT}, loại đông đảo nhất, làm việc trong thời gian liên tục, có thể tính được gì.

\section{Bài tập}

\begin{exercise}[\lvlA]
\label{ex:01:1}
\emph{Các đống và dây dẫn.} Theo bảng~\ref{tab:types} (giữa khoảng theo thang logarit; \nt{ACET} có $10^5$ đầu ra): (a)~nếu mỗi nơ-ron là một người, nơ-ron \nt{GLUT} bằng bao nhiêu lần dân số Trái Đất, và tất cả nơ-ron quảng bá bằng một thành phố cỡ nào? (b)~Tỷ lệ đầu tận cùng của \nt{DOPA}, \nt{SERO}, \nt{NORA}, \nt{ACET} lớn hơn tỷ lệ nơ-ron của chúng bao nhiêu lần?
\end{exercise}

\begin{exercise}[\lvlA]
\label{ex:01:2}
\emph{Đưa về không.} Glutamate trong khe giảm theo $c_0e^{-t/\tau_c}$, $c_0=1\mM$, $\tau_c=1\ms$; nơi nhận nhận biết nồng độ trên $10\,\mu\text{M}$. (a)~Đường truyền ``bận'' bao lâu sau một lần nhả? (b)~Các bơm bị chặn một phần, $\tau_c=10\ms$. Đến tần số nhả nào đường truyền vẫn kịp giải phóng?
\end{exercise}

\begin{exercise}[\lvlB]
\label{ex:01:3}
\emph{Chùm xung dời đi đâu.} Trong thí nghiệm hình~\ref{fig:rpe}, sau đèn là các trạng thái $s_0\to\dots\to s_{K-1}$ với bước $\Delta$, $T=K\Delta$; phần thưởng $r$ đến ở lần chuyển cuối, thời điểm đèn sáng không dự đoán được. Tìm các $V(s_k)$ đã học và đáp ứng $\delta$ với đèn. Đặt $\gamma=e^{-\Delta/\tau_\gamma}$, tìm $\tau_\gamma$ khi $T=1$~s, $\Delta=0.1$~s, $\gamma=0.95$.
\end{exercise}

\begin{exercise}[\lvlB]
\label{ex:01:4}
\emph{Mô phỏng: học theo sai số.} Mô phỏng chuỗi của bài tập~\ref{ex:01:3} với $K=10$, $\gamma=0.95$, $r=1$ theo quy tắc $V(s_t)\leftarrow V(s_t)+\alpha\delta_t$. Ở lượt thử $n^*$ nào đáp ứng với đèn lần đầu vượt đáp ứng với nước ép khi $\alpha=0.05$, $0.1$, $0.2$, $0.5$? $n^*$ phụ thuộc vào $\alpha$ thế nào?
\end{exercise}

\begin{exercise}[\lvlB]
\label{ex:01:5}
\emph{Thiết kế: phát một con số.} Số $\delta$ cần được đưa tới cả $10^{10}$ nơ-ron; mỗi nơi nhận, kể cả các bộ chuyển tiếp, phải nhận $10$ đầu tận cùng từ lớp trước, mỗi mắt xích thêm $1\ms$. Cây chuyển tiếp tiết kiệm nhất có bao nhiêu nơ-ron và mắt xích khi mỗi nơ-ron có $10^4$ đầu ra (\nt{GLUT}) và $10^{5.5}$ (\nt{DOPA})?
\end{exercise}

\begin{exercise}[\lvlB]
\label{ex:01:6}
\emph{Vì sao trạng thái cân bằng nhạy.} Mạng có $80\,\%$ \nt{GLUT} và $20\,\%$ \nt{GABA}, mỗi nơ-ron nối với mọi nơ-ron bằng trọng số $J_E/N$ và $-J_I/N$; $\tau\dot r=-r+Wr+h$. Với $J_E=10$, trị riêng khác không duy nhất của $W$ bằng $0.5$. Tìm tỷ số dòng ức chế trên dòng kích thích, và cần làm yếu ức chế bao nhiêu phần trăm để mạng rơi vào thác hoạt động.
\end{exercise}

\begin{exercise}[\lvlC]
\label{ex:01:7}
\emph{Nghiên cứu: \nt{SERO} có phải là $\gamma$?} Theo bài tập~\ref{ex:01:3}, đáp ứng của nơ-ron \nt{DOPA} với đèn bằng $re^{-T/\tau_\gamma}$. Độ trễ $T=1,\dots,5$~s, mỗi mức $n$ lần trình bày, $\ln\delta$ đo với độ phân tán $0.5$. Cần $n$ bao nhiêu để phân biệt $\tau_\gamma=2$~s với $2.4$~s (mức $0.05$, lực kiểm định $0.8$)? Cơ chế nào ngoài $\gamma$ cho cùng sự suy giảm theo $T$?
\end{exercise}

\chapter{Nơ-ron \nt{GLUT} tính được gì}
\label{sec:glutcomputer}
\begin{chaptergoals}
\item cách nơ-ron tích phân rò tạo ra OR và bộ phát hiện trùng khớp, tức AND theo thời gian;
\item vì sao riêng nơ-ron \nt{GLUT} chỉ tính được hàm đơn điệu, nên hoạt động không thể dập hoạt động;
\item cách cặp dây bật---tắt giúp nơ-ron \nt{GLUT} tính mọi hàm logic và báo kết quả đã sẵn sàng;
\item cách đường trễ biến chênh lệch thời gian thành vị trí, và hồi tiếp tạo ra bộ nhớ và bộ định thời.
\end{chaptergoals}

\section{Luật chơi}

Ở chương~\ref{sec:neurontypes} ta đã thấy tuyệt đại đa số nơ-ron của não là \nt{GLUT}: chúng chỉ kích thích, trọng số của chúng chỉ dương. Hãy lấy bao nhiêu nơ-ron như thế cũng được và hỏi: một mạng như vậy tính được gì, nếu chỉ cho phép những gì có trong cơ thể sống?

Mà trong cơ thể sống không có bộ tạo xung nhịp chuyển mạch mọi phần tử cùng lúc. Mỗi nơ-ron làm việc độc lập, trong thời gian liên tục: xung đến và đi lúc nào tùy ý, với những độ trễ khác nhau. Có \emph{đầu vào} --- các nơ-ron cảm giác, phát xung khi gặp ánh sáng, âm thanh hay áp lực --- và \emph{đầu ra} --- các nơ-ron điều khiển cơ. Ở giữa là mạng, và không ai bảo nó khi nào bắt đầu và khi nào kết thúc phép tính. Với người làm điện tử, đây là một mạch \emph{bất đồng bộ} gồm các phần tử có độ trễ không biết trước chính xác.

Ta lấy mô hình nơ-ron đơn giản nhất trong số những mô hình thực sự làm việc trong thời gian liên tục. Điện áp $V$ trên màng nơ-ron lúc nghỉ bằng không. Mỗi xung đến từ nơ-ron $j$ làm nó nhảy lên một bậc $w_j\ge0$, sau đó điện áp tự rò về không với hằng số thời gian $\tau$:
\begin{empheq}[box=\keybox]{equation}
\tau\,\frac{\dd V}{\dd t} \;=\; -V \;+\; \tau\sum_j w_j \sum_k \delta\bigl(t-t_j^{(k)}-d_j\bigr), \qquad w_j\ge 0 .
\label{eq:glutneuron}
\end{empheq}
Ở đây $t_j^{(k)}$ là các thời điểm phát xung của nơ-ron $j$, $d_j$ là độ trễ trên đường từ nó tới, $\delta$ là hàm delta, biểu thị bước nhảy tức thời. Khi $V$ đạt ngưỡng $\theta$, nơ-ron phát một xung, $V$ được đặt lại về không, và trong khoảng một mili giây nơ-ron không thể phát xung lần nữa. Các con số điển hình: $\tau$ từ một phần nhỏ của mili giây đến hai mươi mili giây tùy nơ-ron, độ trễ từ một phần nhỏ của mili giây đến hàng chục mili giây. Đây là \emph{nơ-ron tích phân rò}: một mạch RC có ngưỡng. Đây là sự đơn giản hóa có chủ ý: ở nơ-ron vỏ não, chính các nhánh đầu vào cũng phát xung cục bộ~\cite{Smith2013}, và một nơ-ron hành xử giống một mạng nhiều lớp nhỏ hơn (chương~\ref{sec:synthesis}). Với các câu hỏi của chương này, một mạch RC là đủ.

Khác biệt chính so với cổng logic là sự rò: nơ-ron không nhớ xung mãi mãi mà chỉ trong khoảng $\tau$, và chỉ những gì đến trong cửa sổ đó mới được cộng lại.

\section{OR và AND}

Bắt đầu từ các cổng. Nếu một xung nâng điện áp vượt ngưỡng, $w\ge\theta$, nơ-ron phát xung với mỗi xung của bất kỳ đầu vào nào. Đó là \emph{OR}.

AND thú vị hơn. Cho $w<\theta<2w$: một xung là không đủ, cần hai. Nhưng giờ câu hỏi ``cả hai đã đến chưa'' vô nghĩa chừng nào chưa nói \emph{trong khoảng thời gian nào}. Giả sử xung thứ nhất đến lúc $0$, còn xung thứ hai sau thời gian $\Delta$. Khi xung thứ hai đến, từ xung thứ nhất còn lại $we^{-\Delta/\tau}$, và nơ-ron phát xung nếu $we^{-\Delta/\tau}+w\ge\theta$. Từ đó có cửa sổ trùng khớp:
\begin{empheq}[box=\keybox]{equation}
\Delta \;\le\; \Delta_{\max} \;=\; \tau\,\ln\frac{w}{\theta-w} .
\label{eq:window}
\end{empheq}
Đây không phải AND logic mà là \emph{bộ phát hiện trùng khớp}: AND theo thời gian (hình~\ref{fig:coincidence}). Khi $\theta=1.5w$, cửa sổ bằng $\tau\ln2\approx0.7\tau$. Với nơ-ron vỏ não có $\tau=10\ms$, đó là khoảng bảy mili giây. Ở các nơ-ron hệ thính giác, nơi $\tau$ nhỏ hơn một mili giây, cửa sổ co lại chỉ còn một phần nhỏ của mili giây.

\begin{figure}[t]
\centering
\begin{tikzpicture}[>=Stealth,x=0.9cm,y=1.6cm]
\foreach \dx/\lab/\c [count=\i] in {0.5/{trùng khớp}/red!70!black, 2.6/{không trùng khớp}/blue!60!black}{
 \begin{scope}[xshift={(\i-1)*7.2cm}]
  \draw[->,gray!70!black] (0,0) -- (6.3,0) node[below left,black] {\small $t$};
  \draw[->,gray!70!black] (0,0) -- (0,1.75) node[left,black] {\small $V$};
  \draw[densely dashed,gray!60!black] (0,1.5) -- (6.0,1.5) node[right,black] {\small $\theta$};
  \draw[very thick,\c] (0,0) -- (1,0) -- (1,1)
      plot[domain=1:1+\dx,samples=30] (\x,{exp(-(\x-1)/1.5)});
  \pgfmathsetmacro{\v}{exp(-\dx/1.5)+1}
  \draw[very thick,\c] (1+\dx,{exp(-\dx/1.5)}) -- (1+\dx,{min(\v,1.5)});
  \ifdim\v pt<1.5pt
    \draw[very thick,\c] plot[domain=1+\dx:6,samples=40] (\x,{\v*exp(-(\x-1-\dx)/1.5)});
  \else
    \draw[very thick,\c] (1+\dx,1.5) -- (1+\dx,0) -- (6,0);
    \draw[->,thick,\c] (1+\dx,1.55) -- (1+\dx,1.72) node[above,font=\scriptsize] {xung};
  \fi
  \draw[->,thick] (1,-0.35) -- (1,-0.08); \draw[->,thick] (1+\dx,-0.35) -- (1+\dx,-0.08);
  \draw[decorate,decoration={brace,amplitude=3pt,mirror}] (1,-0.42) -- (1+\dx,-0.42) node[midway,below=3pt] {\scriptsize $\Delta$};
  \node[\c] at (3.5,-0.95) {\small \lab};
 \end{scope}}
\end{tikzpicture}
\caption{Bộ phát hiện trùng khớp, ngưỡng $\theta=1.5w$. Bên trái, xung thứ hai đến sau $\Delta<\Delta_{\max}$, tổng đạt ngưỡng, nơ-ron phát xung. Bên phải, $\Delta>\Delta_{\max}$: xung thứ nhất gần như không còn lại gì, và nơ-ron im lặng.}
\label{fig:coincidence}
\end{figure}

Một cách khác để có AND là dựa vào thời lượng chứ không phải thời điểm chính xác. Chừng nào đèn còn sáng, nơ-ron cảm giác phát xung dày và đều; nếu một đầu vào như vậy chưa đủ vượt ngưỡng mà hai thì đủ, nơ-ron hoạt động chừng nào cả hai còn hoạt động. Như vậy mạng tính theo \emph{mức}, và thời điểm đến chính xác của xung không quan trọng. Phần lớn những gì dưới đây sẽ theo đúng cách này.

\section{Điều không làm được}

Giờ thử làm NOT: một nơ-ron hoạt động khi đầu vào im lặng. Không được: không có xung vào thì trong~\eqref{eq:glutneuron} không có bước nhảy nào, và điện áp đứng ở không, dưới ngưỡng. Còn mạng nhiều nơ-ron thì sao? Cũng không, và điều này được chứng minh một lần cho mọi mạng.

Tiện nhất là nói về tần số. Gọi $r_i$ là tần số xung của nơ-ron $i$, $I_i$ là dòng vào từ các đầu vào, còn $f_i$ là đặc tuyến truyền đạt của nó: tần số phụ thuộc vào tổng dòng vào như thế nào. Ở nơ-ron tích phân, đặc tuyến này không giảm: dòng vào càng lớn, xung càng dày. Tần số trong mạng đã đạt cân bằng thỏa mãn các phương trình
\begin{equation}
r_i \;=\; f_i\Bigl(\sum_j w_{ij}\,r_j + I_i\Bigr), \qquad w_{ij}\ge 0 .
\label{eq:ratenet}
\end{equation}

\begin{keyblock}
\begin{proposition}[Tính đơn điệu]
\label{prop:monotone}
Giả sử mạng~\eqref{eq:ratenet} chỉ gồm các nơ-ron \nt{GLUT} và bắt đầu từ im lặng hoàn toàn. Nếu tăng cường các đầu vào, tần số xác lập của không nơ-ron nào giảm đi. Mọi thứ một mạng như vậy tính được đều là hàm \emph{đơn điệu} của các đầu vào.
\end{proposition}
\end{keyblock}

Chứng minh. Bắt đầu từ im lặng, $r^{(0)}=0$, và lần lượt thế tần số vào vế phải của~\eqref{eq:ratenet}: $r^{(k+1)}=F\bigl(r^{(k)}\bigr)$. Vế phải $F$ không giảm theo bất kỳ đối số nào, vì trọng số không âm và các $f_i$ không giảm. Vì $r^{(1)}\ge0=r^{(0)}$, theo quy nạp $r^{(k+1)}\ge r^{(k)}$: tần số chỉ tăng và hội tụ về nghiệm nhỏ nhất của~\eqref{eq:ratenet}. Nếu thay các đầu vào $I$ bằng những đầu vào lớn hơn $I'$, thì ở mỗi bước $r^{(k)}(I)\le r^{(k)}(I')$, và ở giới hạn cũng vậy. Mạng thật đạt cân bằng không theo từng bước mà liên tục, nhưng điều tương tự vẫn đúng: với liên kết dương, bất đẳng thức giữa hai lần chạy, nếu đúng lúc ban đầu, sẽ được giữ suốt thời gian.

Với từng xung riêng lẻ, mệnh đề không đúng theo nghĩa đen: một xung vào thừa có thể khiến nơ-ron phát xung sớm hơn một chút, và vì một mili giây trơ, xung tiếp theo của nó sẽ mất. Nhưng hiệu ứng này yếu và không tin cậy, không thể dựa vào nó để tính toán. Với tần số, tính đơn điệu là chặt chẽ.

Các hệ quả giống như ở mọi mạch không có bộ đảo. Không có NOT. Không có XOR: thêm đầu vào thứ hai không thể dập tắt đầu ra. Và khó chịu nhất: \emph{không thể dừng hoạt động bằng hoạt động}, chỉ có thể lấy đi cái đang duy trì nó.

\section{Hai dây dẫn: bật và tắt}

Lối thoát do chính cơ thể gợi ý. Ở nhiều giác quan, kích thích được truyền bởi hai tập nơ-ron: trong thị giác, một số tế bào đáp ứng khi ánh sáng mạnh lên, số khác khi ánh sáng yếu đi~\cite{Schiller1992}. Ta gọi chúng là nơ-ron \emph{bật} và nơ-ron \emph{tắt}. Thay vì một dây $x$, mạng nhận một cặp $(x,\bar x)$, và sự vắng mặt mà mạng không tính được thì chính cảm biến báo cho nó.

Với cặp dây dẫn, NOT là miễn phí: chỉ cần đổi chỗ hai dây. AND và OR mỗi phép cần hai nơ-ron:
\begin{empheq}[box=\keybox]{equation}
\begin{aligned}
x\wedge y \;&\to\; \bigl(\;x\wedge y,\;\; \bar x\vee\bar y\;\bigr),\\
x\vee y \;&\to\; \bigl(\;x\vee y,\;\; \bar x\wedge\bar y\;\bigr).
\end{aligned}
\label{eq:dualrail}
\end{empheq}
Ở vế phải mọi phép toán đều đơn điệu, nên mệnh đề~\ref{prop:monotone} không bị vi phạm.

Với mạch bất đồng bộ, mã hai dây còn có ưu điểm thứ hai, thậm chí quan trọng hơn. Một cặp dây có ba trạng thái: ``có'', ``không'' và, khi cả hai im lặng, \emph{``chưa có câu trả lời''}. Nơi nhận biết phép tính đã xong không phải nhờ đồng hồ mà nhờ chính tín hiệu: ở mỗi cặp có đúng một dây đã hoạt động. Giữa các lần kích thích, cảm biến im lặng, và mạng trở về trạng thái im lặng, sẵn sàng cho lần sau. Trong điện tử bất đồng bộ, thủ thuật này được dùng để xây dựng những mạch làm việc đúng với mọi độ trễ của phần tử~\cite{Sparso2001}.

\begin{keyblock}
\begin{proposition}[Tính toán bất đồng bộ]
\label{prop:dualrail}
Nếu mỗi đầu vào đến dưới dạng một cặp dây ``bật---tắt'', thì mạng nơ-ron \nt{GLUT} có thể tính bất kỳ hàm logic nào của các đầu vào. Kết quả cũng đến dưới dạng các cặp dây, và việc nó đã sẵn sàng thấy được qua chỗ ở mỗi cặp có một dây đã hoạt động. Không cần đồng hồ cho việc này. Cái giá là số nơ-ron tăng khoảng gấp đôi.
\end{proposition}
\end{keyblock}

Ví dụ là XOR: ``đúng một trong hai kích thích đã thay đổi''. Dây thuận của kết quả là $(x\wedge\bar y)\vee(\bar x\wedge y)$, dây nghịch là $(x\wedge y)\vee(\bar x\wedge\bar y)$. Đó là sáu nơ-ron, và không nơ-ron nào cần đến ức chế (hình~\ref{fig:dualxor}).

\begin{figure}[t]
\centering
\begin{tikzpicture}[>=Stealth,x=1cm,y=1cm,
  nrn/.style={circle,draw,very thick,fill=blue!6,minimum size=0.75cm,inner sep=0pt,font=\small},
  inp/.style={font=\small}]
\node[inp] (X)  at (0,3.3) {$x$};
\node[inp] (Xb) at (0,2.2) {$\bar x$};
\node[inp] (Y)  at (0,1.1) {$y$};
\node[inp] (Yb) at (0,0)   {$\bar y$};
\node[nrn] (A1) at (3.2,3.3) {$2$};
\node[nrn] (A2) at (3.2,2.2) {$2$};
\node[nrn] (A3) at (3.2,1.1) {$2$};
\node[nrn] (A4) at (3.2,0)   {$2$};
\node[nrn] (O)  at (7.2,2.75) {$1$};
\node[nrn] (Ob) at (7.2,0.55) {$1$};
\foreach \s/\t in {X/A1,Yb/A1,Xb/A2,Y/A2,X/A3,Y/A3,Xb/A4,Yb/A4}{\draw[->,thick,blue!60!black] (\s) -- (\t);}
\draw[->,thick,blue!60!black] (A1) -- node[pos=0.45,above,sloped,font=\scriptsize,black] {$x\wedge\bar y$} (O);
\draw[->,thick,blue!60!black] (A2) -- node[pos=0.45,below,sloped,font=\scriptsize,black] {$\bar x\wedge y$} (O);
\draw[->,thick,blue!60!black] (A3) -- node[pos=0.45,above,sloped,font=\scriptsize,black] {$x\wedge y$} (Ob);
\draw[->,thick,blue!60!black] (A4) -- node[pos=0.45,below,sloped,font=\scriptsize,black] {$\bar x\wedge\bar y$} (Ob);
\draw[->,very thick,red!70!black] (O) -- (8.6,2.75) node[right,black] {$x\oplus y$: ``có''};
\draw[->,very thick,red!70!black] (Ob) -- (8.6,0.55) node[right,black] {$\overline{x\oplus y}$: ``không''};
\node[below,font=\small,gray!40!black] at (3.2,-0.55) {AND: ngưỡng $2$};
\node[below,font=\small,gray!40!black] at (7.2,-0.55) {OR: ngưỡng $1$};
\node[below,font=\small,gray!40!black] at (0,-0.55) {cặp dây dẫn};
\end{tikzpicture}
\caption{XOR trên mã hai dây, chỉ gồm các nơ-ron \nt{GLUT}. Số trong vòng tròn là ngưỡng, tính theo đơn vị trọng số của một đầu vào. Bốn nơ-ron AND nhận ra bốn tổ hợp đầu vào, hai nơ-ron OR gom chúng thành cặp dây kết quả.}
\label{fig:dualxor}
\end{figure}

\section{Độ trễ thay cho đồng hồ}

Để tính toán với thời gian, chỉ cần độ trễ trên dây dẫn và các bộ phát hiện trùng khớp. Ví dụ hay nhất là cách não xác định âm thanh đến từ đâu.

Âm thanh từ một nguồn ở bên cạnh đến tai gần sớm hơn tai xa. Hiệu thời gian phụ thuộc vào hướng: từ không, khi nguồn ở ngay phía trước, đến $0.22\,\text{m}/343\,\text{m/s}\approx0.6\ms$ ở người, khi nguồn ở đúng bên cạnh. Não phân biệt được hiệu cỡ mười \emph{micro giây}~\cite{KlumppEady1956,Grothe2010}, dù bản thân nơ-ron phát xung với độ chính xác không tốt hơn một phần nhỏ của mili giây. Bằng cách nào?

Kéo một dây từ tai trái dọc theo một hàng bộ phát hiện trùng khớp từ trái sang phải, và một dây từ tai phải từ phải sang trái (hình~\ref{fig:itd}). Tín hiệu chạy trên dây với vận tốc $v$. Tới bộ phát hiện đặt cách mép trái của hàng dài $L$ một khoảng $x$, tín hiệu trái đi mất thời gian $x/v$, tín hiệu phải --- $(L-x)/v$. Nếu âm thanh đến tai phải muộn hơn $\delta t$, hai tín hiệu sẽ gặp nhau đồng thời tại điểm mà
\begin{empheq}[box=\keybox]{equation}
\frac{x}{v} \;=\; \frac{L-x}{v} + \delta t
\quad\Longrightarrow\quad
x \;=\; \frac{L}{2} + \frac{v\,\delta t}{2} .
\label{eq:itd}
\end{empheq}
Chỉ bộ phát hiện nào nhận được cả hai tín hiệu trong phạm vi cửa sổ~\eqref{eq:window} mới phát xung. Số thứ tự của bộ phát hiện đã phát xung chính là hướng tới nguồn. Hiệu thời gian đã biến thành \emph{vị trí}.

\begin{figure}[t]
\centering
\begin{tikzpicture}[>=Stealth,
  nrn/.style={circle,draw,very thick,fill=blue!6,minimum size=0.7cm,inner sep=0pt}]
\foreach \k in {0,...,6}{\node[nrn] (D\k) at (1.4*\k,0) {\scriptsize $2$};}
\node[nrn,fill=red!15] at (D4) {\scriptsize $2$};
\draw[->,thick,blue!60!black] (-1.4,0.9) node[left,black] {\small tai trái} -- (9.2,0.9);
\draw[->,thick,orange!85!black] (9.8,-0.9) node[right,black] {\small tai phải} -- (-0.8,-0.9);
\foreach \k in {0,...,6}{
  \draw[->,blue!60!black] (1.4*\k,0.9) -- (D\k.north);
  \draw[->,orange!85!black] (1.4*\k,-0.9) -- (D\k.south);}
\draw[->,thick,red!70!black] (D4.north east) ++(0.1,0.05) -- ++(0.5,0.45) node[above right,font=\small,black] {phát xung};
\node[font=\small,gray!40!black] at (4.2,-1.5) {độ trễ tăng $\longrightarrow$ \hspace{2.2cm} $\longleftarrow$ độ trễ tăng};
\pic[scale=1.4] at (-2.3,1.75) {speaker};
\node[font=\scriptsize,align=center] at (-2.3,2.35) {nguồn bên trái};
\end{tikzpicture}
\caption{Xác định hướng âm thanh. Tín hiệu từ hai tai chạy ngược chiều nhau; mỗi vòng tròn là một bộ phát hiện trùng khớp có ngưỡng $2$. Nếu âm thanh đến tai phải muộn hơn, các tín hiệu gặp nhau ở bên phải điểm giữa, theo công thức~\eqref{eq:itd}. Đầu ra của mạng là số thứ tự của nơ-ron đã phát xung. Ở đây nguồn ở bên trái: âm thanh đến tai trái sớm hơn.}
\label{fig:itd}
\end{figure}

Ở đây không có đồng hồ, không có ức chế, thậm chí không có mã hai dây: mạng không trả lời ``có'' hay ``không'' mà chỉ ra một nơ-ron trong số nhiều nơ-ron. Một mạch như vậy gồm đường trễ và bộ phát hiện trùng khớp dường như thực sự hoạt động ở thân não thính giác của chim~\cite{Carr1990}. Độ chính xác micro giây không đến từ độ chính xác của từng nơ-ron, mà từ chỗ có nhiều bộ phát hiện và mỗi bộ theo dõi một độ trễ riêng. Ở động vật có vú, người ta không tìm thấy hàng độ trễ như vậy: ở đó cùng bài toán này dường như được giải bằng các bộ phát hiện trùng khớp được tinh chỉnh bởi sự ức chế đến đúng thời điểm từ nơ-ron \nt{GLYC}~\cite{Brand2002,Grothe2010}. Ức chế thu hẹp cửa sổ trùng khớp như thế nào, ta sẽ phân tích ở chương~\ref{sec:gaba} với ví dụ \nt{GABA}; glycine ức chế theo cùng cách, bằng cách mở kênh clo ở nơi nhận.

\section{Bộ nhớ}

Máy tính cần bộ nhớ. Trong một mạng như thế, bộ nhớ là hoạt động tự duy trì. Lấy một nhóm nơ-ron \nt{GLUT} liên kết chặt với nhau và mô tả nó bằng một tần số $r$ (hình~\ref{fig:recurrentgroup}a). Ở trạng thái cân bằng $r=f(wr+I)$, trong đó $w$ là độ mạnh hồi tiếp của nhóm lên chính nó, còn $I$ là dòng vào từ bên ngoài. Giải phương trình này bằng đồ thị, lấy giao điểm của đường cong $f(wr+I)$ với đường thẳng $r$ (hình~\ref{fig:bistable}).

\begin{figure}[t]
\centering
% (b): tau dr/dt = -r + f(w r + I), f as in fig:bistable, w=1, I=0.6; Euler dt=0.001 tau.
\begin{tikzpicture}[>=Stealth,
  nrn/.style={circle,draw,thick,fill=blue!6,minimum size=0.5cm,inner sep=0pt}]
\begin{scope}
\node[font=\small] at (-0.5,1.55) {(a)};
\foreach \k in {1,...,6}{\node[nrn] (n\k) at ({1.4+1.0*cos(60*\k)},{1.0*sin(60*\k)}) {};}
\foreach \a/\b in {1/2,2/3,3/4,4/5,5/6,6/1,1/4,2/5,3/6}{\draw[<->,blue!60!black] (n\a) -- (n\b);}
\foreach \k in {2,3,4}{\draw[->,thick,green!45!black] ($(n\k)+(-0.95,0)$) -- (n\k);}
\node[font=\small,green!45!black] at (-0.35,-0.45) {$I$};
\node[font=\Large] at (3.1,0) {$\approx$};
\node[nrn,minimum size=0.85cm,font=\small] (R) at (4.35,-0.1) {$r$};
\draw[->,thick,blue!60!black] (R) edge[loop above,looseness=6] node[above,font=\small] {$w$} (R);
\draw[->,thick,green!45!black] (3.55,-0.95) node[left,font=\small] {$I$} -- (R);
\end{scope}
\begin{scope}[xshift=6.3cm,y=2.6cm,x=1.05cm]
\node[font=\small] at (-0.6,1.25) {(b)};
\draw[->,gray!70!black] (0,0) -- (7.3,0) node[below left,font=\small,black] {$t/\tau$};
\draw[->,gray!70!black] (0,0) -- (0,1.12) node[left,font=\small,black] {$r$};
\foreach \t in {0,2,4,6}{\draw[gray!60!black] (\t,-0.02) -- (\t,0.02); \node[below,font=\scriptsize] at (\t,0) {\t};}
\draw[densely dotted,gray!60!black] (0,0.5) -- (7,0.5) node[right,font=\scriptsize,black,align=left] {ngưỡng\\ghi};
\fill[green!45!black,opacity=0.25] (1,-0.2) rectangle (2,-0.13);
\fill[gray!60!black,opacity=0.35] (1,-0.29) rectangle (1.5,-0.22);
\draw[very thick,blue!60!black] plot coordinates {(0.0,0.002) (0.1,0.002) (0.2,0.002) (0.3,0.002) (0.4,0.002) (0.5,0.002) (0.6,0.002) (0.7,0.002) (0.8,0.002) (0.9,0.002) (1.0,0.002) (1.1,0.083) (1.2,0.165) (1.3,0.242) (1.4,0.313) (1.5,0.378) (1.6,0.437) (1.7,0.491) (1.8,0.539) (1.9,0.583) (2.0,0.623) (2.1,0.643) (2.2,0.665) (2.3,0.688) (2.4,0.71) (2.5,0.732) (2.6,0.753) (2.7,0.773) (2.8,0.792) (2.9,0.81) (3.0,0.826) (3.2,0.855) (3.4,0.88) (3.6,0.9) (3.8,0.917) (4.0,0.931) (4.4,0.953) (4.8,0.967) (5.2,0.977) (5.6,0.984) (6.0,0.988) (6.5,0.992) (7.0,0.994)};
\draw[very thick,gray!60!black,densely dashed] plot coordinates {(0.0,0.002) (0.1,0.002) (0.2,0.002) (0.3,0.002) (0.4,0.002) (0.5,0.002) (0.6,0.002) (0.7,0.002) (0.8,0.002) (0.9,0.002) (1.0,0.002) (1.1,0.083) (1.2,0.165) (1.3,0.242) (1.4,0.313) (1.5,0.378) (1.6,0.358) (1.7,0.336) (1.8,0.313) (1.9,0.291) (2.0,0.269) (2.2,0.228) (2.4,0.191) (2.6,0.159) (2.8,0.133) (3.0,0.11) (3.2,0.091) (3.4,0.076) (3.6,0.063) (3.8,0.052) (4.0,0.043) (4.4,0.03) (4.8,0.021) (5.2,0.015) (5.6,0.011) (6.0,0.008) (6.5,0.006) (7.0,0.004)};
\node[font=\scriptsize,blue!60!black,anchor=west] at (3.3,0.8) {dòng vào $1\tau$: đã ghi};
\node[font=\scriptsize,gray!50!black,anchor=west] at (2.3,0.3) {dòng vào $0.5\tau$: tắt};
\end{scope}
\end{tikzpicture}
\caption{Một nhóm liên kết chặt với chính nó. (a)~Nhiều nơ-ron \nt{GLUT} kích thích lẫn nhau hành xử như một nơ-ron có tần số $r$, tự kích thích chính nó với độ mạnh $w$; từ bên ngoài có dòng vào $I$. (b)~Tần số của nhóm theo thời gian khi $w=1$ và cùng đường cong $f$ như ở hình~\ref{fig:bistable}. Dòng vào $I=0.6$ được cấp trong khoảng thời gian đánh dấu bằng dải dưới trục. Nếu nó kéo dài $1\tau$, nhóm kịp vượt qua điểm không bền $r=0.5$, bùng lên và tiếp tục cháy khi dòng vào đã hết. Nếu chỉ $0.5\tau$, nhóm không với tới điểm đó và tắt trở lại: để ghi một bit, dòng vào phải kéo dài hơn khoảng $0.7\tau$.}
\label{fig:recurrentgroup}
\end{figure}

\begin{figure}[t]
\centering
\begin{tikzpicture}[>=Stealth,x=5cm,y=3.2cm]
\draw[->,gray!70!black] (0,0) -- (1.1,0) node[below left,black] {\small $r$};
\draw[->,gray!70!black] (0,0) -- (0,1.1);
\draw[thick,gray!60!black] (0,0) -- (1.05,1.05) node[right,black] {\small $r$};
\draw[very thick,blue!60!black] plot[domain=0:1.05,samples=80] (\x,{1/(1+exp(-(\x-0.5)/0.08))});
\draw[very thick,red!70!black,densely dashed] plot[domain=0:1.05,samples=80] (\x,{1/(1+exp(-(\x+0.3-0.5)/0.08))});
\fill[blue!60!black] (0.002,0.002) circle (0.018);
\draw[blue!60!black,fill=white,thick] (0.5,0.5) circle (0.018);
\fill[blue!60!black] (0.998,0.998) circle (0.018);
\node[below right,font=\small] at (0.02,0.0) {tắt};
\node[right,font=\small] at (0.52,0.46) {không bền};
\node[below,font=\small] at (0.99,0.96) {bật};
\node[anchor=east,font=\small,red!70!black] at (0.19,0.62) {$f(wr+I)$};
\node[anchor=west,font=\small,blue!60!black] at (0.45,0.1) {$f(wr)$};
\end{tikzpicture}
\caption{Bộ nhớ từ một nhóm nơ-ron \nt{GLUT} có hồi tiếp mạnh. Đường liền là khi không có dòng vào từ bên ngoài: hai giao điểm bền với đường thẳng $r$, ``tắt'' và ``bật'', và một giao điểm không bền ở giữa. Một dòng vào ngắn $I$ (đường nét đứt) chỉ để lại giao điểm trên.}
\label{fig:bistable}
\end{figure}

Nếu hồi tiếp mạnh, có ba giao điểm: dưới là im lặng, trên là nhóm cháy hết cỡ, giữa là không bền. Ta được một phần tử có hai trạng thái bền, một flip-flop. Ghi số một vào nó rất dễ: một dòng vào ngắn nâng đường cong lên, giao điểm dưới biến mất, và nhóm bùng lên --- rồi vẫn cháy khi dòng vào đã hết (hình~\ref{fig:recurrentgroup}b). Dòng vào quá ngắn không kịp đưa nhóm tới điểm không bền, và nhóm tắt trở lại. Hoạt động tự duy trì như vậy, kéo dài hàng giây sau khi kích thích đã biến mất, thực sự được quan sát trong não và được coi là nền tảng của trí nhớ ngắn hạn~\cite{Wang2001}. Nhưng đó không phải cách duy nhất để nhớ trong vài giây. Bản ghi cũng có thể được lưu trong một biến chậm của chính các khớp thần kinh: sau một chùm xung, sự tạo thuận, như ở khớp thần kinh ``gạch'' của chương~\ref{sec:code}, kéo dài khoảng một giây, và giữa những đợt bùng ngắn mạng có thể gần như im lặng --- không phải flip-flop mà là một tụ điện đã nạp~\cite{Mongillo2008}. Những khoảng ``lặng'' như vậy khi giữ thông tin trong trí nhớ đã được quan sát~\cite{Stokes2015}, và lưu trữ không cần xung liên tục tốn ít năng lượng hơn. Vỏ não dùng cách nào trong trường hợp nào vẫn còn đang được bàn cãi.

Vậy xóa thế nào? Theo mệnh đề~\ref{prop:monotone}, bằng hoạt động thì không thể; chỉ còn cách \emph{lấy đi} một thứ gì đó. Cách thứ nhất là lấy đi sự hỗ trợ: nếu nhóm chỉ cháy khi có dòng vào từ một nhóm khác, bộ nhớ tắt cùng với dòng vào đó. Cách thứ hai là sự mỏi. Ở khớp thần kinh thật, khi làm việc lâu, lượng chất dẫn truyền dự trữ cạn dần~\cite{Tsodyks1997}, còn ở nơ-ron, các dòng chậm làm giảm tính kích thích tăng dần lên. Hồi tiếp yếu đi, giao điểm trên biến mất, và nhóm tự tắt sau vài giây. Ta được một bộ nhớ bật bằng tín hiệu nhưng chỉ tắt theo thời gian.

\section{Chuỗi}

Thay cho đồng hồ có thể có một \emph{bộ định thời khởi động theo sự kiện}. Nối các nhóm nơ-ron \nt{GLUT} thành chuỗi: nhóm thứ nhất kích thích nhóm thứ hai, nhóm thứ hai kích thích nhóm thứ ba, v.v. Một tín hiệu bên ngoài làm nhóm thứ nhất bùng lên, và một làn sóng hoạt động chạy dọc chuỗi, mỗi mắt xích sau mắt xích trước một hai độ trễ và chỉ một lần: ngay sau xung, nơ-ron ở trạng thái trơ, còn làn sóng đã đi tiếp.

Chuỗi như vậy đo thời gian tính từ sự kiện: mắt xích thứ $k$ đã phát xung --- nghĩa là từ lúc khởi động đã qua $k$ độ trễ. Đầu ra của chuỗi có thể đưa tới cơ để có một chuỗi cử động tự triển khai. Ở chim biết hót, trong lúc hót, từng nơ-ron \nt{GLUT} của một vùng não phát xung lần lượt nghiêm ngặt, mỗi nơ-ron ở một thời điểm riêng của bài hót --- rất giống một chuỗi như vậy~\cite{Hahnloser2002}.

Khác với đồng hồ, chuỗi im lặng cho đến khi được khởi động, chạy một lần và chỉ đo thời gian cho những ai được nối với nó. Mạng vẫn không có nhịp chung.

\section{Còn thiếu gì}

Chỉ với nơ-ron \nt{GLUT}, không có ức chế, đã làm được nhiều thứ. Nhưng còn thiếu ba điều, và cả ba đều do mệnh đề~\ref{prop:monotone}.

\emph{Lựa chọn.} Mạng phải chọn một hướng, một hành động. Nếu hai kích thích gần bằng nhau, trong mạng hình~\ref{fig:itd} cả hai đầu ra sẽ phát xung, và không có gì để dập đầu ra yếu nhường cho đầu ra mạnh.

\emph{Đặt lại.} Không thể xóa bộ nhớ bằng tín hiệu.

\emph{Ổn định.} Trong một mạng mà liên kết không hình thành theo thiết kế của kỹ sư, các vòng hồi tiếp dương là không tránh khỏi, và hoạt động trong đó có thể tăng như thác đổ (chương~\ref{sec:neurontypes}). Cái phanh duy nhất là sự mỏi nói trên, chậm và thô.

Cả ba vấn đề được chữa bằng một liều thuốc --- một nơ-ron dập xung bằng xung. Đó là \nt{GABA}, và nó cần không phải để tính toán, mà để lựa chọn, đặt lại và giữ phép tính trong tầm kiểm soát.

\section{Bài tập}

\begin{exercise}[\lvlA]
\label{ex:02:1}
\emph{Hàng bộ phát hiện cho âm thanh.} Tín hiệu từ hai tai chạy trên các dây của hình~\ref{fig:itd} với vận tốc $5$~m/s, hiệu thời gian đến lớn nhất là $0.64\ms$. Các bộ phát hiện --- nơ-ron~\eqref{eq:glutneuron} với $\tau=0.5\ms$, $\theta=1.5w$ --- được đặt sao cho hai bộ kề nhau phân biệt được $10\,\mu\text{s}$. Hàng có bao nhiêu bộ phát hiện và bao nhiêu bộ phát xung với một âm thanh?
\end{exercise}

\begin{exercise}[\lvlB]
\label{ex:02:2}
\emph{Đầu vào không đều làm lệch cửa sổ.} Bộ phát hiện trùng khớp~\eqref{eq:glutneuron} với $\tau=0.5\ms$, $\theta=1.5$ nhận mỗi đầu vào một xung, trọng số $1$ và $0.8$. Tìm cửa sổ trùng khớp và tâm của nó. Hàng ở hình~\ref{fig:itd} sẽ sai bao nhiêu micro giây nếu đầu vào từ một tai yếu hơn tai kia $20\,\%$?
\end{exercise}

\begin{exercise}[\lvlB]
\label{ex:02:3}
\emph{Những mạng không dao động.} Mạng $\tau\dot r_i=-r_i+f_i\bigl(\textstyle\sum_jw_{ij}r_j+I_i\bigr)$ với $f_i$ không giảm và $w_{ij}\ge0$ khi $i\ne j$ là đơn điệu và không dao động bền vững. Chứng minh: qua phép thay $r_i\to\pm r_i$, mạng có số liên kết ức chế chẵn trong mỗi chu trình quy được về nó. Ức chế lẫn nhau giữa hai nhóm có thể dao động không? Vòng \nt{GLUT}--\nt{GABA}?
\end{exercise}

\begin{exercise}[\lvlB]
\label{ex:02:4}
\emph{Thiết kế: bộ cộng hai dây.} Mỗi bit được truyền bằng một cặp dây $(x,\bar x)$. Hãy xây bộ cộng đủ ba bit, cho ra các cặp $(S,\bar S)$ của tổng và $(C,\bar C)$ của số nhớ, từ nơ-ron \nt{GLUT}, chỉ rõ trọng số và ngưỡng. Hãy dùng không quá tám nơ-ron. Một mạch gồm các phần tử AND và OR như ở hình~\ref{fig:dualxor} sẽ cần bao nhiêu?
\end{exercise}

\begin{exercise}[\lvlB]
\label{ex:02:5}
\emph{Mô phỏng: ghi mất bao lâu.} Nhóm ở hình~\ref{fig:recurrentgroup}: $\tau\dot r=-r+f(r+I)$, $f(u)=1/\bigl(1+e^{-(u-0.5)/0.08}\bigr)$, trước khi ghi ở trạng thái dưới; dòng vào $I$ được cấp trong thời gian $T$. Tìm bằng số $T$ nhỏ nhất ghi được bit khi $I=0.6$, và ngưỡng $I_c$ mà dưới nó bit không được ghi với bất kỳ $T$ nào.
\end{exercise}

\begin{exercise}[\lvlA]
\label{ex:02:6}
\emph{Chuỗi làm bộ định thời.} Mỗi mắt xích của chuỗi là $100$ nơ-ron \nt{GLUT}, độ trễ $2\ms$ với độ phân tán độc lập $0.2\ms$. (a)~Cần bao nhiêu nơ-ron để đo $1$~s, và sai số tương đối bằng bao nhiêu? Nó phụ thuộc vào khoảng thời gian thế nào? (b)~Bộ định thời thật có sai số $5\,\%$ với mọi khoảng. Nguồn phân tán nào cho ra điều đó?
\end{exercise}

\begin{exercise}[\lvlC]
\label{ex:02:7}
\emph{Nghiên cứu: flip-flop hay tụ điện.} $100$ nơ-ron giữ một bit trong $10$~s. Flip-flop: mỗi nơ-ron phát $20$ xung mỗi giây. Tụ điện: sự tạo thuận $u$ giảm với $\tau_F=1$~s, bit đọc được khi $u\ge u_{\max}/2$, làm tươi là một chùm ba xung ở mỗi nơ-ron. Tụ điện tiết kiệm hơn bao nhiêu lần, và bit sẽ ra sao sau khi nhóm bị làm câm $200\ms$?
\end{exercise}

\chapter{\nt{GLUT} và \nt{GABA} cùng nhau}
\label{sec:gaba}
\begin{chaptergoals}
\item cách phép cấm $a\wedge\bar b$ làm \nt{GLUT} và \nt{GABA} đầy đủ với một dây cho mỗi bit;
\item vì sao ức chế sun chia chứ không trừ, và cách ức chế thu hẹp cửa sổ trùng khớp;
\item cách ức chế lẫn nhau với $\beta>1$ chọn ra một kẻ thắng, và cách một tín hiệu đặt lại bộ nhớ;
\item khi nào ức chế nhanh giữ mạng ổn định, và khi nào ức chế chậm biến nó thành nhịp.
\end{chaptergoals}

\section{Luật chơi}

Mạng chỉ gồm nơ-ron \nt{GLUT} không biết lựa chọn, không đặt lại được bộ nhớ và không giữ được hoạt động khỏi thác đổ, và tất cả là do tính đơn điệu (mệnh đề~\ref{prop:monotone}). Hãy thêm vào đó loại nơ-ron đông thứ hai --- \nt{GABA}. Luật chơi vẫn như cũ: chỉ những gì có trong cơ thể sống, và mọi thứ diễn ra trong thời gian liên tục.

Mô hình nơ-ron vẫn là~\eqref{eq:glutneuron}, nhưng xung từ nơ-ron \nt{GABA} không nâng điện áp của nơi nhận mà hạ nó xuống. Giờ đây trọng số có hai dấu, và dấu do nơi gửi quyết định, theo quy tắc~\eqref{eq:dalesign}.

Vài tham số của mạch. Nơ-ron \nt{GABA} chiếm từ một phần năm đến một phần tư số nơ-ron vỏ não~\cite{Hendry1987}. Đầu ra của chúng ngắn: chúng ức chế láng giềng chứ không phải những vùng xa. Ức chế đến sau một mili giây và kéo dài vài mili giây. Không có đầu vào, nơ-ron \nt{GABA} im lặng: để ức chế ai đó, chính nó phải nhận được kích thích, thường là từ các nơ-ron \nt{GLUT} của cùng mạng.

Vì vậy liên kết âm từ nơ-ron \nt{GLUT} $A$ tới nơ-ron $B$ phải làm qua trung gian: $A$ kích thích một nơ-ron \nt{GABA}, và nơ-ron này ức chế $B$. Đổi dấu tốn một nơ-ron và một độ trễ, khoảng một mili giây.

Với ức chế, không chỉ quan trọng liên kết đến \emph{từ ai} mà cả nó đậu \emph{vào đâu}: vị trí đậu phần lớn quyết định liên kết làm gì~\cite{Markram2004}. Đầu ra \nt{GLUT} chủ yếu đậu vào các nhánh đầu vào của nơi nhận, tức \emph{đuôi gai}, và mang dữ liệu: mỗi đầu vào như vậy là một số hạng của tổng. Đầu ra đậu vào thân nơ-ron tác động lên cả tổng: đó là cách nhiều nơ-ron \nt{GABA} nhanh chia đáp ứng của nơi nhận và quyết định khi nào nó phát xung. Đầu ra đậu vào chỗ sinh ra xung của nơi nhận là lời cuối cùng, một lệnh cấm lên toàn bộ đầu ra cùng lúc. Còn đầu ra đậu vào đầu tận cùng của dây dẫn khác thay đổi xác suất nhả $p$ của một đường vào duy nhất, không động đến các đường khác~\cite{Rudomin1999}; chẳng hạn, cổng đau ở chương~\ref{sec:pain} hoạt động như vậy. Bên ngoài não, đầu ra đậu vào sợi cơ (chương~\ref{sec:muscles}) và vào các cơ quan (chương~\ref{sec:chemoutput}), còn một số không đậu vào đâu cả mà nhả chất dẫn truyền vào thể tích mô hoặc vào máu (chương~\ref{sec:serocomputer} và~\ref{sec:chemoutput}).

\section{NOT và phép cấm}

Giờ có làm được NOT không? Gần như được. Một nơ-ron hoạt động khi đầu vào im lặng phải lấy kích thích từ đâu đó. Trong não sống thì có: mỗi nơ-ron liên tục nhận hoạt động nền từ hàng nghìn nơ-ron khác. Để nền giữ nơ-ron $Y$ hoạt động, còn đầu vào $x$ ức chế nó qua một trung gian \nt{GABA}. Đó là NOT.

Tuy vậy, hữu ích hơn là \emph{phép cấm}: ``$a$, nhưng không khi có $b$'':
\begin{empheq}[box=\keybox]{equation}
y \;=\; a \wedge \bar b .
\label{eq:veto}
\end{empheq}
Nơ-ron $Y$ nhận kích thích từ $a$ và ức chế từ nơ-ron \nt{GABA} do $b$ kích thích. Ở đây không cần nền: chính $a$ cung cấp kích thích. Từ phép cấm và OR có thể lắp mọi thứ. Chẳng hạn XOR: $x\oplus y=(x\wedge\bar y)\vee(\bar x\wedge y)$ --- hai nơ-ron cấm, hai trung gian \nt{GABA} và một nơ-ron OR.

\begin{keyblock}
\begin{proposition}[Tính đầy đủ của \nt{GLUT}+\nt{GABA}]
\label{prop:gabacomplete}
Mạng nơ-ron \nt{GLUT} và \nt{GABA} có thể tính bất kỳ hàm logic nào của các đầu vào, và mỗi bit được truyền bằng một dây. Không còn cần các cảm biến ``bật---tắt'' của chương~\ref{sec:glutcomputer} cho việc này. Nếu hàm phải cho giá trị một khi mọi đầu vào im lặng, cần một nguồn hoạt động nền.
\end{proposition}
\end{keyblock}

Chứng minh: phép cấm~\eqref{eq:veto} cùng với OR là đầy đủ về mặt chức năng nếu có hằng một, vì NOT $x$ chính là phép cấm ``một, nhưng không khi có $x$''. Còn các hàm cho giá trị không khi mọi đầu vào im lặng thì lắp được từ phép cấm và OR mà không cần hằng một.

\section{Hướng chuyển động}

Phép cấm theo thời gian, cũng như AND theo thời gian ở chương~\ref{sec:glutcomputer}, cho ta một bộ phát hiện hướng chuyển động.

Cho hai cảm biến kề nhau, $A$ và $B$, đặt cách nhau $\Delta x$. Nơ-ron đầu ra $Y$ nhận kích thích từ $B$ và ức chế từ $A$ qua một trung gian \nt{GABA}, với độ trễ $d$ (hình~\ref{fig:direction}). Một đốm sáng chạy với vận tốc $v$ từ $A$ sang $B$ đi qua $A$ lúc $0$, và ức chế tới $Y$ lúc $d$; nó đi qua $B$ lúc $\Delta x/v$, và kích thích tới cùng lúc đó. Nếu hai thời điểm này trùng nhau trong phạm vi cửa sổ~\eqref{eq:window}, ức chế dập tắt kích thích, và $Y$ im lặng. Điều này xảy ra ở vận tốc khoảng
\begin{empheq}[box=\keybox]{equation}
v^* \;=\; \frac{\Delta x}{d} .
\label{eq:vstar}
\end{empheq}
Khi chuyển động từ $B$ sang $A$, kích thích từ $B$ đến lúc $0$, còn ức chế từ $A$ phải đến lúc $\Delta x/v+d$ mới tới, khi $Y$ đã phát xung.

\begin{figure}[t]
\centering
\begin{tikzpicture}[>=Stealth,
  nrn/.style={circle,draw,very thick,fill=blue!6,minimum size=0.9cm,inner sep=0pt},
  gab/.style={nrn,fill=red!10},
  inh/.style={-{Bar[width=7pt]},thick,red!70!black}]
\node[nrn] (A) at (0,2) {$A$};
\node[nrn] (B) at (3,2) {$B$};
\node[gab] (G) at (0.9,0.6) {\small \nt{GABA}};
\node[nrn] (Y) at (3,-0.6) {$Y$};
\draw[->,thick] (A) -- (G);
\draw[inh] (G) -- node[below left=-2pt] {\scriptsize độ trễ $d$} (Y);
\draw[->,thick] (B) -- (Y);
\draw[->,very thick,red!70!black] (Y) -- (4.6,-0.6) node[right,black] {\small đầu ra};
\draw[->,thick,orange!85!black] (-0.6,2.9) -- (3.6,2.9) node[right,black] {\small im lặng};
\draw[<-,thick,orange!85!black] (-0.6,3.4) -- (3.6,3.4) node[right,black] {\small đáp ứng};
\foreach \yy/\dd in {2.9/-1,3.4/1}{
  \foreach \k in {1,2,3}{\draw[orange!70!yellow,line width=0.8pt] (1.5+\dd*0.12+\dd*0.13*\k,\yy+0.1-0.1*\k+0.1) -- ++(\dd*0.25,0);}
  \fill[yellow!70!orange,draw=orange!70!black] (1.5,\yy) circle (0.14);}
\node[font=\scriptsize,align=center] at (1.5,3.95) {đốm sáng};
\end{tikzpicture}
\caption{Bộ phát hiện hướng chuyển động. $B$ kích thích $Y$, $A$ ức chế nó qua một trung gian \nt{GABA} với độ trễ $d$. Khi chuyển động từ $A$ sang $B$ với vận tốc khoảng $\Delta x/d$, ức chế dập tắt kích thích; khi chuyển động ngược lại, nó đến muộn. Đốm vàng có vạch là ánh sáng chạy phía trên các cảm biến.}
\label{fig:direction}
\end{figure}

Trong võng mạc, tính chọn lọc hướng chuyển động dường như được tạo ra đúng theo cách này: bằng ức chế bất đối xứng từ phía mà chuyển động từ đó không được gây đáp ứng~\cite{Barlow1965}.

\section{Phép trừ hay phép chia}

Cho đến giờ, ức chế ở ta làm phép trừ. Thực ra nó tinh tế hơn.

Khớp thần kinh mở một kênh, tức là bật một độ dẫn. Ở khớp thần kinh kích thích, điện thế cân bằng nằm cao hơn ngưỡng nhiều, khoảng $70\mV$ trên mức nghỉ: kênh mở kéo điện áp lên. Ở khớp thần kinh ức chế \nt{GABA}, kênh cho clo đi qua, điện thế cân bằng của clo nằm gần mức nghỉ, và kênh mở không kéo điện áp xuống mà kéo nó \emph{về mức nghỉ}. Nếu tính điện áp từ mức nghỉ, điện áp xác lập của màng có độ dẫn rò $g_L$, độ dẫn kích thích $g_E$ và độ dẫn ức chế $g_I$ bằng
\begin{empheq}[box=\keybox]{equation}
V_\infty \;=\; \frac{g_E\,E_E}{g_L+g_E+g_I} ,
\label{eq:shunt}
\end{empheq}
trong đó $E_E\approx70\mV$ là điện thế cân bằng của khớp thần kinh kích thích. Đây là định luật Ohm cho một bộ chia áp: $g_E$ kéo về $E_E$, còn $g_L$ và $g_I$ kéo về không.

$g_I$ không đứng ở tử số với dấu trừ mà ở mẫu số: ức chế không \emph{trừ} khỏi kích thích mà \emph{chia} nó (hình~\ref{fig:shunt}). Kiểu ức chế này được gọi là ức chế sun: các kênh clo mở làm ngắn mạch màng. Với mạng, đó là điều chỉnh hệ số khuếch đại. Nếu $g_I$ tỷ lệ với tổng hoạt động của các láng giềng, mỗi nơ-ron đáp ứng không phải với độ mạnh tuyệt đối của đầu vào của nó, mà với tỷ phần của nó trong tổng hoạt động. Phép chia cho tổng hoạt động như vậy gặp ở nhiều vùng não và được coi là một trong các phép toán tiêu chuẩn của não~\cite{Carandini2012}: cùng một mạng làm việc được cả với đầu vào yếu lẫn đầu vào mạnh.

Lưu ý: công thức~\eqref{eq:shunt} nói về nơ-ron không phát xung. Ở nơ-ron đang phát xung, điện áp luôn giữ gần ngưỡng, dòng qua độ dẫn ức chế gần như không đổi, và với \emph{tần số} xung, sun tác động chủ yếu bằng phép trừ~\cite{HoltKoch1997}. Tần số bị chia nếu đồng thời thêm cho nơ-ron các độ dẫn nền kích thích và ức chế cân bằng nhau, như trong chế độ cân bằng, dao động liên tục của vỏ não sống~\cite{Chance2002}. Vậy phép chia não có được không phải từ riêng một sun, mà từ ức chế cùng với nhiễu nền~\cite{Carandini2012}.

\begin{figure}[t]
\centering
\begin{tikzpicture}[>=Stealth,x=1.25cm,y=0.05cm]
\draw[->,gray!70!black] (0,0) -- (5.4,0) node[right,black] {\small $g_E/g_L$};
\draw[->,gray!70!black] (0,0) -- (0,68) node[above,black] {\small $V_\infty$, mV};
\foreach \v in {20,40,60}{\draw[gray!60!black] (-0.05,\v) -- (0.05,\v) node[left=3pt,font=\scriptsize] {\v};}
\foreach \x in {1,2,3,4,5}{\draw[gray!60!black] (\x,-1.5) -- (\x,1.5) node[below=5pt,font=\scriptsize] {\x};}
\draw[very thick,blue!60!black] plot[domain=0:5,samples=60] (\x,{70*\x/(1+\x)});
\draw[very thick,red!70!black] plot[domain=0:5,samples=60] (\x,{70*\x/(3+\x)});
\draw[thick,densely dashed,gray!50!black] plot[domain=0.56:5,samples=60] (\x,{70*\x/(1+\x)-25});
\node[right,font=\small,blue!60!black] at (5.02,58.3) {không ức chế};
\node[right,font=\small,red!70!black] at (5.02,45.5) {chia: $g_I=2g_L$};
\node[right,font=\small,gray!40!black] at (5.02,32.5) {trừ $25\mV$};
\end{tikzpicture}
\caption{Ức chế sun chia điện áp chứ không trừ khỏi nó. Đường xanh là điện áp xác lập~\eqref{eq:shunt} khi không có ức chế, đường đỏ là khi có độ dẫn ức chế $g_I=2g_L$. Đường đỏ chính là đường xanh bị kéo giãn gấp ba theo chiều ngang: như thể đầu vào bị chia cho $1+g_I/g_L=3$. Đường nét đứt là ức chế trừ đi một lượng không đổi $25\mV$: cả đường cong hạ xuống, độ dốc không đổi.}
\label{fig:shunt}
\end{figure}

Còn một hệ quả thứ hai. Hằng số thời gian hiệu dụng của màng là $\tau_{\text{hd}}=C/(g_L+g_E+g_I)$, và ức chế làm nó ngắn lại. Mà cửa sổ trùng khớp~\eqref{eq:window} tỷ lệ với $\tau$, nên ức chế đến cùng với kích thích \emph{thu hẹp cửa sổ}. Mạch trong đó đầu vào kích thích nơ-ron trực tiếp và đồng thời, với độ trễ một mili giây, ức chế nó qua một trung gian \nt{GABA} là một trong những mạch phổ biến nhất trong não: nơ-ron chỉ kịp đáp ứng với những gì đến trong một hai mili giây đầu tiên~\cite{Pouille2001}.

\section{Lựa chọn}

Giờ đến điều chính mà mạng \nt{GLUT} còn thiếu: chọn một trong nhiều. Lấy hai nhóm nơ-ron \nt{GLUT} với đầu vào $I_1$ và $I_2$. Mỗi nhóm, qua các trung gian \nt{GABA} của mình, ức chế nhóm kia với độ mạnh $\beta$ (hình~\ref{fig:wta}). Với các tần số $r_1,r_2$ ta có
\begin{equation}
\tau\,\frac{\dd r_1}{\dd t} = -r_1 + \bigl[I_1-\beta r_2\bigr]_+ ,\qquad
\tau\,\frac{\dd r_2}{\dd t} = -r_2 + \bigl[I_2-\beta r_1\bigr]_+ ,
\label{eq:wta}
\end{equation}
trong đó $[u]_+=\max(u,0)$: tần số không thể âm.

\begin{figure}[t]
\centering
\begin{tikzpicture}[>=Stealth,
  nrn/.style={circle,draw,very thick,fill=blue!6,minimum size=0.9cm,inner sep=0pt},
  gab/.style={nrn,fill=red!10,minimum size=0.75cm},
  inh/.style={-{Bar[width=7pt]},thick,red!70!black}]
\node[nrn] (E1) at (0,0) {$r_1$};
\node[nrn] (E2) at (4,0) {$r_2$};
\node[gab] (G1) at (1.4,1.1) {};
\node[gab] (G2) at (2.6,-1.1) {};
\draw[->,thick] (E1) -- (G1); \draw[inh] (G1) -- (E2);
\draw[->,thick] (E2) -- (G2); \draw[inh] (G2) -- (E1);
\draw[->,thick,blue!60!black] (-1.6,0) node[left,black] {$I_1$} -- (E1);
\draw[->,thick,blue!60!black] (5.6,0) node[right,black] {$I_2$} -- (E2);
\end{tikzpicture}
\caption{Chọn một trong hai. Hai nhóm nơ-ron \nt{GLUT} (màu xanh) ức chế lẫn nhau qua các trung gian \nt{GABA} (màu đỏ).}
\label{fig:wta}
\end{figure}

Chừng nào cả hai nơ-ron cùng hoạt động, hệ~\eqref{eq:wta} là tuyến tính, và hành vi của nó được quyết định bởi hai trị riêng: $-(1+\beta)/\tau$ cho các độ lệch cùng chiều của hai tần số và $-(1-\beta)/\tau$ cho các độ lệch ngược chiều. Với ức chế yếu, $\beta<1$, cả hai trị riêng đều âm: cả hai nơ-ron hoạt động, ức chế chỉ làm nơ-ron yếu hơn nhỏ tiếng đi. Với ức chế mạnh, $\beta>1$, trị riêng thứ hai trở thành dương: mọi chênh lệch giữa $r_1$ và $r_2$ đều tăng lên cho đến khi một nơ-ron im lặng hẳn.

\begin{keyblock}
\begin{proposition}[Kẻ thắng lấy tất]
\label{prop:wta}
Nếu ức chế lẫn nhau mạnh hơn một, $\beta>1$, trạng thái trong đó cả hai nhóm cùng hoạt động là không bền. Mạng đi tới trạng thái chỉ một nhóm hoạt động, nhóm kia im lặng. Kẻ thắng với tần số $r_1=I_1$ giữ được chiến thắng chừng nào $I_2<\beta I_1$.
\end{proposition}
\end{keyblock}

Chúng tôi đã kiểm tra điều này trên mô hình. Với $\beta=0.5$ và các đầu vào $1.0$ và $0.9$, cả hai nhóm hoạt động, với tần số $0.73$ và $0.53$. Với $\beta=2$ và cùng các đầu vào đó, nhóm thứ hai im lặng hoàn toàn. Lựa chọn như vậy có bộ nhớ: nếu sau đó đổi chỗ các đầu vào, kẻ thắng cũ vẫn là kẻ thắng, vì $1.0<2\cdot0.9$.

Với một trăm nhóm cũng vậy: mỗi nhóm ức chế các nhóm còn lại, chỉ một nhóm sống sót.

\section{Đặt lại bộ nhớ}

Bộ nhớ ở chương~\ref{sec:glutcomputer} chỉ tắt được nhờ sự mỏi. Thêm vào đó một nơ-ron \nt{GABA} được kích thích bởi tín hiệu ``đặt lại'' và ức chế nhóm. Ức chế hạ đường cong $f(wr+I)$ trên hình~\ref{fig:bistable}, giao điểm trên biến mất, và nhóm tắt --- rồi ở lại trạng thái dưới khi tín hiệu đặt lại đã hết. Giờ bộ nhớ được ghi bằng một tín hiệu và xóa bằng một tín hiệu khác, như flip-flop có đầu vào đặt và đặt lại.

Hay hơn nữa là nối hai nhóm như vậy bằng ức chế lẫn nhau, như ở hình~\ref{fig:wta}, và cho mỗi nhóm hồi tiếp lên chính nó. Tín hiệu vào một trong hai nhóm chuyển ô nhớ sang trạng thái của nhóm đó, và ô nhớ ghi nhớ quyết định sau cùng. Đây là tương tự trực tiếp của mạch chốt gồm hai phần tử NOR nối chéo.

\section{Ổn định và nhịp}

Còn lại tai họa thứ ba của mạng \nt{GLUT} --- thác hoạt động. Lấy một nhóm nơ-ron \nt{GLUT} có hồi tiếp lên chính nó $w_{EE}$ và một nhóm nơ-ron \nt{GABA} được nhóm thứ nhất kích thích với độ mạnh $w_{IE}$ và ức chế nhóm thứ nhất với độ mạnh $w_{EI}$. Với các tần số $E$ và $I$, ta có
\begin{equation}
\tau_E\frac{\dd E}{\dd t} = -E + w_{EE}E - w_{EI}I + h, \qquad
\tau_I\frac{\dd I}{\dd t} = -I + w_{IE}E ,
\label{eq:ei}
\end{equation}
trong đó $h$ là đầu vào bên ngoài~\cite{WilsonCowan1972}. Không có ức chế, khi $w_{EE}>1$ hoạt động tăng không giới hạn: mỗi xung sinh ra hơn một xung tiếp theo. Có ức chế, tần số xác lập
\begin{equation}
E^* \;=\; \frac{h}{1-w_{EE}+w_{EI}w_{IE}}
\label{eq:eistar}
\end{equation}
là hữu hạn nếu mẫu số dương. Nhưng như vậy chưa đủ: trạng thái cân bằng còn phải hút. Phân tích tuyến tính~\eqref{eq:ei} cho hai điều kiện.

\begin{keyblock}
\begin{proposition}[Điều kiện ổn định]
\label{prop:ei}
Trạng thái cân bằng~\eqref{eq:eistar} ổn định nếu ức chế
\begin{enumerate}
\item[(i)] \emph{đủ mạnh}: $w_{EI}\,w_{IE} > w_{EE}-1$;
\item[(ii)] \emph{đủ nhanh}: $\tau_I < \tau_E/(w_{EE}-1)$.
\end{enumerate}
Nếu (i) bị vi phạm, hoạt động tăng như thác đổ. Nếu (i) được thỏa mãn nhưng (ii) bị vi phạm, ức chế đến muộn, và hoạt động bắt đầu dao động.
\end{proposition}
\end{keyblock}

Điều kiện~(i) là định thức của ma trận hệ~\eqref{eq:ei}, điều kiện~(ii) là vết của nó. Ý nghĩa của điều kiện thứ hai rõ với bất kỳ ai từng chỉnh một bộ điều khiển: hồi tiếp đến muộn không dập độ lệch mà làm nó dao động mạnh thêm.

\begin{figure}[t]
\centering
\begin{tikzpicture}[>=Stealth]
\draw[->,gray!70!black] (0,0) -- (10.9,0) node[right,black] {\small $t$, ms};
\draw[->,gray!70!black] (0,0) -- (0,3.3) node[left,black] {\small $E$};
\foreach \t in {0,50,100,150}{\draw[gray!70!black] (\t*0.07,0.06) -- (\t*0.07,-0.06) node[below,font=\scriptsize] {\t};}
\input{figures/ei_traces}
\node[font=\small,gray!40!black,anchor=west] at (1.3,3.45) {không ức chế: thác hoạt động};
\node[font=\small,blue!60!black,anchor=west] at (7.4,1.25) {ức chế nhanh};
\node[font=\small,red!70!black,anchor=west] at (7.4,2.55) {ức chế chậm};
\end{tikzpicture}
\caption{Tần số của một nhóm \nt{GLUT} có hồi tiếp $w_{EE}=1.5$, $\tau_E=10\ms$, sau khi bật đầu vào. Không có ức chế (nét đứt), hoạt động tăng không giới hạn. Với ức chế có độ mạnh $w_{EI}w_{IE}=2$ và $\tau_I=2\ms$ (đường xanh), nó đạt mức $E^*=h/(1-1.5+2)=0.67h$. Với cùng ức chế đó nhưng chậm, $\tau_I=30\ms$ (đường đỏ), hoạt động dao động.}
\label{fig:ei}
\end{figure}

Người ta cho rằng vỏ não làm việc đúng ở chế độ này: hồi tiếp riêng của nó đủ mạnh để rơi vào thác hoạt động nếu không có ức chế, và chỉ có ức chế nhanh giữ nó ở điểm làm việc~\cite{Tsodyks1997isn}.

Chế độ này có một dấu hiệu nghịch lý, nhờ đó có thể nhận ra nó từ bên ngoài~\cite{Tsodyks1997isn}. Thêm vào~\eqref{eq:ei} một đầu vào bên ngoài $h_I$ đưa thẳng tới nhóm \nt{GABA}. Các tần số xác lập trở thành
\begin{equation}
E^* \;=\; \frac{h-w_{EI}\,h_I}{D},\qquad
I^* \;=\; \frac{w_{IE}\,h+(1-w_{EE})\,h_I}{D},\qquad
D \;=\; 1-w_{EE}+w_{EI}w_{IE} .
\label{eq:isn}
\end{equation}
Nếu $w_{EE}>1$, hệ số của $h_I$ trong công thức thứ hai là âm: đẩy thêm các nơ-ron \nt{GABA} --- và tần số của chúng \emph{giảm}. Nguyên nhân nằm ở vòng lặp. Ức chế thừa làm nhóm \nt{GLUT} nhỏ tiếng đi, nhóm này kích thích các nơ-ron \nt{GABA} ít hơn, và chúng mất nhiều hơn phần nhận được. Với các số liệu của hình~\ref{fig:ei} ($w_{EE}=1.5$, $w_{EI}w_{IE}=2$, $D=1.5$), mỗi đơn vị đầu vào thêm vào làm $I^*$ giảm một phần ba đơn vị. Dấu hiệu này đã được tìm thấy trong vỏ não: khi đáp ứng của nơ-ron vỏ thị giác bị kích thích xung quanh ức chế, dòng ức chế mà chúng nhận được không tăng mà giảm cùng với dòng kích thích~\cite{Ozeki2009}. Về sau, cùng dấu hiệu này đã được tìm thấy ở nhiều vùng và nhiều lớp vỏ não khác nhau~\cite{Sanzeni2020}.

Dao động khi điều kiện~(ii) bị vi phạm cũng gặp trong não. Vòng ``\nt{GLUT} kích thích \nt{GABA}, \nt{GABA} ức chế \nt{GLUT}'' với độ trễ vài mili giây sinh ra một nhịp có chu kỳ cỡ $12$--$50\ms$ (tần số $20$--$80$\,Hz), và các nhịp nhanh như vậy trong vỏ não được biết rõ~\cite{Cardin2009,Buzsaki2012}. Đây không phải đồng hồ của máy tính: nhịp mang tính cục bộ và chỉ xuất hiện khi mạng hoạt động. Nhưng trong vài chu kỳ, nó chia thời gian thành những cửa sổ mà trong đó nơ-ron có thể phát xung.

\section{Tổng kết}

\begin{table}[t]
\centering
\small
\begin{tabular}{>{\raggedright}p{3.3cm}>{\raggedright}p{4.4cm}>{\raggedright\arraybackslash}p{4.4cm}}
\toprule
& chỉ \nt{GLUT} & \nt{GLUT} + \nt{GABA} \\
\midrule
NOT & chỉ qua cảm biến ``bật---tắt'' & phép cấm $a\wedge\bar b$; NOT khi có hoạt động nền \\
số dây cho một bit & hai & một \\
hướng chuyển động & không & phép cấm có độ trễ \\
điều chỉnh hệ số khuếch đại & không & phép chia: sun cùng với nhiễu nền \\
cửa sổ trùng khớp & $\sim\tau$ & bị ức chế thu hẹp \\
chọn một trong nhiều & không & ức chế lẫn nhau, $\beta>1$ \\
đặt lại bộ nhớ & chỉ nhờ sự mỏi & bằng tín hiệu qua \nt{GABA} \\
ổn định & chỉ nhờ sự mỏi & hồi tiếp âm nhanh \\
nhịp & không cần & xuất hiện khi ức chế chậm \\
\bottomrule
\end{tabular}
\caption{\nt{GABA} bổ sung gì cho mạng nơ-ron \nt{GLUT}.}
\label{tab:gaba}
\end{table}

\nt{GABA} hầu như không bổ sung cho mạng \emph{phép tính} mới nào --- mọi thứ đó đều tính được với cảm biến hai dây --- mà bổ sung \emph{sự điều khiển} (bảng~\ref{tab:gaba}): lựa chọn, đặt lại, điều chỉnh hệ số khuếch đại và ổn định. Kích thích trong não mang dữ liệu, còn ức chế quyết định dữ liệu nào được đi qua, khi nào và to đến mức nào. Với mỗi nơ-ron thứ tư hoặc thứ năm của vỏ não, đó là một sự phân công lao động rất hợp lý.

\section{Bài tập}

\begin{exercise}[\lvlA]
\label{ex:03:1}
\emph{Sun bằng con số.} $E_E=70\mV$. Theo~\eqref{eq:shunt}, tìm $V_\infty$ khi $g_E=g_L$ và $4g_L$, không có ức chế và có sun $g_I=2g_L$. Phép trừ lấy đi ở $g_E=g_L$ đúng bằng lượng mà sun lấy đi sẽ cho $V_\infty$ bằng bao nhiêu khi $g_E=4g_L$? Sun thu hẹp cửa sổ trùng khớp bao nhiêu lần?
\end{exercise}

\begin{exercise}[\lvlB]
\label{ex:03:2}
\emph{Suy ra điều kiện ổn định.} Tìm vết và định thức của ma trận hệ tuyến tính~\eqref{eq:ei} và từ đó suy ra cả hai điều kiện của mệnh đề~\ref{prop:ei}. Ở biên của điều kiện~(ii), trạng thái cân bằng mất ổn định theo kiểu dao động. Tìm chu kỳ của các dao động đó với số liệu của hình~\ref{fig:ei}.
\end{exercise}

\begin{exercise}[\lvlB]
\label{ex:03:3}
\emph{Nhịp từ độ trễ.} Thay ức chế chậm của mô hình~\eqref{eq:ei} bằng ức chế nhanh có độ trễ $d$: $\tau_E\dot E=-E(t)-GE(t-d)+h$. Với $\tau_E=10\ms$, tìm độ trễ nhỏ nhất $d_c$ gây dao động, và chu kỳ của chúng với $G=2$ và $G=5$. Khác với bài tập~\ref{ex:03:2}, chúng có rơi vào các nhịp nhanh của vỏ não, $12$--$50\ms$, không?
\end{exercise}

\begin{exercise}[\lvlB]
\label{ex:03:4}
\emph{Mô phỏng: lựa chọn có bộ nhớ.} Tích phân~\eqref{eq:wta} với $\beta=2$, $\tau=1$. (a)~Với $I_1=1$, tăng chậm $I_2$ từ $0$ đến $3$ rồi giảm trở lại: ở những giá trị $I_2$ nào mạng chuyển trạng thái? (b)~Thời gian ra quyết định từ trạng thái im lặng tăng bao nhiêu khi hiệu các đầu vào $\varepsilon$ giảm mười lần?
\end{exercise}

\begin{exercise}[\lvlB]
\label{ex:03:5}
\emph{Thiết kế: bộ phát hiện cho một dải vận tốc.} Bộ phát hiện ở hình~\ref{fig:direction} phải im lặng khi chuyển động từ $A$ sang $B$ với thời gian chạy $5$--$20\ms$. Trung gian \nt{GABA} với độ trễ $d_k$ cấm $Y$ trên đoạn $[d_k,\,d_k+5\ms]$ sau xung của $A$; kích thích từ $B$ đến ngay. Cần bao nhiêu trung gian và với những độ trễ nào?
\end{exercise}

\begin{exercise}[\lvlB]
\label{ex:03:6}
\emph{Phép cấm và nền.} Hàm chẵn lẻ $x\oplus y\oplus z$ cần bao nhiêu nơ-ron theo sơ đồ chung ``một trung gian \nt{GABA} cho mỗi đầu vào, một phép cấm cho mỗi tổ hợp có $f=1$, một OR chung''? Hãy xây nó từ hai nơ-ron, \nt{GABA} và \nt{GLUT}, nhận đầu vào trực tiếp, chỉ rõ trọng số và ngưỡng.
\end{exercise}

\begin{exercise}[\lvlC]
\label{ex:03:7}
\emph{Thiết kế: chọn một trong một trăm.} Mỗi nhóm trong $100$ nhóm \nt{GLUT} phải ức chế đều mọi nhóm khác, trừ chính nó. Trung gian \nt{GABA} tuyến tính do các nhóm kích thích và ức chế các nhóm; các nhóm có thể kích thích nhau, trừ chính mình. Cần ít nhất bao nhiêu trung gian? Nếu không có liên kết kích thích trực tiếp nào?
\end{exercise}

\begin{exercise}[\lvlC]
\label{ex:03:8}
\emph{Nghiên cứu: khi nào ức chế nghịch lý.} Trong mô hình hình~\ref{fig:ei} ($w_{EI}=2$, $w_{IE}=1$), nhóm \nt{GLUT} có hàm truyền $f(u)=[u]_+^2$, nhóm \nt{GABA} vẫn tuyến tính. Trên đầu vào $h_c$ nào, việc thêm đầu vào cho nhóm \nt{GABA} làm giảm chính tần số của nó? Kiểm tra bằng mô phỏng.
\end{exercise}

\chapter{Mã Morse: những gì ta biết về ngôn ngữ của nơ-ron \nt{GLUT} và \nt{GABA}}
\chaptermark{Mã Morse}
\label{sec:code}

\section{Bảng chữ cái chỉ có một ký hiệu}

Mã Morse có hai ký hiệu, chấm và gạch, cùng các khoảng lặng với ba độ dài giữa chúng. Như ta đã thấy ở chương~\ref{sec:neurontypes}, bảng chữ cái của nơ-ron còn nghèo hơn: chỉ có một ký hiệu. Một xung kéo dài khoảng một mili giây, và mọi xung của một nơ-ron đều giống nhau về độ cao và hình dạng. Vì thế toàn bộ thông điệp nằm trong các khoảng lặng, tức là trong dãy thời điểm $t_1,t_2,t_3,\dots$ Đó là một thứ mã không có gạch, trong đó chỉ nhịp của các chấm mang ý nghĩa (hình~\ref{fig:morse}).

\begin{figure}[t]
\centering
\begin{tikzpicture}[>=Stealth,x=0.08cm,y=1cm]
% Morse letter: dot, dash, dot
\node[left,font=\small] at (-6,2.55) {Morse:};
\fill[black] (0,2.45) rectangle (6,2.65);
\fill[black] (14,2.45) rectangle (32,2.65);
\fill[black] (40,2.45) rectangle (46,2.65);
\node[right,font=\small,gray!40!black] at (52,2.55) {chấm, gạch, chấm --- ý nghĩa ở độ dài và khoảng lặng};
% stimulus onset
\node[left,font=\small] at (-6,0.4) {nơ-ron:};
\draw[->,thick,green!45!black] (0,-0.55) -- (0,-0.05);
\node[below,font=\scriptsize,green!45!black] at (0,-0.55) {kích thích};
\draw[gray!70!black] (0,0) -- (152,0) node[right,black] {\small $t$};
% spikes: single ones blue, the burst red
\foreach \s in {14,26,41,88,112,131}{\draw[very thick,blue!60!black] (\s,0) -- (\s,0.8);}
\foreach \s in {60,63,66,69}{\draw[very thick,red!70!black] (\s,0) -- (\s,0.8);}
% latency
\draw[<->,thick] (0,1.1) -- node[above,font=\scriptsize] {$t_1$: thời điểm xung đầu} (14,1.1);
% burst
\draw[decorate,decoration={brace,amplitude=4pt},red!70!black] (58,0.95) -- (71,0.95) node[midway,above=4pt,font=\scriptsize] {chùm xung --- ``gạch''};
% counting windows
\foreach \a/\b/\n in {0/50/3,50/100/5,100/150/2}{
  \draw[decorate,decoration={brace,amplitude=4pt,mirror},gray!50!black] (\a+1,-0.2) -- (\b-1,-0.2) node[midway,below=4pt,font=\small] {$n=\n$};}
\node[font=\scriptsize,gray!40!black] at (75,-1.05) {mã tần số: số xung $n$ trong mỗi cửa sổ};
\foreach \x in {0,50,100,150}{\draw[gray!60!black] (\x,-0.06) -- (\x,0.06);}
\end{tikzpicture}
\caption{Mã Morse và mã của nơ-ron. Cùng một chuỗi xung có thể đọc theo nhiều cách: đếm số xung trong các cửa sổ $50\ms$ (mục~\ref{sec:rate}), đo độ trễ $t_1$~\eqref{eq:latency}, hoặc nhận ra một chùm xung --- ``gạch''~\eqref{eq:burst}.}
\label{fig:morse}
\end{figure}

Những khoảng lặng nào là có thể? Từ phía dưới, chúng bị giới hạn bởi thời gian trơ: sau một xung, nơ-ron không thể phát xung lại trong khoảng một mili giây, nên không có quá vài trăm xung mỗi giây. Từ phía trên, không có gì giới hạn: nơ-ron có thể im lặng hàng phút. Ở nơ-ron \nt{GLUT} vỏ não, tần số nằm trong khoảng từ một phần nhỏ của một xung đến vài chục xung mỗi giây, và phần lớn nơ-ron phần lớn thời gian làm việc ở gần giới hạn dưới.

Ở các chương~\ref{sec:glutcomputer} và~\ref{sec:gaba}, ta đã tiện thể chấp nhận khi thì cách này, khi thì cách khác để ghi số bằng xung. Giờ hãy hỏi thẳng: nơ-ron \nt{GLUT} và \nt{GABA} nói với nhau bằng ngôn ngữ nào? Chưa có câu trả lời đầy đủ, ngôn ngữ này chưa được giải mã. Nhưng có một số điều về nó đã được biết chắc chắn, và gần như mọi điều đã biết đều có thể suy ra bằng cách lập luận như một kỹ sư truyền thông: dung lượng kênh, nhiễu, máy thu, giá của việc truyền.

\section{Truyền được bao nhiêu}

Bắt đầu từ cận trên. Giả sử nơi nhận phân biệt thời điểm xung với độ chính xác $\Delta t$: với nó, các dịch chuyển nhỏ hơn $\Delta t$ là không phân biệt được. Cắt thời gian thành các ô dài $\Delta t$, nhỏ hơn thời gian trơ, sao cho trong mỗi ô hoặc có một xung, hoặc không có xung nào (hình~\ref{fig:cells}). Với tần số trung bình $r$, xung rơi vào một ô với xác suất $p=r\Delta t$. Dòng xung mang nhiều thông tin nhất khi các ô độc lập, và khi đó mỗi ô cho entropy $-p\log_2p-(1-p)\log_2(1-p)\approx p\log_2(e/p)$ bit khi $p\ll1$. Chia cho $\Delta t$, ta được tốc độ truyền giới hạn và phần của nó trên mỗi xung~\cite{MacKay1952}:
\begin{empheq}[box=\keybox]{equation}
R_{\max} \;\approx\; r\,\log_2\frac{e}{r\,\Delta t}\ \ \text{bit/s},
\qquad
\frac{R_{\max}}{r} \;\approx\; \log_2\frac{e}{r\,\Delta t}\ \ \text{bit trên một xung}.
\label{eq:spikeentropy}
\end{empheq}
Với $r=10$ xung mỗi giây và $\Delta t=1\ms$, đó là khoảng tám bit trên một xung và tám mươi bit mỗi giây.

\begin{figure}[t]
\centering
\begin{tikzpicture}[>=Stealth,x=0.5cm,y=1cm]
% time cut into cells of width Delta t; a cell holds one impulse or none
\foreach \k in {0,...,23}{\draw[gray!55] (\k,0) rectangle (\k+1,0.7);}
\foreach \k in {2,7,8,15,21}{
  \fill[red!10] (\k,0) rectangle (\k+1,0.7);
  \draw[very thick,red!70!black] (\k+0.5,0.05) -- (\k+0.5,0.65);}
\foreach \k in {0,...,23}{
  \pgfmathtruncatemacro{\bit}{(\k==2)||(\k==7)||(\k==8)||(\k==15)||(\k==21)}
  \ifnum\bit=1 \def\bitcol{red!70!black}\else \def\bitcol{gray!60!black}\fi
  \node[font=\scriptsize,text=\bitcol] at (\k+0.5,-0.25) {\bit};}
\draw[<->,gray!60!black] (0,0.95) -- (1,0.95) node[midway,above,font=\scriptsize,black] {$\Delta t$};
\node[font=\small,anchor=east] at (-0.3,0.35) {xung};
\node[font=\small,anchor=east] at (-0.3,-0.25) {bit};
\draw[decorate,decoration={brace,amplitude=5pt,mirror},gray!60!black] (0,-0.55) -- (24,-0.55)
  node[midway,below=6pt,font=\scriptsize,black,align=center] {nơi nhận chỉ biết đếm: ``$n=5$'';\\ nơi nhận thấy được các ô: đúng $5$ ô nào trong $24$ ô --- tức là $\log_2\binom{24}{5}\approx15$ bit};
\end{tikzpicture}
\caption{Dung lượng~\eqref{eq:spikeentropy} đến từ đâu. Thời gian được cắt thành các ô dài $\Delta t$, trong mỗi ô hoặc có xung ($1$), hoặc không ($0$): dòng xung là một chuỗi nhị phân. Xung rơi vào một ô với xác suất $p=r\Delta t$. Ô càng nhỏ, chuỗi càng dài và càng nhiều bit; bộ đếm xung đánh mất gần như mọi thứ được ghi ở chỗ các số một nằm \emph{ở đâu}.}
\label{fig:cells}
\end{figure}

Số bit trên một xung chỉ phụ thuộc logarit vào độ chính xác thời gian: với độ chính xác $10\ms$ vẫn còn khoảng năm bit trên một xung. Nhưng nếu nơi nhận chỉ đếm xung trong một giây, thì với $r=10$ nó nhận được chưa tới hai bit cho cả giây (mục~\ref{sec:rate}).

\begin{keyblock}
\begin{proposition}[Dung lượng của kênh xung]
\label{prop:capacity}
Một dòng xung có tần số trung bình $r$, được đọc với độ chính xác thời gian $\Delta t$, mang không quá~\eqref{eq:spikeentropy}. Gần như toàn bộ dung lượng này nằm ở \emph{thời điểm} của các xung: nơi nhận chỉ đếm xung rút ra từ cùng dòng đó ít hơn hàng chục lần so với nơi nhận phân biệt được thời điểm của chúng với độ chính xác một mili giây.
\end{proposition}
\end{keyblock}

Đây là cận trên, không phải phép đo. Các phép đo trên nơ-ron cảm giác, mà kích thích đã biết, cho từ một đến vài bit trên một xung; trong những trường hợp tốt nhất, đó là tới một nửa giới hạn tính cho độ chính xác của chúng~\cite{Rieke1997}. Vậy ít nhất ở đầu vào của hệ thần kinh, thời điểm xung được sử dụng chứ không bị bỏ phí.

\section{Kênh có nhiễu}

Dung lượng~\eqref{eq:spikeentropy} được tính khi không có nhiễu. Về nơi nhiễu phát sinh, có ba sự thật đã được biết chắc chắn.

\emph{Bản thân bộ tạo xung là chính xác.} Nếu nhiều lần bơm vào một nơ-ron vỏ não, được lấy ra khỏi não cùng một mẩu mô, cùng một dòng điện thay đổi nhanh theo thời gian, các xung lặp lại với độ chính xác khoảng một mili giây~\cite{Mainen1995}. Nếu dòng không đổi, thời điểm xung trôi dạt từ lần này sang lần khác: thời điểm chính xác được tạo ra bởi những biến đổi nhanh của đầu vào chứ không phải bởi chính nơ-ron.

\emph{Khớp thần kinh không đáng tin cậy.} Một xung chạy tới khớp thần kinh \nt{GLUT} chỉ nhả một gói chất dẫn truyền với một xác suất $p$ nào đó. Ở các khớp thần kinh hồi hải mã, hơn một nửa số xung đơn giản là không tới được nơi nhận, và nguyên nhân chính là tính ngẫu nhiên của việc nhả chứ không phải lỗi dẫn truyền~\cite{AllenStevens1994}. Với người làm truyền thông, đây là một kênh xóa; vài khớp thần kinh giữa hai nơ-ron phần nào cứu vãn tình hình.

\emph{Trong vỏ não sống, dòng xung trông ngẫu nhiên.} Các khoảng giữa các xung phân tán gần giống như ở một dòng Poisson ngẫu nhiên, còn khi lặp lại cùng một kích thích, số xung trong một cửa sổ thay đổi sao cho phương sai của nó gần bằng giá trị trung bình~\cite{Softky1993}.

Một phần tử chính xác sinh ra dòng ngẫu nhiên như thế nào? Nơ-ron nhận hàng nghìn đầu vào, mỗi đầu vào không tin cậy, còn kích thích và ức chế cân bằng nhau, như ở chương~\ref{sec:gaba}. Điện áp màng dao động quanh ngưỡng, và chính các dao động đó quyết định thời điểm xung. Đó là nhiễu hay là thông điệp mà ta không biết đọc --- phần lớn vẫn là câu hỏi mở. Một phần ``nhiễu'' đã hóa ra là thông điệp: ở vỏ thị giác, một phần đáng kể độ phân tán bám theo cử động của chính con vật~\cite{Stringer2019}. Khó khăn ở đây mang tính nguyên tắc: với người không biết mã, một thông điệp được nén tốt không phân biệt được với một dòng ngẫu nhiên.

\section{Tần số: chậm nhưng chắc}
\label{sec:rate}

Mã đơn giản nhất là \emph{mã tần số}: con số được ghi trong số xung mà nơ-ron phát ra trong cửa sổ $T$. Cả sự xóa lẫn sự rung của thời điểm đều không làm hại mã này. Nhưng nó có cái giá. Nếu dòng là Poisson, số xung $n$ trong cửa sổ có trung bình $rT$ và độ phân tán $\sqrt{rT}$:
\begin{empheq}[box=\keybox]{equation}
\frac{\sigma_n}{\bar n} \;=\; \frac{1}{\sqrt{r\,T}} .
\label{eq:rateerror}
\end{empheq}
Để biết tần số với độ chính xác mười phần trăm, cần một trăm xung; với hai mươi xung mỗi giây, đó là năm giây. Hai mức kề nhau phân biệt được nếu chúng cách nhau vài độ phân tán, nên đến tần số $r$ trong thời gian $T$ phân biệt được khoảng $\sqrt{rT}$ mức: với $r=10$ và $T=1\,\text{s}$, đó là ba bốn mức, chưa tới hai bit.
Não không làm việc chậm như vậy. Người nhận ra con vật trong một bức ảnh được chiếu chớp nhoáng, và đáp ứng trong não xuất hiện sau khoảng $150\ms$~\cite{Thorpe1996}. Theo ước tính thô, trong thời gian đó tín hiệu đi qua khoảng một chục tầng xử lý nối tiếp, mỗi tầng $10$--$20\ms$. Với tần số thông thường của vỏ não, mỗi nơ-ron kịp phát không hoặc một xung ở mỗi tầng, và tần số của nó trong cửa sổ như vậy không xác định được. Có hai lối ra: lấy nhiều nơ-ron, hoặc nhìn vào thời điểm.

\emph{Nhiều nơ-ron.} Giả sử $N$ nơ-ron mang cùng một con số, và nơi nhận cộng các xung của chúng lại. Nếu nhiễu của chúng độc lập, trong~\eqref{eq:rateerror} thay cho $rT$ là $NrT$: một trăm nơ-ron với $r=20$ trong $15\ms$ cho ba mươi xung và độ chính xác khoảng hai mươi phần trăm. Không gian thay thế thời gian. Tuy nhiên, nhiễu của các nơ-ron vỏ não kề nhau không độc lập: hệ số tương quan của số xung ở một cặp láng giềng thường vào khoảng $0.1$~\cite{Zohary1994}. Nếu nó bằng $c$, phương sai của trung bình trên $N$ nơ-ron bằng
\begin{empheq}[box=\keybox]{equation}
\sigma^2_N \;=\; \sigma^2_1\,\frac{1+(N-1)\,c}{N}
\;\xrightarrow[N\to\infty]{}\; c\,\sigma^2_1 .
\label{eq:correlated}
\end{empheq}

\begin{keyblock}
\begin{proposition}[Giới hạn của phép lấy trung bình]
\label{prop:pooling}
Với nhiễu tương quan, lấy trung bình trên một nhóm làm giảm sai số không quá $1/\sqrt{c}$ lần. Với $c=0.1$, đó là ba lần, và lợi ích này đạt được ngay với vài chục nơ-ron; thêm nữa cũng vô ích.
\end{proposition}
\end{keyblock}

Không phải mọi tương quan đều có hại như nhau. Chỉ phần nhiễu chung \emph{giống chính tín hiệu} mới giới hạn độ chính xác, tức là phần làm dịch cả nhóm giống như sự thay đổi của đại lượng được đo sẽ làm dịch nó~\cite{MorenoBote2014}. Thực ra trong não có bao nhiêu nhiễu như vậy vẫn đang được làm rõ.

Còn một cách khác để dùng nhiều nơ-ron: ghi con số ở chỗ nơ-ron \emph{nào} đã phát xung, như trong bộ xác định hướng âm thanh (hình~\ref{fig:itd}). Đó là mã \emph{vị trí}, đáng tin cậy nhất trong các mã đã biết: ở nơ-ron cảm giác, số thứ tự của dây dẫn cho biết điểm nào trên da bị chạm hay âm thanh cao bao nhiêu, và điều này đã được xác lập nhiều lần. Nó tốn kém về số nơ-ron nhưng nhanh: một thông điệp chỉ tốn một độ trễ.

\section{Thời điểm: nhanh, nhưng tính từ đâu}

Lối ra thứ hai là ghi con số vào thời điểm xung. Lấy nơ-ron~\eqref{eq:glutneuron} và từ lúc $t=0$ cấp cho nó một đầu vào không đổi đủ để nâng điện áp lên mức $U$. Điện áp tăng theo $V=U\bigl(1-e^{-t/\tau}\bigr)$ và đạt ngưỡng $\theta$ vào thời điểm
\begin{empheq}[box=\keybox]{equation}
t_1 \;=\; \tau\,\ln\frac{U}{U-\theta}, \qquad U>\theta .
\label{eq:latency}
\end{empheq}
Đầu vào mạnh --- xung sớm, yếu --- xung muộn, dưới ngưỡng --- không có xung. Với $\tau=10\ms$, đầu vào $U=4\theta$ cho xung đầu tiên sau $2.9\ms$, $U=2\theta$ --- sau $6.9\ms$, $U=1.2\theta$ --- sau $17.9\ms$. Chỉ một xung duy nhất mang một con số.

Nhưng tính $t_1$ từ thời điểm nào? Không có đồng hồ, và nơi nhận không biết kích thích bắt đầu khi nào. Có ba câu trả lời.

\emph{Độ trễ tương đối.} Hai nơ-ron thấy cùng một kích thích, và hiệu độ trễ của chúng $t_1^A-t_1^B$ chỉ phụ thuộc vào đầu vào của chúng chứ không vào thời điểm bắt đầu. Các thời điểm được so sánh bởi bộ phát hiện trùng khớp có độ trễ (hình~\ref{fig:itd}). Nếu $N$ nơ-ron mỗi nơ-ron phát một xung, chính \emph{thứ tự} đến của chúng mang tới $\log_2N!$ bit: với mười nơ-ron là khoảng hai mươi hai bit, trong khi câu trả lời ``đã phát xung hay chưa'' cho không quá mười bit. Ở võng mạc, người ta đã chỉ ra rằng từ độ trễ tương đối của các xung đầu tiên của nơ-ron đầu ra có thể khôi phục hình ảnh nhanh và đáng tin cậy~\cite{Gollisch2008}.

\emph{Đợt bùng phát làm dấu mở đầu từ.} Một thay đổi đột ngột của kích thích khiến rất nhiều nơ-ron phát xung gần như đồng thời. Đợt bùng phát như vậy tự nó là dấu mốc bắt đầu, như khoảng lặng dài giữa các từ trong mã Morse, và nơi nhận tính độ trễ từ đó.

\emph{Nhịp cục bộ.} Ở chương~\ref{sec:gaba}, vòng \nt{GLUT}--\nt{GABA} sinh ra một nhịp cục bộ có chu kỳ $12$--$50\ms$. Đó không phải đồng hồ, nhưng trong một chu kỳ của nó, nơ-ron có đầu vào mạnh kịp phát xung sớm hơn nơ-ron có đầu vào yếu, và~\eqref{eq:latency} biến thành mã \emph{pha}: con số được ghi ở chỗ xung đến vào phần nào của chu kỳ. Mã như vậy đã được quan sát: các nơ-ron phụ trách một vị trí nhất định trong không gian phát xung ngày càng sớm hơn trong chu kỳ của một nhịp cục bộ chậm khi con vật đi qua vị trí ``của mình''~\cite{OKeefe1993}. Não dùng pha rộng rãi đến mức nào thì hiện chưa biết.

\section{Mã do nơi nhận chọn}

Trong một chuỗi xung có cả tần số, thời điểm và thứ tự. Cái nào trong số đó được đọc là do nơi nhận quyết định, và nó quyết định bằng một tham số --- hằng số thời gian $\tau$ của nó.

Lấy trung bình phương trình~\eqref{eq:glutneuron} theo thời gian. Nếu nơ-ron nhận các dòng độc lập với tần số $r_j$, điện áp trung bình và độ dao động tương đối của nó bằng
\begin{empheq}[box=\keybox]{equation}
\bar V \;=\; \tau\sum_j w_j\,r_j ,
\qquad
\frac{\sigma_V}{\bar V} \;\approx\; \frac{1}{\sqrt{2\,r_{\text{vào}}\,\tau}} ,
\label{eq:reader}
\end{empheq}
trong đó đẳng thức thứ hai được viết cho các trọng số bằng nhau, còn $r_{\text{vào}}$ là tổng tần số của mọi đầu vào. Khi $\tau$ lớn, điện áp hầu như không dao động và lặp lại tổng có trọng số của các tần số: nơi nhận là một \emph{bộ tích phân}, nó đọc tần số và quên thời điểm. Khi $\tau$ nhỏ, điện áp kịp giảm giữa các xung, và ngưỡng chỉ đạt được khi vài xung đến trong phạm vi cửa sổ~\eqref{eq:window}: nơi nhận là một \emph{bộ phát hiện trùng khớp}, nó đọc thời điểm (hình~\ref{fig:reader}). Một nghìn đầu vào, mỗi đầu vào năm xung mỗi giây, với $\tau=10\ms$ cho độ dao động khoảng mười phần trăm, gần như tích phân thuần túy. Nếu ức chế, như ở chương~\ref{sec:gaba}, rút ngắn $\tau$ xuống $2\ms$, độ dao động tăng lên hơn hai mươi phần trăm, và nơ-ron bắt đầu nhận ra các sự trùng khớp.

\begin{figure}[t]
\centering
\begin{tikzpicture}[>=Stealth,x=0.115cm,y=1cm]
% common input: the same spike train for both receivers
\foreach \s in {6,21,22,38,52,55.5,59,62.5,66,69.5,73,76.5,80,83.5,87,90.5,94,97.5}{\draw[thick,gray!40!black] (\s,5.25) -- (\s,5.65);}
\draw[gray!70!black] (0,5.25) -- (105,5.25);
\node[left,font=\small] at (-1,5.45) {đầu vào};
\node[above,font=\scriptsize,gray!40!black] at (13.5,5.65) {thưa};
\node[above,font=\scriptsize,gray!40!black] at (21.5,6.0) {cặp};
\node[above,font=\scriptsize,gray!40!black] at (75,5.65) {dày};
% axes and thresholds
\foreach \y/\th in {2.7/2.5,0/1.5}{
  \draw[gray!70!black] (0,\y) -- (106,\y) node[right,black] {\small $t$};
  \draw[densely dashed,gray!60!black] (0,\y+0.6*\th) -- (104,\y+0.6*\th) node[right,black,font=\small] {$\theta$};
}
\node[anchor=south west,font=\small] at (0,4.3) {$\tau=10\ms$: bộ tích phân};
\node[anchor=south east,font=\small] at (104,1.0) {$\tau=2\ms$: bộ phát hiện trùng khớp};
% integrator: tau=10 ms, theta=2.5w
\begin{scope}[yshift=2.7cm]
\foreach \a/\b/\v in {0/6/0.000,6/21/1.000,21/22/1.223,22/38/2.107,38/52/1.425,52/55.5/1.351,55.5/59/1.952,59/62.5/2.376,62.5/66/0.000,66/69.5/1.000,69.5/73/1.705,73/76.5/2.201,76.5/80/0.000,80/83.5/1.000,83.5/87/1.705,87/90.5/2.201,90.5/94/0.000,94/97.5/1.000,97.5/104/1.705}{\draw[very thick,blue!60!black] plot[domain=\a:\b,samples=14] (\x,{0.6*\v*exp(-(\x-\a)/10)});}
\foreach \s/\p/\q in {6/0.000/1.000,21/0.223/1.223,22/1.107/2.107,38/0.425/1.425,52/0.351/1.351,55.5/0.952/1.952,59/1.376/2.376,62.5/1.674/2.500,62.5/2.500/0.000,66/0.000/1.000,69.5/0.705/1.705,73/1.201/2.201,76.5/1.551/2.500,76.5/2.500/0.000,80/0.000/1.000,83.5/0.705/1.705,87/1.201/2.201,90.5/1.551/2.500,90.5/2.500/0.000,94/0.000/1.000,97.5/0.705/1.705}{\draw[very thick,blue!60!black] (\s,0.6*\p) -- (\s,0.6*\q);}
\foreach \s in {62.5,76.5,90.5}{\draw[->,thick,red!70!black] (\s,1.58) -- (\s,1.95);}
\end{scope}
% coincidence: tau=2 ms, theta=1.5w
\begin{scope}[yshift=0cm]
\foreach \a/\b/\v in {0/6/0.000,6/21/1.000,21/22/1.001,22/38/0.000,38/52/1.000,52/55.5/1.001,55.5/59/1.174,59/62.5/1.204,62.5/66/1.209,66/69.5/1.210,69.5/73/1.210,73/76.5/1.210,76.5/80/1.210,80/83.5/1.210,83.5/87/1.210,87/90.5/1.210,90.5/94/1.210,94/97.5/1.210,97.5/104/1.210}{\draw[very thick,blue!60!black] plot[domain=\a:\b,samples=14] (\x,{0.6*\v*exp(-(\x-\a)/2)});}
\foreach \s/\p/\q in {6/0.000/1.000,21/0.001/1.001,22/0.607/1.500,22/1.500/0.000,38/0.000/1.000,52/0.001/1.001,55.5/0.174/1.174,59/0.204/1.204,62.5/0.209/1.209,66/0.210/1.210,69.5/0.210/1.210,73/0.210/1.210,76.5/0.210/1.210,80/0.210/1.210,83.5/0.210/1.210,87/0.210/1.210,90.5/0.210/1.210,94/0.210/1.210,97.5/0.210/1.210}{\draw[very thick,blue!60!black] (\s,0.6*\p) -- (\s,0.6*\q);}
\foreach \s in {22}{\draw[->,thick,red!70!black] (\s,0.98) -- (\s,1.35);}
\end{scope}
\foreach \x in {0,50,100}{\draw[gray!60!black] (\x,-0.06) -- (\x,0.06) node[below,font=\scriptsize] {\x\,ms};}
\end{tikzpicture}
\caption{Cùng một đầu vào --- hai nơi nhận khác nhau. Mỗi xung nâng điện áp thêm $w$, sau đó điện áp rò đi, như trong mô hình nơ-ron~\eqref{eq:glutneuron}; mũi tên đỏ là xung của nơi nhận. Nơi nhận chậm ($\tau=10\ms$, $\theta=2.5w$) phát xung ở chỗ có \emph{nhiều} xung và không nhận ra cặp xung. Nơi nhận nhanh ($\tau=2\ms$, $\theta=1.5w$) chỉ phát xung với cặp xung trong phạm vi cửa sổ~\eqref{eq:window}, và bỏ qua dòng xung dày nhưng không trùng khớp.}
\label{fig:reader}
\end{figure}

Các phép đo cho thấy ở vỏ não đang hoạt động, độ dẫn của màng tăng lên vài lần so với lúc nghỉ~\cite{Destexhe2003}. Vậy ở đó $\tau$ ngắn hơn, và nơ-ron vỏ não ở chế độ làm việc gần với bộ phát hiện trùng khớp hơn là với bộ tích phân.

\begin{keyblock}
\begin{proposition}[Mã do nơi nhận quyết định]
\label{prop:reader}
Cùng một chuỗi xung, với nơi nhận chậm mang tần số, với nơi nhận nhanh mang thời điểm, với thể tích mà chất dẫn truyền được nhả vào mang mức nồng độ được làm trơn trong vài giây (chương~\ref{sec:serocomputer}). Mã nào được dùng không do nơi gửi quyết định, mà do hằng số thời gian của nơi nhận, và hằng số này được ức chế điều chỉnh.
\end{proposition}
\end{keyblock}
\section{Chấm và gạch: khớp thần kinh như một bộ lọc}

Cho đến giờ, khớp thần kinh ở ta là một dây dẫn có trọng số. Thực ra đáp ứng của nó phụ thuộc vào các xung trước đó, và điều này cho ngôn ngữ của nơ-ron ký hiệu thứ hai.

\emph{Cạn kiệt: khớp thần kinh truyền sự thay đổi.} Mỗi lần nhả tiêu tốn một phần $U$ lượng chất dẫn truyền sẵn sàng để nhả~\cite{Tsodyks1997}, còn lượng dự trữ phục hồi với hằng số thời gian $\tau_{\text{ph}}$ cỡ hàng trăm mili giây. Với tần số $r$, biên độ đáp ứng với một xung xác lập xấp xỉ ở mức $A(r)=A_0/(1+U r\tau_{\text{ph}})$, trong đó $A_0$ là đáp ứng của khớp thần kinh đã nghỉ ngơi. Dòng mà nơi nhận nhận được trong một đơn vị thời gian bằng
\begin{empheq}[box=\keybox]{equation}
r\,A(r) \;=\; \frac{A_0\,r}{1+U r\,\tau_{\text{ph}}}
\;\xrightarrow[r\to\infty]{}\; \frac{A_0}{U\tau_{\text{ph}}} .
\label{eq:depression}
\end{empheq}
Với $U=0.5$ và $\tau_{\text{ph}}=0.8\,\text{s}$, sự bão hòa bắt đầu ngay từ $1/(U\tau_{\text{ph}})=2.5$ xung mỗi giây. Nếu tần số tăng từ năm lên hai mươi, dòng xác lập chỉ tăng thêm một phần ba. Bù lại, những xung đầu tiên ở tần số mới đến khi lượng dự trữ vẫn như cũ, và dòng thoáng chốc vọt lên gấp bốn, rồi sau $\tau_{\text{ph}}/(1+Ur\tau_{\text{ph}})\approx90\ms$ thì giảm xuống (hình~\ref{fig:depression}).

Khớp thần kinh như vậy không truyền tần số mà truyền \emph{sự thay đổi} của nó, về thực chất là sự thay đổi tương đối $\Delta r/r$~\cite{Abbott1997depression}. Với người làm điện tử, đây là một bộ lọc thông cao đặt ngay trên dây dẫn.

\begin{figure}[t]
\centering
\begin{tikzpicture}[>=Stealth,x=12.5cm,y=1cm]
% input rate
\draw[->,gray!70!black] (0,2.6) -- (0.9,2.6) node[right,black] {\small $t$};
\draw[->,gray!70!black] (0,2.6) -- (0,4.3) node[above,black] {\small $r$};
\draw[very thick,blue!60!black] (0,2.9) -- (0.2,2.9) -- (0.2,3.8) -- (0.8,3.8);
\node[left,font=\scriptsize] at (0,2.9) {$5$};
\node[left,font=\scriptsize] at (0,3.8) {$20$};
\node[right,font=\small,blue!60!black] at (0.4,4.1) {tần số ở đầu vào khớp thần kinh};
% output current
\draw[->,gray!70!black] (0,0) -- (0.9,0) node[right,black] {\small $t$};
\draw[->,gray!70!black] (0,0) -- (0,2.45) node[left,black] {\small $rA$};
\draw[densely dashed,gray!60!black] (0,0.667) -- (0.8,0.667);
\draw[very thick,red!70!black] (0,0.5) -- (0.2,0.5) -- (0.2,2.0)
   plot[domain=0.2:0.8,samples=60] (\x,{0.667+(2.0-0.667)*exp(-(\x-0.2)/0.089)});
\node[left,font=\scriptsize] at (0,0.5) {$1.7$};
\node[left,font=\scriptsize] at (0,2.0) {$6.7$};
\node[right,font=\scriptsize] at (0.8,0.667) {$2.2$};
\node[right,font=\small,red!70!black] at (0.28,1.6) {dòng tới nơi nhận: vọt lên rồi giảm trong $\approx90\ms$};
\draw[gray!60!black] (0.2,-0.08) -- (0.2,0.08); \node[below,font=\scriptsize] at (0.2,0) {$0.2\,\text{s}$};
\draw[gray!60!black] (0.8,-0.08) -- (0.8,0.08); \node[below,font=\scriptsize] at (0.8,0) {$0.8\,\text{s}$};
\end{tikzpicture}
\caption{Khớp thần kinh cạn kiệt theo công thức~\eqref{eq:depression} với $U=0.5$, $\tau_{\text{ph}}=0.8\,\text{s}$; dòng tính bằng đơn vị $A_0$ mỗi giây, đường nét đứt là dòng xác lập: nó chỉ tăng từ $1.7$ lên $2.2$, còn tại thời điểm nhảy bậc --- lên tới $6.7$. Khớp thần kinh báo cho nơi nhận rằng tần số đã thay đổi, chứ không báo tần số bằng bao nhiêu.}
\label{fig:depression}
\end{figure}

\emph{Tạo thuận: chùm xung như một gạch.} Ở những khớp thần kinh khác, xác suất nhả thấp, nhưng sau mỗi xung nó tăng lên trong hàng chục mili giây. Một xung đơn lẻ qua khớp thần kinh như vậy thường bị mất. Còn một \emph{chùm xung} --- vài xung liên tiếp cách nhau vài mili giây --- gần như chắc chắn đi qua. Ngay cả khi không có tạo thuận, xác suất mất cả $k$ xung của chùm bằng
\begin{empheq}[box=\keybox]{equation}
P_{\text{mất}} \;=\; (1-p)^k ,
\label{eq:burst}
\end{empheq}
với $p=0.2$, đó là $0.8$ cho một xung đơn và $0.41$ cho chùm bốn xung; tạo thuận còn làm nó giảm mạnh hơn nữa. Như vậy nơ-ron có ký hiệu thứ hai: xung đơn là chấm, chùm xung là gạch. Khớp thần kinh cạn kiệt truyền tốt nhất xung đầu tiên sau khoảng lặng; khớp thần kinh tạo thuận cho chùm xung đi qua và dập các chấm đơn lẻ. Việc chùm xung được truyền đáng tin cậy hơn đã được đo; việc não thực sự dùng chúng như một ký hiệu riêng là một giả thuyết hợp lý~\cite{Lisman1997}.

\section{Nơ-ron \nt{GABA} nói thế nào}

Đông nhất trong số chúng ở vỏ não là các nơ-ron nhanh~\cite{Tremblay2016}. Xung của chúng ngắn hơn xung của nơ-ron \nt{GLUT}, chúng có thể phát xung đều và lâu với tần số hàng trăm mỗi giây mà gần như không mỏi. Khớp thần kinh của chúng đáng tin cậy hơn: liên kết với mỗi nơi nhận gồm nhiều điểm nhả, và xung gần như không bao giờ bị mất. Đầu ra của chúng đậu gần chỗ sinh ra xung của nơi nhận, nên chúng không điều khiển cách nơi nhận cộng các đầu vào mà điều khiển việc nó có phát xung không và khi nào.

Bản thân chúng nhận kích thích từ nhiều nơ-ron \nt{GLUT} lân cận và báo về tổng hoạt động xung quanh --- đúng thứ cần cho phép chia ở chương~\ref{sec:gaba}. Cuối cùng, các nơ-ron \nt{GABA} nhanh nối với nhau và dễ phát xung cùng lúc; chính chúng tạo ra nhịp cục bộ đã nói ở trên~\cite{Cardin2009}.

Theo ngôn ngữ của mã Morse, nơ-ron \nt{GABA} không phụ trách các chữ cái mà phụ trách \emph{khoảng lặng}. Chúng cắt dòng xung \nt{GLUT} liên tục thành những cửa sổ mà trong đó nơi nhận có thể phát xung, thu hẹp cửa sổ trùng khớp và làm nhỏ tiếng cả dòng khi nó trở nên quá ồn.

\begin{keyblock}
\begin{proposition}[Phân vai]
\label{prop:punctuation}
Nơ-ron \nt{GLUT} mang nội dung thông điệp: \emph{cái gì} và \emph{bao nhiêu}. Nơ-ron \nt{GABA} nhanh mang phần đánh dấu của nó: \emph{khi nào} được nói và \emph{to đến đâu}; một lớp nữa, bằng cách ức chế các nơ-ron ức chế, quyết định cổng mở \emph{cho ai}. Từng phần của sự phân vai này đã được đo; như một quy tắc chung, đây là giả thuyết làm việc.
\end{proposition}
\end{keyblock}

Có những nơ-ron \nt{GABA} khác, chậm hơn, có đầu ra đậu vào từng đầu vào riêng của nơi nhận. Chúng làm việc như phép cấm~\eqref{eq:veto} ở chương~\ref{sec:gaba}: đóng một dòng xung \nt{GLUT} mà không động đến các dòng khác.

Việc chia nơ-ron \nt{GABA} thành nhanh và chậm là một bức tranh cũ, và nó chưa đầy đủ. Hiện nay trong vỏ não người ta phân biệt ít nhất ba lớp lớn~\cite{Tremblay2016}, và lớp thứ ba được cấu tạo bất ngờ: nó không ức chế nơ-ron \nt{GLUT} mà ức chế các nơ-ron \nt{GABA} khác, trước hết là những nơ-ron đóng các đầu vào~\cite{Pfeffer2013}. Ức chế sự ức chế là phủ định kép. Nơ-ron như vậy \emph{mở} cổng mà không gửi cho nơi nhận một xung kích thích nào. Nó được bật bởi tín hiệu từ những vùng vỏ não khác và bởi các kênh quảng bá, nên nó hoạt động như đầu vào cho phép của cả một vùng mạng: chừng nào nó hoạt động, các lệnh cấm được gỡ bỏ, và dòng xung đi qua~\cite{Pi2013}. Ở phụ lục~\ref{sec:othermediators}, thủ thuật này được gọi là giải ức chế.
\section{Ngôn ngữ tiết kiệm}

Điều cuối cùng được biết chắc chắn về ngôn ngữ của nơ-ron là nó đắt. Một xung, cùng với công việc của khớp thần kinh do nó gây ra, tốn khoảng $10^9$ phân tử ATP (ước tính cho vỏ não loài gặm nhấm~\cite{Attwell2001}), còn một phân tử cho khoảng $10^{-19}$ joule. Vậy một xung có giá cỡ $E_{\text{xung}}\sim10^{-10}$ joule. Cả bộ não tiêu thụ khoảng $20$ oát, và nếu một nửa dành cho việc truyền tín hiệu, $P\sim10$ oát, thì tần số trung bình của một nơ-ron là
\begin{empheq}[box=\keybox]{equation}
\bar r \;\approx\; \frac{P}{N\,E_{\text{xung}}}
\;\sim\; \frac{10\ \text{W}}{8.6\cdot10^{10}\cdot10^{-10}\ \text{J}}
\;\approx\; 1\ \text{xung mỗi giây}.
\label{eq:energy}
\end{empheq}
Các ước tính chi tiết hơn cho vỏ não còn cho con số nhỏ hơn~\cite{Lennie2003}. Mã của não buộc phải \emph{thưa}: chỉ một phần nhỏ nơ-ron được phép làm việc nhanh.

Từ~\eqref{eq:spikeentropy} có thể thấy điều này không tệ lắm. Với độ chính xác $1\ms$, một xung của nơ-ron làm việc với tần số $100$ mỗi giây mang không quá năm bit, còn một xung của nơ-ron làm việc mỗi giây một lần --- tới mười một bit. Nếu phải trả tiền cho từng xung, nói thưa có lợi. Vỏ não quả thực đổi độ chính xác lấy năng lượng: khi thiếu thức ăn, nơ-ron giữ nguyên tần số xung nhưng chi ít hơn cho dòng qua khớp thần kinh, và đáp ứng của chúng trở nên kém chính xác hơn~\cite{Padamsey2022}.

Trong mã Morse, các chữ cái thường gặp thì ngắn: chữ cái thường gặp nhất trong văn bản tiếng Anh, E, chỉ là một chấm. Theo những gì đã biết, mã của nơ-ron tuân theo cùng quy tắc đó dưới một dạng khác: cái gì đã được dự kiến thì tốn ít xung~\cite{Barlow1961}. Với kích thích không đổi, nơ-ron dần giảm tần số, các cảm biến ở chương~\ref{sec:glutcomputer} đáp ứng với sự thay đổi chứ không phải với mức, còn nơ-ron \nt{DOPA} ở chương~\ref{sec:neurontypes} chỉ báo sai số dự đoán phần thưởng.

\begin{keyblock}
\begin{proposition}[Tiết kiệm]
\label{prop:economy}
Năng lượng giới hạn tần số trung bình của nơ-ron ở mức cỡ một xung mỗi giây. Vì thế ngôn ngữ của nơ-ron \nt{GLUT} là thưa và dùng xung cho tin tức: cho những gì đã thay đổi hoặc không được dự đoán. Giới hạn năng lượng đã được đo; việc mã của não được tinh chỉnh cho số bit trên mỗi joule lớn nhất là một giả thuyết có cơ sở vững chắc.
\end{proposition}
\end{keyblock}

\section{Tổng kết}

\begin{table}[t]
\centering
\footnotesize
\setlength{\tabcolsep}{4pt}
\renewcommand{\arraystretch}{1.15}
\begin{tabular}{>{\raggedright}p{2.6cm}>{\raggedright}p{2.3cm}>{\raggedright}p{3.0cm}>{\raggedright}p{2.0cm}>{\raggedright\arraybackslash}p{3.6cm}}
\toprule
Cách ghi & Mang gì & Ai đọc & Thời gian cho một thông điệp & Điều đã biết \\
\midrule
số thứ tự của nơ-ron đã phát xung (mã vị trí) & giá trị nào & mọi nơi nhận; bộ phát hiện trùng khớp & một độ trễ & đã xác lập ở nơ-ron cảm giác và trong việc xác định hướng âm thanh \\
tần số của một nơ-ron & độ lớn & bộ tích phân có $\tau$ lớn; thể tích (ch.~\ref{sec:serocomputer}) & hàng trăm ms -- vài giây & đã xác lập; quá chậm cho một tầng xử lý $10$--$20\ms$ \\
tần số của một nhóm & độ lớn & nơ-ron có hàng nghìn đầu vào & $10$--$50\ms$ & đã xác lập; lợi ích bị tương quan giới hạn~\eqref{eq:correlated} \\
thời điểm xung đầu tiên, thứ tự & độ lớn, thứ tự & bộ phát hiện trùng khớp có độ trễ & một độ trễ & đã xác lập ở đầu vào hệ thần kinh; ở vỏ não còn bàn cãi \\
pha trong nhịp cục bộ & độ lớn, vị trí & các nơ-ron có chung nhịp & chu kỳ $12$--$50\ms$ và chậm hơn & quan sát được ở một số vùng; vai trò chung là giả thuyết \\
chùm xung (``gạch'') & ``quan trọng'': chuyển phát tin cậy & khớp thần kinh tạo thuận & $\sim10\ms$ & độ tin cậy đã đo; ký hiệu riêng là giả thuyết \\
thay đổi tần số & tin tức & khớp thần kinh cạn kiệt & $\sim0.1\,\text{s}$ & đã xác lập \\
xung của nơ-ron \nt{GABA} nhanh & khi nào và to đến đâu & mọi nơ-ron \nt{GLUT} lân cận & $1$--$30\ms$ & nhịp và phép chia đã đo; ``dấu câu'' như một quy tắc là giả thuyết \\
\bottomrule
\end{tabular}
\caption{Những cách nơ-ron \nt{GLUT} và \nt{GABA} ghi thông điệp bằng xung. Cột bên phải tách điều đã đo khỏi điều giả định.}
\label{tab:code}
\end{table}

Bảng~\ref{tab:code} tập hợp những gì ta biết; từ đó cũng thấy được những gì ta chưa biết. Ta không biết vỏ não dùng thời điểm chính xác của xung đến mức nào, chứ không chỉ tần số của các nhóm. Ta không biết nơ-ron có ``từ'' hay không --- những tổ hợp xung ổn định của nhiều nơ-ron được đọc như một khối. Và ta không biết ngôn ngữ có giống nhau ở các phần khác nhau của não không. Mã Morse có thể học trong một buổi tối, vì bảng mã của nó đã biết. Bảng mã của ngôn ngữ nơ-ron còn phải được lập, và nhiều khả năng người lập sẽ là các kỹ sư truyền thông.

\section{Bài tập}

\begin{exercise}[\lvlA]
\label{ex:04:1}
\emph{Dung lượng bằng con số.} $\Delta t=1\ms$. (a)~Một nơ-ron có tần số $1$ xung mỗi giây cho bao nhiêu bit trên một joule, với giá một xung theo~\eqref{eq:energy}? (b)~Với $r=10$, nơi nhận chỉ đếm xung trong một giây rút ra ít hơn giới hạn~\eqref{eq:spikeentropy} bao nhiêu lần?
\end{exercise}

\begin{exercise}[\lvlB]
\label{ex:04:2}
\emph{Tần số tối ưu.} Nơ-ron tiêu thụ công suất $P_0+rE_{\text{xung}}$, trong đó $P_0$ là nửa còn lại của $20$~W chia cho mọi nơ-ron, $E_{\text{xung}}=10^{-10}$~J. Ở tần số $r$ nào, số bit trên một joule $R_{\max}(r)/(P_0+rE_{\text{xung}})$ theo~\eqref{eq:spikeentropy} với $\Delta t=1\ms$ là lớn nhất?
\end{exercise}

\begin{exercise}[\lvlB]
\label{ex:04:3}
\emph{Tương quan nào có hại.} $N=1000$ nơ-ron, phương sai $\sigma^2=1$, tương quan từng cặp $c=0.1$. Tính thông tin Fisher $\mathcal I=f'^{\mathsf T}\Sigma^{-1}f'$ ($\Sigma$ là ma trận hiệp phương sai) cho các độ dốc bằng nhau $f'=\mathbf 1$ và cho các độ dốc $\pm1$ với $\mathbf 1^{\mathsf T}f'=0$. Chúng khác nhau bao nhiêu lần?
\end{exercise}

\begin{exercise}[\lvlB]
\label{ex:04:4}
\emph{Mô phỏng: bộ phát hiện trùng khớp.} Nơ-ron~\eqref{eq:glutneuron} nhận $1000$ đầu vào Poisson, mỗi đầu vào $5$ xung mỗi giây, $w=1$; ngưỡng $\theta$ sao cho điện áp tự do vượt qua nó mỗi giây một lần. Bằng mô phỏng, tìm số đầu vào đồng thời $m$ làm nơ-ron phát xung với xác suất $1/2$ khi $\tau=10$ và $2\ms$.
\end{exercise}

\begin{exercise}[\lvlB]
\label{ex:04:5}
\emph{Thiết kế: bộ phát hiện thay đổi tương đối.} Chọn khớp thần kinh~\eqref{eq:depression} sao cho: dòng xác lập ở $40$ xung mỗi giây cao hơn không quá $10\,\%$ so với ở $10$, và sau bước nhảy lên $40$, dòng đạt mức mới với hằng số thời gian không quá $50\ms$. $\tau_{\text{ph}}$ nhỏ nhất bằng bao nhiêu và với $U$ nào?
\end{exercise}

\begin{exercise}[\lvlA]
\label{ex:04:6}
\emph{Thời điểm xung đầu và chùm xung.} (a)~Theo~\eqref{eq:latency}, tìm đầu vào làm xung đầu tiên đến đúng sau một hằng số thời gian. Nó phụ thuộc vào gì? (b)~Với xác suất nhả $p=0.2$, chùm xung cần bao nhiêu xung để xác suất mất tất cả dưới $5\,\%$?
\end{exercise}

\begin{exercise}[\lvlC]
\label{ex:04:7}
\emph{Nghiên cứu: bao nhiêu bit nằm ở cách sắp xếp.} Nơ-ron là các ô độc lập~\eqref{eq:spikeentropy}, $r=10$, $\Delta t=1\ms$. (a)~Tách entropy của một ``từ'' gồm $20$ ô thành entropy của số xung và phần còn lại. (b)~Mô phỏng ước lượng entropy theo tần suất các từ từ $N=30$ và $10^4$ bản ghi. Nó bị ước lượng thấp đi bao nhiêu?
\end{exercise}

\chapter{Các nơ-ron \nt{SERO} tính toán gì}
\label{sec:serocomputer}

\section{Kênh quảng bá đầu tiên}

Cho đến giờ, mọi nơ-ron trong sách đều liên lạc qua bus địa chỉ: \nt{GLUT} và \nt{GABA} gửi xung tới những đích nhận xác định, và ta đã biết một mạng như vậy tính được gì và nói bằng ngôn ngữ nào. Giờ ta chuyển sang bus thứ hai trong mục~\ref{sec:buses} --- bus quảng bá. Kênh đầu tiên của nó là \nt{SERO}. Như trước, ta chỉ cho phép mình dùng những gì có trong cơ thể sống.

Chương này giới thiệu ba bộ phận mà sau đó mọi kênh quảng bá đều sẽ dùng: một nơ-ron tự hoạt động khi có dòng nền đều, cho tới khi bị chặn lại; một sợi dây thực chất là một thể tích mô; và một nồng độ biến đổi chậm thay cho từng xung riêng lẻ. \nt{DOPA}, \nt{NORA} và \nt{ACET} ở các chương~\ref{sec:dofa}--\ref{sec:acet} sẽ được lắp từ chính những chi tiết này.

Nhìn từ góc độ kỹ sư, nơ-ron \nt{SERO} có bốn điểm khác biệt quan trọng.

\emph{Thứ nhất: bên nhận chọn dấu.} Serotonin có thụ thể của cả hai dấu, và quy tắc~\eqref{eq:dalesign} không áp dụng ở đây: dấu không do bên gửi quyết định mà do thụ thể của bên nhận (mệnh đề~\ref{prop:receptor})~\cite{BarnesSharp1999}.

\emph{Thứ hai: nơ-ron \nt{SERO} tự hoạt động nếu có nền đều.} Khi thức, nó phát xung đều đặn, vài lần mỗi giây, và không cần được kích cho từng xung: đó là một bộ tạo nhịp. Nhưng nó cần một dòng nền liên tục. Trong lát cắt não, nơi không có dòng này, phần lớn các tế bào như vậy im lặng; nhịp đều trở lại nếu đưa noradrenaline vào qua thụ thể $\alpha_1$, còn chặn các thụ thể này thì nhịp tắt~\cite{VandermaelenAghajanian1983}. Trong cơ thể, nền này đến từ \nt{NORA}. Về mặt mạch điện, đây là một bộ dao động cần nguồn nuôi: khi còn nguồn, nó tự chạy cho tới khi bị ức chế chặn lại.

\emph{Thứ ba: nó chậm.} Hầu hết thụ thể serotonin là thụ thể hướng chuyển hóa: chúng không tự mở kênh mà khởi động một chuỗi phản ứng bên trong tế bào. Đáp ứng vì thế chậm: dòng ức chế qua thụ thể chính tăng lên rồi giảm đi trong khoảng một giây~\cite{CourtneyFord2016} --- lâu gấp hàng chục lần đáp ứng của khớp thần kinh \nt{GLUT} nhanh.

\emph{Thứ tư: nó phóng chất dẫn truyền thần kinh không vào dây dẫn mà vào thể tích.} Ở những vùng mà đầu ra của nơ-ron \nt{SERO} đi tới, serotonin thường không thoát vào khe hẹp mà vào khoảng gian bào rồi lan ra xung quanh~\cite{Bunin1998}. Bên nhận cảm nhận nồng độ và không thể biết phân tử đến từ bên gửi nào.

Với những thang thời gian như vậy, từng xung riêng lẻ không còn quan trọng, nên ta sẽ mô tả mạng bằng tần số $r_j$ và nồng độ. Nồng độ $c$ của serotonin trong thể tích mà các nơ-ron $j$ phóng chất dẫn truyền vào tăng theo lượng phóng ra và giảm vì chất dẫn truyền bị bơm ngược trở lại:
\begin{empheq}[box=\keybox]{equation}
\tau_c\,\frac{\dd c}{\dd t} \;=\; -c \;+\; \kappa\sum_j r_j ,
\qquad
r_i \;=\; f_i\bigl(b_i + s_i\,g\,c\bigr), \quad s_i=\pm1 .
\label{eq:seroneuron}
\end{empheq}
Đây là thêm một cách đọc dòng xung, bổ sung cho mệnh đề~\ref{prop:reader}: thể tích không thấy từng xung, cũng không thấy thứ tự của chúng, mà chỉ thấy tần số lấy trung bình trong thời gian $\tau_c$. Phương trình thứ nhất chính là mạch RC giống như màng tế bào: $\tau_c$ là thời gian sống của chất dẫn truyền trong thể tích; nó chưa được đo trực tiếp, và ta coi nó cỡ một giây; $\kappa$ là lượng chất dẫn truyền do một xung tạo ra. Phương trình thứ hai mô tả bên nhận: $b_i$ là dòng vào riêng của nó, dương ở bộ tạo nhịp (gồm cả đầu vào nền đều từ \nt{NORA}); $g$ là độ nhạy của thụ thể; dấu $s_i$ do bên nhận chọn; $f_i$ là đặc tuyến truyền đạt không giảm.

\section{NOT bằng một nơ-ron}

Lấy một bộ tạo nhịp có thụ thể ức chế, $s=-1$. Khi xung quanh chưa có serotonin, nó chạy bằng dòng vào riêng $b$. Khi serotonin xuất hiện, dòng vào giảm đi $gc$, và nếu $gc>b$ thì nơ-ron im lặng. Đó là NOT: có đầu ra khi không có đầu vào. Nó chỉ tốn một nơ-ron (hình~\ref{fig:serogates}). Mệnh đề~\ref{prop:monotone} không áp dụng ở đây: nó đòi hỏi trọng số dương.

Nếu cùng bộ tạo nhịp đó chịu tác động của serotonin từ hai thể tích khác nhau, nó im lặng khi có serotonin ở ít nhất một thể tích. Đó là NOR, mà từ NOR có thể lắp mọi hàm logic. Mã hai dây của mạng \nt{GLUT} không cần đến ở đây: việc không có tín hiệu được tính bởi những nơ-ron tự chạy cho tới khi bị chặn lại.

\begin{figure}[t]
\centering
\begin{tikzpicture}[>=Stealth,
  nrn/.style={circle,draw,very thick,fill=blue!6,minimum size=0.95cm,inner sep=0pt},
  pace/.style={nrn,double,double distance=1.2pt},
  inh/.style={-{Bar[width=7pt]},thick,red!70!black}]
\node[pace] (N1) at (0,0) {};
\draw[inh] (-1.6,0) node[left,black] {$c$} -- (N1);
\draw[->,very thick,red!70!black] (N1) -- (1.3,0) node[right,black] {$\bar c$};
\node[below=0.55cm] at (0,0) {\small NOT};
\node[pace] (N2) at (5,0) {};
\draw[inh] (3.4,0.45) node[left,black] {$c_1$} -- (N2);
\draw[inh] (3.4,-0.45) node[left,black] {$c_2$} -- (N2);
\draw[->,very thick,red!70!black] (N2) -- (6.3,0) node[right,black] {$\overline{c_1\vee c_2}$};
\node[below=0.55cm] at (5,0) {\small NOR};
\end{tikzpicture}
\caption{Logic từ các nơ-ron \nt{SERO}. Vòng tròn kép là bộ tạo nhịp: khi có dòng nền đều, nó tự chạy cho tới khi bị ức chế. Đường có vạch ngang là thụ thể ức chế, $s=-1$; $c$, $c_1$, $c_2$ là nồng độ serotonin trong các thể tích quanh bên nhận. Bên trái là NOT, bên phải là NOR.}
\label{fig:serogates}
\end{figure}

\begin{keyblock}
\begin{proposition}[Tính đầy đủ của logic \nt{SERO}]
\label{prop:serocomplete}
Nếu mạng có các bộ tạo nhịp với thụ thể ức chế và lượng phóng vào các thể tích khác nhau không trộn lẫn, thì mạng~\eqref{eq:seroneuron} có thể tính mọi hàm logic, và mỗi bit được truyền bằng một dây.
\end{proposition}
\end{keyblock}

Về mặt logic, mạng \nt{SERO} mạnh hơn \nt{GLUT}, nhưng điều kiện ``lượng phóng vào các thể tích khác nhau không trộn lẫn'' không được thỏa mãn trong não (mục~\ref{sec:volume}).

\section{Hồi tiếp âm}

Các nơ-ron serotonin có thụ thể ức chế nằm ngay trên chính chúng~\cite{Sprouse1987}. Serotonin do chúng phóng ra quay trở lại chính chúng và hãm chúng lại. Với kỹ sư, đây là sơ đồ quen thuộc: bộ khuếch đại có hồi tiếp âm (hình~\ref{fig:serofeedback}).

Điều gì ở đây đã được đo, điều gì là mô hình. Ngay trong nhân raphe, nơi có các nơ-ron \nt{SERO}, serotonin ức chế các nơ-ron lân cận không qua thể tích chung mà theo kiểu điểm--điểm, qua khớp thần kinh: chất vận chuyển nhanh chóng dọn nó đi và không để nó tích tụ ngoài khe. Một lần phóng đơn lẻ chỉ gây một quãng ngừng ngắn trong chuỗi phát xung, còn tần số trung bình không đổi~\cite{CourtneyFord2016}. Vì vậy vòng lặp dưới đây, khép qua thể tích chung, là một mô hình: ta tính xem hồi tiếp như thế sẽ làm gì nếu nó hoạt động ở mức tần số trung bình.

\begin{figure}[t]
\centering
\begin{tikzpicture}[>=Stealth,
  nrn/.style={circle,draw,very thick,fill=blue!6,minimum size=0.95cm,inner sep=0pt},
  pace/.style={nrn,double,double distance=1.2pt},
  blk/.style={draw,very thick,rounded corners=2pt,fill=orange!10,minimum height=0.9cm,minimum width=2.2cm,align=center,font=\small},
  inh/.style={-{Bar[width=7pt]},thick,red!70!black}]
\node[pace] (N) at (0,0) {};
\node[blk] (V) at (3.6,0) {thể tích\\$\tau_c$};
\draw[->,very thick] (N) -- node[above] {\small $r$} (V);
\draw[->,very thick] (V) -- (6.0,0) node[right] {\small $c$ --- tới mọi bên nhận};
\draw[inh] (V.south) -- ++(0,-0.9) -| node[pos=0.25,below] {\small $-g\,c$} (N.south);
\draw[->,thick,blue!60!black] (-1.7,0) node[left,black] {\small $b$} -- (N);
\end{tikzpicture}
\caption{Nơ-ron \nt{SERO} có hồi tiếp qua chính chất dẫn truyền của nó: nồng độ $c$ đi tới mọi bên nhận và ức chế chính nơ-ron đó. Trong mô hình, đây là bộ điều chỉnh mức; vai trò như vậy của vòng lặp trong não là giả thuyết.}
\label{fig:serofeedback}
\end{figure}

Hãy tính xem vòng lặp này làm gì. Giả sử đặc tuyến truyền đạt tuyến tính, $f(u)=au$ khi $u>0$. Ở trạng thái cân bằng $c=\kappa r$ và $r=a(b-gc)$. Thế biểu thức này vào biểu thức kia, ta được
\begin{empheq}[box=\keybox]{equation}
r \;=\; \frac{a\,b}{1+G}, \qquad G \;=\; a\,g\,\kappa .
\label{eq:feedback}
\end{empheq}
Đại lượng $G$ là hệ số khuếch đại vòng. Đây chính là công thức của mọi bộ khuếch đại có hồi tiếp âm, và nó mang lại đúng hai lợi ích quen thuộc.

\emph{Bền với sai lệch.} Giả sử độ kích thích $a$ của nơ-ron thay đổi một tỷ lệ $\varepsilon$. Không có hồi tiếp, tần số cũng sẽ thay đổi đúng tỷ lệ đó. Có hồi tiếp, khi $G\gg1$, tần số $r\approx b/(g\kappa)$ gần như không phụ thuộc vào $a$: độ thay đổi của nó nhỏ hơn $1+G$ lần. Với $G=9$, sai lệch bị nén mười lần.

\emph{Nhanh hơn.} Thế $r=a(b-gc)$ vào phương trình thứ nhất của~\eqref{eq:seroneuron}: ta được $\tau_c\,\dd c/\dd t=-(1+G)\,c+\kappa ab$. Nồng độ tiến tới mức của nó với hằng số thời gian $\tau_c/(1+G)$ --- nhanh hơn $1+G$ lần so với khi không có hồi tiếp.

Đây là phép tính đầu tiên mà hệ \nt{SERO} có thể thực hiện: giữ mức serotonin chung trong não ở một giá trị đặt, như bộ ổn nhiệt giữ nhiệt độ. Đây là giả thuyết: trong thí nghiệm, sự tự ức chế cho đến nay chỉ thấy được dưới dạng những quãng ngừng ngắn, không làm thay đổi tần số trung bình.

\section{Từ xung thành mức}

Phép tính thứ hai ẩn trong phương trình thứ nhất của~\eqref{eq:seroneuron}. Nơ-ron phát ra từng xung riêng lẻ, còn bên nhận cảm nhận nồng độ, thứ làm trơn các xung trong thời gian $\tau_c$. Nếu tần số của nguồn thay đổi, nồng độ đi theo nó với độ trễ:
\begin{equation}
c(t) \;=\; \frac{\kappa}{\tau_c}\int_{-\infty}^{t} e^{-(t-t')/\tau_c}\,\sum_j r_j(t')\,\dd t' .
\label{eq:lowpass}
\end{equation}
Đây là trung bình trượt của hoạt động các nguồn trong vài $\tau_c$ gần nhất --- một bộ lọc thông thấp (hình~\ref{fig:lowpass}). Bộ chuyển đổi số--tương tự đơn giản nhất cũng cho xung đi qua mạch RC và biến tần số của chúng thành một mức. Hệ \nt{SERO} làm điều tương tự với hàng trăm nghìn nơ-ron cùng lúc.

Đại lượng này đồng thời cũng là bộ nhớ. Sau khi các nguồn đã im lặng, nồng độ còn duy trì thêm một khoảng thời gian cỡ $\tau_c$. Đó không phải là một bit, mà là một vết mờ dần.

\begin{figure}[t]
\centering
\begin{tikzpicture}[>=Stealth,x=1.45cm,y=1cm]
\draw[gray!70!black] (0,3.2) -- (8.1,3.2);
\node[left,font=\small] at (-0.05,3.4) {xung};
\node[above,font=\scriptsize,gray!40!black] at (1,3.65) {$2$ mỗi giây};
\node[above,font=\scriptsize,gray!40!black] at (3.5,3.65) {$8$ mỗi giây};
\node[above,font=\scriptsize,gray!40!black] at (6.5,3.65) {$2$ mỗi giây};
\draw[->,gray!70!black] (0,0) -- (8.3,0) node[right,black] {\small $t$};
\draw[->,gray!70!black] (0,0) -- (0,2.75) node[left,black] {\small $c$};
\foreach \s in {0.136,0.529,0.760,1.223,1.714,1.748,1.755,2.663,2.701,2.734,3.414,3.493,3.719,3.800,3.928,3.948,4.074,4.327,4.420,4.589,4.728,4.736,4.914,5.025,5.205,5.220,6.224,6.544,7.178}{\draw[thick,blue!60!black] (\s,3.2) -- (\s,3.6);}
\foreach \a/\b/\v in {0.000/0.136/0.000,0.136/0.529/1.000,0.529/0.760/1.675,0.760/1.223/2.330,1.223/1.714/2.466,1.714/1.748/2.509,1.748/1.755/3.425,1.755/2.663/4.402,2.663/2.701/2.775,2.701/2.734/3.672,2.734/3.414/4.553,3.414/3.493/3.306,3.493/3.719/4.055,3.719/3.800/4.235,3.800/3.928/4.905,3.928/3.948/5.316,3.948/4.074/6.211,4.074/4.327/6.476,4.327/4.420/6.028,4.420/4.589/6.493,4.589/4.728/6.483,4.728/4.736/6.642,4.736/4.914/7.589,4.914/5.025/7.351,5.025/5.205/7.579,5.205/5.220/7.331,5.220/6.224/8.221,6.224/6.544/4.012,6.544/7.178/3.914,7.178/8.000/3.076}{\draw[very thick,orange!85!black] plot[domain=\a:\b,samples=10] (\x,{0.25*\v*exp(-(\x-\a))});}
\foreach \s/\p/\q in {0.136/0.000/1.000,0.529/0.675/1.675,0.760/1.330/2.330,1.223/1.466/2.466,1.714/1.509/2.509,1.748/2.425/3.425,1.755/3.402/4.402,2.663/1.775/2.775,2.701/2.672/3.672,2.734/3.553/4.553,3.414/2.306/3.306,3.493/3.055/4.055,3.719/3.235/4.235,3.800/3.905/4.905,3.928/4.316/5.316,3.948/5.211/6.211,4.074/5.476/6.476,4.327/5.028/6.028,4.420/5.493/6.493,4.589/5.483/6.483,4.728/5.642/6.642,4.736/6.589/7.589,4.914/6.351/7.351,5.025/6.579/7.579,5.205/6.331/7.331,5.220/7.221/8.221,6.224/3.012/4.012,6.544/2.914/3.914,7.178/2.076/3.076}{\draw[very thick,orange!85!black] (\s,0.25*\p) -- (\s,0.25*\q);}
\draw[thick,densely dashed,black] plot[domain=0:2,samples=30] (\x,{0.25*2*(1-exp(-\x))})
  plot[domain=2:5,samples=40] (\x,{0.25*(8-6.271*exp(-(\x-2)))})
  plot[domain=5:8,samples=40] (\x,{0.25*(2+5.688*exp(-(\x-5)))});
\foreach \x in {0,2,5,8}{\draw[gray!60!black] (\x,-0.05) -- (\x,0.05) node[below=3pt,font=\scriptsize] {\x\,s};}
\end{tikzpicture}
\caption{Từ xung thành mức. Phía trên là các xung của nơ-ron \nt{SERO}: hai giây thưa, ba giây dày, rồi lại thưa. Phía dưới là nồng độ serotonin~\eqref{eq:lowpass} với $\tau_c=1\,\text{s}$. Đường nét đứt là kết quả tương ứng nếu các xung đến hoàn toàn đều đặn.}
\label{fig:lowpass}
\end{figure}

\section{Dây dẫn chính là thể tích}
\label{sec:volume}

Hãy quay lại điều kiện của mệnh đề~\ref{prop:serocomplete}. Serotonin do những bên gửi khác nhau phóng ra ở gần nhau cộng lại thành một nồng độ duy nhất.

\emph{Dây dẫn ở đây không phải là một sợi dẫn mà là một thể tích mô.} Hai tín hiệu chỉ còn tách biệt nếu chúng được phóng vào những vùng cách nhau xa hơn quãng đường chất dẫn truyền kịp lan tới. Trong thời gian $\tau$, một phân tử có hệ số khuếch tán $D$ đi được quãng đường
\begin{empheq}[box=\keybox]{equation}
\ell \;\sim\; \sqrt{D\,\tau}\, .
\label{eq:diffusionlength}
\end{empheq}
Độ ngoằn ngoèo của các khe gian bào làm khuếch tán của một phân tử nhỏ chậm đi khoảng $2.6$ lần~\cite{Sykova2008}, còn hệ số khuếch tán của chính serotonin trong mô chưa được đo; dựa vào kích thước phân tử, ta ước lượng nó cỡ vài đơn vị $10^{-6}$~cm$^2$/s. Nếu tín hiệu tồn tại $\tau\sim1$~s (cũng là ước lượng), thì $\ell\sim\sqrt{3\cdot10^{-6}}$~cm $\approx20\um$. Địa chỉ trong một mạng như vậy là \emph{vị trí}: bên nhận chỉ biết tín hiệu từ đâu tới qua chỗ chính nó đang nằm.

Các nơ-ron \nt{SERO} thật được tổ chức sao cho gần như không còn lại những địa chỉ riêng biệt ấy. Đầu ra của mỗi nơ-ron trải rộng trên những vùng não rất lớn: các nhánh của một tế bào vươn tới nhiều vùng nằm xa nhau~\cite{GagnonParent2014}. Số vị trí phóng trên mỗi tế bào ước tính khoảng $10^5$ (bảng~\ref{tab:types}); chúng chưa được đếm trực tiếp. ``Dây dẫn'' của các nơ-ron \nt{SERO} khác nhau chồng lên nhau và trộn lẫn. Dù vậy, chúng không hoàn toàn gộp thành một dây: các nơ-ron \nt{SERO} chia thành vài phân hệ, gửi đầu ra tới những vùng khác nhau và đáp ứng khác nhau với cùng một sự kiện~\cite{Ren2018}. Nhưng số phân hệ như vậy chỉ tính bằng đơn vị, chứ không phải hàng trăm nghìn.

\begin{keyblock}
\begin{proposition}[Số dây dẫn độc lập]
\label{prop:volumewires}
Trong mạng truyền dẫn thể tích, số tín hiệu độc lập không bị giới hạn bởi số nơ-ron mà bởi số vùng không chồng lấn, kích thước $\ell\sim\sqrt{D\tau}$, mà chúng phóng chất dẫn truyền vào. Nếu đầu ra của mỗi nơ-ron trải trên một thể tích lớn, toàn mạng chỉ truyền được vài biến chung.
\end{proposition}
\end{keyblock}

Vì vậy trong não không có gì để lắp các mạch logic của mệnh đề~\ref{prop:serocomplete}. Còn bộ điều chỉnh~\eqref{eq:feedback} và bộ lọc~\eqref{eq:lowpass} không cần những dây riêng: chúng chỉ cần một đại lượng chung. Bộ lọc suy trực tiếp từ việc phóng vào thể tích đã đo được ở các vùng đích~\cite{Bunin1998}; bộ điều chỉnh mức là giả thuyết.

\section{Tầm nhìn hoạch định}
\label{sec:horizon}

Một đại lượng chung, biến đổi chậm, có thể dùng vào việc gì? Một ứng viên tốt là tham số phải như nhau cho toàn mạng và không thay đổi nhanh hơn trong vòng vài giây hay vài phút. Ở chương~\ref{sec:neurontypes} đã có một tham số như vậy: hệ số $\gamma$ trong công thức sai số~\eqref{eq:rpe}, cho biết phần thưởng tương lai rẻ hơn phần thưởng tức thời bao nhiêu. Trong thời gian liên tục, vị trí của nó được thay bằng \emph{tầm nhìn} $\tau_\gamma$: phần thưởng $R$ phải chờ thời gian $D$ mới đến thì hiện có giá trị $R\,e^{-D/\tau_\gamma}$.

Tầm nhìn quyết định có đáng chờ hay không. Giả sử có thể lấy ngay một phần thưởng nhỏ $r_0$ hoặc chờ một phần thưởng lớn $R$. Chờ là có lợi khi
\begin{empheq}[box=\keybox]{equation}
D \;<\; \tau_\gamma\,\ln\frac{R}{r_0} .
\label{eq:patience}
\end{empheq}
Với $\tau_\gamma=2\,\text{s}$ và phần thưởng hoãn lại lớn gấp ba, đáng chờ tới $2.2\,\text{s}$; nếu tầm nhìn dài gấp đôi thì tới $4.4\,\text{s}$. Sự kiên nhẫn tỷ lệ với tầm nhìn.

Công thức~\eqref{eq:patience} dựa trên chiết khấu hàm mũ. Chiết khấu đo được ở độ trễ cỡ giây gần với hyperbol $R/(1+D/\tau_\gamma)$ hơn: ở khỉ, cả giá trị của phần thưởng hoãn lại trong lựa chọn lẫn đáp ứng của nơ-ron \nt{DOPA} với tín hiệu báo trước phần thưởng đều giảm như vậy~\cite{KobayashiSchultz2008}. Với hyperbol, điều kiện~\eqref{eq:patience} được thay bằng $D<\tau_\gamma\,(R/r_0-1)$: sự phụ thuộc vào tỷ số phần thưởng không còn là logarit mà là tuyến tính, và với cùng các con số, ngưỡng chờ dịch đi gần gấp đôi. Hyperbol thu được từ chính hàm mũ ấy nếu độ trễ là ngẫu nhiên (bài tập~\ref{ex:05:7}), và sự kiên nhẫn vẫn tỷ lệ với thang thời gian.

Theo giả thuyết, chính \nt{SERO} đặt tầm nhìn~\cite{Doya2002}. Nó có đủ mọi thứ cho vai trò này: một mức chung cho toàn mạng, thay đổi trong vài giây. Cũng có những thí nghiệm trực tiếp: nếu kích hoạt nhân tạo các nơ-ron serotonin, con vật chờ phần thưởng hoãn lại lâu hơn trước khi bỏ cuộc~\cite{Miyazaki2014}, và kiên trì lâu hơn trong việc kiếm phần thưởng ở nơi nó đã trở nên hiếm~\cite{Lottem2018}. Nồng độ serotonin chuyển thành $\tau_\gamma$ bên trong mạng như thế nào thì chưa rõ. Ở chương sau, tầm nhìn sẽ đi vào sai số mà \nt{DOPA} phát đi.

\section{Tổng kết}

\begin{table}[t]
\centering
\small
\begin{tabular}{>{\raggedright}p{3.6cm}>{\raggedright}p{4.3cm}>{\raggedright\arraybackslash}p{4.3cm}}
\toprule
& \nt{GLUT} & \nt{SERO} \\
\midrule
thời gian đáp ứng & mili giây & phần giây --- một giây \\
dấu tác động & chỉ $+$ & $+$ và $-$, do bên nhận chọn \\
khi không có đầu vào & im lặng & tự chạy khi có dòng nền đều \\
định địa chỉ & điểm --- điểm & thể tích, địa chỉ là vị trí \\
NOT & chỉ qua cảm biến ``tắt'' & một nơ-ron \\
bộ nhớ & hoạt động tự duy trì, tắt dần vì mỏi & nồng độ, mờ dần trong thời gian $\tau_c$ \\
phép tính chính & trùng khớp, độ trễ, logic, chuỗi & làm trơn; điều chỉnh mức và tầm nhìn $\tau_\gamma$ (giả thuyết) \\
\bottomrule
\end{tabular}
\caption{Những gì mạng chỉ gồm nơ-ron \nt{GLUT} và mạng chỉ gồm nơ-ron \nt{SERO} tính được.}
\label{tab:glutvssero}
\end{table}

Là một phần tử logic (bảng~\ref{tab:glutvssero}), nơ-ron \nt{SERO} phong phú hơn \nt{GLUT}, nhưng là một hệ liên lạc thì nó là kênh phát quảng bá: chậm và gần như không có địa chỉ. Vì vậy nó không cố làm bộ xử lý, mà làm trơn và phát đi vài đại lượng chung đã nói ở chương~\ref{sec:neurontypes} --- theo giả thuyết, trong đó có tầm nhìn hoạch định. Bộ tạo nhịp, thể tích và nồng độ chậm sẽ còn gặp lại ba lần nữa: ở \nt{DOPA}, \nt{NORA} và \nt{ACET}.

\section{Bài tập}

\begin{exercise}[\lvlA]
\label{ex:05:1}
\emph{Dây dẫn là thể tích.} Lấy $D=3\cdot10^{-6}$~cm$^2$/s cho serotonin (mục~\ref{sec:volume}). Theo~\eqref{eq:diffusionlength}, tìm $\ell$ cho tín hiệu tồn tại $1\,\text{s}$, và thời gian để chất dẫn truyền lan xa $5$~cm. Vậy tính quảng bá của \nt{SERO} nhờ khuếch tán hay nhờ dây dẫn phân nhánh với $\sim10^5$ vị trí phóng?
\end{exercise}

\begin{exercise}[\lvlA]
\label{ex:05:2}
\emph{Bộ điều chỉnh mức.} Chứng minh rằng trong~\eqref{eq:seroneuron} với $f(u)=au$, độ nhạy tương đối $S=(a/r)\,\partial r/\partial a$ bằng đúng $1/(1+G)$, chứ không chỉ khi $G\gg1$. Với $G=9$, độ kích thích $a$ tăng $20\,\%$. Không tuyến tính hóa, $r$ thay đổi bao nhiêu phần trăm?
\end{exercise}

\begin{exercise}[\lvlB]
\label{ex:05:3}
\emph{Thụ thể chậm trong vòng lặp.} Thụ thể \nt{SERO} đáp ứng có trễ: $r(t)=a\bigl(b-gc(t-T)\bigr)$, thể tích theo~\eqref{eq:seroneuron}. Chứng minh rằng khi $G>1$, vòng lặp chỉ ổn định nếu $T<T^*=\tau_c\arccos(-1/G)/\sqrt{G^2-1}$. Tìm $T^*$ với $\tau_c=1\,\text{s}$ cho $G=2$ và $G=9$. Với độ trễ hàng trăm mili giây, sai lệch có thể được nén bao nhiêu?
\end{exercise}

\begin{exercise}[\lvlB]
\label{ex:05:4}
\emph{Nhiễu bắn của mức.} Mỗi xung cộng vào $c$ theo~\eqref{eq:lowpass} một hàm mũ với $\tau_c=1\,\text{s}$. Tìm hệ số biến thiên của $c$ nếu cả $3\cdot10^5$ nơ-ron \nt{SERO} cùng phóng xung Poisson vào một thể tích, mỗi nơ-ron $2$ xung mỗi giây. Với một nơ-ron có tần số $4$ xung mỗi giây, $\tau_c$ bằng bao nhiêu thì $\mathrm{CV}=10\,\%$?
\end{exercise}

\begin{exercise}[\lvlB]
\label{ex:05:5}
\emph{Mô phỏng: hồi tiếp chống nhiễu.} Mô phỏng $N=50$ bộ tạo nhịp Poisson với tần số $r_i=a_i[b-gc]_+$, $a_i$ phân bố đều trong $[2.8,\,5.2]$, và thể tích~\eqref{eq:seroneuron} với $\kappa=1/N$, $\tau_c=1\,\text{s}$. Hồi tiếp ($b=1$, $g=2.25$) giảm CV của mức $c$ bao nhiêu lần so với $b=0.1$, $g=0$? Nó có san bằng tần số của các nơ-ron không?
\end{exercise}

\begin{exercise}[\lvlB]
\label{ex:05:6}
\emph{Thiết kế: XOR bằng logic \nt{SERO}.} Mỗi cổng là một bộ tạo nhịp phát $r_0$ nếu $b-g\sum_kc_k>0$ và $0$ nếu ngược lại, vào thể tích riêng $\tau_c\,\dd c/\dd t=-c+\kappa r$; nó tính NOR. Lắp XOR từ năm cổng như vậy và tìm thời gian xác lập của nó với $\tau_c=1\,\text{s}$ và $G'=g\kappa r_0/b=2$.
\end{exercise}

\begin{exercise}[\lvlB]
\label{ex:05:7}
\emph{Kiên nhẫn khi độ trễ chưa biết.} (a)~Con vật chịu chờ phần thưởng lớn gấp ba nếu thời gian chờ không quá $6\,\text{s}$. Tầm nhìn $\tau_\gamma$ tương ứng là bao nhiêu? (b)~Giả sử độ trễ $D$ phân bố mũ với trung bình $\bar D$. Chứng minh rằng chờ là có lợi khi $\bar D<\tau_\gamma(R/r_0-1)$, và so sánh với độ trễ cố định khi $\tau_\gamma=2\,\text{s}$, $R=3r_0$.
\end{exercise}

\begin{exercise}[\lvlC]
\label{ex:05:8}
\emph{Nồng độ đặt tầm nhìn thế nào.} Nơ-ron \nt{DOPA} nhận $V$ trực tiếp với trọng số $k_1$ và qua một trung gian \nt{GABA} làm trơn với $\tau_d=0.1\,\text{s}$, với trọng số $-k_2$; với $V$ chậm, tổng bằng $(k_1-k_2)V+k_2\tau_d\dot V$. Chọn $k_1$, $k_2$ khớp với~\eqref{eq:ctd} khi $\tau_\gamma=2\,\text{s}$. \nt{SERO} nhân $k_1$ với $m$: phải đổi $m$ bao nhiêu để tầm nhìn tăng gấp đôi?
\end{exercise}

\chapter{Mạng nơ-ron biết học: thêm \nt{DOPA}}
\chaptermark{Mạng biết học}
\label{sec:dofa}

\section{Luật chơi}

Trong mọi sơ đồ ở các chương trước, trọng số do kỹ sư đặt. Trong cơ thể không có kỹ sư. Trọng số phải tự thiết lập trong quá trình hoạt động, và sao cho mạng làm được điều gì đó hữu ích. Đó chính là học.

Ta thêm một quy tắc nữa vào các quy tắc cũ. Khớp thần kinh tự thay đổi trọng số của mình, và nó chỉ có thể biết những gì xảy ra \emph{ngay cạnh nó}: hoạt động của bên gửi $x$, hoạt động của bên nhận $y$ và những chất đến được chỗ nó. Không ai có thể gửi cho khớp thần kinh một hiệu chỉnh riêng.

Điều này loại bỏ phương pháp chính của mạng nơ-ron nhân tạo --- lan truyền ngược sai số. Ở đó, sai số đầu ra đi ngược qua mạng theo chính các trọng số ấy, trong một pha riêng sau lượt truyền thuận, và mỗi trọng số nhận được hiệu chỉnh riêng của mình. Muốn vậy cần có dây ngược lặp lại chính xác dây thuận, và một nhịp đồng hồ chung. Trong não chưa tìm thấy cả hai thứ này, ít nhất là ở dạng như trong mạng nhân tạo~\cite{Lillicrap2020,WhittingtonBogacz2019}.

Ta sẽ viết sự thay đổi trọng số như một quá trình liên tục chậm:
\begin{equation}
\frac{\dd w}{\dd t} \;=\; \eta\,\Phi(\text{những gì khớp thần kinh biết}),
\label{eq:plasticity}
\end{equation}
trong đó $\eta$ là tốc độ học, nhỏ tới mức trong thời gian một xung trọng số gần như không đổi. Cả chương này bàn về việc hàm $\Phi$ có thể như thế nào.

\section{Trọng số được lưu ở đâu}
\label{sec:weightstore}

Một khớp thần kinh ``thay đổi trọng số của mình'' là gì? Liên kết có hai phía: đầu tận cùng của dây bên gửi (phía \emph{trước khớp}) và vùng màng của bên nhận có các thụ thể (phía \emph{sau khớp}), ở giữa là một khe hẹp. Ở các liên kết \nt{GLUT}, màng bên nhận thường nằm trên một \emph{gai} --- một chỗ nhô nhỏ trên nhánh đầu vào của bên nhận. Tiện nhất là tách trọng số của liên kết thành ba thừa số~\cite{delCastillo1954}:
\begin{empheq}[box=\keybox]{equation}
w \;=\; N_s\,p\,A_1 .
\label{eq:quantal}
\end{empheq}
Ở đây $N_s$ là số vị trí phóng giữa hai nơ-ron, $p$ là xác suất để một xung phóng ra một gói chất dẫn truyền thần kinh tại một vị trí (chính là $p$ ở chương~\ref{sec:code}), $A_1$ là đáp ứng của bên nhận với một gói, tức là có bao nhiêu thụ thể bắt được nó. Thừa số $p$ nằm ở bên gửi, $A_1$ ở bên nhận, còn $N_s$ cần cả hai phía: vị trí phóng mới mọc ra, vị trí cũ biến mất, và điều này diễn ra suốt đời.

\emph{Bên nhận quyết định.} Ở một khớp thần kinh \nt{GLUT} điển hình, một trong các thụ thể glutamate chỉ dẫn dòng khi trùng hai điều kiện: glutamate đã đến từ bên gửi, và màng bên nhận đã bị kích thích~\cite{Nowak1984}. Đó là quy tắc Hebb được cài vào một phân tử --- bộ phát hiện trùng khớp ở chương~\ref{sec:glutcomputer}, đặt ngay tại đầu vào của chính khớp thần kinh. Khi được kích hoạt, thụ thể cho calci đi vào bên nhận, và calci khởi động sự thay đổi trọng số. Bên nhận là nơi duy nhất thấy được cùng lúc cả đầu vào lẫn đầu ra, nên quyết định được đưa ra ở đó.

\emph{Phía nào cũng có thể thực thi quyết định.} Dạng tăng cường dài hạn được nghiên cứu kỹ nhất diễn ra ở bên nhận: các thụ thể mới được gắn vào màng~\cite{Malinow2002}, và gai lớn lên~\cite{Matsuzaki2004} --- $A_1$ tăng. Nhưng bên nhận cũng có thể đổi $p$: nó phóng ra một chất đi ngược qua khe về phía bên gửi --- kênh ngược ở phụ lục~\ref{sec:othermediators} --- và bên gửi bắt đầu phóng ít chất dẫn truyền hơn~\cite{WilsonNicoll2001}. Còn những thay đổi trọng số ngắn hạn, suy giảm và tạo thuận ở chương~\ref{sec:code}, là những thay đổi của $p$ mà bên gửi tự điều khiển theo hoạt động gần đây của nó.

Với kỹ sư, thanh ghi trọng số được chia giữa hai vi mạch. Quy tắc học do bên nhận thực thi, còn kết quả thì nó có thể ghi cả ở phía mình lẫn ở bên gửi, qua kênh ngược. Vì vậy chữ ``cục bộ'' trong các quy tắc của chương này không có nghĩa là một phía của liên kết, mà là toàn bộ liên kết.

\section{Học không cần thầy}

Việc đơn giản nhất mà khớp thần kinh có thể làm là mạnh lên khi bên gửi và bên nhận cùng hoạt động:
\begin{equation}
\frac{\dd w}{\dd t} \;=\; \eta\,x\,y .
\label{eq:hebb}
\end{equation}
Đây là quy tắc Hebb~\cite{Hebb1949}. Nó cục bộ và không cần đồng hồ. Nhưng nó mắc cùng một vấn đề như mạng \nt{GLUT} ở chương~\ref{sec:glutcomputer}: hồi tiếp dương. Trọng số mạnh làm $y$ lớn hơn, $y$ lớn lại làm trọng số mạnh thêm, và trọng số tăng không giới hạn. Cách chữa là hồi tiếp âm theo trọng số: cho trọng số đồng thời giảm tỷ lệ với $y^2$. Với nơ-ron có các đầu vào $\mathbf x$ và đầu ra tuyến tính $y=\mathbf w\cdot\mathbf x$, ta được
\begin{empheq}[box=\keybox]{equation}
\frac{\dd\mathbf w}{\dd t} \;=\; \eta\,y\,\bigl(\mathbf x - y\,\mathbf w\bigr) .
\label{eq:oja}
\end{empheq}
Quy tắc vẫn cục bộ: khớp thần kinh biết đầu vào $x_i$ của nó, đầu ra $y$ và trọng số $w_i$.

\begin{keyblock}
\begin{proposition}[Quy tắc Hebb tìm ra gì]
\label{prop:oja}
Quy tắc~\eqref{eq:oja} đưa vectơ trọng số về độ dài đơn vị theo hướng mà các đầu vào biến thiên mạnh nhất --- về vectơ riêng chính của ma trận tương quan đầu vào $C=\langle\mathbf x\mathbf x^{\mathsf T}\rangle$~\cite{Oja1982}. Nơ-ron học cách xuất ra một đại lượng giàu thông tin nhất mà các đầu vào của nó có thể nén vào.
\end{proposition}
\end{keyblock}

Chứng minh. Lấy trung bình~\eqref{eq:oja} theo các đầu vào: $\dd\mathbf w/\dd t=\eta\bigl(C\mathbf w-(\mathbf w^{\mathsf T}C\mathbf w)\,\mathbf w\bigr)$. Các điểm dừng là những vectơ riêng của $C$ có độ dài đơn vị. Nếu $\mathbf w$ gần vectơ riêng $\mathbf e_k$ với trị riêng $\lambda_k$, một nhiễu loạn nhỏ theo vectơ khác $\mathbf e_j$ tăng như $e^{\eta(\lambda_j-\lambda_k)t}$. Chỉ điểm mà mọi $\lambda_j<\lambda_k$ là ổn định, tức là vectơ chính.

Cũng có một biến thể của quy tắc Hebb nhìn vào thời điểm của xung. Nếu xung của bên gửi đến $s$ mili giây \emph{trước} xung của bên nhận, trọng số tăng $A_+e^{-s/\tau_+}$; nếu nó đến $s$ mili giây \emph{sau}, trọng số giảm $A_-e^{-s/\tau_-}$, với $\tau_\pm$ cỡ hai mươi mili giây. Quy tắc này đã được đo trên khớp thần kinh của nơ-ron \nt{GLUT}~\cite{BiPoo1998}. Nó làm mạnh những liên kết \emph{có thể là nguyên nhân} của đáp ứng và làm yếu những liên kết đến muộn (hình~\ref{fig:stdp}). Cho một chuỗi kích thích chạy qua mạng nhiều lần, và trong mạng sẽ tự mọc ra các liên kết từ mỗi mắt xích tới mắt xích kế tiếp --- chuỗi ở chương~\ref{sec:glutcomputer}, được lắp mà không cần kỹ sư.

\begin{figure}[t]
\centering
\begin{tikzpicture}[>=Stealth,x=0.07cm,y=1.3cm]
\draw[->,gray!70!black] (-66,0) -- (68,0) node[right,black] {\small $s$};
\draw[->,gray!70!black] (0,-1.2) -- (0,1.35) node[above,black] {\small thay đổi trọng số};
\draw[very thick,blue!60!black] plot[domain=0.5:60,samples=50] (\x,{exp(-\x/20)});
\draw[very thick,red!70!black] plot[domain=-60:-0.5,samples=50] (\x,{-0.6*exp(\x/20)});
\draw[blue!60!black,thick] (0.5,0) -- (0.5,1);
\draw[red!70!black,thick] (-0.5,0) -- (-0.5,-0.6);
\foreach \x in {-40,-20,20,40}{\draw[gray!60!black] (\x,-0.04) -- (\x,0.04);}
\node[below,font=\scriptsize] at (20,-0.04) {$20\ms$};
\node[above,font=\scriptsize] at (-20,0.04) {$-20\ms$};
\node[align=left,font=\small,blue!60!black,anchor=west] at (22,0.95) {bên gửi trước:\\ liên kết mạnh lên};
\node[align=right,font=\small,red!70!black,anchor=east] at (-12,-0.95) {bên gửi sau:\\ liên kết yếu đi};
\end{tikzpicture}
\caption{Quy tắc Hebb theo thời gian. Trục ngang là $s=t_{\text{nhận}}-t_{\text{gửi}}$, xung của bên nhận chậm hơn xung của bên gửi bao nhiêu; $\tau_\pm=20\ms$, $A_-=0.6A_+$. Cửa sổ rộng vài chục mili giây --- cùng bậc với cửa sổ trùng khớp~\eqref{eq:window}.}
\label{fig:stdp}
\end{figure}

Nhưng mọi quy tắc trong mục này đều dạy điều gì \emph{thường xảy ra}, chứ không dạy điều gì \emph{có ích}. Một khớp thần kinh chỉ biết $x$ và $y$ thì không biết hành động đã mang lại nước ép hay cơn đau. Muốn học từ kết quả, khớp thần kinh cần một tín hiệu thứ ba.

\section{Thừa số thứ ba}

Tín hiệu thứ ba do \nt{DOPA} mang đến: các nơ-ron dopamine phát quảng bá một con số duy nhất --- sai số dự đoán phần thưởng $\delta$~\eqref{eq:rpe} (chương~\ref{sec:neurontypes}). Cho sự thay đổi trọng số tỷ lệ với tích của ba thừa số: hoạt động của bên gửi, hoạt động của bên nhận và $\delta$.

Khó khăn là phần thưởng không đến lúc mạng đang chọn hành động, mà sau đó hàng trăm mili giây hay vài giây. Khi $\delta$ đến, tích $xy$ đã bằng không từ lâu, và khớp thần kinh cần một bộ nhớ về việc nó đã tham gia. Bộ nhớ đơn giản nhất là mạch RC quen thuộc:
\begin{empheq}[box=\keybox]{equation}
\tau_e\,\frac{\dd e}{\dd t} \;=\; -e \;+\; x\,y ,
\qquad
\frac{\dd w}{\dd t} \;=\; \eta\,\delta(t)\,e(t) .
\label{eq:threefactor}
\end{empheq}
Đại lượng $e$ được gọi là \emph{vết đủ điều kiện}~\cite{Fremaux2016}. Hoạt động đồng thời để lại trong khớp thần kinh một dấu mờ dần trong thời gian $\tau_e$. Nếu dopamine đến trong thời gian đó, trọng số thay đổi theo dấu của $\delta$; nếu không, dấu biến mất không để lại gì (hình~\ref{fig:trace}). Trên khớp thần kinh \nt{GLUT} đã đo được: dopamine đến trong vòng một--hai giây sau hoạt động đồng thời biến nó thành sự tăng cường liên kết, còn dopamine đến muộn hơn thì không~\cite{Yagishita2014}. Cửa sổ dẻo dài cỡ giây cũng đã được tìm thấy ở nơi tín hiệu thứ ba không phải là dopamine, mà là sự kích thích mạnh của chính bên nhận~\cite{Bittner2017}.

\begin{figure}[t]
\centering
\begin{tikzpicture}[>=Stealth,x=7cm,y=1cm]
\fill[orange!12] (0.7,-0.2) rectangle (0.8,5.2);
% rows: x*y, delta, traces, weights
\foreach \y/\lab in {4.4/{$x\,y$},3.1/{$\delta$},1.4/{vết $e$},0/{trọng số $w$}}{
  \draw[gray!70!black] (0,\y) -- (1.42,\y);
  \node[left,font=\small] at (-0.02,\y+0.3) {\lab};}
\draw[very thick,blue!60!black] (0,4.4) -- (0,4.9) -- (0.2,4.9) -- (0.2,4.4);
\node[right,font=\scriptsize,blue!60!black] at (0.21,4.8) {bên gửi và bên nhận cùng hoạt động};
\draw[very thick,orange!85!black] (0,3.1) -- (0.7,3.1) -- (0.7,3.7) -- (0.8,3.7) -- (0.8,3.1) -- (1.4,3.1);
\node[right,font=\scriptsize,orange!85!black] at (0.81,3.6) {dopamine đến};
% traces, each in units of its own maximum
\draw[very thick,blue!60!black] plot[domain=0:0.2,samples=20] (\x,{1.4+1.1*(1-exp(-\x))/(1-exp(-0.2))})
  plot[domain=0.2:1.4,samples=60] (\x,{1.4+1.1*exp(-(\x-0.2))});
\draw[thick,densely dashed,gray!50!black] plot[domain=0:0.2,samples=30] (\x,{1.4+1.1*(1-exp(-\x/0.05))/(1-exp(-4))})
  plot[domain=0.2:1.4,samples=80] (\x,{1.4+1.1*exp(-(\x-0.2)/0.05)});
\node[right,font=\scriptsize,blue!60!black] at (1.0,2.15) {$\tau_e=1\,\text{s}$};
\node[right,font=\scriptsize,gray!40!black] at (0.3,1.65) {$\tau_e=50\ms$};
% weights
\draw[very thick,blue!60!black] (0,0.25) -- (0.7,0.25) -- (0.8,0.8) -- (1.4,0.8) node[right,font=\scriptsize] {$\tau_e=1\,\text{s}$: trọng số tăng};
\draw[thick,densely dashed,gray!50!black] (0,0.2) -- (1.4,0.2) node[right,font=\scriptsize,gray!40!black] {$\tau_e=50\ms$: không đổi};
% time axis marks
\foreach \x/\l in {0/0,0.2/0.2,0.7/0.7,1.4/1.4}{\draw[gray!60!black] (\x,-0.05) -- (\x,0.05) node[below=3pt,font=\scriptsize] {\l\,s};}
\draw[<->,thick,gray!50!black] (0.2,5.35) -- node[above,font=\scriptsize,black] {phần thưởng trễ $0.5\,\text{s}$} (0.7,5.35);
\end{tikzpicture}
\caption{Vết theo quy tắc~\eqref{eq:threefactor}. Hoạt động đồng thời kéo dài $0.2\,\text{s}$, dopamine đến nửa giây sau khi nó kết thúc. Các vết được vẽ theo tỷ lệ với giá trị cực đại của chính nó. Vết chậm ($\tau_e=1\,\text{s}$) lúc đó vẫn còn, vết nhanh ($\tau_e=50\ms$) đã tan từ lâu.}
\label{fig:trace}
\end{figure}

\begin{keyblock}
\begin{proposition}[Học bằng tín hiệu quảng bá]
\label{prop:threefactor}
Quy tắc~\eqref{eq:threefactor} chỉ thay đổi những khớp thần kinh đã tham gia hoạt động của mạng trong khoảng $\tau_e$ vừa qua, theo hướng do tín hiệu chung $\delta$ quy định. Tín hiệu quảng bá cùng với vết cục bộ cho ra một hiệu chỉnh có địa chỉ. Độ trễ giữa hành động và kết quả là chấp nhận được chừng nào nó không lớn hơn $\tau_e$.
\end{proposition}
\end{keyblock}

\section{Vì sao điều này hiệu quả: nhiễu là sự tìm kiếm}
\label{sec:dofagradient}

Vì sao quy tắc~\eqref{eq:threefactor} dẫn tới kết quả tốt hơn? Câu trả lời nằm ở nhiễu trong chương~\ref{sec:code}. Giả sử đầu ra của nơ-ron là $y=f(u)+\xi$, trong đó $u=\sum_jw_jx_j$, còn $\xi$ là dao động ngẫu nhiên có trung bình bằng không và phương sai $\sigma^2$. Phần thưởng $R$ phụ thuộc vào $y$. Ta lấy làm vết không phải $x_jy$ mà là phần ngẫu nhiên của nó $x_j\xi$, và làm thừa số thứ ba là sai số $\delta=R-\bar R$, trong đó $\bar R$ là phần thưởng kỳ vọng. Khi nhiễu nhỏ $R\approx R\bigl(f(u)\bigr)+R'\,\xi$, nên
\begin{empheq}[box=\keybox]{equation}
\bigl\langle \delta\,\xi\,x_j \bigr\rangle \;=\; \sigma^2\,R'\,x_j \;=\; \frac{\sigma^2}{f'(u)}\,\frac{\partial\langle R\rangle}{\partial w_j} .
\label{eq:perturbation}
\end{empheq}
Bước trung bình đi theo gradient của phần thưởng kỳ vọng~\cite{Williams1992}. Không ai tính đạo hàm: mạng lệch đi một cách ngẫu nhiên, dopamine báo xem kết quả có tốt hơn không, và các khớp thần kinh đã tham gia dịch về phía có lợi (hình~\ref{fig:perturbation}). Nhiễu, thứ ở chương~\ref{sec:code} cản trở việc truyền tin, ở đây trở thành sự tìm kiếm.

\begin{figure}[t]
\centering
\begin{tikzpicture}[>=Stealth,x=2cm,y=3cm]
\draw[->,gray!70!black] (0,0) -- (4.2,0) node[below left,font=\small,black] {đầu ra $y$};
\draw[->,gray!70!black] (0,0) -- (0,1.2) node[left,font=\small,black] {phần thưởng $R$};
% reward hill
\draw[very thick,blue!60!black] plot[domain=0:4,samples=60] (\x,{max(0,1-((\x-3)/2.2)^2)});
\node[font=\scriptsize,blue!60!black,anchor=south] at (3,1.02) {đầu ra tốt nhất};
% noise around f(u)
\draw[thick,gray!60!black] plot[domain=0.6:2.4,samples=50] (\x,{0.28*exp(-((\x-1.5)/0.3)^2/2)});
\draw[densely dashed,gray!60!black] (1.5,0) -- (1.5,0.535);
\node[below,font=\small] at (1.5,0) {$f(u)$};
\node[font=\scriptsize,gray!50!black] at (1.5,-0.2) {dao động $\xi$};
% expected reward
\draw[densely dotted,gray!60!black] (0.5,0.535) node[left,font=\scriptsize,black] {$\bar R$} -- (2.1,0.535);
% two random deviations
\fill[green!45!black] (2.1,0.833) circle (0.035);
\draw[very thick,green!45!black] (2.1,0.535) -- (2.1,0.833);
\node[font=\scriptsize,green!45!black,anchor=west,align=left] at (2.18,0.6) {$\xi>0$: tốt hơn,\\ $\delta>0$};
\fill[red!70!black] (0.9,0.089) circle (0.035);
\draw[very thick,red!70!black] (0.9,0.535) -- (0.9,0.089);
\node[font=\scriptsize,red!70!black,anchor=east,align=right] at (0.83,0.27) {$\xi<0$: tệ hơn,\\ $\delta<0$};
% resulting step
\draw[->,very thick,orange!85!black] (1.55,0.4) -- (1.95,0.4) node[right,font=\scriptsize,black] {bước};
\end{tikzpicture}
\caption{Nhiễu là sự tìm kiếm~\eqref{eq:perturbation}. Đầu ra của nơ-ron dao động quanh $f(u)$. Độ lệch ngẫu nhiên sang phải (màu xanh lá) cho nhiều hơn kỳ vọng, $\delta>0$; sang trái (màu đỏ) cho ít hơn, $\delta<0$. Trong cả hai trường hợp tích $\delta\,\xi$ đều dương, và trọng số dịch sang phải, về phía phần thưởng tăng. Trung bình, bước tỷ lệ với độ dốc $R'$: trên đỉnh đồi, độ lệch về cả hai phía đều tệ như nhau, và các bước triệt tiêu lẫn nhau.}
\label{fig:perturbation}
\end{figure}

\begin{keyblock}
\begin{proposition}[Đi theo gradient mà không cần dây ngược]
\label{prop:gradient}
Nếu trong mạng có những độ lệch hoạt động ngẫu nhiên, còn tín hiệu quảng bá bằng sai số dự đoán phần thưởng và trong mỗi bối cảnh có trung bình bằng không, thì quy tắc~\eqref{eq:threefactor} trung bình thay đổi trọng số theo gradient của phần thưởng kỳ vọng. Không cần dây ngược hay một pha học riêng.
\end{proposition}
\end{keyblock}

Giờ thì rõ vì sao dopamine không báo về phần thưởng, mà về \emph{sai số} dự đoán nó. Nếu thừa số thứ ba là chính phần thưởng $R\ge0$, trung bình trong~\eqref{eq:perturbation} sẽ không đổi, nhưng sẽ nảy ra hai vấn đề. Thứ nhất, phần hữu ích bị cộng thêm số hạng $\bar R\,\xi x_j$: trung bình bằng không, nhưng là nhiễu mạnh. Thứ hai, và tệ hơn, khớp thần kinh thật không biết tách phần ngẫu nhiên khỏi phần không ngẫu nhiên trong vết $x_jy$. Với các đầu vào cho trước, $x_jf(u)$ là cùng một con số từ lượt thử này sang lượt thử khác; nếu $\delta$ trong bối cảnh này trung bình bằng không, phần đó của vết không thay đổi gì, còn khi $\delta\ge0$ nó hết lần này tới lần khác làm mạnh mọi thứ mạng đã làm --- đúng cũng như sai. Vì vậy $\bar R$ phải là kỳ vọng \emph{trong bối cảnh cụ thể}, chứ không phải trung bình cả đời. Tính nó là việc của bộ phê bình (critic), sẽ nói ở dưới.

Chúng tôi đã kiểm tra điều này trên mô hình. Hai nhóm nơ-ron \nt{GLUT} chọn một trong hai hành động qua ức chế lẫn nhau, như ở hình~\ref{fig:wta}; trọng số từ tín hiệu ``bắt đầu'' tới các nhóm lúc đầu bằng nhau, và nhiễu quyết định nhóm nào thắng. Hành động $A$ mang lại nước ép với xác suất $0.8$, hành động $B$ với xác suất $0.2$, và nước ép đến nửa giây sau lựa chọn. Với $\tau_e=1\,\text{s}$ và $\delta=R-\bar R$, tỷ lệ chọn $A$ qua năm trăm lượt thử tăng từ $0.65$--$0.8$ trong trăm lượt đầu lên $0.96$--$1.0$ trong trăm lượt cuối, còn trọng số vẫn nằm trong khoảng $0.85$--$1.35$. Với $\tau_e=50\ms$, nhỏ hơn độ trễ phần thưởng, mạng không học được gì: tỷ lệ chọn $A$ vẫn quanh một nửa. Với $\delta=R$ không trừ kỳ vọng, mạng cũng bắt đầu chọn $A$, nhưng trọng số của bên thắng tăng một nghìn lần trong cùng năm trăm lượt thử và vẫn tiếp tục tăng; còn với cùng $\delta=R$ và tốc độ học lớn gấp mười, trong một trên ba lần chạy mạng đã cố định vĩnh viễn lựa chọn ngẫu nhiên đầu tiên của nó --- lựa chọn tệ hơn.

\section{Cái giá của một dây dẫn}

Tín hiệu $\delta$ chỉ có một cho tất cả, còn nhiễu mà nó đánh giá là tổng các dao động của cả $N$ nơ-ron ảnh hưởng tới kết quả. Mỗi khớp thần kinh thấy trong $\delta$ phần hữu ích của mình và nhiễu của $N-1$ nơ-ron lân cận. Để tách phần của mình, phải lấy trung bình qua nhiều lượt thử, và thời gian học tăng tỷ lệ với $N$~\cite{Werfel2005}. Lan truyền ngược không phải trả cái giá này.

\begin{keyblock}
\begin{proposition}[Cái giá của quảng bá]
\label{prop:broadcastcost}
Thời gian học bằng một tín hiệu chung tăng tỷ lệ với số nơ-ron có nhiễu ảnh hưởng tới kết quả. Cách học này phù hợp với những mạch nhỏ chọn giữa vài phương án, nhưng không phù hợp để tinh chỉnh hàng tỷ trọng số.
\end{proposition}
\end{keyblock}

Có vẻ như não phân công đúng như vậy: đầu ra của nơ-ron \nt{DOPA} đi trước hết tới những vùng chọn hành động (chương~\ref{sec:neurontypes}), còn mạng \nt{GLUT} khổng lồ của vỏ não, nơi chứa phần lớn áp đảo các trọng số, chủ yếu học bằng các quy tắc tại chỗ, không có thầy bên ngoài. Trước đây người ta quy chúng về các quy tắc Hebb~\eqref{eq:oja}. Nay ngày càng có nhiều dữ liệu cho thấy vỏ não có mục tiêu riêng --- \emph{dự đoán} đầu vào tiếp theo của mình, và khi đó sai số được tính tại chỗ, là hiệu giữa cái được dự đoán và cái thực sự đến~\cite{Keller2018}: người thầy được cài sẵn trong chính mạng. Những cửa sổ dẻo dài cỡ giây, nơi tín hiệu thứ ba là sự kích thích mạnh của chính bên nhận~\cite{Bittner2017}, cho thấy tín hiệu sai số tại chỗ ở khớp thần kinh thực sự tồn tại. Việc đây có phải là quy tắc chung của vỏ não hay không là giả thuyết.

\section{Sai số từ đâu ra: bộ phê bình trong thời gian liên tục}

Còn lại câu hỏi chính các nơ-ron \nt{DOPA} tính $\delta$ thế nào. Trong các chương trình học tăng cường, sai số dự đoán được tính theo bước $t,t+1,\dots$ Trong cơ thể không có bước, nên ta viết nó trong thời gian liên tục~\cite{Doya2000}. Cho $r(t)$ là dòng phần thưởng, còn $V(t)$ là kỳ vọng của toàn bộ phần thưởng tương lai, trong đó phần càng xa càng được coi nhẹ:
\begin{equation}
V(t) \;=\; \int_t^{\infty} e^{-(s-t)/\tau_\gamma}\,r(s)\,\dd s .
\end{equation}
Lấy đạo hàm, ta được $\dd V/\dd t=V/\tau_\gamma-r$. Phần dư của đẳng thức này chính là sai số dự đoán:
\begin{empheq}[box=\keybox]{equation}
\delta(t) \;=\; r(t) \;+\; \frac{\dd V}{\dd t} \;-\; \frac{V}{\tau_\gamma} .
\label{eq:ctd}
\end{empheq}
Hằng số $\tau_\gamma$ là tầm nhìn hoạch định, bản tương tự liên tục của hệ số $\gamma$; theo giả thuyết, nó do \nt{SERO} đặt (mục~\ref{sec:horizon}), và khi đó~\eqref{eq:patience} là quy tắc cho biết đáng chờ bao lâu. Hãy kiểm tra ba trường hợp (hình~\ref{fig:ctdcases}). Bóng đèn báo có nước ép bật sáng: $V$ tăng vọt, $\dd V/\dd t$ lớn, có đỉnh vọt. Nước ép đã hứa đến: $r>0$, nhưng kỳ vọng cũng cạn ngay lúc đó, $\dd V/\dd t\approx-r$, và không có đỉnh vọt. Nước ép không đến: kỳ vọng biến mất mà không có phần thưởng, $\dd V/\dd t<0$, có hõm sụt.

\begin{figure}[t]
\centering
\begin{tikzpicture}[>=Stealth,x=1.15cm,y=1cm]
\foreach \col/\ttl/\lampon/\juice in {0/{nước ép bất ngờ}/0/1, 1/{nước ép đã hứa}/1/1, 2/{nước ép không đến}/1/0}{
 \begin{scope}[xshift=\col*4.6cm]
  \node[font=\small] at (1.5,3.55) {\ttl};
  \foreach \yy in {2.4,1.2,0}{\draw[gray!60!black] (0,\yy) -- (3.1,\yy);}
  % reward r
  \ifnum\juice=1
    \fill[green!45!black] (1.95,2.4) rectangle (2.05,2.85);
    \pic[scale=0.55] at (2.0,3.1) {juice};
  \else
    \pic[scale=0.55] at (2.0,3.1) {nojuice};
  \fi
  % expectation V and error delta
  \ifnum\lampon=1
    \pic[scale=0.55] at (1.0,3.1) {lamp};
    \draw[very thick,blue!60!black] (0,1.2) -- (1,1.2) -- (1,1.45)
      plot[domain=1:2,samples=20] (\x,{1.2+0.25*exp((\x-1)/1.5)}) -- (2,{1.2+0.25*exp(1/1.5)}) -- (2,1.2) -- (3.1,1.2);
    \draw[very thick,red!70!black] (0,0) -- (0.95,0) -- (0.95,0.5) -- (1.05,0.5) -- (1.05,0) -- (3.1,0);
    \ifnum\juice=0
      \draw[very thick,red!70!black] (1.95,0) -- (1.95,-0.45) -- (2.05,-0.45) -- (2.05,0);
      \node[font=\scriptsize,red!70!black,anchor=west] at (2.1,-0.35) {hõm sụt};
    \else
      \node[font=\scriptsize,gray!50!black,anchor=north,align=center] at (2.0,-0.08) {$r$ và mức giảm $V$\\ triệt tiêu nhau};
    \fi
  \else
    \draw[very thick,blue!60!black] (0,1.2) -- (3.1,1.2);
    \draw[very thick,red!70!black] (0,0) -- (1.95,0) -- (1.95,0.5) -- (2.05,0.5) -- (2.05,0) -- (3.1,0);
  \fi
  \ifnum\col=0
    \node[font=\small,anchor=east] at (-0.1,2.55) {$r$};
    \node[font=\small,anchor=east] at (-0.1,1.35) {$V$};
    \node[font=\small,anchor=east] at (-0.1,0.1) {$\delta$};
  \fi
 \end{scope}}
\end{tikzpicture}
\caption{Ba trường hợp của sai số~\eqref{eq:ctd}; từ trên xuống --- phần thưởng $r$, kỳ vọng $V$ và sai số $\delta$. Bên trái không có kỳ vọng, và nước ép hoàn toàn là bất ngờ. Ở giữa, bóng đèn làm $V$ tăng vọt: đó là một bước nhảy của $\dd V/\dd t$, nên đỉnh vọt của $\delta$ rơi vào lúc đèn sáng. Sau đó $V$ tăng đúng như tầm nhìn đòi hỏi, và $\delta=0$; lúc nước ép đến, phần thưởng $r$ và mức giảm của $V$ triệt tiêu nhau. Bên phải $V$ giảm mà không có phần thưởng --- hõm sụt. Đây vẫn là đỉnh vọt, sự dịch chuyển và hõm sụt như ở hình~\ref{fig:rpe}, nhưng giờ ta thấy chúng từ đâu ra.}
\label{fig:ctdcases}
\end{figure}

Từ những sơ đồ đã có, cần gì để tính~\eqref{eq:ctd}? Ba thứ.

\emph{Chính ước lượng $V$.} Nó do một nhóm nơ-ron \nt{GLUT} --- bộ phê bình --- xuất ra, và trọng số của nhóm này học theo cùng quy tắc~\eqref{eq:threefactor} với cùng $\delta$: nếu $\delta>0$ thì kỳ vọng đã bị đánh giá thấp, và các liên kết hoạt động ngay trước đó được làm mạnh lên.

\emph{Đạo hàm $\dd V/\dd t$.} Bất kỳ sơ đồ nào đáp ứng với sự thay đổi chứ không phải với mức đều tính được nó: kích thích trực tiếp cùng với ức chế có trễ qua một trung gian \nt{GABA}, như ở bộ phát hiện hướng trong chương~\ref{sec:gaba}, hoặc một khớp thần kinh suy giảm~\eqref{eq:depression} ở chương~\ref{sec:code}.

\emph{Dấu trong một dây.} Sai số có thể dương hoặc âm, còn tần số xung chỉ dương. Lời giải giống như ở nơ-ron \nt{SERO} trong chương~\ref{sec:serocomputer}: nơ-ron \nt{DOPA} là một bộ tạo nhịp. Vài xung đều đặn mỗi giây của nó nghĩa là $\delta=0$, một chùm xung nghĩa là $\delta>0$, một quãng ngừng nghĩa là $\delta<0$. Sai số âm có ít chỗ hơn: tần số không thể xuống dưới không. Ở nơ-ron \nt{DOPA} thật, hõm sụt quả thực nông hơn đỉnh vọt~\cite{Bayer2005}.

Việc dopamine hành xử như $\delta$ đã được đo. Não tính $\dd V/\dd t$ chính xác thế nào là giả thuyết.

\section{Bộ hành động và bộ phê bình}

Hãy ghép tất cả lại (hình~\ref{fig:actorcritic}). Các nơ-ron bối cảnh kích thích những nhóm hành động chọn ra bên thắng qua \nt{GABA} (mệnh đề~\ref{prop:wta}) --- đó là \emph{bộ hành động} (actor) --- và bộ phê bình xuất ra $V$~\cite{SuttonBarto2018}. Nơ-ron \nt{DOPA} nhận phần thưởng $r$ và $V$, tính~\eqref{eq:ctd} và phát $\delta$ tới mọi khớp thần kinh của bộ hành động và bộ phê bình.

\begin{figure}[t]
\centering
\begin{tikzpicture}[>=Stealth,
  nrn/.style={circle,draw,very thick,fill=blue!6,minimum size=0.9cm,inner sep=0pt,font=\small},
  gab/.style={circle,draw,very thick,fill=red!10,minimum size=0.6cm,inner sep=0pt},
  dofa/.style={circle,draw,very thick,double,double distance=1.2pt,fill=orange!15,minimum size=1.35cm,inner sep=0pt,font=\scriptsize},
  inh/.style={-{Bar[width=7pt]},thick,red!70!black},
  bc/.style={->,thick,densely dashed,orange!85!black}]
\node[nrn] (S1) at (0,2.4) {$x_1$};
\node[nrn] (S2) at (0,1.2) {$x_2$};
\node[nrn] (S3) at (0,0) {$x_3$};
\node[left=0.15cm,align=right,font=\small] at (-0.5,1.2) {bối cảnh};
\node[nrn] (A) at (4,2.6) {$A$};
\node[nrn] (B) at (4,1.0) {$B$};
\node[gab] (G) at (5.2,1.8) {};
\draw[->,thick] (A) -- (G); \draw[->,thick] (B) -- (G);
\draw[inh] (G) to[bend left=35] (A);
\draw[inh] (G) to[bend right=35] (B);
\node[above,font=\small] at (4.6,3.05) {bộ hành động: lựa chọn};
\node[nrn] (V) at (4,-1.1) {$V$};
\node[below,font=\small] at (4,-1.6) {bộ phê bình};
\foreach \s in {S1,S2,S3}{\foreach \t in {A,B,V}{\draw[->,gray!60!black] (\s) -- (\t);}}
\node[dofa] (D) at (8.0,-1.1) {\nt{DOPA}};
\draw[->,thick] (V) -- node[above,font=\scriptsize] {$V,\ \dd V/\dd t$} (D);
\draw[->,thick,green!45!black] (10.0,-1.1) node[right,black,font=\small] {phần thưởng $r$} -- (D);
\pic[scale=1.1] at (10.6,-0.5) {juice};
\draw[bc] (D.north) -- (8.0,4.0) -- node[above,font=\small,black] {$\delta$ --- tới mọi khớp thần kinh của cả hai bộ} (2.2,4.0) -- (2.2,3.0);
\draw[bc] (D.south) -- (8.0,-2.4) -- (2.2,-2.4) -- (2.2,-0.6);
\end{tikzpicture}
\caption{Mạng học cách chọn hành động. Mũi tên xám là các khớp thần kinh \nt{GLUT} với trọng số thay đổi được; vòng tròn đỏ là nơ-ron \nt{GABA}; vòng tròn kép là bộ tạo nhịp \nt{DOPA}, tính $\delta$ theo~\eqref{eq:ctd}; mũi tên nét đứt là việc phát quảng bá $\delta$.}
\label{fig:actorcritic}
\end{figure}

Dấu tác động của chất dẫn truyền do bên nhận chọn (mệnh đề~\ref{prop:receptor}), và trong vùng chọn hành động, một phần nơ-ron mang thụ thể dopamine thuộc một họ, phần còn lại thuộc họ khác. Trong thí nghiệm trên lát cắt, dopamine cùng với hoạt động đồng thời làm mạnh khớp thần kinh ở nhóm thứ nhất và làm yếu khớp thần kinh ở nhóm thứ hai~\cite{Shen2008}; trong mô hình, điều này có nghĩa là ở nhóm thứ hai, dấu của $\delta$ trong~\eqref{eq:threefactor} bị đảo. Nhóm thứ nhất học điều gì \emph{nên làm}, nhóm thứ hai học điều gì \emph{không nên làm}.

\section{Tổng kết}

\begin{table}[t]
\centering
\footnotesize
\setlength{\tabcolsep}{4pt}
\renewcommand{\arraystretch}{1.15}
\begin{tabular}{>{\raggedright}p{2.6cm}>{\raggedright}p{2.7cm}>{\raggedright}p{3.1cm}>{\raggedright\arraybackslash}p{5.0cm}}
\toprule
Quy tắc & Khớp thần kinh biết gì & Mạng học gì & Điều đã biết \\
\midrule
Hebb theo tần số~\eqref{eq:oja} & $x$, $y$, trọng số của nó & nén đầu vào: các hướng chính & sự tăng cường khi cùng hoạt động đã được đo; cơ chế giới hạn trọng số đang được nghiên cứu \\
Hebb theo thời gian & thứ tự xung $x$ và $y$ trong phạm vi $\sim20\ms$ & liên kết nhân quả, chuỗi & đã đo trên khớp thần kinh \nt{GLUT} \\
ba thừa số~\eqref{eq:threefactor} & $x$, $y$, vết $e$, tín hiệu chung $\delta$ & những hành động mang lại phần thưởng & tác động của dopamine trong vài giây sau hoạt động đã được đo; vai trò cơ chế chính để học lựa chọn là giả thuyết có cơ sở vững \\
bộ phê bình~\eqref{eq:ctd} & như trên & kỳ vọng phần thưởng $V$ & sự giống nhau giữa dopamine và $\delta$ đã được đo; sơ đồ tính $\dd V/\dd t$ là giả thuyết \\
lan truyền ngược sai số & hiệu chỉnh riêng cho từng trọng số & mọi thứ, nhưng cần dây ngược và nhịp đồng hồ & chưa tìm thấy trong não ở dạng này; người ta tìm các dạng xấp xỉ của nó \\
\bottomrule
\end{tabular}
\caption{Các quy tắc học trong mạng gồm nơ-ron \nt{GLUT}, \nt{GABA} và \nt{DOPA}.}
\label{tab:learning}
\end{table}

Các quy tắc học được tóm tắt trong bảng~\ref{tab:learning}. \nt{GLUT} mang dữ liệu, \nt{GABA} điều khiển luồng, còn \nt{DOPA} quyết định giữ lại những thay đổi nào trong mạng.

\section{Bài tập}

\begin{exercise}[\lvlA]
\label{ex:06:1}
\emph{Vết sống bao lâu.} Hoạt động đồng thời $xy=1$ kéo dài $T_a=0.2\,\text{s}$, bắt đầu từ $e=0$, còn dopamine đến $d=0.5\,\text{s}$ sau khi nó kết thúc. (a)~Tại thời điểm đó, vết~\eqref{eq:threefactor} với $\tau_e=1\,\text{s}$ lớn hơn bao nhiêu lần so với $\tau_e=50\ms$? (b)~Với $\tau_e$ nào thì vết lớn nhất?
\end{exercise}

\begin{exercise}[\lvlA]
\label{ex:06:2}
\emph{Độ trôi của quy tắc Hebb.} Bên gửi và bên nhận phát các dòng Poisson độc lập, mỗi bên $10$ xung mỗi giây, và mỗi cặp xung thay đổi trọng số theo cửa sổ ở hình~\ref{fig:stdp} với $\tau_\pm=20\ms$, $A_-=0.6A_+$. Sau một giờ hoạt động hoàn toàn ngẫu nhiên, trọng số tăng bao nhiêu $A_+$? Với $\tau_-$ nào thì độ trôi biến mất?
\end{exercise}

\begin{exercise}[\lvlB]
\label{ex:06:3}
\emph{Tương quan, không phải hiệp phương sai.} Quy tắc của mệnh đề~\ref{prop:oja} tìm vectơ riêng chính của $C=\langle\mathbf x\mathbf x^{\mathsf T}\rangle$. Tần số là dương: $\mathbf x=\mathbf m+\mathbf n$, $\mathbf m=(\mu,0)$, hiệp phương sai của nhiễu $\operatorname{diag}(s_1^2,s_2^2)$. Với điều kiện nào nơ-ron học được $\mathbf e_1$? Nó học được gì khi $\mu=1$, $s_1=0.1$, $s_2=0.5$?
\end{exercise}

\begin{exercise}[\lvlA]
\label{ex:06:4}
\emph{Tầm nhìn trong đỉnh vọt.} Bóng đèn bật sáng vào một thời điểm không đoán trước được, nước ép lượng $R$ đến sau $L=1\,\text{s}$. Chứng minh rằng với kỳ vọng đã học theo~\eqref{eq:ctd}, giữa lúc đèn sáng và lúc có nước ép $\delta\equiv0$, và tìm diện tích đỉnh vọt lúc đèn sáng với $\tau_\gamma=2\,\text{s}$ và $\tau_\gamma=4\,\text{s}$.
\end{exercise}

\begin{exercise}[\lvlB]
\label{ex:06:5}
\emph{Cái giá của quảng bá từ đâu ra.} $N$ nơ-ron dao động độc lập, $\xi_i\sim\mathcal N(0,\sigma^2)$, sai số $\delta=a\sum_i\xi_i$, khớp thần kinh $j$ ước lượng gradient là $\delta\,\xi_jx_j$. Tìm tỷ số tín hiệu trên nhiễu của một lượt thử. Một nơ-ron học xong sau $100$ lượt thử, mỗi lượt $5\,\text{s}$: việc học mất bao lâu khi $N=10^4$ và $N=10^{10}$?
\end{exercise}

\begin{exercise}[\lvlB]
\label{ex:06:6}
\emph{Thiết kế: sơ đồ tính $\delta$.} Nơ-ron \nt{DOPA} nhận $r$, nhận $V$ với trọng số $k_1$ và một bản sao của $V$ được làm trơn với $\tau_d$, với trọng số $-k_2$. (a)~Chọn $k_1$, $k_2$ để với $V$ chậm, đầu ra là~\eqref{eq:ctd}. (b)~Với $V=V_0e^{t/\tau_\gamma}$, $\delta$ thật bằng $0$. Với $\tau_d$ nào sai số giả của sơ đồ không vượt quá $5\,\%$ của $V/\tau_\gamma$, nếu $\tau_\gamma=2\,\text{s}$?
\end{exercise}

\begin{exercise}[\lvlB]
\label{ex:06:7}
\emph{Mô phỏng: độ trễ phần thưởng tới hạn.} Thêm độ trễ phần thưởng $d$ vào \texttt{run()} trong một bản sao của \texttt{models/\allowbreak dofa\_learning.py}. Tìm tỷ lệ chọn $A$ trong trăm lượt cuối của $500$ lượt thử khi $d=0.5\,\text{s}$ và $\tau_e=0.05$, $0.2$, $1$, $4\,\text{s}$, và khi $\tau_e=1\,\text{s}$, $d=3\,\text{s}$. Khi phần vết còn lại tới lúc có phần thưởng (bài tập~\ref{ex:06:1}) là bao nhiêu thì việc học biến mất?
\end{exercise}

\begin{exercise}[\lvlC]
\label{ex:06:8}
\emph{Có thể né cái giá của quảng bá không.} $N$ nơ-ron: $y_i=w_i+\xi_i$, $\xi_i\sim\mathcal N(0,\sigma^2)$, phần thưởng $R=-\sum_i(y_i-w_i^*)^2$, quy tắc $\Delta w_i=\eta\,(R-\langle R\rangle)\,\xi_i$. Với $\eta$ tốt nhất, sau bao nhiêu lượt thử sai số $\sum_i(w_i-w_i^*)^2$ giảm $e$ lần? Còn nếu chỉ $m$ nơ-ron ngẫu nhiên trong số $N$ dao động? Tìm kiếm thưa có giúp ích không?
\end{exercise}

\chapter{Các lý thuyết hiện đại khác về học dựa trên dopamine}
\chaptermark{Các lý thuyết khác về \nt{DOPA}}
\label{sec:dofaalt}

\section{Thực ra thứ gì chạy trên dây}

Ở chương~\ref{sec:dofa}, hệ \nt{DOPA} của ta là một dây duy nhất, trên đó chạy một con số --- sai số dự đoán phần thưởng $\delta$~\eqref{eq:ctd}. Bức tranh này giải thích được nhiều điều, nhưng dopamine được đo càng chính xác thì người ta càng tìm thấy nhiều thứ không khớp với nó. Một số nơ-ron đáp ứng với điều khó chịu như với phần thưởng; nồng độ dopamine tăng trong lúc con vật chạy tới đích, dù không có gì bất ngờ xảy ra; các nơ-ron \nt{DOPA} khác nhau hành xử khác nhau.

Với kỹ sư, mọi phát hiện này quy về một câu hỏi: \emph{tín hiệu trên bus quảng bá có định dạng gì?} Một con số hay một vectơ? Một tín hiệu trên dây hay nhiều tín hiệu tách theo tần số? Nó nói về tương lai, như $\delta$, hay về quá khứ? Dưới đây là sáu câu trả lời hiện đại; với mỗi câu, ta sẽ tách điều đã đo khỏi điều được giả định.

\section{Không phải một con số mà là một phân bố}

Sai số~\eqref{eq:ctd} dạy bộ phê bình kỳ vọng phần thưởng \emph{trung bình}. Nhưng nửa giọt nước ép chắc chắn và một giọt với xác suất một nửa có cùng trung bình. Để phân biệt chúng, cần cả phân bố của phần thưởng.

Lấy bài toán đơn giản nhất: sau tín hiệu báo trước là một phần thưởng có kích thước ngẫu nhiên $r$. Giả sử có nhiều nơ-ron \nt{DOPA}, và mỗi nơ-ron, nơ-ron thứ $i$, phụ trách ước lượng $V_i$ của riêng nó, sao cho lúc có phần thưởng, sai số của nó là $\delta_i=r-V_i$. Và giả sử mỗi nơ-ron tính sai số dương với trọng số $\alpha_i^+$, sai số âm với trọng số $\alpha_i^-$. Sau mỗi phần thưởng, ước lượng dịch đi một lượng
\begin{equation}
\Delta V_i \;=\; \eta\,\bigl(\alpha_i^+\,[\delta_i]_+ \;-\; \alpha_i^-\,[-\delta_i]_+\bigr).
\label{eq:asymmetric}
\end{equation}
Ước lượng ngừng thay đổi khi trung bình các hiệu chỉnh dương cân bằng các hiệu chỉnh âm:
\begin{empheq}[box=\keybox]{equation}
\alpha_i^+\,\bigl\langle[r-V_i]_+\bigr\rangle \;=\; \alpha_i^-\,\bigl\langle[V_i-r]_+\bigr\rangle ,
\qquad
\kappa_i \;=\; \frac{\alpha_i^+}{\alpha_i^++\alpha_i^-} .
\label{eq:expectile}
\end{empheq}
Khi $\kappa_i=1/2$, đây là trung bình thông thường. Khi $\kappa_i>1/2$, nơ-ron là một ``người lạc quan'': ước lượng của nó lệch về đầu trên của phân bố; khi $\kappa_i<1/2$ --- một ``người bi quan''.

Ví dụ: nước ép đến với xác suất $1/2$, kích thước giọt là $1$. Khi đó~\eqref{eq:expectile} cho $\kappa_i\cdot\tfrac12(1-V_i)=(1-\kappa_i)\cdot\tfrac12V_i$, tức là $V_i=\kappa_i$ (hình~\ref{fig:expectiles}). Một tập nơ-ron với các $\kappa_i$ khác nhau ghi lại phân bố phần thưởng, giống như một tập phân vị ghi lại một phân bố trong thống kê.

\begin{figure}[t]
\centering
\begin{tikzpicture}[>=Stealth,x=7cm,y=3cm]
\draw[->,gray!70!black] (-0.08,0) -- (1.15,0) node[right,black] {\small phần thưởng};
\draw[->,gray!70!black] (-0.08,0) -- (-0.08,0.75) node[left,black] {\small xác suất};
\fill[blue!25] (-0.03,0) rectangle (0.03,0.5);
\fill[blue!25] (0.97,0) rectangle (1.03,0.5);
\node[below,font=\small] at (0,0) {$0$};
\node[below,font=\small] at (1,0) {$1$};
\node[left,font=\scriptsize] at (-0.08,0.5) {$1/2$};
\foreach \k/\c/\lab/\h in {0.2/blue!60!black/{bi quan, $\kappa=0.2$}/0.62, 0.5/gray!50!black/{trung bình, $\kappa=0.5$}/0.42, 0.8/red!70!black/{lạc quan, $\kappa=0.8$}/0.22}{
  \draw[very thick,\c] (\k,0) -- (\k,\h);
  \node[above,font=\scriptsize,\c] at (\k,\h) {\lab};}
\end{tikzpicture}
\caption{Phân bố phần thưởng được ghi lại bởi một tập nơ-ron \nt{DOPA}. Phần thưởng bằng $0$ hoặc $1$ với xác suất $1/2$ (các cột xanh); nơ-ron có tỷ lệ $\kappa$ đi tới ước lượng $V=\kappa$~\eqref{eq:expectile}.}
\label{fig:expectiles}
\end{figure}

Điều đã đo: các nơ-ron \nt{DOPA} khác nhau có ``điểm đảo chiều'' khác nhau --- kích thước phần thưởng mà tại đó đáp ứng chuyển từ đỉnh vọt sang hõm sụt --- và những nơ-ron có độ bất đối xứng lớn hơn giữa đáp ứng với sai số dương và sai số âm thì đảo chiều ở phần thưởng lớn hơn, đúng như~\eqref{eq:expectile} đòi hỏi~\cite{Dabney2020}. Từ đáp ứng của một nhóm nơ-ron có thể khôi phục hình dạng của phân bố. Việc đây là chức năng chính của sự phân tán, chứ không phải hiệu ứng phụ, là giả thuyết.

\begin{keyblock}
\begin{proposition}[Phân bố từ sự bất đối xứng]
\label{prop:expectile}
Nếu các nơ-ron \nt{DOPA} khác nhau ở tỷ lệ $\kappa_i$ mà chúng dùng để tính sai số dương, thì các ước lượng của chúng hội tụ về những điểm khác nhau~\eqref{eq:expectile} của phân bố phần thưởng. Nhóm nơ-ron truyền đi không phải giá trị trung bình mà là phân bố, và qua đó --- rủi ro.
\end{proposition}
\end{keyblock}

\section{Không chỉ sai số mà còn tầm quan trọng}

Theo công thức sai số~\eqref{eq:ctd}, một sự kiện khó chịu --- cú sốc, vị đắng --- phải gây ra hõm sụt: kết quả tệ hơn kỳ vọng. Một phần nơ-ron \nt{DOPA} đúng là như vậy, nhưng phần khác lại đáp ứng với điều khó chịu bằng \emph{đỉnh vọt}, như với phần thưởng~\cite{Matsumoto2009}. Những nơ-ron này không báo dấu của sai số, mà báo rằng đã xảy ra điều gì đó \emph{quan trọng}: bất ngờ, mạnh, cần chú ý~\cite{BrombergMartin2010}.

Với kỹ sư, tiện nhất là coi đây là hai tín hiệu. Tín hiệu thứ nhất là sai số có dấu $\delta$: thay đổi trọng số theo hướng nào. Tín hiệu thứ hai là độ lớn $|\delta|$ của nó hoặc độ mới của sự kiện: thay đổi \emph{mạnh tới mức nào}, tức là tốc độ học; sau một sự kiện bất ngờ thì phải học nhanh hơn. Về ý nghĩa, nó gần với hệ số khuếch đại do noradrenaline đặt ở chương~\ref{sec:neurontypes}.

Còn có phương án thứ ba: một dây riêng cho một trục riêng. Các nơ-ron \nt{DOPA} đi tới một trong các vùng chọn hành động không đáp ứng với phần thưởng, mà với độ mới và cường độ của kích thích, và cần thiết để con vật học cách tránh điều nguy hiểm~\cite{Menegas2018}: trên dây không chạy ``tốt hơn hay tệ hơn'' mà là ``nguy hiểm hay không''.

Điều đã đo: các nhóm nơ-ron \nt{DOPA} khác nhau đáp ứng khác nhau với điều khó chịu và điều mới, và các đầu ra khác nhau dạy những điều khác nhau. Não dùng tín hiệu tầm quan trọng chính xác thế nào --- như tốc độ học, như tín hiệu chú ý hay như một sai số riêng theo trục nguy hiểm --- vẫn đang được bàn cãi.

\section{Không chỉ dạy mà còn thúc giục}

Dopamine còn có tác động tức thời: càng nhiều dopamine, con vật hành động càng sẵn lòng và càng nhanh. Làm sao dung hòa điều này với việc học?

Câu trả lời của kỹ sư là \emph{phân tách theo tần số}. Các đỉnh vọt và hõm sụt nhanh, kéo dài hàng trăm mili giây, mang $\delta$ và dạy. Mức chậm, thay đổi trong vài phút, mang một đại lượng khác --- phần thưởng trung bình trên một đơn vị thời gian $\bar R$~\cite{Niv2007}. Giả sử một hành động thực hiện trong thời gian $\tau_a$ đòi hỏi nỗ lực $C/\tau_a$: càng nhanh càng đắt. Bù lại, mỗi giây chậm trễ tốn $\bar R$ --- lượng phần thưởng có thể nhận được trong giây đó. Tổng chi phí $C/\tau_a+\bar R\,\tau_a$ nhỏ nhất khi
\begin{empheq}[box=\keybox]{equation}
\tau_a^{*} \;=\; \sqrt{\frac{C}{\bar R}} .
\label{eq:vigor}
\end{empheq}
Môi trường càng giàu, thời gian càng đắt và hành động nhanh càng có lợi: nếu $\bar R$ tăng gấp đôi, thời gian hành động giảm $\sqrt2\approx1.4$ lần. Lưu ý rằng $\bar R$ chính là phần thưởng kỳ vọng đã được trừ đi ở chương~\ref{sec:dofa}.

Lượng dopamine phóng ra ở nơi đến có thể thay đổi cả khi tần số xung không đổi: nó được điều chỉnh bởi các tín hiệu tại chỗ ngay ở đầu ra~\cite{Mohebi2019}. Nhưng thành phần chậm cũng có trong chính các xung: trong các thí nghiệm ở mục sau, sự tăng dần khi tiến gần phần thưởng cũng thấy được trong tần số của nơ-ron \nt{DOPA}~\cite{Kim2020}. Như vậy, mức chậm được tạo từ hai nguồn --- tần số xung và việc điều chỉnh lượng phóng tại chỗ --- và nguồn nào là chính dường như tùy thuộc vào đầu ra và vào nhiệm vụ.

\begin{keyblock}
\begin{proposition}[Hai tín hiệu trên một dây]
\label{prop:multiplex}
Hệ \nt{DOPA} truyền ít nhất hai đại lượng, tách nhau theo thời gian: sai số nhanh $\delta$ để dạy, và mức chậm đặt giá của thời gian~\eqref{eq:vigor}, qua đó đặt tốc độ hành động. Việc thành phần nhanh và thành phần chậm hành xử khác nhau đã được đo; việc thành phần chậm đúng bằng $\bar R$ là giả thuyết.
\end{proposition}
\end{keyblock}

\section{Tăng dần khi tiến lại gần}

Nếu con vật chạy dọc hành lang tới một phần thưởng ở xa, nồng độ dopamine trong vùng chọn hành động tăng dần đều trong nhiều giây khi đích đến gần hơn~\cite{Hamid2016}. Theo chương~\ref{sec:dofa}, điều này không nên xảy ra: mọi thứ diễn ra đúng kế hoạch, và $\delta$ phải bằng không.

Cách giải thích thứ nhất: sự tăng này không phải là $\delta$ mà là chính kỳ vọng $V$, tín hiệu ``đích đã gần, đáng cố gắng'' ở mục trước~\cite{Hamid2016}. Cách giải thích thứ hai giữ lại $\delta$~\cite{Kim2020}. Nếu kỳ vọng đúng, thì với tầm nhìn $\tau_\gamma$, trên đường tới phần thưởng $R$ nó tăng như $V=Re^{-(T-t)/\tau_\gamma}$, và theo công thức sai số~\eqref{eq:ctd} $\dd V/\dd t-V/\tau_\gamma=0$. Nhưng giả sử từ xa kỳ vọng bị đánh giá thấp, còn khi đến gần thì ước lượng chính xác dần. Khi đó kỳ vọng tăng dốc hơn, chẳng hạn như $V=Re^{-(T-t)/\tau_v}$ với $\tau_v<\tau_\gamma$, và
\begin{empheq}[box=\keybox]{equation}
\delta(t) \;=\; \frac{\dd V}{\dd t}-\frac{V}{\tau_\gamma} \;=\; V\Bigl(\frac{1}{\tau_v}-\frac{1}{\tau_\gamma}\Bigr) \;>\; 0 .
\label{eq:ramp}
\end{empheq}
Sai số dương và tăng cùng với $V$ --- chính là sự tăng dần đều ấy (hình~\ref{fig:ramp}). Với $\tau_\gamma=4\,\text{s}$ và $\tau_v=1\,\text{s}$, ta được $\delta=0.75\,V$ mỗi giây. Phiên bản này đã được kiểm tra trong một hành lang ảo bằng cách tức thời đưa con vật lại gần đích: dopamine đáp ứng bằng một đỉnh vọt, đúng như với một đạo hàm, chứ không phải như với một mức biến thiên đều~\cite{Kim2020}.

\begin{figure}[t]
\centering
\begin{tikzpicture}[>=Stealth,x=1.6cm,y=2.6cm]
\draw[->,gray!70!black] (-4.2,0) -- (0.5,0) node[below left,black] {\small $t$};
\draw[->,gray!70!black] (-4.2,0) -- (-4.2,1.2);
\foreach \x/\l in {-4/{$T-4$},-2/{$T-2$},0/{$T$}}{\draw[gray!60!black] (\x,-0.03) -- (\x,0.03); \node[below,font=\scriptsize] at (\x,0) {\l\,s};}
\draw[thick,densely dashed,gray!60!black] plot[domain=-4:0,samples=40] (\x,{exp(\x/4)});
\draw[very thick,blue!60!black] plot[domain=-4:0,samples=60] (\x,{exp(\x)});
\draw[very thick,red!70!black] plot[domain=-4:0,samples=60] (\x,{0.75*exp(\x)});
\node[anchor=south east,font=\small,gray!50!black] at (-1.3,0.77) {$V$ khi $\tau_v=\tau_\gamma$: $\delta=0$};
\node[right,font=\small,blue!60!black] at (0.05,1.0) {$V$ khi $\tau_v=1\,\text{s}$};
\node[right,font=\small,red!70!black] at (0.05,0.72) {$\delta$ theo~\eqref{eq:ramp}};
\end{tikzpicture}
\caption{Dopamine tăng dần khi tiến gần phần thưởng, xem như sai số dự đoán, $\tau_\gamma=4\,\text{s}$. Đường nét đứt là kỳ vọng tăng đúng như tầm nhìn đòi hỏi ($\delta=0$); đường xanh là kỳ vọng tăng dốc hơn; đường đỏ là sai số~\eqref{eq:ramp}.}
\label{fig:ramp}
\end{figure}

Điều đã đo: sự tăng dần đều có thật --- cả trong nồng độ dopamine lẫn trong tần số xung của nơ-ron \nt{DOPA}~\cite{Kim2020} --- và những bước nhảy tức thời của bối cảnh gây ra bước nhảy dopamine. Cách giải thích nào trong hai cách là đúng --- hay cả hai đều đúng ở những đầu ra khác nhau --- vẫn đang được bàn cãi.

\section{Nhìn lại phía sau}

Mọi lý thuyết ở trên đều nhìn \emph{về phía trước}. Một lý thuyết hiện đại đảo ngược hướng này~\cite{Jeong2022}. Giả sử não không học cách dự đoán phần thưởng từ tín hiệu báo trước, mà sau khi nhận phần thưởng thì \emph{nhìn lại}: sự kiện nào thường đi trước nó hơn các sự kiện khác? Thước đo của mối liên hệ đó là hiệu của hai xác suất:
\begin{empheq}[box=\keybox]{equation}
M \;=\; P(\text{tín hiệu trước phần thưởng}) \;-\; P(\text{tín hiệu trong cửa sổ ngẫu nhiên}) .
\label{eq:retro}
\end{empheq}
Nếu tín hiệu báo trước thường xuất hiện cả khi không có phần thưởng, số hạng thứ hai lớn, và mối liên hệ yếu. Trong lý thuyết này, dopamine không báo sai số dự đoán, mà báo sự kiện đã cho có khả năng là \emph{nguyên nhân} của phần thưởng đến mức nào.

Cơ chế nhìn lại đòi hỏi đúng những gì quy tắc ba thừa số ở chương~\ref{sec:dofa} đòi hỏi: mỗi sự kiện phải để lại một vết mờ dần, để lúc có phần thưởng có thể đọc được điều gì đã xảy ra trước đó. Khác biệt không nằm ở khớp thần kinh mà ở điều \nt{DOPA} truyền đi.

Trong những thí nghiệm mà lý thuyết~\eqref{eq:retro} và sai số~\eqref{eq:ctd} dự đoán khác nhau, lượng dopamine phóng ra trong một số trường hợp tuân theo lý thuyết thứ nhất~\cite{Jeong2022}. Nhưng trong những thí nghiệm khác, trong quá trình học, đáp ứng dopamine dịch dần từ thời điểm có phần thưởng sang thời điểm có tín hiệu báo trước --- đúng như việc học theo sai số~\eqref{eq:ctd} đòi hỏi~\cite{Amo2022}. Lý thuyết nào mô tả đúng não thì hiện chưa rõ.

\section{Sai số không chỉ về phần thưởng}

Sự mở rộng cuối cùng liên quan tới việc sai số là \emph{về cái gì}. Nếu thay nước ép được kỳ vọng bằng một loại khác, giá trị như nhau nhưng khác vị, thì giá trị không đổi, và sai số~\eqref{eq:ctd} bằng không. Thế nhưng nơ-ron \nt{DOPA} vẫn đáp ứng với sự thay thế đó~\cite{Takahashi2017}. Còn những đỉnh vọt dopamine được gây ra nhân tạo có thể dạy con vật mối liên hệ giữa hai kích thích trung tính, hoàn toàn không có phần thưởng~\cite{Sharpe2017}. Có vẻ như \nt{DOPA} báo sai số dự đoán không chỉ của phần thưởng, mà cả của việc \emph{điều gì} sẽ xảy ra.

Về hình thức, đây là một tổng quát hóa đơn giản của~\eqref{eq:ctd}: thay cho một dòng phần thưởng $r$, ta lấy một tập đặc trưng của những gì đang diễn ra $\varphi_k$ --- vị, vị trí, âm thanh --- và với mỗi đặc trưng có một kỳ vọng $M_k$ và một sai số riêng~\cite{Gardner2018}:
\begin{empheq}[box=\keybox]{equation}
\delta_k(t) \;=\; \varphi_k(t) \;+\; \frac{\dd M_k}{\dd t} \;-\; \frac{M_k}{\tau_\gamma} ,
\qquad k=1,\dots,K .
\label{eq:vectortd}
\end{empheq}
Phần thưởng chỉ là một trong các đặc trưng, và vectơ sai số không chỉ dạy điều gì là tốt, mà còn dạy điều gì theo sau điều gì --- tức một mô hình thế giới.

Tính không đồng nhất của các nơ-ron \nt{DOPA} phù hợp với điều này: trong cùng một nhiệm vụ, các nơ-ron khác nhau mang những đại lượng khác nhau --- một số liên quan tới phần thưởng, số khác tới chuyển động, số khác nữa tới hướng của sự chú ý~\cite{Engelhard2019}.

\begin{keyblock}
\begin{proposition}[Bó dây thay cho một dây]
\label{prop:bundle}
Tín hiệu của hệ \nt{DOPA} không phải là một con số. Nó được phân bố theo nơ-ron (phân bố~\eqref{eq:expectile}, các đặc trưng khác nhau~\eqref{eq:vectortd}, một trục nguy hiểm riêng) và theo thời gian (sai số nhanh và mức chậm~\eqref{eq:vigor}). Sai số vô hướng~\eqref{eq:ctd} là xấp xỉ đầu tiên của tín hiệu này, giống như bộ cộng có ngưỡng là xấp xỉ đầu tiên của nơ-ron.
\end{proposition}
\end{keyblock}

\section{Tổng kết}

\begin{table}[t]
\centering
\footnotesize
\setlength{\tabcolsep}{4pt}
\renewcommand{\arraystretch}{1.15}
\begin{tabular}{>{\raggedright}p{2.5cm}>{\raggedright}p{3.2cm}>{\raggedright}p{3.6cm}>{\raggedright\arraybackslash}p{4.2cm}}
\toprule
Lý thuyết & Thứ chạy trên dây & Phép đo chính & Điều đã biết \\
\midrule
sai số dự đoán (ch.~\ref{sec:dofa}) & đại lượng vô hướng $\delta$~\eqref{eq:ctd} & đỉnh vọt, dịch sang tín hiệu báo trước, hõm sụt & đã đo; xấp xỉ đầu tiên \\
phân bố & tập $\delta_i$ với độ bất đối xứng khác nhau & các nơ-ron khác nhau có điểm đảo chiều khác nhau & phù hợp với thí nghiệm; vai trò của sự phân tán là giả thuyết \\
tầm quan trọng & $|\delta|$, độ mới; trục nguy hiểm & một phần nơ-ron vọt lên với điều khó chịu & đã đo; cách sử dụng đang được bàn cãi \\
giá của thời gian & mức chậm $\bar R$ & dopamine làm hành động nhanh hơn; thành phần chậm có cả trong lượng phóng ở đầu ra lẫn trong tần số xung & sự tách nhanh và chậm đã được đo; $\bar R$ là giả thuyết \\
tăng dần khi tiến gần & $V$ hoặc $\delta$ khi ước lượng chính xác dần~\eqref{eq:ramp} & tăng dần đều, bước nhảy khi được dời chỗ trong hành lang & sự tăng đã được đo; cách giải thích đang được bàn cãi \\
nhìn lại phía sau & thước đo liên hệ nhân quả~\eqref{eq:retro} & thí nghiệm mà dự đoán của hai lý thuyết khác nhau & các kết quả mâu thuẫn nhau \\
vectơ sai số & $\delta_k$ theo đặc trưng~\eqref{eq:vectortd} & đáp ứng khi đổi vị; học không có phần thưởng & đã đo; dạng vectơ là giả thuyết \\
\bottomrule
\end{tabular}
\caption{Hệ \nt{DOPA} truyền gì: bức tranh chuẩn ở chương~\ref{sec:dofa} và các mở rộng hiện đại của nó.}
\label{tab:dofaalt}
\end{table}

Các lý thuyết trong bảng~\ref{tab:dofaalt} không loại trừ nhau: hầu hết đều giữ sai số dự đoán làm nền tảng và mở rộng định dạng của nó. Chỉ có lối nhìn lại phía sau thực sự cạnh tranh với nó, và cuộc tranh luận này sẽ do thí nghiệm phân xử. Với kỹ sư, kết luận đơn giản: cái giá của quảng bá ở mệnh đề~\ref{prop:broadcastcost} giảm xuống nếu dây dẫn thực ra là nhiều dây với những đích nhận khác nhau.

\section{Bài tập}

\begin{exercise}[\lvlA]
\label{ex:07:1}
\emph{Các điểm của phân bố cho một giọt.} Nước ép kích thước $1$ đến với xác suất $p=0.2$, nếu không thì phần thưởng bằng $0$. Theo~\eqref{eq:expectile}, tìm $V_i$ cho $\kappa_i=0.2$, $0.5$, $0.8$. Trung vị của phần thưởng này bằng bao nhiêu, và điểm~\eqref{eq:expectile} khác phân vị ở chỗ nào?
\end{exercise}

\begin{exercise}[\lvlB]
\label{ex:07:2}
\emph{Phân bố từ hai nơ-ron.} Phần thưởng bằng $0$ hoặc $x$, xác suất của $x$ là $p$. Hai nơ-ron với quy tắc~\eqref{eq:asymmetric} báo $V(1/2)=0.25$ và $V(0.8)=0.5$. Khôi phục $p$, $x$ và độ lệch chuẩn của phần thưởng. Nơ-ron có $\kappa=0.2$ sẽ báo gì?
\end{exercise}

\begin{exercise}[\lvlB]
\label{ex:07:3}
\emph{Mô phỏng: phân bố từ sự bất đối xứng.} Mô phỏng chín nơ-ron \nt{DOPA} với $\kappa_i=0.1,\dots,0.9$ và quy tắc~\eqref{eq:asymmetric}, $\alpha_i^+=2\kappa_i$, $\alpha_i^-=2(1-\kappa_i)$, với phần thưởng phân bố đều trên $[0,1]$. Độ phân tán của $V_i$ phụ thuộc vào $\eta$ thế nào, và cần $\eta$ nào để phân biệt phân bố này với hai giọt đồng khả năng $0.25$ và $0.75$?
\end{exercise}

\begin{exercise}[\lvlA]
\label{ex:07:4}
\emph{Giá của thời gian.} (a)~Suy ra~\eqref{eq:vigor} và tìm tổng chi phí tại điểm tối ưu. (b)~Não ước lượng $\bar R$ sai $k$ lần. Tìm mức trả dư tương đối khi $k=2$ và $k=4$. Mức dopamine chậm phải báo $\bar R$ chính xác tới đâu?
\end{exercise}

\begin{exercise}[\lvlB]
\label{ex:07:5}
\emph{Đỉnh vọt khi dời chỗ.} Kỳ vọng tăng theo~\eqref{eq:ramp} với $\tau_v=1\,\text{s}$, $\tau_\gamma=4\,\text{s}$; con vật được dời tức thời từ điểm $T-2\,\text{s}$ tới $T-1\,\text{s}$. Tìm diện tích đỉnh vọt $\delta$ theo tỷ lệ của $R$, và nó thay cho bao nhiêu giây của đoạn dốc trước khi dời. Việc dời chỗ sẽ cho thấy gì nếu dopamine báo chính $V$?
\end{exercise}

\begin{exercise}[\lvlB]
\label{ex:07:6}
\emph{Thiết kế: hai tín hiệu trên một dây.} Trên dây \nt{DOPA} có một mức tăng dần tới $s=1/60$ xung mỗi giây sau mỗi giây, và các sự kiện, mỗi sự kiện $A=3$ xung. Khớp thần kinh có vết $\tau_e=1\,\text{s}$ trừ khỏi tần số bản sao của nó được làm trơn với $\tau_f$. Với $\tau_f$ nào sự kiện mất không quá $10\,\%$, còn đóng góp của mức trong thời gian vết nhỏ hơn $10\,\%$ đóng góp của sự kiện?
\end{exercise}

\begin{exercise}[\lvlB]
\label{ex:07:7}
\emph{Thiết kế thí nghiệm: nhìn lại hay sai số.} Trong một cửa sổ, tín hiệu báo trước có xác suất $0.1$, sau nó có phần thưởng với xác suất $0.8$. So sánh $\Delta P=P(\text{thưởng}\mid\text{có báo})-P(\text{thưởng}\mid\text{không})$ và $M$ trong~\eqref{eq:retro}, nếu (a)~thêm phần thưởng không báo trước, $0.08$ mỗi cửa sổ; (b)~tăng gấp đôi tần suất tín hiệu không kèm thưởng.
\end{exercise}

\begin{exercise}[\lvlC]
\label{ex:07:8}
\emph{Bó dây và cái giá của quảng bá.} Trong mô hình của bài tập~\ref{ex:06:8}, $N$ nơ-ron được chia thành $K$ nhóm, mỗi nhóm có dây riêng, và vào mỗi dây rò rỉ một tỷ lệ $\rho$ sai số của nhóm khác. Chứng minh rằng hằng số học bằng $N/K+2+\rho^2N(1-1/K)$. Nó cho bao nhiêu lượt thử khi $K\to\infty$, $N=10^{10}$ và $\rho=0.01$?
\end{exercise}

\chapter{Khi nào khám phá, khi nào quyết định: thêm \nt{NORA}}
\chaptermark{Thêm \nt{NORA}}
\label{sec:nora}
\begin{chaptergoals}
\item cách một núm vặn hệ số khuếch đại biến nơ-ron từ bộ khuếch đại tương tự thành bộ so sánh;
\item vì sao tăng hệ số khuếch đại chỉ có ích với nhiễu phát sinh sau bộ khuếch đại;
\item cách hệ số khuếch đại chuyển mạng lựa chọn từ cân nhắc sang đã quyết mà không đổi trọng số nào;
\item vì sao hệ số khuếch đại là nghịch đảo nhiệt độ, và vì sao phần thưởng theo nó là chữ U ngược.
\end{chaptergoals}

\section{Còn thiếu gì}

Ở chương~\ref{sec:dofa}, mạng học được nhờ nhiễu: những độ lệch hoạt động ngẫu nhiên là sự tìm kiếm, còn \nt{DOPA} báo xem chúng có làm kết quả tốt hơn không. Nhưng cũng ở đó còn một tham số chưa được đặt --- tìm kiếm \emph{bao nhiêu}. Nếu nhiễu quá mạnh, mạng chạy lung tung và không dùng tới những gì nó đã biết. Nếu quá yếu, nó hết lần này tới lần khác chọn cùng một thứ và không nhận ra rằng đâu đó gần bên đã có gì tốt hơn. Trong học tăng cường, đây được gọi là sự lựa chọn giữa khai thác và khám phá, và thuật toán nào cũng có một núm vặn cho nó. Ở chương~\ref{sec:neurontypes}, ta đã giả định rằng trong não núm vặn này do \nt{NORA} điều khiển.

Ở người có vài chục nghìn nơ-ron noradrenaline (bảng~\ref{tab:types}), gần như tất cả tập trung trong một nhóm nhỏ, còn đầu ra của chúng tỏa ra hầu như khắp não, dù như người ta đã phát hiện, các phần khác nhau của nhóm này phần nào phục vụ những vùng khác nhau~\cite{Poe2020}. Đây là một kênh quảng bá còn hẹp hơn cả \nt{DOPA}. Nó hoạt động ở hai chế độ~\cite{AstonJones2005}. Ở chế độ \emph{nền}, các nơ-ron phát xung đều, với tần số từ dưới một tới vài xung mỗi giây, và tần số này thay đổi chậm theo mức tỉnh táo. Ở chế độ \emph{từng đợt}, chúng đáp ứng bằng một chùm xung ngắn với sự kiện quan trọng cho nhiệm vụ hiện tại, vào khoảng thời điểm ra quyết định. Một điểm đo kiểm tra tiện lợi nhưng không chính xác là đường kính đồng tử: nó thay đổi theo hoạt động của các nơ-ron này, nhưng cũng theo hoạt động của các vùng khác~\cite{Joshi2016} và của các đầu ra cholinergic trong vỏ não~\cite{Reimer2016}. Đồng tử cho thấy mức hưng phấn chung, chứ không riêng \nt{NORA}.

\section{Hệ số khuếch đại}

Điều đã đo là: noradrenaline ức chế hoạt động nền của nơ-ron mạnh hơn đáp ứng của chúng với kích thích, và tỷ số tín hiệu trên nền tăng lên~\cite{Foote1975NE}. Từ thí nghiệm không suy ra được rằng độ dốc đặc tuyến của từng nơ-ron bị nhân lên đúng theo nghĩa đen với một hệ số. Như ở chương~\ref{sec:neurontypes}, ta lấy một mô hình đơn giản thể hiện hiệu ứng này: noradrenaline thay đổi độ dốc đặc tuyến truyền đạt của bên nhận~\cite{ServanSchreiber1990}. Khi đó các đầu vào yếu, dưới ngưỡng (nền) cho ra còn ít hơn, còn các đầu vào mạnh cho ra còn nhiều hơn. Giả sử nơ-ron đáp ứng bằng tần số
\begin{empheq}[box=\keybox]{equation}
r \;=\; F\bigl(g\,(u-\theta)\bigr), \qquad F(z)=\frac{r_{\max}}{1+e^{-z}} ,
\label{eq:gain}
\end{empheq}
trong đó $u$ là dòng đầu vào, $\theta$ là ngưỡng, còn $g$ là hệ số khuếch đại, thứ mà trong mô hình do \nt{NORA} điều khiển. Khi $g$ nhỏ, tần số đi theo đầu vào một cách trơn tru trong một dải rộng: nơ-ron là một \emph{bộ khuếch đại tương tự}. Khi $g$ lớn, hầu như mọi đầu vào trên ngưỡng đều cho tần số tối đa, dưới ngưỡng thì bằng không: nơ-ron là một \emph{bộ so sánh}, gần như một phần tử số (hình~\ref{fig:gaincurves}). Một núm vặn chuyển cùng một nơ-ron giữa hai chế độ mà trong điện tử phải dùng những mạch khác nhau.

\begin{figure}[t]
\centering
\begin{tikzpicture}[>=Stealth,x=1.1cm,y=2.4cm]
\draw[->,gray!70!black] (-3.3,0) -- (3.4,0) node[below left,black] {\small $u$};
\draw[->,gray!70!black] (-3.3,0) -- (-3.3,1.2) node[left,black] {\small $r$};
\draw[densely dashed,gray!60!black] (0,0) -- (0,1.1);
\node[below,font=\small] at (0,0) {$\theta$};
\draw[densely dotted,gray!60!black] (-3.3,1) -- (3.2,1) node[right,black,font=\small] {$r_{\max}$};
\draw[very thick,blue!60!black] plot[domain=-3.2:3.2,samples=60] (\x,{1/(1+exp(-0.8*\x))});
\draw[very thick,orange!85!black] plot[domain=-3.2:3.2,samples=60] (\x,{1/(1+exp(-2*\x))});
\draw[very thick,red!70!black] plot[domain=-3.2:3.2,samples=120] (\x,{1/(1+exp(min(-8*\x,9)))});
\node[blue!60!black,font=\small,anchor=west] at (0.5,0.12) {$g$ nhỏ: bộ khuếch đại};
\node[red!70!black,font=\small,anchor=east] at (-0.3,0.85) {$g$ lớn: bộ so sánh};
\end{tikzpicture}
\caption{Đặc tuyến truyền đạt~\eqref{eq:gain} với ba giá trị của hệ số khuếch đại $g$.}
\label{fig:gaincurves}
\end{figure}

Khuếch đại lớn có giúp chống nhiễu không? Không phải lúc nào cũng vậy. Giả sử hai đầu vào khác nhau một lượng $\Delta u$, và cần nói cái nào lớn hơn. Nếu nhiễu $\sigma_{\text{trước}}$ cộng vào đầu vào \emph{trước} bộ khuếch đại, thì cả tín hiệu lẫn nhiễu đều được khuếch đại, và tỷ số của chúng không đổi. Còn nếu nhiễu $\sigma_{\text{sau}}$ cộng vào \emph{sau} bộ khuếch đại --- từ các khớp thần kinh không tin cậy và các xung ngẫu nhiên của chính mạng, như ở chương~\ref{sec:code} --- thì tỷ số tín hiệu trên nhiễu
\begin{empheq}[box=\keybox]{equation}
d' \;=\; \frac{g\,\Delta u}{\sqrt{g^2\sigma_{\text{trước}}^2+\sigma_{\text{sau}}^2}}
\label{eq:gainsnr}
\end{empheq}
tăng theo $g$ cho tới khi $g\sigma_{\text{trước}}$ bằng $\sigma_{\text{sau}}$~\cite{ServanSchreiber1990}.

\begin{keyblock}
\begin{proposition}[Khi nào khuếch đại có ích]
\label{prop:gainnoise}
Khi tăng hệ số khuếch đại, \nt{NORA} chỉ cải thiện việc phân biệt các đầu vào trước loại nhiễu phát sinh \emph{sau} bộ khuếch đại, trong chính mạng. Nhiễu đến cùng với đầu vào thì khuếch đại không loại bỏ được.
\end{proposition}
\end{keyblock}

\section{Khuếch đại trong lựa chọn giữa hai phương án}

Quay lại bài toán lựa chọn ở chương~\ref{sec:gaba}: hai nhóm nơ-ron \nt{GLUT} với các đầu vào $I_1$, $I_2$ ức chế lẫn nhau với cường độ $\beta$. Giờ mỗi nhóm có một hệ số khuếch đại $g$: trong các phương trình lựa chọn~\eqref{eq:wta}, biểu thức trong ngoặc được nhân với $g$. Khi cả hai nhóm cùng hoạt động, tần số xác lập được tìm từ $r_1=g(I_1-\beta r_2)$, $r_2=g(I_2-\beta r_1)$. Cộng và trừ hai đẳng thức này, ta được tổng $r_1+r_2=g(I_1+I_2)/(1+g\beta)$ và hiệu $r_1-r_2=g(I_1-I_2)/(1-g\beta)$. Tỷ số của chúng cho biết mạng ưu tiên một đầu vào hơn đầu vào kia rõ rệt tới mức nào:
\begin{empheq}[box=\keybox]{equation}
\frac{r_1-r_2}{r_1+r_2} \;=\; \varepsilon\,\frac{1+g\beta}{1-g\beta}, \qquad \varepsilon=\frac{I_1-I_2}{I_1+I_2}, \quad g\beta<1 .
\label{eq:wtagain}
\end{empheq}
Chênh lệch nhỏ $\varepsilon$ giữa các đầu vào được khuếch đại $(1+g\beta)/(1-g\beta)$ lần, và thừa số này tăng không giới hạn khi $g\beta$ tiến tới một. Khi $g\beta>1$, trạng thái cả hai nhóm cùng hoạt động là không ổn định, như trong mệnh đề~\ref{prop:wta}: một nhóm thắng, nhóm kia im lặng (hình~\ref{fig:pitchfork}).

\begin{figure}[t]
\centering
\begin{tikzpicture}[>=Stealth,x=3.2cm,y=1.5cm]
\draw[->,gray!70!black] (0,0) -- (2.15,0) node[below left,black] {\small $g\beta$};
\draw[->,gray!70!black] (0,-1.25) -- (0,1.35) node[above,black] {\small $(r_1-r_2)/(r_1+r_2)$};
\draw[gray!60!black] (1,-0.04) -- (1,0.04); \node[below right,font=\small] at (1,0) {$1$};
\node[left,font=\scriptsize] at (0,1) {$1$}; \node[left,font=\scriptsize] at (0,-1) {$-1$};
\draw[very thick,blue!60!black] plot[domain=0:0.905,samples=60] (\x,{0.05*(1+\x)/(1-\x)}) -- (1,1) -- (2.05,1);
\draw[very thick,blue!60!black] (1.105,-1) -- (2.05,-1);
\draw[thick,densely dashed,gray!60!black] (1,0) -- (2.05,0);
\node[font=\small,align=center] at (0.45,-0.62) {đang cân nhắc:\\cả hai nhóm\\hoạt động};
\node[font=\small,align=center] at (1.55,0.45) {đã quyết:\\một nhóm hoạt động};
\node[font=\small,blue!60!black,anchor=south west] at (1.05,-0.97) {đầu vào yếu đã thắng trước --- giữ nguyên};
\end{tikzpicture}
\caption{Lựa chọn giữa hai phương án với các hệ số khuếch đại khác nhau, các đầu vào khác nhau $\varepsilon=0.05$. Khi $g\beta>1$, cả hai lời giải đều ổn định (đường liền; đầu vào yếu thắng khi $g\beta>(1+\varepsilon)/(1-\varepsilon)\approx1.1$), và mạng giữ lời giải nó đã chọn; trạng thái đối xứng (đường nét đứt) không ổn định.}
\label{fig:pitchfork}
\end{figure}

\begin{keyblock}
\begin{proposition}[Khuếch đại chuyển chế độ lựa chọn]
\label{prop:gainwta}
Trong mạng lựa chọn có ức chế lẫn nhau $\beta$, hệ số khuếch đại $g$ đưa mạng vượt qua ranh giới $g\beta=1$. Dưới ranh giới, mạng ``cân nhắc'': cả hai nhóm hoạt động, còn chênh lệch đầu vào được khuếch đại theo~\eqref{eq:wtagain}. Trên ranh giới, mạng ``đã quyết'': một nhóm hoạt động, và quyết định được giữ nguyên. Bằng cách thay đổi $g$, \nt{NORA} có thể chuyển mạng giữa hai chế độ này mà không đụng tới một trọng số nào.
\end{proposition}
\end{keyblock}

\section{Khuếch đại như nhiệt độ}

Giờ ta thêm vào lựa chọn nhiễu phát sinh trong chính mạng, sau bộ khuếch đại. Nhóm nào thắng được quyết định bởi dấu của đại lượng $g(u_A-u_B)+\eta$, trong đó $u_A-u_B$ là chênh lệch đầu vào, tức là chênh lệch giữa các trọng số đã học ở chương~\ref{sec:dofa}, còn $\eta$ là nhiễu với thang đo $\sigma$. Nếu nhiễu có phân bố logistic (với phân bố Gauss đường cong gần như y hệt), xác suất chọn $A$ bằng
\begin{empheq}[box=\keybox]{equation}
P(A) \;=\; \frac{1}{1+e^{-g\,(u_A-u_B)/\sigma}} .
\label{eq:softmax}
\end{empheq}
Lập trình viên sẽ nhận ra ở đây phép chọn softmax với nhiệt độ $T=\sigma/g$. Hệ số khuếch đại là nghịch đảo của nhiệt độ. Khi $g$ nhỏ, ngay cả phương án tốt hơn rõ rệt cũng chỉ được chọn thường xuyên hơn phương án tệ một chút: mạng khám phá nhiều. Khi $g$ lớn, phương án tốt nhất đã biết gần như luôn được chọn: mạng khai thác những gì nó biết.

Cần nói rõ $g$ trong~\eqref{eq:wtagain} và~\eqref{eq:softmax} là gì: đó là mức khuếch đại \emph{hiệu dụng} của chênh lệch đầu vào của nhiệm vụ hiện tại tại thời điểm lựa chọn. Hai chế độ của \nt{NORA} tác động lên nó theo hướng ngược nhau. Chùm xung từng đợt vào lúc quyết định làm nó tăng, và mạng củng cố lựa chọn. Ngược lại, hoạt động nền cao đi kèm với khám phá: ở người, trước một lựa chọn ngẫu nhiên, đồng tử nền giãn rộng hơn~\cite{JepmaNieuwenhuis2011}, còn ở chuột cống, đầu vào \nt{NORA} tới vỏ não đai trước chuyển lựa chọn sang chế độ ngẫu nhiên~\cite{Tervo2014stochastic}. Như vậy, \nt{NORA} nền cao tương ứng với $g$ hiệu dụng \emph{nhỏ}, còn chùm xung tương ứng với $g$ lớn.

\section{Khám phá bao nhiêu}

Mức khuếch đại nào là tốt nhất? Chúng tôi đã kiểm tra trên mô hình. Mạng chọn một trong hai hành động theo~\eqref{eq:softmax}; mỗi hành động mang lại phần thưởng với một xác suất nào đó, và mạng ước lượng xác suất này từ phần thưởng của hành động được chọn, như ở chương~\ref{sec:dofa}. Thỉnh thoảng thế giới thay đổi: cả hai xác suất được rút thăm lại. Có thể thay đổi cả hành động đang được chọn lẫn hành động mà mạng đã lâu không thử; về hành động thứ hai, mạng chỉ có thể biết được bằng cách khám phá. Hình~\ref{fig:invertedU} cho thấy mạng nhận được bao nhiêu phần của phần thưởng tốt nhất có thể với các giá trị $g$ cố định khác nhau.

\begin{figure}[t]
\centering
\begin{tikzpicture}[>=Stealth,x=1.2cm,y=1cm]
\draw[->,gray!70!black] (0,0) -- (7.6,0) node[below left,black] {\small $g$};
\draw[->,gray!70!black] (0,0) -- (0,4.4) node[above,black,font=\small] {tỷ lệ so với phần thưởng tối ưu};
\foreach \x/\l in {0/0.5,1/1,2/2,3/4,4/8,5/16,6/32,7/64}{\draw[gray!60!black] (\x,-0.06) -- (\x,0.06); \node[below,font=\scriptsize] at (\x,0) {\l};}
\foreach \v/\y in {0.8/0.8,0.9/2.4,1.0/4.0}{\draw[gray!60!black] (-0.06,\y) -- (0.06,\y); \node[left,font=\scriptsize] at (0,\y) {\v};}
\draw[very thick,blue!60!black] plot[mark=*,mark size=1.3pt] coordinates {(0,0.368) (1,0.880) (2,1.728) (3,2.784) (4,3.456) (5,3.616) (6,3.488) (7,3.264)};
\draw[very thick,red!70!black] plot[mark=*,mark size=1.3pt] coordinates {(0,0.368) (1,0.672) (2,1.184) (3,1.840) (4,2.176) (5,1.920) (6,1.456) (7,1.104)};
\node[blue!60!black,font=\small,anchor=south] at (5,3.7) {thế giới ổn định};
\node[red!70!black,font=\small,anchor=north] at (5.1,1.25) {thế giới thay đổi nhanh};
\draw[->,thick,blue!60!black] (5,3.25) -- (5,3.55);
\draw[->,thick,red!70!black] (4,1.8) -- (4,2.1);
\node[font=\small\itshape,gray!35!black,anchor=west] at (0.15,2.3) {mò mẫm};
\node[font=\small\itshape,gray!35!black] at (6.45,2.55) {không thấy thay đổi};
\end{tikzpicture}
\caption{Tỷ lệ so với phần thưởng tốt nhất có thể khi hệ số khuếch đại $g$ cố định (thang $g$ là thang logarit). Đường xanh: thế giới thay đổi trung bình một lần sau mỗi nghìn lượt thử; đường đỏ: một lần sau mỗi hai mươi lượt. Mũi tên đánh dấu $g$ tốt nhất: trong thế giới thay đổi nhanh, nó chỉ bằng một nửa. Trung bình của 60 lần chạy, mỗi lần 4000 lượt thử.}
\label{fig:invertedU}
\end{figure}

Với $g=0.5$, mạng hầu như không dùng tới hiểu biết của mình và nhận được khoảng ba phần tư mức có thể; với $g$ rất lớn, nó bỏ lỡ các thay đổi. Giá trị tốt nhất: $g=16$ khi thế giới thay đổi một lần sau mỗi nghìn lượt thử (tỷ lệ $0.976$), và $g=8$ khi một lần sau mỗi hai mươi lượt ($0.886$).

\begin{keyblock}
\begin{proposition}[Chữ U ngược]
\label{prop:explore}
Phần thưởng mà mạng thu được phụ thuộc vào hệ số khuếch đại theo dạng chữ U ngược: $g$ quá nhỏ thì phí lượt thử vào khám phá, $g$ quá lớn thì không nhận ra thay đổi. Thế giới thay đổi càng nhanh thì $g$ tốt nhất càng nhỏ.
\end{proposition}
\end{keyblock}

Chữ U ngược cũng được quan sát thấy trong thí nghiệm: con vật giải nhiệm vụ tốt nhất khi hoạt động nền của nơ-ron \nt{NORA} ở mức vừa phải và có đáp ứng từng đợt rõ nét, còn khi hoạt động nền cao thì chúng bị phân tâm và chuyển qua lại giữa các nhiệm vụ~\cite{AstonJones2005}. Điều này khớp với sự khác biệt giữa các chế độ ở mục trước: chùm xung từng đợt lúc quyết định cho $g$ hiệu dụng lớn đúng vào lúc cần chọn, còn hoạt động nền cao không có chùm xung thì khuếch đại mọi thứ, kể cả các đầu vào không liên quan, nên $g$ hiệu dụng cho nhiệm vụ hiện tại giảm xuống --- đó là khám phá.

\section{Ai vặn núm}

Vì mức khuếch đại tốt nhất phụ thuộc vào tốc độ thay đổi của thế giới, nên sẽ tốt nếu điều chỉnh được nó. Nhưng làm sao nơ-ron \nt{NORA} biết được thế giới đã thay đổi? Dấu hiệu tự nhiên là \emph{sự bất ngờ}: sai số dự đoán trở nên lớn hơn bình thường. Điều quan trọng là phân biệt hai loại bất định. Nếu phần thưởng đơn giản là ngẫu nhiên, sai số luôn lớn, và điều đó được dự liệu. Còn nếu sai số đột nhiên tăng so với độ lớn thường ngày của nó, nghĩa là có gì đó đã thay đổi. Theo một giả thuyết, chính sự bất định không được dự liệu này là thứ \nt{NORA} báo đi, còn sự bất định được dự liệu thì do \nt{ACET} báo~\cite{YuDayan2005}.

Làm gì với tín hiệu này? Có thể giảm khuếch đại và khám phá nhiều hơn. Có thể tăng tốc độ học để nhanh chóng quên đi các ước lượng đã lỗi thời: thí nghiệm cho thấy sau một thay đổi đột ngột, người ta học nhanh hơn, và đồng tử khi đó giãn ra~\cite{Nassar2012}; xin nhắc lại rằng đồng tử là chỉ báo không chỉ của \nt{NORA} mà cả của \nt{ACET}~\cite{Reimer2016}. Và cũng có thể làm cả hai cùng lúc: theo một giả thuyết khác nữa, chùm xung từng đợt của \nt{NORA} hoạt động như một \emph{ngắt} --- tín hiệu để mạng xóa trạng thái hiện tại và tái cấu hình cho bối cảnh mới~\cite{BouretSara2005}, như đường ngắt chung của bộ xử lý.

Xin nói thẳng điều chúng tôi chưa tìm ra. Trong mô hình, chúng tôi đã thử cả hai cách điều chỉnh đơn giản: giảm $g$ hoặc tăng tốc độ học khi sai số đã làm trơn vượt quá trung bình dài hạn của nó. Trong một thế giới lúc thì đứng yên rất lâu, lúc thì thay đổi thường xuyên, thiết lập tốt nhất của cả hai cách cho $0.926$--$0.927$ so với phần thưởng tối ưu --- gần như bằng $g$ cố định tốt nhất ($0.925$ khi $g=8$) hay một tốc độ học cố định nhưng đủ lớn ($0.928$), còn thiết lập kém thì cho ít hơn. Các đường cong trên hình~\ref{fig:invertedU} thoải ở gần đỉnh, và trong một thế giới như vậy hầu như chẳng có gì để điều chỉnh. Việc điều chỉnh phải được đền đáp ở nơi các chế độ khác nhau đòi hỏi những thiết lập rất khác nhau. \nt{NORA} thực hiện luật điều khiển cụ thể nào thì hiện chưa được xác định.

\section{Tổng kết}

\begin{table}[t]
\centering
\footnotesize
\setlength{\tabcolsep}{4pt}
\renewcommand{\arraystretch}{1.15}
\begin{tabular}{>{\raggedright}p{3.0cm}>{\raggedright}p{4.4cm}>{\raggedright\arraybackslash}p{6.0cm}}
\toprule
\nt{NORA} làm gì & Bằng cách nào & Điều đã biết \\
\midrule
thay đổi độ dốc đáp ứng của nơ-ron & hệ số khuếch đại $g$ trong~\eqref{eq:gain} & sự tăng tỷ số tín hiệu trên nhiễu đã được đo; mô hình nhân là một sự đơn giản hóa \\
chuyển chế độ lựa chọn & đưa mạng vượt qua $g\beta=1$ & suy ra từ mô hình~\eqref{eq:wtagain}; chưa có phép đo trực tiếp ngưỡng này \\
đặt mức khám phá & $g$ là nghịch đảo nhiệt độ~\eqref{eq:softmax}; nền --- $g$ hiệu dụng nhỏ (khám phá), chùm xung --- lớn (quyết định) & chữ U ngược và hai chế độ hoạt động đã được đo; mối liên hệ với khám phá là giả thuyết có cơ sở vững \\
báo về thay đổi & sự bất định không được dự liệu & đồng tử và tốc độ học tăng sau thay đổi --- đã đo; việc đây là tín hiệu của \nt{NORA} là giả thuyết \\
ngắt và xóa & chùm xung từng đợt với sự kiện bất ngờ & chùm xung với các sự kiện quan trọng đã được đo; ``ngắt'' là giả thuyết \\
\bottomrule
\end{tabular}
\caption{\nt{NORA} bổ sung gì cho mạng gồm \nt{GLUT}, \nt{GABA} và \nt{DOPA}.}
\label{tab:nora}
\end{table}

\nt{DOPA} cho mạng biết phải thay đổi trọng số \emph{theo hướng nào}. \nt{NORA} không thay đổi một trọng số nào, nhưng quyết định mạng dùng những gì nó đã biết \emph{ra sao} (bảng~\ref{tab:nora}): một núm vặn --- hệ số khuếch đại --- biến nơ-ron từ bộ khuếch đại thành bộ so sánh, biến mạng lựa chọn từ ``đang cân nhắc'' thành ``đã quyết'', và biến lựa chọn từ khám phá thành khai thác. \nt{NORA} biết thế giới thay đổi nhanh tới mức nào bằng cách nào --- đó là một câu hỏi còn mở.

\section{Bài tập}

\begin{exercise}[\lvlA]
\label{ex:08:1}
\emph{Một núm vặn cho cả bộ não.} (a)~Dựa vào bảng~\ref{tab:types}, ước lượng mỗi nơ-ron \nt{NORA} ứng với bao nhiêu nơ-ron của não. (b)~Nhiễu trước bộ khuếch đại $\sigma_{\text{trước}}=0.5$, sau bộ khuếch đại $\sigma_{\text{sau}}=4$. Với $g$ nào đại lượng $d'$ trong~\eqref{eq:gainsnr} đạt $90\,\%$ giới hạn của nó?
\end{exercise}

\begin{exercise}[\lvlB]
\label{ex:08:2}
\emph{Lựa chọn có hệ số khuếch đại.} $\tau\,\dd r_1/\dd t=-r_1+g[I_1-\beta r_2]_+$, $\tau\,\dd r_2/\dd t=-r_2+g[I_2-\beta r_1]_+$, $\tau=20\ms$. (a)~Sau bao lâu $r_1-r_2$ tắt dần khi $g\beta=0.5$ và $0.9$? (b)~Với $\varepsilon=0.05$, tìm $g\beta$ mà dưới đó cả hai nhóm hoạt động và trên đó chiến thắng của đầu vào yếu là ổn định.
\end{exercise}

\begin{exercise}[\lvlA]
\label{ex:08:3}
\emph{Softmax từ nhiễu.} Mỗi nhóm trong $K$ nhóm nhận nhiễu Gumbel riêng, $P(\eta_k<z)=\exp(-e^{-z/\sigma})$, và giá trị $gu_k+\eta_k$ lớn nhất thắng. Tìm $P(k)$. Với $g$ nào phương án tốt hơn trong hai phương án có $u_A-u_B=0.1$, $\sigma=1$ được chọn trong $90\,\%$ trường hợp?
\end{exercise}

\begin{exercise}[\lvlB]
\label{ex:08:4}
\emph{Ước lượng mức khuếch đại tốt nhất.} Trong mô hình ở hình~\ref{fig:invertedU}, xác suất phần thưởng sau một thay đổi phân bố đều trên $[0,1]$. Mạng thử hành động tệ hơn một lần sau mỗi $e^{g\Delta}$ lượt; cho giá trị này bằng $1/h$, ta được $g^*\approx\ln(1/h)/\bar\Delta$. Tìm $\bar\Delta=\langle|p_A-p_B|\rangle$ và $g^*$ khi $h=0.001$, $0.01$, $0.05$.
\end{exercise}

\begin{exercise}[\lvlB]
\label{ex:08:5}
\emph{Mô phỏng: khuếch đại và tốc độ thay đổi.} Tìm $g^*(h)$ với $h$ từ $0.001$ tới $0.2$ bằng một bản sao của \texttt{models/\allowbreak nora\_explore.py} (nó ghi tệp vào thư mục hiện tại). Ước lượng ở bài tập~\ref{ex:08:4} đúng ở đâu, và vì sao khi thay đổi nhanh $g^*$ thôi giảm?
\end{exercise}

\begin{exercise}[\lvlB]
\label{ex:08:6}
\emph{Thiết kế: ngắt cho mạng lựa chọn.} Mạng của bài tập~\ref{ex:08:2} với $\beta=2$, $\tau=20\ms$, $g=1$ giữ $r_1=0.9$, $r_2=0$, dù đầu vào giờ là $I_1=0.9$, $I_2=1.0$. \nt{NORA} hạ $g$ xuống $g_{\text{lo}}$ trong thời gian $T_p$. Tìm $T_p$ nhỏ nhất khi $g_{\text{lo}}=0$ bằng giải tích, và bằng mô phỏng khi $g_{\text{lo}}=0.25$.
\end{exercise}

\begin{exercise}[\lvlC]
\label{ex:08:7}
\emph{Khi nào điều chỉnh được đền đáp.} Một nhà tiên tri biết thời kỳ là ổn định hay biến động, và đặt $g$ riêng cho mỗi thời kỳ. (a)~Bằng mô phỏng, tìm phần lợi của nó so với $g$ cố định tốt nhất trong thế giới của chương (thời kỳ $1000$ lượt, $h=0.001$ và $0.05$). (b)~Dựng một thế giới có phần lợi lớn hơn, và tìm phần lợi mà điều chỉnh theo sự bất ngờ giành được.
\end{exercise}

\chapter{Khi nào ghi, khi nào đọc: thêm \nt{ACET}}
\chaptermark{Thêm \nt{ACET}}
\label{sec:acet}

\section{Còn thiếu gì}

Một vi mạch nhớ, ngoài các dây địa chỉ và dữ liệu, còn có thêm một đường nữa --- cho phép ghi. Khi đường này tắt, bộ nhớ chỉ được đọc; khi nó bật, những gì đang có trên các đường dữ liệu được ghi vào. Không có đường này thì bộ nhớ không hoạt động: nếu mỗi lần đọc đồng thời lại ghi một thứ gì đó, thì những gì đã đọc, kể cả những gì đọc sai, lập tức bị ghi ngược trở lại.

Trong não, vấn đề này còn gay gắt hơn trong vi mạch. Ở chương~\ref{sec:glutcomputer}, bộ nhớ là một nhóm nơ-ron \nt{GLUT} được chính các liên kết của nó giữ ở trạng thái bật (hình~\ref{fig:bistable}). Ở chương~\ref{sec:dofa}, ta đã thấy các liên kết này học theo quy tắc Hebb ngay trong lúc hoạt động. Như vậy, cùng những khớp thần kinh đó vừa lưu cái cũ vừa tiếp nhận cái mới, và không có pha ghi riêng, cũng như không có đồng hồ. Mạng phân biệt lúc cần ghi nhớ điều đang thấy với lúc cần nhớ lại điều đã biết bằng cách nào? Ở chương~\ref{sec:neurontypes}, ta đã giả định rằng đường này do \nt{ACET} giữ. Trong chương này, ta sẽ kiểm tra một đường như vậy phải hoạt động thế nào và vì sao không thể thiếu nó.

\section{Ba tác động của một dây}

Các nơ-ron acetylcholine hoạt động bên trong não không nhiều, và chúng tập trung trong vài nhóm nhỏ, còn đầu ra của chúng tỏa khắp vỏ não. Đây là một kênh quảng bá, giống như \nt{DOPA} và \nt{NORA}. Mức acetylcholine trong vỏ não cao khi thức và tập trung chú ý, và thấp trong giấc ngủ sóng chậm~\cite{Hasselmo1999}. Như với dopamine, ý nghĩa của tín hiệu do thụ thể của bên nhận quyết định (mệnh đề~\ref{prop:receptor}): acetylcholine có cả thụ thể nhanh, mở kênh, lẫn thụ thể chậm, khởi động phản ứng bên trong tế bào. Với ta, có ba tác động quan trọng.

\emph{Thứ nhất: làm yếu các liên kết nội bộ.} Thí nghiệm trên lát cắt vỏ não khứu giác cho thấy acetylcholine ức chế mạnh sự truyền dẫn qua các liên kết giữa các nơ-ron trong cùng một vùng, và gần như không động tới các đầu vào đến từ bên ngoài~\cite{HasselmoBower1992}.

\emph{Thứ hai: tăng cường đầu vào.} Acetylcholine làm tăng tính dễ kích thích của nơ-ron và tạo thuận cho sự truyền dẫn qua một số đường vào; tín hiệu từ các giác quan trở nên lớn hơn hoạt động riêng của mạng~\cite{Hasselmo2006}.

\emph{Thứ ba: làm trọng số dễ thay đổi hơn.} Khi mức acetylcholine cao, các liên kết thay đổi dễ hơn~\cite{Hasselmo2006}.

Với kỹ sư, cả ba tác động là một thao tác: chuyển từ chế độ ``đọc'' sang chế độ ``ghi''. Mạng thôi lắng nghe chính nó, bắt đầu lắng nghe đầu vào và ghi nhớ những gì nó nghe được.

\section{Mạng tự điền nốt}

Lấy $N$ nơ-ron \nt{GLUT} nối với nhau. Ta lưu trong các liên kết của chúng vài mẫu --- mỗi mẫu là một tập $pN$ nơ-ron cùng hoạt động --- bằng quy tắc Hebb~\cite{Hebb1949}. Nếu đưa vào mạng như vậy một nửa mẫu, các liên kết sẽ kích thích nửa còn thiếu: mạng \emph{điền nốt} mẫu từ một phần, như nhóm ở hình~\ref{fig:bistable} tự duy trì chính nó. Đó là bộ nhớ liên kết: địa chỉ là một phần của nội dung. \nt{GABA}, như ở chương~\ref{sec:gaba}, giữ tổng số nơ-ron hoạt động quanh $pN$, để hai mẫu không thể cùng sáng lên.

Ký hiệu mức acetylcholine là $a$, từ $0$ tới $1$, và đưa trực tiếp ba tác động của nó vào mô hình:
\begin{empheq}[box=\keybox]{equation}
\tau\frac{\dd r_i}{\dd t} = -r_i + F\Bigl(\tfrac{1+a}{2}\,x_i + (1-a)\sum_j W_{ij}\,r_j - g\Bigr),
\qquad
\frac{\dd W_{ij}}{\dd t} = \eta\,a\,(r_i-p)(r_j-p) .
\label{eq:acetnet}
\end{empheq}
Ở đây $r_i$ là tần số của nơ-ron $i$, $x_i$ là đầu vào của nó từ các giác quan, $F$ là đặc tuyến truyền đạt có ngưỡng, $g$ là ức chế chung từ \nt{GABA}, tăng lên khi số nơ-ron hoạt động nhiều hơn $pN$. Thừa số $\frac{1+a}{2}$ là sự tăng cường đầu vào, $1-a$ là sự làm yếu liên kết nội bộ, $\eta a$ là tốc độ học. Phương trình thứ hai là quy tắc Hebb~\eqref{eq:hebb} ở dạng tiện cho các mẫu thưa~\cite{Tsodyks1988}: trọng số tăng khi cả hai nơ-ron hoạt động mạnh hơn bình thường, và giảm khi chỉ một nơ-ron hoạt động. Phần trọng số âm, như mọi khi, do các trung gian \nt{GABA} đảm nhận. Không có đồng hồ: cả nơ-ron lẫn trọng số đều thay đổi liên tục.

\begin{figure}[t]
\centering
\begin{tikzpicture}[>=Stealth,
  nrn/.style={circle,draw,very thick,fill=blue!6,minimum size=0.8cm,inner sep=0pt,font=\small},
  acet/.style={circle,draw,very thick,double,double distance=1.2pt,fill=orange!15,minimum size=1.35cm,inner sep=0pt,font=\scriptsize},
  bc/.style={->,thick,densely dashed,orange!85!black}]
\foreach \i/\y in {1/2.4,2/1.2,3/0}{
  \node[nrn] (N\i) at (3,\y) {$r_\i$};
  \draw[->,thick,blue!60!black] (0,\y) node[left,black] {\small $x_\i$} -- (N\i);}
\node[left,align=right,font=\small] at (-0.5,1.2) {từ các\\giác quan};
\draw[->,thick,gray!60!black] (N1) to[bend left=30] (N2);
\draw[->,thick,gray!60!black] (N2) to[bend left=30] (N1);
\draw[->,thick,gray!60!black] (N2) to[bend left=30] (N3);
\draw[->,thick,gray!60!black] (N3) to[bend left=30] (N2);
\draw[->,thick,gray!60!black] (N1.east) .. controls (4.6,2.0) and (4.6,0.4) .. (N3.east);
\node[right,font=\small,gray!40!black,align=left] at (4.7,1.2) {liên kết\\nội bộ $W$};
\node[acet] (A) at (1.2,-1.9) {\nt{ACET}};
\draw[bc] (A) -- (1.6,-0.15);
\node[font=\small,orange!60!black,anchor=east] at (0.95,-0.9) {$+$ đầu vào};
\draw[bc] (A) -- (4.3,0.6);
\node[font=\small,orange!60!black,anchor=west] at (4.3,0.1) {$-$ liên kết nội bộ, $+$ tốc độ học};
\node[right,font=\small,align=left] at (2.2,-1.9) {mức cao: ``ghi''\\mức thấp: ``đọc''};
\end{tikzpicture}
\caption{\nt{ACET} là đường cho phép ghi. Các nơ-ron $r_i$ nhận đầu vào $x_i$ từ các giác quan (mũi tên xanh) và kích thích lẫn nhau qua các liên kết nội bộ $W$ (màu xám), nơi lưu các mẫu. Acetylcholine (mũi tên nét đứt) tăng cường đầu vào, làm yếu liên kết nội bộ và làm trọng số dễ thay đổi hơn, như trong~\eqref{eq:acetnet}.}
\label{fig:acetnet}
\end{figure}

\section{Ghi và đọc cản trở nhau}

Hãy kiểm tra mô hình~\eqref{eq:acetnet} trên một bài toán mà ở đó ghi và đọc thực sự xung đột. Một mạng $200$ nơ-ron đã lưu năm mẫu, mỗi mẫu $20$ nơ-ron hoạt động. Giờ cần ghi nhớ mẫu thứ sáu, trùng một nửa với một mẫu cũ: chúng có chung mười nơ-ron. Chuyện này xảy ra liên tục: căn phòng mới giống một căn phòng quen, khuôn mặt mới giống một khuôn mặt cũ.

\emph{Ghi khi $a$ thấp.} Trước hết, ta tách hai tác động đầu của acetylcholine khỏi tác động thứ ba: giữ tốc độ học ở mức đầy đủ, và chỉ để sự tăng cường đầu vào cùng sự làm yếu liên kết nội bộ thay đổi. Khi $a=0$, đầu vào làm sáng mười nơ-ron chung, và chúng, qua các liên kết nội bộ, làm sáng nửa còn lại của mẫu \emph{cũ}. \nt{GABA} không cho cả hai cùng sáng, và mẫu cũ, được chính các liên kết của nó nâng đỡ, đẩy lùi mẫu mới, vốn chỉ đứng được nhờ đầu vào. Mạng không thấy cái đang ở trước mắt mà thấy cái nó nhớ --- và quy tắc Hebb cần mẫn ghi lại những gì mạng thấy.

Kết quả là mẫu mới hoàn toàn không được ghi nhớ: từ nửa đặc trưng của nó, mạng không nhớ lại được gì. Còn mẫu cũ thì bị hỏng: khi được gợi bằng nửa đặc trưng của mình, nó kéo theo trung bình sáu trong mười nơ-ron lạ. Khi $a=0.1$, mẫu mới đã được ghi nhớ, nhưng nó hòa lẫn với mẫu cũ, và hư hỏng còn nặng hơn: mỗi mẫu trong hai mẫu, khi được gợi bằng nửa của mình, kéo theo khoảng tám nơ-ron lạ; từ $a=0.2$ trở lên, việc ghi là sạch: số nơ-ron thừa khi nhớ lại ít hơn một (trung bình qua mười lần chạy).

\emph{Ghi với tốc độ học thực.} Giờ để acetylcholine, như trong~\eqref{eq:acetnet}, đặt cả tốc độ học. Sau một lần trình bày, mẫu mới chỉ được ghi nhớ trọn vẹn khi $a\ge0.9$; khi $a=0.75$ --- được tám phần mười, khi $a\le0.5$ hoàn toàn không được ghi nhớ.

\emph{Đọc.} Đưa vào mạng có năm mẫu được ghi sạch một nửa mẫu, rồi bỏ nó đi. Khi $a\le0.2$, mạng điền nốt toàn bộ mẫu và tự giữ nó; khi $a=0.3$ --- gần như toàn bộ; khi $a\ge0.4$, các liên kết nội bộ bị làm yếu quá mức, và sau khi gợi ý biến mất, hoạt động tắt dần (hình~\ref{fig:acetsim}).

\begin{figure}[t]
\centering
\begin{tikzpicture}[>=Stealth,x=9cm,y=3.2cm]
\draw[->,gray!70!black] (0,0) -- (1.1,0) node[right,black] {\small $a$};
\draw[->,gray!70!black] (0,0) -- (0,1.2);
\foreach \x in {0,0.2,0.4,0.6,0.8,1.0}{\draw[gray!60!black] (\x,-0.02) -- (\x,0.02); \node[below,font=\scriptsize] at (\x,0) {$\x$};}
\foreach \y in {0,0.5,1}{\draw[gray!60!black] (-0.01,\y) -- (0.01,\y); \node[left,font=\scriptsize] at (0,\y) {$\y$};}
\node[rotate=90,font=\small] at (-0.1,0.6) {tỷ lệ mẫu được khôi phục};
\draw[very thick,blue!60!black] plot[mark=*,mark size=1.6pt] coordinates {(0,1.00) (0.1,1.00) (0.2,1.00) (0.3,0.96) (0.4,0.04) (0.5,0.00) (0.75,0.00) (1.0,0.00)};
\draw[very thick,red!70!black] plot[mark=*,mark size=1.6pt] coordinates {(0.1,0.00) (0.2,0.00) (0.3,0.00) (0.5,0.00) (0.75,0.80) (0.9,1.00) (1.0,1.00)};
\node[font=\small,blue!60!black,anchor=west,align=left] at (0.03,0.5) {đọc:\\mẫu từ\\một nửa};
\node[font=\small,red!70!black,anchor=east,align=right] at (1.0,0.35) {ghi:\\mẫu mới\\sau một lần};
\end{tikzpicture}
\caption{Mô hình~\eqref{eq:acetnet}: $200$ nơ-ron, mỗi mẫu $20$ nơ-ron hoạt động, trung bình qua mười lần chạy. Các điểm xanh là đọc: tỷ lệ mẫu được khôi phục từ một nửa của nó sau khi bỏ gợi ý. Các điểm đỏ là ghi: tỷ lệ của mẫu mới, giống một nửa với mẫu cũ, được nhớ lại sau một lần trình bày, nếu acetylcholine đặt cả tốc độ học. Đọc hoạt động khi $a\le0.3$, ghi --- khi $a\ge0.9$.}
\label{fig:acetsim}
\end{figure}

\begin{keyblock}
\begin{proposition}[Ghi và đọc đòi hỏi những chế độ khác nhau]
\label{prop:writeread}
Trong mô hình~\eqref{eq:acetnet}, nơi cùng những liên kết vừa lưu cái cũ vừa học cái mới, còn acetylcholine đặt cả tốc độ học, không có mức $a$ nào mà ở đó cả ghi lẫn đọc đều tốt như nhau. Để ghi, các liên kết nội bộ phải được làm yếu, nếu không mạng sẽ ghi không phải đầu vào mà là điều nó nhớ lại; để đọc, chúng phải được tăng cường, nếu không mạng không điền nốt được mẫu từ một phần. Cần một công tắc, và một dây quảng bá có thể làm công tắc đó.
\end{proposition}
\end{keyblock}

Nếu acetylcholine không thay đổi tốc độ học, thì cho cả hai việc sẽ có một dải hẹp $0.2\le a\le0.3$. Nhưng khi đó mạng sẽ học cả trong mỗi lần đọc, tức là ghi lại những gì đã đọc cùng với các lỗi của nó. Chính ý tưởng --- ức chế các liên kết nội bộ trong lúc học --- được lấy từ các mô hình xây dựng dựa trên phép đo trên lát cắt~\cite{HasselmoAndersonBower1992}. Các con số ở trên thuộc về mô hình đơn giản hóa của chúng tôi; ranh giới giữa các chế độ trong não nằm chính xác ở đâu thì chưa rõ.

\section{Tin vào mắt bao nhiêu}

Công tắc ``ghi/đọc'' là trường hợp cực đoan của một câu hỏi tổng quát hơn: nên tin vào đầu vào bao nhiêu và tin vào bộ nhớ bao nhiêu? Kỹ sư có câu trả lời chính xác cho câu hỏi này, và nó làm sáng tỏ vì sao một dây điều khiển cùng lúc cả chế độ lẫn tốc độ học.

Giả sử mạng theo dõi một đại lượng $s(t)$ trôi dạt chậm: mỗi giây phương sai của nó tăng thêm $q$. Các giác quan báo về $y(t)=s(t)+$nhiễu với cường độ $R$. Ước lượng tốt nhất $\hat s$ trong thời gian liên tục được cho bởi bộ lọc Kalman--Bucy~\cite{KalmanBucy1961}:
\begin{empheq}[box=\keybox]{equation}
\frac{\dd\hat s}{\dd t} \;=\; \frac{P}{R}\,\bigl(y-\hat s\bigr),
\qquad
\frac{\dd P}{\dd t} \;=\; q-\frac{P^2}{R} ,
\label{eq:kalman}
\end{empheq}
trong đó $P$ là phương sai sai số của ước lượng, tức là mức độ mà chính mạng không chắc chắn về điều nó biết. Phương trình thứ nhất là cùng mạch RC như màng trong mô hình nơ-ron~\eqref{eq:glutneuron}, nhưng với hằng số thời gian $R/P$ thay đổi. Khi mạng không chắc chắn, $P$ lớn, hằng số thời gian nhỏ, và ước lượng bám nhanh theo đầu vào: đó là ``ghi''. Khi mạng chắc chắn, $P$ nhỏ, và ước lượng gần như không nghe đầu vào, bám vào bộ nhớ: đó là ``đọc''.

Ở chế độ xác lập $P=\sqrt{qR}$ và hằng số thời gian của bộ nhớ bằng $\sqrt{R/q}$: nếu thế giới trôi đi một đơn vị trong một trăm giây, $q=0.01$, còn nhiễu giác quan là $R=1$, thì bộ nhớ nên nhớ khoảng mười giây.

Điều chính ở đây là một con số $P/R$ đặt cùng lúc hai thứ: cả trọng số của đầu vào so với bộ nhớ, lẫn tốc độ thay đổi của ước lượng được lưu. Đây đúng là điều acetylcholine làm trong~\eqref{eq:acetnet}. Theo giả thuyết~\cite{YuDayan2005}, acetylcholine báo cho mạng một đại lượng giống như $P$ --- sự bất định \emph{được dự liệu}, loại mà mạng biết trước. Kênh này không phải lúc nào cũng chậm: một phần nơ-ron \nt{ACET} đáp ứng với phần thưởng và hình phạt trong khoảng hai chục mili giây, càng mạnh khi sự củng cố càng bất ngờ~\cite{Hangya2015}.

\begin{keyblock}
\begin{proposition}[Một núm vặn cho trọng số đầu vào và tốc độ học]
\label{prop:kalman}
Trong bộ lọc tối ưu~\eqref{eq:kalman}, trọng số dùng để trộn đầu vào với bộ nhớ và tốc độ học là cùng một đại lượng $P/R$. Vì vậy điều khiển chúng bằng một tín hiệu là hợp lý. Khi các tính chất của thế giới không đổi, tín hiệu này đạt một mức ổn định $\sqrt{q/R}$, và chỉ cần điều chỉnh nó khi các tính chất của thế giới thay đổi.
\end{proposition}
\end{keyblock}

Câu thứ hai của mệnh đề giải thích kết quả ở chương~\ref{sec:nora}, nơi việc điều chỉnh tốc độ học theo sự bất ngờ không thắng được tốc độ cố định tốt nhất: trong một thế giới mà tính chất của các thay đổi không đổi, tốc độ học tốt nhất chính là một hằng số. Lợi ích xuất hiện khi thế giới đột ngột trở nên khác. Khi đó ước lượng bỗng ở xa đầu vào, còn $P$ thì không biết điều đó. Hành động đúng là tăng mạnh $P$ và qua đó chuyển sang chế độ ghi. Theo giả thuyết mà ta đã bàn ở chương~\ref{sec:nora}, thay đổi \emph{không được dự liệu} như vậy do \nt{NORA} báo~\cite{YuDayan2005}. Sự phân công khi đó là tự nhiên: \nt{NORA} nhận ra thế giới đã thay đổi, \nt{ACET} giữ mạng ở chế độ ghi cho tới khi bối cảnh mới được học xong.

\section{Giấc ngủ là chế độ đọc}

Nếu mức acetylcholine cao là chế độ ghi, thì mạng làm gì khi mức thấp? Trong giấc ngủ sóng chậm, mức acetylcholine giảm, các liên kết nội bộ thoát khỏi sự ức chế, đầu vào từ các giác quan gần như không có. Mạng ở chế độ đọc mà không có gợi ý sẽ phát lại những gì đã được ghi trong nó. Ghi nhận hoạt động cho thấy những nơ-ron cùng hoạt động ban ngày lại cùng phát xung trong giấc ngủ sóng chậm~\cite{WilsonMcNaughton1994}, thậm chí theo cùng thứ tự~\cite{Skaggs1996}. Theo giả thuyết~\cite{Hasselmo1999}, những lần phát lại này chép những gì học nhanh ban ngày sang bộ nhớ dài hạn (kéo dài nhân tạo các đợt phát lại ngắn làm cải thiện trí nhớ~\cite{FernandezRuiz2019}): ban ngày --- ghi từ đầu vào, ban đêm --- ghi lại từ bên trong. Với kỹ sư, điều này giống như ghi vào một bộ đệm nhanh ban ngày và đổ nó sang bộ nhớ lâu dài chậm vào ban đêm.

\section{Tổng kết}

\begin{table}[t]
\centering
\footnotesize
\setlength{\tabcolsep}{4pt}
\renewcommand{\arraystretch}{1.15}
\begin{tabular}{>{\raggedright}p{3.2cm}>{\raggedright}p{4.2cm}>{\raggedright\arraybackslash}p{6.0cm}}
\toprule
\nt{ACET} làm gì & Bằng cách nào & Điều đã biết \\
\midrule
làm yếu liên kết nội bộ & thừa số $1-a$ trong~\eqref{eq:acetnet} & đã đo trên lát cắt \\
tăng cường đầu vào & thừa số $\frac{1+a}{2}$ & sự tăng tính dễ kích thích và truyền dẫn qua các đường vào đã được đo \\
tạo thuận cho việc học & tốc độ $\eta a$ & đã đo \\
chuyển ``ghi/đọc'' & mệnh đề~\ref{prop:writeread} & suy ra từ các mô hình; chưa có phép đo trực tiếp các chế độ \\
báo sự bất định được dự liệu & $P$ trong~\eqref{eq:kalman} & giả thuyết \\
giải phóng mạng cho những lần phát lại trong giấc ngủ & $a$ thấp trong giấc ngủ sóng chậm & mức thấp khi ngủ và việc phát lại đã được đo; việc chép sang bộ nhớ dài hạn là giả thuyết \\
\bottomrule
\end{tabular}
\caption{\nt{ACET} bổ sung gì cho mạng gồm \nt{GLUT}, \nt{GABA}, \nt{DOPA} và \nt{NORA}.}
\label{tab:acet}
\end{table}

\nt{DOPA} cho mạng biết thay đổi trọng số theo hướng nào, \nt{NORA} --- dùng những gì nó biết quyết đoán tới mức nào, còn \nt{ACET} --- nên nghe thế giới hay nghe chính mình (bảng~\ref{tab:acet}). Đây là đường điều khiển quảng bá cuối cùng trong bảng~\ref{tab:types}: tín hiệu cho phép ghi, thiếu nó thì một bộ nhớ lưu các mẫu trong chính những liên kết mà nó dùng để đọc chúng sẽ tự làm hỏng mình sau mỗi lần đọc.

\section{Bài tập}

\begin{exercise}[\lvlA]
\label{ex:09:1}
\emph{Nhớ bao lâu.} Bộ lọc~\eqref{eq:kalman} với $q=0.01$, $R=1$ nhớ khoảng $10$~s. Tìm $P_\infty$ và thời gian nhớ $\sqrt{R/q}$, nếu (a)~biên độ nhiễu cảm biến tăng gấp đôi; (b)~thế giới trôi dạt nhanh gấp một trăm lần. (c)~Nhiễu phải tăng bao nhiêu lần để mạng nhớ được một phút?
\end{exercise}

\begin{exercise}[\lvlB]
\label{ex:09:2}
\emph{Chế độ ghi sau một thay đổi.} Thế giới đã thay đổi, và độ bất định của bộ lọc~\eqref{eq:kalman} với $q=0.01$, $R=1$ nhảy vọt lên $P_0\to\infty$. (a)~$P$ trở về $P_\infty$ với hằng số thời gian nào? (b)~Sau bao nhiêu giây tốc độ học vượt mức xác lập không quá $10\,\%$?
\end{exercise}

\begin{exercise}[\lvlB]
\label{ex:09:3}
\emph{Tốc độ học cố định.} Mạng học theo $\dd\hat s/\dd t=(y-\hat s)/T$; thế giới trôi dạt, $\dd s/\dd t=w$, $y=s+n$, trong đó $w$ và $n$ là nhiễu trắng với cường độ $q$ và $R$. Tìm $T$ tốt nhất và phương sai sai số khi đó. Nó tăng bao nhiêu lần nếu $T$ sai $2$ lần và $10$ lần?
\end{exercise}

\begin{exercise}[\lvlB]
\label{ex:09:4}
\emph{Ranh giới ghi ở đâu.} Trong \texttt{models/\allowbreak acet\_memory.py}, ngưỡng $\theta=0.35$, trọng số bên trong mẫu $w=0.041$ (lý tưởng là $0.045$). Mẫu mới có chung $m=10$ nơ-ron với mẫu cũ. Với $a$ nào trong~\eqref{eq:acetnet} các nơ-ron còn lại của mẫu cũ không bị $m$ nơ-ron này làm sáng lên (ngưỡng sắc)?
\end{exercise}

\begin{exercise}[\lvlB]
\label{ex:09:5}
\emph{Mô phỏng: dung lượng bộ nhớ.} Trong \texttt{models/\allowbreak acet\_memory.py}, sửa hàm \texttt{read\_sweep}: với $a=1$, ghi $M$ mẫu ($M$ từ $5$ tới $100$), sau đó với $a=0$ đọc mười mẫu đầu tiên từ một nửa của chúng. Với $M$ nào bắt đầu xuất hiện lỗi, và giới hạn $M_{\max}$ phụ thuộc vào $N$ và $p$ thế nào?
\end{exercise}

\begin{exercise}[\lvlB]
\label{ex:09:6}
\emph{Thiết kế: ghi không rò rỉ.} Với tốc độ học $\eta a$, mạng học cả khi đọc. Đề xuất một hàm đơn điệu $\eta(a)\le\eta$, bằng không khi $a\le0.3$ và không kém $\eta a$ khi $a\ge0.9$. Kiểm tra: sau $20$ lần đọc với $a=0.2$ và gợi ý có ba nơ-ron lạ, bao nhiêu nơ-ron lạ đã lọt vào mẫu?
\end{exercise}

\begin{exercise}[\lvlC]
\label{ex:09:7}
\emph{Chép bộ đệm vào ban đêm thế nào.} (a)~Khi ngủ $a=0.05$, mỗi lần phát lại dài $0.1\,\text{s}$, ban ngày $0.2\,\text{s}$ với $a=1$. Bao nhiêu lần phát lại tích lũy đủ thay đổi trọng số của một lần trình bày ban ngày, với tốc độ $\eta a$ và với $\eta(a)$ ở bài tập~\ref{ex:09:6}? (b)~Kho lưu trữ học chậm hơn bộ đệm $100$ lần và nghe mẫu $20$ lần mỗi đêm, mỗi lần $0.1\,\text{s}$. Cần bao nhiêu đêm?
\end{exercise}

\chapter{Toàn bộ cỗ máy}
\chaptermark{Toàn bộ cỗ máy}
\label{sec:synthesis}

\section{Chúng ta đã lắp được gì}

Chúng ta lần lượt lấy từng loại nơ-ron trong bảng~\ref{tab:types} và hỏi nó bổ sung gì cho một cỗ máy tính toán. Lần nào luật chơi cũng như nhau: chỉ dùng những gì có trong cơ thể sống, và không có xung nhịp chung. Giờ hãy ráp các phần lại thành một cỗ máy và xem chúng phối hợp với nhau ra sao.

Bức tranh ở chương~\ref{sec:neurontypes} đã được khẳng định và trở nên chính xác hơn. Một mạng địa chỉ khổng lồ gồm các nơ-ron \nt{GLUT} mang dữ liệu. Các nơ-ron \nt{GABA} điều khiển luồng dữ liệu đó. Còn phía trên mạng là bốn kênh quảng bá, mỗi kênh có núm vặn riêng: \nt{DOPA} cho biết nên giữ lại những thay đổi trọng số nào, \nt{SERO} giữ mức chung và, theo giả thuyết, cả tầm nhìn kế hoạch, \nt{NORA} chọn giữa tìm kiếm và quyết định, \nt{ACET} chọn giữa ghi và đọc (hình~\ref{fig:machine}).

\begin{figure}[t]
\centering
\begin{tikzpicture}[>=Stealth,
  blk/.style={draw,very thick,rounded corners=3pt,align=center,font=\small},
  reg/.style={draw,very thick,double,double distance=1.2pt,circle,minimum size=1.25cm,inner sep=0pt,font=\scriptsize,fill=orange!15},
  bc/.style={->,thick,densely dashed,orange!85!black}]
% sensors, bus, actions
\node[blk,fill=green!8,text width=1.7cm] (S) at (0.9,0) {cảm biến\\\scriptsize bật/tắt};
\draw[very thick,rounded corners=4pt,fill=blue!5] (2.6,-1.35) rectangle (9.0,1.35);
\node[font=\small] at (5.8,0.95) {bus \nt{GLUT}: dữ liệu};
\node[font=\scriptsize,align=center] at (4.3,0.25) {trùng khớp, trễ,\\logic (ch.~\ref{sec:glutcomputer})};
\node[font=\scriptsize,align=center] at (7.4,0.25) {bộ nhớ, mã\\(ch.~\ref{sec:code}, \ref{sec:acet})};
\draw[thick,red!70!black,fill=red!8,rounded corners=2pt] (2.8,-1.15) rectangle (8.8,-0.55);
\node[font=\scriptsize,red!60!black] at (5.8,-0.85) {\nt{GABA} (ch.~\ref{sec:gaba}): cấm, chọn, chia, nhịp};
\node[blk,fill=blue!5,text width=2.0cm] (A) at (10.9,0) {chọn\\hành động};
\draw[->,very thick] (S) -- (2.6,0);
\draw[->,very thick] (9.0,0) -- (A);
\draw[->,very thick] (A) -- (12.6,0) node[right,font=\small] {cơ};
\pic[scale=1.1] at (13.2,-0.5) {spring};
\pic[scale=1.15] at (0.4,0.85) {eyepic};
\pic[scale=1.05] at (1.35,0.85) {speaker};
% registers
\node[reg] (D) at (3.4,3.2) {\nt{DOPA}};
\node[reg] (C) at (5.6,3.2) {\nt{SERO}};
\node[reg] (N) at (7.8,3.2) {\nt{NORA}};
\node[reg] (H) at (10.0,3.2) {\nt{ACET}};
\node[font=\scriptsize,align=center,above=1pt] at (D.north) {sai số $\delta$};
\node[font=\scriptsize,align=center,above=1pt] at (C.north) {mức, $\tau_\gamma$};
\node[font=\scriptsize,align=center,above=1pt] at (N.north) {khuếch đại $g$};
\node[font=\scriptsize,align=center,above=1pt] at (H.north) {ghi $a$};
\draw[bc] (D.south) -- (3.4,1.35);
\draw[bc] (C.south) -- (5.6,1.35);
\draw[bc] (N.south) -- (7.8,1.35);
\draw[bc] (H.south) -- (10.0,2.1) -- (8.8,1.35);
\draw[bc] (H.south east) to[bend left=20] (A.north east);
\draw[->,thick] (2.85,1.35) -- (2.85,2.75) -- (D.west) node[pos=0.25,left,font=\scriptsize] {$V$};
\draw[->,thick,green!45!black] (0.4,3.2) node[left,black,font=\small] {thưởng $r$} -- (2.2,3.2) -- (D.west);
\pic[scale=0.9] at (-0.45,3.7) {juice};
\draw[->,thick,densely dotted,gray!50!black] ([yshift=0.45cm]D.north) to[bend left=18] node[above,font=\scriptsize,black] {bất ngờ} ([yshift=0.45cm]N.north);
\end{tikzpicture}
\caption{Toàn bộ cỗ máy. Bus \nt{GLUT} mang dữ liệu từ cảm biến đến khâu chọn hành động; \nt{GABA} điều khiển luồng bên trong nó. Các vòng tròn kép là kênh quảng bá; mũi tên nét đứt là việc phát tín hiệu của chúng tới toàn mạng, mỗi kênh có núm vặn riêng. \nt{DOPA} nhận phần thưởng $r$ và kỳ vọng $V$ từ bộ phê bình nằm trong bus (ch.~\ref{sec:dofa}); mũi tên chấm là giả thuyết rằng \nt{NORA} biết về sự bất ngờ nhờ các sai số lớn bất thường (ch.~\ref{sec:nora}).}
\label{fig:machine}
\end{figure}

\section{Một công thức cho mọi núm vặn}

Có thể gom các núm vặn vào một công thức sơ đồ. Đây không phải mô hình mới mà là bản tóm tắt các mô hình của chương~\ref{sec:glutcomputer}--\ref{sec:acet}: mỗi ký hiệu lấy từ chương của nó.
\begin{empheq}[box=\keybox]{equation}
\begin{aligned}
\tau\,\frac{\dd r_i}{\dd t} \;&=\; -r_i + F\Bigl(g\,\Bigl[\tfrac{1+a}{2}\,x_i + (1-a)\sum_j W_{ij}\,r_j - \sum_k M_{ik}\,q_k - \theta\Bigr]\Bigr),\\
\frac{\dd W_{ij}}{\dd t} \;&=\; \eta\,a\,\delta(t)\,e_{ij}(t),
\qquad \tau_e\,\frac{\dd e_{ij}}{\dd t} = -e_{ij} + r_i\,r_j,
\qquad \delta = r + \frac{\dd V}{\dd t} - \frac{V}{\tau_\gamma} .
\end{aligned}
\label{eq:machine}
\end{empheq}
Ở đây $r_i$ là tần số phát xung của các nơ-ron \nt{GLUT}, $x_i$ là đầu vào từ cảm biến, $W_{ij}\ge0$ là trọng số giữa chúng (chương~\ref{sec:glutcomputer}); $q_k$ là tần số của các nơ-ron \nt{GABA}, $M_{ik}\ge0$ là cường độ ức chế của chúng (chương~\ref{sec:gaba}); $e_{ij}$ là vết đủ điều kiện, $\delta$ là sai số \nt{DOPA} (chương~\ref{sec:dofa}); $\tau_\gamma$ là tầm nhìn, theo giả thuyết do \nt{SERO} đặt; $g$ là hệ số khuếch đại \nt{NORA} (chương~\ref{sec:nora}); $a$ là mức \nt{ACET} (chương~\ref{sec:acet}).

Hãy xem những gì không có trong~\eqref{eq:machine}. Không có nhịp đồng hồ: mọi thứ được viết bằng phương trình vi phân, và mỗi nơ-ron tự thay đổi. Không có bộ nhớ riêng: chính các $W_{ij}$ dùng để tính toán cũng làm nhiệm vụ ghi nhớ, cùng với chính hoạt động $r_i$. Không có hiệu chỉnh riêng cho tuyệt đại đa số khớp thần kinh: việc học của chúng là tích của một vết cục bộ và một con số chung duy nhất; dây sai số riêng tới từng đầu ra chỉ có trong mạch học vận động (chương~\ref{sec:motorlearning}). Và chỉ có bốn loại tín hiệu điều khiển, $\delta$, $\tau_\gamma$, $g$, $a$, cho tám mươi tỷ nơ-ron. Thật ra mỗi tín hiệu không phải một con số mà là một bó nhỏ các đường song song (mệnh đề~\ref{prop:bundle}), nhưng mỗi bó chỉ có vài đường.

\begin{keyblock}
\begin{proposition}[Kiến trúc]
\label{prop:architecture}
Cỗ máy não, trong chừng mực hiểu biết hiện nay, có thể mô tả bằng ba lớp. Dữ liệu chạy trên bus địa chỉ \nt{GLUT}. Các nơ-ron \nt{GABA} có địa chỉ định hướng, giới hạn và chuẩn hóa luồng dữ liệu. Chế độ làm việc do vài kênh quảng bá đặt ra, mỗi kênh là một bó nhỏ các đường song song, dùng chung cho những vùng mạng lớn. Hai lớp đầu đã được xác lập vững chắc; việc phân vai giữa các kênh quảng bá phần lớn vẫn là giả thuyết.
\end{proposition}
\end{keyblock}

\section{Mọi loại nơ-ron trong một bảng}

Bảng~\ref{tab:synthesis} tổng hợp mọi loại nơ-ron trong sách. Bốn loại không có chương riêng mỗi loại chiếm một dòng; hai trong số đó đã tìm được vai trò trong các chương về đầu ra của cỗ máy: \nt{GLYC} trong các vòng tủy sống (chương~\ref{sec:muscles}), \nt{ADRE} trong đầu ra hóa học (chương~\ref{sec:chemoutput}). Vai trò tính toán của hai loại còn lại hoặc đơn giản, hoặc chưa rõ. Các chất không quy định loại nơ-ron được tập hợp ở phụ lục~\ref{sec:othermediators}.

\begin{table}[t]
\centering
\footnotesize
\setlength{\tabcolsep}{3pt}
\renewcommand{\arraystretch}{1.15}
\begin{tabular}{>{\raggedright}p{1.3cm}>{\raggedright}p{1.0cm}>{\raggedright}p{3.9cm}>{\raggedright}p{1.5cm}>{\raggedright}p{1.8cm}>{\raggedright\arraybackslash}p{3.8cm}}
\toprule
Loại & Chương & Tính gì & Thời gian & Bus & Đã biết gì \\
\midrule
\nt{GLUT} & \ref{sec:glutcomputer}, \ref{sec:code}, \ref{sec:sensors}, \ref{sec:motorlearning}, \ref{sec:dendrites} & dữ liệu: trùng khớp, trễ, logic, bộ nhớ, chuỗi & ms & địa chỉ & đã xác lập \\
\nt{GABA} & \ref{sec:gaba}, \ref{sec:code}, \ref{sec:pain}, \ref{sec:motorlearning} & điều khiển luồng: cấm, chọn, chia, ổn định, nhịp; cổng đau & ms & địa chỉ & đã xác lập; chia tần số: cùng với nhiễu nền \\
\nt{SERO} & \ref{sec:serocomputer}, \ref{sec:pain}, \ref{sec:sleep} & điều chỉnh mức, làm trơn; tầm nhìn $\tau_\gamma$; độ lớn của cơn đau; các chế độ ngủ & giây & thể tích & tự ức chế đã đo; tầm nhìn là giả thuyết \\
\nt{DOPA} & \ref{sec:dofa}, \ref{sec:dofaalt} & sai số dự đoán $\delta$; mức chậm là giá của thời gian & $0.1$\,s và phút & quảng bá & $\delta$ đã đo; mở rộng: ch.~\ref{sec:dofaalt} \\
\nt{NORA} & \ref{sec:nora}, \ref{sec:pain}, \ref{sec:chemoutput}, \ref{sec:sleep} & khuếch đại $g$: tìm kiếm hay quyết định; bất ngờ; ``chân ga'' cho các cơ quan & giây & quảng bá & tăng tỷ số tín hiệu/nhiễu đã đo; vai trò trong tìm kiếm là giả thuyết \\
\nt{ACET} & \ref{sec:acet}, \ref{sec:muscles}, \ref{sec:chemoutput}, \ref{sec:sleep} & ghi hay đọc; tốc độ học; đầu ra tới cơ; ``phanh'' cho các cơ quan & giây & cả hai & ba hiệu ứng đã đo; chuyển chế độ là giả thuyết \\
\nt{GLYC} & \ref{sec:muscles}, \ref{sec:pain} & trọng số âm ở tủy sống và thân não: cấm cơ đối vận & ms & địa chỉ & đã xác lập \\
\nt{HIST} & --- & tín hiệu toàn cục ``hãy thức'' & chậm & quảng bá & vai trò tính toán chưa được nghiên cứu \\
\nt{ADRE} & \ref{sec:chemoutput} & ``chân ga'' cho toàn thân qua máu; trong não chưa rõ & chậm & quảng bá & ngoài não đã xác lập; vai trò trong não chưa rõ \\
\nt{ASPA} & --- & trọng số dương yếu & ms & địa chỉ & vai trò còn tranh cãi \\
\bottomrule
\end{tabular}
\caption{Mọi loại nơ-ron trong sách: mỗi loại tính gì và điều đó đã được biết đến mức nào.}
\label{tab:synthesis}
\end{table}

\section{Các núm vặn phối hợp ra sao}

Cho đến giờ, mỗi chương vặn một núm. Bây giờ hãy theo dõi một sự kiện đi qua toàn bộ cỗ máy (hình~\ref{fig:timeline}). Một con vật từ lâu đã học được rằng hành động $A$ mang lại nước quả. Đến một lúc thế giới thay đổi: $A$ không còn mang lại nước quả nữa, mà hành động $B$ mới làm được điều đó, trong khi con vật hầu như không thử $B$.

\emph{Sai số.} Sau $A$ không có nước quả, và các nơ-ron \nt{DOPA} đáp lại bằng một cú sụt, $\delta<0$ theo công thức sai số~\eqref{eq:ctd}. Các khớp thần kinh vừa chọn $A$ đang mang vết và bị suy yếu.

\emph{Bất ngờ.} Một cú sụt chỉ là chuyện ngẫu nhiên bình thường. Nhưng các cú sụt lặp lại, và sai số trở nên lớn hơn thường lệ. Theo giả thuyết của chương~\ref{sec:nora}, đó chính là tín hiệu cho \nt{NORA}: thế giới đã thay đổi. Hệ số khuếch đại $g$ giảm, lựa chọn trở nên mềm, và nhiễu thường xuyên dẫn tới việc thử $B$ hơn.

\emph{Ghi.} Theo giả thuyết của chương~\ref{sec:acet}, sau thay đổi \nt{ACET} tăng lên và đưa mạng vào chế độ ghi: các liên kết bên trong lưu giữ thói quen cũ bị làm yếu, đầu vào từ cảm biến được tăng cường, trọng số thay đổi dễ dàng.

\emph{Thói quen mới.} Lần thử $B$ mang lại nước quả mà không ai chờ đợi: một đỉnh vọt, $\delta>0$, và các khớp thần kinh đã chọn $B$ được tăng cường. Vài lần thử như vậy là kỳ vọng $V$ bắt kịp thế giới mới, sai số lại nhỏ.

\emph{Trở lại.} Sự bất ngờ qua đi, \nt{NORA} trả lại hệ số khuếch đại cao, lựa chọn lại cứng; \nt{ACET} hạ xuống, mạng ở chế độ đọc và dùng những gì đã học. Suốt thời gian đó, \nt{SERO}, theo giả thuyết, đặt tầm nhìn $\tau_\gamma$: phần nước quả ở xa đến mức nào thì vẫn còn đáng chờ.

\begin{figure}[t]
\centering
\begin{tikzpicture}[>=Stealth,x=1.1cm,y=0.95cm]
\foreach \y/\lab in {4/{$\delta$ (\nt{DOPA})},3/{$g$ (\nt{NORA})},2/{$a$ (\nt{ACET})},1/{chọn $B$}}{
  \draw[gray!60!black] (0,\y-0.35) -- (10,\y-0.35);
  \node[left,font=\scriptsize] at (0,\y) {\lab};}
\draw[densely dashed,gray!60!black] (2,0.4) -- (2,4.55) node[above,font=\scriptsize,black] {thế giới đổi};
\draw[densely dashed,gray!60!black] (5,0.4) -- (5,4.55) node[above,font=\scriptsize,black] {nước quả đầu tiên nhờ $B$};
\pic[scale=0.85] at (2,5.25) {nojuice};
\pic[scale=0.85] at (5,5.25) {juice};
% delta: dips after the change, burst at the first juice for B, then back to zero
\draw[thick,red!70!black] (0,3.65) -- (2.3,3.65) -- (2.3,3.3) -- (2.5,3.3) -- (2.5,3.65) -- (3.2,3.65) -- (3.2,3.3) -- (3.4,3.3) -- (3.4,3.65) -- (4.1,3.65) -- (4.1,3.35) -- (4.3,3.35) -- (4.3,3.65) -- (5.0,3.65) -- (5.0,4.35) -- (5.2,4.35) -- (5.2,3.65) -- (5.8,3.65) -- (5.8,4.1) -- (6.0,4.1) -- (6.0,3.65) -- (6.6,3.65) -- (6.6,3.85) -- (6.8,3.85) -- (6.8,3.65) -- (10,3.65);
% gain: high, drops after surprise, recovers
\draw[thick,blue!60!black] (0,3.2) -- (3.0,3.2) -- (3.4,2.75) -- (5.8,2.75) -- (7.2,3.2) -- (10,3.2);
% ACh: low, rises after the change, falls after learning
\draw[thick,green!45!black] (0,1.75) -- (2.8,1.75) -- (3.3,2.25) -- (6.8,2.25) -- (7.8,1.75) -- (10,1.75);
% choice of B
\draw[thick,black] (0,0.7) -- (3.0,0.7) plot[domain=3:10,samples=40] (\x,{0.7+0.6/(1+exp(-(\x-5.3)*1.8))});
\draw[->,gray!70!black] (0,0.2) -- (10.2,0.2) node[below left,font=\scriptsize,black] {thời gian};
\end{tikzpicture}
\caption{Một sự kiện được theo dõi qua toàn bộ cỗ máy. Đây là sơ đồ định tính dựa trên các mô hình và giả thuyết của chương~\ref{sec:dofa}--\ref{sec:acet}, không phải kết quả tính toán. Sau thay đổi, \nt{DOPA} đáp lại bằng các cú sụt; theo giả thuyết, các cú sụt lặp lại làm tăng tín hiệu bất ngờ, \nt{NORA} giảm hệ số khuếch đại, \nt{ACET} bật chế độ ghi; nước quả đầu tiên nhờ $B$ tạo một đỉnh vọt, lựa chọn chuyển sang $B$, và các núm vặn trở về vị trí cũ.}
\label{fig:timeline}
\end{figure}

Lưu ý rằng không núm vặn nào biết các núm khác đang làm gì. Mỗi núm phản ứng với dấu hiệu riêng của nó (sai số, độ lớn bất thường của sai số, sự bất định), và trình tự hành động tự hình thành, như trong một mạch bất đồng bộ, nơi mỗi khối kích hoạt khi các đầu vào của nó sẵn sàng. Theo những gì chúng tôi biết, sự phối hợp trọn vẹn như vậy chưa từng được kiểm chứng: từng núm vặn riêng lẻ đã được đo, còn hoạt động chung của chúng là giả thuyết.

\section{Cỗ máy này khác máy tính ở đâu}

\begin{table}[t]
\centering
\footnotesize
\setlength{\tabcolsep}{4pt}
\renewcommand{\arraystretch}{1.15}
\begin{tabular}{>{\raggedright}p{2.2cm}>{\raggedright}p{4.2cm}>{\raggedright\arraybackslash}p{6.6cm}}
\toprule
& Máy tính thông thường & Não \\
\midrule
thời gian & xung nhịp chung, khoảng một gigahertz & không có đồng hồ; mỗi nơ-ron tự thay đổi, các cửa sổ do sự kiện và nhịp cục bộ đặt ra (ch.~\ref{sec:glutcomputer}, \ref{sec:gaba}, \ref{sec:code}) \\
phần tử & cổng nhị phân tin cậy & phần tử ngưỡng tương tự; khớp thần kinh làm mất phần lớn số xung (ch.~\ref{sec:code}) \\
dấu của liên kết & tùy ý & do nơ-ron gửi quyết định (ch.~\ref{sec:neurontypes}) \\
bộ nhớ & tách khỏi bộ xử lý & nằm trong chính các liên kết dùng để tính toán, và trong hoạt động tự duy trì (ch.~\ref{sec:glutcomputer}, \ref{sec:acet}) \\
học & lan truyền ngược sai số & quy tắc cục bộ, một con số quảng bá $\delta$ (ch.~\ref{sec:dofa}) và dây sai số tới các đầu ra của mạch vận động (ch.~\ref{sec:motorlearning}) \\
điều khiển & thanh ghi và bus lệnh & bốn kênh quảng bá, mỗi kênh là một bó vài đường (ch.~\ref{sec:serocomputer}, \ref{sec:dofa}--\ref{sec:acet}) \\
năng lượng & hàng chục đến hàng trăm oát cho một bộ xử lý & khoảng $20$ oát cho cả bộ não, trung bình khoảng một xung mỗi giây trên một nơ-ron (ch.~\ref{sec:code}) \\
\bottomrule
\end{tabular}
\caption{Não và máy tính thông thường.}
\label{tab:computer}
\end{table}

Bảng~\ref{tab:computer} tổng hợp các khác biệt chính. Mỗi khác biệt không phải là khuyết điểm mà tự nhiên chưa kịp sửa, mà là một phần của cùng một giải pháp. Không có xung nhịp chung thì cần cảm biến ``bật--tắt'' và bộ phát hiện trùng khớp. Phần tử không tin cậy thì cần sự dư thừa của nhiều nơ-ron. Bộ nhớ gộp với tính toán thì cần \nt{ACET} để việc ghi không làm hỏng việc đọc. Học mà không có dây ngược thì cần \nt{DOPA} và nhiễu. Còn giới hạn năng lượng một xung mỗi giây buộc cả ngôn ngữ phải tằn tiện và chỉ tiêu xung cho tin mới.

\section{Những điều chúng ta chưa biết}

Cách trung thực nhất để khép lại bản tổng kết này là liệt kê những gì chúng ta chưa biết. Mỗi mục là một bài toán cho kỹ sư.

\emph{Mã.} Chúng ta biết dung lượng của kênh xung và biết rằng bên nhận chọn đọc phần nào của thông điệp (chương~\ref{sec:code}). Nhưng vỏ não tận dụng thời điểm chính xác của xung đến mức nào và nơ-ron có ``từ'' hay không thì chưa biết.

\emph{Học qua nhiều lớp.} Tín hiệu quảng bá dạy chậm, và càng chậm hơn với mỗi nơ-ron có ảnh hưởng tới kết quả (mệnh đề~\ref{prop:broadcastcost}). Vậy vỏ não chỉnh hàng tỷ trọng số trong các mạch nhiều lớp như thế nào thì chưa biết; có những mô hình trong đó sai số được truyền giữa các lớp bằng các chùm xung~\cite{Payeur2021}. Có thể một phần câu trả lời nằm ở tính toán bên trong các nhánh của chính nơ-ron, nơi một nơ-ron hoạt động như một mạng nhiều lớp thu nhỏ: để lặp lại đáp ứng của mô hình một nơ-ron vỏ não, mạng nhân tạo cần năm đến tám lớp~\cite{Beniaguev2021}, còn các nhánh nơ-ron người thậm chí tính được cả phép HOẶC loại trừ~\cite{Gidon2020}. Một ứng viên khác là tự học: nếu vỏ não học dự đoán chính đầu vào của mình, sai số nảy sinh ngay tại chỗ, trong từng lớp, và không cần tín hiệu chung nào cho việc đó (chương~\ref{sec:dofa})~\cite{Keller2018}.

\emph{Các kênh quảng bá thực sự truyền gì.} Với \nt{DOPA} có bức tranh chuẩn và sáu mở rộng của nó (chương~\ref{sec:dofaalt}); với \nt{NORA} và \nt{ACET} có những giả thuyết hợp lý (chương~\ref{sec:nora}, \ref{sec:acet}); với \nt{SERO} phần lớn chỉ là phỏng đoán.

\emph{Đồng bộ theo thời gian.} Chúng ta đã thấy cách tính một hàm, cách đo thời gian kể từ một sự kiện và cách cắt luồng thành các cửa sổ bằng nhịp cục bộ. Nhưng một mạng khổng lồ không có xung nhịp chung phối hợp công việc của các vùng xa nhau như thế nào thì mới chỉ hiểu được một phần.

\emph{Năng lượng.} Hai mươi oát cho tất cả. Chúng ta hiểu vì sao mã phải thưa (mệnh đề~\ref{prop:economy}), nhưng chưa ai biết cách chế tạo một cỗ máy làm được điều tương tự với cùng công suất~\cite{MehonicKenyon2022}.

Não là một cỗ máy tính toán không do kỹ sư lắp ráp, nhưng tuân theo những định luật mà kỹ sư nào cũng nhận ra: mạch RC, ngưỡng, phản hồi, bộ lọc, bus và thanh ghi quảng bá. Sơ đồ của nó chúng ta mới đọc được một phần. Đọc nốt phần còn lại sẽ là việc của những người biết đọc sơ đồ.

\section{Phần tiếp theo của cuốn sách}

Chương này đã ráp xong lõi tính toán của cỗ máy. Các chương còn lại vượt ra ngoài lõi đó. Chương~\ref{sec:sensors} và~\ref{sec:recognition} nói về đầu vào: cảm biến biến ánh sáng, âm thanh và phân tử thành xung ra sao, và cỗ máy nhận ra đồ vật trong đó như thế nào. Chương~\ref{sec:muscles}--\ref{sec:chemoutput} nói về đầu ra và những gì phục vụ chúng: cơ, cơn đau như tín hiệu báo động, học vận động qua các dây sai số riêng, và đầu ra hóa học chậm vào cơ thể. Chương~\ref{sec:debugging} nói về cách gỡ lỗi cỗ máy từ bên ngoài, kể cả bằng giao diện não--máy tính. Chương~\ref{sec:eetypes} và~\ref{sec:dendrites} phân loại nơ-ron thêm một lần nữa, lần này theo chức năng điện và theo số đầu vào. Chương~\ref{sec:sleep} nói về những gì cỗ máy làm khi ngắt kết nối với thế giới, còn chương~\ref{sec:livingmachine} và~\ref{sec:dishbrain} nói về các cỗ máy lắp từ nơ-ron sống trên chip.

\section{Bài tập}

Bài tập của chương này kết nối kết quả của nhiều chương: mỗi bài dựa vào ít nhất hai chương.

\begin{exercise}[\lvlA]
\label{ex:10:1}
\emph{Bus và thanh ghi.} Bốn kênh quảng bá là các bó, mỗi bó khoảng năm đường (mệnh đề~\ref{prop:bundle}); mỗi đường thay đổi mỗi giây một lần và phân biệt được bốn mức. Khâu điều khiển cần bao nhiêu năm để truyền lượng thông tin mà bus \nt{GLUT} ($8.6\cdot10^{10}$ nơ-ron với tần số~\eqref{eq:energy}, độ chính xác $1\ms$ theo~\eqref{eq:spikeentropy}) truyền trong một giây?
\end{exercise}

\begin{exercise}[\lvlB]
\label{ex:10:2}
\emph{Hai núm vặn trên một vòng.} Trong vòng~\eqref{eq:ei}, các núm vặn của~\eqref{eq:machine} tác động lên kích thích: $\tau_E\dot E=-E+g[(1-a)w_{EE}E-w_{EI}I+h]$, còn \nt{GABA} không chịu chúng điều khiển. Với số liệu của hình~\ref{fig:ei}, một chùm xung \nt{NORA} nâng $g$ lên $6$. Với $a$ nhỏ nhất bằng bao nhiêu thì vòng ổn định? Có thể vặn các núm \nt{NORA} và \nt{ACET} độc lập không?
\end{exercise}

\begin{exercise}[\lvlB]
\label{ex:10:3}
\emph{Giá của quảng bá từ đâu ra.} Đầu ra của $N$ nơ-ron là $y_k=f(u_k)+\xi_k$ với nhiễu Gauss độc lập phương sai $\sigma^2$; $R\approx R_0+\sum_kR'_k\xi_k$ với các $R'_k$ như nhau, \nt{DOPA} phát $\delta=R-R_0$. Tìm tỷ số tín hiệu/nhiễu của hiệu chỉnh $\delta\,\xi_jx_j$ trong một lần thử. Số lần thử tăng theo $N$ ra sao?
\end{exercise}

\begin{exercise}[\lvlB]
\label{ex:10:4}
\emph{Gắn âm thanh với tia chớp.} Âm thanh đến bus sau $5\ms$, tia chớp sau $30\ms$ cộng độ trễ của xung đầu tiên~\eqref{eq:latency} ($\tau=10\ms$, $U$ từ $1.2\theta$ đến $4\theta$); nền ở mỗi đầu vào là $5$ xung mỗi giây. Hãy chọn độ trễ cho đường âm thanh và cửa sổ trùng khớp nhỏ nhất cho mọi độ tương phản. Nền gây ra bao nhiêu lần gắn nhầm mỗi giây?
\end{exercise}

\begin{exercise}[\lvlB]
\label{ex:10:5}
\emph{Mô phỏng: ai nhận ra thay đổi.} Bộ phê bình thấy $y=s+\text{nhiễu}$ ($R=1$); với xác suất $1/200$, $s$ nhảy tới giá trị mới có phương sai $J^2=4$. Mô phỏng bộ lọc Kalman với $q=0$, nâng $P$ lên $J^2$ khi $E\leftarrow0.9E+\delta^2/(P+R)$ vượt $20$. Sai số của nó nhỏ hơn bao nhiêu so với khi dùng $\alpha$ hằng tốt nhất?
\end{exercise}

\begin{exercise}[\lvlA]
\label{ex:10:6}
\emph{Bao nhiêu lần thử để đổi thói quen.} Mạng đã học $Q_A=0.8$, $Q_B=0.2$ và chọn theo~\eqref{eq:softmax}, $\sigma=1$, $g=8$. Giờ $A$ cho nước quả với xác suất $0.2$; $Q_A\leftarrow Q_A+\alpha(r-Q_A)$, $Q_B$ không đổi. Sau bao nhiêu lần chọn $A$ thì xác suất chọn $A$ giảm xuống $0.6$ khi $\alpha=0.1$ và $0.5$? Còn nếu \nt{NORA} hạ $g$ xuống $2$?
\end{exercise}

\begin{exercise}[\lvlC]
\label{ex:10:7}
\emph{Bó thay cho dây.} Đầu ra $y_k=w_k+\xi_k$, phần thưởng $R=-\sum_ky_k^2$; một đường của bó phát tới nhóm $G$ nơ-ron sai số của các đầu ra trong nhóm, $w_k\leftarrow w_k+\eta\,\delta_l\xi_k$. Chứng minh rằng với $\eta$ tốt nhất, để đạt $\sum w_k^2=\varepsilon\sum w_k^2(0)$ cần $\approx(G+2)\ln(1/\varepsilon)$ lần thử. Với $\varepsilon=0.01$, mạch $10^6$ nơ-ron trên năm đường, mỗi lần thử cách nhau ba giây, học trong bao lâu?
\end{exercise}

\chapter{Đầu vào của cỗ máy: mắt, tai và các cảm biến khác được nối vào ra sao}
\chaptermark{Đầu vào của cỗ máy}
\label{sec:sensors}
\begin{chaptergoals}
\item cách mọi giác quan làm việc như bộ chuyển đổi: chốt thụ thể, điều chỉnh khuếch đại, bộ mã hóa xung;
\item vì sao nhiễu hạt photon giới hạn thị giác trong bóng tối, và vì sao lúc chạng vạng ta nhìn chậm hơn;
\item cách tự động điều chỉnh khuếch đại nén dải đầu vào khổng lồ vào xung, nên cảm biến báo thay đổi;
\item cách mắt nén hình ảnh một trăm lần và tai làm việc như một dãy bộ lọc.
\end{chaptergoals}

\section{Cảm biến như một bộ chuyển đổi}

Trên hình~\ref{fig:machine}, dữ liệu đi vào cỗ máy từ bên trái, từ khối ``cảm biến''. Giờ hãy mở khối đó ra. Với kỹ sư, mỗi giác quan là một \emph{bộ chuyển đổi} có tầng đầu vào tương tự và bộ chuyển đổi tương tự--số ở cuối, chỉ có điều ``chữ số'' ở đầu ra không phải là con số mà là xung.

Mọi giác quan đều có cùng một sơ đồ (hình~\ref{fig:transducer}). Một đại lượng vật lý (ánh sáng, áp lực, dao động, phân tử) tác động lên protein thụ thể trong màng tế bào cảm giác. Protein này hoạt động như một cái chốt: tự nó hoặc qua một chuỗi phản ứng ngắn, nó mở một kênh ion. Xuất hiện \emph{dòng thụ thể}, kéo theo một điện áp tương tự trên màng. Tiếp đó điện áp đi qua khâu điều chỉnh hệ số khuếch đại rồi vào bộ mã hóa xung, chính là nơ-ron tích phân quen thuộc~\eqref{eq:glutneuron}, biến mức thành tần số và thời điểm của xung. Đầu ra hầu như luôn giống nhau: các nơ-ron cảm giác của mắt, tai, khứu giác và da đều giải phóng glutamate, nên đầu vào được nối thẳng vào bus \nt{GLUT}.

\begin{figure}[t]
\centering
\begin{tikzpicture}[>=Stealth,
  blk/.style={draw,very thick,rounded corners=2pt,align=center,font=\footnotesize,minimum height=1.0cm,text width=1.9cm,fill=blue!5}]
\node[align=center,font=\small] (P) at (0.1,0) {ánh sáng, âm,\\áp lực,\\phân tử};
\node[blk] (R) at (3.0,0) {thụ thể\\làm chốt};
\node[blk] (G) at (6.5,0) {điều chỉnh\\khuếch đại};
\node[blk] (E) at (10.0,0) {bộ mã hóa\\xung};
\node[align=center,font=\small] (B) at (13.4,0) {bus\\\nt{GLUT}};
\draw[->,very thick] (P) -- (R);
\foreach \ic/\xx in {sun/-1.1,speaker/-0.3,press/0.5,molecule/1.3}{\pic[scale=0.85] at (\xx,1.05) {\ic};}
\draw[->,very thick] (R) -- node[above,font=\scriptsize] {dòng} (G);
\draw[->,very thick] (G) -- node[above,font=\scriptsize] {mức} (E);
\draw[->,very thick,red!70!black] (E) -- node[above,font=\scriptsize,black] {xung} (B);
\draw[decorate,decoration={brace,amplitude=5pt,mirror},gray!60!black] (1.8,-0.8) -- (7.7,-0.8) node[midway,below=5pt,font=\scriptsize,black] {phần tương tự};
\draw[decorate,decoration={brace,amplitude=5pt,mirror},gray!60!black] (8.8,-0.8) -- (11.2,-0.8) node[midway,below=5pt,font=\scriptsize,black] {xung};
\end{tikzpicture}
\caption{Mọi giác quan đều là một chuỗi biến đổi. Protein thụ thể mở kênh và tạo ra dòng điện; khâu điều chỉnh khuếch đại làm dòng đó thích nghi với hoàn cảnh; bộ mã hóa~\eqref{eq:glutneuron} biến mức thành xung đi lên bus \nt{GLUT}. Ở mắt và tai, các tầng đầu tiên hoàn toàn không phát xung: tín hiệu vẫn là tương tự cho tới khi phải truyền đi xa.}
\label{fig:transducer}
\end{figure}

Có một chi tiết mà người làm điện tử rất quen. Ở mắt và tai, các tế bào đầu tiên hoàn toàn không phát xung: chúng truyền cho tế bào bên cạnh một điện áp biến thiên liên tục, như một tầng tương tự trên bảng mạch. Xung chỉ xuất hiện ở nơi tín hiệu phải đi xa. Tín hiệu tương tự đi vài milimét đã suy giảm và nhiễm nhiễu, còn xung, như ta đã thấy ở chương~\ref{sec:neurontypes}, tới cuối dây vẫn giữ nguyên độ cao. Số hóa diễn ra ở nơi bắt đầu đường truyền dài.

\section{Ở giới hạn của vật lý}

Cảm biến của não làm việc sát giới hạn vật lý của độ nhạy. Trong bóng tối, cảm biến ánh sáng đáp ứng với \emph{một photon}: hấp thụ một hạt ánh sáng tạo ra dòng khoảng một picoampe, đủ nổi lên trên nhiễu của chính nó~\cite{Baylor1979}. Cảm biến âm thanh nhận ra độ dịch chuyển của bó lông cảm giác nhỏ hơn một nanomét~\cite{Hudspeth1989}.

Khi ánh sáng yếu, độ chính xác bị giới hạn không phải bởi cảm biến mà bởi chính ánh sáng: photon đến ngẫu nhiên, theo dòng Poisson. Nếu trong thời gian tích lũy $T$, cảm biến nhận trung bình $N=\Phi T$ photon, thì số photon dao động cỡ $\sqrt N$. Để phần tăng $\Delta N$ nhận ra được, nó phải lớn hơn dao động này vài lần, và tỷ lệ thay đổi độ sáng phân biệt được là
\begin{empheq}[box=\keybox]{equation}
\frac{\Delta I}{I} \;\approx\; \frac{k}{\sqrt{\Phi\,T}} .
\label{eq:photonnoise}
\end{empheq}
Đây chính là công thức cho số xung~\eqref{eq:rateerror}: nhiễu hạt của photon và nhiễu hạt của xung tuân theo cùng một định luật. Trong bóng tối, mắt chỉ có một cách tăng độ chính xác là tích lũy photon lâu hơn, và thời gian tích lũy quả thật tăng tới vài phần mười giây. Vì thế lúc chạng vạng ta nhìn chậm hơn.

\section{Dải động: điều chỉnh khuếch đại}

Bài toán kỹ thuật khó nhất của cảm biến là dải đo. Độ rọi từ đêm đầy sao đến tuyết dưới nắng thay đổi khoảng mười bậc độ lớn, cường độ âm từ ngưỡng nghe đến ngưỡng đau thay đổi mười hai bậc. Còn tần số xung chỉ thay đổi từ không đến vài trăm mỗi giây, chưa tới ba bậc. Không thể xếp tuyến tính một đầu vào như thế vào một đầu ra như thế.

Giải pháp giống như ở mọi máy thu: \emph{tự động điều chỉnh khuếch đại}. Đáp ứng của các tế bào cảm giác ở mắt được mô tả tốt bằng công thức
\begin{empheq}[box=\keybox]{equation}
r \;=\; r_{\max}\,\frac{I}{I+\sigma} ,
\label{eq:nakarushton}
\end{empheq}
trong đó $\sigma$ là độ sáng tại đó đáp ứng đạt một nửa~\cite{NakaRushton1966}. Tự thân~\eqref{eq:nakarushton} chỉ bao được hai đến ba bậc. Nhưng $\sigma$ không cố định: khi làm việc lâu trên nền sáng, nó tăng theo độ sáng trung bình. Đường cong đáp ứng dịch chuyển dọc trục logarit của độ sáng và luôn đặt đoạn dốc của mình vào chỗ đầu vào đang ở (hình~\ref{fig:adaptation}). Đây vẫn là phép chia cho hoạt động chung như ức chế sun ở chương~\ref{sec:gaba}~\cite{Carandini2012}, chỉ khác là mẫu số ở đây là độ sáng trung bình trong vài giây vừa qua.

\begin{figure}[t]
\centering
\begin{tikzpicture}[>=Stealth,x=1.25cm,y=2.8cm]
\draw[->,gray!70!black] (-2,0) -- (6.4,0) node[below left,font=\small,black] {$\log_{10} I$};
\draw[->,gray!70!black] (-2,0) -- (-2,1.15) node[left,font=\small,black] {$r/r_{\max}$};
\foreach \u in {-2,0,2,4,6}{\draw[gray!60!black] (\u,-0.02) -- (\u,0.02); \node[below,font=\scriptsize] at (\u,0) {$\u$};}
\draw[densely dashed,gray!60!black] (-2,0.5) -- (6.2,0.5);
\foreach \s/\c/\lab in {0/blue!60!black/{nền tối},2/green!45!black/{vừa},4/red!70!black/{sáng}}{
  \pgfmathsetmacro{\lo}{max(-2,\s-3)}
  \draw[very thick,\c] (-2,0) -- (\lo,0) plot[domain=\lo:6.2,samples=80] (\x,{1/(1+10^(\s-\x))});
  \node[font=\scriptsize,\c,anchor=west] at (\s+0.15,0.42) {\lab};}
\foreach \s/\ic in {0/moon,2/cloud,4/sun}{\pic at (\s+0.6,0.64) {\ic};}
\end{tikzpicture}
\caption{Điều chỉnh khuếch đại của cảm biến ánh sáng theo~\eqref{eq:nakarushton}. Mỗi đường cong bao hai đến ba bậc độ sáng; khi chuyển sang nền sáng hơn, $\sigma$ tăng và đường cong dịch theo trục logarit. Các đường cong dịch chuyển cùng nhau phủ kín toàn bộ dải.}
\label{fig:adaptation}
\end{figure}

Điều chỉnh này có cái giá quen thuộc từ chương~\ref{sec:code}: cảm biến không báo độ sáng, mà báo độ sáng \emph{so với nền}. Đầu vào không đổi dần dần thôi gây ra đáp ứng, còn mỗi thay đổi thì gây ra đáp ứng. Ngay từ đầu, các đầu vào của cỗ máy đã nói về thay đổi chứ không về mức, và đó cũng là cách tiết kiệm xung như trong mệnh đề~\ref{prop:economy}~\cite{Barlow1961}. Cũng từ đây mà có các cảm biến bật và cảm biến tắt của chương~\ref{sec:glutcomputer}~\cite{Schiller1992}: loại thứ nhất báo rằng trời sáng hơn, loại thứ hai báo rằng tối hơn.

Kỹ sư đã lặp lại mẹo này trong \emph{camera sự kiện}. Camera như vậy không chụp khung hình theo đồng hồ: mỗi điểm ảnh của nó tự mình, không cần xung nhịp chung, phát ra sự kiện ``sáng hơn'' hoặc ``tối hơn'' khi logarit độ sáng thay đổi một lượng định trước~\cite{Lichtsteiner2008}. Đó đúng là mã bật--tắt hai dây của chương~\ref{sec:glutcomputer}, hiện thực trên silicon, và nó cho camera dải động và tốc độ mà cảm biến ảnh thông thường không đạt được~\cite{Gallego2022}.

\section{Mắt: nén xuống một sợi cáp}

Mắt có khoảng $10^8$ cảm biến ánh sáng~\cite{Curcio1990}, còn số nơ-ron đầu ra đưa hình ảnh lên não chỉ khoảng một triệu. Trước khi rời mắt, tín hiệu bị nén khoảng một trăm lần. Các kỹ thuật nén chính đều quen thuộc với chúng ta. Mỗi nơ-ron đầu ra đáp ứng với hiệu giữa độ sáng ở một đốm nhỏ và độ sáng xung quanh, tức phép trừ láng giềng qua ức chế như ở chương~\ref{sec:gaba}. Những vùng đồng đều của ảnh gần như không tốn gì, còn các cạnh và thay đổi thì được truyền đi. Các nơ-ron đầu ra khác nhau được chuyên môn hóa: loại báo sáng lên, loại báo tối đi, loại báo chuyển động theo một hướng nhất định (hình~\ref{fig:direction}); ở chuột nhắt, người ta đếm được hơn ba mươi kênh đầu ra như vậy~\cite{Baden2016}.

Sợi cáp thu được có băng thông bao nhiêu? Ở chuột lang, từ bản ghi các nơ-ron đầu ra, người ta đo được khoảng $875$ nghìn bit mỗi giây trên $\sim10^5$ nơ-ron như vậy; quy đổi cho một triệu nơ-ron đầu ra ở người được cỡ $10^7$ bit mỗi giây~\cite{Koch2006}, bằng một sợi cáp mạng mười megabit đời cũ. Với tần số trung bình vài xung mỗi giây trên một nơ-ron, đó là vài bit mỗi xung, đúng như kết quả ở chương~\ref{sec:code}.

\section{Tai: dãy bộ lọc}

Tai giải một bài toán khác: phải phân tích âm thanh theo tần số. Kỹ sư sẽ đặt một dãy bộ lọc thông dải, và tai làm đúng điều đó bằng cơ học. Dao động âm chạy dọc một vách ngăn đàn hồi có tính chất thay đổi dần theo chiều dài, và mỗi đoạn cộng hưởng ở tần số riêng. Các cảm biến nằm trên một đoạn chỉ đáp ứng với dải tần của đoạn đó (hình~\ref{fig:filterbank}). Tần số được sắp theo vị trí gần như theo logarit: mỗi quãng tám nhận một phần cảm biến tương tự nhau. Số thứ tự của cảm biến bị kích hoạt chính là tần số, tức mã vị trí của chương~\ref{sec:code}.

\begin{figure}[t]
\centering
\begin{tikzpicture}[>=Stealth,x=1.35cm,y=2.2cm]
\draw[->,gray!70!black] (-0.3,0) -- (7.6,0) node[above left,font=\small,black] {tần số, Hz};
\draw[->,gray!70!black] (-0.3,0) -- (-0.3,1.15) node[left,font=\small,black] {đáp ứng};
\foreach \k/\f in {0/125,1/250,2/500,3/1000,4/2000,5/4000,6/8000}{
  \draw[gray!60!black] (\k,-0.02) -- (\k,0.02); \node[below,font=\scriptsize] at (\k,0) {\f};
  \draw[thick,blue!60!black] plot[domain={max(\k-1.2,-0.3)}:{\k+0.7},samples=60] (\x,{1/(1+((\x-\k)/(\x<\k?0.45:0.2))^4)});}
% a piano keyboard: seven white keys per octave, like the octaves on the axis
\foreach \i in {0,...,48}{\draw[gray!50!black,fill=white,line width=0.3pt] ({\i/7},-0.44) rectangle ({(\i+1)/7},-0.26);}
\foreach \o in {0,...,6}{\foreach \j in {0,1,3,4,5}{\fill[black] ({\o+(\j+1)/7-0.035},-0.35) rectangle ({\o+(\j+1)/7+0.035},-0.26);}}
\end{tikzpicture}
\caption{Tai như một dãy bộ lọc thông dải, sơ đồ. Mỗi đường cong là đáp ứng của các cảm biến trên một đoạn với âm có tần số khác nhau; các tần số trung tâm cách nhau một quãng tám, như phím đàn trên thang logarit. Bộ lọc thật có sườn tần số cao dốc và sườn tần số thấp thoải. Phía dưới là bàn phím: trên đó, cũng như trong tai, mỗi quãng tám chiếm chỗ như nhau.}
\label{fig:filterbank}
\end{figure}

Các bộ cộng hưởng của tai không thụ động. Một phần cảm biến tự đẩy vách ngăn dao động theo nhịp âm thanh và khuếch đại các dao động yếu lên hàng nghìn lần về biên độ ($66$--$76$\,dB), còn khi âm to lên thì hệ số khuếch đại giảm~\cite{Ruggero1997,Robles2001}. Với kỹ sư, đó là bộ khuếch đại có phản hồi dương, đặt sát ranh giới tự kích, cũng là kiểu sống bên bờ vực như mạng ở chương~\ref{sec:gaba}: gần ranh giới, hệ đặc biệt nhạy với tín hiệu nhỏ. Việc nén âm lượng ở đây do chính bộ khuếch đại đó đảm nhận: âm nhỏ được khuếch đại mạnh, âm to được khuếch đại yếu.

Tai còn có mã thứ hai. Ở tần số thấp, xung của nơ-ron đi cùng pha với âm: nơ-ron phát xung ở một đoạn nhất định của mỗi sóng, dù không phải ở mọi sóng. Ở chuột lang, khóa pha bắt đầu yếu đi từ khoảng $600$\,Hz và vẫn còn thấy được tới khoảng $3.5$\,kHz~\cite{PalmerRussell1986}. Nhờ đó thời điểm xung giữ được thời gian chính xác của âm, và từ đó là chênh lệch thời điểm âm đến hai tai, dựa vào đó mạng ở chương~\ref{sec:glutcomputer} tìm ra hướng~\cite{Grothe2010}. Ở tần số cao, xung không theo kịp sóng, và chỉ còn mã vị trí. Một cái tai dùng cả hai mã của chương~\ref{sec:code}: thời gian ở chỗ nơ-ron theo kịp, và vị trí ở chỗ không theo kịp.

\section{Các cảm biến khác}

\begin{table}[t]
\centering
\footnotesize
\setlength{\tabcolsep}{3pt}
\renewcommand{\arraystretch}{1.15}
\begin{tabular}{>{\raggedright}p{1.9cm}>{\raggedright}p{2.6cm}>{\raggedright}p{4.0cm}>{\raggedright\arraybackslash}p{4.6cm}}
\toprule
Giác quan & Đo gì & Đầu vào được tổ chức ra sao & Kỹ thuật trong sách \\
\midrule
thị giác & ánh sáng & $\sim10^8$ cảm biến, nén $\sim100$ lần, $\sim10^7$ bit/s & điều chỉnh khuếch đại, trừ láng giềng, hai dây, hướng chuyển động \\
thính giác & áp suất không khí & dãy bộ lọc cơ học, bộ khuếch đại chủ động & mã vị trí và mã thời gian, độ trễ \\
xúc giác & áp lực, rung & cảm biến thích nghi nhanh và chậm & cặp ``bộ lọc thông cao -- bộ lọc thông thấp'' \\
nhiệt độ & nhiệt & cảm biến nóng và lạnh riêng biệt & mã hai dây \\
khứu giác & phân tử & $\sim400$ loại thụ thể, mỗi cảm biến một loại & mã tổ hợp: mùi là tập các loại bị kích hoạt \\
vị giác & chất trong thức ăn & năm vị: ngọt, đắng, mặn, chua, umami & đường dây gắn nhãn; một số tế bào giải phóng ATP hoặc serotonin chứ không phải glutamate \\
tư thế cơ thể & độ giãn của cơ & cảm biến độ dài và tốc độ giãn & phản hồi cho bộ điều khiển vận động \\
\bottomrule
\end{tabular}
\caption{Các cảm biến được nối vào ra sao. Mọi cơ chế ở cột thứ ba đều đã được đo; cột bên phải nối chúng với các kỹ thuật ở những chương trước.}
\label{tab:sensors}
\end{table}

Các giác quan còn lại được tập hợp trong bảng~\ref{tab:sensors}, và hầu như dòng nào cũng có một kỹ thuật quen thuộc từ các chương trước. Xúc giác dùng một cặp bộ lọc: một số cảm biến áp lực đáp ứng khi chạm vào rồi im ngay, số khác giữ đáp ứng chừng nào còn áp lực. Nhiệt độ được truyền qua hai dây, riêng cho nóng và cho lạnh. Đầu vào khác thường nhất là khứu giác. Người có khoảng bốn trăm loại thụ thể khứu giác, và mỗi nơ-ron khứu giác chỉ mang một loại~\cite{Mombaerts2004}. Một mùi kích hoạt không phải một loại mà cả một tập hợp, và thông điệp chính là \emph{những loại nào} đã bị kích hoạt. Với kỹ sư, đó là một vectơ nhị phân dài khoảng bốn trăm, với số giá trị khả dĩ lớn khủng khiếp.

\begin{keyblock}
\begin{proposition}[Đầu vào của cỗ máy]
\label{prop:sensors}
Mỗi giác quan là một bộ chuyển đổi làm việc ở giới hạn vật lý của độ nhạy, nén dải động khổng lồ bằng tự động điều chỉnh khuếch đại, và truyền trên bus \nt{GLUT} trước hết là \emph{các thay đổi}. Để làm vậy, cảm biến dùng chính những kỹ thuật của cỗ máy: mã hai dây, mã vị trí, mã thời gian và phép chia cho nền. Cấu tạo của cảm biến đã được đo; việc mã của chúng được tinh chỉnh để mang nhiều thông tin nhất trên mỗi xung là một giả thuyết có cơ sở vững.
\end{proposition}
\end{keyblock}

Đầu vào, như mọi thứ khác, làm việc bất đồng bộ. Mỗi cảm biến phát xung khi nó có điều cần báo, và các giác quan khác nhau đến với độ trễ khác nhau: âm thanh sau vài mili giây, ánh sáng sau vài chục mili giây. Không có xung nhịp chung thì cỗ máy gộp chúng thành một bức tranh thế giới như thế nào: đó là một trong các bài toán đồng bộ theo thời gian ở chương~\ref{sec:synthesis}.

\section{Bài tập}

\begin{exercise}[\lvlA]
\label{ex:11:1}
\emph{Tích lũy photon bao lâu.} Lúc chạng vạng, cảm biến nhận $100$ photon mỗi giây; phần tăng nhận ra được nếu lớn gấp $3$ lần dao động~\eqref{eq:photonnoise}. (a)~Một cảm biến cần bao lâu để nhận ra độ sáng thay đổi $10\,\%$? (b)~Phải cộng bao nhiêu cảm biến để kịp trong $0.2\,\text{s}$, và độ phân giải không gian giảm bao nhiêu lần?
\end{exercise}

\begin{exercise}[\lvlA]
\label{ex:11:2}
\emph{Mỗi xung mang bao nhiêu bit.} (a)~Mắt truyền $\sim10^7$ bit/s qua $10^6$ nơ-ron đầu ra với tần số $3$--$5$ xung mỗi giây. Nó dùng bao nhiêu phần dung lượng~\eqref{eq:spikeentropy} với $\Delta t=1\ms$? (b)~Một mùi bật $20$ trong $\approx400$ loại thụ thể. Tập đó mang bao nhiêu bit, và trung bình hai mùi ngẫu nhiên có bao nhiêu loại chung?
\end{exercise}

\begin{exercise}[\lvlB]
\label{ex:11:3}
\emph{Một đường cong~\eqref{eq:nakarushton} làm được gì.} (a)~Trong dải độ sáng nào đáp ứng tăng từ $10\,\%$ đến $90\,\%$ của $r_{\max}$? Đó là bao nhiêu bậc? (b)~Cần bao nhiêu vị trí của đường cong dịch chuyển để phủ mười bậc độ sáng? mười hai bậc âm lượng?
\end{exercise}

\begin{exercise}[\lvlB]
\label{ex:11:4}
\emph{Giới hạn của mã thời gian ở tai.} Chênh lệch thời điểm âm đến hai tai tối đa $0.22\,\text{m}/343\,\text{m/s}$. (a)~Xung đi cùng pha với âm: tới tần số nào thì hướng xác định từ chúng là duy nhất? (b)~Xung rung $0.5\ms$, còn cần phân biệt $20$~\textmu s. Phải lấy trung bình bao nhiêu nơ-ron, mỗi nơ-ron $100$ xung mỗi giây, trong $50\ms$?
\end{exercise}

\begin{exercise}[\lvlB]
\label{ex:11:5}
\emph{Camera sự kiện trên cáp của mắt.} Camera $10^6$ điểm ảnh, đầu ra $10^7$ bit/s, gửi một sự kiện (địa chỉ và dấu) khi $\ln I$ ở điểm ảnh thay đổi $\ln1.2$. Một cạnh dài $500$ điểm ảnh, chênh sáng $4$ lần, chạy nhanh nhất được bao nhiêu? Camera thường, $8$ bit mỗi điểm ảnh, qua cùng đầu ra đó cho bao nhiêu khung hình mỗi giây?
\end{exercise}

\begin{exercise}[\lvlB]
\label{ex:11:6}
\emph{Mô phỏng: điều chỉnh khuếch đại.} Cảm biến~\eqref{eq:nakarushton} chỉnh $\sigma$ với $\tau=1\,\text{s}$ theo logarit, $\tau\,\dd\ln\sigma/\dd t=\ln I-\ln\sigma$, hoặc tuyến tính, $\tau\,\dd\sigma/\dd t=I-\sigma$; nền nhảy $1\to100\to1$. Với mỗi luật, sau bước nhảy lên và xuống, bao nhiêu giây thì đáp ứng với phép thử $+10\,\%$ hồi phục được một nửa mức đã thích nghi?
\end{exercise}

\begin{exercise}[\lvlC]
\label{ex:11:7}
\emph{Đường cong khớp với thế giới nào.} Với nhiễu đầu ra nhỏ và không đổi, thông tin về độ sáng lớn nhất khi $r(I)=r_{\max}\Phi_I(I)$, với $\Phi_I$ là hàm phân phối của độ sáng. (a)~Đường cong~\eqref{eq:nakarushton} tối ưu cho phân phối nào, và độ tản của $\log_{10}I$ bằng bao nhiêu? (b)~Với nhiễu đầu ra Poisson thì đường cong nào tối ưu?
\end{exercise}

\chapter{Nhận dạng hình ảnh và video}
\chaptermark{Nhận dạng}
\label{sec:recognition}
\begin{chaptergoals}
\item vì sao nhận dạng phải bỏ qua phép dịch, tỷ lệ và độ sáng, và vì sao so với mẫu thất bại;
\item cách dây chuyền xen kẽ VÀ theo bộ phận và HOẶC theo vị trí nhận dạng trong $150\ms$;
\item cách quy tắc Hebb có vết học tính bất biến từ video không có nhãn;
\item não khác mạng tích chập ở đâu, và ảo giác cho thấy gì về nó.
\end{chaptergoals}

\section{Bài toán: cùng một vật với những điểm ảnh khác nhau}

Chương~\ref{sec:sensors} đã đưa hình ảnh tới đầu vào của cỗ máy: khoảng một triệu nơ-ron đầu ra của mắt, một luồng đã nén, trong đó chủ yếu truyền các cạnh và các thay đổi. Nhưng giữa cảm biến và hành động còn một bước mà ta chưa nhắc tới. Con mèo trong ảnh không phải là một tập điểm ảnh. Dịch nó đi nửa khung hình, xoay nó, đẩy nó ra xa, tắt một nửa ánh sáng: gần như mọi điểm ảnh đều khác đi, vậy mà câu trả lời ``con mèo'' phải giữ nguyên. Kỹ sư gọi đó là \emph{tính bất biến}: câu trả lời của bộ phân loại không được đổi khi dịch chuyển, co giãn, xoay và thay đổi ánh sáng.

Cái khó thấy rõ qua một thí nghiệm đơn giản. Lấy một ``võng mạc'' một chiều gồm $24$ điểm và hai hình, mỗi hình năm điểm, có thể nằm ở bất kỳ đâu. Bộ phân loại so sánh đầu vào với một mẫu ghi nhớ tại một vị trí đoán đúng $49$--$52\%$ số trường hợp, ngang với tung đồng xu. Phép dịch chuyển phá hỏng hoàn toàn việc so sánh theo điểm ảnh. Cần một sơ đồ khác.

\section{Dây chuyền nhiều tầng}

Não làm việc này nhanh đến đâu, ta đã biết từ chương~\ref{sec:code}: con vật trong một bức ảnh chiếu chớp nhoáng được nhận ra sau khoảng $150\ms$, và tín hiệu đi qua chừng mười tầng, mỗi tầng $10$--$20\ms$~\cite{Thorpe1996}. Ở mỗi tầng, nơ-ron chỉ kịp phát không hoặc một xung. Vậy nhận dạng nhanh là một lượt đi qua dây chuyền, không có suy nghĩ lâu và lặp lại. Câu hỏi là mỗi tầng làm gì.

Câu trả lời mà các phép đo ở những tầng thị giác đầu tiên gợi ý~\cite{HubelWiesel1962} gồm hai phép toán xen kẽ (hình~\ref{fig:conveyor}).

\emph{Tính chọn lọc.} Nơ-ron của tầng $S$ nhìn một vùng nhỏ của đầu vào và chỉ phát xung nếu ở đó có một tổ hợp bộ phận nhất định: cạnh này, và cạnh kia bên cạnh. Đó là phép VÀ theo các bộ phận, tức nơ-ron ngưỡng của chương~\ref{sec:glutcomputer} với ngưỡng cao hơn mức mà bất kỳ bộ phận đơn lẻ nào tạo ra.

\emph{Tính bất biến.} Nơ-ron của tầng $C$ gom đầu ra của nhiều nơ-ron $S$ cùng tìm một tổ hợp ở những vị trí khác nhau, và phát xung nếu tổ hợp đó có ở bất cứ đâu. Đó là phép HOẶC theo vị trí, hay chính xác hơn là chọn giá trị lớn nhất.

Với mô hình một chiều, hai phép toán được viết như sau:
\begin{empheq}[box=\keybox]{equation}
s_{f,p} \;=\; \Bigl[\sum_i t_{f,i}\,x_{p+i} - \theta\Bigr]_+ ,
\qquad
c_f \;=\; \max_p\, s_{f,p} ,
\label{eq:scstage}
\end{empheq}
trong đó $x$ là đầu vào, $t_{f,i}$ là mẫu của đặc trưng $f$, $p$ là vị trí, $\theta$ là ngưỡng. Biểu thức thứ nhất làm câu trả lời có tính chọn lọc, biểu thức thứ hai làm nó không phụ thuộc vị trí. Mạng lấy được giá trị lớn nhất trong nhiều đầu vào bằng các công cụ của chương~\ref{sec:gaba}: ức chế lẫn nhau giữ lại đầu vào mạnh nhất (mệnh đề~\ref{prop:wta}), còn phép chia cho hoạt động chung san bằng đáp ứng khi độ sáng khác nhau~\cite{Carandini2012}.

\begin{figure}[t]
\centering
\begin{tikzpicture}[>=Stealth,
  px/.style={draw,minimum size=0.36cm,inner sep=0pt},
  on/.style={px,fill=blue!55!black},
  snode/.style={circle,draw,very thick,minimum size=0.62cm,inner sep=0pt,font=\scriptsize,fill=blue!6},
  cnode/.style={circle,draw,very thick,minimum size=0.8cm,inner sep=0pt,font=\scriptsize,fill=red!10}]
% two inputs: the same shape at two positions
\foreach \row/\sh/\lab in {0/1/{khung 1},1/7/{khung 2}}{
  \begin{scope}[yshift=-\row*1.0cm]
  \node[left,font=\scriptsize] at (-0.25,0) {\lab};
  \foreach \i in {0,...,13}{\node[px] at (\i*0.4,0) {};}
  \foreach \d in {0,1,3,4}{\node[on] at ({(\sh+\d)*0.4},0) {};}
  \end{scope}}
% S stage: detectors of the shape at several positions
\foreach \k in {0,...,5}{
  \node[snode] (S\k) at ({0.8+0.8*\k},1.7) {$S$};}
\node[font=\scriptsize,left] at (-0.25,1.7) {tầng $S$: VÀ theo bộ phận};
\draw[->,thick,blue!60!black] (1.2,0.25) -- (S1.south);
\draw[->,thick,orange!85!black] (3.6,0.25) -- (S4.south);
\node[font=\scriptsize,blue!60!black,anchor=east] at (1.25,0.85) {khung 1};
\node[font=\scriptsize,orange!85!black,anchor=west] at (3.9,0.85) {khung 2};
% C stage
\node[cnode] (C) at (2.8,3.25) {$C$};
\node[font=\scriptsize,left] at (-0.25,3.25) {tầng $C$: HOẶC theo vị trí};
\foreach \k in {0,...,5}{\draw[->,gray!60!black] (S\k.north) -- (C);}
\draw[->,very thick,red!70!black] (C.north) -- ++(0,0.6) node[above,font=\scriptsize,black] {``có hình'' ở cả hai khung};
% next stages
\draw[decorate,decoration={brace,amplitude=5pt},gray!60!black] (6.3,1.4) -- (6.3,-1.3);
\node[right,font=\scriptsize,align=left] at (6.5,0.05) {cùng một\\vật, điểm ảnh\\khác nhau};
\draw[->,thick,densely dashed,gray!60!black] (3.6,3.25) -- (5.9,3.25) node[right,font=\scriptsize,black,align=left] {các cặp $S$, $C$ tiếp theo:\\lớn hơn, phức tạp hơn,\\bất biến hơn};
\end{tikzpicture}
\caption{Một cặp tầng của dây chuyền. Các nơ-ron $S$ tìm cùng một tổ hợp bộ phận ở những vị trí khác nhau; nơ-ron $C$ đáp ứng nếu tổ hợp đó có ở bất cứ đâu~\eqref{eq:scstage}. Hình ở khung 1 và ở khung 2 kích thích các nơ-ron $S$ khác nhau, nhưng cùng một nơ-ron $C$. Các cặp tầng tiếp theo lặp lại điều đó với những tổ hợp ngày càng lớn và phức tạp.}
\label{fig:conveyor}
\end{figure}

Cũng trong mô hình một chiều ấy, cặp tầng~\eqref{eq:scstage} nhận ra hình ở bất kỳ vị trí nào mà không sai, còn nếu mỗi điểm của đầu vào bị lật ngẫu nhiên với xác suất $0.1$ thì đúng $86\%$ số trường hợp, với xác suất $0.2$ thì $70\%$. Đồng xu vẫn chỉ cho một nửa. Chồng vài cặp như vậy lên nhau: cặp đầu tìm cạnh, cặp sau tìm tổ hợp các cạnh, cao hơn nữa là bộ phận của đồ vật, trên cùng là toàn bộ đồ vật. Qua mỗi cặp, tổ hợp lớn hơn, còn câu trả lời ngày càng ít phụ thuộc vào việc đồ vật nằm ở đâu và nằm thế nào~\cite{Riesenhuber1999}.

\begin{keyblock}
\begin{proposition}[Dây chuyền chọn lọc và bất biến]
\label{prop:conveyor}
Nhận dạng nhanh là một lượt đi qua chừng mười tầng. Ở mỗi cặp tầng, phép VÀ theo bộ phận~\eqref{eq:scstage} làm câu trả lời có tính chọn lọc, còn phép HOẶC theo vị trí làm nó không phụ thuộc phép dịch chuyển. Xen kẽ hai phép toán này cho ra câu trả lời phân biệt được đồ vật mà gần như không phụ thuộc vào việc chúng được nhìn thấy ở đâu và thế nào. Cấu tạo của các tầng đầu và tốc độ của một lượt đã được đo; việc cả chuỗi quy về đúng hai phép toán này là một sự đơn giản hóa.
\end{proposition}
\end{keyblock}

Nếu cấu tạo này có vẻ quen thuộc thì đúng là như vậy. Chính sự xen kẽ này (tìm đặc trưng ở mọi vị trí, rồi gộp theo các vị trí lân cận) đã được kỹ sư chuyển vào mạng nhân tạo \emph{tích chập}~\cite{Fukushima1980}. Gần đây thì đi theo chiều ngược lại: mạng tích chập được huấn luyện để nhận dạng đồ vật hóa ra là mô hình tốt nhất đã biết cho đáp ứng của nơ-ron ở các tầng thị giác cao~\cite{Yamins2014}. Kết quả ở trên cùng được đọc ra đơn giản: từ tổng tần số của chừng một trăm nơ-ron ở tầng cuối, một khối quyết định tuyến tính nhận ra đồ vật sau một phần mười giây kể từ khi hiện hình~\cite{Hung2005}. Đó là mã tần số của nhóm ở chương~\ref{sec:code}.

\section{Người thầy không có mặt: học từ video}

Mạng tích chập được huấn luyện trên hàng triệu bức ảnh có nhãn. Não thì hầu như không ai cho nhãn: đứa trẻ nhìn đồ vật hàng giờ, còn từ ``cái cốc'' thì thỉnh thoảng mới nghe. Vậy các tầng $C$ biết gom những nơ-ron $S$ nào từ đâu? Muốn gộp ``hình này ở đây'' với ``hình này ở kia'' thì phải biết sẵn đó là cùng một hình.

Gợi ý đến từ chính video. Thế giới thay đổi liên tục: đồ vật trong trường nhìn vẫn là đồ vật đó khi nó dịch chuyển, xoay và đổi ánh sáng, và chỉ thỉnh thoảng ánh mắt mới chuyển sang vật khác. Mọi thứ thay đổi \emph{chậm} nhiều khả năng thuộc về đồ vật, còn mọi thứ thay đổi nhanh thì thuộc về vị trí và dáng vẻ của nó~\cite{WiskottSejnowski2002}. Vậy nơ-ron $C$ chỉ cần quy tắc Hebb có vết của chương~\ref{sec:dofa}: nối những gì đến liên tiếp, chứ không chỉ những gì đến đồng thời~\cite{Foldiak1991}:
\begin{empheq}[box=\keybox]{equation}
\tau_{\text{tr}}\,\frac{\dd\bar y_k}{\dd t} \;=\; -\bar y_k + y_k ,
\qquad
\frac{\dd\mathbf w_k}{\dd t} \;=\; \eta\,\bar y_k\,\bigl(\mathbf x - \mathbf w_k\bigr) .
\label{eq:traceinvariance}
\end{empheq}
Ở đây $y_k$ là hoạt động của nơ-ron $C$ thứ $k$, nơ-ron đã thắng trong cuộc cạnh tranh qua ức chế, $\bar y_k$ là vết của hoạt động đó, cũng là mạch RC như trong quy tắc~\eqref{eq:threefactor}, còn $\mathbf x$ là đầu ra của các nơ-ron $S$. Nơ-ron thắng ở dáng vẻ thứ nhất của đồ vật vẫn còn nhớ điều đó khi dáng vẻ thứ hai tới, và kéo luôn cả dáng vẻ này về phía mình.

\begin{figure}[t]
\centering
\begin{tikzpicture}[>=Stealth,x=1.0cm,y=1.0cm]
\foreach \i/\lab in {0/1,1/2,2/3,3/4}{
  \node[draw,thick,fill=blue!8,minimum width=0.9cm,minimum height=0.55cm,font=\scriptsize] at (\i+0.5,2.2) {$A$ ở \lab};}
\foreach \i/\lab in {0/4,1/3,2/2,3/1}{
  \node[draw,thick,fill=orange!12,minimum width=0.9cm,minimum height=0.55cm,font=\scriptsize] at (\i+6.5,2.2) {$B$ ở \lab};}
\node[font=\scriptsize,gray!40!black] at (5.0,2.2) {nghỉ};
\draw[->,gray!70!black] (0,0) -- (10.4,0) node[below left,font=\scriptsize,black] {thời gian};
% traces
\draw[thick,blue!60!black] (0,0.05) .. controls (0.4,1.2) and (0.8,1.3) .. (1.2,1.35) -- (4,1.4) .. controls (4.6,0.5) and (5.2,0.15) .. (6,0.08) -- (10,0.05);
\draw[thick,orange!85!black] (0,0.05) -- (6,0.05) .. controls (6.4,1.2) and (6.8,1.3) .. (7.2,1.35) -- (10,1.4);
\node[font=\scriptsize,blue!60!black,anchor=west] at (1.4,1.65) {vết $\bar y_1$ giữ suốt lượt đi qua của $A$};
\node[font=\scriptsize,orange!85!black,anchor=west] at (7.2,1.65) {vết $\bar y_2$};
\draw[decorate,decoration={brace,amplitude=4pt},gray!60!black] (4.0,2.6) -- (0,2.6);
\draw[decorate,decoration={brace,amplitude=4pt,mirror},gray!60!black] (0,-0.35) -- (4,-0.35) node[midway,below=4pt,font=\scriptsize,black] {mọi dáng vẻ của $A$ được nối với nơ-ron 1};
\end{tikzpicture}
\caption{Video dạy tính bất biến như thế nào, sơ đồ. Đồ vật $A$ đi qua bốn vị trí, sau đó là quãng nghỉ, rồi đến đồ vật $B$. Vết~\eqref{eq:traceinvariance} của nơ-ron thắng ở dáng vẻ đầu tiên của $A$ được giữ suốt lượt đi qua, và mọi dáng vẻ của $A$ đều được nối với cùng một nơ-ron. Quãng nghỉ xóa vết, và $B$ thuộc về một nơ-ron khác.}
\label{fig:tracevideo}
\end{figure}

Chúng tôi đã kiểm tra điều này trên mô hình (hình~\ref{fig:tracevideo}). Hai nơ-ron $C$ cạnh tranh giành đầu vào từ $16$ nơ-ron $S$: hai hình ở tám vị trí. Hình chạy qua mọi vị trí, mỗi vị trí $50\ms$, sau đó nghỉ $0.3\,\text{s}$ rồi đến một hình ngẫu nhiên tiếp theo. Không có nhãn nào. Nếu vết ngắn, $\tau_{\text{tr}}=5\ms$, các nơ-ron chia đầu vào theo \emph{vị trí}, và câu trả lời trùng với hình không hơn tung đồng xu: trung bình $52\%$ qua mười lần chạy. Nếu vết dài cỡ thời gian hiện một vị trí, từ $\tau_{\text{tr}}\approx20\ms$, mỗi nơ-ron gom \emph{mọi} vị trí của hình của mình (với $20$, $50$ và $200\ms$ là ở cả mười lần chạy): tính bất biến được học chỉ từ video. Quãng nghỉ giữa các đồ vật là quan trọng: không có nó, vết tràn sang đồ vật tiếp theo và làm lẫn chúng.

Điều tương tự thấy được trong thực nghiệm. Nếu trong thí nghiệm, đồ vật bị đánh tráo một cách kín đáo đúng lúc mắt nhảy từ chỗ này sang chỗ khác, thì sau một hai giờ, nơ-ron ở các tầng cao bắt đầu đáp ứng với đồ vật bị tráo như với chính đồ vật cũ: chúng nối những gì đến liên tiếp~\cite{LiDiCarlo2008}. Tính bất biến trong não quả thật được học từ sự liên tục của thời gian. Quy tắc này giải thích được bao nhiêu phần của toàn bộ việc học thị giác thì vẫn là giả thuyết.

\section{Video là chuyển động}

Video không chỉ là luồng khung hình để học, mà bản thân nó là bài toán: cái gì đang chuyển động, về đâu và nhanh thế nào. Bước đầu tiên ta đã quen: bộ phát hiện hướng của chương~\ref{sec:gaba} (hình~\ref{fig:direction}), trong đó ức chế đến trễ dập tắt chuyển động theo một chiều. Các bộ phát hiện như vậy phủ khắp trường nhìn, và sau đó chúng được xử lý giống như đặc trưng hình dạng: các tầng gom nhiều bộ phát hiện cùng hướng ở những vị trí khác nhau sẽ đáp ứng với chuyển động của một vật lớn bất kể nó ở đâu. Đó vẫn là cặp VÀ và HOẶC~\eqref{eq:scstage}, chỉ có đặc trưng không phải ``cạnh'' mà là ``cạnh đã dịch đi sau thời gian $d$''.

Video còn có thể dự đoán được: khung hình sau gần như giống khung hình trước. Theo giả thuyết \emph{mã hóa dự đoán}, các tầng trên gửi xuống tầng dưới một dự đoán, còn đi lên chủ yếu là sai số, tức phần mà dự đoán chưa tính tới~\cite{RaoBallard1999}. Đây vẫn là cách tiết kiệm xung như trong mệnh đề~\ref{prop:economy}: trả tiền cho tin mới, không trả cho những gì đã biết. Một số quan sát riêng lẻ phù hợp với giả thuyết này, nhưng như một cấu tạo chung của dây chuyền thì nó chưa được chứng minh.

\section{Não và mạng tích chập}

Sự giống nhau đi xa đến đâu? Với việc nhận dạng nhanh một đồ vật, nó quả thật sâu sắc~\cite{Yamins2014}. Nhưng não còn có những thứ mà mạng tích chập đơn giản không có.

\emph{Phản hồi.} Dây chuyền của não không chỉ đi thẳng: mỗi tầng có liên kết ngược và liên kết ngang. Với các ảnh khó (bị che khuất, thiếu sáng), đáp ứng của các tầng trên hình thành muộn hơn, và những đáp ứng muộn này được mô tả bằng mạng có phản hồi tốt hơn mạng chỉ đi thẳng~\cite{Kar2019}. Một lượt giải quyết được các trường hợp dễ, trường hợp khó cần vài vòng.

\emph{Tính bền vững.} Có thể đánh lừa mạng nhân tạo bằng cách sửa điểm ảnh mà mắt không nhận ra~\cite{Szegedy2014}: một thay đổi con người không thấy biến ``gấu trúc'' thành ``vượn''~\cite{Goodfellow2015}. Thị giác người bền vững hơn nhiều trước những kiểu sửa như vậy, nhưng không phải miễn nhiễm: ảnh đánh lừa được nhiều mạng cùng lúc, khi hiện ngắn, cũng làm lệch lựa chọn của người~\cite{Elsayed2018}. Vì sao độ bền vững khác nhau đến vậy là câu hỏi còn mở; các ứng viên là nhiễu, phép chia cho hoạt động chung và phản hồi.

\emph{Người thầy.} Mạng cần hàng triệu ví dụ có nhãn, não chủ yếu cần video không nhãn~\eqref{eq:traceinvariance} và một chút gợi ý từ \nt{DOPA} về điều gì là quan trọng.

\emph{Năng lượng.} Cả bộ não chỉ khoảng $20$ oát (mệnh đề~\ref{prop:economy}), và nhận dạng chỉ chiếm một phần trong đó; mạng tích chập đã huấn luyện chạy trên bộ tăng tốc đồ họa tiêu tốn hàng trăm oát.

\section{Ảo giác thị giác: cỗ máy sai ở đâu và vì sao}
\label{sec:illusions}

Kỹ sư muốn hiểu một mạch của người khác sẽ đưa vào đó những tín hiệu bất thường và xem mạch sai ở đâu. Với thị giác, những tín hiệu như vậy đã được nghĩ ra từ lâu: đó là các ảo giác thị giác. Ảo giác không phải hỏng hóc. Mỗi ảo giác chỉ tới một cơ chế nhất định của cỗ máy, cơ chế vốn làm việc đúng trong hoàn cảnh bình thường, nhưng lộ ra trên một bức ảnh được chọn đặc biệt.

\emph{Kỳ vọng hợp lý.} Đầu vào có nhiễu, và cỗ máy bổ sung cho nó bằng những gì thường gặp trong thế giới. Chuyển động chậm gặp nhiều hơn chuyển động nhanh; khi đầu vào không đáng tin (chẳng hạn ảnh có độ tương phản thấp), ước lượng tốc độ lệch về phía chậm, và vật mờ nhạt có vẻ chuyển động chậm hơn vật sáng rõ~\cite{Weiss2002}. Nhiều ảo giác độ sáng cũng được giải thích như vậy: ta không thấy độ sáng của điểm ảnh, mà thấy bề mặt nào thường ứng với độ sáng đó nhất trong quá khứ~\cite{Purves2011}. Bức ảnh chiếc váy mà người thấy trắng--vàng, người thấy xanh--đen cũng cùng cơ chế: bức ảnh mơ hồ, và mọi người khác nhau ở chỗ vô thức cho rằng ánh sáng chiếu vào như thế nào~\cite{LaferSousa2015,Wallisch2017}. Sự khác biệt giữa người với người đã được đo; việc não ở đây kết hợp đầu vào với kỳ vọng theo quy tắc của lý thuyết xác suất là một giả thuyết có cơ sở vững.

\emph{Điều chỉnh khuếch đại.} Nhìn thác nước một phút, rồi nhìn sang những hòn đá đứng yên: đá sẽ trôi lên trên. Kênh ``xuống'' và kênh ``lên'' là một cặp dây như trong~\eqref{eq:dualrail}; điều chỉnh khuếch đại~\eqref{eq:nakarushton} đã làm dịu kênh bị mỏi, và cân bằng của cặp bị lệch. Các ảo giác độ nghiêng và độ tương phản, trong đó vùng xung quanh làm thay đổi dáng vẻ vùng trung tâm, được giải thích bằng phép chia cho hoạt động của láng giềng~\cite{Schwartz2009}, như ở chương~\ref{sec:gaba}. Cũng có một ví dụ đáng học về sự thận trọng: các đốm tối ở giao điểm của những dải trắng trong lưới từ lâu được giải thích bằng ức chế láng giềng đơn giản, nhưng thí nghiệm với lưới biến dạng cho thấy cách giải thích đó không đúng~\cite{SchillerCarvey2005}.

\emph{Bù độ trễ.} Đường từ mắt lên các tầng trên mất hàng chục mili giây, và trong thời gian đó vật đang chuyển động kịp dịch đi. Nếu cạnh nó lóe lên một điểm đứng yên, vật có vẻ nằm phía trước điểm lóe, dù thật ra chúng trùng nhau~\cite{Nijhawan1994}. Theo một cách giải thích, cỗ máy ngoại suy chuyển động về phía trước để bù độ trễ, cũng như ở chương~\ref{sec:muscles}: khi phản hồi đến trễ thì phải dự đoán. Ảo giác này có nhiều cách giải thích, và tranh luận chưa ngã ngũ.

\emph{Lỗi trong mã thời gian.} Một bức ảnh tĩnh gồm các vòng tròn với chuyển màu sắc nét trông như đang quay. Vùng tương phản cao được xử lý nhanh hơn vùng nhạt, và chênh lệch độ trễ~\eqref{eq:latency} được các bộ phát hiện hướng đọc thành chuyển động~\cite{Conway2005}. Ảo giác được kích hoạt bởi các cú giật mắt nhỏ không chủ ý và cái chớp mắt: mỗi cú giật làm mới đầu vào, và các bộ phát hiện lại thấy ``chuyển động''~\cite{OteroMillan2012}. Đó là mã thời gian của chương~\ref{sec:code}, bị đọc sai.

\emph{Hai hình trong một.} Khung dây hình lập phương lúc quay mặt này về phía ta, lúc quay mặt kia, hoặc hai hình khác nhau ở hai mắt: tri giác nhảy qua lại vài giây một lần. Mô hình tốt nhất là ức chế lẫn nhau giữa hai cách hiểu, sự mỏi chậm và nhiễu~\cite{MorenoBote2007}, tức chính sơ đồ~\eqref{eq:halfcentre} đã tạo ra nhịp bước chân ở chương~\ref{sec:muscles}. Thời điểm chuyển do nhiễu và sự mỏi quyết định; theo mô hình~\cite{MorenoBote2007}, nhiễu đóng vai trò chính, còn sự mỏi chỉ hỗ trợ.

Còn mạng nhân tạo thì sao? Mạng được huấn luyện để dự đoán khung hình tiếp theo của video thấy chuyển động quay trên chính những vòng tròn tĩnh ấy và không thấy nó trên các ảnh đối chứng~\cite{Watanabe2018}. Mạng được huấn luyện để khử nhiễu và làm nét ảnh mắc cùng những ảo giác màu sắc và độ sáng như con người~\cite{GomezVilla2020}. Vậy một phần ảo giác là hệ quả của chính bài toán mà cỗ máy học, chứ không phải của đặc điểm mô thần kinh. Ảo giác nào chỉ riêng não mới có là phép thử tốt cho bất kỳ mô hình thị giác nào.

\begin{keyblock}
\begin{proposition}[Ảo giác là dấu vết của cơ chế]
\label{prop:illusions}
Ảo giác nảy sinh khi một cơ chế hữu ích trong thế giới bình thường nhận một đầu vào bất thường: kỳ vọng thắng đầu vào không đáng tin, điều chỉnh khuếch đại làm lệch một cặp kênh, bù độ trễ đẩy đồ vật lên trước, chênh lệch độ trễ được đọc thành chuyển động, ức chế lẫn nhau kèm nhiễu và sự mỏi nhảy qua lại. Mỗi ảo giác là một tín hiệu thử chỉ tới cơ chế của nó. Bản thân các ảo giác đã được đo; cách giải thích chúng đi từ có cơ sở vững đến còn tranh cãi.
\end{proposition}
\end{keyblock}

\section{Tổng kết}

\begin{table}[t]
\centering
\footnotesize
\setlength{\tabcolsep}{4pt}
\renewcommand{\arraystretch}{1.15}
\begin{tabular}{>{\raggedright}p{3.0cm}>{\raggedright}p{4.6cm}>{\raggedright\arraybackslash}p{5.2cm}}
\toprule
Cỗ máy làm gì & Bằng cách nào & Đã biết gì \\
\midrule
nhận ra trong $150\ms$ & một lượt qua chừng mười tầng & tốc độ đã đo \\
làm câu trả lời chọn lọc & VÀ theo bộ phận, tầng $S$~\eqref{eq:scstage} & đã đo ở các tầng đầu \\
làm câu trả lời bất biến & HOẶC theo vị trí, tầng $C$; ức chế và phép chia & đã đo ở các tầng đầu; với các tầng trên là đơn giản hóa \\
đưa ra câu trả lời & tần số của chừng một trăm nơ-ron tầng cuối, đọc tuyến tính & đã đo \\
học không cần nhãn & quy tắc Hebb có vết~\eqref{eq:traceinvariance} trên video liên tục & đánh tráo đồ vật làm đổi đáp ứng: đã đo; là quy tắc chính: giả thuyết \\
thấy chuyển động & bộ phát hiện hướng, rồi cùng cặp VÀ/HOẶC & đã đo \\
tiết kiệm trên phần dự đoán được & đi lên là sai số dự đoán & giả thuyết \\
giải trường hợp khó & phản hồi, vài vòng & đáp ứng muộn và vai trò của phản hồi đã đo \\
sai một cách dự đoán được & ảo giác: kỳ vọng, điều chỉnh khuếch đại, độ trễ, ức chế lẫn nhau kèm nhiễu và sự mỏi & ảo giác đã đo; giải thích: từ có cơ sở đến còn tranh cãi \\
\bottomrule
\end{tabular}
\caption{Cỗ máy nhận dạng hình ảnh và video như thế nào.}
\label{tab:recognition}
\end{table}

Nhận dạng được lắp từ những chi tiết ta đã có (bảng~\ref{tab:recognition}): ngưỡng cho phép VÀ theo bộ phận, ức chế lẫn nhau cho phép HOẶC theo vị trí, phép chia cho sự độc lập với độ sáng, quy tắc Hebb có vết thay cho người thầy vắng mặt. Kỹ sư đã lấy từ thị giác sự xen kẽ giữa tính chọn lọc và tính bất biến và dựng trên đó mạng tích chập; còn điều quan trọng nhất thì não vẫn chưa trao lại: khả năng học từ video không nhãn mà gần như không tốn năng lượng.

\section{Bài tập}

\begin{exercise}[\lvlA]
\label{ex:12:1}
\emph{Tần số trong một tầng.} Một tầng của dây chuyền kéo dài $15\ms$, nơ-ron phát $20$ xung mỗi giây, nhiễu tương quan với $c=0.1$~\eqref{eq:correlated}. Sai số tương đối của tần số nhóm một trăm nơ-ron bằng bao nhiêu, và giới hạn của nó khi nhóm lớn tùy ý? Trong một tầng phân biệt được bao nhiêu mức tần số?
\end{exercise}

\begin{exercise}[\lvlA]
\label{ex:12:2}
\emph{Năng lượng của một lần nhận dạng.} Trong một lượt $150\ms$, mười tầng, mỗi tầng $10^7$ nơ-ron, phát không quá một xung, mỗi xung tốn $10^{-10}$~J. (a)~Nhận dạng tốn bao nhiêu phần năng lượng của cả bộ não trong $150\ms$ đó? (b)~Bộ tăng tốc công suất $300$~W nhận dạng một ảnh trong $10\ms$ tốn gấp bao nhiêu lần?
\end{exercise}

\begin{exercise}[\lvlB]
\label{ex:12:3}
\emph{Ngưỡng của tầng $S$.} ``Võng mạc'' $24$ điểm, hình $A=11011$ và $B=10101$, mẫu~\eqref{eq:scstage} bằng $+1$ ở các số một của hình và $-1$ ở các số không. (a)~Với ngưỡng nào thì nơ-ron $S$ bỏ qua một điểm bị lật nhưng im lặng với hình kia? (b)~Cần bao nhiêu nơ-ron $S$ cho mười đặc trưng $5\times5$ trên ảnh $100\times100$?
\end{exercise}

\begin{exercise}[\lvlB]
\label{ex:12:4}
\emph{Một trong hai hình kéo dài bao lâu.} Trong sơ đồ nhảy~\eqref{eq:halfcentre} với $\tau\ll\tau_a$, khi nhóm~1 đang thắng thì $r_2=0$ và $r_1=I-a_1$. (a)~Tìm thời gian thắng với $I=1$, $\beta=2$, $b=2$, $\tau_a=0.4\,\text{s}$. (b)~$\tau_a$ nào cho hình đổi ba giây một lần?
\end{exercise}

\begin{exercise}[\lvlB]
\label{ex:12:5}
\emph{Mô phỏng: giá của tính bất biến.} Dùng hàm \texttt{two\_stage} trong \texttt{models/\allowbreak recognition\_models.py} ($4000$ lần thử), tìm xác suất lật điểm $p$ mà tại đó cặp tầng $S$, $C$ với $N=24$ sai một phần tư số trường hợp. Tỷ lệ trả lời đúng khi $p=0.1$ thay đổi ra sao từ $N=8$ đến $N=192$?
\end{exercise}

\begin{exercise}[\lvlB]
\label{ex:12:6}
\emph{Mô phỏng: cửa sổ của vết.} Mười lần chạy \texttt{trace\_learning} trong \texttt{models/\allowbreak recognition\_models.py} với quãng nghỉ $0.3\,\text{s}$: tìm các $\tau_{\text{tr}}$ trong khoảng $5\ms$ đến $2\,\text{s}$ mà mọi lần chạy đều học tính bất biến không sai. Cận trên của cửa sổ này từ đâu mà có?
\end{exercise}

\begin{exercise}[\lvlB]
\label{ex:12:7}
\emph{Chuyển động ở bất kỳ đâu.} Trên các điểm liền kề của một dãy cách nhau $0.1^\circ$ đặt các bộ phát hiện hướng của chương~\ref{sec:gaba}: ức chế từ $A$ với độ trễ $d$ dập tắt kích thích từ $B$ nếu chúng lệch nhau không quá $5\ms$; phía trên là tầng $C$. Cần độ trễ nào để $C$ im lặng với chuyển động ngược chiều ở tốc độ từ $5$ đến $20^\circ$ mỗi giây?
\end{exercise}

\begin{exercise}[\lvlC]
\label{ex:12:8}
\emph{Sửa điểm ảnh.} Cần lật bao nhiêu điểm để đổi câu trả lời của mô hình \texttt{two\_stage} ($N=24$) trên một hình sạch? Còn bộ phân loại ``đầu vào có hơn $3.5$ số một thì là $A$'', cũng bất biến với phép dịch? Tỷ lệ trả lời đúng của nó khi $p=0.1$ là bao nhiêu, so với $0.86$ của cặp tầng?
\end{exercise}

\chapter{Đầu ra của cỗ máy: não điều khiển cơ bắp như thế nào}
\chaptermark{Đầu ra của cỗ máy}
\label{sec:muscles}
\begin{chaptergoals}
\item cách cơ biến xung thành lực, như một bộ chuyển đổi số--tương tự kiêm bộ lọc thông thấp;
\item vì sao bật đơn vị theo kích thước đặt lực theo dấu phẩy động, với bước tỷ lệ với lực;
\item cách một cặp cơ kéo đặt mô-men bằng hiệu và độ cứng bằng tổng;
\item vì sao phản hồi có độ trễ giới hạn tốc độ hiệu chỉnh, và cách ức chế cùng sự mỏi tạo ra nhịp.
\end{chaptergoals}

\section{Đầu ra nhanh}

Mọi việc cỗ máy làm nhanh trong thế giới, nó đều làm bằng cơ. Lời nói, ánh nhìn, bước chân, nụ cười đều là sự co cơ. Đầu ra thứ hai, chậm, là hóa học của cơ thể, được phân tích ở chương~\ref{sec:chemoutput}. Trên hình~\ref{fig:machine}, đầu ra là mũi tên ``cơ''; giờ hãy xem phía sau mũi tên đó có gì.

Mắt xích cuối cùng của chuỗi là các nơ-ron \nt{ACET}, chính những nơ-ron được ghi trong bảng~\ref{tab:types} là ``đầu ra tới cơ''. Mỗi sợi cơ nhận đầu vào từ đúng một nơ-ron như vậy, còn một nơ-ron điều khiển một nhóm sợi, từ vài trăm đến vài nghìn~\cite{Feinstein1955mu}. Ở những cơ cần tinh chỉnh, đơn vị nhỏ hơn; ở các cơ lớn của chân, đơn vị lớn hơn. Nơ-ron cùng các sợi cơ của nó được gọi là \emph{đơn vị vận động}: đó là mẩu cơ nhỏ nhất mà não có thể bật riêng.

Khớp thần kinh của nơ-ron lên cơ được cấu tạo hoàn toàn khác với khớp thần kinh vỏ não ở chương~\ref{sec:code}. Ở đó phần lớn số xung bị mất. Ở đây mỗi xung giải phóng lượng acetylcholine gấp vài lần mức cần để sợi cơ đáp ứng~\cite{WoodSlater2001}. Đây là khớp thần kinh có \emph{hệ số an toàn}: một xung của nơ-ron là đúng một lần co của các sợi. Bên trong cỗ máy có thể chấp nhận các liên kết không tin cậy và sửa lỗi bằng sự dư thừa; ở đầu ra, một lỗi phải trả giá bằng một cử động, và độ tin cậy được mua thẳng bằng phần cứng.

\section{Lực: tần số và số đơn vị}

Một xung cho một lần co đơn lẻ, gọi là cú giật: lực tăng lên trong vài chục mili giây rồi giảm xuống. Một xấp xỉ tốt cho nó là~\cite{Fuglevand1993}
\begin{equation}
h(t) \;=\; F_1\,\frac{t}{T_c}\,e^{\,1-t/T_c},
\label{eq:twitch}
\end{equation}
trong đó $T_c$ là thời gian tăng, cỡ $50\ms$, còn $F_1$ là lực của một cú giật. Nếu xung đến dày hơn, các cú giật cộng lại:
\begin{empheq}[box=\keybox]{equation}
F(t) \;=\; \sum_k h\bigl(t-t_k\bigr) .
\label{eq:forcesum}
\end{empheq}
Đây chính là bộ lọc thông thấp đã biến xung thành nồng độ ở chương~\ref{sec:serocomputer}, chỉ có điều đầu ra giờ không phải nồng độ mà là lực. Cơ là bộ chuyển đổi số--tương tự từ xung sang lực (hình~\ref{fig:tetanus}). Trong mô hình của chúng tôi, với năm xung mỗi giây, các cú giật hầu như không chồng lên nhau và lực nhảy từ $0.2F_1$ đến $1.1F_1$; với mười lăm xung, lực đã giữ trong khoảng $1.8F_1$ đến $2.2F_1$; với bốn mươi, lực gần như đều, khoảng $5.4F_1$. Cơ thật bão hòa khi xung dày, nên tổng tuyến tính~\eqref{eq:forcesum} chỉ đúng tới vài chục xung mỗi giây.

\begin{figure}[t]
\centering
\begin{tikzpicture}[>=Stealth,x=20cm,y=0.55cm]
\draw[->,gray!70!black] (0,0) -- (0.66,0) node[right,font=\small,black] {$t$, s};
\draw[->,gray!70!black] (0,0) -- (0,6.4) node[left,font=\small,black] {$F/F_1$};
\foreach \t in {0,0.2,0.4,0.6}{\draw[gray!60!black] (\t,-0.1) -- (\t,0.1); \node[below,font=\scriptsize] at (\t,0) {\t};}
\foreach \f in {1,3,5}{\draw[gray!60!black] (-0.004,\f) -- (0.004,\f); \node[left,font=\scriptsize] at (0,\f) {\f};}
\draw[thick,blue!60!black] plot file {figures/muscle_twitch_5.dat};
\draw[thick,green!45!black] plot file {figures/muscle_twitch_15.dat};
\draw[thick,red!70!black] plot file {figures/muscle_twitch_40.dat};
\node[font=\scriptsize,blue!60!black,anchor=west] at (0.605,0.9) {5 xung/s};
\node[font=\scriptsize,green!45!black,anchor=west] at (0.605,2.1) {15 xung/s};
\node[font=\scriptsize,red!70!black,anchor=west] at (0.605,5.4) {40 xung/s};
\end{tikzpicture}
\caption{Cơ như một bộ chuyển đổi số--tương tự: lực~\eqref{eq:forcesum} với các xung đều đặn của một đơn vị vận động, $T_c=50\ms$. Xung thưa cho các cú giật rời rạc, xung dày hòa thành lực đều, giống như xung dày hòa thành nồng độ đều trên hình~\ref{fig:lowpass}.}
\label{fig:tetanus}
\end{figure}

Não có hai núm vặn lực: tần số xung của mỗi đơn vị và số đơn vị được bật. Núm thứ hai được thiết kế rất khéo. Các đơn vị trong một cơ khác nhau: đơn vị yếu nhất yếu hơn đơn vị mạnh nhất cả trăm lần. Khi cần lực ngày càng lớn, các đơn vị được bật \emph{theo thứ tự kích thước}, từ nhỏ đến lớn~\cite{Henneman1957}. Không ai tính ra thứ tự đó: nơ-ron nhỏ đơn giản là dễ bị kích thích hơn bởi cùng một đầu vào chung, nên thứ tự bật do chính phần cứng quy định.

Điều đó mang lại gì? Giả sử mỗi đơn vị kế tiếp theo thứ tự mạnh hơn đơn vị trước $q$ lần, với $q$ lớn hơn một một chút. Lực của mọi đơn vị đang bật là tổng của một cấp số nhân, và gần như toàn bộ rơi vào các đơn vị sau cùng. Khi đó đơn vị được bật tiếp theo thêm vào
\begin{empheq}[box=\keybox]{equation}
\frac{\Delta F}{F} \;\approx\; q-1 \;=\; \text{const} ,
\label{eq:sizeprinciple}
\end{empheq}
nghĩa là bước lực tỉ lệ với chính lực. Khi gắng sức nhẹ, não định liều lực bằng những phần nhỏ, khi gắng sức mạnh thì bằng những phần lớn. Đó là \emph{số dấu phẩy động}: bậc độ lớn do số đơn vị được bật quyết định, còn phần định trị do tần số xung của chúng. Đây cũng là thang logarit như ở các cảm biến của chương~\ref{sec:sensors}: đầu vào và đầu ra của cỗ máy đo thế giới bằng tỷ lệ tương đối.

Dấu phẩy động có mặt trái. Độ tản của lực cũng tăng tỉ lệ với chính lực: cử động mạnh chắc chắn kém chính xác hơn cử động nhẹ. Theo một giả thuyết có ảnh hưởng, chính nhiễu phụ thuộc tín hiệu này giải thích vì sao cử động của chúng ta mượt: một lệnh mượt cho độ tản nhỏ nhất ở điểm cuối~\cite{HarrisWolpert1998}.

\section{Kéo chứ không đẩy: hai dây cho một khớp}

Cơ chỉ biết kéo. Nó không thể đẩy, cũng như nơ-ron \nt{GLUT} không thể phát tín hiệu âm. Giải pháp đã quen từ chương~\ref{sec:glutcomputer}: \emph{hai dây}. Mỗi khớp được một cặp cơ kéo theo hai chiều ngược nhau phục vụ (hình~\ref{fig:antagonists}). Nếu lực của chúng là $F_1$ và $F_2$, còn cánh tay đòn là $\ell$, thì mô-men quay và độ cứng của khớp bằng
\begin{empheq}[box=\keybox]{equation}
M \;=\; \ell\,(F_1-F_2), \qquad K \;\approx\; \kappa_m\,(F_1+F_2),
\label{eq:antagonists}
\end{empheq}
trong đó $\kappa_m$ là hệ số liên hệ lực của cơ với độ cứng của nó: cơ căng thì đàn hồi mạnh hơn cơ chùng. Hiệu các lực quyết định mô-men, tổng các lực quyết định độ cứng~\cite{Hogan1984}. Bằng hai dây, não điều khiển hai đại lượng độc lập: xoay khớp về đâu và giữ nó chắc đến mức nào. Khi cần cầm tách mà không làm sánh, ta căng cả hai cơ cùng lúc: mô-men vẫn thế, độ cứng cao hơn, và một cú va ngẫu nhiên ít làm lệch tay hơn.

\begin{figure}[t]
\centering
\begin{tikzpicture}[>=Stealth,
  nrn/.style={circle,draw,very thick,minimum size=0.95cm,inner sep=0pt,font=\scriptsize},
  inh/.style={-{Bar[width=7pt]},thick,red!70!black}]
% the joint: a bar on a pivot, two springs pulling down on either side
\fill[gray!30] (4.6,-0.25) -- (5.4,-0.25) -- (5,0.15) -- cycle;
\draw[line width=3pt,gray!50!black] (2.4,0.2) -- (7.6,0.2);
\fill[black] (5,0.2) circle (0.08);
\draw[decorate,decoration={coil,segment length=5pt,amplitude=4pt},very thick,blue!60!black] (3.0,0.2) -- (3.0,-1.6);
\draw[decorate,decoration={coil,segment length=5pt,amplitude=4pt},very thick,orange!85!black] (7.0,0.2) -- (7.0,-1.6);
\draw[very thick] (2.5,-1.6) -- (3.5,-1.6); \draw[very thick] (6.5,-1.6) -- (7.5,-1.6);
\node[font=\small,blue!60!black] at (2.35,-0.75) {$F_1$};
\node[font=\small,orange!85!black] at (7.65,-0.75) {$F_2$};
\draw[->,thick] (5.9,0.75) arc (60:120:1.8) node[above,font=\scriptsize,pos=0.5] {$M=\ell(F_1-F_2)$};
% motor neurons and reflex circuit
\node[nrn,fill=green!12] (M1) at (1.6,-3.0) {\nt{ACET}};
\node[nrn,fill=green!12] (M2) at (8.4,-3.0) {\nt{ACET}};
\node[nrn,fill=red!10] (G) at (5.0,-3.0) {\nt{GLYC}};
\node[nrn,fill=blue!6] (S) at (1.6,-5.0) {\nt{GLUT}};
\draw[->,very thick] (M1) -- (3.0,-1.7);
\draw[->,very thick] (M2) -- (7.0,-1.7);
\draw[->,thick] (S) -- (M1);
\draw[->,thick] (S) -- (G);
\draw[inh] (G) -- (M2);
\draw[->,thick,densely dashed,gray!60!black] (3.15,-1.2) to[bend left=25] node[pos=0.35,right,font=\scriptsize,black] {độ giãn} (S.north east);
\draw[->,thick,blue!60!black] (-0.3,-3.0) node[left,font=\scriptsize,black,align=right] {lệnh\\từ não} -- (M1);
\draw[->,thick,blue!60!black] (10.3,-3.0) node[right,font=\scriptsize,black,align=left] {lệnh\\từ não} -- (M2);
\end{tikzpicture}
\caption{Hai cơ cho một khớp và vòng phản hồi nhanh nhất. Các nơ-ron \nt{ACET} nhận lệnh từ não và kéo khớp về hai phía; hiệu các lực làm khớp quay, tổng làm khớp cứng hơn. Cảm biến độ giãn (\nt{GLUT}) của cơ bên trái kích thích nơ-ron \nt{ACET} của chính cơ đó và qua nơ-ron trung gian \nt{GLYC} ức chế nơ-ron của cơ đối vận: phép cấm của chương~\ref{sec:gaba}.}
\label{fig:antagonists}
\end{figure}

\section{Phản hồi có độ trễ}

Cơ không làm việc mù quáng. Trong mỗi cơ đều có cảm biến độ giãn, như đã nói trong bảng~\ref{tab:sensors}, và vòng ngắn nhất khép kín ngay trong tủy sống (hình~\ref{fig:antagonists}). Cơ bị kéo giãn bất ngờ, thì cảm biến độ giãn kích thích nơ-ron \nt{ACET} của chính cơ đó, cơ co lại và đưa khớp về chỗ cũ. Đồng thời, qua một nơ-ron trung gian ức chế, cơ đối vận bị tắt để khỏi cản trở. Nơ-ron trung gian ấy chính là \nt{GLYC}, loại trong bảng~\ref{tab:types} không có chương riêng: chỗ của nó chính là ở đây, trong các vòng nhanh của tủy sống.

Mỗi vòng có một độ trễ $d$: vòng tủy sống khoảng $25$--$30\ms$ với cánh tay, vòng qua vỏ não dài hơn một rưỡi đến ba lần, vòng qua mắt $100$--$200\ms$. Mà phản hồi có trễ, như ta đã thấy ở mệnh đề~\ref{prop:ei}, làm hệ dao động. Với bộ điều khiển đơn giản nhất, sửa độ lệch $x$ với tốc độ $G$ dựa trên thông tin cũ $d$ giây, $\dd x/\dd t=-G\,x(t-d)$, điều kiện ổn định là~\cite{Hayes1950}:
\begin{empheq}[box=\keybox]{equation}
G\,d \;<\; \frac{\pi}{2} ,
\label{eq:delaystable}
\end{empheq}
và tại ranh giới, hệ dao động với tần số $1/(4d)$. Chúng tôi đã kiểm tra điều này trên mô hình: với $d=30\ms$, độ lệch tắt dần khi $G=40$ mỗi giây, còn khi $G=55$ mỗi giây, cao hơn ranh giới $52$ một chút, thì dao động tăng dần.

Từ đó có hai hệ quả. Thứ nhất: vòng càng dài thì càng phải sửa chậm. Qua mắt, với độ trễ một trăm rưỡi mili giây, không thể chỉnh tay nhanh hơn khoảng một phần mười giây, nếu không tay sẽ run. Thứ hai: ranh giới ổn định của vòng tủy sống, $1/(4\cdot30\ms)\approx8$ dao động mỗi giây, gần với tần số run nhẹ mà người khỏe mạnh nào cũng có, tám đến mười hai dao động mỗi giây. Vòng phản xạ giãn là một trong những nguồn có khả năng gây ra cơn run này~\cite{McAuleyMarsden2000}.

Vậy làm sao não cử động nhanh và chính xác khi phản hồi đến trễ? Theo một giả thuyết phổ biến, não \emph{dự đoán}: nó giữ một mô hình bên trong của cơ thể và tính trước kết quả của lệnh mà không chờ cảm biến~\cite{WolpertGhahramani2000}. Cảm biến cần để sửa mô hình, chứ không để sửa từng cử động. Kỹ sư sẽ gọi đó là điều khiển thuận kết hợp bộ quan sát trạng thái.

\section{Bộ tạo nhịp vận động}

Đi bộ, thở, nhai là những cử động có nhịp. Chúng có cần đồng hồ không? Không. Trong tủy sống và thân não có những mạng nhỏ tự sinh ra nhịp, kể cả khi không có gì có nhịp đi vào~\cite{MarderBucher2001}. Mạng đơn giản nhất loại này được lắp từ những chi tiết ta đã có: hai nhóm nơ-ron \nt{GLUT} ức chế nhau qua các nơ-ron trung gian \nt{GABA} hoặc \nt{GLYC}, như trên hình~\ref{fig:wta}, và mỏi dần, như bộ nhớ ở chương~\ref{sec:glutcomputer}:
\begin{empheq}[box=\keybox]{equation}
\tau\,\frac{\dd r_i}{\dd t} = -r_i + \bigl[I - \beta\,r_j - a_i\bigr]_+ ,
\qquad
\tau_a\,\frac{\dd a_i}{\dd t} = -a_i + b\,r_i ,
\label{eq:halfcentre}
\end{empheq}
trong đó $a_i$ là độ mỏi của nhóm $i$, $j$ là nhóm kia. Ức chế lẫn nhau với $\beta>1$ chọn ra kẻ thắng (mệnh đề~\ref{prop:wta}). Kẻ thắng làm việc và mỏi dần; khi độ mỏi ăn hết đầu vào của nó, kẻ thua, lúc này đã nghỉ ngơi, vượt lên và đến lượt mình bắt đầu mỏi. Các nhóm làm việc luân phiên, và mỗi nhóm điều khiển một cơ của cặp (hình~\ref{fig:halfcentre}).

\begin{figure}[t]
\centering
\begin{tikzpicture}[>=Stealth,x=3.2cm,y=2.4cm]
\draw[->,gray!70!black] (0,0) -- (4.2,0) node[right,font=\small,black] {$t$, s};
\draw[->,gray!70!black] (0,0) -- (0,1.2) node[left,font=\small,black] {$r$};
\foreach \t in {0,1,2,3,4}{\draw[gray!60!black] (\t,-0.03) -- (\t,0.03); \node[below,font=\scriptsize] at (\t,0) {\t};}
\draw[thick,blue!60!black] plot file {figures/halfcentre1.dat};
\draw[thick,orange!85!black] plot file {figures/halfcentre2.dat};
\node[font=\scriptsize,blue!60!black,anchor=west] at (1.6,1.28) {nhóm 1: ``tiến''};
\node[font=\scriptsize,orange!85!black,anchor=west] at (2.7,1.28) {nhóm 2: ``lùi''};
\end{tikzpicture}
\caption{Bộ tạo nhịp theo mô hình~\eqref{eq:halfcentre}: $I=1$, $\beta=2$, $b=2$, $\tau=20\ms$, $\tau_a=0.4\,\text{s}$. Hai nhóm ức chế lẫn nhau và mỏi dần làm việc luân phiên với chu kỳ $0.84\,\text{s}$, dù tín hiệu đầu vào là không đổi.}
\label{fig:halfcentre}
\end{figure}

Trong mô hình của chúng tôi, với đầu vào không đổi, các nhóm thay phiên nhau với chu kỳ $0.84\,\text{s}$, khoảng gấp đôi thời gian mỏi $\tau_a$. Một lệnh không đổi từ trên xuống, ``đi'', được biến thành nhịp ngay tại chỗ. Đây là thêm một nhịp cục bộ, như nhịp của vòng \nt{GLUT}--\nt{GABA} ở chương~\ref{sec:gaba}: nó chỉ xuất hiện khi mạng làm việc, và chỉ giữ nhịp cho những cơ mà nó điều khiển.

\begin{keyblock}
\begin{proposition}[Đầu ra của cỗ máy]
\label{prop:muscles}
Não điều khiển cơ như một hệ phân cấp các bộ điều khiển. Ở trên cùng, hành động được chọn (chương~\ref{sec:dofa}); thấp hơn, các bộ tạo nhịp cục bộ biến lệnh không đổi thành nhịp, còn các vòng nhanh của tủy sống giữ các khớp. Lực được đặt theo dấu phẩy động: bậc độ lớn do số đơn vị được bật theo kích thước, phần định trị do tần số xung. Mỗi khớp được điều khiển bởi một cặp cơ kéo, hiệu lực của chúng quyết định mô-men, còn tổng quyết định độ cứng. Các cơ chế này đã được đo; việc não dựa vào một mô hình bên trong của cơ thể là một giả thuyết có cơ sở vững.
\end{proposition}
\end{keyblock}

\section{Tổng kết}

\begin{table}[t]
\centering
\footnotesize
\setlength{\tabcolsep}{4pt}
\renewcommand{\arraystretch}{1.15}
\begin{tabular}{>{\raggedright}p{3.0cm}>{\raggedright}p{4.6cm}>{\raggedright\arraybackslash}p{5.2cm}}
\toprule
Đầu ra làm gì & Bằng cách nào & Đã biết gì \\
\midrule
biến xung thành lực & tổng các cú giật~\eqref{eq:forcesum}, bộ lọc thông thấp & đã đo; tổng tuyến tính là đơn giản hóa \\
định liều lực & bật đơn vị theo kích thước, dấu phẩy động~\eqref{eq:sizeprinciple} & thứ tự bật đã đo \\
đẩy dù chỉ biết kéo & cặp cơ: mô-men là hiệu, độ cứng là tổng~\eqref{eq:antagonists} & đã đo \\
giữ khớp & vòng phản xạ giãn, cấm cơ đối vận qua \nt{GLYC} & đã đo \\
không run & $Gd<\pi/2$~\eqref{eq:delaystable} & suy ra từ lý thuyết điều khiển; góp phần gây run là nguyên nhân có khả năng \\
cử động nhanh & dự đoán theo mô hình bên trong & giả thuyết có cơ sở vững \\
đi và thở có nhịp & bộ tạo nhịp~\eqref{eq:halfcentre} & các bộ tạo nhịp đã đo; sơ đồ đơn giản nhất là mô hình \\
\bottomrule
\end{tabular}
\caption{Cỗ máy điều khiển cơ như thế nào.}
\label{tab:muscles}
\end{table}

Đầu ra của cỗ máy được xây bằng chính những kỹ thuật như mọi phần khác (bảng~\ref{tab:muscles}): hai dây thay cho dấu, bộ lọc thông thấp thay cho DAC, thang logarit thay cho thang tuyến tính, ức chế lẫn nhau và sự mỏi thay cho bộ phát xung nhịp. Đầu vào (chương~\ref{sec:sensors}) và đầu ra đo thế giới bằng tỷ lệ tương đối, còn giữa chúng là cỗ máy mà cuốn sách này dành để nói về.

\section{Bài tập}

\begin{exercise}[\lvlA]
\label{ex:13:1}
\emph{Dấu phẩy động bằng con số.} Một cơ có $N=100$ đơn vị vận động với lực $f_n=f_1q^{\,n-1}$; đơn vị mạnh nhất mạnh gấp $100$ lần đơn vị yếu nhất. (a)~Tìm $q$, bước tương đối~\eqref{eq:sizeprinciple} và dải động tính bằng decibel. (b)~Ở lực $10f_1$, ``DAC tuyến tính'' gồm $100$ đơn vị giống nhau có cùng tổng lực có bước bằng bao nhiêu?
\end{exercise}

\begin{exercise}[\lvlA]
\label{ex:13:2}
\emph{Giá của hệ số an toàn.} Với mỗi xung, khớp thần kinh trên sợi cơ giải phóng một số gói tuân theo phân phối Poisson với trung bình $m$, còn sợi cơ đáp ứng khi có từ $n_{\text{ngưỡng}}=20$ gói trở lên; hệ số an toàn $S=m/n_{\text{ngưỡng}}$. Một đơn vị làm việc ở $20$ xung mỗi giây gây bao nhiêu lần hỏng mỗi ngày khi $S=2$ và khi $S=4$?
\end{exercise}

\begin{exercise}[\lvlB]
\label{ex:13:3}
\emph{Cơ như bộ lọc.} Cú giật~\eqref{eq:twitch} với $T_c=50\ms$ là đáp ứng xung của một hệ tuyến tính. (a)~Tìm hàm truyền $H(s)$ và tần số cắt của nó. (b)~Với tần số xung đều bằng bao nhiêu thì biên độ dao động của lực bằng $10\%$ lực trung bình?
\end{exercise}

\begin{exercise}[\lvlB]
\label{ex:13:4}
\emph{Không sửa nhanh hơn độ trễ.} Bộ điều khiển $\dd x/\dd t=-G\,x(t-d)$. (a)~Chứng minh rằng độ lệch tắt nhanh nhất mà không dao động khi $Gd=1/e$, với tốc độ $1/d$. (b)~Tìm $G$ đó và hằng số thời gian tắt dần cho vòng tủy sống, $d=30\ms$, và cho vòng qua mắt, $d=150\ms$.
\end{exercise}

\begin{exercise}[\lvlB]
\label{ex:13:5}
\emph{Chu kỳ của bộ tạo nhịp.} Khi $\tau\ll\tau_a$, chu kỳ $T$ của mô hình~\eqref{eq:halfcentre} cho bởi phương trình $(\beta-1)(1-E^{2+b})/(\beta-E)=b\,(1-E^{1+b})/(1+b)$, với $E=e^{-T/(2\tau_a)}$. (a)~Tìm $T$ khi $\beta=b=2$, $\tau_a=0.4$~s. (b)~Chứng minh rằng $T$ không phụ thuộc đầu vào $I$, và nhịp chỉ có thể có khi $b>\beta-1$.
\end{exercise}

\begin{exercise}[\lvlB]
\label{ex:13:6}
\emph{Mô phỏng bộ tạo nhịp.} Nhập hàm \texttt{halfcentre} từ \texttt{models/\allowbreak muscle\_models.py} (đừng chạy chính tệp đó) và đo chu kỳ, bỏ hai chu kỳ đầu, khi $b=0.8$, $1.05$, $1.2$, $1.5$, $2$, $4$ và khi $I=0.5$, $1$, $2$. Lệnh ``đi nhanh hơn'' có thể đổi nhịp độ thông qua $I$ không?
\end{exercise}

\begin{exercise}[\lvlB]
\label{ex:13:7}
\emph{Độ cứng đấu với phản xạ.} Khớp nằm trong vòng phản xạ giãn: $\dd\theta/\dd t=-a\theta(t)-G\theta(t-d)$, $a=K/B$, $d=30\ms$. Hãy đưa về bài tập~\ref{ex:13:4}, tìm $a$ nhỏ nhất để độ lệch tắt không dao động với hằng số thời gian $15\ms$, và $G$ cần thiết. Nên đổi $G$ thế nào khi cơ đồng co mạnh hơn?
\end{exercise}

\begin{exercise}[\lvlC]
\label{ex:13:8}
\emph{Núm vặn nhịp độ.} Theo bài tập~\ref{ex:13:5} và~\ref{ex:13:6}, đầu vào $I$ của bộ tạo nhịp~\eqref{eq:halfcentre} chỉ đổi biên độ, không đổi nhịp độ. Hãy đề xuất thay đổi nhỏ nhất của mô hình, bằng chi tiết trong sách, để dưới một $I_{\min}$ nào đó không có nhịp, còn trên đó chu kỳ giảm khi $I$ tăng, ít nhất hai lần khi $I$ tăng gấp ba. Tìm $I_{\min}$ khi $\tau\ll\tau_a$ và kiểm tra bằng mô phỏng.
\end{exercise}

\chapter{Đau: tín hiệu báo động}
\chaptermark{Đau}
\label{sec:pain}
\begin{chaptergoals}
\item vì sao đau là tín hiệu báo động có mức ưu tiên cao chứ không phải kênh đo;
\item cách một dây nhanh và một dây chậm báo tổn thương ở đâu và tệ đến mức nào;
\item cách cổng trong tủy sống cho xúc giác cấm cơn đau, và cách kênh quảng bá đặt âm lượng báo động;
\item cách nhạy cảm hóa tăng hệ số khuếch đại sau tổn thương, và cách cơn đau dạy và ngắt.
\end{chaptergoals}

\section{Báo động, không phải phép đo}

Mọi cảm biến ở chương~\ref{sec:sensors} đều đo thế giới: bao nhiêu ánh sáng, âm cao bao nhiêu, áp lực bao nhiêu. Cơn đau được tổ chức khác. Nó không báo ``có gì ở đó'', mà báo ``nguy hiểm, bỏ hết mọi việc''. Với kỹ sư, đây không phải một kênh đo mà là một \emph{tín hiệu báo động}: một đường dây có mức ưu tiên cao nhất, phải xuyên qua được mọi công việc đang chạy của cỗ máy.

Hệ báo động khác kênh đo ở ba đặc tính, và cơn đau có cả ba. Thứ nhất, khó làm nó dịu đi bằng thích nghi: cảm biến đau thích nghi yếu hơn nhiều so với cảm biến xúc giác, và sau một tổn thương nhẹ còn trở nên nhạy hơn~\cite{LaMotte1982hyper}; sự suy giảm đáp ứng sau khi bị nung nóng mạnh chỉ đo được ở một loại cảm biến đau~\cite{Treede1995heat}. Cảm biến ánh sáng im lặng trên nền không đổi (hình~\ref{fig:adaptation}), còn một hệ báo động im sau một phút thì vô dụng. Thứ hai, nó có ngưỡng cao: cảm biến đau chỉ phát xung khi tác động đã gần mức nguy hiểm, chẳng hạn với nhiệt, ngưỡng của các cảm biến khác nhau từ $41^\circ$ đến $46$--$48^\circ$, tùy loại~\cite{Treede1995heat, Weidner1999cfib}. Thứ ba, và đây là điều chính của chương này, âm lượng của báo động không chỉ do cảm biến quyết định mà còn do chính cỗ máy. Cùng một vết thương đau khác nhau trong những hoàn cảnh khác nhau. Hãy xem nó được lắp từ những chi tiết quen thuộc nào.

Cảm biến đau là nơ-ron \nt{GLUT}, như các nơ-ron cảm giác khác ở chương~\ref{sec:sensors}. Một số trong đó giải phóng cùng với glutamate chất~P ở phụ lục~\ref{sec:othermediators}, chất làm tăng cường truyền dẫn một cách chậm rãi.

\section{Hai dây: đau nhanh và đau chậm}

Đập vào ngón tay, bạn sẽ thấy đau hai lần: trước là đau ngắn và nhói, một giây sau là đau âm ỉ và kéo dài. Đó là hai dây khác nhau. Dây nhanh có lớp cách điện và dẫn xung với tốc độ $v$ từ $5$ đến $30$ mét mỗi giây; dây chậm thì mảnh và trần, $0.5$--$2$ mét mỗi giây. Tín hiệu từ bàn chân đi khoảng một mét, và thời gian đến
\begin{empheq}[box=\keybox]{equation}
t \;=\; \frac{L}{v}
\label{eq:paintime}
\end{empheq}
với dây nhanh là vài chục mili giây, với dây chậm là khoảng một giây. Hai dây mang hai thông điệp. Dây nhanh nói \emph{ở đâu} và \emph{khi nào}: xung của nó đến sớm hơn và chính xác hơn, như trong mã thời gian ở chương~\ref{sec:code}. Dây chậm nói \emph{tệ đến mức nào}, và cơn đau âm ỉ kéo dài của nó giữ cỗ máy ở chế độ báo động cho tới khi tổn thương qua đi.

\section{Cổng trong tủy sống}

Nơi đầu tiên tín hiệu đau tới là tủy sống, và ở đó nó đi qua một \emph{cổng}~\cite{MelzackWall1965}; các nơ-ron cụ thể của cổng này mới được tìm ra tương đối gần đây~\cite{Duan2014}. Trên nơ-ron \nt{GLUT} đầu ra, nơ-ron truyền cơn đau tiếp lên não, hội tụ dây đau và dây xúc giác. Còn bên cạnh là một nơ-ron trung gian ức chế, \nt{GABA} hoặc \nt{GLYC}, được xúc giác kích thích (hình~\ref{fig:gate}). Xúc giác đóng cổng lại. Trong thuyết cổng ban đầu~\cite{MelzackWall1965}, ngược lại, cơn đau tắt nơ-ron trung gian này, và đau mạnh thì mở cổng; đây là một giả định của thuyết, chưa được chứng minh trực tiếp trên các nơ-ron cổng đã tìm thấy.

\begin{figure}[t]
\centering
\begin{tikzpicture}[>=Stealth,
  nrn/.style={circle,draw,very thick,minimum size=1.0cm,inner sep=0pt,font=\scriptsize},
  inh/.style={-{Bar[width=7pt]},thick,red!70!black},
  bc/.style={->,thick,densely dashed,orange!85!black}]
\node[nrn,fill=blue!6] (Z) at (6,0) {\nt{GLUT}};
\node[nrn,fill=red!10] (Q) at (3.6,-1.6) {\nt{GABA}};
\draw[->,very thick,red!70!black] (0,0.6) node[left,font=\small,black,align=right] {đau $\nu$} -- (5.45,0.15);
\draw[->,very thick,blue!60!black] (0,-2.6) node[left,font=\small,black,align=right] {chạm $\mu$} -- (2.9,-2.0);
\draw[->,thick,blue!60!black] (1.8,-2.35) -- (5.5,-0.35) node[pos=0.8,below right=-2pt,font=\scriptsize] {$w_{z\mu}$};
\draw[inh] (1.6,0.5) -- (3.35,-1.1) node[pos=0.5,left=1pt,font=\scriptsize] {$w_{q\nu}$};
\draw[inh] (Q) -- (Z) node[pos=0.5,above left=-1pt,font=\scriptsize] {$\beta$};
\draw[->,very thick] (Z) -- (8.2,0) node[right,font=\small,align=left] {lên não: $z$};
\draw[bc] (6,2.1) node[above,font=\scriptsize,black,align=center] {từ trên: \nt{SERO}, \nt{NORA}, opioid} -- (Z.north) node[pos=0.45,right,font=\scriptsize,black] {khuếch đại $g$};
\end{tikzpicture}
\caption{Cổng đau trong tủy sống theo mô hình~\eqref{eq:gate}. Nơ-ron \nt{GLUT} đầu ra nhận cơn đau $\nu$ và kích thích yếu từ cái chạm $\mu$. Nơ-ron trung gian ức chế (\nt{GABA} hoặc \nt{GLYC}) được cái chạm kích thích và, theo giả định của thuyết cổng, bị cơn đau tắt đi. Từ trên, qua các kênh quảng bá, đến hệ số khuếch đại $g$.}
\label{fig:gate}
\end{figure}

Hãy viết điều đó cho tần số, như ở chương~\ref{sec:gaba}. Tần số $q$ của nơ-ron trung gian và tần số đầu ra $z$ bằng
\begin{empheq}[box=\keybox]{equation}
q \;=\; \bigl[w_{q\mu}\,\mu + q_0 - w_{q\nu}\,\nu\bigr]_+ ,
\qquad
z \;=\; \bigl[\,g\,(\nu + w_{z\mu}\,\mu) - \beta\,q\,\bigr]_+ ,
\label{eq:gate}
\end{empheq}
trong đó $\nu$ là đầu vào đau, $\mu$ là đầu vào xúc giác, $w$ là trọng số liên kết (chỉ số thứ nhất là nơi nhận, chỉ số thứ hai là nơi gửi), $\beta$ là cường độ ức chế, $q_0$ là hoạt động nền riêng của nơ-ron trung gian, $g$ là hệ số khuếch đại được đặt từ trên (sẽ nói ở dưới). Đây là phép cấm~\eqref{eq:veto} của chương~\ref{sec:gaba}: ``đau, nhưng không khi có chạm'', với một sửa đổi lấy từ thuyết cổng làm giả định: đau mạnh tự tắt nơ-ron cấm.

Chúng tôi đã tính mô hình với $w_{q\mu}=1$, $q_0=0.2$, $w_{q\nu}=0.5$, $w_{z\mu}=0.3$, $\beta=1.5$ (hình~\ref{fig:gatecurves}). Không có chạm, cơn đau bắt đầu đi qua từ $\nu\approx0.18$. Có chạm, ngưỡng tăng lên $0.86$, và khi $\nu=1$ đầu ra giảm bốn lần, từ $1$ xuống $0.25$. Còn khi đau mạnh, $\nu=2$, cái chạm không thay đổi được gì nữa: cổng đã mở toang. Ngoài đời cũng vậy: xoa chỗ bị va giúp đỡ đau, còn với chấn thương nặng thì không. Bản thân cái chạm, không có đau, không đi qua được cổng: báo động không nổi lên vì một cái chạm.

\begin{figure}[t]
\centering
\begin{tikzpicture}[>=Stealth,x=5.2cm,y=1.35cm]
\draw[->,gray!70!black] (0,0) -- (2.12,0) node[right,font=\small,black] {đau $\nu$};
\draw[->,gray!70!black] (0,0) -- (0,4.3) node[left,font=\small,black] {$z$};
\foreach \t in {0,0.5,1,1.5,2}{\draw[gray!60!black] (\t,-0.05) -- (\t,0.05); \node[below,font=\scriptsize] at (\t,0) {\t};}
\foreach \f in {1,2,3,4}{\draw[gray!60!black] (-0.01,\f) -- (0.01,\f); \node[left,font=\scriptsize] at (0,\f) {\f};}
\draw[very thick,red!70!black] plot file {figures/pain_gate_notouch.dat};
\draw[very thick,blue!60!black] plot file {figures/pain_gate_touch.dat};
\draw[very thick,orange!85!black,densely dashed] plot file {figures/pain_gate_sens.dat};
\draw[very thick,red!70!black] (0.08,3.9) -- (0.2,3.9) node[right,font=\scriptsize,black] {không chạm};
\draw[very thick,blue!60!black] (0.08,3.45) -- (0.2,3.45) node[right,font=\scriptsize,black] {có chạm};
\draw[very thick,orange!85!black,densely dashed] (0.08,3.0) -- (0.2,3.0) node[right,font=\scriptsize,black] {sau nhạy cảm hóa, $g=2$};
\end{tikzpicture}
\caption{Đầu ra của cổng~\eqref{eq:gate} theo cường độ đau. Cái chạm nâng ngưỡng và làm dịu cơn đau nhẹ và vừa, nhưng không làm dịu cơn đau mạnh. Nhạy cảm hóa (đường nét đứt) tăng hệ số khuếch đại gấp đôi: cùng một cơn đau được truyền to gấp đôi, còn ngưỡng hạ xuống.}
\label{fig:gatecurves}
\end{figure}

\section{Núm âm lượng từ trên}

Cổng không chỉ được điều khiển từ dưới. Từ thân não có các đường đi xuống tới cổng, làm thay đổi hệ số khuếch đại $g$ trong~\eqref{eq:gate}: \nt{SERO}, \nt{NORA} và các peptide opioid ở phụ lục~\ref{sec:othermediators}~\cite{Fields2004}. Đó chính là các kênh quảng bá của chương~\ref{sec:serocomputer}--\ref{sec:acet}, chỉ có điều ở đây núm vặn của chúng là âm lượng báo động. Chúng vặn được cả hai chiều: cùng những mạch đi xuống ấy có thể vừa làm dịu cơn đau, vừa khuếch đại nó.

Vì sao cỗ máy lại giảm báo động? Vì báo động là một ngắt, mà ngắt không phải lúc nào cũng hợp lúc. Con vật đang chạy trốn kẻ săn mồi không được dừng lại vì một chân bị thương: chạy trước, liếm vết thương sau. Sự trông đợi được đỡ đau cũng làm dịu cơn đau, và một phần hiệu ứng này mất đi nếu chặn các thụ thể opioid~\cite{Levine1978}: kỳ vọng tính ra trong não đi xuống theo kênh opioid tới cổng. Bức tranh này tự nó đã được đo; còn não quyết định âm lượng phải bằng bao nhiêu như thế nào thì vẫn là giả thuyết.

\section{Nhạy cảm hóa: hệ báo động biết học}

Sau tổn thương, các nơ-ron đầu ra của cổng trở nên nhạy hơn: cùng đầu vào cho đầu ra lớn hơn, còn ngưỡng hạ xuống~\cite{Woolf1983}. Trong mô hình~\eqref{eq:gate}, đó là sự tăng của hệ số khuếch đại $g$, và mô tả đơn giản nhất của nó là một mạch RC chậm được chính cơn đau nạp:
\begin{empheq}[box=\keybox]{equation}
\tau_s\,\frac{\dd g}{\dd t} \;=\; -(g-1) + k_s\,z ,
\label{eq:sensitization}
\end{empheq}
trong đó $\tau_s$ là thời gian để độ nhạy trở về bình thường; nó dài hơn chính cơn đau, vì nhạy cảm hóa vẫn còn sau khi kích thích gây tổn thương đã ngừng~\cite{Woolf1983}, còn $k_s$ là cường độ nạp. Chừng nào mô còn tổn thương, đầu ra của cổng giữ $g$ trên một, và mọi thứ xảy ra quanh vết thương đều được truyền to hơn (đường nét đứt trên hình~\ref{fig:gatecurves}: khi $g=2$, ngưỡng hạ xuống $0.11$, còn đầu ra tăng gấp đôi). Đó là một cách chỉnh hệ báo động hợp lý: chừng nào vết thương chưa lành, thận trọng quá vẫn hơn.

Với kỹ sư, đây là phản hồi dương có rò chậm: đau làm tăng khuếch đại, khuếch đại làm tăng đau. Chừng nào vòng còn yếu, $g$ trở về bình thường khi vết thương lành. Nhưng nếu vòng mạnh, hệ báo động có thể kẹt ở trạng thái bật cả khi tổn thương không còn nữa, như nhóm có phản hồi mạnh trên hình~\ref{fig:bistable}. Mô hình~\eqref{eq:sensitization} là một sự đơn giản hóa; việc nhạy cảm hóa trung ương tồn tại thì đã được đo.

\section{Cơn đau như người thầy và như một ngắt}

Ngoài việc báo động, cơn đau trong cỗ máy còn hai nhiệm vụ.

\emph{Người thầy.} Cơn đau là phần thưởng âm: trong công thức sai số~\eqref{eq:ctd}, nó xuất hiện dưới dạng $r<0$. Mọi điều đã nói ở chương~\ref{sec:dofa} đều đúng ở đây: tín hiệu báo trước cơn đau bắt đầu gây ra sai số từ sớm, và mạng học cách tránh những gì dẫn tới đau. Thí nghiệm trên người cho thấy hoạt động não khi học từ cơn đau đi theo đúng sai số sai phân thời gian~\cite{Seymour2004}, và một phần nơ-ron \nt{DOPA} cũng đáp ứng với các sự kiện khó chịu (chương~\ref{sec:dofaalt}).

\emph{Ngắt.} Cơn đau kéo sự chú ý khỏi việc đang làm: đó là mục đích của nó, không phải tác dụng phụ~\cite{EcclestonCrombez1999}. Ở chương~\ref{sec:nora}, cùng vai trò ``ngắt'' ấy đã được bàn cho chùm xung dứt khoát của \nt{NORA}. Còn phản ứng nhanh nhất thậm chí không lên tới não: việc rụt tay khỏi vật nóng khép kín ngay trong tủy sống, bằng chính cặp cơ và chính phép cấm cơ đối vận như trên hình~\ref{fig:antagonists}.

\begin{keyblock}
\begin{proposition}[Tín hiệu báo động]
\label{prop:pain}
Cơn đau không phải phép đo, mà là tín hiệu báo động có mức ưu tiên cao. Nó thích nghi yếu hơn nhiều so với các kênh đo, đi theo hai dây (dây nhanh: ``ở đâu'', dây chậm: ``tệ đến mức nào''), đi qua một cổng mà cái chạm đóng lại (còn đau mạnh, theo giả định của thuyết cổng, mở ra), và âm lượng của nó do các kênh quảng bá đặt từ trên xuống. Các cơ chế này đã được đo, trừ việc cơn đau mở cổng; các mô hình đơn giản~\eqref{eq:gate} và~\eqref{eq:sensitization} là sự đơn giản hóa.
\end{proposition}
\end{keyblock}

\section{Tổng kết}

\begin{table}[t]
\centering
\footnotesize
\setlength{\tabcolsep}{4pt}
\renewcommand{\arraystretch}{1.15}
\begin{tabular}{>{\raggedright}p{3.0cm}>{\raggedright}p{4.9cm}>{\raggedright\arraybackslash}p{4.9cm}}
\toprule
Cơn đau làm gì & Bằng cách nào & Đã biết gì \\
\midrule
báo động & ngưỡng cao, thích nghi yếu hơn xúc giác & ngưỡng đã đo; so sánh thích nghi là ước tính \\
báo ``ở đâu'' và ``tệ đến mức nào'' & hai dây với tốc độ khác nhau~\eqref{eq:paintime} & đã đo \\
đi qua cổng & cấm qua \nt{GABA}/\nt{GLYC}~\eqref{eq:gate} & cổng đã đo; cơn đau mở cổng là giả định của thuyết; mô hình là đơn giản hóa \\
đổi âm lượng & \nt{SERO}, \nt{NORA}, opioid đi xuống & đã đo; quy tắc chọn âm lượng là giả thuyết \\
học sau tổn thương & tăng hệ số khuếch đại~\eqref{eq:sensitization} & nhạy cảm hóa đã đo; mô hình là đơn giản hóa \\
dạy cách tránh & phần thưởng âm trong~\eqref{eq:ctd} & sự giống sai số sai phân thời gian đã đo \\
ngắt công việc & chiếm sự chú ý; phản xạ rụt tay & đã đo \\
\bottomrule
\end{tabular}
\caption{Cơn đau như một tín hiệu báo động.}
\label{tab:pain}
\end{table}

Cơn đau được lắp từ những chi tiết ta đã gặp (bảng~\ref{tab:pain}): cảm biến \nt{GLUT}, hai dây, phép cấm qua nơ-ron trung gian ức chế, núm âm lượng quảng bá, mạch RC chậm, sai số dự đoán và ngắt. Điều mới ở nó chỉ có một: đó là đầu vào duy nhất của cỗ máy mà âm lượng do chính cỗ máy đặt, một hệ báo động mà chủ nhân có thể vừa làm dịu, vừa chỉnh lại.

\section{Bài tập}

\begin{exercise}[\lvlA]
\label{ex:14:1}
\emph{Hai dây.} Tín hiệu đi từ bàn chân, $L=1$~m. (a)~Một đợt bùng phát đi trên một bó dây cùng loại, tốc độ của chúng phủ kín dải~\eqref{eq:paintime}. Khi tới tủy sống, đợt bùng phát giãn ra bao nhiêu với loại nhanh và loại chậm? (b)~Hai đợt đau theo dây $15$ và $1$~m/s phân biệt được khi cách nhau từ $0.1$~s. Từ khoảng cách nào thì chúng hòa làm một?
\end{exercise}

\begin{exercise}[\lvlB]
\label{ex:14:2}
\emph{Khi cái chạm gây đau.} Cổng~\eqref{eq:gate} với tham số trong bài, hệ số khuếch đại $g$ tăng khi nhạy cảm hóa. Tới $g$ nào thì cái chạm không kèm đau không qua được cổng với mọi cường độ $\mu$? Với $g$ nào thì cái chạm $\mu=1$ bắt đầu đi qua?
\end{exercise}

\begin{exercise}[\lvlB]
\label{ex:14:3}
\emph{Độ ổn định của nhạy cảm hóa.} Đầu vào không đổi, đầu ra của cổng tuyến tính theo $g$: $z=gA-C$, $A=\nu+w_{z\mu}\mu$, $C=\beta q\ge0$. (a)~Tìm điểm cân bằng $g^*$ của mô hình~\eqref{eq:sensitization} và điều kiện ổn định của nó. (b)~Với $k_s=0.5$, tìm cường độ đau mà trên đó mô hình mất kiểm soát, khi không có chạm và khi có chạm $\mu=1$.
\end{exercise}

\begin{exercise}[\lvlA]
\label{ex:14:4}
\emph{Tín hiệu báo trước cơn đau.} Tín hiệu ở thời điểm $0$ báo trước cơn đau ở thời điểm $T=2$~s, trong~\eqref{eq:ctd} $r(t)=-P\,\delta(t-T)$, $\tau_\gamma=2$~s. (a)~Tìm diện tích của đỉnh vọt $\delta$ tại thời điểm có tín hiệu. (b)~Một hành động ở thời điểm $t_e=1$~s hủy cơn đau. Đỉnh vọt $\delta$ lúc đó bằng bao nhiêu và nó dạy mạng điều gì?
\end{exercise}

\begin{exercise}[\lvlB]
\label{ex:14:5}
\emph{Hệ báo động tự chốt.} Bổ sung cho cổng trong \texttt{models/\allowbreak pain\_gate.py} động học~\eqref{eq:sensitization} và bão hòa $z\le4$. Cái chạm $\mu=1$ luôn có, tổn thương $\nu=2$ kéo dài thời gian $T$, sau đó $\nu=0$. Bằng mô phỏng, tìm $T$ nhỏ nhất (theo đơn vị $\tau_s$) mà sau đó hệ báo động vẫn bật, với $k_s=4$, $4.6$, $5$, $6$, $8$.
\end{exercise}

\begin{exercise}[\lvlB]
\label{ex:14:6}
\emph{Thiết kế cổng.} Với $\beta=1.5$, $w_{z\mu}=0.3$, $g=1$, hãy chọn $q_0$, $w_{q\nu}$, $w_{q\mu}$ trong~\eqref{eq:gate} để ngưỡng đau là $0.2$ khi không chạm và $1.0$ khi có chạm $\mu=1$, còn cái chạm thôi làm dịu đau khi $\nu_\times=2.5$. Tới hệ số khuếch đại $g$ nào thì cái chạm không kèm đau không qua được cổng?
\end{exercise}

\begin{exercise}[\lvlC]
\label{ex:14:7}
\emph{Núm âm lượng.} Công việc mang lại giá trị $V$ trong một đơn vị thời gian, làm việc với tổn thương $\nu$ tốn $c\nu$, nên ngắt có lợi khi $\nu>V/c$; cơn đau ngắt khi $z>\theta=0.5$. (a)~Hệ số khuếch đại $g^*$ nào của cổng~\eqref{eq:gate} không có chạm cho ngưỡng này khi $V/c=0.25$, $0.5$, $2$? (b)~Cỗ máy có thể ước lượng $V$ và $c$ từ những tín hiệu nào trong sách?
\end{exercise}

\chapter{Cỗ máy học vận động như thế nào}
\chaptermark{Cỗ máy học vận động}
\label{sec:motorlearning}
\begin{chaptergoals}
\item vì sao cử động cần người thầy thứ ba, một dây sai số tới mỗi đầu ra, và cái giá của nó;
\item cách lớp mở rộng \nt{GLUT} khổng lồ và một dây sai số mỗi đầu ra cho quy tắc bình phương tối thiểu;
\item cách mạch đã học trở thành mô hình bên trong, một bộ lọc thích nghi dự đoán cơ thể;
\item vì sao một quá trình học nhanh và một chậm giải thích hậu hiệu ứng và sự quay lại tự phát.
\end{chaptergoals}

\section{Ba người thầy}

Ở chương~\ref{sec:muscles}, ta đã thấy cỗ máy điều khiển cơ ra sao, và dừng lại ở một giả thuyết: cử động nhanh và chính xác đòi hỏi một mô hình bên trong của cơ thể, dự đoán kết quả của lệnh~\cite{WolpertGhahramani2000}. Mô hình này từ đâu ra? Nó phải được học, và học theo cách khác với những gì ta đã học cho tới giờ.

Sách đã có hai người thầy. Người thứ nhất là \emph{không ai cả}: quy tắc Hebb (chương~\ref{sec:dofa}) tăng cường những gì hay xuất hiện cùng nhau và nén đầu vào. Người thứ hai là \emph{phần thưởng}: \nt{DOPA} phát cho mọi nơi một con số, ``tốt hơn hay tệ hơn kỳ vọng''. Vận động cần một người thầy thứ ba: \emph{sai số}. Bàn tay với trượt cái tách hai centimét về bên trái: đó không phải là ``tệ hơn'', mà là một vectơ cho mỗi đầu ra biết nó sai về hướng nào và bao nhiêu. Kỹ sư sẽ gọi đó là học có giám sát.

Khác biệt giữa người thầy thứ hai và thứ ba là khác biệt giữa một dây dùng chung cho tất cả và một dây riêng cho mỗi đầu ra. Trong mệnh đề~\ref{prop:broadcastcost}, học theo một con số chung chậm đi với mỗi nơ-ron có nhiễu ảnh hưởng tới kết quả. Còn nếu mỗi đầu ra $i$ nhận sai số riêng $e_i$ của nó, trọng số có thể thay đổi theo quy tắc đơn giản nhất:
\begin{empheq}[box=\keybox]{equation}
\frac{\dd w_{ij}}{\dd t} \;=\; -\eta\,e_i(t)\,x_j(t) .
\label{eq:lms}
\end{empheq}
Đây là quy tắc bình phương tối thiểu, đã quen thuộc từ lâu trong kỹ thuật lọc thích nghi~\cite{WidrowHoff1960}: trọng số thay đổi tỉ lệ với tích của đầu vào của nó và sai số của đầu ra của nó, và trung bình đi theo gradient của bình phương sai số. Không cần nhiễu để dò tìm như ở chương~\ref{sec:dofa}: hướng hiệu chỉnh đến thẳng qua dây sai số. Và tốc độ không giảm theo số đầu ra, vì sai số của đầu ra khác không tới khớp thần kinh này.

\begin{keyblock}
\begin{proposition}[Ba người thầy]
\label{prop:teachers}
Quy tắc Hebb dạy không cần thầy những gì thường gặp; tín hiệu \nt{DOPA} dạy bằng một con số cho cả mạng những gì có lợi, và chậm dần với mỗi nơ-ron; dây sai số tới mỗi đầu ra dạy theo quy tắc~\eqref{eq:lms} những gì chính xác, với tốc độ không phụ thuộc số đầu ra. Giá của cách thứ ba là một dây riêng cho mỗi đầu ra và một ai đó biết câu trả lời đúng.
\end{proposition}
\end{keyblock}

\section{Dây sai số}

Ai có thể biết câu trả lời đúng cho một cử động? Chính các cảm biến. Nếu tay lệch sang trái so với mục tiêu, mắt và các cảm biến độ giãn ở chương~\ref{sec:sensors} sẽ báo điều đó. Cần một sơ đồ đưa sai số này tới đúng các đầu ra.

Trong não có một vùng dường như được cấu tạo đúng cho việc này~\cite{Marr1969,Albus1971}. Sơ đồ của nó đều đặn khác thường, gần như tinh thể, và dễ nhận ra trong đó quy tắc~\eqref{eq:lms} (hình~\ref{fig:errorwire}).

\emph{Lớp mở rộng đầu vào.} Tín hiệu về lệnh và về trạng thái cơ thể đi tới một lớp khổng lồ gồm các nơ-ron \nt{GLUT} nhỏ. Có khoảng bảy mươi tỷ nơ-ron như vậy, nhiều hơn toàn bộ phần còn lại của não cộng lại~\cite{Azevedo2009}. Mỗi nơ-ron trộn vài đầu vào, và tín hiệu được trải ra thành một vectơ rất dài. Kỹ sư quen với mẹo này: nếu nâng số chiều của đầu vào, thì trong không gian mới hầu như mọi quan hệ cần thiết đều có thể thu được bằng một tổng có trọng số đơn giản~\cite{Cover1965}.

\emph{Nơ-ron đầu ra.} Tổng có trọng số được tính bởi khoảng ba mươi triệu nơ-ron đầu ra lớn~\cite{Andersen1992cereb}, mỗi nơ-ron có cỡ một trăm nghìn đầu vào từ lớp mở rộng. Và đó là các nơ-ron \nt{GABA}: kết quả học được truyền đi không bằng kích thích mà bằng ức chế, như trong phép cấm ở chương~\ref{sec:gaba}.

\emph{Dây sai số.} Mỗi nơ-ron đầu ra còn nhận thêm một đầu vào đặc biệt: đúng một dây \nt{GLUT}, phát xung thưa, khoảng mỗi giây một lần, nhưng mỗi lần đều gây ra trong nơ-ron đầu ra một đợt bùng mạnh~\cite{Ito2001}. Đó chính là $e_i$: xác suất xảy ra đợt bùng phụ thuộc vào hướng của sai số vận động~\cite{Herzfeld2018}. Sự trùng khớp giữa đợt bùng sai số và hoạt động của một đầu vào làm yếu đầu vào đó, đúng như số hạng $-\eta\,e_i x_j$ trong~\eqref{eq:lms}.

\begin{figure}[t]
\centering
\begin{tikzpicture}[>=Stealth,
  gran/.style={circle,draw,thick,fill=blue!6,minimum size=0.32cm,inner sep=0pt},
  outn/.style={circle,draw,very thick,fill=red!10,minimum size=1.1cm,inner sep=0pt,font=\scriptsize},
  inh/.style={-{Bar[width=7pt]},very thick,red!70!black}]
% inputs
\foreach \y/\lab in {1.6/{lệnh},0.4/{trạng thái cơ thể}}{
  \draw[->,thick,blue!60!black] (-0.6,\y) node[left,font=\scriptsize,black] {\lab} -- (0.35,\y);}
% expansion layer
\foreach \i in {0,...,11}{\node[gran] (g\i) at ({0.8+0.28*mod(\i,2)},{-0.3+0.23*\i}) {};}
\foreach \i in {0,...,11}{\pgfmathsetmacro{\yy}{mod(\i,3)>0 ? 1.6 : 0.4}\draw[gray!60!black] (0.35,\yy) -- (g\i);}
\node[font=\scriptsize,align=center] at (0.95,-1.05) {lớp mở rộng\\$\sim7\cdot10^{10}$ \nt{GLUT}};
% parallel wires to the output neuron
\foreach \i in {0,...,11}{\draw[->,gray!70!black] (g\i.east) -- (4.1,{1.0+0.07*(\i-5.5)});}
\node[outn] (P) at (4.6,1.0) {\nt{GABA}};
\node[font=\scriptsize,align=center,above=2pt] at (P.north) {nơ-ron đầu ra\\$\sim10^5$ đầu vào $x_j$, trọng số $w_{ij}$};
\draw[inh] (P) -- (6.8,1.0) node[right,font=\scriptsize,align=left] {hiệu chỉnh\\cử động};
% error wire
\draw[->,very thick,red!70!black,densely dashed] (4.6,-1.4) node[below,font=\scriptsize,black,align=center] {một dây sai số $e_i$\\\nt{GLUT}, $\sim1$ lần/s} -- (P.south);
\end{tikzpicture}
\caption{Sơ đồ học có giám sát, theo cách mà não dường như hiện thực. Lệnh và trạng thái cơ thể được một lớp khổng lồ các nơ-ron \nt{GLUT} trải thành một vectơ dài $x_j$; nơ-ron \nt{GABA} đầu ra cộng nó với các trọng số $w_{ij}$. Dây sai số duy nhất mang $e_i$ tới; sự trùng khớp giữa sai số và một đầu vào đang hoạt động làm yếu trọng số của đầu vào đó, như trong quy tắc~\eqref{eq:lms}.}
\label{fig:errorwire}
\end{figure}

Có một khó khăn ta đã gặp ở chương~\ref{sec:dofa}: sai số đến muộn hơn đầu vào đã gây ra nó. Cú trượt được thấy sau lệnh vài chục đến vài trăm mili giây. Vậy khớp thần kinh phải nhớ rằng nó đã hoạt động: cần một vết, như trong quy tắc~\eqref{eq:threefactor}. Và theo các thí nghiệm, độ dài của vết này được chỉnh khớp với độ trễ phản hồi trong từng nhiệm vụ: ở nơi sai số đến sau một trăm mili giây, đầu vào bị làm yếu mạnh nhất là đầu vào đã hoạt động một trăm mili giây trước đó~\cite{Suvrathan2016}.

\begin{keyblock}
\begin{proposition}[Dây sai số]
\label{prop:errorwire}
Trong sơ đồ ở hình~\ref{fig:errorwire}, mỗi nơ-ron đầu ra nhận đúng một dây sai số, và sự trùng khớp giữa sai số và một đầu vào làm yếu đầu vào đó. Đây là quy tắc bình phương tối thiểu~\eqref{eq:lms} áp dụng cho đầu vào đã được mở rộng tới số chiều khổng lồ. Cấu tạo của sơ đồ và sự làm yếu khi trùng khớp đã được đo; việc cả sơ đồ hoạt động như học có giám sát là giả thuyết hàng đầu.
\end{proposition}
\end{keyblock}

\section{Mô hình bên trong như một bộ lọc thích nghi}

Sơ đồ như vậy học được gì? Câu trả lời tự nhiên nhất là học dự đoán. Giả sử đầu vào là lệnh $u(t)$ đi qua một bộ các bộ lọc với những hằng số thời gian khác nhau, như trong lớp mở rộng: $x_j(t)=(g_j*u)(t)$. Đầu ra là tổng có trọng số của chúng, dự đoán những gì cảm biến sẽ báo:
\begin{empheq}[box=\keybox]{equation}
\hat y(t) \;=\; \sum_j w_j\,(g_j*u)(t), \qquad e(t) \;=\; \hat y(t) - y(t) ,
\label{eq:adaptivefilter}
\end{empheq}
còn trọng số học theo~\eqref{eq:lms}. Đây là \emph{bộ lọc thích nghi}: một sơ đồ tự chỉnh hàm truyền của mình để lặp lại một hệ chưa biết~\cite{Fujita1982}. Ở đây hệ chưa biết là chính cơ thể với khối lượng, độ đàn hồi và độ trễ của nó. Bộ lọc đã học chính là mô hình bên trong của chương~\ref{sec:muscles}: nó cho biết cảm biến sẽ báo gì mà không cần chờ chúng.

Mô hình như vậy có ích theo hai cách. Thứ nhất, dự đoán của nó có thể được đưa vào thay cho phản hồi đến trễ, và khi đó điều kiện ổn định~\eqref{eq:delaystable} không còn trói tay ta nữa: có thể sửa nhanh, theo mô hình. Thứ hai, hiệu giữa điều được dự đoán và điều nhận được là tín hiệu về cái \emph{bất ngờ}. Mọi thứ cỗ máy tự làm đều đã được mô hình dự đoán và trừ đi; phần còn lại là việc thế giới làm. Vì vậy, theo một giả thuyết, ta không thể tự cù mình~\cite{Blakemore1998}.

\section{Khi thế giới thay đổi: nhanh và chậm}

Hãy kiểm tra sơ đồ bằng một thí nghiệm dễ làm. Một người với tay tới mục tiêu, và thế giới bất ngờ thay đổi: chẳng hạn một lực ngang tác động lên tay. Những cử động đầu tiên bị trượt, sau đó độ trượt giảm dần. Nếu bỏ lực đi, tay ban đầu trượt về phía \emph{ngược lại}: mô hình đã học lực đó và tiếp tục bù cho nó. Đây là bằng chứng trực tiếp rằng cái được học là một mô hình, chứ không đơn giản là ``đã làm được''.

Thí nghiệm còn cho thấy một điều tinh tế hơn: có hai quá trình cùng học một lúc, một quá trình nhanh, học nhanh và quên nhanh, và một quá trình chậm, học chậm và nhớ lâu~\cite{Smith2006}. Các hiệu chỉnh được cập nhật sau mỗi cử động, theo sự kiện:
\begin{empheq}[box=\keybox]{equation}
e = p - (x_f + x_s), \qquad x_f \leftarrow A_f\,x_f + B_f\,e, \qquad x_s \leftarrow A_s\,x_s + B_s\,e ,
\label{eq:tworate}
\end{empheq}
trong đó $p$ là nhiễu loạn, $x_f$ và $x_s$ là hiệu chỉnh đã học của hai quá trình, $A$ là mức hiệu chỉnh được giữ lại tới cử động sau, $B$ là mức học từ sai số. Ở quá trình nhanh, $B$ lớn còn $A$ nhỏ hơn một khá rõ; ở quá trình chậm thì ngược lại.

\begin{figure}[t]
\centering
\begin{tikzpicture}[>=Stealth,x=0.034cm,y=1.9cm]
\draw[->,gray!70!black] (0,0) -- (355,0) node[right,font=\small,black] {cử động};
\draw[->,gray!70!black] (0,-1.15) -- (0,1.25);
\foreach \v in {-1,1}{\draw[gray!60!black] (-3,\v) -- (3,\v); \node[left,font=\scriptsize] at (0,\v) {$\v$};}
\foreach \n in {100,200,300}{\draw[gray!60!black] (\n,-0.04) -- (\n,0.04); \node[below,font=\scriptsize] at (\n,0) {\n};}
\draw[thin,gray!70!black] plot file {figures/tworate_p.dat};
\draw[thick,orange!85!black] plot file {figures/tworate_fast.dat};
\draw[thick,blue!60!black] plot file {figures/tworate_slow.dat};
\draw[very thick,black] plot file {figures/tworate_net.dat};
\node[font=\scriptsize,gray!50!black] at (120,1.12) {nhiễu loạn $p$};
\node[font=\scriptsize,black,anchor=west] at (238,0.52) {tổng: quay lại};
\node[font=\scriptsize,blue!60!black,anchor=west] at (150,0.39) {chậm $x_s$};
\node[font=\scriptsize,orange!85!black,anchor=west] at (60,0.08) {nhanh $x_f$};
\draw[densely dashed,gray!60!black] (226,-1.15) -- (226,1.25);
\node[font=\scriptsize,align=center,anchor=south west] at (228,-1.1) {hết\\sai số};
\end{tikzpicture}
\caption{Hai quá trình học theo mô hình~\eqref{eq:tworate}, $A_f=0.92$, $B_f=0.1$, $A_s=0.996$, $B_s=0.02$. Hai trăm cử động với nhiễu loạn $p=1$, sau đó sáu cử động với nhiễu loạn ngược, $p=-1$, cho tới khi tổng về không. Sau đường nét đứt, sai số không còn được đưa vào, và tổng tự quay về hiệu chỉnh cũ: quá trình nhanh đã quên nhiễu loạn ngược, quá trình chậm vẫn nhớ nhiễu loạn thuận.}
\label{fig:tworate}
\end{figure}

Mô hình~\eqref{eq:tworate} giải thích một thí nghiệm khó hiểu nếu không có nó (hình~\ref{fig:tworate}). Trong tính toán của chúng tôi, sau hai trăm cử động có nhiễu loạn, tổng hiệu chỉnh đạt $0.84$, trong đó phần lớn do quá trình chậm giữ, $0.63$. Tiếp đó sáu cử động với nhiễu loạn ngược đưa tổng về không: quá trình nhanh đã xuống âm, $-0.53$, quá trình chậm vẫn dương, $0.46$. Nếu bây giờ thôi báo sai số, quá trình nhanh quên hiệu chỉnh của nó sau vài chục cử động, quá trình chậm lâu hơn nhiều, và tổng \emph{tự} tăng lên $0.37$ sau bốn mươi cử động: kỹ năng cũ quay lại mà không cần luyện tập gì. Sự hồi phục tự phát như vậy đã được đo ở người; mô hình một quá trình không thể cho ra điều đó: không có sai số, nó chỉ đơn giản tắt dần về không.

Tại sao lại có hai quá trình? Một câu trả lời kỹ thuật hợp lý đến từ bộ lọc tối ưu ở chương~\ref{sec:acet}. Cơ thể thay đổi vì những lý do khác nhau với tốc độ khác nhau: cơ mỏi trong vài phút, lớn lên và yếu đi trong vài tháng. Nếu nhiễu có loại nhanh và loại chậm, ước lượng tốt nhất gồm hai bộ lọc với hằng số thời gian khác nhau, mỗi bộ bắt một loại~\cite{Kording2007}. Quá trình nhanh là hiệu chỉnh tạm thời, ``tay mỏi''; quá trình chậm là ``tay đã khác đi''. Đây vẫn là giả thuyết, nhưng nó giải thích vì sao một kỹ năng học trong một ngày bị mất một phần sang ngày hôm sau, và vì sao lần thứ hai học nhanh hơn.

\section{Tổng kết}

\begin{table}[t]
\centering
\footnotesize
\setlength{\tabcolsep}{4pt}
\renewcommand{\arraystretch}{1.15}
\begin{tabular}{>{\raggedright}p{3.0cm}>{\raggedright}p{4.6cm}>{\raggedright\arraybackslash}p{5.2cm}}
\toprule
Sơ đồ làm gì & Bằng cách nào & Đã biết gì \\
\midrule
học có giám sát & dây sai số tới mỗi đầu ra, quy tắc~\eqref{eq:lms} & cấu tạo sơ đồ và sự làm yếu khi trùng khớp đã đo; học có giám sát là giả thuyết hàng đầu \\
mở rộng đầu vào & bảy mươi tỷ nơ-ron \nt{GLUT} & số nơ-ron đã đo; vai trò của việc mở rộng là lý thuyết \\
tính đến độ trễ của sai số & vết được chỉnh theo độ trễ & đã đo \\
xây mô hình bên trong & bộ lọc thích nghi~\eqref{eq:adaptivefilter} & giả thuyết có cơ sở vững \\
học với hai nhịp độ & quá trình nhanh và chậm~\eqref{eq:tworate} & hậu hiệu ứng và sự quay lại tự phát đã đo; hai quá trình là mô hình \\
\bottomrule
\end{tabular}
\caption{Cỗ máy học vận động như thế nào.}
\label{tab:motorlearning}
\end{table}

Giờ cỗ máy có ba người thầy (bảng~\ref{tab:motorlearning}). Hebb nén những gì thường gặp; \nt{DOPA} dạy chọn điều có lợi bằng một con số; dây sai số dạy sự chính xác, và dạy nhanh, vì mỗi đầu ra có sai số riêng của nó. Não dường như dùng cả ba, và dùng mỗi cách ở nơi nó rẻ nhất: tín hiệu quảng bá cho ít quyết định, dây riêng cho việc tinh chỉnh cơ thể, nơi câu trả lời đúng do chính cảm biến báo.

\section{Bài tập}

\begin{exercise}[\lvlA]
\label{ex:15:1}
\emph{Dung lượng của sơ đồ có dây sai số.} Sơ đồ có $3\cdot10^7$ nơ-ron đầu ra, mỗi nơ-ron $10^5$ đầu vào và một dây sai số riêng ($1$~xung/s). Nơ-ron có $n$ đầu vào học được mọi cách gán nhãn cho $\approx2n$ vectơ~\cite{Cover1965}. (a)~Sơ đồ lưu được bao nhiêu liên kết? (b)~Theo~\eqref{eq:spikeentropy} với $\Delta t=10\ms$, ước lượng thời gian nhỏ nhất để lấp đầy dung lượng một nơ-ron.
\end{exercise}

\begin{exercise}[\lvlB]
\label{ex:15:2}
\emph{Tốc độ của quy tắc bình phương tối thiểu.} Các mode của quy tắc~\eqref{eq:lms} tắt dần với tốc độ $\eta\lambda_k(C)$. Với năm bộ lọc~\eqref{eq:adaptivefilter} trên nhiễu trắng, $C_{jk}=2\sqrt{\tau_j\tau_k}/(\tau_j+\tau_k)$. Tìm tỷ số giữa tốc độ nhanh nhất và chậm nhất nếu các $\tau_j$ tăng theo bội $2$ ($10$--$160\ms$) và theo bội $4$.
\end{exercise}

\begin{exercise}[\lvlB]
\label{ex:15:3}
\emph{Phân tích hai quá trình.} Viết mô hình~\eqref{eq:tworate} với tham số của hình~\ref{fig:tworate} dưới dạng $\mathbf x_{n+1}=M\mathbf x_n+\mathbf b\,p$. (a)~Từ trị riêng của $M$, tìm các hằng số thời gian tính theo số cử động. (b)~Tìm $x_f$, $x_s$ ở trạng thái dừng và tổng của chúng khi $p=1$ không đổi.
\end{exercise}

\begin{exercise}[\lvlB]
\label{ex:15:4}
\emph{Mô phỏng: lần thứ hai nhanh hơn.} Theo mô hình~\eqref{eq:tworate} với tham số trong \texttt{models/\allowbreak motor\_learning.py}, tìm số cử động cho tới khi $x_f+x_s\ge0.8$ khi $p=1$: (a)~với cỗ máy chưa học; (b)~sau $200$ cử động với $p=1$ và nhiễu loạn ngược cho tới khi tổng bằng không, như trên hình~\ref{fig:tworate}. Mô hình một quá trình có tăng tốc như vậy được không?
\end{exercise}

\begin{exercise}[\lvlB]
\label{ex:15:5}
\emph{Bộ điều khiển có mô hình bên trong.} Đối tượng $\dd x/\dd t=ku$ được quan sát với độ trễ $d$; bộ điều khiển $u=-G\hat x$ dự đoán $\hat x(t)=x(t-d)+\hat k\int_{t-d}^{t}u\,\dd s$. (a)~Với $\hat k=k=1$, $d=150\ms$, tìm $G$ cho hằng số thời gian $20\ms$ và so với giới hạn~\eqref{eq:delaystable}. (b)~Sơ đồ có ổn định khi $\hat k=0.4k$ không?
\end{exercise}

\begin{exercise}[\lvlB]
\label{ex:15:6}
\emph{Bộ lọc thích nghi học cơ.} Đối tượng là cú giật~\eqref{eq:twitch} có diện tích đơn vị, $T_c=50\ms$; các bộ lọc~\eqref{eq:adaptivefilter} là mạch RC với $\tau_j=12.5$, $25$, $50$, $100$, $200\ms$; đầu vào là một số chuẩn mới sau mỗi $10\ms$. Sau bao nhiêu giây, quy tắc~\eqref{eq:lms} với $\eta=50$ và $200$ giảm sai số xuống $0.5\%$? Vì sao sau $300$~s trọng số vẫn còn xa giá trị tốt nhất?
\end{exercise}

\begin{exercise}[\lvlC]
\label{ex:15:7}
\emph{Hai quá trình như một bộ lọc tối ưu.} Nhiễu loạn là tổng các quá trình $p^{f,s}_{n+1}=A_{f,s}\,p^{f,s}_n+w^{f,s}_n$ với phương sai nhiễu $q_f$, $q_s$, sai số được quan sát với nhiễu phương sai $R$; bộ lọc Kalman cho ra~\eqref{eq:tworate}. Với $q_f/R$, $q_s/R$ nào thì từ $A_f=0.92$, $A_s=0.996$ thu được $B_f=0.1$, $B_s=0.02$? $B_f$, $B_s$ thay đổi ra sao nếu tăng $q_s$ lên bốn lần?
\end{exercise}

\chapter{Đầu ra thứ hai: bộ não điều khiển hóa học của cơ thể như thế nào}
\chaptermark{Đầu ra hóa học}
\label{sec:chemoutput}

\section{Đầu ra nhanh và đầu ra chậm}

Ở chương~\ref{sec:muscles} ta đã thấy đầu ra nhanh của cỗ máy: cơ bắp. Nhưng cỗ máy còn một đầu ra thứ hai, chậm hơn. Bộ não giữ nhiệt độ cơ thể, nhịp tim, lượng nước và lượng đường trong giới hạn cần thiết, chuẩn bị cho cơ thể chạy hoặc ngủ. Để làm việc đó, nó không cử động các khớp mà phát ra tín hiệu hóa học. Có hai con đường. Thứ nhất là các dây thần kinh đi thẳng tới nội tạng: tim, mạch máu, ruột, các tuyến. Thứ hai là máu: một số nơ-ron và tuyến không thả chất vào khe tới tế bào bên cạnh mà thả thẳng vào dòng máu.

Máu là bus quảng bá rộng nhất trong tất cả những gì cuốn sách đã gặp. Các kênh quảng bá ở chương~\ref{sec:serocomputer}--\ref{sec:acet} gửi tín hiệu tới những vùng lớn của não; máu gửi nó tới mọi tế bào của cơ thể. Máu đi hết một vòng cơ thể trong khoảng một phút, nên thông điệp tới được mọi nơi sau vài chục giây và tồn tại hàng phút, hàng giờ cho đến khi chất bị phân hủy. Kiểu phát này hoàn toàn không có địa chỉ. Giống như trong mệnh đề~\ref{prop:receptor}, ý nghĩa của thông điệp do bên nhận chọn: chỉ những tế bào có thụ thể với chất đó mới đáp ứng.

\section{Ga và phanh: hai dây dẫn tới nội tạng}

Ví dụ rõ nhất là tim. Tim có bộ tạo nhịp riêng, giống như nơ-ron \nt{SERO} ở chương~\ref{sec:serocomputer}: nó tự đặt nhịp, và nếu cắt hết dây thần kinh thì tim đập khoảng một trăm lần mỗi phút. Bộ não không tạo ra nhịp mà chỉ tinh chỉnh nó, bằng hai dây dẫn trái dấu (hình~\ref{fig:gasbrake}). Một dây mang \nt{NORA}: đó là \emph{chân ga}, nó làm nhịp nhanh lên. Dây kia mang \nt{ACET}: đó là \emph{phanh}, nó làm nhịp chậm lại. Một cách thô,
\begin{empheq}[box=\keybox]{equation}
f \;\approx\; f_0 \;+\; a_S\,S \;-\; a_P\,P ,
\label{eq:gasbrake}
\end{empheq}
trong đó $f_0\approx100$ nhịp mỗi phút là nhịp riêng của bộ tạo nhịp, $S$ và $P$ là mức hoạt động của ga và phanh. Lúc nghỉ, phanh đang được đạp, và tim thường đập khoảng sáu mươi đến bảy mươi lần mỗi phút; khi gắng sức, phanh được nhả ra và ga được đạp xuống.

\begin{figure}[t]
\centering
\begin{tikzpicture}[>=Stealth,
  blk/.style={draw,very thick,rounded corners=2pt,align=center,font=\small},
  pace/.style={circle,draw,very thick,double,double distance=1.2pt,minimum size=1.8cm,inner sep=0pt,font=\scriptsize,fill=red!8,align=center},
  inh/.style={-{Bar[width=7pt]},thick,red!70!black}]
\node[blk,fill=blue!5,text width=1.6cm] (B) at (0,0) {lệnh\\của não};
\node[pace] (H) at (6.2,0) {bộ tạo\\nhịp};
\draw[->,very thick,orange!85!black] (B.north east) to[bend left=20] node[above,font=\scriptsize,black,align=center] {\nt{NORA}: ga, vài giây} (H.north west);
\draw[inh,very thick,blue!60!black] (B.south east) to[bend right=20] node[below,font=\scriptsize,black,align=center] {\nt{ACET}: phanh, dưới một giây} (H.south west);
\draw[->,very thick] (H) -- (8.4,0) node[right,font=\small] {tần số $f$};
\node[blk,fill=orange!10,text width=1.6cm] (A) at (0,-2.6) {\nt{ADRE}\\vào máu};
\draw[->,thick,orange!85!black] (B) -- (A);
\foreach \y/\lab in {-2.0/{tim},-2.6/{mạch máu},-3.2/{gan, cơ}}{
  \draw[->,thick,densely dashed,orange!85!black] (A.east) -- (4.3,\y) node[right,font=\scriptsize,black] {\lab};}
\end{tikzpicture}
\caption{Hai dây dẫn trái dấu tới bộ tạo nhịp tim: ga \nt{NORA} chậm, phanh \nt{ACET} nhanh. Phía dưới là cũng chân ga ấy nhưng đi qua bus quảng bá: các tế bào \nt{ADRE} thả adrenaline vào máu, và mọi cơ quan có thụ thể phù hợp đều nhận được tín hiệu (mũi tên nét đứt).}
\label{fig:gasbrake}
\end{figure}

Đây chính là mã hai dây của chương~\ref{sec:glutcomputer} và cặp cơ của chương~\ref{sec:muscles}, chỉ khác là giờ đây các dây không điều khiển khớp mà điều khiển nhịp độ của bộ tạo nhịp. Và hai dây có tốc độ khác nhau. Cả hai đều hoạt động như bộ lọc thông thấp, nhưng phanh truyền các thay đổi nhanh hơn ga vài lần~\cite{Berger1989}: đường đi của \nt{ACET} ngắn, protein gắn với thụ thể tự mở kênh kali, không cần chất truyền tin thứ hai bên trong tế bào~\cite{Logothetis1987gbg}, còn \nt{NORA} tác động qua một chuỗi phản ứng bên trong tế bào. Kỹ sư sẽ nhận ra ở đây thủ pháp dùng hai bộ truyền động có tốc độ khác nhau: bộ nhanh lo tinh chỉnh chính xác tức thời, bộ chậm lo những thay đổi lớn và lâu dài.

Cũng ở đây \nt{ADRE} tìm thấy vai trò của mình; bảng~\ref{tab:types} từng ghi rằng vai trò của nó trong não chưa rõ. Bên ngoài não thì vai trò ấy rõ ràng. Một phần tế bào của đường ga không gửi dây dẫn tới cơ quan mà thả adrenaline thẳng vào máu. Đó vẫn là chân ga, nhưng đi qua bus quảng bá: sau vài chục giây nó tới tim, mạch máu, gan và cơ cùng một lúc và tồn tại hàng phút. Chân ga có địa chỉ \nt{NORA} tinh chỉnh một cơ quan; chân ga quảng bá \nt{ADRE} chuyển cả cơ thể sang chế độ chạy.

\section{Chuỗi tầng có phản hồi}

Nhiều hormone không do một tuyến tiết ra mà do một chuỗi tầng. Một nhóm nhỏ nơ-ron ở đáy não thả vào máu tín hiệu thứ nhất; tín hiệu này khiến một tuyến trung gian tiết ra tín hiệu thứ hai; tín hiệu thứ hai lại khiến tuyến cuối tiết ra hormone thực sự tác động lên cơ thể. Còn hormone cuối quay ngược về các tầng đầu và ức chế chúng. Kết quả là một bộ khuếch đại ba tầng được bao bởi phản hồi âm, tức một bộ điều chỉnh mức, giống như hệ \nt{SERO} trong công thức~\eqref{eq:feedback}.

Ta mô tả mỗi tầng bằng một mạch RC với hằng số thời gian $\tau$ và hệ số khuếch đại $g$, còn phản hồi bằng hệ số $k$:
\begin{equation}
\tau\,\frac{\dd x_1}{\dd t} = -x_1 + g\bigl[u - k\,x_3\bigr]_+ ,\qquad
\tau\,\frac{\dd x_2}{\dd t} = -x_2 + g\,x_1 ,\qquad
\tau\,\frac{\dd x_3}{\dd t} = -x_3 + g\,x_2 ,
\label{eq:cascade}
\end{equation}
trong đó $u$ là lệnh của não, $x_3$ là mức hormone cuối. Hệ số khuếch đại vòng là $K=g^3k$. Ở trạng thái cân bằng $x_3=g^3u/(1+K)$, và như mọi bộ điều chỉnh có phản hồi, vòng lớn làm cho mức gần như không nhạy với sự tản mát hệ số khuếch đại của các tầng. Nhưng mỗi tầng lại làm lệch pha: ở tần số $\omega=\sqrt3/\tau$, ba tầng giống nhau cộng lại cho độ lệch nửa vòng và suy giảm tám lần. Từ đó có kết quả kinh điển của lý thuyết điều khiển:
\begin{empheq}[box=\keybox]{equation}
K \;<\; 8 ,
\label{eq:cascadestable}
\end{empheq}
nếu không, bộ điều chỉnh bắt đầu dao động, với chu kỳ khoảng $2\pi\tau/\sqrt3\approx3.6\,\tau$.

\begin{keyblock}
\begin{proposition}[Độ chính xác đổi lấy độ ổn định]
\label{prop:cascade}
Với chuỗi ba tầng có phản hồi tỉ lệ, sai số của mức bằng $1/(1+K)$, còn chuỗi chỉ ổn định khi $K<8$. Vì vậy một bộ điều chỉnh như thế không giữ mức chính xác hơn khoảng mười phần trăm; cố nâng hệ số khuếch đại sẽ biến nó thành bộ tạo dao động.
\end{proposition}
\end{keyblock}

Chúng tôi đã kiểm tra điều này trên mô hình~\eqref{eq:cascade} (hình~\ref{fig:cascade}). Với $K=4$, mức dao động nhẹ sau khi bật rồi ổn định. Với $K=7$, nó cũng ổn định, nhưng chậm. Với $K=12$, dao động không tắt dần mà cũng không tăng vô hạn: một tầng không thể xuất ra tín hiệu âm, nên chuỗi đi vào các xung đều đặn giữa $0.83$ và $1.24$ lần mức trung bình, với chu kỳ $3.7$--$3.9\,\tau$.

\begin{figure}[t]
\centering
\begin{tikzpicture}[>=Stealth,x=0.3cm,y=1.2cm]
\draw[->,gray!70!black] (0,0) -- (41.5,0) node[right,font=\small,black] {$t/\tau$};
\draw[->,gray!70!black] (0,0) -- (0,2.8) node[left,font=\small,black] {$x_3/x_3^*$};
\foreach \t in {0,10,20,30,40}{\draw[gray!60!black] (\t,-0.05) -- (\t,0.05); \node[below,font=\scriptsize] at (\t,0) {\t};}
\foreach \f in {1,2}{\draw[gray!60!black] (-0.3,\f) -- (0.3,\f); \node[left,font=\scriptsize] at (0,\f) {\f};}
\draw[densely dashed,gray!60!black] (0,1) -- (40,1);
\draw[thick,blue!60!black] plot file {figures/cascade_K4.dat};
\draw[thick,red!70!black] plot file {figures/cascade_K12.dat};
\node[font=\scriptsize,blue!60!black,anchor=west] at (16,0.55) {$K=4$: ổn định};
\node[font=\scriptsize,red!70!black,anchor=west] at (16,1.75) {$K=12$: tạo xung};
\end{tikzpicture}
\caption{Chuỗi ba tầng có phản hồi theo~\eqref{eq:cascade}; mức hormone cuối tính theo tỉ lệ với mức cân bằng $x_3^*$. Khi hệ số khuếch đại vòng dưới tám, mức ổn định; trên tám, chuỗi tự tạo ra các xung đều đặn.}
\label{fig:cascade}
\end{figure}

Các hormone của những chuỗi tầng thật quả thực được tiết theo xung: chẳng hạn mức hormone stress dao động với chu kỳ khoảng một giờ. Theo một giả thuyết, các xung này sinh ra từ chính phản hồi có trễ giữa các tầng, và không cần một bộ tạo nhịp riêng cho chúng~\cite{Walker2010}. Đó cũng là nguyên nhân dao động của vòng \nt{GLUT}--\nt{GABA} trong mệnh đề~\ref{prop:ei} và của vòng phản xạ căng cơ trong điều kiện~\eqref{eq:delaystable}: phản hồi âm bị trễ.

\section{Mã xung của hormone}

Xung không chỉ là tác dụng phụ mà còn có thể là mã. Các nơ-ron điều khiển sinh sản thả tín hiệu của chúng vào máu thành từng đợt ngắn, khoảng một giờ một lần. Nếu cấp cùng tín hiệu ấy cho tuyến nhận một cách liên tục thay vì từng đợt, tuyến không làm việc hết công suất như khi nhận xung mà lại im bặt; nếu lại cấp từng đợt mỗi giờ một lần, hoạt động được phục hồi~\cite{Belchetz1978}. Nghĩa là bên nhận không đọc mức mà đọc \emph{tần số} các đợt.

Với kỹ sư, ở đây không có gì bí ẩn. Một thụ thể mệt đi và thôi đáp ứng khi chất tác động lâu là một bộ lọc thông cao, giống như khớp thần kinh suy giảm~\eqref{eq:depression} ở chương~\ref{sec:code}. Tín hiệu liên tục bị nó dập tắt, còn các thay đổi thì được cho qua. Với bên nhận như vậy, chỉ có thể truyền thông tin bằng xung, và thông tin chuyển từ mức sang tần số và hình dạng các đợt. Ngoài ra, các tần số đợt khác nhau tác động khác nhau lên các loại tế bào khác nhau trong cùng một tuyến, nên qua một đường hóa học có thể truyền nhiều hơn một đại lượng, giống như trên một dây \nt{DOPA} ở chương~\ref{sec:dofaalt} có cả thành phần nhanh lẫn thành phần chậm.

\section{Nhịp ngày đêm: bộ tạo chế độ chậm}

Nhịp chậm nhất của cơ thể là nhịp ngày đêm. Nó được sinh ra bởi phản hồi âm có trễ bên trong tế bào: các gen tạo ra protein, và sau độ trễ vài giờ, các protein này ức chế chính sự sản xuất của mình~\cite{TakahashiJS2017}. Lại là phản hồi âm có trễ, lần thứ tư trong cuốn sách này, và lại là một nhịp, lần này với chu kỳ khoảng một ngày. Bộ tạo nhịp chính là một nhóm nhỏ gồm vài chục nghìn nơ-ron ở đáy não; nhịp của chúng đồng bộ các tế bào còn lại của cơ thể.

Chu kỳ riêng của bộ tạo nhịp này ở người không đúng một ngày mà trung bình là $24.18$ giờ~\cite{Czeisler1999}. Mỗi ngày nó được chỉnh lại bởi ánh sáng đến từ các cảm biến của mắt (chương~\ref{sec:sensors}). Với kỹ sư điện tử, đây là \emph{vòng khóa pha}: một bộ dao động có tần số riêng $\omega_0$ bám theo tín hiệu ngoài có tần số $\omega$, và độ lệch pha $\varphi$ tuân theo phương trình
\begin{empheq}[box=\keybox]{equation}
\frac{\dd\varphi}{\dd t} \;=\; (\omega-\omega_0) \;-\; K_\varphi\,\sin\varphi .
\label{eq:pll}
\end{empheq}
Khóa pha được giữ nếu $|\omega-\omega_0|<K_\varphi$. Bộ tạo nhịp có chu kỳ $24.18$ giờ cần dịch mỗi ngày khoảng mười một phút; ánh sáng buổi sáng thường đủ cho mức dịch đó. Trong đêm vùng cực hay trong hang tối, không có khóa pha, và nhịp trôi theo chu kỳ riêng.

Bộ tạo nhịp ngày đêm không phải là bộ tạo xung nhịp của máy tính. Nó không đếm các bước tính toán và không đồng bộ các xung của nơ-ron. Nó chuyển đổi \emph{chế độ}: ngủ và thức, mức hormone, nhiệt độ cơ thể, sự sẵn sàng ăn uống. Đó là một thanh ghi quảng bá chậm cùng loại với các thanh ghi ở chương~\ref{sec:serocomputer}--\ref{sec:acet}, chỉ khác là giá trị của nó thay đổi theo lịch trình được chỉnh theo mặt trời.

\begin{keyblock}
\begin{proposition}[Đầu ra hóa học]
\label{prop:chemoutput}
Đầu ra chậm của cỗ máy được lắp từ cùng những chi tiết như đầu ra nhanh. Các cơ quan nhận hai dây dẫn trái dấu và khác tốc độ, ga \nt{NORA} và phanh \nt{ACET}, còn toàn bộ cơ thể cùng lúc nhận các tín hiệu quảng bá qua máu, trong đó có adrenaline \nt{ADRE}. Mức hormone được giữ bởi các chuỗi tầng có phản hồi âm, với độ chính xác bị giới hạn bởi độ ổn định; một phần hormone được mã hóa bằng tần số các đợt; nhịp ngày đêm là một bộ tạo nhịp chậm có khóa pha theo ánh sáng. Cấu trúc các đường dẫn, các thí nghiệm với các đợt tiết và với chu kỳ đã được đo; việc giải thích xung của chuỗi tầng bằng chính phản hồi là giả thuyết.
\end{proposition}
\end{keyblock}

\section{Tóm tắt}

\begin{table}[t]
\centering
\footnotesize
\setlength{\tabcolsep}{4pt}
\renewcommand{\arraystretch}{1.15}
\begin{tabular}{>{\raggedright}p{3.1cm}>{\raggedright}p{4.8cm}>{\raggedright\arraybackslash}p{4.9cm}}
\toprule
Đầu ra hóa học làm gì & Bằng cách nào & Điều đã biết \\
\midrule
tinh chỉnh các cơ quan & ga \nt{NORA} và phanh \nt{ACET}~\eqref{eq:gasbrake} & đã đo; phanh nhanh hơn ga \\
chuyển cả cơ thể sang chế độ chạy & adrenaline \nt{ADRE} vào máu & đã đo \\
giữ mức hormone & chuỗi tầng có phản hồi, $K<8$~\eqref{eq:cascadestable} & chuỗi tầng và phản hồi đã đo; xung do chính vòng phản hồi là giả thuyết \\
truyền bằng tần số các đợt & thụ thể mệt đi như bộ lọc thông cao & sự phụ thuộc vào các đợt đã đo \\
chuyển chế độ theo ngày & bộ tạo nhịp có khóa pha theo ánh sáng~\eqref{eq:pll} & chu kỳ và việc chỉnh bằng ánh sáng đã đo \\
\bottomrule
\end{tabular}
\caption{Cỗ máy điều khiển hóa học của cơ thể như thế nào.}
\label{tab:chemoutput}
\end{table}

Cỗ máy có hai đầu ra (bảng~\ref{tab:chemoutput}). Đầu ra nhanh là cơ bắp: mili giây, có địa chỉ, tới từng khớp. Đầu ra chậm là hóa học của cơ thể: giây, phút và ngày, qua hai dây dẫn tới các cơ quan và qua máu tới tất cả cùng lúc. Chi tiết của cả hai đầu ra là như nhau: hai dây dẫn trái dấu, bộ tạo nhịp, phản hồi và độ trễ của nó; và cũng chính là những chi tiết đã dựng nên cỗ máy bên trong.

\section{Bài tập}

\begin{exercise}[\lvlA]
\label{ex:16:1}
\emph{Máu như một bộ lọc.} Hormone được đào thải khỏi máu với chu kỳ bán rã $t_{1/2}=10$~phút và được tiết thành các đợt ngắn mỗi $T=60$~phút. Nồng độ cực đại lớn hơn cực tiểu bao nhiêu lần, và hài cơ bản của các đợt bị suy giảm bao nhiêu lần? Với $t_{1/2}$ nào thì cực đại chỉ gấp đôi cực tiểu?
\end{exercise}

\begin{exercise}[\lvlA]
\label{ex:16:2}
\emph{Khóa pha ngày đêm.} Chu kỳ riêng của bộ tạo nhịp~\eqref{eq:pll} là $24.18$~giờ, ánh sáng dịch pha không quá một giờ mỗi ngày: $K_\varphi=2\pi/24$~rad/ngày. Tìm độ lệch pha xác lập $\varphi^*$ theo phút và thời gian hồi phục. Sau bao nhiêu ngày thì độ lệch sau khi tín hiệu ngoài nhảy $\pm6$~giờ sẽ nhỏ hơn một giờ?
\end{exercise}

\begin{exercise}[\lvlB]
\label{ex:16:3}
\emph{Độ chính xác đổi lấy độ ổn định.} Chuỗi~\eqref{eq:cascade} đã tuyến tính hóa được kéo dài thành $n$ tầng giống nhau. Tìm hệ số khuếch đại tới hạn $K_c(n)$ và sai số nhỏ nhất đạt được $1/(1+K_c)$ khi $n=3$ và $n=6$. Sai số này tiến tới đâu khi $n\to\infty$?
\end{exercise}

\begin{exercise}[\lvlB]
\label{ex:16:4}
\emph{Ga và phanh.} Trong~\eqref{eq:gasbrake}, phần đóng góp của ga và phanh là đầu ra của các mạch RC với $\tau_S=2$~s và $\tau_P=0.4$~s, lệnh không vượt quá $80$ và $60$ nhịp mỗi phút. Nhịp $f_0=100$; lúc nghỉ phanh $35$, ga $0$, $f=65$. Cần bao lâu để đạt $f=120$? Còn nếu không nhả phanh mà chỉ dùng ga thì sao?
\end{exercise}

\begin{exercise}[\lvlB]
\label{ex:16:5}
\emph{Mô phỏng: chuỗi tầng có trễ.} Thêm độ trễ $d$ vào phản hồi của hàm \texttt{cascade} trong \texttt{models/\allowbreak chem\_models.py} (hãy sao chép hàm, đừng chạy tệp): tầng đầu thấy $x_3(t-d)$. Tìm $K_c$ và chu kỳ tại biên khi $d/\tau=0$, $0.5$, $1$, $2$, rồi kiểm tra bằng mô phỏng tại $0.9K_c$ và $1.1K_c$.
\end{exercise}

\begin{exercise}[\lvlB]
\label{ex:16:6}
\emph{Bên nhận đọc các đợt.} Thụ thể mệt đi: $y=[c-a]_+$, $\tau_a\,\dd a/\dd t=c-a$, $\tau_a=30$~phút. Hormone đến thành các đợt dài $6$~phút với chu kỳ $T$, cùng liều trung bình $\bar c$. Tìm đáp ứng trung bình khi $T=15$, $60$, $120$~phút và $T\to\infty$. Đáp ứng bằng bao nhiêu khi nồng độ không đổi $\bar c$?
\end{exercise}

\begin{exercise}[\lvlC]
\label{ex:16:7}
\emph{Bộ tạo nhịp hay phản hồi?} Đề xuất một thí nghiệm phân biệt các xung sinh ra bởi phản hồi có trễ của chuỗi tầng~\eqref{eq:cascade} với một bộ tạo nhịp riêng. Có thể phân biệt hai giả thuyết qua sự dịch chu kỳ không, nếu $K$ chỉ thay đổi được trong phạm vi gấp rưỡi và chu kỳ đo được với độ chính xác $5\%$?
\end{exercise}

\chapter{Gỡ lỗi cỗ máy}
\chaptermark{Gỡ lỗi cỗ máy}
\label{sec:debugging}

\section{Gỡ lỗi một cỗ máy không có sơ đồ như thế nào}

Một kỹ sư được giao một bo mạch đang chạy mà không có sơ đồ sẽ gỡ lỗi nó bằng bốn thủ pháp. Anh ta đặt \emph{đầu dò} và xem cái gì đang chạy trên dây dẫn. Anh ta \emph{đưa vào tín hiệu thử} và xem cái gì thay đổi. Anh ta \emph{ngắt} một phần mạch và xem cái gì ngừng hoạt động. Và anh ta \emph{đọc các thanh ghi điều khiển} để biết bo mạch đang làm việc ở chế độ nào. Cả cuốn sách này là một nỗ lực đọc sơ đồ của bộ não, và trong chương này ta sẽ xem các kỹ sư đã biết dùng những thủ pháp nào trong bốn thủ pháp ấy, và các chương trước nói gì về những điều nên chờ đợi.

Ví dụ nổi bật nhất hiện nay của việc gỡ lỗi như vậy là các giao diện não--máy tính: những thiết bị đọc xung của nơ-ron và biến chúng thành lệnh cho các máy bên ngoài, hoặc ngược lại, đưa tín hiệu từ bên ngoài vào não. Phần lớn chương này dành cho chúng, và trước hết cho những gì công ty Neuralink đang làm.

\section{Đầu dò: điện cực nghe thấy gì}

Đầu dò của não là một điện cực mảnh cắm vào mô. Nó nghe thấy xung của các nơ-ron nằm trong bán kính cỡ một trăm micromét, thường là từ một đến vài nơ-ron. Một xung trên điện cực là một đỉnh nhọn vài chục đến vài trăm microvolt, dài khoảng một mili giây, và theo hình dạng của nó có thể phân biệt các nơ-ron lân cận với nhau. Nghe rõ nhất là các nơ-ron \nt{GLUT} lớn: chúng chiếm đa số trong vỏ não (bảng~\ref{tab:types}), và dòng điện của chúng mạnh hơn.

Khó khăn chính lộ ra ngay. Một đầu dò tốt ngày nay có hàng trăm hoặc hàng nghìn điện cực, trong khi não có $8.6\cdot10^{10}$ nơ-ron. Ta nghe lén chừng một phần trăm triệu của cỗ máy. Vậy tại sao từ một mẫu nhỏ đến thế lại có thể hiểu được điều gì? Câu trả lời đến từ chương~\ref{sec:code}. Các nơ-ron lân cận nhiễu một cách tương quan, và việc lấy trung bình theo nhóm chạm tới một giới hạn (mệnh đề~\ref{prop:pooling}): một trăm nơ-ron mang gần bằng một nghìn. Ngoài ra, bài toán mà vỏ não vận động giải có số chiều thấp: cánh tay có ít khớp (chương~\ref{sec:muscles}), và hoạt động của hàng trăm nơ-ron điều khiển nó nằm gọn trong một không gian chỉ gồm mười đến hai mươi biến chung~\cite{Gallego2017}. Một đầu dò trên một trăm nơ-ron nghe được gần như tất cả các biến ấy. Nếu bộ não lưu một thông điệp độc lập trong mỗi nơ-ron, việc nghe lén như vậy sẽ vô ích.

\section{Luồng dữ liệu: nén ngay trên đầu dò}

Đầu dò tạo ra một luồng khổng lồ. Để thấy hình dạng xung, mỗi điện cực phải được số hóa khoảng $20$ nghìn lần mỗi giây với độ chính xác khoảng mười bit. Một nghìn điện cực cho
\begin{empheq}[box=\keybox]{equation}
\begin{aligned}
R_{\text{thô}} \;&=\; N\,f_s\,b \;\approx\; 10^3\cdot2\cdot10^4\cdot10 \;\approx\; 2\cdot10^8\ \text{bit/s},\\
R_{\text{xung}} \;&\approx\; N\,r\,\log_2\frac{e}{r\,\Delta t} \;\approx\; 8\cdot10^4\ \text{bit/s}.
\end{aligned}
\label{eq:probedata}
\end{empheq}
Ước tính thứ hai là dung lượng của luồng xung~\eqref{eq:spikeentropy} với $r=10$ xung mỗi giây và độ chính xác $\Delta t=1\ms$: mọi thứ thực sự có trong luồng chỉ chiếm ít hơn hai nghìn rưỡi lần. Với thiết bị cấy ghép, đây là khác biệt quyết định: không thể truyền luồng thô bằng sóng vô tuyến qua da với mức công suất cho phép bên trong đầu mà không làm nóng mô. Vì vậy thiết bị cấy ghép phải tách xung ngay trên chip cạnh các điện cực và chỉ gửi ra ngoài các sự kiện: số kênh và thời điểm xung. Mô tả đầu tiên về hệ thống Neuralink chưa đạt tới mức đó: các chip khuếch đại và số hóa tín hiệu, toàn bộ luồng thô đi qua cáp, còn xung được một chương trình bên ngoài tách ra; nén trên chip và kết nối không dây được nêu ở đó như kế hoạch~\cite{Musk2019}.

Đây là một thủ pháp quen thuộc. Mắt nén tín hiệu một trăm lần trước khi gửi nó qua sợi cáp dài (chương~\ref{sec:sensors}); camera sự kiện không truyền khung hình mà truyền các thay đổi. Để vừa với đường truyền, đầu dò buộc phải làm đúng điều bộ não tự làm: nói bằng xung và chỉ nói về cái mới.

\section{Neuralink làm gì}

Thiết bị của Neuralink là một đầu dò được thiết kế theo những giới hạn này~\cite{Musk2019}. Thay cho mảng kim cứng là nhiều sợi mềm mảnh hơn sợi tóc với các điện cực dọc theo mỗi sợi: trong mô tả đầu tiên của hệ thống là tới $3072$ điện cực trên $96$ sợi, mỗi sợi $32$ điện cực, còn trong thiết bị đã cấy cho người, theo công bố của công ty, khoảng một nghìn kênh trên $64$ sợi. Sợi mềm ít làm tổn thương mô hơn và chịu tốt hơn việc não trong hộp sọ luôn xê dịch nhẹ. Các sợi được một robot cắm vào, tránh các mạch máu. Trong mô tả đầu tiên, như đã nói ở trên, luồng thô~\eqref{eq:probedata} đi ra ngoài qua cáp. Theo công bố của công ty, thiết bị cho người được làm khác: vỏ thiết bị nằm hoàn toàn dưới da, xung được tách bên trong, dữ liệu đi ra bằng sóng vô tuyến, năng lượng được cấp không dây.

Nhiệm vụ của thiết bị đầu tiên là đọc. Các sợi được đặt vào vùng vỏ não lập kế hoạch cử động bàn tay, và bộ giải mã biến xung thành chuyển động của con trỏ trên màn hình. Một người bị liệt tay điều khiển máy tính bằng cách hình dung cử động. Bản thân ý tưởng không mới: các giao diện dùng mảng cứng khoảng một trăm điện cực từ lâu đã cho phép điều khiển con trỏ~\cite{Hochberg2006}, viết chữ bằng cách hình dung cử động tay khi viết~\cite{Willett2021}, và thậm chí chuyển những cố gắng nói thành văn bản trên màn hình với tốc độ khoảng sáu mươi từ mỗi phút~\cite{Willett2023}. Cái mới của Neuralink là kỹ thuật: số kênh, không dây, vỏ kín và việc cắm bằng robot, tức là nỗ lực tạo ra một đầu dò có thể mang nhiều năm bên ngoài phòng thí nghiệm. Theo chỗ chúng tôi biết, kết quả thử nghiệm trên người chủ yếu do chính công ty công bố, chứ không phải trong các bài báo được bình duyệt; công ty có thông báo rằng một phần các sợi đã xê dịch sau khi cắm và số kênh hoạt động giảm đi. Với việc gỡ lỗi, đây là một chuyện phiền toái bình thường: đầu dò dịch chuyển so với bo mạch thì đã nghe những dây dẫn khác.

Hướng thứ hai của công ty là ghi: một thiết bị cho thị giác, đưa tín hiệu vào vỏ não thị giác của người đã mất thị lực. Về nó xem bên dưới, trong phần về ghi.

\section{Đọc: bộ giải mã là bên nhận}

Bộ giải mã của giao diện là bên nhận theo nghĩa của mệnh đề~\ref{prop:reader}: nó tự quyết định đọc phần nào của thông điệp. Hầu như mọi giao diện đều đọc \emph{tần số của nhóm}: đếm xung của từng kênh trong các cửa sổ vài chục mili giây và cộng chúng với các trọng số được chọn sao cho ra cử động mong muốn. Đây là mã tần số của nhóm ở mục~\ref{sec:rate}, và nó được chọn không phải ngẫu nhiên: khi mỗi cửa sổ chỉ có vài xung, tần số của một nơ-ron không xác định được, còn tổng theo một trăm nơ-ron thì đã dùng được.

Rút ra được bao nhiêu thông tin? Viết chữ ``bằng bàn tay tưởng tượng'' đạt tốc độ khoảng chín mươi ký tự mỗi phút~\cite{Willett2021}, tức một ký tự rưỡi mỗi giây: ngay cả với năm bit mỗi ký tự, như thế là dưới tám bit mỗi giây. Điều khiển con trỏ, theo công bố của Neuralink, cho khoảng mười bit mỗi giây. Trong khi đó giới hạn trên~\eqref{eq:spikeentropy} cho \emph{một} nơ-ron là vài chục bit mỗi giây, cho một trăm nơ-ron là hàng nghìn. Giao diện chỉ rút ra từ các nơ-ron bị nghe lén một phần nhỏ những gì chúng có thể mang. Nút thắt không phải là dây dẫn, mà là số biến độc lập mà nhóm mang (mệnh đề~\ref{prop:pooling}) và mức độ bộ giải mã hiểu mã, thứ mà như ta đã thấy ở chương~\ref{sec:code}, vẫn chưa được giải hết.

Còn một đặc điểm thứ hai, quen thuộc từ chương~\ref{sec:motorlearning}. Bộ giải mã và bộ não học lẫn nhau. Bộ não thấy con trỏ đi đâu, nhận sai số và chỉnh lệnh của mình, cũng chính là việc học vận động qua dây dẫn sai số. Còn các kỹ sư thỉnh thoảng lại chọn lại trọng số của bộ giải mã, vì các nơ-ron dưới các sợi thay đổi và các liên kết trong mạng học lại. Kết quả là hai người học điều chỉnh theo nhau; để cặp này không mất ổn định, nên tách tốc độ học của chúng ra: một bên học nhanh, bên kia học chậm.

\begin{keyblock}
\begin{proposition}[Vì sao đầu dò trên một nghìn nơ-ron hoạt động được]
\label{prop:probe}
Một giao diện nghe lén một phần cực nhỏ số nơ-ron vẫn có thể điều khiển cử động, vì hoạt động của nhóm có số chiều thấp và nhiễu của nó tương quan: vài trăm nơ-ron mang gần như mọi biến chung. Cũng vì lý do đó, giao diện chỉ rút ra được vài bit mỗi giây, bằng số biến độc lập trong nhiệm vụ, chứ không bằng dung lượng của luồng xung. Hoạt động của các giao diện đã được đo; bộ giải mã sẽ tốt hơn bao nhiêu khi mã được hiểu là một câu hỏi mở.
\end{proposition}
\end{keyblock}

\section{Ghi: khó hơn đọc}

Đưa tín hiệu vào cỗ máy khó hơn đọc. Một điện cực có dòng điện chạy qua kích thích không phải một nơ-ron, mà mọi nơ-ron và dây dẫn quanh nó, bất kể loại và dấu. Không thể dùng nó để ghi một mã mà trong đó điều quan trọng là nơ-ron \emph{nào} và \emph{khi nào} phát xung (chương~\ref{sec:code}): chỉ có thể đặt một cách thô \emph{ở đâu} và \emph{mạnh đến mức nào}.

Với thị giác, như thế là đủ phần nào. Dòng điện qua một điện cực trong vỏ não thị giác được người đó thấy như một chấm sáng ở một vị trí xác định trong thị trường; khi các chấm được bật cùng lúc, tới mười sáu chấm một lần, một người đã mất thị lực nhận ra được một số chữ cái và phân biệt được đường biên của đồ vật~\cite{Fernandez2021}. Đó là mã theo vị trí, và nó có thể ghi được. Nhưng hãy nhớ lại chương~\ref{sec:sensors}: mắt không gửi vào não các điểm ảnh, mà gửi một mô tả đã nén: các cạnh, các thay đổi, sáng lên và tối đi, trên những dây dẫn riêng. Ghi ``các chấm sáng'' vào vỏ não là ghi điểm ảnh vào một cỗ máy đang chờ ở đầu vào một mã hoàn toàn khác. Vì vậy thị giác nhân tạo hiện nay còn thô hơn mức có thể kỳ vọng theo số điện cực. Cũng có những tín hiệu thử không cần điện cực: ảo giác thị giác đưa qua mắt một đầu vào bất thường và đẩy cỗ máy ra khỏi chế độ bình thường, và mỗi lỗi chỉ ra một cơ chế riêng (mục~\ref{sec:illusions}).

\section{Điều đầu dò không thấy}

Điện cực nghe được bus địa chỉ: xung của các nơ-ron \nt{GLUT}, và kém hơn, của các nơ-ron \nt{GABA} nhỏ. Nhưng ở các chương~\ref{sec:serocomputer}--\ref{sec:acet} ta đã thấy rằng chế độ làm việc của cả cỗ máy do các thanh ghi quảng bá đặt: nồng độ \nt{SERO}, \nt{DOPA}, \nt{NORA}, \nt{ACET}. Điện cực không nghe được chúng: nơ-ron nguồn rất ít, và tín hiệu của chúng không phải là xung cạnh đầu dò mà là nồng độ trong một thể tích. Một người gỡ lỗi thấy bus dữ liệu nhưng không thấy các thanh ghi điều khiển sẽ luôn ngạc nhiên: cùng một đầu vào lại cho các đáp ứng khác nhau, vì ở đâu đó $g$, $a$ hay $\tau_\gamma$ đã thay đổi. Nồng độ có thể đo bằng cảm biến hóa học ngay trong mô~\cite{Bunin1998}, nhưng các giao diện hiện chưa dùng những cảm biến như vậy. Hoàn toàn nằm ngoài tầm nhìn của đầu dò còn có đầu ra thứ hai, chậm, của cỗ máy: các hormone trong máu (chương~\ref{sec:chemoutput}); chúng được đo trong mẫu máu chứ không bằng điện cực.

Đầu dò cũng không thấy các trọng số, mà chính trong chúng lưu gần như toàn bộ trí nhớ và mọi điều đã học. Chỉ có thể đoán chúng qua cách các đáp ứng thay đổi. Và nó không thấy cơn đau theo cách con người cảm thấy: ở chương~\ref{sec:pain} ta đã thấy rằng độ mạnh của tín hiệu báo động do chính cỗ máy đặt ra, bởi cổng ở tủy sống và các bộ điều chỉnh từ phía trên, nên không thể đọc từ cảm biến xem đau đến mức nào.

\section{Tóm tắt: gỡ lỗi và các câu hỏi mở}

\begin{table}[t]
\centering
\footnotesize
\setlength{\tabcolsep}{4pt}
\renewcommand{\arraystretch}{1.15}
\begin{tabular}{>{\raggedright}p{2.2cm}>{\raggedright}p{3.6cm}>{\raggedright}p{3.4cm}>{\raggedright\arraybackslash}p{4.0cm}}
\toprule
Thủ pháp gỡ lỗi & Giao diện làm gì & Cuốn sách giải thích gì & Điều đã biết \\
\midrule
đầu dò & nghe hàng trăm đến hàng nghìn nơ-ron & số chiều thấp và nhiễu tương quan (ch.~\ref{sec:code}) & đã đo \\
nén trên đầu dò & truyền sự kiện thay cho tín hiệu thô & dung lượng luồng xung~\eqref{eq:spikeentropy} & đã giải quyết về kỹ thuật \\
đọc & tần số nhóm $\to$ cử động, văn bản, lời nói & bộ giải mã là bên nhận, mã tần số của nhóm & hoạt động; vài bit mỗi giây \\
học cùng nhau & não và bộ giải mã điều chỉnh theo nhau & dây dẫn sai số (ch.~\ref{sec:motorlearning}) & quan sát được; quy tắc điều chỉnh đang được nghiên cứu \\
ghi & dòng điện bật các chấm trong thị trường & mã vị trí ghi được, mã thời gian thì không (ch.~\ref{sec:sensors}) & các chấm đã đo; thị giác hữu ích thì chưa \\
đọc thanh ghi & hầu như chưa làm & các kênh quảng bá (ch.~\ref{sec:serocomputer}--\ref{sec:acet}) & cảm biến hóa học có trong thí nghiệm \\
\bottomrule
\end{tabular}
\caption{Gỡ lỗi cỗ máy: giao diện não--máy tính làm được gì và các chương của sách nói gì về chúng.}
\label{tab:debugging}
\end{table}

Bảng~\ref{tab:debugging} tóm tắt chương. Giao diện não--máy tính đã biết đặt đầu dò lên bus địa chỉ và đọc từ đó vài bit mỗi giây, và việc điều này hoạt động được là nhờ số chiều thấp và nhiễu tương quan ở chương~\ref{sec:code}. Chúng mới chỉ biết ghi vào cỗ máy một cách thô, vì mã đầu vào của nó đã được nén và gửi tới từng dây dẫn riêng. Còn các thanh ghi điều khiển và trọng số, những thứ làm cho cỗ máy học được và điều chỉnh được, thì đầu dò hầu như chưa thấy.

Mỗi câu hỏi mở ở chương~\ref{sec:synthesis} cũng là một bài toán cho người gỡ lỗi. Nên đọc mã nào sẽ được giải quyết khi đầu dò dày hơn và bộ giải mã biết dùng thời điểm xung chứ không chỉ tần số. Mạng nhiều lớp học như thế nào có thể kiểm tra bằng cách đặt đầu dò đồng thời vào nhiều lớp và xem đáp ứng của chúng thay đổi ra sao khi học. Các kênh quảng bá truyền gì sẽ lộ ra khi bên cạnh đầu dò điện có thêm đầu dò hóa học. Cuốn sách này là một nỗ lực đọc sơ đồ của cỗ máy qua hành vi của nó và qua những định luật mà kỹ sư nào cũng biết. Những công cụ gỡ lỗi đang được xây dựng hiện nay sẽ cho phép kiểm tra sơ đồ ấy bằng đầu dò.

\section{Bài tập}

\begin{exercise}[\lvlA]
\label{ex:17:1}
\emph{Ngân sách kênh vô tuyến.} Truyền qua da tốn $1$~nJ mỗi bit, bộ phát được tiêu thụ không quá $10$~mW. Cần công suất bao nhiêu cho luồng thô~\eqref{eq:probedata} với $N=1000$ điện cực và cho luồng xung? Luồng thô sẽ đòi hỏi năng lượng mỗi bit là bao nhiêu?
\end{exercise}

\begin{exercise}[\lvlA]
\label{ex:17:2}
\emph{Một trăm nơ-ron chứa bao nhiêu bit.} Bộ giải mã cộng xung của $N$ nơ-ron trong các cửa sổ $50\ms$ với $r=20$~xung/s, tương quan nhiễu từng cặp $c=0.1$, tín hiệu làm tần số nhóm thay đổi $30\%$ (trung bình bình phương). Ước tính tốc độ $\tfrac12\log_2(1+\text{SNR})$ mỗi cửa sổ khi $N=10$, $100$, $1000$ và $N\to\infty$. Còn khi $c=0$, $N=1000$?
\end{exercise}

\begin{exercise}[\lvlB]
\label{ex:17:3}
\emph{Tốc độ của giao diện bàn phím.} Giao diện chọn một trong $N=32$ ký tự, một lần rưỡi mỗi giây: đúng ký tự dự định với xác suất $P$, các lỗi đồng xác suất. Nó truyền bao nhiêu bit mỗi giây khi $P=0.99$, $0.94$, $0.8$? Khi $P=0.94$ cần bao nhiêu ký tự mỗi giây để đạt $10$~bit/s?
\end{exercise}

\begin{exercise}[\lvlB]
\label{ex:17:4}
\emph{Mô phỏng: bộ giải mã nhóm.} $N$ nơ-ron với $\lambda_i=[b+m\,\mathbf v\cdot\mathbf p_i]_+$, $b=20$, $m=15$~xung/s, cửa sổ $50\ms$, $\mathbf v$ có độ lệch chuẩn $0.5$ mỗi thành phần; tất cả cùng thấy nhiễu chung, $\mathbf v+\boldsymbol\xi$, $\sigma_\xi=0.2$. Mô phỏng bộ giải mã vectơ nhóm không chệch. Với $N$ nào hai loại nhiễu bằng nhau, và $N\to\infty$ cho bao nhiêu bit mỗi thành phần mỗi cửa sổ?
\end{exercise}

\begin{exercise}[\lvlB]
\label{ex:17:5}
\emph{Hai người học.} Não xuất lệnh $m=b\,v^*$, bộ giải mã xuất $v=d\,m$, $\langle v^{*2}\rangle=1$. Sau mỗi lần thử, cả hai đi một bước gradient theo $\tfrac12(v-v^*)^2$ với tốc độ $\eta$. Mô phỏng từ $b=1.2$, $d=1$ khi $\eta=0.8$, $1.1$, $1.5$, $2$. Riêng mỗi bên ổn định khi $\eta<2$; cặp ổn định khi $\eta$ bằng bao nhiêu?
\end{exercise}

\begin{exercise}[\lvlB]
\label{ex:17:6}
\emph{Thiết kế: đường ống của đầu dò.} $1024$ kênh, mỗi kênh $20$~nghìn mẫu mỗi giây, nơ-ron $10$~xung/s, vô tuyến $1$~nJ/bit. Ngưỡng nào theo nhiễu trắng của mẫu cho không quá $0.1$~Hz xung giả mỗi kênh? Truyền số xung trong mỗi $20\ms$ bằng $2$~bit: công suất là bao nhiêu và nhỏ hơn bao nhiêu lần so với mẫu thô $10$ bit?
\end{exercise}

\begin{exercise}[\lvlC]
\label{ex:17:7}
\emph{Giới hạn tốc độ của giao diện.} Ước tính $R\le DW\log_2(1+\text{SNR})$ cho $D=10$ biến với băng thông $W=5$~Hz khi $\text{SNR}\approx0.9$ (nhiễu giống tín hiệu) và $90$ (nhiễu độc lập của $1000$ nơ-ron). Đo được khoảng $10$~bit/s. Thí nghiệm nào phân biệt được hai cách giải thích khoảng cách này: nhiễu giống tín hiệu hay chưa hiểu mã?
\end{exercise}

\chapter{Các loại nơ-ron theo chức năng điện}
\chaptermark{Nơ-ron như linh kiện}
\label{sec:eetypes}
\begin{chaptergoals}
\item nơ-ron làm những linh kiện nào: bộ tích phân, bộ lọc, bộ chuyển đổi, bộ dao động, mạch chốt;
\item cách một tế bào trở thành bộ cộng hưởng thông dải, hay mạch chốt do \nt{SERO} bật lên;
\item cách hồi tiếp dựng bộ tích phân không rò từ nơ-ron rò, và vì sao nó cần tinh chỉnh chính xác;
\item vì sao đây là chế độ của một tế bào chứ không phải loại tế bào, do kênh ion và kênh quảng bá đặt.
\end{chaptergoals}

Ở chương~\ref{sec:neurontypes} ta đã phân loại nơ-ron theo chất dẫn truyền thần kinh mà đầu ra của chúng giải phóng. Một kỹ sư điện tử sẽ phân loại khác: theo việc nơ-ron làm gì với tín hiệu của nó, tức nó hoạt động như linh kiện nào. Linh kiện chính của cuốn sách ta đã biết từ chương~\ref{sec:glutcomputer}: nơ-ron tích phân rò~\eqref{eq:glutneuron}, một mạch RC với bộ so sánh. Nhưng trong não còn có những linh kiện khác. Bảng~\ref{tab:eetypes} gom chúng thành bốn nhóm, và hầu như linh kiện nào cũng đã xuất hiện trong sách.

\begin{center}
\footnotesize
\setlength{\tabcolsep}{3pt}
\renewcommand{\arraystretch}{1.15}
\begin{longtable}{>{\raggedright}p{3.1cm}>{\raggedright}p{3.3cm}>{\raggedright}p{4.7cm}>{\raggedright\arraybackslash}p{2.4cm}}
\caption{Nơ-ron như linh kiện: tên gọi, linh kiện tương ứng trong điện tử, việc linh kiện làm và nơi nó xuất hiện trong sách.}\label{tab:eetypes}\\
\toprule
Nơ-ron & Linh kiện & Làm gì & Ở đâu trong sách \\
\midrule
\endfirsthead
\toprule
Nơ-ron & Linh kiện & Làm gì & Ở đâu trong sách \\
\midrule
\endhead
\bottomrule
\endfoot
\multicolumn{4}{l}{\emph{Bộ lọc và bộ tích phân}} \\
nơ-ron tích phân rò & bộ tích phân RC với bộ so sánh & cộng các đầu vào trong cửa sổ $\tau$ và phát xung khi chạm ngưỡng & ch.~\ref{sec:glutcomputer}, \eqref{eq:glutneuron} \\
bộ phát hiện trùng khớp & cổng AND có cửa sổ thời gian & chỉ phát xung nếu các đầu vào đến gần như cùng lúc; $\tau$ rất nhỏ & ch.~\ref{sec:glutcomputer}, \ref{sec:code}, \eqref{eq:window} \\
nơ-ron pha (phasic) & bộ lọc thông cao, bộ phát hiện sườn & đáp ứng với thay đổi của đầu vào rồi im lặng & ch.~\ref{sec:sensors} \\
nơ-ron thích nghi & bộ lọc thông cao có tự động điều chỉnh hệ số khuếch đại & tần số giảm khi đầu vào không đổi & ch.~\ref{sec:sensors}, \eqref{eq:nakarushton} \\
bộ cộng hưởng & mạch dao động, bộ lọc thông dải & đáp ứng tốt nhất với đầu vào có tần số của riêng nó & mục~\ref{sec:resonator} \\
bộ tích phân không rò & bộ tích phân dùng khuếch đại thuật toán & biến vận tốc thành vị trí và giữ vị trí đó & mục~\ref{sec:perfectint} \\
\midrule
\multicolumn{4}{l}{\emph{Bộ chuyển đổi}} \\
nơ-ron trương lực (tonic) & bộ chuyển đổi điện áp--tần số & tần số bám tuyến tính theo đầu vào và hầu như không mệt & ch.~\ref{sec:code} \\
nơ-ron phân cấp (không phát xung) & bộ khuếch đại tương tự, bộ đệm & truyền điện áp biến thiên liên tục đi một quãng ngắn & ch.~\ref{sec:sensors} \\
nơ-ron chỉnh lưu & bộ chỉnh lưu nửa chu kỳ & tần số không bao giờ âm, $[u]_+$ & ch.~\ref{sec:glutcomputer}, \ref{sec:gaba} \\
cặp ``bật--tắt'' & cặp bộ chỉnh lưu, mã hai dây & một bên đáp ứng với ``nhiều hơn'', bên kia với ``ít hơn'' & ch.~\ref{sec:glutcomputer}, \eqref{eq:dualrail} \\
nơ-ron có ức chế sun & bộ chia tương tự, điều chỉnh hệ số khuếch đại & chia đầu vào chứ không trừ khỏi nó & ch.~\ref{sec:gaba}, \eqref{eq:shunt} \\
\midrule
\multicolumn{4}{l}{\emph{Bộ tạo dao động và thời gian}} \\
bộ tạo nhịp & bộ dao động tích thoát, mạch đa hài & tự phát xung với tần số không đổi & ch.~\ref{sec:serocomputer}, \ref{sec:dofa} \\
bộ tạo nhịp có đầu vào & bộ dao động điều khiển bằng điện áp & có nhịp riêng, đầu vào làm nó nhanh lên hoặc chậm lại & ch.~\ref{sec:chemoutput} \\
nơ-ron phát chùm xung & bộ tạo chùm xung có cổng & phát các chùm xung dày, tức các ``gạch'' & ch.~\ref{sec:code}, \eqref{eq:burst} \\
nơ-ron gắn với pha & bộ tách sóng pha, vòng khóa pha & phát xung tại một pha xác định của sóng đầu vào & ch.~\ref{sec:sensors}, \ref{sec:chemoutput} \\
sợi trục có trễ & đường dây trễ & làm trễ tín hiệu một khoảng thời gian định trước & ch.~\ref{sec:glutcomputer}, hình~\ref{fig:itd} \\
\midrule
\multicolumn{4}{l}{\emph{Bộ nhớ và chuyển mạch}} \\
nơ-ron hai trạng thái bền & trigger Schmitt, mạch chốt & một đầu vào ngắn đưa nó vào một trạng thái kéo dài & mục~\ref{sec:bistable} \\
nơ-ron có các nhánh đuôi gai & bộ dồn kênh, một mạng nhiều lớp nhỏ & nhánh chỉ cho đầu vào đi qua nếu có đầu vào cho phép & ch.~\ref{sec:synthesis} \\
\end{longtable}
\end{center}

Một lưu ý quan trọng hơn cả bảng: đây không phải là các tế bào khác nhau, mà là các \emph{chế độ} khác nhau. Cùng một nơ-ron có thể là bộ tích phân khi đầu vào chậm và là bộ phát hiện trùng khớp khi sự ức chế đã rút ngắn $\tau$ của nó (chương~\ref{sec:code}); là bộ tạo nhịp khi thức và là bộ thu im lặng khi ngủ. Linh kiện nào được tạo thành là do các kênh ion của màng quyết định, cùng với các kênh quảng bá điều chỉnh chúng.

\section{Bốn linh kiện chưa từng có trong sách}

\subsection{Bộ cộng hưởng}
\label{sec:resonator}

Một số nơ-ron có một dòng điện chậm chống lại mọi thay đổi điện áp: điện áp tăng thì dòng kéo nó xuống, điện áp giảm thì kéo lên, nhưng có trễ. Với kỹ sư điện tử, dòng như vậy hoạt động như một cuộn cảm, và cùng với điện dung của màng, chúng tạo thành một mạch dao động. Nơ-ron đáp ứng mạnh nhất với đầu vào có một tần số xác định, thường từ vài hertz đến một vài chục hertz, và yếu hơn với đầu vào chậm hơn hay nhanh hơn~\cite{HutcheonYarom2000}. Đó là một bộ lọc thông dải trong một tế bào: từ đầu vào nhiễu, nó chọn ra nhịp có tần số của riêng nó, chẳng hạn nhịp cục bộ của chương~\ref{sec:gaba}.

\subsection{Bộ tích phân không rò}
\label{sec:perfectint}

Bộ tích phân rò chỉ nhớ đầu vào trong thời gian $\tau$, tức vài chục mili giây. Nhưng để đứng yên, mắt cần nhớ vị trí của nó hàng giây: lệnh ``xoay'' đến dưới dạng một vận tốc ngắn ngủi, còn mắt thì phải được giữ ở vị trí mới. Mạng giải quyết việc này bằng chính phản hồi đã tạo nên bộ nhớ ở chương~\ref{sec:glutcomputer}. Nếu một nhóm nơ-ron có hằng số thời gian $\tau$ tự kích thích chính nó với trọng số $w$, thì $\tau\,\dd r/\dd t=-r+w\,r+I$, và hằng số thời gian của nhóm bằng
\begin{empheq}[box=\keybox]{equation}
\tau_{\text{hd}} \;=\; \frac{\tau}{1-w} .
\label{eq:perfectint}
\end{empheq}
Với $\tau=0.1\,\text{s}$ và $w=0.995$ ta được $20\,\text{s}$: một bộ tích phân gần như lý tưởng làm từ các nơ-ron mà mỗi nơ-ron tự nó quên sau một phần mười giây~\cite{Seung1996}. Cái giá là phải tinh chỉnh chính xác: với $w$ lớn hơn một chút so với một, bộ tích phân đi vào bão hòa, với $w$ nhỏ hơn một chút thì nó quên nhanh. Vì vậy một bộ tích phân như thế cần được học liên tục để tự điều chỉnh, giống như cách được mô tả ở chương~\ref{sec:motorlearning}.

\subsection{Nơ-ron hai trạng thái bền}
\label{sec:bistable}

Ở các chương~\ref{sec:glutcomputer} và~\ref{sec:gaba}, flip-flop được lắp từ một nhóm nơ-ron. Nhưng cũng có flip-flop trong một tế bào. Một số nơ-ron có một dòng điện mà khi đã bật lên thì tự duy trì sự kích thích: một đầu vào ngắn đưa nơ-ron vào trạng thái hoạt động kéo dài, gọi là \emph{bình nguyên}, còn sự ức chế đưa nó trở về trạng thái nghỉ. Đó là một trigger Schmitt: nơ-ron có hai trạng thái bền và vòng trễ giữa chúng. Các nơ-ron \nt{ACET} điều khiển cơ chẳng hạn hoạt động như vậy, và tính chất này được bật bởi một kênh quảng bá: bình nguyên xuất hiện khi serotonin tới được nơ-ron~\cite{Hounsgaard1988}. Cùng một tế bào, có \nt{SERO} thì là mạch chốt, không có thì là bộ tích phân thông thường.

\subsection{Bộ tích phân và bộ cộng hưởng}

Lý thuyết chia các tế bào kích thích được thành hai lớp~\cite{Izhikevich2007}. \emph{Bộ tích phân} giống như một bộ dao động điều khiển bằng điện áp bắt đầu từ không: ngay trên ngưỡng một chút, nơ-ron có thể phát xung chậm tùy ý, và tần số tăng dần đều theo đầu vào. Một nơ-ron như vậy cộng mọi thứ đến với nó và truyền tốt độ lớn của đầu vào bằng tần số. \emph{Bộ cộng hưởng} giống như một bộ dao động có tần số tối thiểu: chúng không biết phát chậm, khi đầu vào ở ngưỡng thì phát ngay xung với tần số hữu hạn, và ưa đầu vào có tần số của riêng chúng. Loại thứ nhất là linh kiện tự nhiên của mã tần số ở chương~\ref{sec:code}, loại thứ hai là của mã thời gian.

\begin{keyblock}
\begin{proposition}[Một nơ-ron, nhiều linh kiện]
\label{prop:eetypes}
Theo chức năng điện, nơ-ron hoạt động như bộ tích phân rò và không rò, bộ phát hiện trùng khớp, bộ lọc thông cao và thông dải, bộ chuyển đổi điện áp--tần số, bộ chỉnh lưu, bộ chia, bộ dao động, đường dây trễ và flip-flop. Một tế bào cho trước trở thành linh kiện nào là do các kênh ion của nó quyết định, cùng với các kênh quảng bá điều chỉnh chúng. Bản thân các chế độ đã được đo; mức độ bộ não dùng sự tái cấu hình này để tính toán phần lớn vẫn là giả thuyết.
\end{proposition}
\end{keyblock}

Với kỹ sư, điều này giống một mảng logic lập trình được: phần cứng chỉ có một, còn linh kiện do cấu hình quyết định. Có điều ở đây cấu hình không được thay đổi lúc nạp chương trình, mà ngay trong khi chạy, bởi các kênh quảng bá chậm đã nói đến ở các chương~\ref{sec:serocomputer}--\ref{sec:acet}.

\section{Bài tập}

\begin{exercise}[\lvlA]
\label{ex:18:1}
\emph{Dung sai của bộ tích phân không rò.} Các nơ-ron với $\tau=0.1$~s giữ vị trí mắt theo~\eqref{eq:perfectint} trong $T=1$~s với sai số không quá $5\%$ về mỗi phía. (a)~Tìm khoảng cho phép của $w$. (b)~$w$ là tổng trọng số của $K$ khớp thần kinh, mỗi khớp có sai số độc lập $10\%$. Với $K$ nào độ tản mát của tổng nằm trong dung sai?
\end{exercise}

\begin{exercise}[\lvlB]
\label{ex:18:2}
\emph{Bộ cộng hưởng như một mạch dao động.} Mô tả dòng chậm ở mục~\ref{sec:resonator} bằng biến $W$: $C\,\dd V/\dd t=-g_LV-g_wW+I$, $\tau_w\,\dd W/\dd t=V-W$. Chứng minh đây là mạch $RC$ có nhánh song song $R_w=1/g_w$, $L_w=\tau_w/g_w$, và với $C/g_L=10\ms$, $\tau_w=100\ms$, $\alpha=g_w/g_L=4$, tìm tần số cộng hưởng và độ lợi của $|Z|$ tại đó so với $|Z(0)|$.
\end{exercise}

\begin{exercise}[\lvlB]
\label{ex:18:3}
\emph{Bộ tích phân bắt đầu hoạt động như thế nào.} (a)~Nơ-ron $\tau\,\dd V/\dd t=-V+\mu$, $\tau=10\ms$, với ngưỡng $\theta$ và đặt lại về không. Với $\mu/\theta-1$ bằng bao nhiêu thì tần số là $1$~Hz? (b)~Mô hình $\tau\,\dd v/\dd t=v^2+\mu$ (xung là khi $v$ chạy tới $+\infty$, đặt lại về $-\infty$) cho tần số bao nhiêu khi $\mu=0.01$?
\end{exercise}

\begin{exercise}[\lvlB]
\label{ex:18:4}
\emph{Mô phỏng: bộ tích phân và bộ cộng hưởng.} Mô hình của bài~\ref{ex:18:2} với $g_L=1$, $\tau_m=10\ms$, $\tau_w=100\ms$, ngưỡng $1$ và đặt lại $V\to0$ ($W$ không đổi) nhận $\mu=1.5\sin2\pi ft$, $f=1$--$40$~Hz. Nơ-ron phát xung trong dải tần nào khi $\alpha=0$ và khi $\alpha=4$? So sánh với điều kiện $1.5\,|Z(f)|>1$.
\end{exercise}

\begin{exercise}[\lvlB]
\label{ex:18:5}
\emph{Mạch chốt từ một tế bào.} Nơ-ron có dòng bình nguyên: $\tau\,\dd r/\dd t=-r+F(wr+I)$, $F(x)=\min(1,\max(0,x))$, $\tau=10\ms$, $w=1.5$, nền $I=-0.2$. (a)~Xung $I=0.3$ phải kéo dài ít nhất bao lâu để đưa nơ-ron từ $r=0$ lên bình nguyên? (b)~Xung dài $I=-0.2-B$ có biên độ $B$ nhỏ nhất bằng bao nhiêu thì đưa nó về trạng thái nghỉ?
\end{exercise}

\begin{exercise}[\lvlB]
\label{ex:18:6}
\emph{Bộ tích phân tự điều chỉnh.} Khi mắt định thị, đầu vào $I=0$, cảm biến trượt cho tốc độ trôi $e=\dd r/\dd t$, và $\dd w/\dd t=-\eta\,e\,r$. Với $w=0.95$, $\tau=0.1$~s, $\langle r^2\rangle=1$, tìm $\eta$ để sau $60$~s định thị, $|1-w|$ lọt vào dung sai của bài~\ref{ex:18:1}. Điều gì hỏng nếu cũng học cả trong lúc mắt xoay?
\end{exercise}

\begin{exercise}[\lvlC]
\label{ex:18:7}
\emph{Tái cấu hình linh kiện để làm gì?} Một kênh quảng bá chuyển $\alpha$ trong mô hình của bài~\ref{ex:18:2} giữa $0$ và $4$. Bộ cộng hưởng hơn bộ tích phân bao nhiêu lần về tỉ số tín hiệu/nhiễu với một sóng sin yếu trong nhiễu trắng dòng điện ở $11$~Hz và ở $1$~Hz? Hãy nghĩ ra một bài toán mà chính việc chuyển chế độ đem lại lợi ích.
\end{exercise}

\chapter{Nơ-ron theo số đuôi gai}
\chaptermark{Số đuôi gai}
\label{sec:dendrites}

\section{Số đầu vào như một tham số của mạch}

Ở chương~\ref{sec:neurontypes} ta phân loại nơ-ron theo thứ mà đầu ra của chúng giải phóng, còn ở chương~\ref{sec:eetypes} theo linh kiện mà chúng đóng vai. Còn một dấu hiệu thứ ba mà kỹ sư sẽ nhận ra ngay: \emph{nơ-ron có bao nhiêu đầu vào}. Các đầu vào được gom bởi đuôi gai, những nhánh mà sợi trục của nơ-ron khác bám vào (chương~\ref{sec:dofa}, mục~\ref{sec:weightstore}). Có nơ-ron hoàn toàn không có đuôi gai, có nơ-ron có một hoặc hai, có nơ-ron có cả một cái cây gom tới một trăm nghìn đầu vào. Với kỹ sư điện tử, đây là hệ số ghép đầu vào, fan-in, và nó hầu như luôn tương ứng với nhiệm vụ.

Vì sao số lượng và kích thước đuôi gai quan trọng đến thế có thể thấy từ một ước tính đơn giản. Màng là một tụ điện với điện dung riêng $c_m\approx1$~\textmu F/cm$^2$, và diện tích $A$ của nơ-ron cùng các đuôi gai càng lớn thì càng phải mang tới nhiều điện tích để nâng điện áp lên $\Delta V$. Thời gian để dòng đầu vào $I$ làm việc đó, nếu bỏ qua rò, là
\begin{empheq}[box=\keybox]{equation}
t \;\approx\; \frac{c_m\,A\,\Delta V}{I} .
\label{eq:dendriteload}
\end{empheq}
Một nơ-ron nhỏ với vài đuôi gai ngắn có diện tích cỡ $2\cdot10^{-5}$~cm$^2$ và điện dung khoảng $20$~pF. Một khớp thần kinh lớn với dòng vài nanoampe nâng nó lên $20\mV$ trong chưa tới một phần mười mili giây. Nơ-ron có cây lớn có diện tích lớn hơn khoảng mười lăm lần và điện dung khoảng $300$~pF; một khớp thần kinh thông thường với dòng khoảng $0.1$~nA sẽ nạp nó trong $60\ms$, lâu hơn nhiều so với hằng số thời gian rò, tức là sẽ không bao giờ nạp được. Nơ-ron như vậy cần hàng chục đầu vào đồng thời. Các con số ở đây chỉ là bậc độ lớn, nhưng kết luận không phụ thuộc vào chúng: \emph{ít đầu vào thì nhanh và chính xác, nhiều đầu vào thì chậm và theo tổng}.

\begin{figure}[t]
\centering
\begin{tikzpicture}[>=Stealth,
  soma/.style={circle,draw,very thick,fill=blue!6,minimum size=0.55cm,inner sep=0pt},
  ax/.style={->,very thick,red!70!black},
  den/.style={very thick,blue!60!black}]
% 0 dendrites: sensor on a line
\begin{scope}[xshift=0cm]
\draw[den,decorate,decoration={snake,amplitude=1pt,segment length=4pt}] (-0.9,0) -- (0,0);
\node[font=\scriptsize] at (-0.9,0.3) {cảm biến};
\draw[ax] (0,0) -- (1.0,0);
\draw[thick] (0.1,0) -- (0.1,0.45);
\node[soma,minimum size=0.4cm] at (0.1,0.65) {};
\node[font=\scriptsize,align=center] at (0.05,-0.75) {$0$\\đường có cảm biến};
\end{scope}
% 1 dendrite
\begin{scope}[xshift=3.3cm]
\draw[den] (-0.9,0) -- (-0.28,0);
\node[soma] at (0,0) {};
\draw[ax] (0.28,0) -- (0.9,0);
\node[font=\scriptsize,align=center] at (0,-0.75) {$1$\\bộ đệm, bộ lặp};
\end{scope}
% 2 dendrites: one per ear
\begin{scope}[xshift=6.2cm]
\draw[den] (-0.28,0) -- (-0.9,0); \draw[den] (0.28,0) -- (0.9,0);
\node[soma] at (0,0) {};
\draw[ax] (0,-0.28) -- (0,-0.6);
\node[font=\scriptsize] at (-0.9,0.25) {trái}; \node[font=\scriptsize] at (0.9,0.25) {phải};
\node[font=\scriptsize,align=center] at (0,-1.0) {$2$\\AND theo thời gian};
\end{scope}
% ~4 short dendrites: mixer
\begin{scope}[xshift=9.0cm]
\foreach \a in {45,135,225,315}{\draw[den] (\a:0.28) -- (\a:0.7); \fill[blue!60!black] (\a:0.72) circle (0.06);}
\node[soma] at (0,0) {};
\draw[ax] (0.28,0) -- (1.0,0);
\node[font=\scriptsize,align=center] at (0,-1.0) {$\sim4$\\bộ trộn};
\end{scope}
% many: summing tree
\begin{scope}[xshift=12.0cm]
\foreach \a in {60,90,120}{\draw[den] (0,0.28) -- ++(\a:0.55) coordinate (b); \foreach \d in {-20,20}{\draw[den] (b) -- ++(\a+\d:0.35);}}
\foreach \a in {-120,-60}{\draw[den] (0,-0.28) -- ++(\a:0.45);}
\node[soma] at (0,0) {};
\draw[ax] (0.28,0) -- (1.0,0);
\node[font=\scriptsize,align=center] at (0,-1.0) {$10^4$--$10^5$\\bộ cộng};
\end{scope}
\end{tikzpicture}
\caption{Năm lớp theo số đầu vào. Màu xanh là đuôi gai (đầu vào), màu đỏ là sợi trục (đầu ra). Đầu vào càng ít, nơ-ron truyền một tín hiệu riêng lẻ càng nhanh và chính xác; càng nhiều, nó cộng và phân lớp càng tốt. Đây là sơ đồ, không phải hình vẽ hình dạng tế bào.}
\label{fig:dendriteclasses}
\end{figure}

\section{Không đuôi gai: cảm biến ở đầu đường truyền}

Các nơ-ron cảm giác của xúc giác, cơn đau và độ căng cơ (các chương~\ref{sec:sensors}, \ref{sec:pain}, \ref{sec:muscles}) không có một đuôi gai nào trên thân tế bào. Sợi trục rời thân thành một nhánh rồi lập tức chia đôi: một nhánh đi tới da hoặc cơ, và đầu mút của nó tự làm cảm biến; nhánh kia đi vào tủy sống. Xung chạy thẳng từ cảm biến vào tủy sống, bỏ qua thân tế bào. Với kỹ sư, đây là một đường truyền có cảm biến ở đầu xa, còn thân tế bào treo bên cạnh như một bộ nguồn: nó cung cấp cho đường truyền nhưng không tham gia truyền. Lợi ích là tốc độ và độ tin cậy: trên đường đi không có chỗ nào phải cộng tín hiệu rồi quyết định lại có truyền tiếp hay không.

Cảm biến ánh sáng và âm thanh là trường hợp còn cực đoan hơn: chúng không phải nơ-ron gom tín hiệu mà chính là bộ chuyển đổi, trong đó đầu vào không phải khớp thần kinh mà là một protein thụ thể (hình~\ref{fig:transducer}).

\section{Một đuôi gai: bộ đệm}

Nơ-ron có một đuôi gai và một sợi trục nhận một đường đầu vào và chuyển tiếp nó. Các nơ-ron khứu giác được cấu tạo như vậy: đuôi gai duy nhất vươn ra bề mặt mũi và mang một loại thụ thể duy nhất~\cite{Mombaerts2004}; các tế bào trung chuyển của võng mạc nằm giữa cảm biến ánh sáng và các nơ-ron đầu ra của mắt cũng vậy; các nơ-ron nối cảm biến của tai với não cũng vậy. Kỹ sư sẽ gọi chúng là bộ đệm hay bộ lặp: chúng không tính toán mà chuyển một tín hiệu, đôi khi thay đổi dạng của nó, chẳng hạn biến điện áp biến thiên liên tục của cảm biến thành xung, như ở chương~\ref{sec:sensors}.

\section{Hai đuôi gai: AND theo thời gian}

Các bộ phát hiện trùng khớp xác định hướng của âm thanh (hình~\ref{fig:itd}) ở động vật có vú thường có đúng hai đuôi gai hướng về hai phía ngược nhau: một đuôi gai nhận đầu vào từ tai trái, đuôi gai kia từ tai phải~\cite{Grothe2010}. Tại sao phải tách các đầu vào ra? Hãy nhớ lại công thức~\eqref{eq:shunt}: khi trên một đoạn màng có nhiều kênh kích thích mở, điện áp ở đó tiến gần đến điện thế cân bằng, và mỗi đầu vào tiếp theo đóng góp ngày càng ít. Vì vậy nhiều đầu vào từ một tai, dồn trên một đuôi gai, làm nó bão hòa và một mình không thể đưa nơ-ron tới ngưỡng. Còn mỗi đuôi gai một đầu vào thì cộng lại gần như tuyến tính. Hai đuôi gai biến nơ-ron thành một cổng AND chặt chẽ hơn: chỉ phát xung nếu tín hiệu đến từ cả hai phía và cùng lúc. Các mô hình đã cho thấy việc tách như vậy thực sự cải thiện khả năng phát hiện trùng khớp~\cite{AgmonSnir1998}.

\section{Vài đuôi gai ngắn: bộ trộn và bộ chuyển tiếp}

\emph{Bộ trộn.} Trong mạch học vận động ở chương~\ref{sec:motorlearning}, khoảng bảy mươi tỉ nơ-ron \nt{GLUT} nhỏ mở rộng đầu vào thành một vectơ rất dài. Mỗi nơ-ron chỉ có khoảng bốn đuôi gai ngắn, và mỗi đuôi gai kết thúc bằng một ``móng vuốt'' ôm lấy đầu tận cùng của một đầu vào. Như vậy, mỗi bộ trộn chỉ thấy khoảng bốn đầu vào trong số rất nhiều, và hoạt động như một cổng AND nhỏ trên bộ bốn của mình. Từ $M$ đường đầu vào có thể chọn ra $\binom{M}{4}$ bộ bốn khác nhau: với $M=100$ là gần bốn triệu. Cần nhiều như vậy để làm gì? Một phần tử ngưỡng có $n$ đầu vào có thể phân đúng thành hai lớp không quá
\begin{empheq}[box=\keybox]{equation}
P_{\max} \;\approx\; 2n
\label{eq:cover}
\end{empheq}
mẫu ngẫu nhiên~\cite{Cover1965}. Nơ-ron đầu ra của cùng mạch đó nhận khoảng $10^5$ đầu vào từ các bộ trộn, và theo~\eqref{eq:cover}, giới hạn trên của nó là cỡ $2\cdot10^5$ phép ánh xạ khác nhau. Đó là giới hạn cho một phần tử lý tưởng: nếu mọi trọng số đều dương, như ở các khớp thần kinh kích thích, và cần có dự trữ cho nhiễu, thì ước tính cho chính nơ-ron này là khoảng $4\cdot10^4$ phép ánh xạ~\cite{Brunel2004}. Các bộ trộn với ít đầu vào tồn tại chính là để bộ cộng lớn có nhiều đầu vào khác nhau và gần như độc lập~\cite{Marr1969,Albus1971}.

\emph{Bộ chuyển tiếp.} Trên đường từ tai tới các bộ phát hiện hướng có những nơ-ron với ít đuôi gai ngắn, chỉ nhận vài khớp thần kinh khổng lồ, và một số về thực chất chỉ nhận một. Theo~\eqref{eq:dendriteload}, đây đúng là điều cần thiết: điện dung nhỏ và dòng lớn, nên nơ-ron phát xung chỉ sau một phần nhỏ của mili giây kể từ mỗi xung đầu vào, hầu như không mất độ chính xác về thời gian. Một trong các bộ chuyển tiếp ấy nhận \nt{GLUT} nhưng xuất \nt{GLYC}: đó là một bộ đảo với độ chính xác vài chục micro giây, cung cấp ức chế nhanh cho các bộ phát hiện hướng~\cite{BorstSoria2012}.

\section{Hàng nghìn đầu vào: bộ cộng và bộ phân lớp}

Ở đầu kia của thang đo là các nơ-ron \nt{GLUT} vỏ não với khoảng mười nghìn đầu vào và các nơ-ron \nt{GABA} đầu ra của mạch học vận động với một trăm nghìn. Theo~\eqref{eq:dendriteload}, nơ-ron như vậy chậm và không thể đáp ứng với một đầu vào riêng lẻ: nó cộng. Đó là nơ-ron tích phân của chương~\ref{sec:glutcomputer} và cũng là bộ cộng dùng để dựng các bộ phân lớp. Cây lớn còn thêm một điều mới: các nhánh của nó hoạt động như những mạch con riêng với ngưỡng riêng, nên một nơ-ron hành xử như một mạng nhiều lớp nhỏ~\cite{Beniaguev2021,Gidon2020}.

\section{Đuôi gai kiêm đầu ra}

Có những nơ-ron hoàn toàn không có sợi trục, và đuôi gai vừa là đầu vào vừa là đầu ra: chúng nhận tín hiệu và tự giải phóng chất dẫn truyền thần kinh. Ví dụ được nghiên cứu kỹ nhất là các nơ-ron \nt{GABA} của võng mạc tham gia xác định hướng chuyển động (chương~\ref{sec:gaba}, hình~\ref{fig:direction}). Ở nơ-ron như vậy, mỗi nhánh tự đáp ứng với chuyển động từ thân tế bào về phía đầu mút của nó và phát ra ức chế từ đầu mút ấy~\cite{Euler2002}. Một nơ-ron là vài bộ phát hiện hướng độc lập, mỗi nhánh một bộ. Với kỹ sư, đây không phải là một phần tử, mà là một vi mạch nhiều kênh trong một vỏ.

\section{Tóm tắt}

\begin{table}[t]
\centering
\footnotesize
\setlength{\tabcolsep}{3pt}
\renewcommand{\arraystretch}{1.15}
\begin{tabular}{>{\raggedright}p{1.9cm}>{\raggedright}p{3.4cm}>{\raggedright}p{2.6cm}>{\raggedright}p{1.8cm}>{\raggedright\arraybackslash}p{3.8cm}}
\toprule
Số đuôi gai & Ví dụ & Linh kiện & Số đầu vào & Điều đã biết \\
\midrule
$0$ & cảm biến xúc giác, đau, độ căng & đường truyền có cảm biến ở đầu & --- & đã xác lập \\
$1$ & nơ-ron khứu giác, tế bào trung chuyển của võng mạc, nơ-ron tai & bộ đệm, bộ lặp & $1$ đến vài & đã xác lập \\
$2$ & bộ phát hiện hướng âm thanh & AND theo thời gian & hai bó & cấu tạo đã xác lập; lợi ích của việc tách đầu vào được chỉ ra trên mô hình \\
$\sim4$ & bộ trộn của mạch học vận động & cổng AND nhỏ & $\sim4$ & đã xác lập; vai trò mở rộng đầu vào là giả thuyết có cơ sở vững \\
ít, ngắn & bộ chuyển tiếp trên đường từ tai & bộ chuyển tiếp, bộ đảo & $1$ đến vài, khổng lồ & đã xác lập \\
hàng nghìn & nơ-ron \nt{GLUT} vỏ não, nơ-ron đầu ra học vận động & bộ cộng, bộ phân lớp & $10^4$--$10^5$ & đã xác lập; nhánh như mạch con đã đo ở một số ví dụ \\
đuôi gai kiêm đầu ra & nơ-ron \nt{GABA} hướng ở võng mạc & vi mạch nhiều kênh & nhiều & đã đo \\
\bottomrule
\end{tabular}
\caption{Nơ-ron theo số đuôi gai.}
\label{tab:dendrites}
\end{table}

\begin{keyblock}
\begin{proposition}[Số đầu vào đi theo nhiệm vụ]
\label{prop:dendrites}
Ở nơi cần tốc độ và độ chính xác của từng tín hiệu riêng lẻ, tức ở cảm biến, bộ chuyển tiếp và bộ phát hiện trùng khớp, nơ-ron có ít đuôi gai, ít đầu vào và khớp thần kinh lớn: điện dung nhỏ~\eqref{eq:dendriteload} được nạp nhanh. Ở nơi cần cộng và phân lớp, nơ-ron có cây lớn với hàng nghìn đầu vào. Cấu tạo của các lớp đã được đo; lời giải thích về tính toán có cơ sở vững, nhưng với nhiều lớp vẫn là giả thuyết.
\end{proposition}
\end{keyblock}

Bảng~\ref{tab:dendrites} bổ sung cho hai cách phân loại trước: theo chất dẫn truyền thần kinh (bảng~\ref{tab:types}) và theo linh kiện (bảng~\ref{tab:eetypes}). Cả ba nhìn cùng một phần tử từ các phía khác nhau: nó xuất ra gì, nó tính gì và nó nghe bao nhiêu.

\section{Bài tập}

\begin{exercise}[\lvlA]
\label{ex:19:1}
\emph{Nơ-ron lớn cần bao nhiêu đầu vào.} Nơ-ron có $C=300$~pF và $\tau_m=20\ms$ nhận các khớp thần kinh là nguồn dòng $0.1$~nA mỗi khớp, ngưỡng cao hơn mức nghỉ $20\mV$. (a)~Cần bao nhiêu khớp hoạt động đồng thời để tới được ngưỡng; tới ngưỡng trong $10\ms$; trong $5\ms$? (b)~Nơ-ron có $C=20$~pF tới ngưỡng sau bao lâu với một khớp $4$~nA?
\end{exercise}

\begin{exercise}[\lvlA]
\label{ex:19:2}
\emph{Vì sao là bốn.} Bộ trộn là cổng AND trên $k$ đường trong số $M=100$, trong một mẫu có một nửa số đường hoạt động, có $7\cdot10^{10}$ bộ trộn. (a)~Bao nhiêu bộ đáp ứng với một mẫu khi $k=2$; $4$; $40$? (b)~Bộ trộn tăng số phép ánh xạ~\eqref{eq:cover} của nơ-ron đầu ra có $10^5$ đầu vào lên bao nhiêu lần so với $M=100$ đường trực tiếp?
\end{exercise}

\begin{exercise}[\lvlB]
\label{ex:19:3}
\emph{$2n$ từ đâu ra.} Một siêu phẳng qua gốc tọa độ chia $P$ điểm ở vị trí tổng quát trong không gian $n$ chiều thành hai lớp theo $C(P,n)=2\sum_{k=0}^{n-1}\binom{P-1}{k}$ cách. (a)~Chứng minh rằng khi $P=2n$, đúng một nửa trong $2^P$ cách phân chia là tách được. (b)~Tỉ lệ tách được là bao nhiêu khi $n=100$ và $P=180$; $220$?
\end{exercise}

\begin{exercise}[\lvlB]
\label{ex:19:4}
\emph{Hai đuôi gai như một cổng AND chặt.} Đuôi gai là một đoạn có rò $g_L$, mỗi đầu vào mở trên nó độ dẫn $g_0=0.5\,g_L$ với điện thế $E_E$, thân thấy $\kappa(V_1+V_2)$, ngưỡng $\theta=1.1\,\kappa E_E$. Tại sao không có số đầu vào nào trên một đuôi gai kích hoạt được nơ-ron, và cần bao nhiêu đầu vào trên mỗi đuôi gai?
\end{exercise}

\begin{exercise}[\lvlB]
\label{ex:19:5}
\emph{Thiết kế bộ chuyển tiếp.} Độ rung pha của xung là $\sigma_t=\sigma_V/\dot V$, trong đó $\dot V$ là độ dốc điện áp gần ngưỡng. Cần $\sigma_t\le20$~\textmu s với nhiễu $\sigma_V=1\mV$; bỏ qua rò. Nơ-ron có $C=300$~pF cần bao nhiêu khớp đồng bộ $0.1$~nA mỗi khớp, và độ rung pha của nơ-ron có $C=20$~pF với một khớp $4$~nA là bao nhiêu?
\end{exercise}

\begin{exercise}[\lvlB]
\label{ex:19:6}
\emph{Mô phỏng: nạp một đuôi gai dài.} Mô phỏng cáp thụ động $\tau_m\partial_tV=\lambda^2\partial_x^2V-V+i/g$ có độ dài $L=\ell/\lambda$ với hai đầu kín ($40$ đoạn, phương pháp Euler). Một xung dòng ngắn được đưa vào đầu xa: khi nào đầu gần đạt cực đại, và đạt bao nhiêu phần so với điện áp do cùng điện tích đó trải đều trên cả đuôi gai, khi $L=0.2$; $1$; $2$?
\end{exercise}

\begin{exercise}[\lvlC]
\label{ex:19:7}
\emph{Nghiên cứu: số đầu vào tối ưu.} Điện dung $C=C_0+nc_s$, có $m$ đầu vào hoạt động đồng thời với dòng $I_s$ mỗi đầu vào, nơ-ron ghi nhớ $P\approx2n$ phép ánh xạ~\eqref{eq:cover}. Thời gian đáp ứng phụ thuộc vào $P$ thế nào khi $m$ không đổi và khi tỉ lệ $m=\varphi n$ không đổi? Đề xuất một tiêu chí chi phí và tìm $n$ tối ưu cho bộ chuyển tiếp và cho bộ phân lớp.
\end{exercise}

\chapter{Nơ-ron làm gì khi ngủ}
\chaptermark{Giấc ngủ}
\label{sec:sleep}
\begin{chaptergoals}
\item vì sao giấc ngủ là tập hợp chế độ do \nt{ACET}, \nt{NORA} và \nt{SERO} chuyển, không phải tắt máy;
\item cách áp lực ngủ và hai ngưỡng tạo thành rơ-le có vòng trễ khóa theo nhịp ngày đêm;
\item cách sự mỏi biến vỏ não thành bộ dao động tích thoát, và cách \nt{ACET} dừng các dao động;
\item cách một ngày được phát lại tua nhanh, và vì sao trọng số có thể được thu nhỏ khi ngủ.
\end{chaptergoals}

\section{Ngủ không phải là tắt máy, mà là một chế độ khác}

Nhìn một chiếc máy tính đã tắt, kỹ sư sẽ nói: nó không làm gì cả. Với bộ não đang ngủ thì không thể nói như vậy. Khi ngủ, nơ-ron không im lặng: trong giấc ngủ sâu, vỏ não vẫn phát xung, trong giấc ngủ REM (giấc ngủ chuyển động mắt nhanh) thì gần bằng ban ngày. Cái thay đổi không phải là nơ-ron \emph{có làm việc} hay không, mà là chúng làm việc \emph{như thế nào}. Giấc ngủ là một tập hợp các chế độ của cỗ máy, và chúng được chuyển đổi bởi chính những kênh quảng bá đã hoạt động ban ngày.

Với mỗi chế độ có thể ghi ra ``trạng thái các thanh ghi'' (bảng~\ref{tab:sleepmodes}). Ban ngày \nt{ACET}, \nt{NORA} và \nt{SERO} hoạt động, các cảm biến được kết nối, cơ bắp nghe lệnh. Trong giấc ngủ sóng chậm, tức giấc ngủ sâu, cả ba kênh lắng xuống, còn đầu vào từ các giác quan bị yếu đi~\cite{Hasselmo1999}: lượng acetylcholine giải phóng trong vỏ não giảm, còn các nơ-ron \nt{NORA} và \nt{SERO} phát xung thưa hơn ban ngày~\cite{Marrosu1995,AstonJonesBloom1981,McGintyHarper1976}. Trong giấc ngủ REM, bức tranh thật lạ: \nt{ACET} lại tăng lên mức ban ngày~\cite{Marrosu1995}, còn \nt{NORA} và \nt{SERO} gần như im bặt hoàn toàn~\cite{AstonJonesBloom1981,McGintyHarper1976,JacobsAzmitia1992}. Trong giấc ngủ REM, cơ của cơ thể bị tắt: các nơ-ron \nt{ACET} điều khiển chúng bị ức chế mạnh qua \nt{GABA} và \nt{GLYC}~\cite{BrooksPeever2012}, chính là lệnh cấm ở chương~\ref{sec:gaba}, chỉ có điều được áp lên toàn bộ đầu ra cùng lúc.

\begin{table}[t]
\centering
\footnotesize
\setlength{\tabcolsep}{4pt}
\renewcommand{\arraystretch}{1.15}
\begin{tabular}{>{\raggedright}p{3.1cm}>{\raggedright}p{2.9cm}>{\raggedright}p{3.2cm}>{\raggedright\arraybackslash}p{3.4cm}}
\toprule
& thức & ngủ sóng chậm & ngủ REM \\
\midrule
\nt{ACET} (ghi, ch.~\ref{sec:acet}) & cao & thấp & cao \\
\nt{NORA} (khuếch đại, ch.~\ref{sec:nora}) & trung bình, có đợt bùng & thấp & gần như im lặng \\
\nt{SERO} (ch.~\ref{sec:serocomputer}) & trung bình & thấp & gần như im lặng \\
đầu vào từ cảm biến & kết nối & yếu đi & yếu đi \\
đầu ra tới cơ & kết nối & yếu đi & bị tắt bằng ức chế \\
hoạt động vỏ não & đều & dao động chậm: làm việc --- im lặng & đều, như ban ngày \\
\bottomrule
\end{tabular}
\caption{Trạng thái các thanh ghi của cỗ máy trong ba chế độ. Mọi dòng đều đã được đo; mức \nt{ACET}, \nt{NORA} và \nt{SERO} được đo bằng các bản ghi trực tiếp theo từng giai đoạn ngủ~\cite{Marrosu1995,AstonJonesBloom1981,McGintyHarper1976}.}
\label{tab:sleepmodes}
\end{table}

\section{Khi nào ngủ: rơ-le có vòng trễ}

Điều gì quyết định khi nào cỗ máy chuyển sang ngủ? Một mạch đơn giản gồm hai phần hoạt động khá tốt~\cite{Borbely1982}. Phần thứ nhất là \emph{áp lực ngủ} $S$, tích lũy trong khi ta thức và được xả trong khi ta ngủ. Đó là một tụ điện nạp vào ban ngày và phóng vào ban đêm:
\begin{empheq}[box=\keybox]{equation}
\frac{\dd S}{\dd t} \;=\; \frac{1-S}{\tau_r}\ \ \text{(thức)},\qquad
\frac{\dd S}{\dd t} \;=\; -\frac{S}{\tau_d}\ \ \text{(ngủ)} .
\label{eq:sleeppressure}
\end{empheq}
Phần thứ hai là hai ngưỡng: cỗ máy ngủ khi $S$ chạm ngưỡng trên $H(t)$ và thức dậy khi $S$ hạ xuống ngưỡng dưới $L(t)$. Cả hai ngưỡng dao động nhẹ theo nhịp ngày đêm ở chương~\ref{sec:chemoutput}. Kết quả là một rơ-le có vòng trễ, tức trigger Schmitt ở mục~\ref{sec:bistable}, và cùng với tụ điện nó tạo thành một bộ dao động tích thoát, được nhịp ngày đêm khóa pha~\eqref{eq:pll}.

\begin{figure}[t]
\centering
\begin{tikzpicture}[>=Stealth,x=0.16cm,y=4.2cm]
\foreach \a/\b in {0/4.3,21.2/29.3,45.4/53.4,69.5/72}{\fill[blue!8] (\a,0) rectangle (\b,1);}
\draw[->,gray!70!black] (0,0) -- (75,0) node[right,font=\small,black] {$t$, giờ};
\draw[->,gray!70!black] (0,0) -- (0,1.08) node[left,font=\small,black] {$S$};
\foreach \t in {0,24,48,72}{\draw[gray!60!black] (\t,-0.02) -- (\t,0.02); \node[below,font=\scriptsize] at (\t,0) {\t};}
\draw[thick,densely dashed,red!70!black] plot file {figures/sleep_H.dat};
\draw[thick,densely dashed,green!45!black] plot file {figures/sleep_L.dat};
\draw[very thick,blue!60!black] plot file {figures/sleep_S.dat};
\node[font=\scriptsize,red!70!black,anchor=west] at (73,0.72) {$H$: ngủ};
\node[font=\scriptsize,green!45!black,anchor=west] at (73,0.24) {$L$: thức dậy};
\node[font=\scriptsize,blue!60!black] at (25.25,0.93) {ngủ};
\node[font=\scriptsize,blue!60!black] at (49.4,0.93) {ngủ};
\end{tikzpicture}
\caption{Áp lực ngủ $S$~\eqref{eq:sleeppressure} với $\tau_r=18.2$~giờ, $\tau_d=4.2$~giờ. Ban ngày $S$ nạp tới ngưỡng trên $H$, ban đêm (vùng tô màu) phóng xuống ngưỡng dưới $L$; cả hai ngưỡng dao động theo nhịp ngày đêm. Cỗ máy tự đi tới chế độ khoảng $16$ giờ thức và $8$ giờ ngủ.}
\label{fig:twoprocess}
\end{figure}

Trong mô hình của chúng tôi, với các giá trị thông thường $\tau_r=18.2$~giờ và $\tau_d=4.2$~giờ, cỗ máy tự đi tới $16.1$ giờ thức và $8.0$ giờ ngủ (hình~\ref{fig:twoprocess}). Mô hình còn giải thích được điều có vẻ lạ: sau bốn mươi giờ không ngủ, áp lực lên tới $0.90$, nhưng giấc ngủ bù trong mô hình chỉ kéo dài $8.8$ giờ, chứ không dài hơn bình thường cả một ngày. Sự phóng điện có dạng hàm mũ, và mức nạp cao mất đi nhanh; ``món nợ'' được trả không bằng độ dài mà bằng độ sâu của giấc ngủ.

Về mặt vật lý, cái gì đóng vai trò $S$? Ứng viên mạnh là adenosine, ``bộ đếm tải'' ở phụ lục~\ref{sec:othermediators}: nồng độ của nó tăng trong thời gian thức và giảm khi ngủ~\cite{PorkkaHeiskanen1997,Dunwiddie2001}. Việc áp lực ngủ chính là adenosine chứ không gì khác là giả thuyết.

\section{Ngủ sóng chậm: cả mạng như một bộ dao động tích thoát}

Trong giấc ngủ sâu, các nơ-ron vỏ não làm việc luân phiên theo cả nhóm: chưa tới một lần mỗi giây, chúng cùng nhau bật lên trong vài trăm mili giây (\emph{trạng thái làm việc}) rồi cùng nhau im lặng (\emph{trạng thái im lặng}): tần số của dao động này dưới $1$~Hz, khi gây mê thường là $0.3$--$0.4$~Hz~\cite{Steriade1993}. Ta có thể lắp dao động chậm này từ những chi tiết đã có. Lấy một nhóm nơ-ron \nt{GLUT} có phản hồi mạnh về chính nó, tức bộ nhớ ở chương~\ref{sec:glutcomputer} với hai trạng thái bền, và thêm vào một sự mệt mỏi chậm, như trong bộ tạo nhịp ở chương~\ref{sec:muscles}:
\begin{empheq}[box=\keybox]{equation}
\tau\,\frac{\dd r}{\dd t} = -r + F\bigl(w\,r + I - a\bigr),
\qquad
\tau_a\,\frac{\dd a}{\dd t} = -a + b\,r .
\label{eq:updown}
\end{empheq}
Nhóm bật lên, làm việc và mệt đi; sự mệt mỏi $a$ ăn mòn trạng thái trên, và nhóm rơi vào im lặng; trong im lặng, sự mệt mỏi qua đi, trạng thái dưới biến mất, và nhóm lại bật lên. Kết quả là cùng một bộ dao động tích thoát như áp lực ngủ, chỉ nhanh hơn khoảng tám mươi nghìn lần.

\begin{figure}[t]
\centering
\begin{tikzpicture}[>=Stealth,x=2.9cm,y=1.0cm]
\begin{scope}[yshift=1.9cm]
\draw[->,gray!70!black] (0,0) -- (4.1,0);
\draw[->,gray!70!black] (0,0) -- (0,1.25) node[left,font=\small,black] {$r$};
\draw[thick,blue!60!black] plot file {figures/updown_sleep.dat};
\node[font=\scriptsize,anchor=west] at (4.15,0.5) {ngủ: $b=5$};
\end{scope}
\draw[->,gray!70!black] (0,0) -- (4.1,0) node[right,font=\small,black] {$t$, s};
\draw[->,gray!70!black] (0,0) -- (0,1.25) node[left,font=\small,black] {$r$};
\draw[thick,red!70!black] plot file {figures/updown_wake.dat};
\node[font=\scriptsize,anchor=west] at (4.15,0.5) {ngày: $b=1$};
\foreach \t in {0,1,2,3,4}{\draw[gray!60!black] (\t,-0.05) -- (\t,0.05); \node[below,font=\scriptsize] at (\t,0) {\t};}
\end{tikzpicture}
\caption{Cùng một nhóm~\eqref{eq:updown} ở hai chế độ: $w=5$, $I=2$, $\theta=2.5$, $\tau=10\ms$, $\tau_a=0.3$~s. Với mệt mỏi mạnh ($b=5$, như trong giấc ngủ sóng chậm), nhóm dao động giữa làm việc và im lặng với chu kỳ $1.1$~s (giá trị của mô hình; trong vỏ não chu kỳ dài hơn một giây, khi gây mê khoảng $3$~s). Làm yếu sự mệt mỏi ($b=1$), đúng như tác dụng của acetylcholine, và cùng nhóm ấy làm việc đều đặn.}
\label{fig:updown}
\end{figure}

Vậy cái gì chuyển mạng từ dao động sang làm việc đều đặn ban ngày? Acetylcholine ức chế trong nơ-ron các dòng chậm gây ra mệt mỏi~\cite{McCormick1992}. Trong mô hình của chúng tôi, khi mệt mỏi mạnh, $b=5$, nhóm dao động với chu kỳ $1.1$~s, nhanh hơn dao động đo được nhưng cùng bậc, và làm việc khoảng $35\%$ thời gian nếu lấy trung bình trên một số nguyên chu kỳ; khi mệt mỏi yếu, $b=1$, cùng nhóm ấy làm việc đều đặn (hình~\ref{fig:updown}). Một thanh ghi, tức mức \nt{ACET}, chuyển bộ dao động thành bộ khuếch đại. Chồng lên dao động chậm trong giấc ngủ còn có những đợt bùng nhanh hơn của nhịp $7$--$15$ dao động mỗi giây; chúng sinh ra từ vòng \nt{GLUT}--\nt{GABA} giữa vỏ não và một nút chuyển tiếp trung gian~\cite{SteriadeMcCormickSejnowski1993}, họ hàng với nhịp ở chương~\ref{sec:gaba}.

\section{Phát lại: tua nhanh một ngày}

Ở chương~\ref{sec:acet} ta đã thấy rằng trong giấc ngủ sóng chậm, các nơ-ron từng làm việc cùng nhau ban ngày lại phát xung cùng nhau và theo cùng thứ tự. Ta thêm một chi tiết kỹ thuật: việc phát lại diễn ra \emph{có tua nhanh}. Một chuỗi ban ngày kéo dài vài giây được phát lại trong giấc ngủ trong vài phần mười giây~\cite{Nadasdy1999}: ở hồi hải mã nhanh hơn khoảng hai mươi lần~\cite{LeeWilson2002}, ở vỏ não nhanh hơn sáu đến bảy lần~\cite{Euston2007}. Và nó không được phát lại tùy lúc: các đợt phát lại ngắn ở một vùng khớp nhịp với dao động làm việc--im lặng và với các đợt bùng nhịp nhanh ở vỏ não. Nếu tăng cường nhân tạo sự phối hợp này, trí nhớ được cải thiện~\cite{Maingret2016}; đây đã là quan hệ nhân quả chứ không còn là trùng hợp.

Vì sao cỗ máy cần tua nhanh một ngày? Kỹ sư quen với một khó khăn gọi là \emph{quên thảm khốc}: một mạng học điều mới liên tiếp, cái này rồi cái khác, sẽ làm hỏng cái cũ. Cách cứu là \emph{trộn lẫn}: học cái mới xen kẽ với cái cũ, từng phần nhỏ, nhiều lần. Giả thuyết về hai hệ thống học~\cite{McClelland1995} cho rằng bộ nhớ nhanh ghi lại cả ngày, rồi trong giấc ngủ, từng chút một, xen kẽ và nhiều lần, phát lại nó cho vỏ não chậm. Đây chính là cặp ``bộ đệm nhanh --- kho lưu trữ chậm'' ở chương~\ref{sec:acet}, và cũng là cặp người học nhanh và chậm ở chương~\ref{sec:motorlearning}. Việc phát lại và sự phối hợp của nó đã được đo; việc mục đích của nó là học trộn lẫn cho bộ nhớ chậm là một giả thuyết có cơ sở vững.

\section{Chuẩn hóa lại trọng số}

Ban ngày, các quy tắc Hebb ở chương~\ref{sec:dofa} chủ yếu làm mạnh các liên kết. Nếu chỉ làm mạnh, trọng số tăng lên, cùng với đó chi phí năng lượng tăng và độ nhạy giảm: một nơ-ron có mọi đầu vào đều mạnh sẽ không còn phân biệt được chúng. Theo giả thuyết \emph{cân bằng nội môi khớp thần kinh}~\cite{TononiCirelli2014}, giấc ngủ cần thiết một phần là để đưa trọng số về mức làm việc: mọi liên kết yếu đi cùng một số lần, nên \emph{tỉ lệ} giữa chúng, tức những gì đã học, được giữ nguyên. Với kỹ sư, đây là phép chuẩn hóa, như trong quy tắc~\eqref{eq:oja}, chỉ có điều không làm trong khi chạy mà làm vào một thời điểm riêng.

Các phép đo trực tiếp không mâu thuẫn với điều này: ở chuột nhắt, diện tích tiếp xúc tại các khớp thần kinh vỏ não sau khi ngủ nhỏ hơn khoảng $18\%$ so với sau khi thức, mức giảm tỉ lệ với kích thước, còn các tiếp xúc lớn nhất và bền nhất hầu như không đổi~\cite{deVivo2017}. Việc co giãn trọng số đã được đo; việc đó là nhiệm vụ chính của giấc ngủ là giả thuyết.

\section{Ngủ REM: cỗ máy bị ngắt khỏi đầu vào và đầu ra}

Trong giấc ngủ REM, vỏ não làm việc gần như ban ngày, nhưng đầu vào từ giác quan bị yếu đi, còn đầu ra tới cơ bị ức chế tắt hẳn. Với kỹ sư, đây là một chế độ quen thuộc: \emph{thử nghiệm trên bàn thử}. Mạch chạy hết công suất, nhưng đầu vào được nối với bộ tạo tín hiệu bên trong, còn đầu ra được tách khỏi cơ cấu chấp hành để không làm hỏng gì. Theo một giả thuyết, giấc mơ chính là một lần chạy thử bên trong như vậy của mô hình thế giới được tích lũy trong ngày; mạng khi đó làm gì và có học hay không thì chưa biết. Ở đây chỉ những gì ghi trong bảng~\ref{tab:sleepmodes} là đã đo: kênh nào bật, kênh nào im, và đầu ra bị chặn.

\section{Bảo trì và giấc ngủ cục bộ}

Từ lâu người ta cho rằng trong khi ngủ, các khe gian bào giãn rộng và não được rửa sạch các sản phẩm trao đổi chất nhanh hơn~\cite{Xie2013}. Các thí nghiệm gần đây đo việc rửa bằng một cách khác lại thu được kết quả ngược lại: khi ngủ, việc rửa chậm hơn~\cite{Miao2024}. Câu hỏi này hiện còn mở, và chúng tôi để nó mở.

Một sự thật đáng tin cậy hơn: giấc ngủ là tính chất của các mạng cục bộ, chứ không phải của toàn bộ cỗ máy cùng lúc. Nếu một con vật bị giữ thức lâu, thì từng vùng riêng lẻ của vỏ não nó, trong khi con vật vẫn thức và cử động, thỉnh thoảng rơi vào im lặng trong chốc lát như trong giấc ngủ sóng chậm, và vào những lúc ấy con vật mắc lỗi nhiều hơn~\cite{Vyazovskiy2011}. Mỗi mạng tự mệt và tự chuyển chế độ; chế độ chung xuất hiện khi các kênh quảng bá chuyển tất cả các mạng cùng lúc.

\begin{keyblock}
\begin{proposition}[Giấc ngủ như một tập hợp chế độ]
\label{prop:sleep}
Khi ngủ, nơ-ron không tắt: các kênh quảng bá chuyển cỗ máy sang những chế độ khác. Khi nào ngủ do một rơ-le có vòng trễ~\eqref{eq:sleeppressure} quyết định, được nhịp ngày đêm chỉnh pha. Trong giấc ngủ sóng chậm, mạng trở thành một bộ dao động tích thoát~\eqref{eq:updown} và phát lại ngày đã qua ở tốc độ tua nhanh; trong giấc ngủ REM, nó làm việc trong khi bị ngắt khỏi đầu vào và đầu ra. Các chế độ, việc phát lại và việc co giãn trọng số đã được đo; mục đích của chúng, tức học trộn lẫn và chuẩn hóa lại, là những giả thuyết có cơ sở vững.
\end{proposition}
\end{keyblock}

\section{Tóm tắt}

\begin{table}[t]
\centering
\footnotesize
\setlength{\tabcolsep}{4pt}
\renewcommand{\arraystretch}{1.15}
\begin{tabular}{>{\raggedright}p{3.2cm}>{\raggedright}p{4.4cm}>{\raggedright\arraybackslash}p{5.2cm}}
\toprule
Điều gì diễn ra & Bằng cách nào & Điều đã biết \\
\midrule
đổi chế độ & các thanh ghi \nt{ACET}, \nt{NORA}, \nt{SERO} (bảng~\ref{tab:sleepmodes}) & đã đo \\
khi nào ngủ & tụ điện $S$ và rơ-le có vòng trễ~\eqref{eq:sleeppressure} & mô hình mô tả tốt các thí nghiệm; adenosine là $S$ là giả thuyết \\
dao động chậm & nhóm có phản hồi và mệt mỏi~\eqref{eq:updown} & dao động đã đo; mô hình là mạch đơn giản nhất \\
tua nhanh một ngày & phát lại nhanh hơn $6$--$20$ lần, khớp nhịp với dao động & đã đo; tăng cường sự phối hợp cải thiện trí nhớ \\
phát lại để làm gì & học trộn lẫn cho bộ nhớ chậm & giả thuyết có cơ sở vững \\
chuẩn hóa lại trọng số & co giãn các liên kết trong khi ngủ & mức giảm tiếp xúc đã đo; vai trò chính của giấc ngủ là giả thuyết \\
ngủ REM & làm việc với đầu vào và đầu ra bị ngắt & chế độ đã đo; mục đích chưa biết \\
rửa sạch & giãn rộng khe gian bào & kết quả mâu thuẫn \\
giấc ngủ cục bộ & từng vùng rơi vào im lặng & đã đo \\
\bottomrule
\end{tabular}
\caption{Nơ-ron làm gì khi ngủ.}
\label{tab:sleep}
\end{table}

Bảng~\ref{tab:sleep} tóm tắt chương. Giấc ngủ hóa ra không phải là một quãng nghỉ trong công việc của cỗ máy, mà là việc bảo trì nó ngay trong khi chạy: chuyển chế độ bằng chính các kênh quảng bá, một bộ dao động làm từ cùng những chi tiết như nhịp bước đi, tua nhanh một ngày cho bộ nhớ chậm và chuẩn hóa lại trọng số. Ban ngày cỗ máy ghi, ban đêm nó chép lại và sắp xếp những gì đã ghi.

\section{Bài tập}

\begin{exercise}[\lvlA]
\label{ex:20:1}
\emph{Rơ-le không có nhịp ngày đêm.} Cho các ngưỡng không đổi: $H=0.67$, $L=0.17$ (giá trị trung bình của các ngưỡng trong mô hình ở hình~\ref{fig:twoprocess}). Từ~\eqref{eq:sleeppressure} suy ra thời gian thức, thời gian ngủ và chu kỳ của bộ dao động khi $\tau_r=18.2$~giờ, $\tau_d=4.2$~giờ. Tại sao mô hình với ngưỡng dao động vẫn đi tới $16.1+8.0=24$~giờ? Hiệu ứng này thuộc về linh kiện nào ở chương~\ref{sec:chemoutput}?
\end{exercise}

\begin{exercise}[\lvlA]
\label{ex:20:2}
\emph{Nợ ngủ.} Giấc ngủ bắt đầu ở áp lực $S_0$ kéo dài $t_s=\tau_d\ln(S_0/L)$, $\tau_d=4.2$~giờ, $L$ là ngưỡng dưới lúc thức dậy. (a)~Sau $40$~giờ không ngủ $S_0=0.90$, và giấc ngủ kéo dài $8.8$~giờ. $L$ bằng bao nhiêu? Với $L\approx0.092$ thông thường thì giấc ngủ kéo dài bao lâu? (b)~Qua một đêm, diện tích tiếp xúc khớp thần kinh giảm $18\%$. Trọng số phải tăng bao nhiêu trong ngày?
\end{exercise}

\begin{exercise}[\lvlB]
\label{ex:20:3}
\emph{Thiết kế ngưỡng.} Chọn các ngưỡng không đổi $H$ và $L$ sao cho rơ-le~\eqref{eq:sleeppressure} với $\tau_r=18.2$~giờ và $\tau_d=4.2$~giờ cho đúng $16$~giờ thức và $8$~giờ ngủ. Cỗ máy như vậy kém hơn cỗ máy có ngưỡng dao động ở điểm nào, nếu một đêm nó bị đánh thức sau $4$~giờ?
\end{exercise}

\begin{exercise}[\lvlB]
\label{ex:20:4}
\emph{Chu kỳ của dao động chậm.} Trong mô hình~\eqref{eq:updown}, $F(x)=1/\bigl(1+e^{-(x-\theta)/k}\bigr)$, $w=5$, $I=2$, $\theta=2.5$, $k=0.25$, $b=5$, $\tau\ll\tau_a=0.3$~s. (a)~Tìm các điểm ngoặt $a_{\text{trên}}$ và $a_{\text{dưới}}$ của đường cong $r=F(wr+I-a)$, nơi trạng thái làm việc và im lặng biến mất. (b)~Lấy $r\approx1$ khi làm việc và $r\approx0$ khi im lặng, tìm chu kỳ và tỉ lệ thời gian làm việc.
\end{exercise}

\begin{exercise}[\lvlB]
\label{ex:20:5}
\emph{Khi nào dao động tắt.} Trong mô hình của bài~\ref{ex:20:4}, điểm cân bằng nằm ở giao của đường cong $a=wr+I-F^{-1}(r)$ với đường thẳng $a=br$. Dao động xảy ra khi giao điểm nằm trên nhánh giữa, giữa hai điểm ngoặt. Điều này đúng với những $b$ nào? Bằng cách thay đổi $b$, \nt{ACET} chuyển bộ dao động thành bộ khuếch đại ra sao?
\end{exercise}

\begin{exercise}[\lvlB]
\label{ex:20:6}
\emph{Mô phỏng: nợ hay pha.} Trong \texttt{sleep\_models.py}, lấy lần thức dậy đầu tiên sau $48$~giờ và giữ cỗ máy thức trong $D=24$, $32$, $40$, $48$, $64$~giờ (\texttt{stay\_awake}, \texttt{days=8}). Giấc ngủ bù kéo dài bao lâu với mỗi $D$? Điều gì quyết định độ dài của nó?
\end{exercise}

\begin{exercise}[\lvlB]
\label{ex:20:7}
\emph{Mô phỏng: quên và trộn lẫn.} Nơ-ron tuyến tính $y=\mathbf w\cdot\mathbf x$, $n=40$, học theo quy tắc $\mathbf w\leftarrow\mathbf w-0.5\,(y-y^*)\,\mathbf x$. Các nhiệm vụ A và B có $10$ mẫu mỗi nhiệm vụ, $x_j\sim\mathcal N(0,1/n)$, $y^*=\pm1$. Sai số trung bình bình phương trên A là bao nhiêu sau $200$ epoch A rồi $200$ epoch B? Còn sau $200$ epoch A và B trộn lẫn?
\end{exercise}

\begin{exercise}[\lvlC]
\label{ex:20:8}
\emph{Nghiên cứu: trộn lẫn hay chuẩn hóa lại?} Nơ-ron của bài~\ref{ex:20:7} có ngưỡng; hai quy trình ban đêm: (i)~mọi trọng số nhân với $0.82$; (ii)~các mẫu cũ được lặp lại xen kẽ với mẫu mới. Mỗi quy trình làm gì với tỉ lệ giữa các trọng số và với đáp ứng của nơ-ron? Làm sao phân biệt hai giả thuyết theo kích thước khớp thần kinh sau khi ngủ?
\end{exercise}

\chapter{Cỗ máy từ nơ-ron sống}
\chaptermark{Cỗ máy từ nơ-ron sống}
\label{sec:livingmachine}
\begin{chaptergoals}
\item vì sao mạng nuôi cấy không có kênh quảng bá phát đợt bùng phát, và khoảng cách giữa chúng do đâu;
\item cách mạng sống thay đổi đáp ứng dưới phản hồi, và điều gì còn chưa được chứng minh;
\item cách organoid làm hồ chứa có lớp đọc ra tuyến tính được huấn luyện bên ngoài, trên silicon;
\item người ta đã dạy mẫu nuôi cấy ra sao, và vì sao người thầy vẫn đứng ngoài đĩa nuôi.
\end{chaptergoals}

\section{Bàn thử có sơ đồ hở}

Ở chương~\ref{sec:debugging}, ta gỡ lỗi một cỗ máy không có sơ đồ: đầu dò nghe được một nghìn nơ-ron trong số tám mươi tỷ, phần còn lại phải đoán. Ở vào tình thế đó, kỹ sư sẽ làm một việc đơn giản: lắp một mẩu mạch nhỏ trên bo cắm thử, nơi thấy rõ từng linh kiện. Với nơ-ron cũng làm được như vậy.

Nơ-ron vỏ não được tách rời rồi nuôi trên một con chip có mảng điện cực, từ vài chục đến hàng chục nghìn điện cực. Sau vài tuần, tế bào mọc nhánh và nối với nhau thành một mạng gồm hàng nghìn hoặc hàng trăm nghìn nơ-ron, phần lớn là \nt{GLUT}, cùng một tỷ lệ nhỏ hơn \nt{GABA} tùy theo mẫu: trong mẫu nuôi cấy hồi hải mã, con số này vào khoảng $6\%$ số nơ-ron~\cite{Benson1994}. Mỗi điện cực vừa nghe được, như đầu dò ở chương~\ref{sec:debugging}, vừa bơm được dòng điện, như một tín hiệu thử. Loại bàn thử thứ hai là \emph{organoid}: những khối mô thần kinh ba chiều cỡ khoảng một milimét, nuôi từ tế bào gốc người~\cite{Lancaster2013}. Trên các bàn thử này, mạng sống hoạt động được hàng tháng; có những nền tảng cho phép làm thí nghiệm trên organoid từ xa, qua mạng, liên tục hơn một trăm ngày~\cite{Jordan2024}. Lĩnh vực này được gọi là \emph{tính toán sinh học} hay \emph{trí tuệ organoid}~\cite{Smirnova2023}.

Bàn thử khác bộ não ở hai điểm quan trọng. Nó rất nhỏ: số nơ-ron ít hơn một trăm nghìn lần. Và nó không có kênh quảng bá: trong mẫu nuôi cấy từ vỏ não không có nơ-ron \nt{DOPA}, \nt{SERO}, \nt{NORA} hay \nt{ACET}. Đây là một cỗ máy có bus dữ liệu và điều khiển luồng, nhưng không có thanh ghi điều khiển. Hãy xem nó làm được gì.

\section{Mạng tự làm gì}

Để yên, mẫu nuôi cấy không im lặng, cũng không nhiễu đều. Thỉnh thoảng gần như cả mạng bùng lên cùng lúc, thành một \emph{đợt bùng phát}, rồi lại lặng đi; khoảng cách giữa các đợt từ một giây đến vài phút, tùy mẫu~\cite{Wagenaar2006}. Với kỹ sư, đây là hình ảnh quen thuộc: một bộ khuếch đại có hồi tiếp dương mạnh và không tải sẽ biến thành bộ dao động.

Cơ chế thì ta đã biết. Liên kết qua lại mạnh giữa các nơ-ron \nt{GLUT} khởi phát một trận thác, như ở chương~\ref{sec:gaba} khi các điều kiện của mệnh đề~\ref{prop:ei} bị vi phạm. Trận thác nhanh chóng làm cạn kho chất dẫn truyền thần kinh ở các khớp thần kinh~\eqref{eq:depression}, và mạng im lặng. Kho hồi phục với hằng số thời gian $\tau_{\text{ph}}$, và khi đã hồi được một tỷ lệ $x^*$ đủ cho trận thác mới thì một đợt bùng phát mới nổ ra. Khoảng cách giữa các đợt là
\begin{empheq}[box=\keybox]{equation}
T_{\text{đợt}} \;\approx\; \tau_{\text{ph}}\,\ln\frac{1}{1-x^*} .
\label{eq:burstinterval}
\end{empheq}
Với $\tau_{\text{ph}}=0.8\,\text{s}$ từ chương~\ref{sec:code} và $x^*=0.9$, ta được khoảng hai giây, với $x^*=0.99$ thì khoảng bốn giây. Đây là ước tính mô hình với $\tau_{\text{ph}}$ của cuốn sách, và nó rơi vào đầu ngắn của dải khoảng cách đã đo. Phép đo trực tiếp trên vi mẫu nuôi cấy cho thấy kho chất dẫn truyền hồi phục sau một đợt bùng phát trong vài chục giây~\cite{CohenSegal2011}, tức là ở đó $\tau_{\text{ph}}$ lớn hơn; với $\tau_{\text{ph}}$ như vậy, công thức~\eqref{eq:burstinterval} cho vài chục giây đến vài phút, như ở những mẫu chậm. Đây là một bộ dao động tích thoát, họ hàng với bộ tạo nhịp ở chương~\ref{sec:muscles}: ở đó nó được nạp lại nhờ sự mỏi của các nhóm, còn ở đây là nhờ kho chất dẫn truyền.

Các đợt bùng phát là việc mạng làm khi không có việc gì để làm. Trong bộ não sống, hình ảnh tương tự xuất hiện trong giấc ngủ sâu (chương~\ref{sec:sleep}) và trong bộ não đang phát triển, trước khi có đầu vào thật. Sự giống nhau này gợi ý nhiều điều, nhưng việc cơ chế là một vẫn chỉ là giả thuyết.

\section{Dạy mạng khi không có thầy dopamine}

Vì mẫu nuôi cấy không có \nt{DOPA}, tín hiệu ``tốt hơn hay tệ hơn'' phải tự ta đưa vào, bằng điện cực. Thí nghiệm cho thấy mạng trên bàn thử quả thật thay đổi đáp ứng, và điều thú vị nhất là \emph{dùng gì} để dạy nó.

\emph{Người thầy im lặng.} Mạng được kích thích qua một điện cực, cứ một đến ba giây một lần, cho đến khi ở điện cực khác xuất hiện đáp ứng mong muốn sau kích thích $50\pm10\ms$. Ngay khi đáp ứng ấy xuất hiện, kích thích bị tắt trong năm phút. Sau vài chu kỳ như vậy, đáp ứng mong muốn xuất hiện ngay khi có kích thích~\cite{Shahaf2001}. Phần thưởng ở đây là sự im lặng: kích thích làm mạng lung lay, còn việc ngừng kích thích giữ nó lại ở trạng thái đã cho đáp ứng đúng. Đây là phép xuống dốc gradient qua các nhiễu loạn ngẫu nhiên ở chương~\ref{sec:dofa} (mệnh đề~\ref{prop:gradient}), trong đó vai trò của sai số do chính việc rung lắc đảm nhận: cứ lắc cho đến khi đúng thì thôi.

\emph{Mạng chơi quần vợt.} Trong một thí nghiệm trên bàn thử có mảng điện cực dày, khoảng một triệu nơ-ron được nối với một trò chơi quần vợt trên máy tính có một cây vợt~\cite{Kagan2022}. Vị trí quả bóng được đưa vào bằng dòng điện qua các điện cực ``cảm giác'': bóng ở đâu thì mã hóa theo vị trí, bóng xa bao nhiêu thì mã hóa theo tần số. Hoạt động ở hai vùng ``vận động'' đẩy vợt lên hoặc xuống. Quy tắc học rất khác thường. Nếu vợt đỡ trúng bóng, mạng nhận một kích thích ngắn \emph{đoán trước được}; nếu trượt, nó nhận vài giây dòng điện ngẫu nhiên \emph{không đoán trước được}. Sau hai mươi phút chơi, các chuỗi cú đỡ trúng dài ra. Mức tăng có ý nghĩa thống kê trong một phiên chỉ có khi phản hồi đầy đủ; khi thay kích thích bằng im lặng, cuối phiên mẫu nuôi cấy cũng chơi tốt hơn so với không có phản hồi, nhưng hiệu ứng nhỏ hơn, còn việc học bền vững từ phiên này sang phiên khác qua nhiều ngày thì không có. Phân tích chi tiết thí nghiệm này nằm ở chương~\ref{sec:dishbrain}.

Các tác giả giải thích kết quả bằng giả thuyết rằng mạng hành động sao cho đầu vào của nó dễ đoán~\cite{Friston2010}. Đây cũng là ý tưởng tiết kiệm xung ở chương~\ref{sec:code} và dự đoán khung hình ở chương~\ref{sec:recognition}: cỗ máy tốn ít hơn khi đầu vào đúng như chờ đợi. Nhưng cả thí nghiệm, và nhất là lời các tác giả về sự ``có tri giác'' của mẫu nuôi cấy, đã bị phê phán gay gắt: hiệu ứng nhỏ, và có thể giải thích đơn giản hơn~\cite{Balci2023}. Đánh giá công bằng là thế này: mẫu nuôi cấy trên bàn thử thay đổi hành vi dưới tác động của phản hồi, điều này đã đo được; còn việc nó học chính là để dự đoán đầu vào của mình thì là giả thuyết.

\emph{Mạng điều khiển cơ thể.} Theo cách tương tự, mẫu nuôi cấy từng được nối với một ``cơ thể'' ảo đơn giản và được dạy di chuyển tới mục tiêu, bằng cách chọn mẫu kích thích huấn luyện theo không gian và thời gian~\cite{Bakkum2008}. Không thí nghiệm nào trong số này có dopamine: thầy dạy là mẫu dòng điện do kỹ sư chọn. Điều gì thay đổi khi thầy dạy là một chương trình hoặc chính dopamine được bàn ở các mục~\ref{sec:dishdopamine}--\ref{sec:cartpole}.

\begin{figure}[t]
\centering
\begin{tikzpicture}[>=Stealth,
  blk/.style={draw,very thick,rounded corners=2pt,align=center,font=\footnotesize,minimum height=0.8cm,fill=blue!5}]
% (a) closed loop
\node[font=\small] at (1.5,4.3) {(a) vòng kín};
\node[blk,text width=2.6cm] (G) at (1.5,3.3) {trò chơi: bóng và vợt};
\draw[very thick,fill=orange!10,rounded corners=3pt] (0.5,0.2) rectangle (2.5,1.6);
\foreach \x in {0.7,1.0,...,2.35}{\foreach \y in {0.4,0.7,1.0,1.3}{\fill[gray!60] (\x,\y) circle (0.04);}}
\draw[->,thick,blue!60!black] (G.west) -| (-0.5,0.9) -- (0.5,0.9);
\node[font=\scriptsize,blue!60!black,anchor=east] at (-0.6,2.1) {bóng $\to$ điện};
\draw[->,thick,red!70!black] (2.5,0.9) -- (3.5,0.9) |- (G.east);
\node[font=\scriptsize,red!70!black,anchor=west] at (3.6,2.1) {xung $\to$ vợt};
\node[font=\scriptsize] at (1.5,-0.2) {nơ-ron trên mảng điện cực};
\node[font=\scriptsize,align=center] at (1.5,-0.85) {trúng: kích thích đoán trước được\\trượt: dòng điện ngẫu nhiên};
% (b) reservoir
\begin{scope}[xshift=7.9cm]
\node[font=\small] at (1.6,4.3) {(b) hồ chứa};
\node[blk,text width=1.1cm] (X) at (-0.4,1.0) {vào\\$u(t)$};
\draw[very thick,fill=orange!10] (1.6,1.0) circle (0.8);
\foreach \a/\r in {20/0.35,80/0.5,150/0.3,210/0.55,280/0.4,330/0.2,110/0.15}{\fill[red!60!black] ({1.6+\r*cos(\a)},{1.0+\r*sin(\a)}) circle (0.05);}
\node[font=\scriptsize] at (1.6,-0.2) {organoid: không có thầy};
\node[blk,text width=1.8cm,fill=green!8] (W) at (4.0,1.0) {lớp đọc ra\\$W_{\text{ra}}$};
\draw[->,thick] (X) -- (0.8,1.0);
\draw[->,thick] (2.4,1.0) -- node[above,font=\scriptsize] {$\mathbf r(t)$} (W);
\draw[->,thick] (W) -- (4.0,2.3) node[above,font=\scriptsize] {$y(t)$};
\node[font=\scriptsize,green!40!black] at (4.0,0.3) {học có thầy};
\end{scope}
\end{tikzpicture}
\caption{Hai cách buộc mạng sống tính toán. (a) Vòng kín: vị trí bóng được đưa vào bằng dòng điện, xung của các vùng vận động đẩy vợt, còn kích thích đoán trước được hoặc ngẫu nhiên sau cú đánh đóng vai thầy dạy. (b) Hồ chứa~\eqref{eq:reservoir}: organoid không nhận tín hiệu dạy nào, nó tách đầu vào thành nhiều đáp ứng phi tuyến có trí nhớ; chỉ một lớp đọc ra tuyến tính bên ngoài học theo sai số. Trong lúc đó, liên kết của chính organoid cũng thay đổi, không cần thầy.}
\label{fig:livingmachine}
\end{figure}

\section{Hồ chứa sống}

Còn một cách khác để dùng mạng sống: hoàn toàn không dạy nó. Trong kỹ thuật, cách này gọi là \emph{tính toán hồ chứa}~\cite{Maass2002,JaegerHaas2004}. Lấy một hệ phi tuyến có sẵn và có trí nhớ, hệ nào cũng được, miễn là đáp ứng của nó với đầu vào phong phú và phụ thuộc vào quá khứ gần. Đầu vào $u(t)$ kích động hệ, tạo ra vectơ đáp ứng $\mathbf r(t)$, và chỉ có lớp đọc ra tuyến tính là được huấn luyện:
\begin{empheq}[box=\keybox]{equation}
y(t) \;=\; W_{\text{ra}}\,\mathbf r(t) ,
\label{eq:reservoir}
\end{empheq}
với trọng số được chọn theo quy tắc bình phương tối thiểu~\eqref{eq:lms} ở chương~\ref{sec:motorlearning}. Hồ chứa làm đúng việc mà các nơ-ron \nt{GLUT} nhỏ ở chương~\ref{sec:motorlearning} đã làm: tách đầu vào thành một vectơ dài, trong đó các lớp cần phân biệt trở nên tách được bằng một mặt phẳng (chương~\ref{sec:dendrites}).

Một organoid trên mảng điện cực tỏ ra dùng được làm hồ chứa như vậy. Âm thanh tiếng nói, đưa vào organoid dưới dạng các mẫu dòng điện, sau khi huấn luyện một lớp đọc ra tuyến tính đã được phân biệt theo người nói với độ chính xác khoảng $78\%$~\cite{Cai2023}. Độ chính xác $78\%$ là con số do chính các tác giả công bố và báo chí nhắc lại. Kết quả hữu ích, nhưng cần hiểu ``trí khôn'' nằm ở đâu: organoid cung cấp tính phi tuyến và trí nhớ tắt dần, còn việc học theo sai số được làm bên ngoài, trên silicon (hình~\ref{fig:livingmachine}). Tuy vậy, organoid không đứng yên: theo các tác giả, trong quá trình huấn luyện, độ liên kết chức năng của chính nó thay đổi, và họ gọi đó là học không giám sát bên trong organoid~\cite{Cai2023}. Phần độ chính xác nào do thay đổi này, phần nào do lớp đọc ra, thì chưa được tách bạch.

\section{Bàn thử nói gì về cỗ máy}

\emph{Năng lượng.} Theo ước tính~\eqref{eq:energy}, một xung tốn khoảng $10^{-10}$~J. Một mẫu nuôi cấy một triệu nơ-ron, với tần số một xung mỗi giây, tiêu tốn cho việc truyền tín hiệu
\begin{empheq}[box=\keybox]{equation}
P \;\approx\; 10^{6}\times1\,\text{s}^{-1}\times10^{-10}\,\text{J} \;=\; 10^{-4}\,\text{W},
\label{eq:dishpower}
\end{empheq}
tức một phần mười miliwatt. Phần điện tử kích thích, ghi và xử lý các xung này tốn nhiều hơn vài bậc độ lớn. Cũng như ở chương~\ref{sec:debugging}, nút thắt cổ chai không phải là tính toán, mà là việc kết nối với nó.

\emph{Những linh kiện còn thiếu.} Bàn thử cho thấy bus \nt{GLUT} cùng điều khiển luồng \nt{GABA}, khi không có thanh ghi điều khiển, làm được bao nhiêu: tự tổ chức, tạo ra các đợt bùng phát, thay đổi đáp ứng dưới tác động của phản hồi và làm một hồ chứa tốt. Điều nó không làm được khi thiếu chúng là tự quyết định thế nào là thành công. Trên bàn thử, tín hiệu ấy do kỹ sư đưa vào, còn trong não, như ta đã thấy ở các chương~\ref{sec:dofa}--\ref{sec:acet}, là do các kênh quảng bá.

\emph{Đạo đức.} Organoid được nuôi từ tế bào người, và chúng càng phức tạp thì câu hỏi liệu mô như vậy có cảm nhận được gì không càng nghiêm túc. Góc nhìn kỹ thuật gợi ý một phần câu trả lời: đau ở chương~\ref{sec:pain} không phải là tín hiệu của một cảm biến, mà là cả một hệ thống báo động có cổng, núm chỉnh âm lượng và các kênh quảng bá, những thứ bàn thử không có. Nhưng đây là lập luận chứ không phải chứng minh, và chính lĩnh vực này đề xuất tiến hành đánh giá đạo đức song song với thí nghiệm, ngay từ đầu~\cite{Smirnova2023}.

\begin{keyblock}
\begin{proposition}[Mạng sống trên bàn thử]
\label{prop:livingmachine}
Một mạng gồm các nơ-ron \nt{GLUT} và \nt{GABA} trên mảng điện cực tự tạo ra các đợt bùng phát~\eqref{eq:burstinterval}, thay đổi đáp ứng dưới tác động của phản hồi và có thể làm hồ chứa~\eqref{eq:reservoir} cho một lớp đọc ra được huấn luyện bên ngoài. Tất cả những điều này đã được đo. Việc mạng như vậy học theo một quy tắc chung nào đó, chẳng hạn dự đoán đầu vào của mình, là giả thuyết, còn những lời lớn về sự ``có tri giác'' của nó thì phần đông chuyên gia không chấp nhận.
\end{proposition}
\end{keyblock}

\section{Dopamine theo chớp sáng}
\label{sec:dishdopamine}

Thí nghiệm trực tiếp duy nhất mà ta biết, trong đó mẫu nuôi cấy được dùng dopamine làm thầy dạy, được thực hiện trên một nền tảng organoid mà các nhà nghiên cứu kết nối từ xa~\cite{Jordan2024}. Vào dung dịch bao quanh organoid, người ta thêm dopamine nằm trong một ``lồng'' nhạy sáng: phân tử bị nhốt không tác động lên thụ thể. Nồng độ dopamine trong lồng là $1$~mM. Khi cần thưởng cho mạng, người ta chiếu tia cực tím: lồng vỡ ra, và dopamine thoát vào thể tích.

Vòng lặp được bố trí như sau. Cùng một kích thích được đưa vào lặp đi lặp lại; nếu nó gây ra nhiều xung hơn trước, đèn chớp bật lên và dopamine được giải phóng. Qua $240$ lần lặp, số xung được gây ra nhìn chung tăng lên. Chính các tác giả viết rằng chỉ từ một thí nghiệm như vậy thì không thể kết luận mức tăng là do dopamine~\cite{Jordan2024}.

Với kỹ sư, đây là lần thử đầu tiên lắp quy tắc ba thừa số~\eqref{eq:threefactor} trong đĩa nuôi: hoạt động của bên gửi và bên nhận để lại vết, còn tín hiệu chung quyết định có củng cố thay đổi hay không. Nhưng tín hiệu ở đây chỉ có thể dương: có thưởng hoặc không có, còn sai số âm $\delta<0$ thì thiết bị không tạo ra được. Hơn nữa, dopamine lan ra và được thu hồi chậm, giống nồng độ trong~\eqref{eq:lowpass}, nên các chớp sáng dày có thể chỉ đơn giản nâng mức chung của nó lên. Để phân biệt việc học với hiện tượng này, cần hai đối chứng: cùng số chớp sáng ở những thời điểm ngẫu nhiên, không gắn với đáp ứng của mạng, và cùng những chớp sáng ấy nhưng không có dopamine trong lồng. Cũng cần kiểm tra sau khi nghỉ: thay đổi có còn lại khi dopamine đã hết hay không.

\section{Dopamine dạy silicon}
\label{sec:biohybrid}

Trong thí nghiệm ngược lại, tế bào sống dạy thiết bị điện tử~\cite{Keene2020}. Các tế bào tiết dopamine được nuôi trong một vi kênh phía trên một phần tử neuromorphic hữu cơ, tức một transistor mà độ dẫn của kênh đóng vai trò trọng số khớp thần kinh. Dopamine thoát ra từ tế bào bị oxy hóa ở điện cực của phần tử, và điều đó làm thay đổi độ dẫn. Thiết bị còn mô phỏng cả cơ chế dọn chất dẫn truyền khỏi khe khớp, nên trọng số không chỉ thay đổi được lâu dài mà còn đưa về lại được.

Về mặt mạch, đây là một bộ tích phân rò với đầu vào hóa học: trọng số tích lũy dòng dopamine, còn việc ``thu hồi'' quyết định nó bị quên nhanh đến đâu, cũng là mạch RC như~\eqref{eq:lowpass}. Điều quan trọng ở đây là chiều kết nối. Đầu dò ở chương~\ref{sec:debugging} không thấy các thanh ghi quảng bá, vì chúng là hóa học. Phần tử này đọc đúng một thanh ghi hóa học và biến nó thành trọng số điện. Với giao diện não--máy tính, đây là con đường tới một cảm biến nghe được không chỉ bus dữ liệu mà cả các thanh ghi điều khiển.

\section{Organoid có nơ-ron dopamine riêng}
\label{sec:midbrainorg}

Từ tế bào gốc người, ta đã biết cách nuôi những organoid giống vùng não nơi có các nơ-ron \nt{DOPA}. Trong đó có những nơ-ron dopamine hoạt động và tiết ra dopamine~\cite{Jo2016}. Những organoid như vậy được nuôi chủ yếu để nghiên cứu bệnh. Chúng tôi không tìm thấy thí nghiệm nào dùng dopamine của chúng làm thầy dạy cho một mạng khác.

Với cỗ máy từ nơ-ron sống, đây chính là linh kiện còn thiếu. Bàn thử ở chương~\ref{sec:livingmachine} là bus dữ liệu và điều khiển luồng không có thanh ghi điều khiển; ở đây nguồn tín hiệu quảng bá đã có. Cái còn thiếu là bộ phê bình: một mạng tính sai số dự đoán~\eqref{eq:ctd} và điều khiển các nơ-ron này. Nối một organoid vỏ não với một organoid như vậy sao cho dopamine được tiết ra theo sai số là một bài toán mở, và hiện nó vẫn là giả thuyết, chưa phải thí nghiệm.

\section{Organoid giữ thăng bằng cây sào không cần dopamine}
\label{sec:cartpole}

Thí nghiệm đầy đủ nhất trong số các thí nghiệm gần đây hoàn toàn không dùng dopamine~\cite{Robbins2026}. Organoid vỏ não chuột được nối vào vòng kín với bài toán ``cây sào trên xe đẩy'': hoạt động của organoid đẩy xe, còn vị trí cây sào được đưa trở lại bằng kích thích. Giữa các lần thử, organoid nhận những kích thích huấn luyện ngắn, tần số cao. Chọn kích thích nào là việc của một chương trình học tăng cường.

Các kết quả đã được đo:
\begin{itemize}
\item với phần lớn organoid, tín hiệu huấn luyện do chương trình chọn cho kết quả tốt hơn so với tín hiệu chọn ngẫu nhiên hoặc không có tín hiệu;
\item sau $45$ phút nghỉ, sự cải thiện không còn;
\item chặn cả hai loại thụ thể glutamate, AMPA và NMDA, làm mất sự cải thiện.
\end{itemize}

Điểm cuối cho biết cái được học nằm ở đâu: trong các khớp thần kinh \nt{GLUT}, và qua chính thụ thể mà ở mục~\ref{sec:weightstore} hoạt động như bộ phát hiện trùng khớp giữa đầu vào và đầu ra. Điểm thứ hai cho thấy cái còn thiếu: có thay đổi nhanh, nhưng không có củng cố, giống người học nhanh ở chương~\ref{sec:motorlearning} mà thiếu người học chậm. Còn vai trò thầy dạy đã đổi: chương trình đã thay kỹ sư chọn mẫu dòng điện, nhưng thầy vẫn đứng ngoài đĩa nuôi.

Trong các thí nghiệm sớm hơn về việc dạy mẫu nuôi cấy trên mảng điện cực, chúng tôi cũng không tìm thấy thí nghiệm nào dùng dopamine: tín hiệu dạy ở mọi nơi đều là kích thích điện.

\section{Ai dạy mẫu nuôi cấy}

\begin{table}[t]
\centering
\footnotesize
\setlength{\tabcolsep}{4pt}
\renewcommand{\arraystretch}{1.15}
\begin{tabular}{>{\raggedright}p{2.7cm}>{\raggedright}p{2.5cm}>{\raggedright}p{3.2cm}>{\raggedright\arraybackslash}p{4.4cm}}
\toprule
Thí nghiệm & Ai dạy & Tín hiệu dạy & Điều đã biết \\
\midrule
ngừng kích thích khi có đáp ứng mong muốn & kỹ sư & kết thúc kích thích & đáp ứng được củng cố --- đã đo \\
DishBrain (ch.~\ref{sec:dishbrain}) & kỹ sư & dòng điện đoán trước được hoặc ngẫu nhiên & mức tăng có ý nghĩa của lượt đánh trong một phiên chỉ khi phản hồi đầy đủ, với im lặng thì hiệu ứng nhỏ hơn --- đã đo; không bền vững từ phiên này sang phiên khác; nguyên nhân còn tranh luận \\
cây sào trên xe đẩy & chương trình học tăng cường & kích thích tần số cao do chương trình chọn & tốt hơn kích thích ngẫu nhiên nhưng không giữ được sau khi nghỉ --- đã đo \\
dopamine theo chớp sáng & quy tắc ``nhiều xung hơn --- có thưởng'' & dopamine được ánh sáng giải phóng & số xung tăng --- quan sát; vai trò của dopamine chưa được chứng minh \\
khớp thần kinh lai sinh học & tế bào sống & dopamine từ tế bào & thay đổi lâu dài và đưa trọng số của phần tử về lại --- đã đo \\
organoid có nơ-ron \nt{DOPA} & --- & --- & có nguồn dopamine; chưa được dùng làm thầy dạy \\
\bottomrule
\end{tabular}
\caption{Ai dạy nơ-ron sống trên chip, và dạy bằng gì.}
\label{tab:dishteachers}
\end{table}

Trong mọi thí nghiệm mà chính mẫu nuôi cấy học, thầy dạy cho đến nay vẫn ở bên ngoài (bảng~\ref{tab:dishteachers}): kỹ sư, chương trình, hoặc một quy tắc do kỹ sư đặt ra. Dopamine mới vào đĩa nuôi đúng một lần, và vai trò của nó vẫn còn phải chứng minh. Lắp trong đĩa nuôi một kênh quảng bá thực thụ, có bộ phê bình, có sai số và có vết như ở chương~\ref{sec:dofa}, là bước tiếp theo mà chưa ai làm được.


\section{Tổng kết}

\begin{table}[t]
\centering
\footnotesize
\setlength{\tabcolsep}{4pt}
\renewcommand{\arraystretch}{1.15}
\begin{tabular}{>{\raggedright}p{2.8cm}>{\raggedright}p{4.2cm}>{\raggedright\arraybackslash}p{5.8cm}}
\toprule
Thí nghiệm & Mạng sống làm gì & Điều đã biết \\
\midrule
mạng không có đầu vào & đợt bùng phát cách nhau từ một giây đến vài phút~\eqref{eq:burstinterval} & đã đo; cơ chế qua sự cạn kiệt khớp thần kinh là mô hình được chấp nhận rộng rãi \\
người thầy im lặng & đáp ứng mong muốn được củng cố khi tắt kích thích & đã đo \\
chơi quần vợt & chuỗi cú đỡ trúng dài ra; mức tăng có ý nghĩa trong một phiên chỉ có khi phản hồi đầy đủ & thay đổi hành vi đã đo; việc học từ phiên này sang phiên khác không bền vững; độ lớn hiệu ứng và cách giải thích bằng dự đoán đầu vào còn gây tranh cãi \\
cơ thể ảo & di chuyển tới mục tiêu khi chọn đúng mẫu kích thích & đã đo \\
hồ chứa & tính phi tuyến và trí nhớ cho lớp đọc ra~\eqref{eq:reservoir} & đã đo; việc học theo sai số nằm bên ngoài, trên silicon; liên kết của organoid cũng thay đổi \\
năng lượng & khoảng $10^{-4}$~W cho một triệu nơ-ron~\eqref{eq:dishpower} & ước tính theo~\eqref{eq:energy} \\
\bottomrule
\end{tabular}
\caption{Những gì cỗ máy từ nơ-ron sống đã cho thấy.}
\label{tab:livingmachine}
\end{table}

Cỗ máy từ nơ-ron sống (bảng~\ref{tab:livingmachine}) là một bàn thử cho thấy bus dữ liệu và điều khiển luồng làm được gì khi không có thanh ghi điều khiển. Chúng làm được nhiều, nhưng không tự quyết định được thế nào là tốt. Trên bàn thử, vai trò đó thuộc về kỹ sư với mẫu dòng điện của mình; trong não, thuộc về số ít dây dẫn quảng bá mà phần giữa cuốn sách đã dành để bàn.

\section{Bài tập}

\begin{exercise}[\lvlA]
\label{ex:21:1}
\emph{Khoảng cách giữa các đợt bùng phát.} Một đợt làm cạn kho về không, sau đó $\tau_{\text{ph}}\,\dd x/\dd t=1-x$, và đợt mới bắt đầu khi $x=x^*$; $\tau_{\text{ph}}=0.8$~s. (a)~Tìm khoảng cách khi $x^*=0.9$ và $0.99$. (b)~Ngưỡng $x^*$ dao động ngẫu nhiên $\pm0.005$ quanh $0.99$. Các khoảng cách nằm trong dải nào?
\end{exercise}

\begin{exercise}[\lvlA]
\label{ex:21:2}
\emph{Năng lượng của bàn thử.} (a)~Ước tính~\eqref{eq:dishpower} cho công suất bao nhiêu với cả bộ não ($8.6\cdot10^{10}$ nơ-ron)? (b)~Bàn thử ghi $1024$ điện cực ở tần số $20$~kHz, $10$~bit mỗi mẫu. Truyền luồng dữ liệu này với $1$~nJ mỗi bit tốn gấp bao nhiêu lần công suất của chính mẫu nuôi cấy?
\end{exercise}

\begin{exercise}[\lvlB]
\label{ex:21:3}
\emph{Thiết kế: dập tắt các đợt bùng phát.} Kích thích thưa qua nhiều điện cực khiến mỗi khớp thần kinh giải phóng chất dẫn truyền với tần số $r_b$; kho tuân theo $\tau_{\text{ph}}\,\dd x/\dd t=1-x-Ur_b\tau_{\text{ph}}x$, $U=0.5$, $\tau_{\text{ph}}=0.8$~s. Với $r_b$ nhỏ nhất bằng bao nhiêu thì kho không hồi tới $x^*=0.9$; $0.99$, và khi đó khớp yếu đi bao nhiêu?
\end{exercise}

\begin{exercise}[\lvlB]
\label{ex:21:4}
\emph{Trí nhớ của hồ chứa không lớn hơn số nơ-ron.} Một hồ chứa có $N$ đầu ra $\mathbf r(t)$ nhận đầu vào trắng $u(t)$ có trung bình bằng không và phương sai đơn vị; $m_k$ là bình phương tương quan lớn nhất của lớp đọc ra $\mathbf w_k\cdot\mathbf r(t)$ với $u(t-k)$. Chứng minh rằng dung lượng nhớ $\mathrm{MC}=\sum_{k\ge0}m_k\le N$.
\end{exercise}

\begin{exercise}[\lvlB]
\label{ex:21:5}
\emph{Mô phỏng: hồ chứa trên silicon.} Hồ chứa $\mathbf r(t+1)=\tanh(W\mathbf r(t)+\mathbf v\,u(t)+\mathbf c)$, $N=30$, $W$ ngẫu nhiên với bán kính phổ $0.9$, $v_i\sim U(-0.5,0.5)$, $c_i\sim U(-0.2,0.2)$, $u\sim U(-1,1)$, lớp đọc ra là hồi quy ridge. Tìm dung lượng nhớ của bài tập~\ref{ex:21:4} khi có và không có $\tanh$. Hồ chứa nào nhớ nhiều hơn?
\end{exercise}

\begin{exercise}[\lvlB]
\label{ex:21:6}
\emph{Lớp đọc ra cho hồ chứa sống.} Organoid cho $1024$ kênh và $P=240$ bản ghi huấn luyện thuộc hai lớp. (a)~Có thể giữ lại bao nhiêu đặc trưng $n$ để, theo công thức của bài tập~\ref{ex:19:3}, một cách gán nhãn ngẫu nhiên tách được với xác suất không quá $1\%$? (b)~Độ chính xác $78\%$ trên $100$ bản ghi kiểm tra cao hơn mức đoán mò bao nhiêu sigma?
\end{exercise}

\begin{exercise}[\lvlC]
\label{ex:21:7}
\emph{Nghiên cứu: thầy dạy không có kênh quảng bá.} Hãy đề xuất một giao thức kích thích hiện thực quy tắc ba thừa số của chương~\ref{sec:dofa} trên bàn thử qua $8$--$1024$ điện cực, và một đối chứng với lớp đọc ra bị đóng băng, phân biệt việc học trọng số với sự dịch chuyển trạng thái của mạng. Kết quả nào sẽ bác bỏ giả thuyết về việc học?
\end{exercise}

\chapter{DishBrain}
\chaptermark{DishBrain}
\label{sec:dishbrain}

\section{Mạng trong đĩa nuôi chơi quần vợt}

DishBrain là tên các tác giả đặt cho bàn thử, trên đó một mẫu nuôi cấy nơ-ron sống, mọc trên con chip có điện cực, chơi một trò quần vợt đơn giản với một cây vợt~\cite{Kagan2022}. Đây là thí nghiệm được bàn luận nhiều nhất với cỗ máy từ nơ-ron sống (tổng quan về những cỗ máy như vậy nằm ở chương~\ref{sec:livingmachine}), và đặc biệt tiện cho cuốn sách này. Ở đây, cả mạng nhỏ đều nằm trong tầm tay: mọi đầu vào do người làm thí nghiệm đặt, mọi đầu ra đều được ghi. Mọi câu hỏi của các chương trước, như tín hiệu được mã hóa ra sao, ai dạy, đầu dò thấy gì, đều có thể đặt cho một mạng không có mắt, không có cơ và không có nơ-ron \nt{DOPA}. Ta sẽ mổ xẻ thí nghiệm như kỹ sư mổ xẻ bàn thử của người khác: sơ đồ, đầu vào, đầu ra, thầy dạy, kết quả, tranh cãi.

\section{Bàn thử}

Nơ-ron được lấy từ vỏ não phôi chuột, khoảng $800\,000$ tế bào mỗi chip, hoặc được nuôi từ tế bào gốc người, khoảng một triệu mỗi chip. Trong vài tuần, chúng lớn lên trên chip, mọc đan vào nhau và bắt đầu tự phát ra xung và chùm xung. Trên một diện tích $8$~mm$^2$ bên dưới chúng có $26\,000$ điện cực bạch kim; có thể ghi đồng thời tối đa $1024$ điện cực, còn kích thích trên thực tế thì qua tám điện cực.

Quanh mẫu nuôi cấy là một vòng kín (hình~\ref{fig:dishloop}). Máy tính mô phỏng trò chơi: bóng bay, dội vào tường, vợt di chuyển lên xuống. Vị trí bóng được chuyển thành dòng điện trên tám điện cực ``cảm giác''. Xung từ hai vùng ``vận động'' được đếm theo cửa sổ $10\ms$ và đẩy vợt. Sau mỗi cú đỡ trúng hoặc trượt, mẫu nuôi cấy nhận thêm một kích thích nữa, tức phản hồi. Cửa sổ $10\ms$ là nhịp của máy tính trong bàn thử, không phải của mẫu nuôi cấy: bản thân mạng vẫn làm việc bất đồng bộ, như mọi mạng trong cuốn sách này.

\begin{figure}[t]
\centering
\begin{tikzpicture}[>=Stealth,
  blk/.style={draw,very thick,rounded corners=3pt,align=center,font=\small,minimum height=1.2cm},
  lab/.style={font=\scriptsize,align=center}]
\node[blk,fill=blue!6,text width=3.4cm] (C) at (0,0) {mẫu nuôi cấy trên chip\\\scriptsize $\sim10^6$ nơ-ron, $26\,000$ điện cực};
\node[blk,fill=orange!10,text width=3.0cm] (G) at (7.2,0) {máy tính:\\trò quần vợt};
\draw[->,very thick,blue!60!black] (G.north west) to[bend right=25] node[lab,above] {bóng: \emph{vị trí} --- điện cực nào trong 8,\\\emph{tần số} --- 4--40 xung/s} (C.north east);
\draw[->,very thick,red!70!black] (C.south east) to[bend right=25] node[lab,below] {vợt: số xung của các vùng\\``lên'' và ``xuống'' trong $10\ms$} (G.south west);
\draw[->,thick,densely dashed,green!45!black] (G.east) -- ++(0.9,0) |- ++(0,-2.6) -| node[lab,pos=0.25,below] {trúng: kích thích đoán trước được; trượt: không đoán trước được} (C.south);
\end{tikzpicture}
\caption{Vòng lặp DishBrain. Mũi tên xanh lam là đầu vào: vị trí bóng được ghi bằng mã vị trí và mã tần số (chương~\ref{sec:code}). Mũi tên đỏ là đầu ra: hiệu số đếm của hai vùng đẩy vợt. Đường nét đứt xanh lục là phản hồi sau mỗi lượt đánh, ``thầy dạy'' duy nhất của bàn thử.}
\label{fig:dishloop}
\end{figure}

\section{Đầu vào và đầu ra}

\emph{Đầu vào} được ghi cùng lúc bằng hai mã của chương~\ref{sec:code}. Điện cực nào trong tám điện cực được kích thích cho biết bóng ở độ cao nào: đó là mã vị trí, như ở các cảm biến của chương~\ref{sec:sensors}. Xung kích thích đến dày hay thưa cho biết bóng còn xa bao nhiêu: $4$ xung mỗi giây khi bóng ở bức tường xa, và tăng tuyến tính tới $40$ xung mỗi giây khi bóng ở bức tường phía vợt. Đó là mã tần số; biên độ xung của bóng là $75$~mV. Mẫu nuôi cấy không ``biết'' trước mã nào: mã do người làm thí nghiệm đặt ra, và mạng phải học cách đọc nó.

\emph{Đầu ra} được đọc giống như ở các giao diện của chương~\ref{sec:debugging}. Xung của một vùng vận động đẩy vợt lên, của vùng kia đẩy xuống; ở dạng đơn giản nhất, với mỗi cửa sổ
\begin{empheq}[box=\keybox]{equation}
\Delta p \;=\; c\,\bigl(n_\uparrow - n_\downarrow\bigr),
\label{eq:dishreadout}
\end{empheq}
trong đó $n_\uparrow$, $n_\downarrow$ là số xung của hai vùng trong $10\ms$, còn $c$ là hệ số của bàn thử. Đây là mã hai dây của chương~\ref{sec:glutcomputer}: chiều chuyển động không nằm ở dấu của tín hiệu, mà ở chỗ dây nào trong hai dây ``to'' hơn.

Hãy xem nhiễu của đầu ra như vậy. Nếu một vùng phát tổng cộng $R$ xung mỗi giây, thì trong cửa sổ $T=10\ms$ trung bình có $RT$ xung, và theo~\eqref{eq:rateerror}, độ tản mạn tương đối bằng $1/\sqrt{RT}$. Với $R=1000$ xung mỗi giây, con số này khoảng $30$ phần trăm ở mỗi cửa sổ; với $R=100$ thì hơn một trăm. Vợt chắc chắn sẽ rung, và mạng chỉ có một cách để học: dịch hiệu số đếm \emph{trung bình} về đúng phía. Đây cũng là những định luật giới hạn mọi giao diện ở chương~\ref{sec:debugging}, chỉ khác là giờ chúng tác động ở cả hai phía của vòng lặp.

\section{Thầy dạy không có dopamine}

Điều thú vị nhất ở bàn thử là thầy dạy. Mẫu nuôi cấy nơ-ron vỏ não không có nơ-ron \nt{DOPA}, và bàn thử không gửi tín hiệu ``tốt hơn mong đợi'' nào cả. Phản hồi thì khác. Mạng đỡ trúng bóng thì cả tám điện cực cảm giác cùng nhận một kích thích ngắn \emph{đoán trước được}: xung $75$~mV, $100$ xung mỗi giây trong $0.1\,\text{s}$. Mạng đỡ trượt thì trong bốn giây có một kích thích mà các tác giả gọi là \emph{không đoán trước được}: xung $150$~mV, $5$ xung mỗi giây, vào những thời điểm ngẫu nhiên, trên các điện cực chọn ngẫu nhiên; sau đó là bốn giây nghỉ không kích thích, rồi bóng xuất phát lại~\cite{Kagan2022}.

Các tác giả xuất phát từ giả thuyết rằng mạng hành động sao cho đầu vào của nó trở nên dễ đoán~\cite{Friston2010}. Với kỹ sư, ý nghĩa rất đơn giản: mẫu nuôi cấy có đúng một cách để thoát khỏi bốn giây nhiễu ngẫu nhiên, đó là đỡ trúng bóng. Ta đã gặp ý tưởng này ở việc tiết kiệm xung trong chương~\ref{sec:code} và ở dự đoán khung hình trong chương~\ref{sec:recognition}.

Nhưng có thật cần đến tính dễ đoán ở đây không? Cũng có thể có một cơ chế thô hơn. Giả sử kích thích không đoán trước được sau cú trượt chỉ đơn giản \emph{xáo trộn} những thay đổi gần đây trong mạng, còn kích thích êm sau cú trúng thì không đụng đến chúng. Hãy mô tả cấu hình của mạng bằng một vectơ $\boldsymbol\psi$, và giả sử ở cấu hình $\boldsymbol\psi$ mạng đỡ trượt với xác suất $P_{\text{trượt}}(\boldsymbol\psi)$. Nếu các cú hích là đối xứng, tức từ $\boldsymbol\psi$ sang $\boldsymbol\psi'$ có xác suất bằng chiều ngược lại, thì ở trạng thái dừng, dòng rời khỏi mỗi cấu hình cân bằng với dòng đi vào, và tỷ lệ thời gian mạng ở trong cấu hình $\boldsymbol\psi$ bằng
\begin{empheq}[box=\keybox]{equation}
\rho(\boldsymbol\psi) \;\propto\; \frac{1}{P_{\text{trượt}}(\boldsymbol\psi)} .
\label{eq:dishdensity}
\end{empheq}
Cấu hình trượt ít hơn mười lần thì được giữ lâu hơn mười lần. Không ai báo cho mạng biết phải đổi trọng số theo hướng nào: các cấu hình tốt chỉ đơn giản là ít bị phá hơn.

Chúng tôi đã kiểm tra điều này trên một mô hình, rút gọn mẫu nuôi cấy về hai con số: vợt đến điểm $a\,y+b$ khi bóng đến độ cao $y$, và đỡ trúng nếu lệch ít hơn $0.1$. Sau cú trượt, cả hai số nhận một cú hích ngẫu nhiên, sau cú trúng thì giữ nguyên. Qua $2000$ lượt đánh, tỷ lệ bóng đỡ trúng tăng từ $0.21$ ở hai trăm lượt đầu lên $0.38$ ở năm trăm lượt cuối (bốn mươi ``mẫu nuôi cấy'', hình~\ref{fig:dishmodel}). Nếu cú hích đến sau \emph{mỗi} lượt đánh, phản hồi không mang thông tin gì, và tỷ lệ ở quanh $0.12$; nếu hoàn toàn không có cú hích, nó đứng ở $0.19$. Thí nghiệm trên bàn thử phù hợp với điều này: mức tăng có ý nghĩa trong một phiên chỉ có khi phản hồi đầy đủ. Khi thay kích thích bằng im lặng, cuối phiên mẫu nuôi cấy vẫn chơi tốt hơn so với không có phản hồi, nhưng hiệu ứng nhỏ hơn~\cite{Kagan2022}; lưu ý rằng ngay trong điều kiện này, sau cú trượt là một khoảng nghỉ dài, còn sau cú trúng là khoảng nghỉ ngắn, nên sự khác biệt giữa hai kết cục không mất hẳn.

\begin{figure}[t]
\centering
\begin{tikzpicture}[x=1cm,y=5cm]
\draw[->,gray!70!black] (0,0) -- (10.6,0);
\draw[->,gray!70!black] (0,0) -- (0,0.5) node[above,font=\small,black] {tỷ lệ đỡ trúng};
\foreach \v in {0.1,0.2,0.3,0.4}{\draw[gray!60!black] (-0.06,\v) -- (0.06,\v); \node[left,font=\scriptsize] at (0,\v) {\v};}
\foreach \x/\f/\l/\lab in {1.8/0.21/0.38/{nhiễu sau\\cú trượt},5.2/0.13/0.11/{nhiễu sau\\mỗi lượt đánh},8.6/0.19/0.19/{không nhiễu}}{
  \fill[blue!35] (\x-0.8,0) rectangle (\x-0.05,\f);
  \fill[red!60!black] (\x+0.05,0) rectangle (\x+0.8,\l);
  \node[below,font=\scriptsize,align=center] at (\x,-0.01) {\lab};
  \node[above,font=\scriptsize] at (\x-0.425,\f) {\f};
  \node[above,font=\scriptsize] at (\x+0.425,\l) {\l};}
\fill[blue!35] (7.4,0.47) rectangle (7.7,0.495); \node[right,font=\scriptsize] at (7.75,0.482) {200 lượt đầu};
\fill[red!60!black] (7.4,0.41) rectangle (7.7,0.435); \node[right,font=\scriptsize] at (7.75,0.422) {500 lượt cuối};
\end{tikzpicture}
\caption{Mô hình ``xáo trộn khi trượt'' (\texttt{dishbrain\_model.py}): tỷ lệ bóng đỡ trúng ở đầu và cuối $2000$ lượt đánh, trung bình trên bốn mươi mô hình mẫu nuôi cấy. Việc học chỉ diễn ra nếu cú hích đi sau cú trượt chứ không đi sau cú trúng, cũng như trong thí nghiệm, nơi mức tăng có ý nghĩa trong một phiên chỉ có khi phản hồi đầy đủ.}
\label{fig:dishmodel}
\end{figure}

\begin{keyblock}
\begin{proposition}[Thầy dạy bằng cách xáo trộn]
\label{prop:dishscramble}
Nếu thất bại làm mạng bị lắc, còn thành công để yên cho nó, thì mạng tự trôi về những cấu hình ít sai hơn: thời gian ở trong một cấu hình tỷ lệ nghịch với xác suất thất bại~\eqref{eq:dishdensity}. Việc học như vậy không cần tín hiệu thưởng, không cần dấu của nó, cũng không cần tính dễ đoán tự thân; chỉ cần phản hồi sau thành công và sau thất bại khác nhau. Sự thay đổi hành vi của mẫu nuôi cấy đã được đo; còn cơ chế nào hoạt động trong đĩa nuôi, dự đoán đầu vào hay xáo trộn, thì chưa được xác lập.
\end{proposition}
\end{keyblock}

Hãy so sánh với chương~\ref{sec:dofa}. Ở đó nhiễu cũng dùng để tìm kiếm, nhưng dopamine có dấu: sai số $\delta$ cho biết phải dịch trọng số về phía nào, và bước đi theo gradient~\eqref{eq:perturbation}. Ở đây không có dấu, chỉ có ``giữ lại'' hoặc ``lắc''. Cách tìm kiếm này thô hơn và chậm hơn, nhưng đòi hỏi ở mạng ít hơn: không cần nơ-ron \nt{DOPA}, không cần bộ phê bình, không cần vết kéo dài hàng giây.

\section{Đã đo được gì và tranh cãi về điều gì}

Mỗi ván chơi kéo dài $20$ phút; người ta so sánh $5$ phút đầu với $15$ phút cuối. Ở các mẫu nuôi cấy có phản hồi đầy đủ, chuỗi cú đỡ trúng dài ra, còn số lần trượt ngay lập tức giảm đi; ở các đối chứng không có tế bào, không có trò chơi và với ``vợt'' do máy tính điều khiển thì không thấy gì tương tự. Mẫu nuôi cấy từ tế bào người trung bình đỡ lâu hơn mẫu của chuột. Sự cải thiện đáng tin cậy về thống kê, trên chín mẫu chuột và mười một mẫu người, với hơn một trăm ván cho mỗi nhóm, nhưng không lớn: mạng không trở thành người chơi giỏi, nó chỉ tốt hơn ngẫu nhiên một chút. Sự cải thiện đáng tin cậy trong một phiên; việc học bền vững từ phiên này sang phiên khác qua vài ngày thì các tác giả không đạt được~\cite{Kagan2022}. Trong công trình tiếp theo của cùng nhóm, người ta thấy khi chơi, hoạt động của mẫu nuôi cấy tiến gần tới chế độ \emph{tới hạn}, tức ranh giới giữa tắt dần và trận thác, chính là kiểu sống ở bờ vực trong chương~\ref{sec:gaba}, và càng gần thì chơi càng tốt~\cite{Habibollahi2023}.

Tranh cãi chính không nằm ở con số, mà ở lời lẽ. Tiêu đề bài báo khẳng định rằng nơ-ron trong đĩa nuôi ``thể hiện tri giác'' (sentience), và một nhóm nhà nghiên cứu đã phản bác: thay đổi hành vi trong vòng kín chưa phải là trí tuệ hay cảm giác, và kết luận cần được phát biểu khiêm tốn hơn~\cite{Balci2023}. Từ góc nhìn kỹ thuật, còn hai câu hỏi mở, và cả hai đều về cơ chế. Thứ nhất: mạng có học không, tức trọng số của nó có thay đổi lâu dài không, hay phản hồi chỉ dịch chuyển trạng thái hiện tại của nó? Thứ hai: có thật cần đến tính dễ đoán, hay chỉ cần bất kỳ sự khác biệt nào giữa đáp ứng với thành công và với thất bại, như trong mệnh đề~\ref{prop:dishscramble}? Một bàn thử mà mọi điện cực đều trong tầm tay cho phép trả lời cả hai, và có lẽ đó là giá trị chính của nó.

\section{Đặc tả bàn thử: cách dựng lại}
\label{sec:dishspec}

Dưới đây là mọi thứ cần để dựng lại một bàn thử như vậy: thiết bị, vòng lặp, mã hóa đầu vào, đọc đầu ra, phản hồi và giao thức thí nghiệm. Mỗi tham số đều ghi rõ nguồn. ``Bài báo'' nghĩa là giá trị lấy từ văn bản công trình~\cite{Kagan2022}. ``Lựa chọn'' nghĩa là bài báo không có, và để dựng lại thì phải tự đặt; chúng tôi đề xuất giá trị mặc định và nói cách kiểm tra nó.

\subsection*{Thiết bị và mẫu nuôi cấy}

\begin{center}
\footnotesize
\setlength{\tabcolsep}{4pt}
\renewcommand{\arraystretch}{1.15}
\begin{tabular}{>{\raggedright}p{4.0cm}>{\raggedright}p{6.9cm}>{\raggedright\arraybackslash}p{1.6cm}}
\toprule
Tham số & Giá trị & Nguồn \\
\midrule
chip & ma trận MaxOne: $26\,000$ điện cực bạch kim trên diện tích $8$~mm$^2$ & bài báo \\
ghi & tối đa $1024$ kênh đồng thời, $20\,000$ mẫu mỗi giây & bài báo \\
kích thích & về kỹ thuật tối đa $32$ điện cực, trong thí nghiệm là $8$ điện cực độc lập & bài báo \\
tế bào & vỏ não phôi chuột; nơ-ron người từ tế bào gốc (hai cách nuôi); tế bào HEK293T (không phải nơ-ron) làm đối chứng & bài báo \\
cấy & khoảng $10^6$ tế bào mỗi chip & bài báo \\
tuổi mẫu nuôi cấy khi làm thí nghiệm & chuột: từ $\sim10$ ngày; người: $14$--$24$ ngày hoặc $\sim80$ ngày tùy cách nuôi & bài báo \\
\bottomrule
\end{tabular}
\end{center}

\subsection*{Vòng lặp và thời gian}

Vòng lặp chạy theo cửa sổ $10\ms$ ($200$ mẫu): trong mỗi cửa sổ, xung trên các điện cực của vùng vận động được đếm, trò chơi được cập nhật và kích thích tiếp theo được phát ra. Độ trễ từ một xung trong mẫu nuôi cấy đến kích thích đáp lại là khoảng $5\ms$. Khi chơi, tín hiệu của mỗi điện cực đi qua bộ lọc thông cao Bessel bậc 2 với tần số cắt $100$~Hz; giá trị tuyệt đối của tín hiệu được làm trơn bằng bộ lọc thông thấp Bessel bậc 1 ở $1$~Hz, và ngưỡng xung tỷ lệ với giá trị tuyệt đối đã làm trơn này. Ngưỡng sáu độ lệch chuẩn của nền được bài báo nêu cho các phép đo hoạt động riêng lẻ, không phải cho lúc chơi. Để kích thích không bị tính là đáp ứng của mẫu nuôi cấy, mọi tín hiệu bị xóa khi có hơn $15$ xung lớn, trên $75$~mV, đến cùng lúc. Tất cả đều theo bài báo.

\subsection*{Mã hóa quả bóng}

\begin{center}
\footnotesize
\setlength{\tabcolsep}{4pt}
\renewcommand{\arraystretch}{1.15}
\begin{tabular}{>{\raggedright}p{3.4cm}>{\raggedright}p{7.5cm}>{\raggedright\arraybackslash}p{1.6cm}}
\toprule
Tham số & Giá trị & Nguồn \\
\midrule
vùng cảm giác & $8$ điện cực, bố trí sao cho các điện cực kề nhau ứng với các vị trí bóng kề nhau & bài báo \\
mã vị trí & số thứ tự điện cực $k=1,\dots,8$ cho biết vị trí bóng so với tâm vợt & bài báo \\
tọa độ nào & vị trí theo chiều dọc của bóng; để dựng lại: so với vợt, $k=\lceil 8(y-p+\tfrac12)\rceil$ khi $y-p\in[-\tfrac12,\tfrac12]$ & bài báo; công thức là lựa chọn \\
mã tần số & tần số kích thích từ $4$ đến $40$~xung/s mang khoảng cách tới vợt & bài báo \\
chiều của thang & $4$~xung/s khi bóng ở bức tường đối diện, tăng tuyến tính tới $40$~xung/s khi bóng ở bức tường phía vợt: $f=4+36\,(1-x/L)$, $x$ là khoảng cách tới tường phía vợt, $L$ là chiều rộng sân & bài báo \\
dạng xung & xung chữ nhật hai pha ngắn: pha dương trước, pha âm sau & bài báo \\
biên độ & $75$~mV; được chọn để không buộc các tế bào đang bị ức chế phải phát xung & bài báo \\
độ dài các pha & không nêu; lấy thiết lập chuẩn của hệ thống kích thích và giữ nguyên suốt thí nghiệm & lựa chọn \\
\bottomrule
\end{tabular}
\end{center}

\subsection*{Đọc vợt}

Bài báo có hai vùng vận động: xung ở vùng thứ nhất đẩy vợt lên, ở vùng thứ hai đẩy xuống; đóng góp của các vùng được lấy trọng số để tính đến mật độ tế bào và mức hoạt động khác nhau~\cite{Kagan2022}. Công thức chính xác không được đưa ra. Để dựng lại, ta lấy quy tắc~\eqref{eq:dishreadout}, $\Delta p=c\,(n_\uparrow-n_\downarrow)$ cho mỗi cửa sổ $10\ms$, với trọng số cân bằng các vùng theo hoạt động trung bình lúc nghỉ (lựa chọn): $n_\uparrow$ và $n_\downarrow$ được chia cho số đếm trung bình mỗi cửa sổ của vùng mình, đo trước khi chơi. Hệ số $c$ được chọn sao cho chênh lệch một số đếm trung bình dịch vợt $1\%$ chiều cao sân mỗi cửa sổ (lựa chọn); khi đó, một vùng trội hơn liên tục sẽ đưa vợt qua cả sân trong $1$~s.

Vị trí các vùng trên chip không được cho bằng tọa độ trong văn bản bài báo; chỉ có sơ đồ. Để dựng lại, chỉ cần ba nhóm điện cực không chồng lên nhau: tám điện cực cảm giác ở giữa và mỗi bên một nhóm điện cực ghi, một nhóm ``lên'', một nhóm ``xuống'' (lựa chọn).

\subsection*{Trò chơi}

Kích thước sân, chiều dài vợt và tốc độ bóng không được nêu trong bài báo; chỉ nói rằng bóng bay với tốc độ không đổi và dội lại từ các bức tường. Để dựng lại (lựa chọn): sân $1\times1$, vợt ở tường bên phải dài $0.2$ (trúng khi $|y-p|<0.1$, như trong mô hình~\texttt{models/dishbrain\_model.py}), bóng băng qua sân trong $1$~s. Cú trượt mà trong lượt đánh không có lần đỡ trúng nào được tính là ``ace''; lượt đánh có hơn ba cú đỡ liên tiếp là lượt ``dài'' (bài báo).

\subsection*{Phản hồi}

\begin{center}
\footnotesize
\setlength{\tabcolsep}{4pt}
\renewcommand{\arraystretch}{1.15}
\begin{tabular}{>{\raggedright}p{2.6cm}>{\raggedright}p{8.3cm}>{\raggedright\arraybackslash}p{1.6cm}}
\toprule
Sự kiện & Đưa vào gì & Nguồn \\
\midrule
trúng & kích thích đoán trước được: cả $8$ điện cực cảm giác cùng lúc, $75$~mV, $100$~xung/s, $100\ms$; trong thời gian này, kích thích vị trí bóng thông thường bị ngắt & bài báo \\
trượt & kích thích không đoán trước được: $150$~mV, $5$~xung/s, $4$~s, vào các thời điểm ngẫu nhiên trên các điện cực chọn ngẫu nhiên trong số $8$; sau đó $4$~s không kích thích, và bóng xuất phát theo hướng ngẫu nhiên & bài báo \\
luật ngẫu nhiên & không nêu; để dựng lại: thời điểm xung là dòng Poisson với tần số trung bình $5$~xung/s & lựa chọn \\
\bottomrule
\end{tabular}
\end{center}

\subsection*{Giao thức và đối chứng}

Một phiên kéo dài $20$~phút; so sánh $5$ phút đầu với $15$~phút cuối (bài báo). Ba thước đo kết quả: độ dài trung bình của lượt đánh, tỷ lệ ace và tỷ lệ lượt dài. Các điều kiện thí nghiệm (bài báo):
\begin{itemize}
\item \emph{kích thích}: vòng lặp đầy đủ, như mô tả ở trên;
\item \emph{im lặng}: thay cho kích thích khi trúng hoặc trượt là cùng khoảng thời gian ấy không có kích thích nào, sau đó bóng xuất phát theo hướng ngẫu nhiên;
\item \emph{không phản hồi}: sau cú trượt, bóng chỉ dội lại và bay tiếp, kích thích vị trí bóng vẫn tiếp tục;
\item \emph{nghỉ}: trò chơi vẫn chạy và vợt di chuyển theo hoạt động, nhưng hoàn toàn không có kích thích;
\item \emph{đối chứng}: chip chỉ có môi trường nuôi; tế bào HEK293T; ``mẫu nuôi cấy trong máy tính'', tức mô phỏng thay cho tế bào.
\end{itemize}
Cỡ mẫu trong loạt chính: $9$ mẫu chuột ($101$ phiên), $11$ mẫu người ($138$), $6$ chip chỉ có môi trường ($80$), $20$ mẫu ở trạng thái nghỉ ($42$), $3$ mô phỏng ($38$), tổng cộng $399$ phiên. Trong loạt về phản hồi: $486$ phiên trong $3$ ngày, mỗi ngày $3$ phiên, với các điều kiện ``kích thích'', ``im lặng'' và ``không phản hồi'' ($n=27$, $17$ và $15$). Độ dài trung bình của lượt đánh tăng từ $5$ phút đầu đến $15$ phút cuối chỉ ở các mẫu sống. Trong loạt về phản hồi, mức tăng có ý nghĩa theo thời gian chỉ có ở điều kiện ``kích thích''; cuối phiên, điều kiện ``im lặng'' vẫn hơn ``không phản hồi'', nhưng với hiệu ứng nhỏ hơn, còn việc học bền vững từ phiên này sang phiên khác trong ba ngày thì không có~\cite{Kagan2022}.

\subsection*{Chu trình của bàn thử từng bước}

Cứ mỗi $10\ms$, máy tính của bàn thử làm như sau.
\begin{enumerate}
\item Đọc số xung $n_\uparrow$, $n_\downarrow$ ở hai vùng vận động trong cửa sổ vừa qua và chuẩn hóa chúng theo số đếm trung bình của vùng lúc nghỉ.
\item Dịch vợt một đoạn $\Delta p=c\,(n_\uparrow-n_\downarrow)$ và giới hạn nó trong mép sân.
\item Dịch bóng một bước; cho bóng dội lại từ tường trên, tường dưới và tường trái.
\item Nếu bóng đã tới tường phải: khi $|y-p|<0.1$ là trúng, số đếm của lượt đánh tăng một, phát kích thích trúng ($100\ms$); nếu không thì là trượt, ghi lại lượt đánh, phát kích thích trượt ($4$~s), sau đó nghỉ $4$~s, và bóng xuất phát lại.
\item Nếu không, chọn điện cực $k$ theo vị trí dọc của bóng và tần số $f$ theo khoảng cách tới tường phía vợt, rồi đưa vào điện cực $k$ các xung hai pha $75$~mV với tần số $f$ cho đến cửa sổ tiếp theo.
\end{enumerate}
Như vậy là đủ để dựng một bàn thử tương thích với bàn thử đã công bố và kiểm tra câu hỏi chính của chương: dạy mẫu nuôi cấy là tính dễ đoán, hay chỉ là sự khác biệt giữa đáp ứng với cú trúng và với cú trượt. Muốn vậy, nên thêm vào ba điều kiện của bài báo một điều kiện thứ tư mà bài báo không có: sau cú trượt là kích thích \emph{đoán trước được}, cùng độ dài và năng lượng với kích thích không đoán trước được. Nếu mẫu nuôi cấy vẫn học được, thì vấn đề nằm ở sự khác biệt, không phải ở tính dễ đoán.


\section{Tổng kết}

\begin{table}[t]
\centering
\footnotesize
\setlength{\tabcolsep}{4pt}
\renewcommand{\arraystretch}{1.15}
\begin{tabular}{>{\raggedright}p{2.1cm}>{\raggedright}p{4.2cm}>{\raggedright}p{1.9cm}>{\raggedright\arraybackslash}p{4.6cm}}
\toprule
Khối của bàn thử & Cấu tạo & Chương & Điều đã biết \\
\midrule
đầu vào & vị trí (điện cực nào trong 8) và tần số (4--40 xung/s) & \ref{sec:code}, \ref{sec:sensors} & do người làm thí nghiệm đặt \\
đầu ra & hiệu số đếm của hai vùng trong $10\ms$~\eqref{eq:dishreadout} & \ref{sec:glutcomputer}, \ref{sec:debugging} & do người làm thí nghiệm đặt; nhiễu theo~\eqref{eq:rateerror} \\
thầy dạy & kích thích đoán trước được sau cú trúng, không đoán trước được sau cú trượt; không có dopamine & \ref{sec:dofa} & mức tăng có ý nghĩa trong một phiên chỉ khi phản hồi đầy đủ; với im lặng thì hiệu ứng nhỏ hơn; không bền vững từ phiên này sang phiên khác \\
cơ chế & dự đoán đầu vào hoặc xáo trộn khi trượt~\eqref{eq:dishdensity} & \ref{sec:dofa}, \ref{sec:recognition} & chưa xác lập; cả hai đều giải thích được kết quả \\
chế độ của mạng & tiến gần chế độ tới hạn khi chơi & \ref{sec:gaba} & mối liên hệ với chất lượng chơi đã đo \\
``tri giác'' & khẳng định của tác giả & --- & chưa được chứng minh; bị phản bác \\
\bottomrule
\end{tabular}
\caption{DishBrain qua con mắt kỹ sư.}
\label{tab:dishbrain}
\end{table}

DishBrain là cỗ máy từ nơ-ron sống ở dạng đơn giản nhất: đầu vào được đặt bằng mã, đầu ra được đọc bằng phép đếm, thầy dạy là sự khác biệt giữa đáp ứng với thành công và với thất bại (bảng~\ref{tab:dishbrain}). Một mạng không có mắt, cơ và dopamine đã học chơi được một chút, và chỉ riêng điều đó đã biến nó thành bàn thử cho các ý tưởng của cuốn sách. Dù cơ chế hóa ra là gì, dự đoán, xáo trộn hay một thứ thứ ba, câu trả lời sẽ được tìm theo cách chương~\ref{sec:debugging} khuyên: bằng đầu dò, tín hiệu thử và việc tắt từng phần trên một cỗ máy mà cuối cùng ta nhìn thấy được trọn vẹn.

\section{Bài tập}

\begin{exercise}[\lvlA]
\label{ex:22:1}
\emph{Cần dịch hiệu số đếm bao nhiêu.} Hai vùng cùng phát $R_\uparrow+R_\downarrow=2000$ xung mỗi giây, các dòng xung là Poisson. (a)~Độ tản mạn tương đối của số đếm một vùng trong $10\ms$ bằng bao nhiêu khi $R=1000$ và $R=100$? (b)~Trong $1$~s bóng bay, các độ dịch~\eqref{eq:dishreadout} cộng dồn. Với $R_\uparrow-R_\downarrow$ nhỏ nhất bằng bao nhiêu thì độ dịch trung bình lớn gấp đôi độ tản mạn?
\end{exercise}

\begin{exercise}[\lvlA]
\label{ex:22:2}
\emph{Cần bao nhiêu lượt đánh để thấy việc học.} Các lượt đánh độc lập, tỷ lệ đỡ trúng có phân phối nhị thức. Mỗi giai đoạn trong hai giai đoạn cần bao nhiêu lượt để mức tăng tỷ lệ đỡ trúng vượt hai độ lệch chuẩn của hiệu: (a)~từ $0.20$ lên $0.25$; (b)~từ $0.21$ lên $0.38$, như trong mô hình?
\end{exercise}

\begin{exercise}[\lvlB]
\label{ex:22:3}
\emph{Suy ra~\eqref{eq:dishdensity}.} Cấu hình $\boldsymbol\psi$ chạy trên một tập hữu hạn; khi trượt (xác suất $P(\boldsymbol\psi)>0$), mạng chuyển sang $\boldsymbol\psi'$ với xác suất đối xứng $K(\boldsymbol\psi,\boldsymbol\psi')$, khi trúng thì giữ nguyên. (a)~Chứng minh rằng $\rho\propto1/P$ là phân phối dừng. (b)~Tỷ lệ trượt bằng bao nhiêu khi có hai cấu hình với $P=0.9$ và $P=0.1$?
\end{exercise}

\begin{exercise}[\lvlB]
\label{ex:22:4}
\emph{Vùng hấp thụ của mô hình.} Vợt đến $p=\min\bigl(1,\max(0,ay+b)\bigr)$, bóng đến độ cao $y\sim U(0,1)$, trúng nếu $|p-y|<h=0.1$. (a)~Chứng minh rằng khi $|b|<h$ và $|a-1+b|<h$ mạng đỡ trúng mọi quả bóng, và tìm diện tích vùng này. (b)~Tìm bằng số tỷ lệ đỡ trúng trung bình khi $a\sim U(-1,1)$, $b\sim U(0,1)$.
\end{exercise}

\begin{exercise}[\lvlB]
\label{ex:22:5}
\emph{Mô phỏng: hàng rào quanh các cấu hình.} Trong một bản sao của \texttt{dishbrain\_model.py}, cho $200$ mẫu nuôi cấy mỗi mẫu $20\,000$ lượt đánh. Tỷ lệ đỡ trúng trong $500$ lượt cuối là bao nhiêu? Còn nếu sau mỗi cú hích, cắt cấu hình vào hình chữ nhật $a\in[-1,2]$, $b\in[-1,1]$? Giải thích sự khác biệt bằng bài tập~\ref{ex:22:3}.
\end{exercise}

\begin{exercise}[\lvlB]
\label{ex:22:6}
\emph{Thiết kế mã đầu vào.} Bóng được đỡ trúng nếu sai số về độ cao nhỏ hơn $h=0.1$; mạng đọc đầu vào trong $T=0.25$~s, với tần số đến $r$ phân biệt được khoảng $\sqrt{rT}$ mức, kích thích không dày quá $40$ xung mỗi giây. (a)~Mã vị trí trên tám điện cực có đủ không? (b)~Có thể truyền độ cao bằng tần số trên một điện cực không?
\end{exercise}

\begin{exercise}[\lvlC]
\label{ex:22:7}
\emph{Nghiên cứu: dự đoán hay xáo trộn?} Tổng quát hóa mô hình lên $d$ số điều chỉnh được ($p=\sum_{i<d}a_iy^i$). Số lượt đánh cần để đạt tỷ lệ đỡ trúng $0.5$ tăng theo $d$ như thế nào khi xáo trộn và khi dùng quy tắc có dấu của sai số~\eqref{eq:perturbation}? Thí nghiệm nào trên bàn thử sẽ phân biệt được hai giả thuyết?
\end{exercise}

\appendix
\renewcommand{\chaptername}{\appendixname}
\chapter{Những chất không quyết định loại nơ-ron}
\label{sec:othermediators}

Ở chương~\ref{sec:neurontypes}, ta đã phân loại nơ-ron theo chất dẫn truyền thần kinh chính và được mười loại (bảng~\ref{tab:types}). Giờ là lúc thêm vào những điều phức tạp. Nơ-ron không chỉ nhả chất dẫn truyền chính, và trong não, ngoài các chất dẫn truyền chính còn có những chất khác hoạt động. Chúng thay đổi cách mạng truyền và lưu tín hiệu.

Ngoài bus địa chỉ và bus quảng bá (mục~\ref{sec:buses}) còn có một bus thứ ba: \emph{kênh ngược}. Một số chất không do bên gửi mà do \emph{bên nhận} nhả ra, và chúng đi ngược qua khe về phía đầu ra của bên gửi, làm thay đổi lượng chất dẫn truyền mà bên gửi sẽ nhả ở lần sau. Trong mạng truyền thông, điều này giống tín hiệu xác nhận đã nhận mà bên gửi dùng để điều chỉnh việc truyền của mình. Nơ-ron dùng chính kênh này khi thay đổi trọng số liên kết: qua nó, đầu ra của bên gửi biết được hoạt động của bên nhận, tức thừa số thứ hai trong các quy tắc học của chương~\ref{sec:dofa}.

Danh sách đầy đủ các chất dẫn truyền đã biết rất dài: hơn một trăm chất, phần lớn là những chuỗi protein ngắn, gọi là neuropeptide. Nhưng tất cả đều xếp được vào vài lớp, và trong mỗi lớp các chất hoạt động tương tự nhau. Những chất không bao giờ là chất dẫn truyền chính được gom trong bảng~\ref{tab:othermediators}, với đúng ba câu hỏi mà kỹ sư sẽ đặt ra: tín hiệu đi theo bus nào, tác động nhanh hay chậm, và dấu thường gặp là gì. Chúng không quyết định loại nơ-ron: chúng được nhả cùng chất dẫn truyền chính, hoặc không do nơ-ron nhả ra, hoặc đi theo chiều ngược lại.

\begin{table}[t]
\centering
\footnotesize
\setlength{\tabcolsep}{4pt}
\renewcommand{\arraystretch}{1.12}
\begin{tabular}{>{\raggedright}p{2.25cm}>{\raggedright}p{2.8cm}>{\raggedright}p{1.8cm}>{\raggedright}p{1.5cm}>{\centering}p{0.8cm}>{\raggedright\arraybackslash}p{4.3cm}}
\toprule
Lớp & Chất & Bus & Tốc độ & Dấu & Vai trò trong tính toán \\
\midrule
Amino acid & D-serine & cục bộ & chậm & VÀ & chìa khóa thứ hai cho thụ thể glutamate: điều kiện ``VÀ'' \\
\midrule
Monoamine & amine vết & quảng bá & chậm & $\pm$ & tinh chỉnh các kênh monoamine \\
\midrule
\multirow{2}{2.25cm}{Purine}
 & ATP & địa chỉ & nhanh và chậm & $+$ & chất dẫn truyền đi kèm; tín hiệu giữa nơ-ron và tế bào thần kinh đệm \\
 & adenosine & thể tích & chậm & $-$ & tích tụ khi hoạt động: bộ đếm tải~\cite{Dunwiddie2001} \\
\midrule
\multirow{8}{2.25cm}{Neuropeptide (hơn một trăm)}
 & enkephalin, endorphin, dynorphin & thể tích & chậm & $-$ & kìm hãm truyền dẫn \\
 & chất P & thể tích & chậm & $+$ & tăng cường chậm \\
 & neuropeptide Y & thể tích & chậm & $-$ & ức chế chậm \\
 & somatostatin & thể tích & chậm & $-$ & ức chế chậm, đi kèm GABA \\
 & peptide ruột hoạt mạch (VIP) & thể tích & chậm & $+$ & giải ức chế: ức chế các nơ-ron ức chế \\
 & cholecystokinin & thể tích & chậm & $\pm$ & đi kèm GABA ở một phần nơ-ron ức chế \\
 & orexin & quảng bá & chậm & $+$ & bộ ổn định chế độ thức \\
 & oxytocin, vasopressin & quảng bá & chậm & $\pm$ & các thiết lập hành vi toàn cục \\
\midrule
\multirow{2}{2.25cm}{Khí}
 & nitơ oxit NO & ngược, thể tích & chậm & $\pm$ & đi xuyên qua màng; tín hiệu ngược khi thay đổi trọng số~\cite{Garthwaite2008} \\
 & CO, H$_2$S & thể tích & chậm & $\pm$ & vai trò còn ít được nghiên cứu \\
\midrule
Lipid & endocannabinoid (anandamide, 2-AG) & ngược & chậm & $-$ & bên nhận giảm lượng nhả của bên gửi: phản hồi theo trọng số~\cite{WilsonNicoll2001} \\
\bottomrule
\end{tabular}
\caption{Những chất không quyết định loại nơ-ron. \emph{Bus}, \emph{tốc độ} và \emph{dấu} như trong bảng~\ref{tab:types}; ngoài ra, bus thể tích --- chất dẫn truyền lan ra trong khoảng gian bào quanh chỗ nhả, ngược --- đi từ bên nhận ngược về bên gửi, cục bộ --- tác động gần chỗ nhả. ``VÀ'' là tín hiệu cho phép: tự nó không gây ra dòng điện, nhưng thiếu nó thì một đầu vào khác không kích hoạt được, như đầu vào thứ hai của cổng logic ``VÀ''. Neuropeptide hầu như luôn được nhả cùng chất dẫn truyền chính và chủ yếu khi tần số xung cao~\cite{vandenPol2012}.}
\label{tab:othermediators}
\end{table}

\chapter{Ký hiệu}
\label{sec:notation}

Các loại nơ-ron được ký hiệu bằng bốn chữ cái đầu trong tên tiếng Anh của chất dẫn truyền thần kinh chính: \nt{GLUT} --- glutamate, \nt{GABA} --- axit gamma-aminobutyric, \nt{GLYC} --- glycine, \nt{ACET} --- acetylcholine, \nt{DOPA} --- dopamine, \nt{SERO} --- serotonin, \nt{HIST} --- histamine, \nt{NORA} --- noradrenaline, \nt{ADRE} --- adrenaline, \nt{ASPA} --- aspartate (bảng~\ref{tab:types}).

Chữ cái ít hơn đại lượng, nên ở các chương khác nhau đôi khi cùng một chữ mang nghĩa khác nhau. Trong bảng dưới đây, mỗi nghĩa có một dòng riêng; cột bên phải là công thức nơi nó được đưa ra.

\begin{center}
\footnotesize
\setlength{\tabcolsep}{4pt}
\renewcommand{\arraystretch}{1.12}
\begin{longtable}{>{\raggedright}p{2.0cm}>{\raggedright}p{9.6cm}>{\raggedright\arraybackslash}p{2.4cm}}
\toprule
Ký hiệu & Ý nghĩa & Được đưa ra ở \\
\midrule
\endfirsthead
\toprule
Ký hiệu & Ý nghĩa & Được đưa ra ở \\
\midrule
\endhead
\bottomrule
\endfoot
$a$ & độ dốc của đặc tuyến truyền tuyến tính, $f(u)=au$ & \eqref{eq:feedback} \\
$a$ & mức \nt{ACET}: $0$ --- đọc, $1$ --- ghi & \eqref{eq:acetnet} \\
$a_i$, $a$ & sự mỏi của nhóm: trong bộ tạo nhịp; trong dao động của giấc ngủ sóng chậm & \eqref{eq:halfcentre}, \eqref{eq:updown} \\
$a_S$, $a_P$ & độ nhạy của nhịp tim với chân ga và với phanh & \eqref{eq:gasbrake} \\
$A(r)$, $A_0$ & đáp ứng của khớp thần kinh với một xung ở tần số $r$; của khớp đã nghỉ & \eqref{eq:depression} \\
$A_1$ & đáp ứng của bên nhận với một gói chất dẫn truyền & \eqref{eq:quantal} \\
$A$ & diện tích màng của nơ-ron tính cả đuôi gai & \eqref{eq:dendriteload} \\
$A_f$, $A_s$ & phần hiệu chỉnh được giữ lại đến lần chuyển động sau, ở quá trình nhanh và quá trình chậm & \eqref{eq:tworate} \\
$b_i$ & dòng vào riêng của bộ tạo nhịp & \eqref{eq:seroneuron} \\
$b$ & công việc làm mỏi nhóm nhanh đến đâu, trong bộ tạo nhịp hoặc trong dao động của giấc ngủ & \eqref{eq:halfcentre}, \eqref{eq:updown} \\
$b$ & số bit cho một mẫu số hóa & \eqref{eq:probedata} \\
$B_f$, $B_s$ & mức độ phần hiệu chỉnh học từ sai số, ở quá trình nhanh và quá trình chậm & \eqref{eq:tworate} \\
$c$ & nồng độ chất dẫn truyền trong thể tích & \eqref{eq:seroneuron} \\
$c$ & hệ số tương quan giữa nhiễu của các nơ-ron lân cận & \eqref{eq:correlated} \\
$c$ & hệ số theo đó hiệu số đếm đẩy vợt & \eqref{eq:dishreadout} \\
$c_f$ & đáp ứng của nơ-ron $C$: giá trị lớn nhất của $s_{f,p}$ trên mọi vị trí & \eqref{eq:scstage} \\
$c_m$ & điện dung riêng của màng & \eqref{eq:dendriteload} \\
$C$ & giá của nỗ lực: hành động trong thời gian $\tau_a$ tốn $C/\tau_a$ & \eqref{eq:vigor} \\
$d'$ & độ phân biệt của hai đầu vào: hiệu số chia cho nhiễu & \eqref{eq:gainsnr} \\
$d$, $d_j$ & độ trễ trên đường đi của tín hiệu & \eqref{eq:glutneuron}, \eqref{eq:vstar} \\
$d$ & độ trễ của vòng phản hồi với cơ & \eqref{eq:delaystable} \\
$D$ & hệ số khuếch tán & \eqref{eq:diffusionlength} \\
$D$ & mẫu số trong tần số của mạng \nt{GLUT}--\nt{GABA} & \eqref{eq:isn} \\
$D$ & thời gian chờ phần thưởng bị trì hoãn & \eqref{eq:patience} \\
$e$, $e_{ij}$ & vết ở khớp thần kinh & \eqref{eq:threefactor} \\
$e$, $e_i$ & sai số của đầu ra hoặc của chuyển động, đến qua dây dẫn sai số & \eqref{eq:lms}, \eqref{eq:adaptivefilter}, \eqref{eq:tworate} \\
$E_{\text{xung}}$ & năng lượng của một xung tính cả công việc của khớp thần kinh & \eqref{eq:energy} \\
$E$ & tần số của nhóm \nt{GLUT} & \eqref{eq:ei} \\
$E_E$ & điện thế cân bằng của khớp thần kinh kích thích & \eqref{eq:shunt} \\
$f$, $F$ & đặc tuyến truyền: tần số như một hàm của đầu vào & \eqref{eq:ratenet}, \eqref{eq:gain} \\
$f$, $f_0$ & tần số nhịp tim; nhịp riêng của bộ tạo nhịp & \eqref{eq:gasbrake} \\
$f_s$ & tần số số hóa của một điện cực & \eqref{eq:probedata} \\
$F(t)$, $F_1$ & lực của cơ; lực của một cú giật & \eqref{eq:forcesum} \\
$F_1$, $F_2$ & lực của hai cơ kéo khớp theo hai hướng ngược nhau & \eqref{eq:antagonists} \\
$g$ & độ nhạy của thụ thể với nồng độ & \eqref{eq:seroneuron} \\
$g_L$, $g_E$, $g_I$ & độ dẫn rò, độ dẫn kích thích và độ dẫn ức chế & \eqref{eq:shunt} \\
$g$ & hệ số khuếch đại do \nt{NORA} đặt & \eqref{eq:gain}, \eqref{eq:machine} \\
$g$ & ức chế chung từ \nt{GABA} trong mạng nhớ & \eqref{eq:acetnet} \\
$g$ & hệ số khuếch đại của cổng đau, được đặt từ trên xuống & \eqref{eq:gate} \\
$g$ & hệ số khuếch đại của một tầng trong thác hormone & \eqref{eq:cascade} \\
$g_j$ & bộ lọc $j$ với hằng số thời gian riêng ở đầu vào bộ lọc thích nghi & \eqref{eq:adaptivefilter} \\
$G$ & hệ số khuếch đại của vòng phản hồi & \eqref{eq:feedback} \\
$G$ & tốc độ sửa sai trong vòng có trễ & \eqref{eq:delaystable} \\
$h$, $h_I$ & đầu vào ngoài của nhóm \nt{GLUT} và nhóm \nt{GABA} & \eqref{eq:ei}, \eqref{eq:isn} \\
$h(t)$ & một cú giật đơn của cơ & \eqref{eq:twitch} \\
$H(t)$ & ngưỡng trên của áp lực ngủ: chạm tới đó thì cỗ máy ngủ & \eqref{eq:sleeppressure} \\
$I$, $I_i$ & dòng vào từ bên ngoài tới nơ-ron hoặc nhóm & \eqref{eq:ratenet}, \eqref{eq:wta}, \eqref{eq:updown} \\
$I$ & tần số của nhóm \nt{GABA} & \eqref{eq:ei} \\
$I$ & độ sáng hoặc cường độ của kích thích & \eqref{eq:photonnoise}, \eqref{eq:nakarushton} \\
$I$ & dòng điện đầu vào nạp cho màng & \eqref{eq:dendriteload} \\
$k$ & số xung trong một chùm xung & \eqref{eq:burst} \\
$k$ & phần tăng thêm phải lớn hơn nhiễu bao nhiêu lần & \eqref{eq:photonnoise} \\
$k$ & hệ số phản hồi của thác hormone & \eqref{eq:cascade} \\
$k_s$ & mức độ cơn đau bơm thêm vào sự nhạy cảm hóa & \eqref{eq:sensitization} \\
$K$ & số đặc trưng của kết quả & \eqref{eq:vectortd} \\
$K$ & độ cứng của khớp & \eqref{eq:antagonists} \\
$K$ & hệ số khuếch đại vòng của thác hormone, $K=g^3k$ & \eqref{eq:cascade}, \eqref{eq:cascadestable} \\
$K_\varphi$ & độ mạnh của việc chỉnh bộ tạo nhịp ngày đêm theo tín hiệu ngoài & \eqref{eq:pll} \\
$L$ & độ dài đường dây trễ & \eqref{eq:itd} \\
$L$ & độ dài dây dẫn từ cảm biến đau & \eqref{eq:paintime} \\
$L(t)$ & ngưỡng dưới của áp lực ngủ: chạm tới đó thì cỗ máy thức dậy & \eqref{eq:sleeppressure} \\
$M$ & thước đo mối liên hệ nhân quả giữa tín hiệu báo trước và phần thưởng & \eqref{eq:retro} \\
$M_k$ & kỳ vọng của đặc trưng $k$ & \eqref{eq:vectortd} \\
$M_{ik}$ & độ mạnh ức chế từ nơ-ron \nt{GABA} $k$ & \eqref{eq:machine} \\
$M$ & mô-men quay của khớp & \eqref{eq:antagonists} \\
$M$ & số đường đầu vào của các bộ trộn & \eqref{eq:cover} \\
$n$, $\bar n$ & số xung trong một cửa sổ; giá trị trung bình của nó & \eqref{eq:rateerror} \\
$n$ & số đầu vào của phần tử ngưỡng & \eqref{eq:cover} \\
$n_\uparrow$, $n_\downarrow$ & số đếm xung của các vùng ``lên'' và ``xuống'' trong một cửa sổ & \eqref{eq:dishreadout} \\
$N$ & số nơ-ron & \eqref{eq:correlated}, \eqref{eq:energy} \\
$N$ & số điện cực của đầu dò & \eqref{eq:probedata} \\
$N_s$ & số vị trí nhả giữa hai nơ-ron & \eqref{eq:quantal} \\
$p$ & xác suất nhả chất dẫn truyền ứng với một xung & \eqref{eq:burst} \\
$p$ & tỷ lệ nơ-ron hoạt động trong mạng nhớ & \eqref{eq:acetnet} \\
$p$ & nhiễu loạn mà hệ học cách bù trừ & \eqref{eq:tworate} \\
$P$ & công suất tiêu tốn cho tín hiệu & \eqref{eq:energy}, \eqref{eq:dishpower} \\
$P$ & độ bất định của ước tính đang lưu & \eqref{eq:kalman} \\
$P$ & hoạt động của ``phanh'': dây dẫn \nt{ACET} tới tim & \eqref{eq:gasbrake} \\
$P_{\max}$ & số mẫu ngẫu nhiên tối đa mà phần tử ngưỡng có thể chia thành hai lớp & \eqref{eq:cover} \\
$P_{\text{trượt}}$ & xác suất trượt & \eqref{eq:dishdensity} \\
$q$ & tốc độ thay đổi của thế giới & \eqref{eq:kalman} \\
$q_k$ & tần số của nơ-ron \nt{GABA} $k$ & \eqref{eq:machine} \\
$q_0$ & hoạt động nền của nơ-ron trung gian ức chế & \eqref{eq:gate} \\
$q$ & tần số của nơ-ron trung gian ức chế trong cổng đau & \eqref{eq:gate} \\
$q$ & đơn vị vận động kế tiếp mạnh hơn đơn vị trước bao nhiêu lần & \eqref{eq:sizeprinciple} \\
$r$, $r_i$ & tần số xung của nơ-ron & \eqref{eq:ratenet} \\
$r$, $r_t$ & phần thưởng (dòng phần thưởng) & \eqref{eq:rpe}, \eqref{eq:ctd} \\
$\mathbf r(t)$ & vectơ đáp ứng của hồ chứa & \eqref{eq:reservoir} \\
$R$, $\bar R$ & phần thưởng nhận được; kỳ vọng của nó hoặc giá trị trung bình trên một đơn vị thời gian & \eqref{eq:perturbation}, \eqref{eq:vigor} \\
$R$ & phương sai của nhiễu ở đầu vào bộ lọc & \eqref{eq:kalman} \\
$R$, $r_0$ & phần thưởng bị trì hoãn và phần thưởng tức thì & \eqref{eq:patience} \\
$R_{\max}$ & tốc độ truyền giới hạn, bit/s & \eqref{eq:spikeentropy} \\
$r_{\max}$ & tần số xung lớn nhất & \eqref{eq:gain}, \eqref{eq:nakarushton} \\
$R_{\text{thô}}$, $R_{\text{xung}}$ & luồng dữ liệu của đầu dò: thô và chỉ các xung, bit/s & \eqref{eq:probedata} \\
$s_j$ & dấu của các đầu ra của nơ-ron $j$ & \eqref{eq:dalesign} \\
$s_i$ & dấu của thụ thể ở bên nhận & \eqref{eq:seroneuron} \\
$s$, $\hat s$ & trạng thái của thế giới; ước tính của nó & \eqref{eq:rpe}, \eqref{eq:kalman} \\
$s_{f,p}$ & đáp ứng của nơ-ron $S$ với đặc trưng $f$ ở vị trí $p$ & \eqref{eq:scstage} \\
$S$ & hoạt động của ``chân ga'': dây dẫn \nt{NORA} tới tim & \eqref{eq:gasbrake} \\
$S$ & áp lực ngủ & \eqref{eq:sleeppressure} \\
$t_1$ & thời điểm của xung đầu tiên & \eqref{eq:latency} \\
$t_k$ & thời điểm của xung thứ $k$ & \eqref{eq:forcesum} \\
$t_{f,i}$ & mẫu của đặc trưng $f$ & \eqref{eq:scstage} \\
$T$ & cửa sổ đếm xung; thời gian tích lũy photon & \eqref{eq:rateerror}, \eqref{eq:photonnoise} \\
$T_c$ & thời gian tăng của cú giật cơ & \eqref{eq:twitch} \\
$T_{\text{đợt}}$ & khoảng cách giữa các đợt bùng phát của mẫu nuôi cấy & \eqref{eq:burstinterval} \\
$u$ & tổng đầu vào của nơ-ron & \eqref{eq:gain} \\
$u$, $u(t)$ & lệnh của não: ở đầu vào bộ lọc thích nghi; tới thác hormone & \eqref{eq:adaptivefilter}, \eqref{eq:cascade} \\
$u(t)$ & đầu vào của hồ chứa & \eqref{eq:reservoir} \\
$U$ & mức của đầu vào không đổi & \eqref{eq:latency} \\
$U$ & phần của kho chất dẫn truyền được nhả ra ứng với một xung & \eqref{eq:depression} \\
$v$ & tốc độ tín hiệu trong dây dẫn; tốc độ di chuyển của vệt sáng & \eqref{eq:itd}, \eqref{eq:vstar} \\
$V$ & điện áp trên màng & \eqref{eq:glutneuron}, \eqref{eq:shunt} \\
$V$ & kỳ vọng phần thưởng tương lai & \eqref{eq:rpe}, \eqref{eq:ctd} \\
$w$, $w_{ij}$, $W_{ij}$ & trọng số liên kết & \eqref{eq:glutneuron}, \eqref{eq:ratenet}, \eqref{eq:acetnet}, \eqref{eq:lms}, \eqref{eq:adaptivefilter} \\
$\mathbf w_k$ & trọng số các đầu vào của nơ-ron $C$ thứ $k$ & \eqref{eq:traceinvariance} \\
$w_{q\mu}$, $w_{q\nu}$, $w_{z\mu}$ & trọng số các liên kết trong cổng đau & \eqref{eq:gate} \\
$W_{\text{ra}}$ & trọng số của lớp đọc ra tuyến tính của hồ chứa & \eqref{eq:reservoir} \\
$x$, $y$ & các đầu vào và đầu ra logic & \eqref{eq:dualrail}, \eqref{eq:veto} \\
$x$, $y$ & hoạt động của bên gửi và bên nhận & \eqref{eq:hebb} \\
$x$ & vị trí của bộ phát hiện trên đường dây trễ & \eqref{eq:itd} \\
$x_i$ & đầu vào từ các giác quan & \eqref{eq:acetnet} \\
$x_p$ & đầu vào ở vị trí $p$ & \eqref{eq:scstage} \\
$\mathbf x$ & các đầu ra của nơ-ron $S$, đầu vào của nơ-ron $C$ & \eqref{eq:traceinvariance} \\
$x_j$ & đầu vào $j$ của bộ lọc thích nghi: lệnh sau bộ lọc $g_j$ & \eqref{eq:lms}, \eqref{eq:adaptivefilter} \\
$x_f$, $x_s$ & phần hiệu chỉnh của quá trình nhanh và quá trình chậm & \eqref{eq:tworate} \\
$x_1$, $x_2$, $x_3$ & các mức ở ba tầng của thác hormone; $x_3$ --- hormone cuối cùng & \eqref{eq:cascade} \\
$x^*$ & tỷ lệ kho chất dẫn truyền đã hồi phục mà tại đó một trận thác mới bắt đầu & \eqref{eq:burstinterval} \\
$y$ & đầu vào có nhiễu của bộ lọc & \eqref{eq:kalman} \\
$y$, $\hat y$ & tín hiệu cảm biến; dự đoán của nó bởi bộ lọc thích nghi & \eqref{eq:adaptivefilter} \\
$y_k$, $\bar y_k$ & hoạt động của nơ-ron $C$ thứ $k$; vết của nó & \eqref{eq:traceinvariance} \\
$y(t)$ & đầu ra của lớp đọc ra của hồ chứa & \eqref{eq:reservoir} \\
$z$ & tần số đầu ra của cổng đau & \eqref{eq:gate} \\
\midrule
$\alpha_i^\pm$ & trọng số của sai số dương và sai số âm & \eqref{eq:asymmetric} \\
$\beta$ & độ mạnh của ức chế lẫn nhau & \eqref{eq:wta} \\
$\beta$ & độ mạnh ức chế đầu ra của cổng đau & \eqref{eq:gate} \\
$\gamma$ & hệ số chiết khấu & \eqref{eq:rpe} \\
$\delta(\cdot)$ & hàm delta: bước nhảy tức thời & \eqref{eq:glutneuron} \\
$\delta$, $\delta_t$ & sai số dự đoán phần thưởng & \eqref{eq:rpe}, \eqref{eq:ctd} \\
$\delta t$ & chênh lệch thời điểm âm thanh tới hai tai & \eqref{eq:itd} \\
$\Delta$, $\Delta_{\max}$ & khoảng cách giữa các xung; cửa sổ trùng khớp & \eqref{eq:window} \\
$\Delta t$ & độ chính xác mà bên nhận phân biệt được thời gian & \eqref{eq:spikeentropy} \\
$\Delta F$ & lực tăng thêm từ đơn vị vận động kế tiếp & \eqref{eq:sizeprinciple} \\
$\Delta I$ & phần tăng độ sáng phân biệt được & \eqref{eq:photonnoise} \\
$\Delta p$ & độ dịch của vợt trong một cửa sổ & \eqref{eq:dishreadout} \\
$\Delta u$ & hiệu đầu vào của hai nhóm & \eqref{eq:gainsnr} \\
$\Delta V$ & cần nâng điện áp màng lên bao nhiêu & \eqref{eq:dendriteload} \\
$\Delta x$ & khoảng cách giữa các cảm biến lân cận & \eqref{eq:vstar} \\
$\varepsilon$ & hiệu tương đối của các đầu vào & \eqref{eq:wtagain} \\
$\eta$ & tốc độ học & \eqref{eq:plasticity}, \eqref{eq:lms}, \eqref{eq:traceinvariance} \\
$\eta$ & nhiễu trong mạng sau bộ khuếch đại & \eqref{eq:softmax} \\
$\mu$ & đầu vào xúc giác tới cổng đau & \eqref{eq:gate} \\
$\nu$ & đầu vào đau & \eqref{eq:gate} \\
$\theta$ & ngưỡng kích hoạt & \eqref{eq:window}, \eqref{eq:scstage} \\
$\kappa$ & lượng chất dẫn truyền ứng với một xung & \eqref{eq:seroneuron} \\
$\kappa_i$ & tỷ số giữa trọng số của sai số dương và sai số âm & \eqref{eq:expectile} \\
$\kappa_m$ & liên hệ giữa lực của cơ và độ cứng của nó & \eqref{eq:antagonists} \\
$\ell$ & khoảng cách mà chất dẫn truyền lan ra & \eqref{eq:diffusionlength} \\
$\ell$ & cánh tay đòn mà cơ kéo khớp & \eqref{eq:antagonists} \\
$\xi$ & dao động ngẫu nhiên ở đầu ra của nơ-ron & \eqref{eq:perturbation} \\
$\rho(\boldsymbol\psi)$ & tỷ lệ thời gian ở trong cấu hình $\boldsymbol\psi$ & \eqref{eq:dishdensity} \\
$\sigma$ & độ tản mạn, thang của nhiễu & \eqref{eq:rateerror}, \eqref{eq:softmax} \\
$\sigma$ & độ sáng mà tại đó đáp ứng bằng một nửa giá trị lớn nhất & \eqref{eq:nakarushton} \\
$\sigma_{\text{trước}}$, $\sigma_{\text{sau}}$ & nhiễu trước bộ khuếch đại và sau nó & \eqref{eq:gainsnr} \\
$\tau$ & hằng số thời gian của màng & \eqref{eq:glutneuron} \\
$\tau$ & hằng số thời gian của một tầng trong thác hormone & \eqref{eq:cascade} \\
$\tau_{\text{hd}}$ & hằng số thời gian của nhóm tự kích thích chính nó & \eqref{eq:perfectint} \\
$\tau_c$ & thời gian sống của chất dẫn truyền trong thể tích & \eqref{eq:seroneuron} \\
$\tau_E$, $\tau_I$ & hằng số thời gian của nhóm \nt{GLUT} và nhóm \nt{GABA} & \eqref{eq:ei} \\
$\tau_e$ & thời gian sống của vết & \eqref{eq:threefactor} \\
$\tau_{\text{tr}}$ & thời gian sống của vết hoạt động của nơ-ron $C$ & \eqref{eq:traceinvariance} \\
$\tau_s$ & thời gian phục hồi độ nhạy sau tổn thương & \eqref{eq:sensitization} \\
$\tau_r$, $\tau_d$ & thời gian tích lũy áp lực ngủ khi thức và thời gian xả nó khi ngủ & \eqref{eq:sleeppressure} \\
$\tau_{\text{ph}}$ & thời gian phục hồi kho chất dẫn truyền & \eqref{eq:depression} \\
$\tau_\gamma$ & tầm nhìn lập kế hoạch & \eqref{eq:ctd} \\
$\tau_v$ & tốc độ chính xác hóa kỳ vọng & \eqref{eq:ramp} \\
$\tau_a$ & thời gian mỏi trong bộ tạo nhịp hoặc trong dao động của giấc ngủ & \eqref{eq:halfcentre}, \eqref{eq:updown} \\
$\tau_a^*$ & thời gian thực hiện hành động tốt nhất & \eqref{eq:vigor} \\
$\Phi$ & vế phải của quy tắc học & \eqref{eq:plasticity} \\
$\Phi$ & dòng photon & \eqref{eq:photonnoise} \\
$\varphi_k$ & đặc trưng $k$ của kết quả nhận được & \eqref{eq:vectortd} \\
$\varphi$ & độ lệch pha giữa bộ tạo nhịp ngày đêm và tín hiệu ngoài & \eqref{eq:pll} \\
$\boldsymbol\psi$ & cấu hình của mạng trong mô hình DishBrain & \eqref{eq:dishdensity} \\
$\omega$, $\omega_0$ & tần số của tín hiệu ngoài (ánh sáng); tần số riêng của bộ tạo nhịp ngày đêm & \eqref{eq:pll} \\
\end{longtable}
\end{center}

\chapter{Đáp án}
\label{sec:answers}

Đáp án ngắn cho các bài tập cuối chương. Lời giải đầy đủ nằm trong một tài liệu riêng dành cho giảng viên.

\answer{ex:01:1}{(a)~$8\cdot10^{10}$ \nt{GLUT} --- $10$ lần dân số Trái Đất; nơ-ron quảng bá $\approx1.9\cdot10^6$ --- một thành phố cỡ Viên. (b)~$\approx8\cdot10^{14}$ đầu tận cùng, trong đó của nơ-ron quảng bá $\approx3\cdot10^{11}$, tỷ lệ $\approx4\cdot10^{-4}$ --- lớn hơn $\approx16$ lần tỷ lệ các nơ-ron này ($2.2\cdot10^{-5}$).}
\answer{ex:01:2}{(a)~$\tau_c\ln(c_0/c_*)\approx4.6\ms$. (b)~Đường truyền rảnh khi $T\ge\tau_c\ln101$: tới $\approx22$ lần nhả mỗi giây thay vì $\approx220$.}
\answer{ex:01:3}{$V(s_k)=\gamma^{K-1-k}r$; $\delta_{\text{đèn}}=\gamma^Kr=re^{-T/\tau_\gamma}\approx0.60\,r$ với mọi $\Delta$, $\tau_\gamma\approx1.95\,\text{s}$.}
\answer{ex:01:4}{$n^*=93$, $47$, $25$, $12$: $n^*\approx5/\alpha$ khi $\alpha$ nhỏ.}
\answer{ex:01:5}{\nt{GLUT}: $1\to10\to10^4\to10^7$ --- $\approx10^7$ nơ-ron, $4$ mắt xích, $4\ms$. \nt{DOPA}: $1\to10\to3.2\cdot10^5$ --- $\approx3\cdot10^5$ nơ-ron, $3$ mắt xích, ít hơn $30$ lần; cùng bậc với số nơ-ron \nt{DOPA}.}
\answer{ex:01:6}{$\lambda=f_EJ_E-f_IJ_I$, suy ra $J_I=37.5$; tỷ số dòng $f_IJ_I/f_EJ_E=0.94$; thác hoạt động ($\lambda\ge1$) xảy ra khi ức chế chỉ yếu đi $6.7\,\%$.}
\answer{ex:01:7}{$n\ge57$ lần trình bày cho mỗi độ trễ trong mỗi điều kiện. Cùng sự suy giảm đó có thể đến, chẳng hạn, từ sai số của ước lượng thời gian bên trong, tăng theo $T$.}

\answer{ex:02:1}{Hàng dài $3.2$~mm với bước $25\um$: $\approx129$ bộ phát hiện. Phát xung là một dải rộng $v\,\tau\ln2=1.7$~mm --- $\approx69$ bộ phát hiện, hơn nửa hàng; hướng là tâm của dải.}
\answer{ex:02:2}{Cửa sổ $-\tau\ln\frac{w_2}{\theta-w_1}\le\delta\le\tau\ln\frac{w_1}{\theta-w_2}$, tức là $[-0.235,\,0.178]\ms$; tâm $c=\frac\tau2\ln\frac{w_1(\theta-w_1)}{w_2(\theta-w_2)}=-28\,\mu\text{s}$. Hướng lệch $28\,\mu\text{s}$ về phía tai mạnh hơn --- gần ba bước phân giải.}
\answer{ex:02:3}{Cần các $s_i$ với $s_iw_{ij}s_j\ge0$; chúng tồn tại khi và chỉ khi mọi chu trình đều chẵn. Ức chế lẫn nhau quy được (không có dao động bền vững), vòng \nt{GLUT}--\nt{GABA} thì không (có thể dao động).}
\answer{ex:02:4}{Sáu nơ-ron $g_k=[x+y+z\ge k]$ và $\bar g_k=[\bar x+\bar y+\bar z\ge4-k]$, $k=1,2,3$; $C=g_2$, $\bar C=\bar g_2$, $S=[g_1+\bar g_2+g_3\ge2]$, $\bar S=[\bar g_1+g_2+\bar g_3\ge2]$: tám nơ-ron. Bằng AND và OR --- $12$.}
\answer{ex:02:5}{$T_{\min}(0.6)=0.718\,\tau$; $I_c=0.225$, gần ngưỡng $T_{\min}\propto(I-I_c)^{-1/2}$.}
\answer{ex:02:6}{(a)~$500$ mắt xích, $5\cdot10^4$ nơ-ron; độ phân tán $4.5\ms$ ($0.45\,\%$), sai số tương đối $\propto T^{-1/2}$. (b)~Một hệ số ngẫu nhiên của độ trễ, chung cho mọi mắt xích, với độ phân tán $5\,\%$.}
\answer{ex:02:7}{Làm tươi mỗi $\tau_F\ln2=0.69\,\text{s}$ một lần: $4.3$ xung mỗi giây trên một nơ-ron so với $20$, tiết kiệm hơn $4.6$ lần ($4\cdot10^{-7}$ so với $2\cdot10^{-6}$~J với $10^{-10}$~J mỗi xung). Flip-flop mất bit; tụ điện giữ được bit nếu việc làm câm bắt đầu không muộn hơn $0.49\,\text{s}$ sau lần làm tươi (xác suất $0.71$).}

\answer{ex:03:1}{$35\to17.5\mV$ khi $g_E=g_L$; $56\to40\mV$ khi $g_E=4g_L$, phép trừ sẽ cho $38.5\mV$: sun chia $g_E$ cho $3$. Cửa sổ hẹp đi một nửa ($\tau_{\text{hd}}$ giảm $2$ lần).}
\answer{ex:03:2}{$\operatorname{tr}=\frac{w_{EE}-1}{\tau_E}-\frac1{\tau_I}$, $\det=\frac{1-w_{EE}+w_{EI}w_{IE}}{\tau_E\tau_I}$; ở biên $\omega=\sqrt{\det}$, $\tau_I=20\ms$, chu kỳ $73\ms$.}
\answer{ex:03:3}{$d_c=\tau_E\arccos(-1/G)/\sqrt{G^2-1}$, chu kỳ $2\pi\tau_E/\sqrt{G^2-1}\in(2d_c,4d_c)$. $G=2$: $d_c=12.1\ms$, chu kỳ $36\ms$; $G=5$: $d_c=3.6\ms$, chu kỳ $12.8\ms$ --- đúng dải nhịp nhanh.}
\answer{ex:03:4}{(a)~Ở $I_2=\beta I_1=2$ khi tăng và $I_2=I_1/\beta=0.5$ khi giảm, vòng trễ rộng $1.5$. (b)~Thêm $\tau\ln10/(\beta-1)=2.3\tau$ cho mỗi bậc thập phân của $\varepsilon$.}
\answer{ex:03:5}{Ba trung gian, $d_k=5$, $10$, $15\ms$: các phép cấm phải phủ $15\ms$, mỗi phép $5\ms$.}
\answer{ex:03:6}{$n+M+1=3+4+1=8$ nơ-ron. Hai nơ-ron: \nt{GABA} $=[x+y+z\ge2]$, đầu ra $=[x+y+z-2\,\nt{GABA}\ge1]$.}
\answer{ex:03:7}{$\binom84=70<100\le\binom94=126$: $9$ trung gian (định lý Sperner). Không có liên kết trực tiếp --- $100$: hạng của $\beta(J-I)$ bằng $n$.}
\answer{ex:03:8}{$h_c=\frac1{2w_{EE}}+\frac{w_{EI}w_{IE}-w_{EE}}{4w_{EE}^2}=\frac7{18}\approx0.39$; mô phỏng cho thấy đổi dấu giữa $h=0.38$ và $0.40$.}

\answer{ex:04:1}{(a)~$\approx1.1\cdot10^{11}$ bit/J. (b)~$\log_2\sqrt{10}\approx1.7$ bit so với $81$: ít hơn $\approx50$ lần.}
\answer{ex:04:2}{$r^*=(P_0/E_{\text{xung}})\ln\bigl(1/(r^*\Delta t)\bigr)$, $P_0/E_{\text{xung}}=1.16\,\text{s}^{-1}$: $r^*\approx6$ xung mỗi giây, $7.4\cdot10^{10}$ bit/J.}
\answer{ex:04:3}{$\mathcal I=\frac{N}{1+(N-1)c}=9.9$ (giới hạn $10$) so với $\mathcal I=\frac{N}{1-c}=1111$: khác $\approx110$ lần; tương quan chỉ có hại khi nó giống tín hiệu.}
\answer{ex:04:4}{$\theta\approx68.7$ và $19.9$ (theo đơn vị $w$), $m=19$ và $m=10$: nơi nhận nhanh cần số đầu vào trùng khớp ít hơn gần một nửa, dù đầu vào trung bình của nó nhỏ hơn năm lần.}
\answer{ex:04:5}{$U\tau_{\text{ph}}\ge0.725\,\text{s}$ và $\tau_{\text{ph}}(1-2U)\le0.05\,\text{s}$, đòi hỏi $U\ge0.483$; $\tau_{\text{ph}}$ nhỏ nhất bằng $0.725\,\text{s}$ khi $U=1$ ($1.45\,\text{s}$ khi $U=0.5$).}
\answer{ex:04:6}{(a)~$\theta\,e/(e-1)\approx1.58\,\theta$ với mọi $\tau$ và $\theta$. (b)~$14$ xung.}
\answer{ex:04:7}{(a)~$1.62$ bit mỗi từ: $0.77$ --- số xung, $0.84$ --- cách sắp xếp của chúng. (b)~$1.13$ và $1.59$ bit: với $30$ từ, ước lượng mất gần một phần ba.}

\answer{ex:05:1}{$\ell\approx17\um$; $\approx100$~ngày: việc phát đi khắp não là nhờ dây dẫn phân nhánh, khuếch tán chỉ làm nhòe mỗi vị trí phóng ra khoảng $\sim20\um$.}
\answer{ex:05:2}{$r=ab/(1+G)$, $S=1/(1+G)$ chính xác; $+1.7\,\%$ thay vì $+20\,\%$.}
\answer{ex:05:3}{$T^*=1.21\,\text{s}$ khi $G=2$ và $0.19\,\text{s}$ khi $G=9$: thực tế $G\sim2$--$4$, sai lệch được nén $3$--$5$ lần.}
\answer{ex:05:4}{$\mathrm{CV}=1/\sqrt{2\lambda\tau_c}\approx9\cdot10^{-4}$ với tổng tần số $\lambda$; một nơ-ron cần $\tau_c=12.5\,\text{s}$.}
\answer{ex:05:5}{$c\approx0.40$ trong cả hai trường hợp; $\mathrm{CV}=0.050$ so với $0.158$, giảm $\sqrt{10}$ lần. Tần số của từng nơ-ron vẫn tỷ lệ với $a_i$: chỉ có tổng được điều chỉnh.}
\answer{ex:05:6}{$N_1=\overline{A\vee B}$, $N_2=\overline{A\vee N_1}$, $N_3=\overline{B\vee N_1}$, $N_4=\overline{N_2\vee N_3}$, XOR $=N_5=\overline{N_4}$; mỗi lần chuyển trạng thái mất $\tau_c\ln2$, đường qua $4$ cổng --- $2.8\,\text{s}$ cho một XOR.}
\answer{ex:05:7}{(a)~$\tau_\gamma=6/\ln3\approx5.5\,\text{s}$. (b)~$\bar D<4\,\text{s}$ so với $D<2.2\,\text{s}$: độ trễ khó đoán làm con vật kiên nhẫn hơn.}
\answer{ex:05:8}{$k_2=1/\tau_d=10\,\text{s}^{-1}$, $k_1=1/\tau_d-1/\tau_\gamma=9.5\,\text{s}^{-1}$; $\tau_\gamma=1/(k_2-mk_1)$: tăng gấp đôi khi $m$ tăng $+2.6\,\%$ (còn khi $+5.3\,\%$ thì tầm nhìn vô hạn).}

\answer{ex:06:1}{$e=(1-e^{-T_a/\tau_e})e^{-d/\tau_e}$. (a)~$0.110$ so với $4.5\cdot10^{-5}$, gấp $\approx2.5\cdot10^3$ lần. (b)~$\tau_e^*=T_a/\ln(1+T_a/d)=0.59\,\text{s}$.}
\answer{ex:06:2}{Tốc độ $\nu_x\nu_y(A_+\tau_+-A_-\tau_-)=0.8A_+$ mỗi giây, $\approx2900A_+$ mỗi giờ; $\tau_-=\tau_+/0.6=33\ms$.}
\answer{ex:06:3}{$\mathbf e_1$ khi $\mu^2+s_1^2>s_2^2$; ở đây $1.01>0.25$, và nơ-ron học được $\mathbf e_1$, dù độ phân tán theo $\mathbf e_2$ lớn hơn $25$ lần: quy tắc tìm vectơ chính của ma trận tương quan chứ không phải ma trận hiệp phương sai.}
\answer{ex:06:4}{$V=Re^{-(t_c+L-t)/\tau_\gamma}$ giữa lúc đèn sáng và lúc có nước ép, nên $\dot V=V/\tau_\gamma$; đỉnh vọt có diện tích $Re^{-L/\tau_\gamma}$: $0.61R$ và $0.78R$.}
\answer{ex:06:5}{Tỷ số tín hiệu trên nhiễu là $1/(N+1)$, số lượt thử $\propto N+1$: $5\cdot10^5$ lượt thử $\approx1$~tháng và $5\cdot10^{11}$ lượt thử $\approx8\cdot10^4$~năm.}
\answer{ex:06:6}{(a)~$k_2=1/\tau_d$, $k_1=1/\tau_d-1/\tau_\gamma$. (b)~Sai số giả $-(V/\tau_\gamma)\,\varepsilon/(1+\varepsilon)$, $\varepsilon=\tau_d/\tau_\gamma$, suy ra $\tau_d\le0.105\,\text{s}$.}
\answer{ex:06:7}{$0.46$, $0.90$, $0.99$, $0.95$ và $0.70$ khi $d=3\,\text{s}$. Các điểm nằm trên cùng một đường cong theo phần vết còn lại $F=(1-e^{-T_{\text{cue}}/\tau_e})e^{-d/\tau_e}$: học được khi $F\gtrsim0.1$, chọn ngẫu nhiên khi $F<10^{-2}$.}
\answer{ex:06:8}{$\eta^*=1/\bigl(2\sigma^2(N+2)\bigr)$, sau $N+2$ lượt thử; khi chỉ $m$ trong $N$ nơ-ron dao động --- sau $N(m+2)/m\ge N+2$: tính thưa không giúp ích.}

\answer{ex:07:1}{$V=\kappa p/\bigl(\kappa p+(1-\kappa)(1-p)\bigr)$: $0.059$, $0.20$, $0.50$; trung vị bằng $0$, còn điểm~\eqref{eq:expectile} thay đổi liên tục theo $\kappa$.}
\answer{ex:07:2}{$p=1/3$, $x=0.75$, $\sigma_r=0.35$, $V(0.2)=0.083$.}
\answer{ex:07:3}{$V=\sqrt\kappa/(\sqrt\kappa+\sqrt{1-\kappa})$, từ $0.25$ tới $0.75$; độ phân tán $\propto\sqrt\eta$; với hai giọt $V=0.25+0.5\kappa$, hai phân bố khác nhau $0.05$ ở các nơ-ron ngoài cùng, cần $\eta\lesssim0.01$.}
\answer{ex:07:4}{(a)~$J^*=2\sqrt{C\bar R}$. (b)~Mức trả dư $\bigl(k^{1/2}+k^{-1/2}\bigr)/2-1$: $6\,\%$ khi $k=2$ và $25\,\%$ khi $k=4$ --- độ chính xác ``trong phạm vi hai lần'' là đủ.}
\answer{ex:07:5}{Đỉnh vọt có diện tích $R(e^{-1}-e^{-2})\approx0.23R$ --- bằng $2.3\,\text{s}$ của đoạn dốc; nếu dopamine báo $V$ thì không có đỉnh vọt, mức chỉ nhảy bậc lên $e$ lần.}
\answer{ex:07:6}{Sự kiện dịch trọng số $\propto A\tau_f/(\tau_e+\tau_f)$, mức gây sai số $s\tau_f$: $9\,\text{s}\le\tau_f\le17\,\text{s}$.}
\answer{ex:07:7}{Ban đầu $\Delta P=0.80$, $M=0.90$. (a)~$\Delta P=0.74$ ($-8\,\%$), $M=0.43$ ($-52\,\%$); (b)~$\Delta P=0.40$ ($-50\,\%$), $M=0.80$ ($-11\,\%$): phân ly kép.}
\answer{ex:07:8}{$N_{\text{hd}}\to2+\rho^2N\approx10^6$ lượt thử: ít hơn $10^4$ lần so với một dây, nhưng rò rỉ $1\,\%$ vẫn tốn một triệu lượt thử.}

\answer{ex:08:1}{(a)~$\approx1.7\cdot10^6$. (b)~$g=\frac{0.9}{\sqrt{0.19}}\,\sigma_{\text{sau}}/\sigma_{\text{trước}}\approx16.5$: đáp số chỉ phụ thuộc vào tỷ số các nhiễu.}
\answer{ex:08:2}{(a)~Trị riêng $-(1\pm g\beta)/\tau$; hiệu tắt dần sau $\tau/(1-g\beta)$: $40\ms$ khi $g\beta=0.5$ (chênh lệch đầu vào được khuếch đại ba lần), $200\ms$ khi $0.9$ ($19$ lần). (b)~$0.905$ và $1.105$; giữa hai ranh giới chỉ đầu vào mạnh thắng.}
\answer{ex:08:3}{$P(k)=e^{gu_k/\sigma}/\sum_le^{gu_l/\sigma}$, tức là~\eqref{eq:softmax}; $g=10\ln9\approx22$.}
\answer{ex:08:4}{$\bar\Delta=1/3$, $g^*\approx3\ln(1/h)$: $21$, $14$, $9$ so với $16$--$23$, $11$ và $8$ trong mô hình.}
\answer{ex:08:5}{$g^*$ từ $16$--$23$ khi $h=0.001$ tới $6$--$8$ khi $h=0.2$; ước lượng $3\ln(1/h)$ đúng trong phạm vi một lần rưỡi khi $h\le0.05$; khi $h\gtrsim\eta/10$, quy tắc $Q\leftarrow Q+\eta(R-Q)$ không kịp tiếp thu những gì khám phá mang lại, và $g^*$ dừng lại quanh $8$.}
\answer{ex:08:6}{$T_p=\tau\ln9=44\ms$ khi $g_{\text{lo}}=0$ (mô hình $44\ms$) và $70\ms$ khi $g_{\text{lo}}=0.25$: một chùm xung dài $2$--$3\tau$ với mức khuếch đại gần bằng không.}
\answer{ex:08:7}{(a)~$g$ cố định tốt nhất $0.925$ ($g=8$), nhà tiên tri $0.929$ ($16/8$) --- hầu như không có gì để giành. (b)~Ví dụ, thời kỳ ổn định với phần thưởng hiếm ($p\in[0,0.3]$) và thời kỳ hay thay đổi với $h=0.1$: $0.850$ so với $0.872$; điều chỉnh giành được khoảng một phần tư khi $\eta=0.2$ và gần như toàn bộ khi $\eta=0.05$.}

\answer{ex:09:1}{(a)~$P_\infty=0.2$, thời gian nhớ $20\,\text{s}$. (b)~$P_\infty=1$, thời gian nhớ $1\,\text{s}$. (c)~Thời gian nhớ tỷ lệ với biên độ nhiễu: sáu lần.}
\answer{ex:09:2}{$P(t)=P_\infty\coth(t\sqrt{q/R})$. (a)~$\frac12\sqrt{R/q}=5\,\text{s}$ --- ngắn bằng một nửa thời gian nhớ. (b)~$t=\frac12\ln21\,\sqrt{R/q}\approx15\,\text{s}$: \nt{ACET} phải giữ ở mức cao chừng ấy thời gian.}
\answer{ex:09:3}{Phương sai $qT/2+R/(2T)$; $T^*=\sqrt{R/q}$, phương sai $\sqrt{qR}=P_\infty$, như ở bộ lọc~\eqref{eq:kalman}; tăng $(\kappa+1/\kappa)/2$ lần: $1.25$ và $5.05$.}
\answer{ex:09:4}{$a>a_w=1-\theta/(mw)$: $0.15$ khi $w=0.041$, $0.22$ khi $w=0.045$.}
\answer{ex:09:5}{Lỗi xuất hiện khi $M\approx40$--$60$, khi $M=80$ --- $1.9$ nơ-ron thừa, khi $M=100$ --- tỷ lệ mẫu $0.86$; nhiễu từ các mẫu khác có phương sai $\approx(M-1)p(1-p)/N$, nên $M_{\max}\approx80$, $M_{\max}\propto N/p$.}
\answer{ex:09:6}{Ví dụ, $\eta(a)=\eta\min\bigl(1,\max(0,(a-0.3)/0.6)\bigr)$. Với $\eta a$, cả ba nơ-ron lạ đều lọt vào mẫu; với $\eta(a)$ --- không nơ-ron nào, việc ghi khi $a\ge0.9$ vẫn như cũ (tỷ lệ $1.00$).}
\answer{ex:09:7}{(a)~$40$ lần phát lại; với $\eta(a)$ ở bài tập~\ref{ex:09:6} --- không bao giờ. (b)~$200$ lần phát lại, $10$ đêm.}

\answer{ex:10:1}{Bus $\approx10^{12}$ bit/s ($11.4$ bit mỗi xung), khâu điều khiển $\approx40$ bit/s; $2.5\cdot10^{10}\,\text{s}$, tức khoảng tám trăm năm.}
\answer{ex:10:2}{$a>1/3$ (với $a=0$ là dao động tăng dần $\approx67$~Hz). Không: độ ổn định do tích $g(1-a)$ quyết định.}
\answer{ex:10:3}{Trung bình $\sigma^2R'x_j$, phương sai $(N+1)\,x_j^2\sigma^4R'^2$; tỷ số $1/(N+1)$, cần $n\approx z^2(N+1)\propto N$ lần thử.}
\answer{ex:10:4}{$L_v\in[32.9,\,47.9]\ms$: độ trễ $35.4\ms$, cửa sổ $\pm7.5\ms$; số lần gắn nhầm $2\Delta_{\max}r_ar_v\approx0.38$ mỗi giây.}
\answer{ex:10:5}{$0.098$ so với $0.180$ khi dùng $\alpha=0.2$ tốt nhất: nhỏ hơn $45\,\%$.}
\answer{ex:10:6}{$n=\ln\bigl(\ln1.5/(0.6g)\bigr)/\ln(1-\alpha)$: $24$ và $4$; khi $g=2$ là $11$ và $2$.}
\answer{ex:10:7}{$S'=S\bigl(1-4\eta\sigma^2+4\eta^2\sigma^4(G+2)\bigr)$, $\eta\sigma^2=1/(2(G+2))$ là tốt nhất; $\approx10^6$ lần thử, khoảng một tháng.}

\answer{ex:11:1}{(a)~$T=k^2/\bigl(\Phi(\Delta I/I)^2\bigr)=9\,\text{s}$. (b)~$45$ cảm biến, đốm $\approx7\times7$, độ phân giải kém đi $\approx7$ lần.}
\answer{ex:11:2}{(a)~$2$--$3.3$ bit mỗi xung, $22$--$34\,\%$ của giới hạn ($9.1$--$9.8$ bit). (b)~$\log_2\binom{400}{20}\approx111$ bit; trung bình $1$ loại chung.}
\answer{ex:11:3}{(a)~$\sigma/9\le I\le9\sigma$, $1.9$ bậc. (b)~$6$ và $7$.}
\answer{ex:11:4}{(a)~$f<343/0.44\approx780$~Hz. (b)~$625$ xung, $125$ nơ-ron.}
\answer{ex:11:5}{$21$ bit mỗi sự kiện, $4.8\cdot10^5$ sự kiện/s, $7$ sự kiện trên mỗi điểm ảnh của cạnh: $136$ điểm ảnh/s. Camera thường: $1.25$ khung hình/s.}
\answer{ex:11:6}{Theo logarit $\approx1.0\,\text{s}$ cả hai chiều (lý thuyết $0.96\tau$); tuyến tính $0.2\,\text{s}$ khi lên và $3.0\,\text{s}$ khi xuống (lý thuyết $0.18\tau$ và $3.0\tau$).}
\answer{ex:11:7}{(a)~$\ln I$ có phân phối logistic với trung vị $\ln\sigma$ và thang $1$: độ tản $\pi/\sqrt3$ theo logarit tự nhiên, tức $0.79$ bậc. (b)~$r=r_{\max}\Phi_I(I)^2$.}

\answer{ex:12:1}{$0.60$ (không có tương quan thì là $0.18$), giới hạn $\sqrt{c/(rT)}\approx0.58$: một trăm nơ-ron chỉ phân biệt được ``có/không''.}
\answer{ex:12:2}{(a)~$10^8$ xung, $\approx10\,\text{mJ}$: $0.3\,\%$ của $3$~J của cả bộ não. (b)~$3$~J, lớn hơn $\approx300$ lần.}
\answer{ex:12:3}{(a)~$\theta_A\in[2,3)$, $\theta_B\in[1,2)$; hình kia cho không quá $2$ và $1$. (b)~$10\cdot96^2\approx9.2\cdot10^4$.}
\answer{ex:12:4}{(a)~$T_d=0.85\,\tau_a=0.34\,\text{s}$. (b)~$T_d\propto\tau_a$ và không phụ thuộc $I$: $\tau_a\approx3.5\,\text{s}$ (mô phỏng: $3.1\,\text{s}$).}
\answer{ex:12:5}{$p\approx0.17$. Tỷ lệ giảm chậm: $0.90$, $0.88$, $0.86$, $0.84$, $0.82$, $0.79$ khi $N=8$, $12$, $24$, $48$, $96$, $192$.}
\answer{ex:12:6}{Cả mười lần chạy không sai khi $\tau_{\text{tr}}=20$, $50$ và $200\ms$ (với $30\ms$, một trong mười lần chạy hỏng); với $5$--$10\ms$ là khoảng $0.5$, với $0.5$--$2\,\text{s}$ chất lượng giảm xuống $0.86$. Cận trên: vết phải tắt trong quãng nghỉ, $e^{-0.3\,\text{s}/\tau_{\text{tr}}}\ll1$.}
\answer{ex:12:7}{Một độ trễ là không đủ: cửa sổ $\pm5\ms$ không phủ được $\Delta x/v$ từ $5$ đến $20\ms$. Hai nhánh ức chế với $5<d_1\le10\ms$ và $15\le d_2\le d_1+10\ms$.}
\answer{ex:12:8}{\texttt{two\_stage}: $\Delta=2$, hòa khi lật một điểm, đổi câu trả lời từ hai điểm. Bộ đếm số một: một điểm lật; $0.57$ khi $p=0.1$.}

\answer{ex:13:1}{(a)~$q=100^{1/99}\approx1.048$, bước $\approx4.8\%$; dải $\approx2.2\cdot10^3$ ($67$~dB). (b)~Bước $22f_1$, tức $220\%$ ở $10f_1$.}
\answer{ex:13:2}{$P_{\text{hỏng}}\approx1.8\cdot10^{-4}$ và $2.8\cdot10^{-16}$: khoảng $300$ lần hỏng mỗi ngày khi $S=2$ và $5\cdot10^{-10}$ lần mỗi ngày khi $S=4$.}
\answer{ex:13:3}{(a)~$H(s)=eF_1T_c/(1+sT_c)^2$, bộ lọc bậc hai với tần số cắt $1/(2\pi T_c)\approx3.2$~Hz. (b)~$r\approx22$~xung/s (theo hài bậc nhất $\approx20$).}
\answer{ex:13:4}{(a)~Khi $Gd=1/e$, nghiệm $s=-1/d$ là nghiệm kép, và không có $G$ nào khác đưa nghiệm ngoài cùng bên phải sang trái xa hơn. (b)~$d=30\ms$: $G\approx12$~s$^{-1}$, $\tau=30\ms$; $d=150\ms$: $G\approx2.5$~s$^{-1}$, $\tau=150\ms$: hằng số thời gian bằng độ trễ.}
\answer{ex:13:5}{(a)~$E\approx0.427$, $T\approx1.70\,\tau_a=0.68$~s (mô hình với $\tau=20\ms$ cho $0.84$~s). (b)~Phương trình không chứa $I$; nhịp tồn tại khi $b>\beta-1$.}
\answer{ex:13:6}{Khi $b=0.8$ không có nhịp; khi $b=1.05$, $1.2$, $1.5$, $2$, $4$ thì $T\approx2.73$, $1.74$, $1.19$, $0.84$, $0.45$~s; với mọi $I$, $T=0.840$~s: đầu vào chỉ đổi biên độ, không đổi nhịp độ.}
\answer{ex:13:7}{Phép thế $s=u-a$ cho tốc độ tắt lớn nhất $a+1/d$ khi $G=e^{-1-ad}/d$; $a\ge33.3$~s$^{-1}$, $G=e^{-2}/d\approx4.5$~s$^{-1}$; đồng co càng mạnh ($a$ càng lớn) thì $G$ càng phải nhỏ.}
\answer{ex:13:8}{Bài toán mở. Ví dụ, một ngưỡng trong việc nạp độ mỏi $b[r_i-r_0]_+$ cho $I_{\min}=\beta b r_0/(b-\beta+1)\approx0.53$ khi $b=4$, $r_0=0.2$, và chu kỳ từ $1.8$~s khi $I=0.6$ xuống $0.52$~s khi $I=3$.}

\answer{ex:14:1}{(a)~$0.17$~s và $1.5$~s. (b)~Từ khoảng cách dưới $11$~cm.}
\answer{ex:14:2}{Cái chạm không kèm đau không qua được với mọi $\mu$ chừng nào $g\le\beta w_{q\mu}/w_{z\mu}=5$; cái chạm $\mu=1$ đi qua khi $g>6$: nhạy cảm hóa mạnh biến một cái chạm bình thường thành cơn đau.}
\answer{ex:14:3}{(a)~$g^*=(1-k_sC)/(1-k_sA)$, ổn định khi $k_sA<1$. (b)~Không chạm: khi $\nu\ge2$; có chạm: ngay từ $\nu\ge1.7$.}
\answer{ex:14:4}{$V(t)=-Pe^{-(T-t)/\tau_\gamma}$. (a)~$-Pe^{-T/\tau_\gamma}\approx-0.37P$. (b)~$+Pe^{-(T-t_e)/\tau_\gamma}\approx+0.61P$; nó tăng theo $t_e$ tới $+P$ và dạy mạng tránh cơn đau.}
\answer{ex:14:5}{Với $k_s=4$ không có chốt (trạng thái ổn định thứ hai có khi $k_s\ge4.58$); $T_{\min}\approx4.35$, $1.40$, $0.65$, $0.32\,\tau_s$ khi $k_s=4.6$, $5$, $6$, $8$.}
\answer{ex:14:6}{$w_{q\nu}=4/9\approx0.444$, $q_0=2/9\approx0.222$, $w_{q\mu}=49/45\approx1.089$; cái chạm không kèm đau an toàn tới $g=49/9\approx5.4$.}
\answer{ex:14:7}{(a)~Ngưỡng $\nu_c(g)=\theta/g$ khi $\nu_c\ge q_0/w_{q\nu}$, ngược lại là $(\theta+\beta q_0)/(g+\beta w_{q\nu})$; $g^*\approx2.45$, $1.0$, $0.25$. (b)~Câu hỏi mở.}

\answer{ex:15:1}{(a)~$2\cdot10^5$ mỗi nơ-ron, $6\cdot10^{12}$ cho cả sơ đồ. (b)~$\approx8$~bit/s mỗi dây, không dưới $2.5\cdot10^4$~s $\approx7$~giờ.}
\answer{ex:15:2}{Bội $2$: tương quan giữa hai bộ lọc kề nhau $0.943$, $\lambda$ từ $4.18$ đến $7.3\cdot10^{-4}$, tỷ số $\approx5.7\cdot10^3$; bội $4$: tương quan $0.8$, tỷ số $\approx94$.}
\answer{ex:15:3}{(a)~$\lambda=0.988$ và $0.808$: $82$ và $4.7$ cử động. (b)~$x_s^*=0.690$, $x_f^*=0.172$, tổng $0.862$.}
\answer{ex:15:4}{(a)~$116$ cử động. (b)~$19$. Mô hình một quá trình với tổng bằng không không cho khoản tiết kiệm nào.}
\answer{ex:15:5}{(a)~$G=50$~s$^{-1}$ so với giới hạn $10.5$~s$^{-1}$ khi không có mô hình. (b)~Không: với $G=50$~s$^{-1}$, độ trễ tới hạn $d_c\approx103\ms$; với $d=150\ms$ cần $\hat k>0.445k$ (với mọi $G$ và $d$ là $\hat k>k/2$).}
\answer{ex:15:6}{Sai số xuống dưới $0.5\%$ sau $6$--$30$~s, đáy $\approx(6$--$9)\cdot10^{-4}$, với trọng số tốt nhất là $7.5\cdot10^{-4}$; với $\eta=50$, trọng số vẫn còn lệch khỏi giá trị tốt nhất tới $0.2$ theo các hướng điều kiện kém của $C$ (bài tập~\ref{ex:15:2}).}
\answer{ex:15:7}{$\mathbf B=A\mathbf K$; $q_f/R\approx0.035$, $q_s/R\approx8.6\cdot10^{-4}$; với $4q_s$ thì $B_f\approx0.088$, $B_s\approx0.047$: quá trình chậm học nhanh hơn quá hai lần, còn quá trình nhanh học chậm hơn.}

\answer{ex:16:1}{Máu là một mạch RC với $\tau=t_{1/2}/\ln2=14.4$~phút; tỉ số cực đại trên cực tiểu là $2^{T/t_{1/2}}=64$; hài cơ bản bị suy giảm còn $0.55$; tỉ số bằng $2$ khi $t_{1/2}=T=60$~phút.}
\answer{ex:16:2}{$\varphi^*=\arcsin\bigl((\omega-\omega_0)/K_\varphi\bigr)\approx41$~phút ($\omega-\omega_0\approx0.047$~rad/ngày), thời gian hồi phục $1/(K_\varphi\cos\varphi^*)\approx3.9$~ngày; sau khi dịch $+6$ và $-6$~giờ lần lượt là $7.7$ và $7.9$~ngày.}
\answer{ex:16:3}{$K_c(n)=\sec^n(\pi/n)$: $8$ và $2.37$, sai số $11\%$ và $30\%$; khi $n\to\infty$ thì $K_c\to1$, sai số $\to50\%$.}
\answer{ex:16:4}{Các lệnh cực trị (ga $80$, phanh $0$) cho $f=120$ sau $0.76$~s; ``chỉ dùng ga'' mất $2.33$~s, lâu gấp ba.}
\answer{ex:16:5}{$3\arctan(\omega\tau)+\omega d=\pi$, $K_c=(1+\omega^2\tau^2)^{3/2}$: $K_c=8$, $3.54$, $2.50$, $1.76$ và chu kỳ $3.63$, $5.46$, $6.86$, $9.27\,\tau$; tại $0.9K_c$ dao động tắt dần, tại $1.1K_c$ dao động xác lập.}
\answer{ex:16:6}{$0.60\bar c$, $0.87\bar c$, $0.90\bar c$, giới hạn $0.91\bar c$; với nồng độ không đổi, đáp ứng bằng không: bên nhận như vậy chỉ đọc các đợt.}
\answer{ex:16:7}{Bài mở. Với phản hồi, nhịp chỉ có khi $K>8$, chu kỳ tỉ lệ với $\tau$ của các tầng và hầu như không phụ thuộc $K$ ($3.64\,\tau$ khi $K=8.5$, $4.23\,\tau$ khi $K=40$): thay đổi $K$ gấp rưỡi chỉ dịch chu kỳ $2$--$4\%$, nhỏ hơn sai số đo. Cách quyết định là mở vòng (giữ hormone cuối không đổi ở đầu vào tầng đầu sẽ triệt tiêu nhịp kiểu này) hoặc đo đáp ứng với một liều hormone ngắn bổ sung ở các pha khác nhau của chu kỳ.}

\answer{ex:17:1}{$0.2$~W so với $81$~\textmu W; luồng thô cần không quá $50$~pJ/bit.}
\answer{ex:17:2}{Khi $c=0.1$: $5.6$, $8.7$, $9.2$~bit/s, giới hạn $9.3$~bit/s; khi $c=0$ và $N=1000$: $65$~bit/s.}
\answer{ex:17:3}{$B=\log_2N+P\log_2P+(1-P)\log_2\frac{1-P}{N-1}=4.87$, $4.37$, $3.29$~bit mỗi ký tự, tức $7.3$, $6.6$, $4.9$~bit/s; để đạt $10$~bit/s cần $2.3$~ký tự mỗi giây.}
\answer{ex:17:4}{Phương sai sai số $2b/(Nm^2T)+\sigma_\xi^2$ mỗi thành phần; hai phần đóng góp bằng nhau khi $N\approx90$; giới hạn $\tfrac12\log_2(1+0.25/0.04)\approx1.4$~bit mỗi thành phần mỗi cửa sổ.}
\answer{ex:17:5}{Cặp hội tụ khi $\eta_bd^2+\eta_db^2<2$, tức khoảng $\eta<1$: với $0.8$ hội tụ, với $1.1$ dao động ổn định chu kỳ $2$, với $1.5$ không tuần hoàn, với $2$ phân kỳ; cần tách tốc độ học của hai bên.}
\answer{ex:17:6}{Ngưỡng $\approx4.4\sigma$; $102$~kbit/s, $0.10$~mW so với $205$~mW cho luồng thô, nhỏ hơn $2000$ lần; chỉ $5.7\cdot10^{-5}$ số giá trị đếm bị bão hòa (hơn $3$ xung).}
\answer{ex:17:7}{Bài mở. Khoảng $46$~bit/s khi $\text{SNR}\approx0.9$ và khoảng $330$~bit/s với nhiễu độc lập. Nếu cản trở là nhiễu giống tín hiệu, thông tin rút ra sẽ bão hòa khi số kênh tăng; nếu là do chưa hiểu mã, nó sẽ tăng với bộ giải mã tốt hơn (chẳng hạn bộ giải mã dùng thời điểm xung).}

\answer{ex:18:1}{(a)~$0.9949\le w\le1.0049$, tức $|1-w|\lesssim0.005$. (b)~$K\ge400$ khớp thần kinh.}
\answer{ex:18:2}{$(\omega_R\tau_w)^2=\sqrt{\alpha(\alpha+2+2\varrho)}/\varrho-1$, $\varrho=\tau_m/\tau_w$; $f_R\approx11.1$~Hz; $|Z(\omega_R)|/|Z(0)|=4.6$.}
\answer{ex:18:3}{(a)~$f=1/\bigl(\tau\ln\frac{\mu}{\mu-\theta}\bigr)$; để có $1$~Hz cần $\mu/\theta-1\approx3.7\cdot10^{-44}$. (b)~$T=\pi\tau/\sqrt\mu$, $f\approx3.2$~Hz: $f\propto\sqrt\mu$.}
\answer{ex:18:4}{Bộ tích phân phát xung khi $f\lesssim17.7$~Hz (bộ lọc thông thấp), bộ cộng hưởng chỉ trong dải $5.5$--$22$~Hz (bộ lọc thông dải).}
\answer{ex:18:5}{(a)~$T_{\min}=\frac{\tau}{w-1}\ln\frac{0.5}{0.3}\approx10.2\ms$. (b)~$B>w-1-0.2=0.3$.}
\answer{ex:18:6}{Hằng số thời gian của việc điều chỉnh là $\tau/(\eta\langle r^2\rangle)$, $\eta=\tau\ln10/60\,\text{s}\approx3.8\cdot10^{-3}$; khi $I\neq0$ trọng số bị lệch, cần trừ khỏi $e$ bản sao của lệnh $I/\tau$.}
\answer{ex:18:7}{Bài mở. Bộ cộng hưởng chỉ hơn $1.35$ lần ở $11$~Hz và thua $17$ lần ở $1$~Hz; lợi ích chính của việc tái cấu hình nằm ở việc chọn dải tần theo ngưỡng (bài~\ref{ex:18:4}).}

\answer{ex:19:1}{(a)~Một khớp cho $6.7\mV$ ($R=66.7$~M$\Omega$), nên để tới ngưỡng nói chung cần ít nhất $4$ khớp (ba khớp chỉ đủ theo tiệm cận); $8$ khớp để tới trong $10\ms$; $14$ khớp để tới trong $5\ms$. (b)~$0.1\ms\ll\tau_m$, hiệu chỉnh do rò $\sim0.25\%$.}
\answer{ex:19:2}{(a)~$1.75\cdot10^{10}$ (một phần tư số bộ trộn đáp ứng với mọi thứ), $4.4\cdot10^9$ và $0.06$ (với $k=40$ hầu như không bao giờ đáp ứng). (b)~$10^3$ lần: $2\cdot10^5$ so với $200$.}
\answer{ex:19:3}{(a)~$C(2n,n)=2^{2n-1}$ theo tính đối xứng của hệ số nhị thức. (b)~$0.93$ và $0.09$; độ rộng vùng chuyển tiếp theo $P$ là $\sim\sqrt{2n}$, tương đối là $\sim1/\sqrt n$.}
\answer{ex:19:4}{Từ một đuôi gai, $V_{\text{thân}}<\kappa E_E<\theta$ với bất kỳ số đầu vào nào; với $g_0=0.5g_L$ cần $n\ge3$ trên mỗi đuôi gai.}
\answer{ex:19:5}{$I\ge C\sigma_V/\sigma_t=15$~nA, tức $150$ khớp đồng bộ, không thực tế; nơ-ron nhỏ chỉ cần $1$~nA, và với một khớp $4$~nA độ rung pha là $5$~\textmu s.}
\answer{ex:19:6}{Cực đại ở đầu gần: $0.03\,\tau_m$ và $0.97$ ($L=0.2$); $0.32\,\tau_m$ và $0.66$ ($L=1$); $0.79\,\tau_m$ và $0.33$ ($L=2$). Ước tính~\eqref{eq:dendriteload} theo điện dung toàn phần chỉ đúng khi $L\ll1$.}
\answer{ex:19:7}{Bài mở. Khi $m$ cố định, thời gian đáp ứng tăng tuyến tính theo $P$ được ghi nhớ; khi tỉ lệ $m=\varphi n$ cố định, nó không còn phụ thuộc $n$, và sự chậm chạp của nơ-ron lớn là hệ quả của việc cộng nhiều đầu vào chứ không phải của kích thước tự thân.}

\answer{ex:20:1}{$t_{\text{thức}}=\tau_r\ln\frac{1-L}{1-H}=16.8$~giờ, $t_{\text{ngủ}}=\tau_d\ln\frac HL=5.8$~giờ, chu kỳ $22.5$~giờ; dao động ngày đêm của các ngưỡng kéo bộ dao động về $24$~giờ, đó là khóa pha~\eqref{eq:pll}.}
\answer{ex:20:2}{(a)~$L=0.111$; với $L=0.092$ giấc ngủ sẽ kéo dài $9.6$~giờ. (b)~Tăng $22\%$ ($1/0.82=1.22$), chứ không phải $18\%$.}
\answer{ex:20:3}{$H=\frac{p-1}{p-1/q}=0.623$, $L=H/q=0.093$, trong đó $p=e^{16/\tau_r}$, $q=e^{8/\tau_d}$. Sau khi bị đánh thức sau $4$~giờ, pha bị dịch vĩnh viễn (ngủ sớm hơn $7.2$~giờ); với ngưỡng dao động, pha trở lại sau hai đến ba ngày.}
\answer{ex:20:4}{(a)~$r_\pm=\frac12\bigl(1\pm\sqrt{1-4k/w}\bigr)=0.947,\ 0.053$; $a_{\text{trên}}=3.51$, $a_{\text{dưới}}=0.49$. (b)~Làm việc $\approx0.33$~s, im lặng $\approx0.59$~s, chu kỳ $\approx0.93$~s, tỉ lệ làm việc $0.36$.}
\answer{ex:20:5}{$a_{\text{trên}}/r_+<b<a_{\text{dưới}}/r_-$: $3.71<b<9.20$. \nt{ACET} làm $b$ giảm xuống dưới $3.71$: giao điểm chuyển lên nhánh trên, và trạng thái làm việc ($r\approx1$) trở nên ổn định.}
\answer{ex:20:6}{$4.9$, $6.8$, $8.8$, $5.5$, $8.9$~giờ với $D=24$, $32$, $40$, $48$, $64$: độ dài do pha ngày đêm lúc đi ngủ quyết định, chứ không phải do món nợ.}
\answer{ex:20:7}{Sau B, sai số trên A trung bình là $0.51$ (từ $0.19$ đến $1.08$ qua $20$ bộ dữ liệu; đáp ứng bằng không cho $1$). Khi trộn lẫn, cả hai sai số đều nhỏ hơn $10^{-4}$.}
\answer{ex:20:8}{Bài mở. Co giãn đều giữ nguyên tỉ lệ giữa các trọng số, nhưng chỉ giữ nguyên đáp ứng của nơ-ron có ngưỡng khi ngưỡng cũng được co giãn cùng tỉ lệ; phát lại trộn lẫn thì làm thay đổi tỉ lệ. Việc các tiếp xúc lớn không đổi mâu thuẫn với sự co giãn hoàn toàn đều.}

\answer{ex:21:1}{(a)~$T=\tau_{\text{ph}}\ln\frac{1}{1-x^*}$: $1.84$~s và $3.68$~s. (b)~Từ $3.36$ đến $4.24$~s: độ nhạy $\dd T/\dd x^*=\tau_{\text{ph}}/(1-x^*)=80$~s, khi $x^*\to1$ các đợt bùng phát không đều.}
\answer{ex:21:2}{(a)~$8.6$~W, như trong~\eqref{eq:energy}. (b)~$205$~Mbit/s, $0.2$~W, gấp $2\cdot10^3$ lần mẫu nuôi cấy ($10^{-4}$~W).}
\answer{ex:21:3}{$x_{\text{dừng}}=1/(1+Ur_b\tau_{\text{ph}})<x^*$ khi $r_b>(1/x^*-1)/(U\tau_{\text{ph}})=0.28$~Hz với $x^*=0.9$ và $0.025$~Hz với $x^*=0.99$; sức mạnh khớp giảm xuống $x_{\text{dừng}}$, tức giảm $10\%$ và $1\%$.}
\answer{ex:21:4}{Các $u(t-k)$ trực chuẩn, $m_k=\|P_Vu(t-k)\|^2$, trong đó $P_V$ là phép chiếu lên bao tuyến tính $N$ chiều của các đầu ra; theo bất đẳng thức Bessel, $\sum_k m_k\le N$.}
\answer{ex:21:5}{$\mathrm{MC}\approx9$--$11$ khi có $\tanh$ và $15$--$18$ với hồ chứa tuyến tính (ba hạt giống ngẫu nhiên), cả hai đều $\le30$: tính phi tuyến phải trả giá bằng trí nhớ.}
\answer{ex:21:6}{(a)~$n\le102$. (b)~$5.6$ sigma.}
\answer{ex:21:7}{Bài mở. Đối chứng then chốt: đóng băng lớp đọc ra và kiểm tra xem sự cải thiện có còn sau một khoảng nghỉ không kích thích hay không; nếu không, cái thay đổi là trạng thái, không phải trọng số.}

\answer{ex:22:1}{(a)~$32\%$ và $100\%$. (b)~$R_\uparrow-R_\downarrow\ge2\sqrt{2000\,\text{Hz}/1\,\text{s}}\approx89$~Hz, chỉ $4.5\%$ tần số tổng.}
\answer{ex:22:2}{(a)~$n\ge4(p_1q_1+p_2q_2)/\Delta^2\approx560$. (b)~$\approx56$.}
\answer{ex:22:3}{(a)~Cân bằng chi tiết: $\rho PK$ đối xứng. (b)~Tỷ lệ trượt bằng trung bình điều hòa $1/\langle1/P\rangle=0.18$, so với $0.5$ khi không có phản hồi.}
\answer{ex:22:4}{(a)~Hình bình hành diện tích $4h^2=0.04$ (tính cả giới hạn của $p$ thì bằng số là $0.045$). (b)~$\approx0.18$.}
\answer{ex:22:5}{$0.49$: một phần mẫu nuôi cấy trôi vào vùng gần như luôn trượt và lang thang ở đó. Có giới hạn thì $1.00$: trên một tập bị chặn, mật độ $\propto1/P$ dồn vào vùng hấp thụ.}
\answer{ex:22:6}{(a)~Cần $\ge5$ mức; $8$ điện cực là đủ. (b)~Không: cần $r_{\max}\ge25/T=100$~Hz.}
\answer{ex:22:7}{Bài mở. Thời gian tìm kiếm ngẫu nhiên tăng theo $d$ xấp xỉ như tỷ số thể tích $(L/h)^d$, còn tìm kiếm có dấu của sai số thì tăng tuyến tính theo $d$. Thí nghiệm hoán đổi sẽ phân biệt hai giả thuyết: sau cú trượt là kích thích đoán trước được nhưng mạnh, sau cú trúng là kích thích không đoán trước được nhưng yếu; cần vài chục mẫu nuôi cấy mỗi nhóm.}


\chapter{Phụ lục thực nghiệm}
\label{sec:supplement}

Với mỗi chương, ở đây tập hợp các khẳng định chính của chương và cơ sở của chúng: thí nghiệm đã đo điều đó, cụ thể đo được gì và được công bố ở đâu. Trạng thái ở cột bên phải cho biết khẳng định vững đến mức nào: \emph{đã đo} --- có thí nghiệm trực tiếp; \emph{lý thuyết} --- hệ quả toán học, được suy ra trong chương hoặc trong một công trình kinh điển; \emph{mô hình} --- kết quả tính toán theo mô hình của sách (các script trong thư mục \texttt{models/}); \emph{ước tính} --- con số theo bậc độ lớn, không có phép đo trực tiếp; \emph{giả thuyết} --- cách giải thích hợp lý nhưng chưa được thí nghiệm chứng minh. Chỗ nào không có thí nghiệm, chúng tôi nói rõ như vậy.

\section{Chương~\ref{sec:neurontypes}. Các loại nơ-ron}

Các con số của chương dựa trên những phép đếm tế bào trực tiếp trong não người ở những nơi có phép đếm như vậy; quy tắc ``một nơ-ron --- một chất dẫn truyền'' dựa trên các bản đồ tế bào và những thí nghiệm đã tìm ra cả ngoại lệ; vai trò của \nt{DOPA} dựa trên các bản ghi xung; còn vai trò của các kênh quảng bá khác dựa trên giả thuyết.

{\footnotesize
\setlength{\tabcolsep}{3pt}\renewcommand{\arraystretch}{1.15}
\begin{longtable}{>{\raggedright}p{0.5cm}>{\raggedright}p{4.0cm}>{\raggedright}p{5.6cm}>{\raggedright}p{2.3cm}>{\raggedright\arraybackslash}p{1.7cm}}
\toprule
TT & Khẳng định & Thí nghiệm và điều đã đo & Nguồn & Trạng thái \\
\midrule\endfirsthead
\toprule
TT & Khẳng định & Thí nghiệm và điều đã đo & Nguồn & Trạng thái \\
\midrule\endhead
1 & Não người có khoảng $8.6\cdot10^{10}$ nơ-ron, và gần như mọi nơ-ron \nt{GLUT} là nơ-ron nhỏ của tiểu não (bảng~\ref{tab:types}) & Phương pháp phân đoạn đẳng hướng: mô được hòa tan thành huyền phù nhân đồng nhất, rồi đếm các nhân mang dấu ấn nơ-ron NeuN. Ở nam giới trưởng thành có $86.1\pm8.1$ tỷ nơ-ron; trong vỏ đại não chỉ có $19\,\%$, phần lớn nằm ở tiểu não & \cite{Azevedo2009} & đã đo \\
2 & Hầu như mỗi nơ-ron có một chất dẫn truyền chính (mệnh đề~\ref{prop:dale}), nhưng có ngoại lệ & Bản đồ phiên mã não chuột: khoảng $4\cdot10^6$ tế bào, $5322$ cụm; mỗi cụm được gán chất dẫn truyền theo các gen tổng hợp và đóng gói. Ngoại lệ được đo trực tiếp: trong lát cắt, nơ-ron \nt{DOPA} nhả cả GABA qua cùng chất vận chuyển túi và ức chế nơ-ron thể vân; ở một phần sợi trục \nt{DOPA}, glutamate và dopamine đi ra từ những đầu tận cùng khác nhau của cùng một dây (kính hiển vi điện tử và quang di truyền) & \cite{Strata1999,Yao2023,Tritsch2012,Zhang2015} & đã đo \\
3 & Cả cột ma trận trọng số cùng một dấu~\eqref{eq:dalesign} & Suy luận trong chương từ mệnh đề~\ref{prop:dale} và từ việc gần như mọi thụ thể glutamate là kích thích, còn thụ thể GABA là ức chế. Chưa có kiểm chứng trực tiếp ``mọi nơi nhận của một nơ-ron đều nhận trọng số cùng dấu'' như một quy tắc chung; phản ví dụ là việc nơ-ron \nt{DOPA} nhả kèm GABA (dòng~2) và những nơi nhận có thụ thể trái dấu (dòng~11) & suy luận trong chương & lý thuyết \\
4 & Từ ba phần tư đến bốn phần năm nơ-ron vỏ não là \nt{GLUT}, từ một phần năm đến một phần tư là \nt{GABA} & Nhuộm miễn dịch GABA trong các cột rộng $50\um$ xuyên suốt độ dày vỏ não ở $10$ vùng vỏ não khỉ: ở $8$ vùng, nơ-ron GABA chiếm khoảng $25\,\%$ tổng số nơ-ron, ở vỏ thị giác sơ cấp --- dưới $20\,\%$ & \cite{Hendry1987} & đã đo \\
5 & Nơ-ron \nt{GLUT} có khoảng $10^4$ đầu ra & Lập thể học trên tử thi: vỏ não mới của nam thanh niên có $1.64\cdot10^{14}$ khớp thần kinh (kính hiển vi điện tử, disector) và khoảng $2.3\cdot10^{10}$ nơ-ron. Tỷ số cho $\sim7\cdot10^3$ khớp thần kinh trên một nơ-ron; tính trung bình, số đầu vào bằng số đầu ra & \cite{Tang2001synapses,Pakkenberg1997} & đã đo \\
6 & Đầu ra của nơ-ron \nt{DOPA} phân nhánh thành $10^5$--$10^6$ đầu tận cùng; đây là ước tính theo độ phủ vùng đích chứ không phải phép đếm & Đánh dấu màng bằng virus ở từng nơ-ron \nt{DOPA} của chuột cống và tái dựng toàn bộ sợi trục: một nơ-ron phủ từ $0.45$ đến $5.7\,\%$ (trung bình $2.7\,\%$) thể tích thể vân. Đại lượng đo được là độ phủ chứ không phải số đầu tận cùng: số này được suy từ độ phủ theo bậc độ lớn, chưa có phép đếm trực tiếp. Với \nt{SERO} và \nt{NORA}, số đầu tận cùng là ước tính thô & \cite{Matsuda2009} & đã đo \\
7 & Nơ-ron quảng bá rất ít: \nt{DOPA} $\sim5\cdot10^5$ (nhóm chính trong hai nhóm, cận dưới), \nt{SERO} $\sim3\cdot10^5$ (tổng của hai phép đếm từng phần), \nt{HIST} $\sim6\cdot10^4$ (bảng~\ref{tab:types}, hình~\ref{fig:typepie}) & Các phép đếm lập thể học ở người: $550\,000$ nơ-ron sắc tố trong chất đen (không tính vùng trần não thất); $235\,000$ nơ-ron trong nhân đường giữa lưng (mọi nơ-ron của nhân) và thêm $125\,000$ nơ-ron serotonin ở trần cầu não; khoảng $64\,000$ nơ-ron histamine ở vùng dưới đồi. Với \nt{NORA}, \nt{GLYC}, \nt{ACET} và \nt{ADRE} không có phép đếm trực tiếp: trong bảng đó là ước tính thô theo bậc độ lớn & \cite{Pakkenberg1991,Baker1990,Baker1991,Airaksinen1991} & đã đo \\
8 & Glutamate trong khe vọt lên khoảng một milimol và giảm trong $\sim1\ms$ & Qua sự đẩy một chất chẹn cạnh tranh tách ra nhanh khỏi thụ thể NMDA trong lúc truyền qua khớp thần kinh (mẫu nuôi cấy nơ-ron hồi hải mã): đỉnh $1.1\mM$, hằng số suy giảm $1.2\ms$ & \cite{Clements1992} & đã đo \\
9 & Trong vỏ não đang làm việc, dòng kích thích và dòng ức chế gần như cân bằng, nơ-ron nằm sát ngưỡng & Lý thuyết: mạng có liên kết mạnh tự đi vào chế độ cân bằng. Thực nghiệm: ở vỏ não trước trán của chồn sương in vivo, trong các trạng thái ``lên'', điện thế đảo chiều của dòng qua khớp thần kinh giữ quanh $-37\mV$ suốt hàng trăm mili giây, dù độ dẫn thay đổi $\sim21$\,nS: kích thích và ức chế cùng tăng và cùng giảm & \cite{vanVreeswijk1996,Haider2006} & đã đo \\
10 & Nơ-ron \nt{DOPA} báo sai số dự đoán phần thưởng~\eqref{eq:rpe} (hình~\ref{fig:rpe}) & Xung của nơ-ron \nt{DOPA} ở khỉ: chùm xung với nước ép bất ngờ; sau khi học --- chùm xung với tín hiệu báo trước và không đáp ứng với nước ép đã được dự đoán; tần số sụt khi nước ép đã hứa không được cho. Trong thí nghiệm định lượng, tần số tỷ lệ với hiệu giữa phần thưởng hiện tại và trung bình có trọng số của các phần thưởng trước, nhưng chỉ với sai số dương. Bản thân phương trình~\eqref{eq:rpe} là thuật toán sai phân thời gian & \cite{Schultz1997,Bayer2005,SuttonBarto2018} & đã đo \\
11 & Dấu và tốc độ do thụ thể quyết định: thụ thể hướng ion --- phần nhỏ của mili giây, thụ thể hướng chuyển hóa --- chậm hơn nhiều (mệnh đề~\ref{prop:receptor}) & Xung glutamate ngắn lên các mảnh màng của nơ-ron hồi hải mã: dòng qua thụ thể hướng ion tăng (từ $20$ đến $80\,\%$) trong $0.2$--$0.6\ms$ và giảm trong $2$--$3\ms$. Điện thế ức chế chậm của nơ-ron tháp bị xóa bởi chất chẹn chọn lọc thụ thể GABA hướng chuyển hóa. Hai họ thụ thể dopamine có tác động ngược nhau; ở thể vân, nơ-ron có thụ thể D1 cần dopamine để tăng cường khớp thần kinh, nơ-ron có D2 --- để làm yếu khớp thần kinh & \cite{Colquhoun1992,DutarNicoll1988,Beaulieu2011,Shen2008} & đã đo \\
12 & Kênh quảng bá không phải một dây dẫn mà là một bó ít đường (mục~\ref{sec:buses}) & Đánh dấu di truyền bằng virus các nơ-ron serotonin của nhân đường giữa lưng ở chuột nhắt: nơ-ron gửi đầu ra tới hạch hạnh nhân và tới vỏ não trán nằm ở những chỗ khác nhau trong nhân, phân nhánh vào những vùng bổ sung cho nhau và đáp ứng ngược chiều với kích thích khó chịu. Tổng quan: các phân nhóm nơ-ron của nhân lục (\nt{NORA}) mã hóa những quá trình khác nhau & \cite{Ren2018,Poe2020} & đã đo \\
13 & \nt{NORA} đặt hệ số khuếch đại $g$ & Giả thuyết khuếch đại thích nghi; mô hình mạng trong đó catecholamine làm tăng độ dốc của đặc tuyến truyền đạt. Chưa có phép đo trực tiếp ``$g$ tăng theo noradrenaline'' trên toàn mạng; đã đo được rằng đường kính đồng tử bám theo xung của nơ-ron nhân lục ở khỉ & \cite{AstonJones2005,ServanSchreiber1990,Joshi2016} & giả thuyết \\
14 & \nt{ACET} chuyển ``ghi/đọc'' và đặt tốc độ học & Giả thuyết. Cơ sở: trong lát cắt vỏ khứu giác chuột cống, chất chủ vận acetylcholine ($100\,\mu\text{M}$) làm yếu các khớp thần kinh bên trong, ``hồi tưởng'', hơn $60\,\%$, còn các khớp thần kinh đầu vào --- dưới $15\,\%$ & \cite{Hasselmo2006,HasselmoBower1992} & giả thuyết \\
15 & \nt{SERO} đặt hệ số chiết khấu $\gamma$; nơ-ron serotonin làm việc đều, vài xung mỗi giây, và gần như im lặng trong một pha của giấc ngủ & Giả thuyết ``serotonin $=\gamma$''. Lập luận mạnh nhất: hoạt hóa quang di truyền các nơ-ron serotonin của nhân đường giữa lưng ở chuột nhắt làm giảm số lần bỏ chờ sớm phần thưởng trì hoãn. Tần số xung và sự im lặng trong giấc ngủ có cử động mắt nhanh đã được đo (tổng quan các bản ghi) & \cite{Doya2002,Miyazaki2014,JacobsAzmitia1992} & giả thuyết \\
\bottomrule
\end{longtable}}

\section{Chương~\ref{sec:glutcomputer}. Nơ-ron \nt{GLUT} tính được gì}

Các khẳng định logic của chương là những định lý về mạch có trọng số không âm, được suy ra ngay trong chương; thực nghiệm xác nhận mô hình nơ-ron và việc các mạch được lắp theo những định lý đó (bộ phát hiện trùng khớp, đường trễ, nhóm tự duy trì, chuỗi) thực sự có trong não.

{\footnotesize
\setlength{\tabcolsep}{3pt}\renewcommand{\arraystretch}{1.15}
\begin{longtable}{>{\raggedright}p{0.5cm}>{\raggedright}p{4.0cm}>{\raggedright}p{5.6cm}>{\raggedright}p{2.3cm}>{\raggedright\arraybackslash}p{1.7cm}}
\toprule
TT & Khẳng định & Thí nghiệm và điều đã đo & Nguồn & Trạng thái \\
\midrule\endfirsthead
\toprule
TT & Khẳng định & Thí nghiệm và điều đã đo & Nguồn & Trạng thái \\
\midrule\endhead
1 & Nơ-ron là mạch RC có ngưỡng và đặt lại~\eqref{eq:glutneuron} & Bơm vào thân nơ-ron tháp vỏ não chuột cống (lát cắt) một dòng nhiễu giống dòng qua khớp thần kinh trong não sống. Sự phụ thuộc của tần số trung bình vào giá trị trung bình và độ phân tán của dòng được mô tả bằng nơ-ron tích phân rò nếu thêm sự thích nghi tần số chậm. Các khoảng giá trị của $\tau$ và độ trễ được nêu trong chương mà không dẫn thí nghiệm & \cite{Rauch2003} & đã đo \\
2 & Mô hình~\eqref{eq:glutneuron} là một sự đơn giản hóa: chính các nhánh đầu vào phát xung cục bộ & Ghi trực tiếp từ đuôi gai của nơ-ron tháp vỏ thị giác chuột nhắt in vivo: kích thích thị giác gây ra xung đuôi gai cục bộ phụ thuộc thụ thể NMDA; ức chế chúng làm nới rộng độ chọn lọc hướng của nơ-ron & \cite{Smith2013} & đã đo \\
3 & Cửa sổ trùng khớp $\Delta_{\max}=\tau\ln\frac{w}{\theta-w}$~\eqref{eq:window}; khi $\theta=1.5w$ nó bằng $0.7\tau$ & Hệ quả của mô hình~\eqref{eq:glutneuron}. Phép đo trực tiếp một cửa sổ cộng hẹp có ở nơ-ron có ức chế nhanh (chương~\ref{sec:gaba}, dòng~3 trong bảng của chương đó) & suy luận trong chương & lý thuyết \\
4 & Mạng chỉ gồm nơ-ron \nt{GLUT} chỉ tính được các hàm đơn điệu của đầu vào (mệnh đề~\ref{prop:monotone}) & Chứng minh trong chương: lặp từ trạng thái im lặng với vế phải không giảm. Đây là khẳng định về mô hình; chưa có thí nghiệm kiểm chứng nó trên mạng sống không có ức chế & suy luận trong chương & lý thuyết \\
5 & Giác quan cho ra một cặp dây: một số tế bào đáp ứng khi kích thích tăng, số khác khi kích thích giảm & Võng mạc: ngay từ tế bào lưỡng cực, tín hiệu của tế bào nón đã tách thành kênh bật và kênh tắt. Chẹn dược lý các tế bào lưỡng cực bật ở khỉ: các kênh đi riêng tới tận vỏ não; sau khi chẹn, khả năng phát hiện ánh sáng mạnh lên bị suy giảm, còn phát hiện ánh sáng yếu đi thì không & \cite{Schiller1992} & đã đo \\
6 & Với mã hai dây, mạng nơ-ron \nt{GLUT} tính mọi hàm logic một cách bất đồng bộ, không có nhịp chung, với cái giá là số nơ-ron gấp đôi (mệnh đề~\ref{prop:dualrail}, \eqref{eq:dualrail}, hình~\ref{fig:dualxor}) & Mã hai dây và việc xác định trạng thái sẵn sàng qua chính tín hiệu là kỹ thuật tiêu chuẩn của mạch bất đồng bộ không nhạy với độ trễ. Mạch XOR sáu nơ-ron được lắp trong chương. Chưa có thí nghiệm cho thấy não tính các hàm logic đúng bằng mã hai dây & \cite{Sparso2001} & lý thuyết \\
7 & Hiệu thời gian âm thanh đến hai tai lớn nhất ở người là $\approx0.6\ms$, còn não phân biệt được khoảng mười micro giây & Người nghe trong thí nghiệm hai phương án trả lời: ngưỡng phân biệt hiệu thời gian (75\,\% đúng) là $9\,\mu\text{s}$ với nhiễu dải $150$--$1700$\,Hz, $11\,\mu\text{s}$ với âm $1$\,kHz, $28\,\mu\text{s}$ với tiếng tách. Giới hạn $0.6\ms$ là phép tính của chương, $0.22\,\text{m}/343\,\text{m/s}$ & \cite{KlumppEady1956} & đã đo \\
8 & Ở chim, hiệu thời gian được biến thành vị trí nhờ đường trễ và bộ phát hiện trùng khớp~\eqref{eq:itd} (hình~\ref{fig:itd}) & Cú lợn lưng xám: sợi trục từ hai tai đi vào nhân của các bộ phát hiện trùng khớp từ hai phía ngược nhau; thời điểm xung đến dọc sợi trục thay đổi có trật tự theo độ sâu, và dải độ trễ qua nhân ($160\,\mu\text{s}$ trên $700\um$) trùng với dải độ trễ giữa hai tai & \cite{Carr1990} & đã đo \\
9 & Ở động vật có vú không tìm thấy hàng độ trễ; cùng bài toán được giải bằng bộ phát hiện trùng khớp với ức chế đúng thời điểm từ nơ-ron \nt{GLYC} & Ghi in vivo ở nhân trám trên giữa của chuột nhảy: đáp ứng không tạo thành bản đồ hướng; đưa glycine và chất chẹn nó qua vi pipet cho thấy sự ức chế \emph{glycine} đến đúng thời điểm quy định dải độ trễ làm việc. Ức chế ở đây là từ \nt{GLYC} chứ không phải từ \nt{GABA} & \cite{Brand2002,Grothe2010} & đã đo \\
10 & Nhóm có hồi tiếp mạnh là một flip-flop; hoạt động tự duy trì kéo dài hàng giây sau kích thích (hình~\ref{fig:bistable}) & Vỏ não trước trán của khỉ trong nhiệm vụ đưa mắt trì hoãn tới một điểm được ghi nhớ (trì hoãn $1$--$6$\,s): $87$ trên $170$ nơ-ron liên quan đến nhiệm vụ thay đổi tần số trong thời gian trì hoãn, ở $79\,\%$ trong số đó điều này phụ thuộc vào hướng tới điểm được ghi nhớ. Việc hoạt động này được duy trì bởi hồi tiếp có hai trạng thái bền là lý thuyết & \cite{Funahashi1989,Wang2001} & đã đo \\
11 & Bit được ghi nếu dòng vào kéo dài hơn $\approx0.7\tau$ (hình~\ref{fig:recurrentgroup}b) & Giải phương trình $\tau\,\dd r/\dd t=-r+f(wr+I)$ bằng phương pháp Euler với $w=1$, $I=0.6$; không có script trong \texttt{models/}. Chưa có thí nghiệm & tính toán trong chương & mô hình \\
12 & Bộ nhớ có thể lưu trong sự tạo thuận của khớp thần kinh, gần như không cần xung & Lý thuyết: calci dư trong đầu tận cùng giữ bản ghi khoảng một giây. Cơ sở: trong vỏ não trước trán giữa của chồn sương (ghi đồng thời nhiều nơ-ron) có một mạng con nơ-ron tháp nối bằng khớp thần kinh tạo thuận có hậu hiệu ứng kéo dài. Tổng quan các quan sát về những khoảng ``lặng'' khi giữ thông tin trong trí nhớ. Vỏ não dùng cách nào khi nào là câu hỏi còn mở & \cite{Mongillo2008,WangMarkram2006,Stokes2015} & giả thuyết \\
13 & Bộ nhớ bị xóa bởi sự mỏi: dự trữ chất dẫn truyền cạn kiệt & Ghi cặp nơ-ron tháp vỏ não: đáp ứng của khớp thần kinh giảm khi xung dày, tốc độ giảm được quy định bởi xác suất nhả. Chưa có bằng chứng trực tiếp rằng chính điều này dập tắt nhóm tự duy trì & \cite{Tsodyks1997} & đã đo \\
14 & Chuỗi các nhóm là bộ định thời khởi động theo sự kiện; ở chim biết hót, nơ-ron phát xung lần lượt nghiêm ngặt & Ghi các nơ-ron đã được nhận dạng ở trung tâm phát âm cao cấp của chim sẻ vằn khi ngủ và khi hót: mỗi nơ-ron gửi đầu ra tới nhân vận động phát một chùm xung ngắn vào một thời điểm chính xác của bài hót, và các nơ-ron phát xung lần lượt & \cite{Hahnloser2002} & đã đo \\
\bottomrule
\end{longtable}}

\section{Chương~\ref{sec:gaba}. \nt{GLUT} và \nt{GABA} cùng nhau}

Các khẳng định logic và động lực học của chương được suy ra ngay trong chương từ phân tích tuyến tính và định luật Ohm; thực nghiệm cho thấy các mạch làm nền cho chúng --- phép cấm có độ trễ, ức chế thuận theo sau kích thích, sự ổn định nhờ ức chế, nhịp của vòng \nt{GLUT}--\nt{GABA} --- thực sự hoạt động trong não.

{\footnotesize
\setlength{\tabcolsep}{3pt}\renewcommand{\arraystretch}{1.15}
\begin{longtable}{>{\raggedright}p{0.5cm}>{\raggedright}p{4.0cm}>{\raggedright}p{5.6cm}>{\raggedright}p{2.3cm}>{\raggedright\arraybackslash}p{1.7cm}}
\toprule
TT & Khẳng định & Thí nghiệm và điều đã đo & Nguồn & Trạng thái \\
\midrule\endfirsthead
\toprule
TT & Khẳng định & Thí nghiệm và điều đã đo & Nguồn & Trạng thái \\
\midrule\endhead
1 & Nơ-ron \nt{GABA} chiếm từ một phần năm đến một phần tư số nơ-ron vỏ não & Nhuộm miễn dịch GABA ở $10$ vùng vỏ não khỉ: khoảng $25\,\%$ nơ-ron ở $8$ vùng, dưới $20\,\%$ ở vỏ thị giác sơ cấp. Tổng quan về sự đa dạng của nơ-ron ức chế vỏ não & \cite{Hendry1987,Markram2004} & đã đo \\
2 & Ức chế làm gì phần lớn do vị trí đậu quyết định: đuôi gai, thân, chỗ sinh xung, đầu tận cùng của dây dẫn khác & Tổng quan: các lớp nơ-ron \nt{GABA} vỏ não khác nhau ở đích trên nơi nhận. Ghi đồng thời từ thân và đuôi gai của nơ-ron tháp: ức chế thuận mạnh hơn nhiều ở thân so với ở đuôi gai. Ức chế trên đầu tận cùng của dây dẫn khác ở tủy sống làm giảm lượng chất dẫn truyền được nhả của một đường vào (tổng quan các phép đo) & \cite{Markram2004,Pouille2001,Rudomin1999} & đã đo \\
3 & Đổi dấu qua trung gian tốn một độ trễ, khoảng một mili giây; ức chế theo sau kích thích thu hẹp cửa sổ trùng khớp & Nơ-ron tháp hồi hải mã chuột cống trong lát cắt: ức chế qua một trung gian đến theo sau kích thích trực tiếp với độ trễ ngắn, và cửa sổ trong đó hai đầu vào kích thích cộng lại tới ngưỡng nhỏ hơn $2\ms$. Ở đuôi gai, nơi ức chế này yếu hơn, cửa sổ rộng hơn & \cite{Pouille2001} & đã đo \\
4 & Phép cấm $a\wedge\bar b$~\eqref{eq:veto} cùng với OR và hằng một là đầy đủ về chức năng (mệnh đề~\ref{prop:gabacomplete}) & Chứng minh trong chương. Kết quả kinh điển: mạng phần tử ngưỡng có đầu vào ức chế hiện thực được mọi biểu thức logic. Khẳng định về mô hình, chưa có thí nghiệm & \cite{McCullochPitts1943} & lý thuyết \\
5 & Phép cấm có độ trễ cho bộ phát hiện hướng chuyển động với vận tốc tối ưu $v^*=\Delta x/d$~\eqref{eq:vstar} & Võng mạc thỏ: tính chọn lọc hướng được tạo ra bởi ức chế, dập tắt kích thích khi chuyển động theo hướng ``không''. Ghi riêng đầu vào kích thích và ức chế của các tế bào chọn lọc cho thấy tế bào amacrine hình sao (\nt{GABA}) ức chế chúng không đối xứng --- chỉ bằng các nhánh hướng về phía ``không''. Công thức $v^*$ là suy luận của chương & \cite{Barlow1965,Fried2002} & đã đo \\
6 & Ức chế sun chia điện áp chứ không trừ~\eqref{eq:shunt} (hình~\ref{fig:shunt}) & Định luật Ohm cho bộ chia gồm ba độ dẫn, với điện thế cân bằng của clo gần mức nghỉ; cho nơ-ron không phát xung & suy luận trong chương & lý thuyết \\
7 & Với tần số xung, sun tác động bằng phép trừ, còn phép chia có được nhờ ức chế cùng với nhiễu nền cân bằng & Mô hình: ở nơ-ron đang phát xung, dòng qua sun gần như không đổi, và hiệu ứng lên tần số là phép trừ. Thực nghiệm: đưa vào nơ-ron tháp vỏ cảm giác thân thể chuột cống (lát cắt) các độ dẫn nền kích thích và ức chế bằng kẹp động; tăng đồng thời và cân bằng cả hai làm chia độ dốc đáp ứng mà không dịch tần số đáng kể. Tổng quan: phép chia cho tổng hoạt động gặp ở nhiều vùng não & \cite{HoltKoch1997,Chance2002,Carandini2012} & đã đo \\
8 & Khi $\beta>1$, ức chế lẫn nhau chọn ra một kẻ thắng (mệnh đề~\ref{prop:wta}, \eqref{eq:wta}) & Các trị riêng $-(1\pm\beta)/\tau$ là suy luận trong chương. Tần số $0.73$ và $0.53$ khi $\beta=0.5$ là điểm bất động $r_{1,2}=(I_{1,2}-\beta I_{2,1})/(1-\beta^2)$; không có script trong \texttt{models/}. Chương không dẫn thí nghiệm trực tiếp & suy luận trong chương & lý thuyết \\
9 & Đặt lại bộ nhớ bằng tín hiệu qua \nt{GABA}; hai nhóm ức chế lẫn nhau là một mạch chốt & Suy ra từ hình~\ref{fig:bistable} và mệnh đề~\ref{prop:wta}. Chưa có thí nghiệm & suy luận trong chương & lý thuyết \\
10 & Trạng thái cân bằng của vòng \nt{GLUT}--\nt{GABA} ổn định nếu ức chế đủ mạnh và đủ nhanh (mệnh đề~\ref{prop:ei}) & Mô hình hai quần thể~\eqref{eq:ei}; các điều kiện (i) và (ii) là định thức và vết của ma trận, được suy ra trong chương & \cite{WilsonCowan1972} & lý thuyết \\
11 & Với $w_{EE}=1.5$ và $w_{EI}w_{IE}=2$, ức chế nhanh ($\tau_I=2\ms$) giữ $E^*=0.67h$, ức chế chậm ($30\ms$) gây dao động (hình~\ref{fig:ei}) & Nghiệm số của~\eqref{eq:ei}, script \texttt{figures/make\_ei.py} & tính toán & mô hình \\
12 & Vỏ não làm việc ở chế độ ổn định nhờ ức chế; dấu hiệu --- đẩy thêm nơ-ron \nt{GABA} thì tần số của chúng giảm~\eqref{eq:isn} & Lý thuyết đã dự đoán đáp ứng nghịch lý. Thực nghiệm: ghi nội bào ở vỏ thị giác sơ cấp của mèo --- khi đáp ứng bị kích thích xung quanh ức chế, độ dẫn ức chế không tăng mà giảm cùng với độ dẫn kích thích. Kích thích quang di truyền các nơ-ron PV của chuột nhắt ở vỏ thị giác, vỏ cảm giác thân thể và vỏ vận động làm giảm tần số của nơ-ron ức chế, kể cả khi không có kích thích và khi gây mê nhẹ. Hệ số $-1/3$ với số liệu hình~\ref{fig:ei} là suy luận của chương & \cite{Tsodyks1997isn,Ozeki2009,Sanzeni2020} & đã đo \\
13 & Vòng ``\nt{GLUT} kích thích \nt{GABA}, \nt{GABA} ức chế \nt{GLUT}'' sinh ra nhịp nhanh với chu kỳ $12$--$50\ms$ & Kích thích quang di truyền các nơ-ron \nt{GABA} nhanh ở vỏ cảm giác thân thể chuột nhắt với tần số $8$--$200$\,Hz tăng cường có chọn lọc nhịp gamma ($20$--$80$\,Hz, chu kỳ $12$--$50\ms$); kích thích nơ-ron \nt{GLUT} chỉ tăng cường các nhịp chậm hơn & \cite{Cardin2009,Buzsaki2012} & đã đo \\
\bottomrule
\end{longtable}}

\section{Chương~\ref{sec:code}. Mã Morse}

Các ước tính giới hạn của chương là lý thuyết thông tin và phép tính; những gì đã đo là nhiễu phát sinh ở đâu trong kênh, não làm việc nhanh đến đâu và một xung tốn bao nhiêu, còn câu hỏi vỏ não thực sự dùng mã nào vẫn để ngỏ.

{\footnotesize
\setlength{\tabcolsep}{3pt}\renewcommand{\arraystretch}{1.15}
\begin{longtable}{>{\raggedright}p{0.5cm}>{\raggedright}p{4.0cm}>{\raggedright}p{5.6cm}>{\raggedright}p{2.3cm}>{\raggedright\arraybackslash}p{1.7cm}}
\toprule
TT & Khẳng định & Thí nghiệm và điều đã đo & Nguồn & Trạng thái \\
\midrule\endfirsthead
\toprule
TT & Khẳng định & Thí nghiệm và điều đã đo & Nguồn & Trạng thái \\
\midrule\endhead
1 & Tần số của nơ-ron \nt{GLUT} vỏ não từ một phần nhỏ đến hàng chục xung mỗi giây, và phần lớn thời gian nơ-ron làm việc gần giới hạn dưới & Ghi bằng pipet áp sát tế bào (việc chọn nơ-ron không phụ thuộc vào hoạt động của nó) ở vỏ thính giác chuột cống đang thức: tại mỗi thời điểm, dưới $5\,\%$ nơ-ron có tần số trên $20$ xung mỗi giây; một phần nơ-ron có tần số nền dưới $0.01$ mỗi giây; phân bố tần số là log-chuẩn & \cite{Hromadka2008} & đã đo \\
2 & Dung lượng của kênh xung $R_{\max}\approx r\log_2\frac{e}{r\Delta t}$~\eqref{eq:spikeentropy}, gần như toàn bộ nằm ở thời điểm xung (mệnh đề~\ref{prop:capacity}) & Entropy của chuỗi nhị phân gồm các ô dài $\Delta t$. Các con số của chương: $8$ bit trên một xung khi $r=10$ mỗi giây và $\Delta t=1\ms$, khoảng $5$ bit khi $\Delta t=10\ms$, dưới $2$ bit mỗi giây với bộ đếm xung & \cite{MacKay1952} & lý thuyết \\
3 & Nơ-ron cảm giác truyền từ một đến vài bit trên một xung, trong trường hợp tốt nhất tới một nửa giới hạn & Khôi phục kích thích từ xung của các nơ-ron cảm giác mà kích thích đã biết; tổng hợp các phép đo trong sách & \cite{Rieke1997} & đã đo \\
4 & Bộ tạo xung chính xác; thời điểm chính xác do những biến đổi nhanh của đầu vào quyết định & Nơ-ron vỏ não chuột cống trong lát cắt: với dòng lặp lại có dao động nhanh, xung lặp lại với độ chính xác dưới $1\ms$; với dòng không đổi, thời điểm trôi dạt & \cite{Mainen1995} & đã đo \\
5 & Khớp thần kinh không đáng tin cậy: hơn một nửa số xung không tới được nơi nhận do tính ngẫu nhiên của việc nhả & Kích thích tối thiểu các đầu vào của nơ-ron tháp hồi hải mã (lát cắt): hơn một nửa số xung không gây đáp ứng. Đã loại trừ dao động ngưỡng kích thích, lỗi dẫn truyền ở chỗ phân nhánh sợi trục và nhiệt độ thí nghiệm thấp; còn lại là sự nhả mang tính xác suất & \cite{AllenStevens1994} & đã đo \\
6 & Dòng xung trong vỏ não sống trông ngẫu nhiên: các khoảng phân tán như ở dòng Poisson, phương sai số xung gần bằng trung bình; một phần ``nhiễu'' là thông điệp về cử động của con vật & Ghi ở các vùng thị giác V1 và MT của khỉ đang thức: hệ số biến thiên của các khoảng xấp xỉ $1$, tỷ số phương sai số xung trên trung bình như ở một quá trình ngẫu nhiên; mô hình cộng các đầu vào nhỏ độc lập cho độ phân tán nhỏ hơn nhiều. Ở vỏ thị giác chuột nhắt, một phần đáng kể độ phân tán dự đoán được từ video ghi cử động của chính con vật & \cite{Softky1993,Stringer2019} & đã đo \\
7 & Mã tần số chậm: sai số $1/\sqrt{rT}$~\eqref{eq:rateerror}, còn não nhận ra hình ảnh trong $\sim150\ms$ & Công thức là thống kê Poisson, suy luận của chương. Thực nghiệm: người tham gia quyết định có con vật trong một bức ảnh lạ được chiếu $20\ms$ hay không; điện thế não với ``có'' và ``không'' tách nhau sau khoảng $150\ms$. Việc chia thời gian này thành khoảng một chục tầng, mỗi tầng $10$--$20\ms$, là ước tính của chương, không phải phép đo & \cite{Thorpe1996} & đã đo \\
8 & Lấy trung bình trên nhóm bị tương quan giới hạn: sai số giảm không quá $1/\sqrt c$ lần~\eqref{eq:correlated} (mệnh đề~\ref{prop:pooling}) & Các cặp nơ-ron vùng MT của khỉ được ghi đồng thời: hệ số tương quan số xung trung bình $0.12$, và phân tích lý thuyết cho thấy ngay cả tương quan như vậy cũng giới hạn lợi ích của nhóm. Lý thuyết: chỉ phần tương quan giống tín hiệu mới đặt ra giới hạn. Mô hình của quan điểm đối lập: tần số của nhóm $50$--$100$ nơ-ron được đọc trong một khoảng giữa hai xung, $10$--$50\ms$, và thời điểm của từng xung không được vỏ não sử dụng & \cite{Zohary1994,MorenoBote2014,ShadlenNewsome1998} & đã đo \\
9 & Con số có thể được ghi vào thời điểm xung đầu tiên~\eqref{eq:latency}, vào thứ tự xung của một nhóm hoặc vào pha của nhịp cục bộ & Võng mạc kỳ giông: cấu trúc không gian của một hình ảnh được chiếu ngắn được ghi trong độ trễ tương đối của các xung đầu tiên của nơ-ron đầu ra, gần như không phụ thuộc độ tương phản. Tế bào vị trí ở hồi hải mã chuột cống: khi đi qua vị trí ``của mình'', xung dịch ngày càng sớm hơn trong chu kỳ nhịp theta ($7$--$12$\,Hz), độ dịch từ $100^\circ$ đến $355^\circ$. Vỏ não dùng thời điểm xung đến mức nào là câu hỏi mở (dòng~8) & \cite{Gollisch2008,OKeefe1993} & đã đo \\
10 & Mã nào được đọc do hằng số thời gian của nơi nhận quyết định~\eqref{eq:reader} (mệnh đề~\ref{prop:reader}); ở vỏ não đang hoạt động, nơ-ron gần với bộ phát hiện trùng khớp hơn & Công thức~\eqref{eq:reader} là suy luận của chương. Tổng quan các bản ghi in vivo: ở vỏ não đang hoạt động, độ dẫn của màng cao hơn vài lần so với lúc nghỉ, và hằng số thời gian ngắn hơn tương ứng. Chưa có bằng chứng trực tiếp rằng vỏ não chuyển mã đúng theo cách này & \cite{Destexhe2003} & giả thuyết \\
11 & Khớp thần kinh cạn kiệt truyền sự thay đổi tần số, về thực chất là $\Delta r/r$~\eqref{eq:depression} (hình~\ref{fig:depression}) & Ghi cặp nơ-ron tháp vỏ não: tốc độ cạn kiệt do xác suất nhả quy định; khi cạn kiệt chậm, đáp ứng phản ánh tần số, khi nhanh --- phản ánh sự trùng khớp. Mô hình dựa trên sự cạn kiệt đo được ở vỏ não chuột cống: những thay đổi tương đối như nhau của tần số ở đầu vào nhanh và chậm cho đáp ứng như nhau, tức là khớp thần kinh truyền $\Delta r/r$. Các con số $1.7\to2.2$, $6.7$ và $90\ms$ là tính toán của chương & \cite{Tsodyks1997,Abbott1997depression} & đã đo \\
12 & Chùm xung là ``gạch'': xác suất nó bị mất $(1-p)^k$~\eqref{eq:burst} là nhỏ, tạo thuận còn làm giảm thêm & Tổng quan: ở các khớp thần kinh không tin cậy, xung đơn thường bị mất, còn chùm xung được truyền đáng tin cậy nhờ tạo thuận; chùm xung gây ra những thay đổi lâu dài của khớp thần kinh. Việc não đọc chùm xung như một ký hiệu riêng là giả thuyết & \cite{Lisman1997,AllenStevens1994} & giả thuyết \\
13 & Nơ-ron \nt{GABA} nhanh tạo nhịp và ``dấu câu''; một lớp khác ức chế các nơ-ron ức chế và mở cổng (mệnh đề~\ref{prop:punctuation}) & Tổng quan tính chất các lớp nơ-ron \nt{GABA} vỏ não. Quang di truyền: kích thích nhịp nhàng các nơ-ron \nt{GABA} nhanh gây ra nhịp gamma, và đáp ứng với kích thích phụ thuộc pha của chu kỳ. Ở vỏ thị giác chuột nhắt, nơ-ron PV ức chế lẫn nhau, nơ-ron SST ức chế mọi lớp khác, nơ-ron VIP chủ yếu ức chế SST. Kích thích nơ-ron VIP ở chuột nhắt đang thức gỡ bỏ ức chế khỏi nơ-ron \nt{GLUT}; chúng được bật bởi tín hiệu củng cố. Sự phân vai như một quy tắc chung là giả thuyết & \cite{Tremblay2016,Cardin2009,Pfeffer2013,Pi2013} & giả thuyết \\
14 & Năng lượng giới hạn tần số trung bình ở mức cỡ một xung mỗi giây~\eqref{eq:energy} & Ngân sách năng lượng của chất xám loài gặm nhấm theo dữ liệu giải phẫu và sinh lý: xung và dòng qua khớp thần kinh chiếm $47\,\%$ và $34\,\%$ năng lượng truyền tín hiệu; cùng lúc chỉ có thể có không quá $\sim15\,\%$ nơ-ron hoạt động. Với vỏ não người: không quá $\sim1\,\%$ nơ-ron có thể hoạt động mạnh cùng lúc. Con số $\sim10^9$ ATP mỗi xung và công suất $\sim10$\,W cho truyền tín hiệu là ước tính của chương dựa trên các tính toán này & \cite{Attwell2001,Lennie2003} & lý thuyết \\
15 & Mã thưa và được tinh chỉnh cho số bit trên mỗi joule lớn nhất; vỏ não đổi độ chính xác lấy năng lượng (mệnh đề~\ref{prop:economy}) & Giả thuyết mã tiết kiệm. Cơ sở: ở chuột nhắt bị hạn chế thức ăn, độ dẫn của thụ thể AMPA và mức tiêu thụ ATP của khớp thần kinh ở vỏ thị giác giảm ($-29\,\%$), tần số xung được giữ nguyên, còn độ chọn lọc hướng nới rộng $32\,\%$ & \cite{Barlow1961,Padamsey2022} & giả thuyết \\
\bottomrule
\end{longtable}}

\section{Chương~\ref{sec:serocomputer}. Nơ-ron \nt{SERO}}

Các tính chất của chính nơ-ron \nt{SERO} (nhịp, tự ức chế, dấu theo thụ thể, các phân hệ) đã được đo trực tiếp; các phép tính của chương --- bộ điều chỉnh, bộ lọc, logic --- suy ra từ các phương trình của nó; $\tau_c$, hệ số khuếch tán của serotonin và số vị trí phóng là ước tính; vai trò bộ điều chỉnh mức của vòng lặp trong não là giả thuyết (trong nhân raphe, tự ức chế mang tính điểm và chỉ gây những quãng ngừng ngắn), còn vai trò tầm nhìn của \nt{SERO} là giả thuyết được các thí nghiệm hành vi ủng hộ.

{\footnotesize
\setlength{\tabcolsep}{3pt}\renewcommand{\arraystretch}{1.15}
\begin{longtable}{>{\raggedright}p{0.5cm}>{\raggedright}p{4.0cm}>{\raggedright}p{5.6cm}>{\raggedright}p{2.3cm}>{\raggedright\arraybackslash}p{1.7cm}}
\toprule
TT & Khẳng định & Thí nghiệm và điều đã đo & Nguồn & Trạng thái \\
\midrule\endfirsthead
\toprule
TT & Khẳng định & Thí nghiệm và điều đã đo & Nguồn & Trạng thái \\
\midrule\endhead
1 & Dấu tác động của serotonin do thụ thể của bên nhận chọn (mệnh đề~\ref{prop:receptor}) & Thụ thể serotonin được chia thành các họ theo dược lý và cơ chế: 5-HT$_{1A}$ qua protein G mở kênh kali và ức chế tế bào, 5-HT$_{2A}$ và thụ thể hướng ion 5-HT$_3$ gây kích thích. Cùng một chất dẫn truyền cho đáp ứng ngược nhau ở các tế bào khác nhau. & \cite{BarnesSharp1999} & đã đo \\
2 & Khi thức, nơ-ron \nt{SERO} phát đều vài xung mỗi giây & Ghi nơ-ron nhân raphe lưng ở mèo di chuyển tự do: phát xung nhịp nhàng chậm $2.8\pm0.2$ xung/s khi thức yên tĩnh; trong giấc ngủ sóng chậm tần số giảm $34$--$68\,\%$, trong giấc ngủ REM giảm $98\,\%$. Ở chuột cống gây mê: $1.7\pm0.5$ xung/s, rất đều. & \cite{TrulsonJacobs1979, GagnonParent2014, JacobsAzmitia1992} & đã đo \\
3 & Nơ-ron \nt{SERO} là bộ tạo nhịp: không cần được kích cho từng xung, nhưng cần dòng nền đều (trong lát cắt là noradrenaline qua thụ thể $\alpha_1$, trong cơ thể là từ \nt{NORA}) & Trong lát cắt não, tế bào phát xung chậm và đều, sau mỗi xung có một pha tăng phân cực sau xung dài $200$--$800\ms$. Trong lát cắt, số tế bào tự phát xung ít hơn trong cơ thể: nhịp đều cần đầu vào noradrenaline liên tục qua thụ thể $\alpha_1$, chặn chúng thì nhịp tắt. & \cite{VandermaelenAghajanian1983} & đã đo \\
4 & Hầu hết thụ thể là hướng chuyển hóa; đáp ứng chậm: tăng và giảm trong khoảng một giây & Mọi họ thụ thể, trừ 5-HT$_3$, đều gắn với protein G. Trong lát cắt, serotonin được phóng ra do kích thích tạo qua 5-HT$_{1A}$ một dòng ức chế tăng và giảm trong vòng một giây. & \cite{BarnesSharp1999, CourtneyFord2016} & đã đo \\
5 & Ở các vùng đích, serotonin thoát vào thể tích chứ không chỉ vào khe; ngay trong nhân raphe, ức chế lân cận đi theo kiểu điểm--điểm & Đo vôn-ampe vòng quét nhanh trong lát cắt: lượng phóng trên một xung như nhau ở vùng có khớp thần kinh và ở nhân raphe, nơi gần như không có khớp thần kinh; không phụ thuộc vào việc chặn tái hấp thu và thụ thể; số phân tử trên mỗi đầu tận cùng ít hơn nhiều so với số chất vận chuyển và thụ thể, còn nồng độ ngoại bào trùng với ái lực của thụ thể chính. Trong nhân raphe, ức chế qua 5-HT$_{1A}$ đi theo kiểu điểm--điểm: chất vận chuyển không để serotonin tích tụ ngoài khe. & \cite{Bunin1998, CourtneyFord2016} & đã đo \\
6 & Nồng độ theo~\eqref{eq:seroneuron} giảm do bị bơm đi; $\tau_c$ cỡ một giây là ước tính & Đo vôn-ampe trong não sống: serotonin ngoại bào bị giới hạn trước hết bởi tái hấp thu và phân hủy, chứ không phải bởi tổng hợp. Các công trình được trích không đo trực tiếp $\tau_c$; bậc độ lớn là ước tính của chương. & \cite{Hashemi2012serotonin} & đã đo; $\tau_c$ là ước tính \\
7 & NOT và NOR trên một bộ tạo nhịp; tính đầy đủ của logic \nt{SERO} (mệnh đề~\ref{prop:serocomplete}, hình~\ref{fig:serogates}) & Suy ra từ~\eqref{eq:seroneuron}: bộ tạo nhịp bị ức chế là NOR, và NOR đầy đủ về chức năng. Chưa có thí nghiệm nào lắp logic từ nơ-ron \nt{SERO}; điều kiện các thể tích tách biệt không được thỏa mãn trong não. & suy ra trong chương & lý thuyết \\
8 & Serotonin ức chế chính nơ-ron \nt{SERO} qua các thụ thể trên nó & Ở chuột cống gây mê, chất chủ vận 5-HT$_{1A}$ tiêm tĩnh mạch và điện di ion ức chế sự phát xung của nơ-ron raphe lưng với cùng quan hệ liều--đáp ứng như chính serotonin; chất chủ vận 5-HT$_{1B}$ gần như không tác dụng. Trong lát cắt, serotonin nội sinh từ các tế bào lân cận tạo qua 5-HT$_{1A}$ một dòng và một quãng ngừng ngắn trong chuỗi xung, không làm đổi tần số trung bình. & \cite{Sprouse1987, CourtneyFord2016} & đã đo \\
9 & Trong mô hình, vòng lặp qua thể tích chung là bộ điều chỉnh mức: sai lệch và hằng số thời gian giảm $1+G$ lần~\eqref{eq:feedback}; việc vòng lặp hoạt động như vậy trong não là giả thuyết & Công thức kinh điển của bộ khuếch đại có hồi tiếp âm; được suy lại trong chương. Hệ số khuếch đại vòng $G$ của hệ \nt{SERO} chưa được đo. Đã đo: trong lát cắt, tự ức chế đi theo kiểu điểm--điểm, gây quãng ngừng ngắn và không đổi tần số trung bình. & \cite{Black1934feedback, CourtneyFord2016} & lý thuyết; trong não là giả thuyết \\
10 & Nồng độ là trung bình trượt của tần số, bộ lọc thông thấp~\eqref{eq:lowpass}, hình~\ref{fig:lowpass} & Nghiệm của phương trình tuyến tính~\eqref{eq:seroneuron}; hình vẽ được tính theo công thức này với $\tau_c=1\,\text{s}$. Không có mô hình riêng trong \texttt{models/}. & suy ra trong chương & lý thuyết \\
11 & Độ ngoằn ngoèo của khe làm khuếch tán của phân tử nhỏ chậm đi khoảng $2.6$ lần & Phương pháp nguồn điểm: điện di ion tetramethylammonium và ghi nồng độ của nó bằng điện cực chọn lọc ion trong lát cắt và trong não sống. Độ ngoằn ngoèo $\lambda\approx1.6$, tức $D$ hiệu dụng nhỏ hơn giá trị tự do $\lambda^2\approx2.6$ lần; tỷ lệ thể tích gian bào khoảng $0.2$. Hệ số khuếch tán của chính serotonin trong mô không được đo trong các công trình này; trong chương nó là ước tính. & \cite{Sykova2008} & đã đo \\
12 & Chiều dài của một ``dây dẫn'' độc lập $\ell\sim\sqrt{D\tau}\approx20\um$~\eqref{eq:diffusionlength}; số dây bị giới hạn bởi số vùng như vậy (mệnh đề~\ref{prop:volumewires}) & Định luật khuếch tán và thế các ước tính $D$, $\tau$ (dòng~6 và~11); mệnh đề~\ref{prop:volumewires} là hệ quả. Chưa có thí nghiệm trực tiếp đếm số kênh độc lập của truyền dẫn thể tích. & suy ra trong chương & lý thuyết \\
13 & Đầu ra của một nơ-ron \nt{SERO} trải trên những vùng não rất lớn; số vị trí phóng ước tính $\sim10^5$ & Tái dựng toàn bộ từng nơ-ron raphe lưng ở chuột cống: mọi sợi trục đi lên đều phân nhánh mạnh, tổng chiều dài sợi trục của một tế bào tới $18.7$~cm, các nhánh của một tế bào vươn tới thể vân, vỏ não trước trán và hạch hạnh nhân. Số vị trí phóng trên mỗi tế bào chưa được đếm trực tiếp. & \cite{GagnonParent2014} & đã đo; $10^5$ là ước tính \\
14 & Các nơ-ron \nt{SERO} chia thành vài phân hệ với đầu ra và đáp ứng khác nhau & Đánh dấu di truyền bằng virus ở chuột nhắt: các nơ-ron chiếu tới hạch hạnh nhân và tới vỏ não trán hầu như không chồng lấn về phân nhánh, nhận đầu vào khác nhau và đáp ứng ngược nhau với kích thích khó chịu; bật và tắt mỗi nhóm làm thay đổi hành vi theo cách khác nhau. & \cite{Ren2018} & đã đo \\
15 & Điều kiện kiên nhẫn $D<\tau_\gamma\ln(R/r_0)$~\eqref{eq:patience} với chiết khấu hàm mũ; với chiết khấu đo được, dạng hyperbol, là $D<\tau_\gamma(R/r_0-1)$ & Suy ra từ chiết khấu $e^{-D/\tau_\gamma}$ hoặc $1/(1+D/\tau_\gamma)$; suy ra trong chương. Khỉ, lựa chọn với độ trễ cỡ giây: chiết khấu gần hyperbol hơn hàm mũ; đáp ứng của nơ-ron \nt{DOPA} với tín hiệu báo trước phần thưởng hoãn lại cũng giảm theo độ trễ như hyperbol. & suy ra trong chương; \cite{KobayashiSchultz2008} & lý thuyết \\
16 & \nt{SERO} đặt tầm nhìn $\tau_\gamma$ (mục~\ref{sec:horizon}) & Kích hoạt quang di truyền các nơ-ron \nt{SERO} raphe lưng ở chuột nhắt kéo dài thời gian chờ phần thưởng hoãn lại và tăng sự kiên trì khi tìm phần thưởng đã trở nên hiếm. Ở người, giảm serotonin (chế độ ăn không có tryptophan) làm chiết khấu phần thưởng hoãn lại dốc hơn. Nồng độ chuyển thành $\tau_\gamma$ thế nào thì chưa rõ. & \cite{Doya2002, Miyazaki2014, Lottem2018, Schweighofer2008} & giả thuyết \\
\bottomrule
\end{longtable}}

\section{Chương~\ref{sec:dofa}. Học với \nt{DOPA}}

Các cơ chế của khớp thần kinh (gói, bộ phát hiện trùng khớp, calci, cửa sổ dẻo) và việc dopamine hành xử như sai số dự đoán đã được đo trực tiếp; việc các quy tắc dẫn tới gradient và cái giá của quảng bá là định lý; kết quả học lựa chọn là tính toán theo mô hình của sách; sơ đồ tính sai số là giả thuyết.

{\footnotesize
\setlength{\tabcolsep}{3pt}\renewcommand{\arraystretch}{1.15}
\begin{longtable}{>{\raggedright}p{0.5cm}>{\raggedright}p{4.0cm}>{\raggedright}p{5.6cm}>{\raggedright}p{2.3cm}>{\raggedright\arraybackslash}p{1.7cm}}
\toprule
TT & Khẳng định & Thí nghiệm và điều đã đo & Nguồn & Trạng thái \\
\midrule\endfirsthead
\toprule
TT & Khẳng định & Thí nghiệm và điều đã đo & Nguồn & Trạng thái \\
\midrule\endhead
1 & Trong não chưa tìm thấy lan truyền ngược sai số ở dạng như trong mạng nhân tạo & Không có thí nghiệm trực tiếp: đây là một khẳng định về sự vắng mặt. Các bài tổng quan phân tích thuật toán đòi hỏi gì (liên kết ngược lặp lại liên kết thuận, một pha riêng) và những xấp xỉ sinh học nào đã được đề xuất. & \cite{Lillicrap2020, WhittingtonBogacz2019} & giả thuyết \\
2 & Trọng số tách thành số vị trí phóng, xác suất phóng và đáp ứng với một gói: $w=N_s\,p\,A_1$~\eqref{eq:quantal} & Khớp thần kinh thần kinh--cơ của ếch khi lượng phóng bị giảm: đáp ứng của bên nhận dao động theo từng bậc, là bội số của đáp ứng thu nhỏ tự phát với một gói; số gói trên một xung tuân theo luật Poisson. & \cite{delCastillo1954} & đã đo \\
3 & Thụ thể của bên nhận là bộ phát hiện trùng khớp; calci khởi động sự thay đổi trọng số & Kẹp miếng màng (patch-clamp) nơ-ron chuột nhắt: ion Mg$^{2+}$ bịt kênh do glutamate mở ở điện thế nghỉ và rời đi khi khử cực. Lát cắt hồi hải mã: chất tạo chelat calci trong bên nhận chặn tăng cường dài hạn, còn giải phóng calci bằng ánh sáng bên trong tế bào tự nó làm tăng truyền dẫn. & \cite{Nowak1984, Malenka1988calcium} & đã đo \\
4 & Phía nào cũng có thể thực thi quyết định: $A_1$ tăng ở bên nhận, $p$ đổi ở bên gửi & Glutamate được giải phóng bằng ánh sáng phía trên một gai gây ra sự tăng trưởng nhanh và bền của gai cùng dòng qua các thụ thể của nó; sự tăng cường đi kèm việc gắn thêm thụ thể AMPA. Khử cực bên nhận tạm thời ức chế sự phóng ở bên gửi; hiệu ứng được truyền bởi endocannabinoid đi ngược qua khe (đã chỉ ra ở các đầu vào \nt{GABA}). Suy giảm ngắn hạn phụ thuộc vào xác suất phóng ở bên gửi. & \cite{Matsuzaki2004, Malinow2002, WilsonNicoll2001, Tsodyks1997} & đã đo \\
5 & Khớp thần kinh mạnh lên khi bên gửi và bên nhận cùng hoạt động~\eqref{eq:hebb} & Thỏ gây mê: sau những loạt xung ngắn theo đường vào của hồi hải mã, đáp ứng được tăng cường trong khoảng từ $30$~phút tới $10$~giờ. Trong lát cắt hồi hải mã, xung của bên nhận chỉ làm mạnh những khớp thần kinh đang hoạt động cùng lúc. & \cite{Bliss1973LTP, Kelso1986hebb} & đã đo \\
6 & Quy tắc~\eqref{eq:oja} tìm vectơ riêng chính của tương quan đầu vào (mệnh đề~\ref{prop:oja}) & Định lý; chứng minh trong chương. Chưa có thí nghiệm nào cho thấy một nơ-ron học được thành phần chính của các đầu vào của nó; cơ chế giới hạn sự tăng trọng số trong khớp thần kinh chưa được xác định. & \cite{Oja1982} & lý thuyết \\
7 & Hebb theo thời gian: cửa sổ khoảng $\pm20\ms$ (hình~\ref{fig:stdp}); từ các chuỗi lặp lại mọc ra các chuỗi liên kết & Các cặp nơ-ron hồi hải mã trong mẫu nuôi cấy: xung của bên nhận trong vòng $20\ms$ sau xung của bên gửi cho tăng cường bền, trong vòng $20\ms$ trước thì cho suy yếu; cả hai phụ thuộc thụ thể NMDA. Tỷ lệ $A_-=0.6A_+$ trên hình chỉ để minh họa. Việc các chuỗi mọc ra trong mạng theo cách này chưa được đo trực tiếp. & \cite{BiPoo1998} & đã đo \\
8 & Vết~\eqref{eq:threefactor}: dopamine đến sau hoạt động đồng thời $1$--$2\,\text{s}$ thì làm mạnh liên kết, muộn hơn thì không (mệnh đề~\ref{prop:threefactor}) & Lát cắt thể vân, kích thích quang di truyền riêng rẽ đầu vào glutamate và dopamine: gai chỉ lớn lên nếu dopamine đến sau đầu vào glutamate $0.3$--$2\,\text{s}$; cửa sổ do sự tăng nhanh và phân rã của cAMP quyết định. Ở vỏ não, việc ghép cặp để lại những vết riêng cho tăng cường và suy yếu, và các thụ thể monoamine biến chúng thành thay đổi dài hạn trong cửa sổ cỡ giây. & \cite{Yagishita2014, He2015eligibility} & đã đo \\
9 & Cửa sổ dài cỡ giây cũng có khi không có dopamine: tín hiệu thứ ba là sự kích thích mạnh của bên nhận & Chuột nhắt sống, vùng CA1: một sự kiện điện thế bình nguyên ở bên nhận làm mạnh các đầu vào đến trong vài giây trước và sau nó, và chỉ sau một lượt đi qua đã tạo ra một trường vị trí. & \cite{Bittner2017} & đã đo \\
10 & Bước trung bình của quy tắc ba thừa số đi theo gradient phần thưởng~\eqref{eq:perturbation} (mệnh đề~\ref{prop:gradient}) & Định lý gradient cho việc học theo độ lệch ngẫu nhiên; trong chương được suy ra cho nhiễu nhỏ. & \cite{Williams1992} & lý thuyết \\
11 & Nhiễu là sự tìm kiếm: mạng học từ những độ lệch ngẫu nhiên của chính nó & Chim sẻ vằn non: tắt nhân đầu ra của mạch hạch nền làm mất đi đột ngột và có hồi phục tính biến thiên của tiếng hót. Chim sẻ vằn trưởng thành: nhiễu được đưa vào những lần hót một âm tiết có cao độ nằm về một phía so với trung bình, và chim nhanh chóng dịch cao độ sang phía kia --- học từ độ phân tán của chính mình. & \cite{Olveczky2005, TumerBrainard2007} & đã đo \\
12 & Lựa chọn giữa hai phương án được học khi $\tau_e=1\,\text{s}$ và $\delta=R-\bar R$; không học được khi $\tau_e=50\ms$; không trừ kỳ vọng thì trọng số tăng không giới hạn, còn với tốc độ học lớn mạng có thể cố định lựa chọn tệ hơn & \texttt{models/dofa\_learning.py}, $\eta=0.05$, $500$ lượt thử, $6$ seed. $\tau_e=1\,\text{s}$: tỷ lệ chọn $A$ $0.65$--$0.79\to0.96$--$1.00$, trọng số $0.85$--$1.33$. $\tau_e=50\ms$: $0.38$--$0.55$, trọng số không đổi. $\delta=R$: trọng số của bên thắng $860$--$1190$. $\delta=R$, $\eta=0.5$, $3$ seed: một lần bị cố định ở $B$ ($0.04\to0$). Script in ra cả bốn trường hợp; tính lại khớp với chương. & \texttt{models/\allowbreak dofa\_\allowbreak learning.py} & mô hình \\
13 & Thời gian học bằng một tín hiệu chung tăng tỷ lệ với số nơ-ron gây nhiễu (mệnh đề~\ref{prop:broadcastcost}) & Đường cong học của mạng tuyến tính: nhiễu loạn nơ-ron chậm hơn gradient chính xác một số lần bằng số nơ-ron đầu ra, nhiễu loạn trọng số còn chậm thêm một số lần bằng số đầu vào. & \cite{Werfel2005} & lý thuyết \\
14 & Đầu ra của \nt{DOPA} đi trước hết tới các vùng chọn hành động & Tái dựng toàn bộ sợi trục của từng nơ-ron \nt{DOPA} đường liềm đen--thể vân ở chuột cống: các nhánh của một tế bào phủ $0.45$--$5.7\,\%$ thể tích thể vân (trung bình $2.7\,\%$), ngoài thể vân gần như không có nhánh. & \cite{Matsuda2009} & đã đo \\
15 & Vỏ não học cách dự đoán đầu vào của nó, sai số được tính tại chỗ & Tổng quan: trong vỏ não cảm giác có những đáp ứng với sự không khớp giữa đầu vào được kỳ vọng và đầu vào thực đến. Việc đây là quy tắc học chung của vỏ não chưa được chứng minh. & \cite{Keller2018} & giả thuyết \\
16 & Dopamine hành xử như sai số~\eqref{eq:ctd}: đỉnh vọt, dịch chuyển sang tín hiệu báo trước, hõm sụt (hình~\ref{fig:ctdcases}) & Nơ-ron \nt{DOPA} ở khỉ: đỉnh vọt khi có nước ép bất ngờ; sau khi học --- đỉnh vọt lúc có tín hiệu báo trước và không đáp ứng với nước ép đã hứa; hõm sụt khi nước ép không đến. Dạng liên tục~\eqref{eq:ctd} là lý thuyết. & \cite{Schultz1997, Doya2000} & đã đo \\
17 & Sơ đồ tính $\dd V/\dd t$ và trừ đi kỳ vọng & Chuột nhắt, ghi và quang di truyền ở vùng mái bụng: nơ-ron \nt{DOPA} trừ kỳ vọng khỏi phần thưởng; các nơ-ron \nt{GABA} lân cận ức chế chúng khi phần thưởng được kỳ vọng, và kích thích hay ức chế các nơ-ron này làm thay đổi sai số. Đạo hàm được tính thế nào thì chưa được chỉ ra. & \cite{Eshel2015arithmetic} & giả thuyết \\
18 & Nơ-ron \nt{DOPA} là bộ tạo nhịp; hõm sụt nông hơn đỉnh vọt & Trong lát cắt, nơ-ron \nt{DOPA} phát xung tự phát và rất đều, trên nền một điện thế tạo nhịp chậm. Ở khỉ, tần số bám theo sai số dương một cách định lượng, còn sai số âm thì chỉ truyền được yếu ớt. & \cite{GraceOnn1989, Bayer2005} & đã đo \\
19 & Bộ hành động và bộ phê bình (hình~\ref{fig:actorcritic}) & Thuật toán lấy từ lý thuyết học tăng cường. Ở người (fMRI), thể vân bụng hành xử như bộ phê bình cả khi có lựa chọn lẫn khi không, thể vân lưng --- chỉ khi phải chọn hành động. & \cite{SuttonBarto2018, ODoherty2004actor} & giả thuyết \\
20 & Dấu của $\delta$ trong quy tắc bị đảo ở các nơ-ron có họ thụ thể thứ hai & Lát cắt thể vân: ở nơ-ron có thụ thể D1, việc ghép cặp cho tăng cường, cần D1; ở nơ-ron có D2 --- suy yếu, cần D2. & \cite{Shen2008} & đã đo \\
\bottomrule
\end{longtable}}

\section{Chương~\ref{sec:dofaalt}. Các lý thuyết khác về \nt{DOPA}}

Mỗi bức tranh mở rộng dựa trên những phép đo dopamine trực tiếp của riêng nó (độ phân tán của các điểm đảo chiều, đáp ứng với điều khó chịu và điều mới, sự tăng chậm, đáp ứng khi đổi vị), nhưng các công thức giải thích chúng là lý thuyết, và giữa hai trong số đó các thí nghiệm hiện vẫn mâu thuẫn nhau.

{\footnotesize
\setlength{\tabcolsep}{3pt}\renewcommand{\arraystretch}{1.15}
\begin{longtable}{>{\raggedright}p{0.5cm}>{\raggedright}p{4.0cm}>{\raggedright}p{5.6cm}>{\raggedright}p{2.3cm}>{\raggedright\arraybackslash}p{1.7cm}}
\toprule
TT & Khẳng định & Thí nghiệm và điều đã đo & Nguồn & Trạng thái \\
\midrule\endfirsthead
\toprule
TT & Khẳng định & Thí nghiệm và điều đã đo & Nguồn & Trạng thái \\
\midrule\endhead
1 & Quy tắc bất đối xứng~\eqref{eq:asymmetric} hội tụ về điểm~\eqref{eq:expectile}; khi $\kappa=1/2$ đó là trung bình, với một giọt có xác suất $1/2$ thì $V=\kappa$ (hình~\ref{fig:expectiles}) & Điểm~\eqref{eq:expectile} là expectile, nghiệm của bài toán bình phương tối thiểu với trọng số khác nhau cho độ lệch dương và âm; với ví dụ trong chương, phép suy ra có ngay trong chương. & \cite{NeweyPowell1987}; suy ra trong chương & lý thuyết \\
2 & Các nơ-ron \nt{DOPA} khác nhau có điểm đảo chiều khác nhau, và chúng được sắp theo độ bất đối xứng (mệnh đề~\ref{prop:expectile}) & Chuột nhắt, nơ-ron \nt{DOPA} vùng mái bụng được đánh dấu quang di truyền, phần thưởng có kích thước khác nhau: điểm mà đỉnh vọt chuyển thành hõm sụt khác nhau ở các nơ-ron; ở nơ-ron có đáp ứng bất đối xứng hơn thì điểm này cao hơn; từ đáp ứng của nhóm có thể khôi phục hình dạng phân bố phần thưởng. Việc đây là chức năng chính của sự phân tán là giả thuyết. & \cite{Dabney2020} & đã đo \\
3 & Một phần nơ-ron \nt{DOPA} đáp ứng với điều khó chịu bằng đỉnh vọt & Khỉ: một số nơ-ron \nt{DOPA} bị kích thích bởi tín hiệu báo trước phần thưởng và bị ức chế bởi tín hiệu báo trước luồng khí thổi vào mắt, số khác bị kích thích bởi cả hai; các nhóm nằm ở những phần khác nhau của não giữa. & \cite{Matsumoto2009, BrombergMartin2010} & đã đo \\
4 & Độ lớn của sai số hoặc độ mới đặt tốc độ học & Các công trình được trích không có thí nghiệm trực tiếp nào trong đó tín hiệu tầm quan trọng của nơ-ron \nt{DOPA} làm thay đổi tốc độ học; bài tổng quan đề xuất vai trò tín hiệu ``chú ý, quan trọng''. & \cite{BrombergMartin2010} & giả thuyết \\
5 & Một dây riêng cho trục nguy hiểm & Chuột nhắt: nơ-ron \nt{DOPA} đi tới phần đuôi của thể vân đáp ứng với kích thích mới và mạnh chứ không với phần thưởng; phá hủy chúng làm yếu đi hành vi tránh một vật mới. & \cite{Menegas2018} & đã đo \\
6 & Thời gian hành động tối ưu $\tau_a^*=\sqrt{C/\bar R}$~\eqref{eq:vigor} & Cực tiểu của tổng $C/\tau_a+\bar R\tau_a$; suy ra trong chương. Trong mô hình gốc, tốc độ hành động do giá của thời gian quyết định, bằng phần thưởng trung bình. & \cite{Niv2007}; suy ra trong chương & lý thuyết \\
7 & Mức dopamine chậm mang $\bar R$ và đặt tốc độ hành động (mệnh đề~\ref{prop:multiplex}) & Chuột cống, vi thẩm tách và đo vôn-ampe ở nhân accumbens: mức dopamine theo từng phút bám theo tần suất phần thưởng và tốc độ hành động. Ở người, L-DOPA làm tăng sự phụ thuộc của thời gian phản ứng vào tần suất phần thưởng trung bình. Việc mức chậm đúng bằng $\bar R$ chưa được chứng minh. & \cite{Hamid2016, Beierholm2013vigor} & giả thuyết \\
8 & Mức chậm được tạo từ hai nguồn: lượng phóng thay đổi ở đầu ra khi tần số xung không đổi, nhưng sự tăng chậm cũng có trong chính tần số & Chuột cống, cùng một nhiệm vụ: lượng phóng ở nhân accumbens bám theo kỳ vọng phần thưởng thay đổi, còn tần số xung của các nơ-ron \nt{DOPA} vùng mái bụng đã được nhận dạng thì không. Chuột nhắt trong hành lang ảo (dòng~10): sự tăng dần đều khi tiến gần phần thưởng có cả trong xung của các nơ-ron \nt{DOPA} đã được nhận dạng. & \cite{Mohebi2019, Kim2020} & đã đo \\
9 & Dopamine tăng dần đều trong khi con vật tiến gần một phần thưởng ở xa & Chuột cống trong mê lộ, đo vôn-ampe ở thể vân: nồng độ dopamine tăng trong nhiều giây khi tiến gần đích, độ dốc phụ thuộc vào khoảng cách và kích thước phần thưởng. & \cite{Howe2013ramp, Hamid2016} & đã đo \\
10 & Sự tăng là sai số $\delta$ khi ước lượng chính xác dần~\eqref{eq:ramp}; bước nhảy trong hành lang cho đỉnh vọt (hình~\ref{fig:ramp}) & Công thức và con số $\delta=0.75\,V$ mỗi giây được suy ra trong chương. Thí nghiệm: chuột nhắt trong hành lang ảo, được dời tức thời lại gần đích; xung ở thân tế bào, calci ở thân và ở đầu ra, và nồng độ trong thể vân đáp ứng như sai số chứ không như kỳ vọng $V$. Cách giải thích nào trong hai cách là đúng ở mọi đầu ra thì vẫn đang được bàn cãi. & \cite{Kim2020} & đã đo \\
11 & Nhìn lại phía sau: dopamine báo thước đo liên hệ nhân quả~\eqref{eq:retro} & Chuột nhắt, lượng dopamine phóng ở nhân accumbens trong những thí nghiệm mà lý thuyết này và sai số~\eqref{eq:ctd} dự đoán khác nhau: trong một số thí nghiệm, lượng phóng tuân theo lý thuyết thứ nhất. & \cite{Jeong2022} & giả thuyết \\
12 & Trong quá trình học, đáp ứng dopamine dịch dần từ phần thưởng sang tín hiệu báo trước & Chuột nhắt, tín hiệu ở đầu ra của nơ-ron \nt{DOPA} trong thể vân bụng theo tiến trình học: đỉnh vọt dần chuyển từ thời điểm có phần thưởng sang thời điểm có tín hiệu báo trước, đúng như việc học theo~\eqref{eq:ctd} đòi hỏi. & \cite{Amo2022} & đã đo \\
13 & Nơ-ron \nt{DOPA} đáp ứng khi vị của phần thưởng đổi mà giá trị vẫn như cũ & Chuột cống, ghi nơ-ron vùng mái bụng: thay nước ép bằng một loại khác có giá trị tương đương gây ra đáp ứng, dù sai số~\eqref{eq:ctd} bằng không. & \cite{Takahashi2017} & đã đo \\
14 & Đỉnh vọt nhân tạo dạy mối liên hệ giữa hai kích thích trung tính & Chuột cống, thí nghiệm ghép cặp trước hai kích thích không có phần thưởng: kích hoạt quang di truyền nơ-ron \nt{DOPA} lúc chuyển từ kích thích thứ nhất sang thứ hai tạo ra mối liên hệ, tắt chúng thì cản trở việc học liên hệ đó. & \cite{Sharpe2017} & đã đo \\
15 & Vectơ sai số theo đặc trưng~\eqref{eq:vectortd} & Lý thuyết về sai số dự đoán đặc trưng; chưa có kiểm tra trực tiếp dạng vectơ. & \cite{Gardner2018} & giả thuyết \\
16 & Trong cùng một nhiệm vụ, các nơ-ron \nt{DOPA} khác nhau mang những đại lượng khác nhau & Chuột nhắt trong hành lang ảo, chụp hai photon nơ-ron \nt{DOPA}: một số liên quan tới phần thưởng, số khác tới tốc độ và hướng rẽ, số khác nữa tới tín hiệu báo trước và độ chính xác của lựa chọn. & \cite{Engelhard2019} & đã đo \\
17 & Bó dây thay cho một dây (mệnh đề~\ref{prop:bundle}) & Tổng hợp các dòng~2--16: mỗi cách phân tách đã được đo riêng; chưa có thí nghiệm cho bức tranh thống nhất rằng tất cả đều là phần của một tín hiệu. & --- & giả thuyết \\
\bottomrule
\end{longtable}}

\section{Chương~\ref{sec:nora}. \nt{NORA}}

Cấu trúc của hệ \nt{NORA}, hai chế độ của nó, mối liên hệ với đồng tử (không chỉ qua \nt{NORA}), sự tăng tỷ số tín hiệu trên nền và chữ U ngược trong hành vi đã được đo trực tiếp; mức khuếch đại nhân $g$ là mô hình; việc khuếch đại chuyển chế độ lựa chọn và hoạt động như nghịch đảo nhiệt độ là những điều suy ra trong chương; lợi ích của việc điều chỉnh khuếch đại là tính toán theo mô hình của sách; ``\nt{NORA} là núm vặn khám phá'' và ``\nt{NORA} là ngắt'' là giả thuyết.

{\footnotesize
\setlength{\tabcolsep}{3pt}\renewcommand{\arraystretch}{1.15}
\begin{longtable}{>{\raggedright}p{0.5cm}>{\raggedright}p{4.0cm}>{\raggedright}p{5.6cm}>{\raggedright}p{2.3cm}>{\raggedright\arraybackslash}p{1.7cm}}
\toprule
TT & Khẳng định & Thí nghiệm và điều đã đo & Nguồn & Trạng thái \\
\midrule\endfirsthead
\toprule
TT & Khẳng định & Thí nghiệm và điều đã đo & Nguồn & Trạng thái \\
\midrule\endhead
1 & Ở người có vài chục nghìn nơ-ron \nt{NORA}, gần như tất cả trong một nhóm & Lát cắt liên tiếp thân não người, nhuộm tyrosine hydroxylase: $53\,900$ tế bào trong nhân lục (locus coeruleus) và thêm $6\,260$ trong phân nhân lân cận. & \cite{Baker1989LC} & đã đo \\
2 & Đầu ra tỏa gần như khắp não, nhưng các phần của nhóm phần nào phục vụ những vùng khác nhau & Đánh dấu di truyền bằng virus ở chuột nhắt: nơ-ron \nt{NORA} của nhân lục nhận đầu vào từ những vùng rộng của não và gửi đầu ra tới những vùng rộng của não và tủy sống. Ở chuột cống: nhóm chiếu tới hạch hạnh nhân tăng cường việc học sợ hãi, còn một nhóm riêng chiếu tới vỏ não trước trán tăng cường việc dập tắt nó. & \cite{Schwarz2015LC, Uematsu2017LC, Poe2020} & đã đo \\
3 & Chế độ nền: xung đều với tần số từ dưới một tới vài xung mỗi giây, thay đổi theo mức tỉnh táo & Chuột cống, ghi nơ-ron nhân lục qua chu kỳ ngủ--thức: tần số cao nhất khi thức, thấp hơn trong giấc ngủ sóng chậm, gần bằng không trong giấc ngủ REM; thay đổi tần số đi trước sự chuyển trạng thái. & \cite{AstonJonesBloom1981} & đã đo \\
4 & Chế độ từng đợt: chùm xung ngắn với sự kiện quan trọng cho nhiệm vụ, vào khoảng thời điểm quyết định & Khỉ, nhiệm vụ tìm một mục tiêu hiếm: mọi nơ-ron nhân lục chỉ đáp ứng với mục tiêu bằng một chùm xung, $\sim90\ms$ sau mục tiêu và $\sim200\ms$ trước phản ứng bằng tay; khi làm kém, chùm xung yếu hơn. Trong nhiệm vụ chọn một trong hai, chùm xung gắn với phản ứng chặt hơn là với kích thích, và có cả khi trả lời đúng lẫn khi trả lời sai. & \cite{AstonJones1994vigilance, Clayton2004LC} & đã đo \\
5 & Đường kính đồng tử là điểm đo kiểm tra \nt{NORA} không chính xác: nó bám theo cả \nt{NORA}, cả \nt{ACET} và cả các vùng khác & Khỉ: xung của nhân lục, cho tới từng xung đơn lẻ của một tế bào, dự đoán sự giãn đồng tử; một mối liên hệ tương tự nhưng muộn hơn có ở các gò não và vỏ não đai. Chuột nhắt: dao động của đồng tử bám theo hoạt động của cả đầu ra noradrenaline lẫn đầu ra cholinergic trong vỏ não. & \cite{Joshi2016, Reimer2016} & đã đo \\
6 & Noradrenaline ức chế hoạt động nền mạnh hơn hoạt động được gợi ra; mức khuếch đại nhân~\eqref{eq:gain} là mô hình thể hiện hiệu ứng này & Khỉ tỉnh, vỏ não thính giác, điện di ion noradrenaline: hoạt động nền của nơ-ron bị ức chế mạnh hơn đáp ứng với tiếng kêu của đồng loại --- tỷ số tín hiệu trên nhiễu tăng. Hàm sigmoid nhân~\eqref{eq:gain} lấy từ một mô hình mạng; việc độ dốc bị nhân lên theo đúng nghĩa đen chưa được đo. & \cite{Foote1975NE, ServanSchreiber1990} & đã đo; $g$ là mô hình \\
7 & Khuếch đại chỉ làm tăng $d'$ trước nhiễu sau bộ khuếch đại~\eqref{eq:gainsnr} (mệnh đề~\ref{prop:gainnoise}) & Suy ra trong chương; trong mô hình mạng, khuếch đại cải thiện việc phát hiện tín hiệu. Chưa có thí nghiệm tách nhiễu trước và sau bộ khuếch đại và kiểm tra sự khác biệt này. & suy ra trong chương; \cite{ServanSchreiber1990} & lý thuyết \\
8 & Ranh giới $g\beta=1$ chuyển mạng lựa chọn từ ``cân nhắc'' sang ``đã quyết''~\eqref{eq:wtagain} (mệnh đề~\ref{prop:gainwta}, hình~\ref{fig:pitchfork}) & Phân tích tuyến tính hai nhóm ức chế lẫn nhau; suy ra trong chương. Chưa có phép đo trực tiếp ngưỡng này trong não. & suy ra trong chương & lý thuyết \\
9 & Hệ số khuếch đại là nghịch đảo nhiệt độ của lựa chọn~\eqref{eq:softmax} & Với nhiễu logistic sau bộ khuếch đại, xác suất lựa chọn là softmax với $T=\sigma/g$; suy ra trong chương. & suy ra trong chương & lý thuyết \\
10 & \nt{NORA} đặt mức khám phá: hoạt động nền cao --- $g$ hiệu dụng nhỏ (khám phá), chùm xung từng đợt lúc quyết định --- lớn (quyết định) & Người: trước một lựa chọn ngẫu nhiên, đồng tử nền rộng hơn so với trước khi chọn phương án tốt nhất đã biết; người có đồng tử rộng hơn thiên về khám phá hơn. Chuột cống: việc chuyển sang chế độ chọn ngẫu nhiên được điều khiển bởi đầu vào \nt{NORA} tới vỏ não đai trước. Việc chùm xung làm tăng $g$ hiệu dụng chưa được đo trực tiếp. & \cite{JepmaNieuwenhuis2011, Tervo2014stochastic, AstonJones2005} & giả thuyết \\
11 & Phần thưởng phụ thuộc vào $g$ theo chữ U ngược; $g$ tốt nhất nhỏ hơn trong thế giới thay đổi nhanh (mệnh đề~\ref{prop:explore}, hình~\ref{fig:invertedU}) & \texttt{models/nora\_explore.py}, $60$ lần chạy, mỗi lần $4000$ lượt thử. Thay đổi một lần sau mỗi $1000$ lượt: $g=16$ tốt nhất, tỷ lệ $0.976$; sau mỗi $20$ lượt: $g=8$, tỷ lệ $0.886$; khi $g=0.5$ --- $0.77$. Tính lại khớp với chương. & \texttt{models/\allowbreak nora\_\allowbreak explore.py} & mô hình \\
12 & Chữ U ngược trong hành vi: tốt nhất khi hoạt động nền vừa phải với chùm xung rõ nét & Khỉ, cùng nhiệm vụ như ở dòng~4: trong những giai đoạn làm tốt, tần số nền vừa phải và chùm xung với mục tiêu rõ nét; khi tần số nền cao thì không có chùm xung và nhiều phản ứng sai hơn. Thay đổi của hoạt động nền và hoạt động được gợi ra bám theo dao động của chất lượng thực hiện. & \cite{Usher1999LC, AstonJones2005} & đã đo \\
13 & \nt{NORA} báo sự bất định không được dự liệu, \nt{ACET} --- sự bất định được dự liệu & Lý thuyết suy luận Bayes với hai loại bất định; các công trình được trích không có thí nghiệm trực tiếp nào tách các tín hiệu này theo chất dẫn truyền. & \cite{YuDayan2005} & giả thuyết \\
14 & Sau một thay đổi đột ngột, người ta học nhanh hơn, và đồng tử giãn ra & Người, nhiệm vụ dự đoán với những thay đổi bất ngờ: những thay đổi ngắn của đồng tử bám theo ước lượng xác suất có thay đổi, và cùng với đồng tử nền dự đoán dữ liệu mới làm thay đổi ước lượng mạnh tới mức nào (tốc độ học); thay đổi đồng tử không liên quan tới ánh sáng và nhiệm vụ làm thay đổi tốc độ này. Việc ở đây đồng tử phản ánh chính \nt{NORA} là suy luận từ dữ liệu gián tiếp ở dòng~5. & \cite{Nassar2012} & đã đo \\
15 & Chùm xung từng đợt hoạt động như một ngắt: mạng xóa trạng thái & Chưa có thí nghiệm trực tiếp nào trong đó chùm xung \nt{NORA} xóa trạng thái của mạng; mới chỉ đo được chùm xung với các sự kiện quan trọng (dòng~4). & \cite{BouretSara2005} & giả thuyết \\
16 & Điều chỉnh $g$ hoặc tốc độ học theo sự bất ngờ hầu như không tốt hơn các thiết lập cố định tốt nhất & \texttt{models/nora\_explore.py}, thế giới với các thời kỳ dài $1000$ lượt thử (thay đổi một lần sau mỗi $1000$ và sau mỗi $20$ lượt). tốt nhất với $g$ cố định là $g=8$: $0.925$; điều chỉnh $g$: tới $0.927$; điều chỉnh tốc độ học: tới $0.926$; tốc độ cố định $0.4$: $0.928$. Tính lại khớp với chương. & \texttt{models/\allowbreak nora\_\allowbreak explore.py} & mô hình \\
\bottomrule
\end{longtable}}

\section{Chương~\ref{sec:acet}. Thêm \nt{ACET}}

Ba tác động của acetylcholine đã được đo trên lát cắt và trong não sống, xung đột giữa ghi và đọc được chỉ ra trên mô hình của sách, còn việc \nt{ACET} làm đường cho phép ghi và báo sự bất định được dự liệu là giả thuyết có một phần cơ sở thực nghiệm.

{\footnotesize
\setlength{\tabcolsep}{3pt}\renewcommand{\arraystretch}{1.15}
\begin{longtable}{>{\raggedright}p{0.5cm}>{\raggedright}p{4.0cm}>{\raggedright}p{5.6cm}>{\raggedright}p{2.3cm}>{\raggedright\arraybackslash}p{1.7cm}}
\toprule
TT & Khẳng định & Thí nghiệm và điều đã đo & Nguồn & Trạng thái \\
\midrule\endfirsthead
\toprule
TT & Khẳng định & Thí nghiệm và điều đã đo & Nguồn & Trạng thái \\
\midrule\endhead
1 & Não có ít nơ-ron \nt{ACET}, chúng tập trung trong vài nhóm, còn đầu ra tỏa khắp vỏ não: một kênh quảng bá & Nhuộm kháng thể với enzyme tổng hợp acetylcholine và đánh dấu ngược dòng ở chuột cống: vỏ não, hồi hải mã, hạch hạnh nhân, hành khứu và đồi thị nhận acetylcholine chủ yếu từ sáu nhóm tế bào đặc của não trước nền và thân não & \cite{Mesulam1983} & đã đo \\
2 & Mức acetylcholine trong vỏ não cao khi thức và thấp trong giấc ngủ sóng chậm & Vi thẩm tách ở vỏ não và hồi hải mã của mèo di chuyển tự do, kèm ghi điện não đồ: lượng acetylcholine phóng ra khi thức yên tĩnh cao hơn $\sim100\%$, khi thức hoạt động cao hơn $\sim175\%$ so với giấc ngủ sóng chậm; trong giấc ngủ REM, ở vỏ não bằng mức khi thức & \cite{Marrosu1995,Hasselmo1999} & đã đo \\
3 & Ý nghĩa của tín hiệu do thụ thể quyết định: acetylcholine có các thụ thể--kênh nhanh và các thụ thể chậm khởi động phản ứng trong tế bào (mệnh đề~\ref{prop:receptor}) & Tổng quan các phép đo: tác động của acetylcholine lên tính dễ kích thích, truyền dẫn và tính dẻo phụ thuộc vào nơi phóng, phân nhóm thụ thể (nicotinic là kênh ion, muscarinic hoạt động qua protein G) và tế bào nhận & \cite{Picciotto2012} & đã đo \\
4 & Tác động thứ nhất: làm yếu liên kết nội bộ, gần như không động tới đầu vào từ bên ngoài (thừa số $1-a$ trong~\eqref{eq:acetnet}) & Lát cắt vỏ não khứu giác chuột cống: với $100$\,\textmu M chất chủ vận acetylcholine, điện thế của liên kết nội bộ (lớp Ib) giảm hơn $60\%$, của đầu vào bên ngoài (lớp Ia) --- dưới $15\%$, trong cùng một tế bào; sự ức chế là ở phía trước khớp thần kinh. Lát cắt hồi hải mã (CA1): $89\%$ so với $40\%$ cho đường nội bộ và đường vào & \cite{HasselmoBower1992,HasselmoSchnell1994} & đã đo \\
5 & Tác động thứ hai: tăng cường đầu vào (thừa số $\frac{1+a}{2}$) & Lát cắt tân vỏ não: thụ thể nicotinic tăng cường các khớp thần kinh của đầu vào từ đồi thị và không tác động lên các khớp thần kinh trong vỏ não; thụ thể muscarinic làm yếu cả hai loại. Acetylcholine làm tăng tính dễ kích thích của nơ-ron (tổng quan) & \cite{Gil1997,Hasselmo2006} & đã đo \\
6 & Tác động thứ ba: làm trọng số dễ thay đổi hơn (tốc độ $\eta a$) & Chuột cống: kích thích điện nhân nền (nguồn acetylcholine cho vỏ não) ghép cặp với một âm gây ra ở con vật trưởng thành sự tái tổ chức mạnh bản đồ vỏ não thính giác theo âm được trình bày & \cite{KilgardMerzenich1998,Hasselmo2006} & đã đo \\
7 & Quy tắc Hebb cho mẫu thưa trong phương trình thứ hai của~\eqref{eq:acetnet} cho ra một bộ nhớ liên kết, điền nốt mẫu từ một phần & Tính dung lượng của mạng với mẫu thưa và quy tắc hiệp phương sai: khi tỷ lệ nơ-ron hoạt động $p$ nhỏ, mạng lưu được nhiều mẫu hơn đáng kể so với khi $p=1/2$ & \cite{Tsodyks1988} & lý thuyết \\
8 & Ghi khi $a$ thấp làm hỏng bộ nhớ: mạng ghi điều nó nhớ lại, không phải điều nó thấy; khi $a=0.1$, hư hỏng không nhẹ hơn khi $a=0$ & \texttt{models/\allowbreak acet\_memory.py}: $200$ nơ-ron, năm mẫu mỗi mẫu $20$, mẫu mới có chung $10$ nơ-ron với một trong số đó, $10$ lần chạy. Khi $a=0$, mẫu mới không được ghi nhớ, mẫu cũ kéo theo $5.9$ trong $10$ nơ-ron lạ; khi $a=0.1$, mẫu mới được nhớ lại ($0.97$), nhưng hai mẫu hòa lẫn: mẫu mới kéo theo $7.3$ nơ-ron lạ, mẫu cũ --- $7.9$; từ $a=0.2$ --- ít hơn một & suy ra trong chương & mô hình \\
9 & Mệnh đề~\ref{prop:writeread}: không có mức $a$ nào mà cả ghi lẫn đọc đều tốt như nhau (hình~\ref{fig:acetsim}) & Cùng script, tốc độ học $\eta a$: mẫu mới được ghi nhớ trọn vẹn sau một lần trình bày khi $a\ge0.9$, được $0.8$ khi $a=0.75$, không được ghi nhớ khi $a\le0.5$. Đọc từ một nửa mẫu: $1.0$ khi $a\le0.2$, $0.96$ khi $a=0.3$, $0.04$ khi $a=0.4$ & suy ra trong chương & mô hình \\
10 & Trong não, một dây quảng bá --- \nt{ACET} --- chuyển giữa ghi và đọc & Ý tưởng lấy từ các mô hình xây dựng dựa trên phép đo trên lát cắt. Thí nghiệm trực tiếp nhất là ở người: chất chẹn thụ thể muscarinic scopolamine làm giảm khả năng học thuộc các cặp từ mới, mạnh nhất với những cặp chồng lấn với cặp cũ, nhưng không làm giảm việc nhớ lại các cặp đã học. Chưa có phép đo trực tiếp hai chế độ của mạng & \cite{HasselmoAndersonBower1992,Atri2004} & giả thuyết \\
11 & Bộ lọc~\eqref{eq:kalman} là tối ưu; trọng số đầu vào và tốc độ học là cùng một đại lượng $P/R$, ở chế độ xác lập $P=\sqrt{qR}$ và thời gian nhớ $\sqrt{R/q}$ (mệnh đề~\ref{prop:kalman}) & Không cần thí nghiệm: tính tối ưu của bộ lọc Kalman--Bucy là kết quả kinh điển của lý thuyết ước lượng; chế độ xác lập được suy ra trong chương & \cite{KalmanBucy1961} & lý thuyết \\
12 & \nt{ACET} báo cho mạng một đại lượng giống $P$ --- sự bất định được dự liệu & Được đề xuất trong một mô hình xác suất về sự chú ý. Chưa có phép đo trực tiếp cho thấy mức acetylcholine bám theo độ bất định của ước lượng & \cite{YuDayan2005} & giả thuyết \\
13 & Một phần nơ-ron \nt{ACET} đáp ứng với phần thưởng và hình phạt trong khoảng hai chục mili giây, càng mạnh khi sự củng cố càng bất ngờ & Chuột nhắt, nơ-ron cholinergic của não trước nền được nhận dạng bằng quang di truyền: đáp ứng với phần thưởng và hình phạt sau $18\pm3$\,ms, độ lớn tăng theo mức bất ngờ của sự củng cố, đáp ứng giống nhau ở nơ-ron của hai nhân & \cite{Hangya2015} & đã đo \\
14 & \nt{NORA} nhận ra thay đổi bất ngờ của thế giới, \nt{ACET} giữ mạng ở chế độ ghi & Sự phân công được đề xuất trong một mô hình. Gián tiếp: ở người, đường kính đồng tử, liên quan tới \nt{NORA}, bám theo các thay đổi và tốc độ học. Chưa có kiểm tra trực tiếp cặp \nt{NORA}--\nt{ACET} & \cite{YuDayan2005,Nassar2012} & giả thuyết \\
15 & Trong giấc ngủ sóng chậm, mạng ở chế độ đọc phát lại những gì được ghi ban ngày, và những lần phát lại này chép chúng sang bộ nhớ dài hạn & Ghi các tập hợp nơ-ron hồi hải mã ở chuột cống: các tế bào cùng hoạt động ban ngày thường cùng phát xung hơn trong giấc ngủ sóng chậm sau đó, và theo cùng thứ tự --- đã đo. Về việc chép: ức chế các sóng gợn hồi hải mã sau khi học làm giảm trí nhớ không gian, còn kéo dài các sóng gợn tự phát thì cải thiện nó; việc nguyên nhân là mức \nt{ACET} thấp chưa được chứng minh & \cite{WilsonMcNaughton1994,Skaggs1996,Girardeau2009,FernandezRuiz2019,Hasselmo1999} & giả thuyết \\
\bottomrule
\end{longtable}}

\section{Chương~\ref{sec:synthesis}. Toàn bộ cỗ máy}

Là chương tổng kết, chương này không đưa ra thí nghiệm mới: mỗi dòng dựa vào chương đã phân tích khẳng định đó và vào cùng các nguồn; bus \nt{GLUT} và việc điều khiển luồng bằng \nt{GABA} đã được xác lập vững chắc, còn vai trò của các kênh quảng bá và sự phối hợp giữa chúng phần lớn là giả thuyết.

{\footnotesize
\setlength{\tabcolsep}{3pt}\renewcommand{\arraystretch}{1.15}
\begin{longtable}{>{\raggedright}p{0.5cm}>{\raggedright}p{4.0cm}>{\raggedright}p{5.6cm}>{\raggedright}p{2.3cm}>{\raggedright\arraybackslash}p{1.7cm}}
\toprule
TT & Khẳng định & Thí nghiệm và điều đã đo & Nguồn & Trạng thái \\
\midrule\endfirsthead
\toprule
TT & Khẳng định & Thí nghiệm và điều đã đo & Nguồn & Trạng thái \\
\midrule\endhead
1 & Với $8.6\cdot10^{10}$ nơ-ron chỉ có bốn loại tín hiệu điều khiển~\eqref{eq:machine} & Đếm nhân tế bào trong dịch đồng thể não (phương pháp phân đoạn đẳng hướng): nam giới trưởng thành có $86.1\pm8.1$ tỷ nơ-ron & \cite{Azevedo2009} & đã đo \\
2 & Mệnh đề~\ref{prop:architecture}, hai lớp đầu: dữ liệu chạy trên bus địa chỉ \nt{GLUT}, dấu của liên kết do bên gửi quyết định, các nơ-ron \nt{GABA} định hướng và chuẩn hóa luồng (chương~\ref{sec:neurontypes}, \ref{sec:gaba}) & Mỗi nơ-ron có một chất dẫn truyền thần kinh chính (tổng quan các phép đo); ức chế thuận thu hẹp cửa sổ trong đó nơ-ron tháp cộng các đầu vào; phép chia cho hoạt động chung đã được đo ở nhiều vùng & \cite{Strata1999,Pouille2001,Carandini2012} & đã đo \\
3 & Phần tử không tin cậy: khớp thần kinh làm mất phần lớn số xung (bảng~\ref{tab:computer}, chương~\ref{sec:code}) & Lát cắt hồi hải mã, kích thích tối thiểu: hơn một nửa số xung không gây đáp ứng ở CA1; đã loại trừ lỗi ngưỡng và lỗi dẫn truyền dọc sợi trục, nguyên nhân là sự giải phóng chất dẫn truyền mang tính xác suất & \cite{AllenStevens1994} & đã đo \\
4 & Não chạy bằng $\sim20$ oát, trung bình khoảng một xung mỗi giây trên một nơ-ron (chương~\ref{sec:code}); kỹ sư chưa chế tạo được cỗ máy như vậy & Ước tính~\eqref{eq:energy} từ chi phí đã đo: khoảng $10^9$ phân tử ATP cho mỗi xung, kể cả công của khớp thần kinh (vỏ não gặm nhấm); ước tính chi tiết cho vỏ não cho tần số còn thấp hơn. Hiện trạng kỹ thuật theo một bài tổng quan & \cite{Attwell2001,Lennie2003,MehonicKenyon2022} & lý thuyết \\
5 & Việc học của phần lớn khớp thần kinh là tích của vết cục bộ và một con số chung $\delta$ (phương trình thứ hai của~\eqref{eq:machine}, chương~\ref{sec:dofa}) & Nơ-ron \nt{DOPA} của khỉ đáp ứng với sai số dự đoán phần thưởng; trong lát cắt thể vân, dopamine làm gai lớn lên chỉ khi đến sau đầu vào glutamate $0.3$--$2$\,s; vết đã được đo & \cite{Schultz1997,Yagishita2014} & đã đo \\
6 & Dây sai số riêng tới từng đầu ra chỉ có trong mạch học vận động (chương~\ref{sec:motorlearning}) & Tiểu não: sự trùng khớp giữa tín hiệu sợi leo và một đầu vào làm yếu đầu vào đó; ở khỉ, các xung phức của tế bào Purkinje mang hướng của sai số vận động và gây ra hiệu chỉnh & \cite{Ito2001,Herzfeld2018} & đã đo \\
7 & Mỗi kênh quảng bá là một bó vài đường (mệnh đề~\ref{prop:bundle}) & Với \nt{DOPA}: trong cùng một nhiệm vụ, các nơ-ron khác nhau mang phần thưởng, cử động và đặc điểm kích thích; sai số nhanh và mức dopamine chậm là khác nhau. Với \nt{SERO}, \nt{NORA}, \nt{ACET}, sự phân chia như vậy ít được nghiên cứu hơn & \cite{Engelhard2019,Hamid2016,Mohebi2019} & đã đo \\
8 & \nt{SERO}: bộ điều chỉnh mức và, theo giả thuyết, tầm nhìn $\tau_\gamma$ (chương~\ref{sec:serocomputer}) & Tự ức chế: chất chủ vận của chính thụ thể 5-HT$_{1A}$ ức chế nơ-ron serotonin của nhân đường giữa, đã đo. Tầm nhìn được đề xuất trong lý thuyết; thí nghiệm mạnh nhất: hoạt hóa quang di truyền các nơ-ron này ở chuột nhắt kéo dài thời gian chờ phần thưởng bị hoãn & \cite{Sprouse1987,Doya2002,Miyazaki2014} & giả thuyết \\
9 & \nt{NORA}: hệ số khuếch đại $g$ chọn giữa tìm kiếm và quyết định (chương~\ref{sec:nora}) & Tăng tỷ số tín hiệu/nhiễu: ở khỉ tỉnh, noradrenaline trong vỏ não thính giác ức chế hoạt động nền mạnh hơn đáp ứng với tiếng kêu của đồng loại, đã đo; khuếch đại nhân là mô hình; vai trò trong việc chọn giữa tìm kiếm và quyết định là lý thuyết & \cite{Foote1975NE,ServanSchreiber1990,AstonJones2005} & giả thuyết \\
10 & \nt{ACET}: chuyển giữa ghi và đọc (chương~\ref{sec:acet}) & Ba tác động (làm yếu liên kết bên trong, tăng cường đầu vào, tạo thuận lợi cho tính dẻo) đã được đo; việc chuyển chế độ được thí nghiệm với scopolamine ở người ủng hộ gián tiếp & \cite{HasselmoBower1992,Gil1997,Atri2004} & giả thuyết \\
11 & Công thức~\eqref{eq:machine} mô tả toàn bộ cỗ máy & Đây là bản tóm tắt các mô hình của chương~\ref{sec:glutcomputer}--\ref{sec:acet}, mỗi phần dựa vào chương của nó; toàn bộ chưa được kiểm chứng bằng thí nghiệm & suy ra trong chương & giả thuyết \\
12 & Các núm vặn phối hợp mà không cần điều khiển chung: sai số, bất ngờ, ghi, trở lại (hình~\ref{fig:timeline}) & Chưa có thí nghiệm: sơ đồ mang tính định tính. Các mắt xích riêng lẻ: lý thuyết phân công giữa \nt{NORA} và \nt{ACET}, và mối liên hệ giữa đồng tử với tốc độ học ở người & \cite{YuDayan2005,Nassar2012} & giả thuyết \\
13 & Học theo một tín hiệu chung càng chậm khi càng nhiều nơ-ron ảnh hưởng tới kết quả (mệnh đề~\ref{prop:broadcastcost}) & Tính đường cong học của mạng tuyến tính: học bằng nhiễu loạn đầu ra chậm hơn hạ gradient một số lần bằng đúng số đầu ra & \cite{Werfel2005} & lý thuyết \\
14 & Vỏ não chỉnh các mạch nhiều lớp thế nào thì chưa biết; các ứng viên: chùm xung, tính toán trong nhánh nơ-ron, tự học bằng dự đoán & Chùm xung làm vật mang sai số là mô hình; chỉ mạng $5$--$8$ lớp mới lặp lại được đáp ứng của mô hình chi tiết một nơ-ron vỏ não, là mô hình; xung calci trong nhánh nơ-ron vỏ não người cho ra phép HOẶC loại trừ, đã đo trên lát cắt; tự học là tổng quan các bằng chứng & \cite{Payeur2021,Beniaguev2021,Gidon2020,Keller2018} & giả thuyết \\
\bottomrule
\end{longtable}}

\section{Chương~\ref{sec:sensors}. Đầu vào của cỗ máy}

Cấu tạo của cảm biến đã được đo trực tiếp: bằng bản ghi dòng điện của từng tế bào, dao động của ốc tai và việc đếm tế bào võng mạc; giới hạn độ nhạy và công thức nhiễu hạt suy ra từ vật lý, còn việc mã của cảm biến được tinh chỉnh để mang nhiều thông tin nhất là giả thuyết có cơ sở vững.

{\footnotesize
\setlength{\tabcolsep}{3pt}\renewcommand{\arraystretch}{1.15}
\begin{longtable}{>{\raggedright}p{0.5cm}>{\raggedright}p{4.0cm}>{\raggedright}p{5.6cm}>{\raggedright}p{2.3cm}>{\raggedright\arraybackslash}p{1.7cm}}
\toprule
TT & Khẳng định & Thí nghiệm và điều đã đo & Nguồn & Trạng thái \\
\midrule\endfirsthead
\toprule
TT & Khẳng định & Thí nghiệm và điều đã đo & Nguồn & Trạng thái \\
\midrule\endhead
1 & Cảm biến là bộ chuyển đổi: thụ thể tạo ra dòng điện, còn đầu ra của tế bào cảm giác ở mắt và tai là glutamate đưa lên bus \nt{GLUT}; các tế bào đầu tiên không phát xung (hình~\ref{fig:transducer}) & Tế bào que và tế bào nón của rùa giải phóng một amino acid kích thích: nó được các kênh glutamate trên một mảnh màng tách rời bắt được. Tế bào lông trong của ốc tai chuột cống: dòng điện ở đầu tận cùng sợi thần kinh đi qua thụ thể glutamate AMPA, tần số giải phóng tăng theo khử cực liên tục của tế bào tới $150$\,Hz & \cite{CopenhagenJahr1989,GlowatzkiFuchs2002} & đã đo \\
2 & Trong bóng tối, cảm biến ánh sáng đáp ứng với một photon bằng dòng khoảng một picoampe & Điện cực hút trên đoạn ngoài tế bào que của cóc: đáp ứng với chớp sáng mờ bị lượng tử hóa, mỗi sự kiện $\approx1$\,pA ($5\%$ dòng tối) với độ tản $0.2$\,pA, hiệu suất $0.5$. Người báo nhận được một photon đúng nhiều hơn mức ngẫu nhiên & \cite{Baylor1979,Tinsley2016} & đã đo \\
3 & Cảm biến âm thanh nhận ra độ dịch chuyển nhỏ hơn một nanomét & Phương pháp Mössbauer, màng đáy ốc tai chuột lang tại vị trí $16$--$19$\,kHz: ở ngưỡng đáp ứng của dây thần kinh thính giác, vận tốc màng khoảng $0.04$\,mm/s, tức độ dịch chuyển $\approx0.35$\,nm (chúng tôi quy đổi). Đo trên màng, không phải trên bó lông của cảm biến & \cite{Sellick1982,Hudspeth1989} & đã đo \\
4 & Trong bóng tối, độ chính xác bị giới hạn bởi nhiễu hạt photon~\eqref{eq:photonnoise}; ánh sáng yếu thì tích lũy lâu hơn & Công thức suy ra từ thống kê Poisson. Ở tế bào que, nhiễu dưới ánh sáng mờ trùng với tổng các sự kiện lượng tử độc lập; trên nền sáng, đáp ứng với chớp sáng nhỏ hơn và ngắn hơn. Con số ``vài phần mười giây'' ở đây chưa được kiểm tra riêng & \cite{Baylor1979} & lý thuyết \\
5 & Dải đầu vào là mười bậc với ánh sáng và mười hai bậc với âm thanh, còn tần số xung chưa tới ba bậc & Với âm thanh: đáp ứng của màng đáy ở tần số đặc trưng tăng với độ dốc tới $0.2$\,dB/dB trong dải $40$--$80$\,dB, sự nén vẫn thấy được trên $100$\,dB. Bản thân số bậc là ước tính chuẩn, trong chương không dẫn nguồn & \cite{Ruggero1997} & đã đo \\
6 & Đáp ứng của cảm biến theo~\eqref{eq:nakarushton}, còn $\sigma$ theo độ sáng trung bình, làm đường cong dịch dọc trục logarit (hình~\ref{fig:adaptation}) & Công thức lần đầu được khớp với điện thế S của võng mạc cá. Tế bào nón của rùa: đáp ứng tuân theo $V=I V_m/(I+\sigma)$ và bao $\sim3.5$ bậc độ sáng; sau $2$--$3$ phút có nền, đường cong dịch ngang và giữ nguyên độ rộng. Liên hệ với phép chia cho hoạt động chung: tổng quan & \cite{NakaRushton1966,NormannPerlman1979,Carandini2012} & đã đo \\
7 & Cảm biến truyền các thay đổi, và mã của chúng được tinh chỉnh để mang nhiều thông tin nhất trên mỗi xung (mệnh đề~\ref{prop:sensors}) & Nguyên lý được đề xuất về lý thuyết. Bằng chứng mạnh nhất: ở ruồi, đặc tuyến ``độ tương phản $\to$ đáp ứng'' của nơ-ron lớp thứ nhất trùng với phân phối tích lũy độ tương phản của cảnh tự nhiên, và nhờ đó đạt thông tin lớn nhất & \cite{Barlow1961,Laughlin1981} & giả thuyết \\
8 & Cảm biến bật và cảm biến tắt là mã hai dây; camera sự kiện lặp lại nó trên silicon & Tổng quan thí nghiệm: kênh bật và kênh tắt của võng mạc tách nhau ngay từ khớp thần kinh đầu tiên và có thể bị vô hiệu hóa riêng rẽ. Camera: $128\times128$ điểm ảnh không có khung hình, dải $120$\,dB, độ trễ $15$\,\textmu s & \cite{Schiller1992,Lichtsteiner2008,Gallego2022} & đã đo \\
9 & Mắt có $\sim10^8$ cảm biến ánh sáng và $\sim10^6$ nơ-ron đầu ra: nén $\sim100$ lần & Đếm trên toàn bộ tiêu bản võng mạc người: trung bình $92$ triệu tế bào que và $4.6$ triệu tế bào nón; tế bào đầu ra (tế bào hạch) từ $0.7$ đến $1.5$ triệu tùy mắt & \cite{Curcio1990,CurcioAllen1990} & đã đo \\
10 & Ở chuột nhắt, mắt có hơn ba mươi kênh đầu ra & Chuột nhắt: chụp calci hai photon hơn $11\,000$ tế bào đầu ra và phân cụm đáp ứng: rõ ràng nhiều hơn $30$ loại chức năng. Ở người con số này chưa được đo & \cite{Baden2016} & đã đo \\
11 & Mắt chuột lang truyền $\sim875\,000$ bit/s; quy đổi cho mắt người được $\sim10^7$ bit/s & Chuột lang, mảng đa điện cực, chuyển động trong cảnh tự nhiên: $6$--$13$ bit/s mỗi tế bào, $\sim875\,000$ bit/s cho $\sim10^5$ tế bào đầu ra. Với người ($\sim10^6$ tế bào), $\sim10^7$ bit/s là quy đổi của các tác giả & \cite{Koch2006} & đã đo \\
12 & Tai là dãy bộ lọc thông dải với bản đồ tần số gần như logarit (hình~\ref{fig:filterbank}) & Vị trí dao động lớn nhất của màng đáy phụ thuộc tần số; hàm ``tần số--vị trí'' gần như hàm mũ, cùng một dạng mô tả ốc tai người và động vật sống & \cite{Greenwood1990,Robles2001} & đã đo \\
13 & Bộ cộng hưởng là chủ động: một phần cảm biến khuếch đại dao động yếu hàng nghìn lần về biên độ ($66$--$76$\,dB), âm to lên thì hệ số khuếch đại giảm & Đo rung bằng laser ở đáy ốc tai sóc chinchilla: với âm yếu ở tần số đặc trưng, hệ số khuếch đại $66$--$76$\,dB so với xương bàn đạp; sau khi con vật chết, độ nhạy giảm $60$--$81$\,dB ($10^3$--$10^4$ lần về biên độ) và sự nén biến mất. Nguồn năng lượng là tế bào lông ngoài & \cite{Ruggero1997,Robles2001} & đã đo \\
14 & Bộ khuếch đại của tai đặt sát ranh giới tự kích & Tính toán: gần điểm rẽ nhánh Hopf, không tránh khỏi nén dải, chọn lọc tần số sắc nét và âm tổ hợp, tức mọi phi tuyến đã biết của thính giác. Việc ốc tai được chỉnh đúng như vậy chưa được chứng minh trực tiếp & \cite{Eguiluz2000} & giả thuyết \\
15 & Ở tần số thấp, xung đi cùng pha với âm (ở chuột lang, khóa pha yếu đi từ $\sim600$\,Hz và còn thấy tới $\sim3.5$\,kHz), và từ đó tìm được hướng; ở tần số cao chỉ còn mã vị trí & Dây thần kinh thính giác chuột lang: khóa pha bắt đầu giảm từ khoảng $600$\,Hz và biến mất trên $3.5$\,kHz; ở mèo và khỉ, ngưỡng này cao hơn khoảng một quãng tám. Chênh lệch thời điểm âm đến hai tai: tổng quan & \cite{PalmerRussell1986,Grothe2010} & đã đo \\
16 & Các cảm biến khác (bảng~\ref{tab:sensors}): khứu giác có $\sim400$ loại thụ thể, mỗi nơ-ron một loại, và mã tổ hợp; vị giác đi một phần qua ATP và serotonin; xúc giác có cảm biến thích nghi nhanh và chậm; nóng và lạnh tách riêng & Chuột nhắt: một thụ thể nhận nhiều mùi, một mùi kích hoạt nhiều thụ thể, các mùi khác nhau cho các tập khác nhau. Không có thụ thể ATP thì dây thần kinh vị giác không đáp ứng; nụ vị giác giải phóng serotonin. Ghi từng sợi thần kinh ở da bàn tay người: bốn loại theo tốc độ thích nghi. Kênh lạnh khác với các kênh nóng & \cite{Mombaerts2004,Malnic1999,Finger2005,Huang2005,JohanssonVallbo1979,McKemy2002} & đã đo \\
\bottomrule
\end{longtable}}

\section{Chương~\ref{sec:recognition}. Nhận dạng}

Tốc độ nhận dạng, cấu tạo của các tầng đầu, việc đọc tuyến tính câu trả lời và việc học từ sự liên tục của thời gian đã được đo ở người và khỉ; việc quy cả chuỗi về hai phép toán và việc học từ video đã được kiểm tra trên các mô hình của sách, còn cách giải thích các ảo giác đi từ có cơ sở vững đến còn tranh cãi.

{\footnotesize
\setlength{\tabcolsep}{3pt}\renewcommand{\arraystretch}{1.15}
\begin{longtable}{>{\raggedright}p{0.5cm}>{\raggedright}p{4.0cm}>{\raggedright}p{5.6cm}>{\raggedright}p{2.3cm}>{\raggedright\arraybackslash}p{1.7cm}}
\toprule
TT & Khẳng định & Thí nghiệm và điều đã đo & Nguồn & Trạng thái \\
\midrule\endfirsthead
\toprule
TT & Khẳng định & Thí nghiệm và điều đã đo & Nguồn & Trạng thái \\
\midrule\endhead
1 & So sánh với mẫu theo điểm ảnh khi có dịch chuyển đoán đúng không hơn đồng xu; cặp tầng~\eqref{eq:scstage} nhận ra hình ở bất kỳ vị trí nào & \texttt{models/\allowbreak recognition\_models.py}, ``võng mạc'' $24$ điểm, $4000$ lần thử: so mẫu cho $0.49$, $0.51$, $0.52$ với tỷ lệ điểm bị lật $0$, $0.1$, $0.2$; cặp $S$ và $C$ cho $1.00$, $0.86$, $0.70$ & suy ra trong chương & mô hình \\
2 & Con vật trong ảnh được nhận ra sau $\sim150$\,ms, một lượt qua chừng mười tầng, mỗi tầng $10$--$20$\,ms & Người, nhiệm vụ ``có con vật hay không'' trên ảnh hiện $20$\,ms: điện thế gợi của hai câu trả lời tách nhau sau $\sim150$\,ms. Khỉ: thời điểm đáp ứng đầu tiên tăng dần từ tầng này sang tầng khác (V1 sớm nhất, rồi V2 và V4). Số tầng và ``không hoặc một xung mỗi tầng'' được suy ra từ các con số này & \cite{Thorpe1996,Schmolesky1998} & đã đo \\
3 & Tầng $S$ là VÀ theo bộ phận, tầng $C$ là giá trị lớn nhất theo vị trí (mệnh đề~\ref{prop:conveyor}) & Mèo, vỏ não thị giác sơ cấp: tế bào đơn giản đáp ứng với cạnh có hướng nhất định ở vị trí nhất định, tế bào phức tạp đáp ứng với cùng hướng đó trong vùng rộng hơn. Ghi nội bào tế bào phức tạp: ở nhiều tế bào, đáp ứng với hai vạch gần bằng đáp ứng lớn hơn trong hai, trung bình gần nhất với phép max. Giá trị lớn nhất qua ức chế lẫn nhau: Mệnh đề~\ref{prop:wta}; phép chia cho hoạt động chung: tổng quan & \cite{HubelWiesel1962,Lampl2004,Carandini2012} & đã đo \\
4 & Càng lên cao trong chuỗi, tổ hợp càng lớn và phức tạp, câu trả lời càng ít phụ thuộc vị trí & Khỉ, đơn giản hóa kích thích tới ``đặc trưng tới hạn'': tế bào V4 và vỏ não thái dương dưới phía sau đáp ứng với tổ hợp phức tạp của hình dạng, màu sắc và kết cấu, với trường tiếp nhận nhỏ; ở vỏ não thái dương dưới phía trước, trường lớn hơn. Sự xen kẽ $S$ và $C$ trong mô hình tái hiện bức tranh này & \cite{KobatakeTanaka1994,Riesenhuber1999} & đã đo \\
5 & Mạng tích chập lặp lại sự xen kẽ này và dự đoán tốt nhất đáp ứng của các tầng trên; đồ vật được đọc tuyến tính từ tần số của một nhóm nơ-ron & Mạng được huấn luyện để nhận dạng, không khớp với não: lớp trên cùng dự đoán đáp ứng của nơ-ron vỏ não thái dương dưới của khỉ, các lớp giữa dự đoán đáp ứng của V4. Bộ phân loại tuyến tính dựa trên hoạt động của $\sim100$ tế bào chọn ngẫu nhiên, trong cửa sổ từ $12.5$\,ms, nhận ra đồ vật và phạm trù ở các vị trí và kích thước khác nhau & \cite{Fukushima1980,Yamins2014,Hung2005} & đã đo \\
6 & Quy tắc Hebb có vết~\eqref{eq:traceinvariance} dạy tính bất biến từ video không nhãn nếu vết không ngắn hơn $\sim20$\,ms & Ý tưởng đặc trưng chậm và quy tắc có vết là lý thuyết. \texttt{models/\allowbreak recognition\_models.py}, hai nơ-ron $C$, $16$ đầu vào, $10$ lần chạy: nghỉ giữa các đồ vật $0.3$\,s. Với $\tau_{\text{tr}}=5$\,ms, trùng với hình trung bình $0.52$, với $10$\,ms là $0.56$; với $20$, $50$ và $200$\,ms là $1.00$ ở mọi lần chạy (với $30$\,ms, một trong mười lần chạy sai) & \cite{Foldiak1991,WiskottSejnowski2002} & mô hình \\
7 & Trong não, tính bất biến được học từ sự liên tục của thời gian & Khỉ: đồ vật bị tráo kín đáo trong lúc mắt nhảy tới một vị trí nhất định của trường nhìn; tính bất biến vị trí của nơ-ron vỏ não thái dương dưới thay đổi theo sự đánh tráo, có ý nghĩa thống kê chỉ sau một giờ. Việc đây là quy tắc chính của học thị giác là giả thuyết & \cite{LiDiCarlo2008} & đã đo \\
8 & Chuyển động: bộ phát hiện hướng, rồi cùng cặp VÀ/HOẶC trên những vùng lớn & Bộ phát hiện hướng trong võng mạc thỏ; ở vùng MT của khỉ, cả $110$ nơ-ron được ghi đều chọn lọc theo hướng, đáp ứng với chuyển động mạnh hơn ở V1 & \cite{Barlow1965,Albright1984} & đã đo \\
9 & Các tầng trên gửi xuống một dự đoán, đi lên là sai số & Mô hình giải thích các hiệu ứng ngoài trường cổ điển ở vỏ não thị giác sơ cấp; một số quan sát riêng lẻ phù hợp (tổng quan), như một cấu tạo chung của dây chuyền thì chưa được chứng minh & \cite{RaoBallard1999,Keller2018} & giả thuyết \\
10 & Ảnh khó cần phản hồi và vài vòng & Khỉ: với ảnh ``khó'', quyết định ở vỏ não thái dương dưới hình thành muộn hơn $\sim30$\,ms; các đáp ứng muộn này được mạng đi thẳng dự đoán kém, còn mạng có phản hồi hoặc mạng sâu hơn dự đoán tốt hơn & \cite{Kar2019} & đã đo \\
11 & Sửa điểm ảnh không nhận ra được đánh lừa mạng; thị giác người bền vững hơn nhiều nhưng không miễn nhiễm & Trên mạng: một nhiễu loạn nhỏ được chọn đặc biệt làm đổi câu trả lời; ví dụ ``gấu trúc $\to$ vượn'' lấy từ công trình thứ hai. Ở người, khi hiện ngắn, các ảnh sửa có khả năng chuyển tốt giữa các mạng làm lệch lựa chọn & \cite{Szegedy2014,Goodfellow2015,Elsayed2018} & đã đo \\
12 & Ảo giác kỳ vọng: đầu vào không đáng tin nhường cho kỳ vọng; vật nhạt chuyển động ``chậm hơn'', chiếc váy được thấy khác nhau (mục~\ref{sec:illusions}) & Cách tử có độ tương phản thấp hơn trông như chuyển động chậm hơn. Khảo sát $1401$ người: chiếc váy được thấy theo ba phương án ổn định, gợi ý về ánh sáng làm đổi màu; ở $\sim13\,000$ người quan sát, giả định về ánh sáng quyết định. Mô hình Bayes với kỳ vọng về chuyển động chậm giải thích một loạt ảo giác chuyển động & \cite{Thompson1982,LaferSousa2015,Wallisch2017,Weiss2002,Purves2011} & giả thuyết \\
13 & Điều chỉnh khuếch đại: thác nước, độ nghiêng và độ tương phản; lưới Hermann không phải ức chế láng giềng & Bộ phát hiện hướng ở võng mạc thỏ giảm tần số khi chuyển động kéo dài, và sau khi dừng thì tụt xuống dưới mức nền. Ảo giác độ nghiêng được tái hiện bằng mô hình chia cho hoạt động của láng giềng. Những biến đổi lưới không làm đổi đáp ứng của nơ-ron đầu ra võng mạc lại làm mất các đốm: cách giải thích bằng ức chế láng giềng bị bác bỏ & \cite{BarlowHill1963,Schwartz2009,SchillerCarvey2005} & đã đo \\
14 & Bù độ trễ: vật đang chuyển động có vẻ nằm trước điểm lóe & Ảo giác đã được đo. Tâm vật lý học mâu thuẫn với cả việc ngoại suy chuyển động lẫn chênh lệch độ trễ; một cách giải thích hồi tố đã được đề xuất, dựa trên các sự kiện $\sim80$\,ms sau điểm lóe. Tranh luận chưa ngã ngũ & \cite{Nijhawan1994,EaglemanSejnowski2000} & giả thuyết \\
15 & ``Rắn'' đứng yên lại quay: chênh lệch độ trễ~\eqref{eq:latency} được bộ phát hiện hướng đọc thành chuyển động & Ảo giác vẫn xảy ra khi không có màu; từng cặp phần tử kề nhau tự tạo ra nó; nơ-ron vỏ não khỉ chọn lọc theo hướng cho đáp ứng có hướng với các cặp đứng yên này. Ở người, chuyển động quay được kích hoạt bởi vi giật mắt và cái chớp mắt & \cite{Conway2005,OteroMillan2012} & đã đo \\
16 & Ảnh hai nghĩa: ức chế lẫn nhau, sự mỏi và nhiễu; sự chuyển chủ yếu do nhiễu gây ra & Mô hình các nhóm nơ-ron cạnh tranh: sự chuyển trong mô hình do \emph{nhiễu} gây ra, không có nhiễu thì không chuyển; mô hình giải thích việc tần số chuyển tăng theo cường độ kích thích. Sự mỏi ở đây đóng vai trò nhỏ hơn & \cite{MorenoBote2007} & giả thuyết \\
17 & Một phần ảo giác là hệ quả của bài toán mà cỗ máy học & Mạng được huấn luyện để dự đoán khung hình tiếp theo của video ``thấy'' các vòng tròn tĩnh quay và không thấy điều đó trên ảnh đối chứng; mạng khử nhiễu và làm nét ảnh mắc các ảo giác màu sắc và độ sáng như con người & \cite{Watanabe2018,GomezVilla2020} & đã đo \\
\bottomrule
\end{longtable}}

\section{Chương~\ref{sec:muscles}. Đầu ra tới cơ}

Cấu tạo của đầu ra (đơn vị vận động, hệ số an toàn của khớp thần kinh, bật theo kích thước, cặp cơ, vòng phản xạ giãn và bộ tạo nhịp) đã được đo trực tiếp; điều kiện ổn định của vòng có trễ là một định lý; mô hình bên trong của cơ thể là giả thuyết có cơ sở vững; các con số trên hình được tính theo mô hình của sách.

{\footnotesize
\setlength{\tabcolsep}{3pt}\renewcommand{\arraystretch}{1.15}
\begin{longtable}{>{\raggedright}p{0.5cm}>{\raggedright}p{4.0cm}>{\raggedright}p{5.6cm}>{\raggedright}p{2.3cm}>{\raggedright\arraybackslash}p{1.7cm}}
\toprule
TT & Khẳng định & Thí nghiệm và điều đã đo & Nguồn & Trạng thái \\
\midrule\endfirsthead
\toprule
TT & Khẳng định & Thí nghiệm và điều đã đo & Nguồn & Trạng thái \\
\midrule\endhead
1 & Đơn vị vận động: một nơ-ron \nt{ACET} điều khiển một nhóm sợi, từ vài trăm đến vài nghìn; ở cơ cần tinh chỉnh, đơn vị nhỏ hơn & Cơ người, đếm sợi trục vận động và sợi cơ: trung bình khoảng $340$ sợi trên một nơ-ron ở cơ gian cốt bàn tay, khoảng $410$ ở cơ cánh tay quay (brachioradialis), khoảng $1930$ ở cơ bụng chân trong. Chương không nêu con số chính xác cho những đơn vị nhỏ nhất. & \cite{Feinstein1955mu} & đã đo \\
2 & Khớp thần kinh trên cơ giải phóng lượng chất dẫn truyền gấp vài lần mức cần để sợi cơ đáp ứng & Tổng quan các phép đo ở khớp thần kinh--cơ: hệ số an toàn (tỷ số giữa lượng chất dẫn truyền được giải phóng và lượng ngưỡng) ở động vật có vú trưởng thành thường là $3$--$5$; ở người, hệ số an toàn dựa vào các nếp gấp của màng bên nhận nhiều hơn là vào lượng giải phóng. & \cite{WoodSlater2001} & đã đo \\
3 & Cú giật~\eqref{eq:twitch} tăng lên trong thời gian $T_c$ cỡ $50\ms$ & Từng đơn vị vận động bàn tay người khi gắng sức tự chủ, lực được lấy trung bình theo xung của đơn vị: thời gian tăng $30$--$100\ms$, ở $80\,\%$ đơn vị dưới $70\ms$. Hàm $(t/T_c)e^{1-t/T_c}$ là xấp xỉ lấy từ mô hình quần thể đơn vị. & \cite{MilnerBrown1973rec, Fuglevand1993} & đã đo \\
4 & Tổng các cú giật~\eqref{eq:forcesum}: với $5$, $15$ và $40$ xung/s, lực là $0.2$--$1.1F_1$, $1.8$--$2.2F_1$ và khoảng $5.4F_1$ (hình~\ref{fig:tetanus}) & Tính theo \texttt{models/\allowbreak muscle\_\allowbreak models.py}, $T_c=50\ms$: $0.21$--$1.10$, $1.76$--$2.19$ và $5.28$--$5.49\,F_1$ trong $0.3\,\text{s}$ cuối. Phép cộng tuyến tính là đơn giản hóa: mô hình không có sự bão hòa của cơ thật. & \texttt{models/\allowbreak muscle\_\allowbreak models.py} & mô hình \\
5 & Các đơn vị được bật theo kích thước; đơn vị yếu nhất yếu hơn đơn vị mạnh nhất cả trăm lần & Rễ trước tủy sống của mèo: khi đầu vào phản xạ tăng dần, những nơ-ron có xung nhỏ trong rễ phóng điện trước, tức là nơ-ron nhỏ. Cơ gian cốt bàn tay người: lực giật của các đơn vị từ $0.1$ đến $10$~g và tăng gần tuyến tính theo mức lực tại đó đơn vị được bật. & \cite{Henneman1957, MilnerBrown1973rec} & đã đo \\
6 & Bước lực tỉ lệ với lực, $\Delta F/F\approx q-1$~\eqref{eq:sizeprinciple}: dấu phẩy động & Suy ra trong chương từ cấp số nhân của các lực. Phù hợp với thực nghiệm: khi gắng sức lớn, cùng một mức tăng lực chỉ cần bật thêm ít đơn vị mới hơn nhiều. & suy ra trong chương; \cite{MilnerBrown1973rec} & lý thuyết \\
7 & Độ tản của lực tăng tỉ lệ với chính lực & Gắng sức đẳng trường tự chủ của cơ ngón cái: độ lệch chuẩn của lực tăng tuyến tính theo lực trung bình; với cùng mức lực tạo bằng kích thích điện, độ tản không đổi. Mô hình quần thể: tăng tuyến tính là do bật đơn vị theo thứ tự lực. & \cite{Jones2002sdn} & đã đo \\
8 & Sự mượt mà của cử động là hệ quả của nhiễu phụ thuộc tín hiệu & Mô hình: quỹ đạo cực tiểu hóa độ tản ở điểm cuối dưới loại nhiễu này tái hiện các quỹ đạo đã đo của mắt và tay. Chưa có thí nghiệm trực tiếp tách nguyên nhân này khỏi các nguyên nhân khác. & \cite{HarrisWolpert1998} & giả thuyết \\
9 & Hiệu lực của cặp cơ quyết định mô-men, tổng quyết định độ cứng~\eqref{eq:antagonists} & Lý thuyết điều khiển trở kháng bằng đồng co. Cánh tay người: độ cứng của mỗi khớp, đo bằng những cú đẩy nhỏ, xấp xỉ tỉ lệ với mô-men tại khớp đó; các tỷ lệ đồng co khác nhau thay đổi hình dạng và hướng của elip độ cứng ở bàn tay. & \cite{Hogan1984, GomiOsu1998stiff} & đã đo \\
10 & Cảm biến độ giãn kích thích nơ-ron \nt{ACET} của nó và qua nơ-ron trung gian ức chế tắt cơ đối vận (hình~\ref{fig:antagonists}) & Tủy sống mèo: một nơ-ron trung gian đơn lẻ, được kích thích bởi sợi của cảm biến độ giãn, tạo điện thế ức chế qua một khớp thần kinh ở các nơ-ron \nt{ACET} của cơ đối vận, độ trễ qua khớp thần kinh $0.28$--$0.42\ms$. Chất dẫn truyền của nơ-ron trung gian (glycine) không được xác định trong thí nghiệm này. & \cite{Jankowska1972Ia} & đã đo \\
11 & Độ trễ của vòng: tủy sống $25$--$30\ms$, qua vỏ não dài hơn một rưỡi đến ba lần, qua mắt $100$--$200\ms$ & Cánh tay người, đẩy vào khớp: đáp ứng ngắn của cơ sau $20$--$45\ms$, đáp ứng dài $45$--$105\ms$, đáp ứng tự chủ $120$--$180\ms$. Dịch vị trí nhìn thấy của ngón tay trong lúc cử động: hiệu chỉnh trung bình sau $160\ms$. & \cite{Pruszynski2008rapid, SaundersKnill2003vis} & đã đo \\
12 & $\dd x/\dd t=-Gx(t-d)$ ổn định khi $Gd<\pi/2$~\eqref{eq:delaystable}; tại ranh giới, tần số là $1/(4d)$ & Định lý về nghiệm của phương trình có trễ. Kiểm tra theo \texttt{models/\allowbreak muscle\_\allowbreak models.py}, $d=30\ms$: tới $3\,\text{s}$, biên độ là $3\cdot10^{-6}$ khi $G=40$, $0.13$ khi $G=50$, $38$ khi $G=55$ mỗi giây (ranh giới $52$). & \cite{Hayes1950} & lý thuyết \\
13 & Run sinh lý $8$--$12$~Hz; vòng phản xạ giãn là một trong các nguồn của nó & Tổng quan các phép đo: run nhẹ trong dải quanh $10$~Hz có ở người khỏe mạnh; nguồn gốc của nó hỗn hợp (dao động trung ương, đặc tính phóng điện của đơn vị, cộng hưởng cơ học và cộng hưởng của vòng phản xạ), và trong những điều kiện khác nhau thì yếu tố khác nhau chiếm ưu thế. & \cite{McAuleyMarsden2000} & giả thuyết \\
14 & Não dự đoán kết quả của lệnh theo mô hình bên trong của cơ thể & Người cầm vật nặng bằng cách kẹp ngón và cử động tay: với tải quán tính, nhớt và đàn hồi, lực kẹp thay đổi đồng thời với tải chứ không trễ theo vòng phản hồi, tức là tải đã được dự đoán. Bài tổng quan tập hợp các thí nghiệm như vậy; bản thân mô hình chưa được quan sát trực tiếp trong nơ-ron. & \cite{FlanaganWing1997grip, WolpertGhahramani2000} & giả thuyết \\
15 & Nhịp đi, thở, nhai do các mạng cục bộ sinh ra khi đầu vào không đổi & Tổng quan thí nghiệm: hệ thần kinh bị cô lập (tủy sống, hạch) tạo ra mẫu vận động có nhịp mà không cần đầu vào có nhịp và không cần phản hồi từ cơ. & \cite{MarderBucher2001} & đã đo \\
16 & Ức chế lẫn nhau kèm sự mỏi~\eqref{eq:halfcentre} cho nhịp với chu kỳ $0.84\,\text{s}\approx2\tau_a$ (hình~\ref{fig:halfcentre}) & Định lý: mạng ức chế lẫn nhau có thích nghi, không có trạng thái ổn định khi đầu vào không đổi, nhất thiết phải dao động. Tính theo \texttt{models/\allowbreak muscle\_\allowbreak models.py} ($I=1$, $\beta=2$, $b=2$, $\tau=20\ms$, $\tau_a=0.4\,\text{s}$): chu kỳ $0.84\,\text{s}$. Việc bộ tạo nhịp thật được cấu tạo đúng như vậy chưa được chứng minh. & \cite{Matsuoka1985osc} & mô hình \\
\bottomrule
\end{longtable}}

\section{Chương~\ref{sec:pain}. Đau}

Cảm biến đau, hai dây, cổng trong tủy sống, núm âm lượng đi xuống, nhạy cảm hóa và việc chiếm sự chú ý đã được đo trực tiếp; các công thức cổng và nhạy cảm hóa là mô hình đơn giản hóa của sách; quy tắc mà cỗ máy dùng để chọn âm lượng báo động là giả thuyết.

{\footnotesize
\setlength{\tabcolsep}{3pt}\renewcommand{\arraystretch}{1.15}
\begin{longtable}{>{\raggedright}p{0.5cm}>{\raggedright}p{4.0cm}>{\raggedright}p{5.6cm}>{\raggedright}p{2.3cm}>{\raggedright\arraybackslash}p{1.7cm}}
\toprule
TT & Khẳng định & Thí nghiệm và điều đã đo & Nguồn & Trạng thái \\
\midrule\endfirsthead
\toprule
TT & Khẳng định & Thí nghiệm và điều đã đo & Nguồn & Trạng thái \\
\midrule\endhead
1 & Ngưỡng cao: ngưỡng nhiệt của các cảm biến đau khác nhau từ $41^\circ$ đến $46$--$48^\circ$ & Ghi từng sợi thần kinh. Da khỉ: ngưỡng trung vị của cảm biến chậm (C) là $41^\circ$, của cảm biến nhanh loại hai là $46^\circ$, loại một trên $53^\circ$. Da người: cảm biến C đáp ứng cả với áp lực là $40.7^\circ$, không đáp ứng với áp lực là $48^\circ$. & \cite{Treede1995heat, Weidner1999cfib} & đã đo \\
2 & Cảm biến đau thích nghi yếu hơn nhiều so với cảm biến xúc giác và sau tổn thương nhẹ trở nên nhạy hơn; suy giảm sau khi bị nung nóng mạnh đã đo ở một loại & Bỏng da nhẹ ($50^\circ$, $100\,\text{s}$): sau $5$--$10$~phút, ngưỡng đau ở người và ngưỡng của cảm biến C ở khỉ giảm $2$--$6^\circ$; các cảm biến da khác không có sự dịch chuyển như vậy. Ở cảm biến nhanh loại hai, đáp ứng sau khi nung nóng mạnh bị ức chế. Việc so sánh sự thích nghi với cảm biến xúc giác là một ước tính chung, ở đây không có thí nghiệm riêng hỗ trợ. & \cite{LaMotte1982hyper, Treede1995heat} & đã đo \\
3 & Hai dây: nhanh $5$--$30$~m/s, chậm $0.5$--$2$~m/s; thời gian đến $t=L/v$~\eqref{eq:paintime} & Tốc độ dẫn truyền: cảm biến đau nhanh của khỉ trung bình $15$ và $25$~m/s (hai loại); cảm biến C của người $0.8$--$1.0$~m/s. Với $L=1$~m, đó là vài chục mili giây và khoảng một giây. & \cite{Treede1995heat, Weidner1999cfib} & đã đo \\
4 & Cơn đau đến hai lần: trước là đau nhanh và chính xác, sau là đau âm ỉ và kéo dài & Thời gian phản ứng với nhiệt ở cẳng tay người: từ $1100$ xuống $400\ms$ khi nhiệt độ tăng; nung nóng mạnh da có lông cho đau kép, lòng bàn tay thì không. Cơn đau thứ nhất chỉ có thể do các sợi nhanh hơn $6$~m/s mang; theo thời gian tiềm, ứng với nó là cảm biến nhanh loại hai, loại không có ở lòng bàn tay. Việc phân vai (``ở đâu'' và ``tệ đến mức nào'') là cách diễn giải. & \cite{CampbellLaMotte1983first, Treede1995heat} & đã đo \\
5 & Cổng: nơ-ron đầu ra của tủy sống nhận đau và chạm, nơ-ron trung gian ức chế được cái chạm kích thích (hình~\ref{fig:gate}) & Sơ đồ cổng được đề xuất như một lý thuyết. Chuột nhắt, đánh dấu di truyền và loại bỏ các nhóm nơ-ron: nơ-ron kích thích chứa somatostatin cần cho đau cơ học; nơ-ron ức chế chứa dynorphin ngăn đầu vào từ cảm biến xúc giác tới được chúng; không có các nơ-ron này, cái chạm nhẹ gây đau. & \cite{MelzackWall1965, Duan2014} & đã đo \\
6 & Cái chạm làm dịu cơn đau & Người, xung laser gây đau trên cánh tay: cái chạm đồng thời làm giảm mức đau được đánh giá và khả năng phân biệt năng lượng xung; hiệu ứng yếu đi theo khoảng cách giữa điểm chạm và điểm đau trong cùng một đoạn da. & \cite{Mancini2014touch} & đã đo \\
7 & Giả định của thuyết cổng: đau mạnh tự tắt nơ-ron trung gian ức chế, và cổng mở toang & Trong chương được ghi rõ là giả định của thuyết ban đầu. Công trình trên chuột nhắt mô tả cách cổng chặn cái chạm đi vào kênh đau; việc cơn đau tắt nơ-ron trung gian không được chứng minh trong đó. & \cite{MelzackWall1965} & giả thuyết \\
8 & Cổng~\eqref{eq:gate}: ngưỡng $0.18$ khi không chạm, $0.86$ khi có chạm, $0.11$ khi $g=2$; khi $\nu=1$, cái chạm giảm đầu ra từ $1$ xuống $0.25$, khi $\nu=2$ không có tác dụng; bản thân cái chạm không đi qua (hình~\ref{fig:gatecurves}) & Tính theo \texttt{models/\allowbreak pain\_\allowbreak gate.py} với trọng số trong bài tái hiện mọi con số này. & \texttt{models/\allowbreak pain\_\allowbreak gate.py} & mô hình \\
9 & Các đường đi xuống từ thân não thay đổi hệ số khuếch đại của cổng theo cả hai chiều & Chuột cống, hành não, ghi khi nung nóng đuôi: một số nơ-ron phóng điện trước khi rụt đuôi (``bật''), số khác im lặng (``tắt''); kích thích yếu vùng này ức chế phản xạ. Tổng quan: cùng các mạch ấy vừa tạo thuận lợi, vừa ức chế việc truyền cơn đau, và opioid tác động thông qua chúng. & \cite{Fields1983rvm, Fields2004} & đã đo \\
10 & Sự trông đợi được đỡ đau làm dịu cơn đau qua kênh opioid & Người bị đau cấp: ở những người mà sự trông đợi làm giảm đau, việc chặn thụ thể opioid đưa cơn đau trở lại mức của những người còn lại; ở những người còn lại, việc chặn không làm đổi cơn đau. & \cite{Levine1978} & đã đo \\
11 & Não chọn âm lượng báo động theo lợi ích của việc ngắt & Chưa có thí nghiệm nào đo được quy tắc chọn âm lượng. & --- & giả thuyết \\
12 & Nhạy cảm hóa: sau tổn thương, đầu ra của cổng tăng, ngưỡng giảm, một phần là do tủy sống; nó vẫn còn sau khi kích thích gây tổn thương đã ngừng & Chuột cống: sau tổn thương da, ngưỡng phản xạ gập giảm, còn đáp ứng tăng; phân tích điện sinh lý cho thấy một phần sự tăng nảy sinh ngay trong tủy sống. Kích thích gây tổn thương gây ra rối loạn cảm giác kéo dài, tồn tại lâu hơn chính kích thích. Chương không nêu thời gian suy giảm chính xác. & \cite{Woolf1983} & đã đo \\
13 & Nhạy cảm hóa~\eqref{eq:sensitization} là phản hồi dương; khi vòng mạnh, báo động có thể tự chốt sau khi vết thương đã lành & Suy ra từ mô hình~\eqref{eq:sensitization} (bài tập cuối chương). Chưa có thí nghiệm cho thấy cơn đau dai dẳng sau khi lành là một vòng tự chốt đúng kiểu này. & suy ra trong chương & giả thuyết \\
14 & Cơn đau là phần thưởng âm, và việc học từ nó đi theo sai số sai phân thời gian & Máy chụp cộng hưởng từ, người, chuỗi tín hiệu báo trước cơn đau: hoạt động của thể vân bụng và vỏ não thùy đảo trước trùng với tín hiệu học theo sai phân thời gian. Khỉ: một phần nơ-ron \nt{DOPA} bị kích thích bởi luồng hơi khó chịu và tín hiệu báo trước nó, phần khác bị ức chế. & \cite{Seymour2004, Matsumoto2009} & đã đo \\
15 & Cơn đau ngắt công việc đang làm & Người, kích thích điện gây đau và các phép thử âm thanh: đáp ứng với phép thử $100\ms$ sau khi cơn đau bắt đầu chậm đi rõ rệt, sau $1.5\,\text{s}$ không còn khác biệt với đối chứng. Bài tổng quan tập hợp các thí nghiệm như vậy thành mô hình chiếm sự chú ý. & \cite{Crombez1996disrupt, EcclestonCrombez1999} & đã đo \\
16 & Phản xạ rụt khép kín trong tủy sống & Chuột cống: nơ-ron ở các lớp sâu của sừng sau lặp lại khuôn mẫu không gian của phản xạ rụt của từng cơ; không thể kích thích ngược chiều chúng từ đoạn cổ, tức đó là các nơ-ron trung gian tủy sống. & \cite{Schouenborg1995wr} & đã đo \\
\bottomrule
\end{longtable}}

\section{Chương~\ref{sec:motorlearning}. Học vận động}

Quy tắc bình phương tối thiểu và lợi ích của dây sai số riêng là các định lý; cấu tạo của sơ đồ có dây sai số, sự làm yếu khi trùng khớp, việc chỉnh vết theo độ trễ, hậu hiệu ứng và sự quay lại tự phát của kỹ năng đã được đo; việc cả sơ đồ hoạt động như học có giám sát và xây mô hình bên trong là giả thuyết hàng đầu; các con số của mô hình hai quá trình được tính theo mô hình của sách.

{\footnotesize
\setlength{\tabcolsep}{3pt}\renewcommand{\arraystretch}{1.15}
\begin{longtable}{>{\raggedright}p{0.5cm}>{\raggedright}p{4.0cm}>{\raggedright}p{5.6cm}>{\raggedright}p{2.3cm}>{\raggedright\arraybackslash}p{1.7cm}}
\toprule
TT & Khẳng định & Thí nghiệm và điều đã đo & Nguồn & Trạng thái \\
\midrule\endfirsthead
\toprule
TT & Khẳng định & Thí nghiệm và điều đã đo & Nguồn & Trạng thái \\
\midrule\endhead
1 & Quy tắc~\eqref{eq:lms} trung bình đi theo gradient của bình phương sai số & Kết quả của lý thuyết lọc thích nghi: hiệu chỉnh tỉ lệ với đầu vào và sai số đầu ra là hạ gradient ngẫu nhiên của sai số bình phương trung bình. & \cite{WidrowHoff1960} & lý thuyết \\
2 & Có dây sai số tới mỗi đầu ra thì tốc độ không phụ thuộc số đầu ra; với một con số chung thì tốc độ giảm theo số nơ-ron có nhiễu (mệnh đề~\ref{prop:teachers}) & Đường cong học của mạng tuyến tính: học bằng nhiễu loạn ngẫu nhiên các nơ-ron chậm hơn gradient chính xác một số lần bằng đúng số nơ-ron bị nhiễu loạn. & \cite{Werfel2005} & lý thuyết \\
3 & Sơ đồ có dây sai số hoạt động như học có giám sát (mệnh đề~\ref{prop:errorwire}) & Các lý thuyết suy ra từ cấu trúc của sơ đồ. Các thí nghiệm ở dòng 7--10 phù hợp với chúng; việc cả sơ đồ hoạt động đúng như vậy chưa được chứng minh. & \cite{Marr1969, Albus1971} & giả thuyết \\
4 & Lớp mở rộng: khoảng $7\cdot10^{10}$ nơ-ron \nt{GLUT} nhỏ, nhiều hơn cả phần còn lại của não & Đếm nhân trong huyền phù mô đã hòa tan: tiểu não người có $69\cdot10^9$ nơ-ron trong tổng số $86\cdot10^9$ của cả não. Phương pháp lập thể học: $101\cdot10^9$ nơ-ron nhỏ trong tiểu não của nam giới cao tuổi. & \cite{Azevedo2009, Andersen1992cereb} & đã đo \\
5 & Nâng số chiều của đầu vào làm quan hệ cần thiết trở thành tuyến tính & Định lý về số cách phân chia: tổng có trọng số của $n$ đầu vào kèm ngưỡng hiện thực được một cách gán nhãn tùy ý cho khoảng $2n$ vectơ ở vị trí tổng quát. Việc tiểu não dùng đúng điều này là một phần của giả thuyết ở dòng~3. & \cite{Cover1965} & lý thuyết \\
6 & Khoảng $3\cdot10^7$ nơ-ron \nt{GABA} đầu ra, mỗi nơ-ron cỡ $10^5$ đầu vào & Lập thể học tiểu não người: $30.5\cdot10^6$ nơ-ron đầu ra. Kính hiển vi điện tử, chuột cống: khoảng $1.75\cdot10^5$ đầu vào từ lớp mở rộng trên một nơ-ron đầu ra. Chất dẫn truyền ức chế của nơ-ron đầu ra: trong bài tổng quan. & \cite{Andersen1992cereb, NapperHarvey1988, Ito2001} & đã đo \\
7 & Đúng một dây sai số cho mỗi nơ-ron đầu ra, khoảng mỗi giây một lần, mỗi lần một đợt bùng mạnh & Tổng quan các bản ghi: mỗi nơ-ron đầu ra nhận một đầu vào \nt{GLUT} từ sợi leo, gây ra xung phức với tần số cỡ $1$~Hz. & \cite{Ito2001} & đã đo \\
8 & Xác suất của đợt bùng phụ thuộc vào hướng của sai số vận động & Khỉ, vùng vận nhãn của tiểu não, thích nghi của vận động giật nhãn cầu: hướng của sai số thị giác quyết định xác suất xung phức, độ lớn quyết định thời điểm của nó; đợt bùng làm đổi các xung đơn trong lần thử tiếp theo và dịch cử động dọc theo sai số ưa thích của tế bào. & \cite{Herzfeld2018} & đã đo \\
9 & Sự trùng khớp giữa đợt bùng sai số và một đầu vào đang hoạt động làm yếu đầu vào đó ($-\eta e_i x_j$) & Tổng quan thí nghiệm: kích thích đồng thời đầu vào từ lớp mở rộng và dây sai số gây suy yếu lâu dài đầu vào đó; kích thích riêng một đầu vào thì không. & \cite{Ito2001} & đã đo \\
10 & Độ dài của vết được chỉnh theo độ trễ của sai số: với độ trễ $100\ms$, đầu vào bị làm yếu mạnh nhất là đầu vào hoạt động trước đó $100\ms$ & Lát cắt và chuột nhắt sống: ở vùng phụ trách học vận nhãn, sự suy yếu được chỉnh hẹp theo khoảng cách khoảng $120\ms$ giữa đầu vào và sai số, đúng bằng thời gian tín hiệu sai số đi tới trong nhiệm vụ này; ở một vùng khác, các tế bào khác nhau được chỉnh theo các khoảng khác nhau. & \cite{Suvrathan2016} & đã đo \\
11 & Sơ đồ là một bộ lọc thích nghi~\eqref{eq:adaptivefilter} học mô hình bên trong của cơ thể & Mô hình suy ra từ cấu trúc của sơ đồ. Chưa có thí nghiệm trực tiếp đo bộ lọc đã học như một hàm truyền của cơ thể. & \cite{Fujita1982} & giả thuyết \\
12 & Dự đoán trừ đi hệ quả của chính cử động của mình, nên không thể tự cù mình & Người: cái chạm của chính mình ít gây nhột hơn cùng cái chạm đó của người khác; trong máy chụp, vỏ não cảm giác thân thể đáp ứng yếu hơn với cái chạm của chính mình, còn tiểu não ít hoạt động hơn khi cử động chạm vào da. Cách giải thích bằng việc trừ đi dự đoán là giả thuyết. & \cite{Blakemore1998} & giả thuyết \\
13 & Khi có lực ngoài, độ trượt giảm dần, còn sau khi bỏ lực, tay trượt về phía ngược lại & Người với tay bằng tay cầm của robot trong trường lực: quỹ đạo trở lại thẳng; sau khi bỏ trường, hậu hiệu ứng gần như là ảnh gương của những biến dạng ban đầu; nó còn chuyển sang cả những vùng không gian không được luyện tập. & \cite{ShadmehrMussaIvaldi1994} & đã đo \\
14 & Có hai quá trình học~\eqref{eq:tworate}; sau nhiễu loạn ngược, kỹ năng tự quay lại & Người, trường lực, sau đó một trường ngược ngắn và các lần thử có sai số bị kẹp: hiệu chỉnh tự quay về hiệu chỉnh cũ; mô hình hai quá trình tái hiện được điều này, mô hình một quá trình thì không. & \cite{Smith2006} & đã đo \\
15 & Các con số của hình~\ref{fig:tworate}: $0.84$ (chậm $0.63$) sau $200$ cử động; $6$ cử động ngược; nhanh $-0.53$, chậm $0.46$; quay lại tới $0.37$ sau $40$ cử động & Tính theo \texttt{models/\allowbreak motor\_\allowbreak learning.py}, $A_f=0.92$, $B_f=0.1$, $A_s=0.996$, $B_s=0.02$: $0.840$ ($0.634$ và $0.205$), $6$ cử động, $-0.531$ và $0.457$, đỉnh $0.371$ ở cử động thứ $40$. & \texttt{models/\allowbreak motor\_\allowbreak learning.py} & mô hình \\
16 & Cần hai quá trình vì cơ thể thay đổi với nhiều tốc độ khác nhau & Lý thuyết: ước lượng tối ưu khi nhiễu có nhiều thời gian sống khác nhau cho ra nhiều bộ lọc với hằng số thời gian khác nhau, và giải thích sự quay lại tự phát cũng như việc học lại nhanh hơn. & \cite{Kording2007} & giả thuyết \\
\bottomrule
\end{longtable}}

\section{Chương~\ref{sec:chemoutput}. Đầu ra hóa học}

Hai dây dẫn tới tim và tốc độ khác nhau của chúng, phản hồi âm trong các chuỗi tầng hormone, mã bằng tần số các đợt, bộ tạo nhịp ngày đêm và việc chỉnh nó bằng ánh sáng đều đã được đo trực tiếp; giới hạn $K<8$ là một định lý của lý thuyết điều khiển, được suy ra trong chương; các xung sinh ra bởi chính phản hồi của chuỗi tầng là giả thuyết có bằng chứng thực nghiệm mạnh ủng hộ.

{\footnotesize
\setlength{\tabcolsep}{3pt}\renewcommand{\arraystretch}{1.15}
\begin{longtable}{>{\raggedright}p{0.5cm}>{\raggedright}p{4.0cm}>{\raggedright}p{5.6cm}>{\raggedright}p{2.3cm}>{\raggedright\arraybackslash}p{1.7cm}}
\toprule
TT & Khẳng định & Thí nghiệm và điều đã đo & Nguồn & Trạng thái \\
\midrule\endfirsthead
\toprule
TT & Khẳng định & Thí nghiệm và điều đã đo & Nguồn & Trạng thái \\
\midrule\endhead
1 & Bộ tạo nhịp tim tự đập khoảng $100$ lần mỗi phút; lúc nghỉ phanh đang được đạp, và tần số thường vào khoảng $60$--$70$ & Người, cắt hoàn toàn cả hai dây dẫn tới tim: tần số riêng cao hơn tần số lúc nghỉ và giảm theo tuổi, ở người trưởng thành khoảng $100$ nhịp mỗi phút. Vậy lúc nghỉ phanh chiếm ưu thế. Tần số nghỉ $60$--$70$ là giá trị điển hình, không trích dẫn riêng. & \cite{JoseCollison1970ihr} & đã đo \\
2 & Ga \nt{NORA} và phanh \nt{ACET} là các bộ lọc thông thấp, phanh nhanh hơn~\eqref{eq:gasbrake} & Chó, kích thích điều tần dây thần kinh phế vị hoặc dây giao cảm của tim và hàm truyền tới tần số tâm nhĩ: cả hai đường đều là bộ lọc thông thấp, đường giao cảm có tần số cắt thấp hơn nhiều và thêm một độ trễ thuần khoảng $1.7\,\text{s}$. Các đặc tính phụ thuộc vào trương lực trung bình, nên công thức tuyến tính~\eqref{eq:gasbrake} chỉ là gần đúng. & \cite{Berger1989} & đã đo \\
3 & Phanh nhanh vì \nt{ACET} mở kênh kali không cần chất truyền tin thứ hai, còn \nt{NORA} thì qua một chuỗi phản ứng & Mảnh màng tế bào tâm nhĩ: kênh kali do acetylcholine mở được mở bởi tiểu đơn vị $\beta\gamma$ của protein G gắn với thụ thể, không cần chất truyền tin thứ hai bên trong tế bào. Đường \nt{NORA} qua cAMP không được trích dẫn ở đây. & \cite{Logothetis1987gbg} & đã đo \\
4 & Máu đi hết một vòng cơ thể trong khoảng một phút; \nt{ADRE} từ dòng máu tới được mọi cơ quan có thụ thể & Ước tính chung về bậc độ lớn; trong chương cũng được phát biểu như vậy, không trích dẫn nguồn. & --- & ước tính \\
5 & Chuỗi tầng hormone được bao bởi phản hồi âm: hormone cuối ức chế các tầng đầu & Tổng quan thí nghiệm: hormone vỏ thượng thận ức chế tuyến yên tiết ACTH và vùng dưới đồi tiết hormone giải phóng trong ba dải thời gian: nhanh, trung gian và chậm. & \cite{KellerWood1984fb} & đã đo \\
6 & Ba tầng giống nhau ổn định khi $K<8$~\eqref{eq:cascadestable}, tại biên chu kỳ là $2\pi\tau/\sqrt3\approx3.6\tau$ & Suy ra trong chương: ở tần số $\sqrt3/\tau$, ba tầng cho lệch pha $180^\circ$ và suy giảm $8$ lần. & suy ra trong chương & lý thuyết \\
7 & Sai số của mức là $1/(1+K)$, nên bộ điều chỉnh như vậy không giữ chính xác hơn khoảng $10\,\%$ (mệnh đề~\ref{prop:cascade}) & Hệ quả của dòng~6 với phản hồi tỉ lệ. Trục thật có phản hồi ở nhiều tầng và trong nhiều dải thời gian (dòng~5), nên giới hạn này thuộc về mô hình~\eqref{eq:cascade}, chứ không phải đã đo. & suy ra trong chương & lý thuyết \\
8 & Khi $K=4$ và $7$ mức ổn định, khi $K=12$ có các xung đều với chu kỳ $3.7$--$3.9\,\tau$ (hình~\ref{fig:cascade}) & Tính theo \texttt{models/\allowbreak chem\_\allowbreak models.py}: khi $K=12$ các chu kỳ liên tiếp là $3.9$, $3.7$, $3.8$, $3.7\,\tau$; khi $K=4$ biên độ ở cuối phép tính là $0.003$. & \texttt{models/\allowbreak chem\_\allowbreak models.py} & mô hình \\
9 & Hormone stress được tiết theo xung khoảng một giờ một lần; các xung do chính phản hồi sinh ra, không cần bộ tạo nhịp riêng & Chuột cống, truyền liên tục hormone giải phóng: ACTH và corticosterone dao động với tần số sinh lý, gần một giờ một lần; không cần đầu vào nhịp nhàng từ vùng dưới đồi để có xung. Mô hình tuyến yên--thượng thận với phản hồi có trễ cho ra các dao động như vậy. & \cite{Walker2012glc, Walker2010} & giả thuyết \\
10 & Bên nhận đọc tần số các đợt chứ không đọc mức & Khỉ không tự tiết hormone giải phóng gonadotropin: truyền liên tục không phục hồi sự tiết hormone tuyến yên, còn các đợt mỗi giờ một lần thì phục hồi; chuyển sang truyền liên tục lại dập tắt đáp ứng. Thụ thể mệt đi như bộ lọc thông cao là cách diễn giải. & \cite{Belchetz1978} & đã đo \\
11 & Các tần số đợt khác nhau tác động khác nhau lên các loại tế bào khác nhau trong cùng một tuyến & Cùng những con khỉ đó: tăng tần số đợt lên $2$--$5$ lần mỗi giờ làm giảm cả hai hormone tuyến yên; giảm xuống một lần mỗi $3$~giờ làm giảm một hormone (LH) và tăng hormone kia (FSH), nên tỉ lệ giữa chúng do tần số quyết định. & \cite{Wildt1981gnrh} & đã đo \\
12 & Nhịp ngày đêm sinh ra từ phản hồi âm nội bào với độ trễ vài giờ & Tổng quan: protein của các gen đồng hồ ức chế có trễ sự phiên mã của chính các gen đó; vòng phiên mã--dịch mã cho chu kỳ khoảng một ngày. & \cite{TakahashiJS2017} & đã đo \\
13 & Bộ tạo nhịp chính là một nhóm vài chục nghìn nơ-ron, đồng bộ phần còn lại của cơ thể & Tổng quan: nhân trên giao thoa thị giác là bộ tạo nhịp ngày đêm chính; từng nơ-ron của nó khi nuôi cấy riêng là những bộ tạo nhịp độc lập, mạng lưới đồng bộ chúng; ở chuột nhắt có khoảng $10^4$ nơ-ron mỗi bên. & \cite{Welsh2010scn} & đã đo \\
14 & Chu kỳ riêng ở người là $24.18$~giờ & Người sống trong ánh sáng mờ theo lịch tách khỏi ngày đêm: chu kỳ của nhịp melatonin, nhiệt độ và cortisol trung bình $24.18$~giờ ở cả người trẻ và người già, với độ tản mát hẹp. & \cite{Czeisler1999} & đã đo \\
15 & Ánh sáng chỉnh bộ tạo nhịp như một vòng khóa pha~\eqref{eq:pll}; cần dịch khoảng $11$~phút mỗi ngày & Người, một lần chiếu sáng mạnh $6.7$~giờ vào các thời điểm khác nhau trong ngày: đường cong đáp ứng pha loại một với biên độ $5$~giờ, làm trễ pha trước điểm cực tiểu nhiệt độ cơ thể và làm sớm pha sau đó. Phương trình~\eqref{eq:pll} là mô hình chuẩn; $11$~phút $=0.18$~giờ là phép tính số học. & \cite{Khalsa2003prc} & đã đo \\
16 & Không có ánh sáng thì không có khóa pha, và nhịp trôi theo chu kỳ riêng & Người mù hoàn toàn, không cảm nhận ánh sáng, sống trong xã hội bình thường: khoảng một nửa có nhịp melatonin và cortisol chạy theo chu kỳ riêng, không trùng với ngày đêm. & \cite{Sack1992blind} & đã đo \\
\bottomrule
\end{longtable}}

\section{Chương~\ref{sec:debugging}. Gỡ lỗi cỗ máy}

Điều điện cực nghe được, số chiều thấp của hoạt động, hiệu quả của các giao diện dùng mảng cứng, việc não và bộ giải mã học cùng nhau và các chấm sáng thị giác do kích thích đều đã được đo trong các thí nghiệm được bình duyệt; ước tính luồng dữ liệu là phép tính số học của chương; về thiết bị Neuralink cấy cho người, theo chỗ chúng tôi kiểm tra được, dữ liệu chủ yếu do chính công ty công bố.

{\footnotesize
\setlength{\tabcolsep}{3pt}\renewcommand{\arraystretch}{1.15}
\begin{longtable}{>{\raggedright}p{0.5cm}>{\raggedright}p{4.0cm}>{\raggedright}p{5.6cm}>{\raggedright}p{2.3cm}>{\raggedright\arraybackslash}p{1.7cm}}
\toprule
TT & Khẳng định & Thí nghiệm và điều đã đo & Nguồn & Trạng thái \\
\midrule\endfirsthead
\toprule
TT & Khẳng định & Thí nghiệm và điều đã đo & Nguồn & Trạng thái \\
\midrule\endhead
1 & Điện cực nghe xung của các nơ-ron trong bán kính cỡ một trăm micromét, thường từ một đến vài nơ-ron; chúng được phân biệt theo hình dạng & Chuột cống, vùng CA1, ghi đồng thời trong tế bào và ngoài tế bào: hình dạng xung ngoài tế bào lặp lại độ rộng và biên độ của xung trong tế bào. Ước tính: một tetrode có thể phân biệt $60$--$100$ nơ-ron, nhưng thực tế chỉ tách được khoảng sáu. & \cite{Henze2000extra} & đã đo \\
2 & Não có $8.6\cdot10^{10}$ nơ-ron; đầu dò nghe khoảng một phần trăm triệu & Đếm nhân tế bào trong huyền phù mô đã hòa tan: $86\cdot10^9$ nơ-ron trong não người. Tỉ lệ là phép tính số học của chương. & \cite{Azevedo2009} & đã đo \\
3 & Hoạt động của hàng trăm nơ-ron vỏ não vận động nằm gọn trong mười đến hai mươi biến chung & Tổng quan các bản ghi ở vỏ não vận động của khỉ: hoạt động của quần thể được mô tả bằng ít biến ẩn, chung cho nhiều nơ-ron. & \cite{Gallego2017} & đã đo \\
4 & Nhiễu của các nơ-ron lân cận tương quan, và một trăm nơ-ron mang gần bằng một nghìn (mệnh đề~\ref{prop:pooling}) & Vỏ não thị giác của khỉ: hệ số tương quan nhiễu của các nơ-ron lân cận khoảng $0.1$, và lợi ích từ việc lấy trung bình bão hòa ở một trăm nơ-ron. Lý thuyết về các tương quan giới hạn thông tin. & \cite{Zohary1994, MorenoBote2014} & đã đo \\
5 & Luồng thô $2\cdot10^8$~bit/s so với $8\cdot10^4$~bit/s trong xung~\eqref{eq:probedata} & Phép tính của chương theo dung lượng luồng xung~\eqref{eq:spikeentropy}. Thông số số hóa là thực: chip Neuralink là $19.3$~kHz, $10$~bit mỗi kênh. & suy ra trong chương; \cite{Rieke1997, Musk2019} & lý thuyết \\
6 & Hệ thống Neuralink đầu tiên: chip khuếch đại và số hóa, luồng thô đi qua cáp, xung do chương trình bên ngoài tách; tách trên chip, vỏ kín, vô tuyến và cấp năng lượng không dây là ở thiết bị cho người, theo công bố của công ty & Bài báo của công ty: toàn bộ luồng thô qua cáp USB-C, xung được một chương trình tách theo thời gian thực bên ngoài chip; nén trên thiết bị, cấp năng lượng và truyền không dây được nêu là kế hoạch cho thiết bị lâm sàng. Với thiết bị cấy cho người, chỉ có công bố của công ty. Việc tách xung trên chip là cần thiết là kết luận của chương từ~\eqref{eq:probedata}. & \cite{Musk2019}; báo cáo của Neuralink, không có bài báo bình duyệt & đã đo (bài báo của công ty); với người là báo cáo của công ty \\
7 & Tới $3072$ điện cực trên $96$ sợi, mỗi sợi $32$; các sợi được robot cắm, tránh mạch máu & Bài báo của công ty: $3072$ điện cực trên $96$ sợi, mỗi sợi $32$; robot cắm $6$ sợi ($192$ điện cực) mỗi phút; tới $70\,\%$ điện cực ở chuột cống được cấy mảng lâu dài nghe được xung. & \cite{Musk2019} & đã đo (bài báo của công ty) \\
8 & Ở người khoảng một nghìn kênh trên $64$ sợi; một phần sợi đã xê dịch; con trỏ khoảng $10$~bit/s & Chỉ có công bố của công ty. Tìm kiếm trong PubMed không thấy bài báo bình duyệt nào có kết quả trên người; điều này khớp với lưu ý trong chương. & báo cáo của Neuralink; không có bài báo bình duyệt & đã đo (báo cáo của công ty) \\
9 & Giao diện dùng mảng khoảng một trăm điện cực điều khiển được con trỏ & Người bị liệt, $96$ điện cực trong vỏ não vận động sơ cấp: ý định cử động tay làm thay đổi xung ba năm sau chấn thương; bộ giải mã cho ``con trỏ thần kinh'' (thư điện tử, tivi), điều khiển bàn tay giả và tay robot. & \cite{Hochberg2006} & đã đo \\
10 & Viết ``bằng bàn tay tưởng tượng'': khoảng $90$ ký tự mỗi phút, dưới $8$~bit/s & Người bị liệt bàn tay, cố gắng viết tay, bộ giải mã hồi quy: $90$ ký tự mỗi phút với độ chính xác $94.1\,\%$ theo thời gian thực, trên $99\,\%$ sau khi tự sửa lỗi. Ước tính theo bit là phép tính của chương. & \cite{Willett2021} & đã đo \\
11 & Những cố gắng nói được chuyển thành văn bản với tốc độ khoảng $60$ từ mỗi phút & Người đã mất khả năng nói rõ, các mảng điện cực trong vỏ não: $62$ từ mỗi phút; $9.1\,\%$ lỗi từ với bộ từ vựng $50$ từ, $23.8\,\%$ với bộ từ vựng $125\,000$ từ. & \cite{Willett2023} & đã đo \\
12 & Giao diện chỉ rút ra một phần nhỏ dung lượng: một nơ-ron có vài chục bit mỗi giây, một trăm nơ-ron có hàng nghìn (mệnh đề~\ref{prop:probe}) & Giới hạn trên~\eqref{eq:spikeentropy} là lý thuyết; so sánh với các dòng~8, 10 là kết luận của chương. Việc nút thắt là số chiều thấp và sự chưa hiểu mã, chứ không phải dây dẫn, chưa được kiểm tra bằng thí nghiệm trực tiếp. & \cite{Rieke1997} & lý thuyết \\
13 & Não và bộ giải mã học lẫn nhau & Khỉ điều khiển con trỏ; bộ giải mã tự điều chỉnh trong vòng kín, còn các nơ-ron đồng thời học lại: độ chính xác tăng, kỹ năng được giữ lại và ít bị ảnh hưởng bởi cử động của chính tay con vật. Lời khuyên tách tốc độ học là quy tắc kỹ thuật, trong chương không có thí nghiệm riêng. & \cite{Orsborn2014coad} & đã đo \\
14 & Dòng điện qua điện cực kích thích các nơ-ron và dây dẫn xung quanh bất kể loại & Vỏ não chuột nhắt và mèo, ghi canxi hai photon: ngay cả dòng yếu cũng kích thích các nơ-ron thưa thớt, rải rác trên khoảng cách tới vài milimét, thông qua kích thích trực tiếp các sợi trục trong một thể tích đường kính vài chục micromét. Tức là kích thích không chọn lọc, nhưng cũng không liền khối. & \cite{Histed2009stim} & đã đo \\
15 & Kích thích vỏ não thị giác làm bật các chấm sáng; tới $16$ chấm cùng lúc cho phép nhận ra một số chữ cái và đường biên đồ vật & Người mù, $96$ điện cực trong vỏ não thị giác trong $6$ tháng: ngưỡng của một điện cực $66.8\pm36.5$~\textmu A; kích thích đồng thời theo nhóm (tới $16$ điện cực, giới hạn của bộ kích thích) hạ ngưỡng và cho các mẫu phân biệt được; một số chữ cái và đường biên đồ vật được nhận ra. & \cite{Fernandez2021} & đã đo \\
16 & Hướng thứ hai của Neuralink là thiết bị đưa tín hiệu vào vỏ não thị giác & Chỉ có thông báo của công ty về việc phát triển; không tìm thấy dữ liệu bình duyệt về hoạt động của nó. & báo cáo của Neuralink & giả thuyết \\
17 & Nồng độ chất dẫn truyền thần kinh quảng bá có thể đo bằng cảm biến hóa học trong mô & Lát não chuột cống, đo volt-ampe vòng quét nhanh trên sợi carbon: đo được sự giải phóng và tái hấp thu serotonin; chất này thoát ra ngoài khớp thần kinh. & \cite{Bunin1998} & đã đo \\
\bottomrule
\end{longtable}}

\section{Chương~\ref{sec:eetypes}. Nơ-ron như linh kiện}

Các chế độ hoạt động của từng tế bào đã được đo trực tiếp, bằng dòng điện qua điện cực và ghi điện áp; toán học của bộ tích phân và của việc phân lớp là lý thuyết kinh điển, còn việc bộ não dùng sự tái cấu hình chế độ để tính toán vẫn là giả thuyết.

{\footnotesize
\setlength{\tabcolsep}{3pt}\renewcommand{\arraystretch}{1.15}
\begin{longtable}{>{\raggedright}p{0.5cm}>{\raggedright}p{4.0cm}>{\raggedright}p{5.6cm}>{\raggedright}p{2.3cm}>{\raggedright\arraybackslash}p{1.7cm}}
\toprule
TT & Khẳng định & Thí nghiệm và điều đã đo & Nguồn & Trạng thái \\
\midrule\endfirsthead
\toprule
TT & Khẳng định & Thí nghiệm và điều đã đo & Nguồn & Trạng thái \\
\midrule\endhead
1 & Bộ cộng hưởng: dòng chậm chống lại thay đổi điện áp biến nơ-ron thành bộ lọc thông dải với cực đại ở vài đến vài chục hertz (mục~\ref{sec:resonator}) & Trong lát não, người ta đưa vào nơ-ron \nt{GLUT} dòng hình sin có tần số tăng dần và đo trở kháng $|Z(f)|$: cực đại ở $2$--$5$~Hz tại $33^\circ$C (khoảng $7$~Hz tại $38^\circ$C); cộng hưởng do các dòng kali và dòng cation chậm tạo ra và được dòng natri dai dẳng khuếch đại. Tổng quan cho thấy cùng thủ pháp ở nhiều lớp tế bào & \cite{HuVervaekeStorm2002,HutcheonYarom2000} & đã đo \\
2 & Dòng chậm hoạt động như cuộn cảm, cùng với điện dung màng tạo thành mạch dao động & Đáp ứng dưới ngưỡng của sợi trục với dòng điện đã được đo và mô tả bằng các phương trình độ dẫn tuyến tính hóa; sơ đồ tương đương có một cuộn cảm ``hiện tượng luận'' với giá trị phụ thuộc điện áp & \cite{Mauro1970} & lý thuyết \\
3 & Hằng số thời gian của nhóm có phản hồi $\tau_{\text{hd}}=\tau/(1-w)$~\eqref{eq:perfectint}; với $\tau=0.1$~s và $w=0.995$ là $20$~s & Suy ra trong chương từ phương trình tuyến tính của nhóm; được đề xuất như cơ chế giữ vị trí mắt trong một công trình lý thuyết & suy ra trong chương; \cite{Seung1996} & lý thuyết \\
4 & Bộ tích phân vị trí mắt giữ đầu vào nhờ mạng, chứ không nhờ tính chất của một tế bào riêng lẻ & Ghi nội bào ở cá, trong nhân giữ vị trí mắt: sau một lần xoay nhanh, tần số của nơ-ron thay đổi theo bậc và được giữ; một dòng ngắn vào một tế bào chỉ gây dịch chuyển nhất thời, còn các bậc điện áp vẫn được duy trì cả dưới ngưỡng, nhờ các đầu vào khớp thần kinh & \cite{Aksay2001} & đã đo \\
5 & Bộ tích phân không rò cần được học liên tục để tự điều chỉnh; khi lệch chỉnh, nó hoặc quên, hoặc đi vào bão hòa & Cá được huấn luyện bằng phản hồi thị giác tỉ lệ với $\pm$ vị trí mắt: việc giữ vị trí trở nên không ổn định hoặc bị rò, hằng số thời gian của nơ-ron giảm hơn một bậc, xuống $<1$~s; phản hồi thị giác bình thường dần dần khôi phục độ ổn định & \cite{Major2004} & đã đo \\
6 & Nơ-ron hai trạng thái bền: đầu vào ngắn bật một bình nguyên kéo dài, ức chế tắt nó; tính chất này được bật bởi serotonin (mục~\ref{sec:bistable}) & Ghi nội bào các nơ-ron \nt{ACET} điều khiển cơ ở mèo: một kích thích khớp thần kinh ngắn đưa nơ-ron vào hoạt động kéo dài, một ức chế ngắn đặt lại nó; ở mèo bị cắt tủy sống, trạng thái này xuất hiện sau khi đưa vào tiền chất của serotonin & \cite{Hounsgaard1988} & đã đo \\
7 & Hai lớp: bộ tích phân (tần số tăng từ không) và bộ cộng hưởng (tại ngưỡng lập tức có tần số hữu hạn) & Các lớp được suy ra từ lý thuyết rẽ nhánh. Thí nghiệm: trong lát vỏ não, nơ-ron \nt{GLUT} bắt đầu phát xung với tần số nhỏ tùy ý (loại~1), nơ-ron \nt{GABA} nhanh thì nhảy vọt từ một tần số hữu hạn (loại~2) & \cite{Izhikevich2007,Tateno2004} & đã đo \\
8 & Một nơ-ron là bộ tích phân hoặc bộ phát hiện trùng khớp: ức chế rút ngắn cửa sổ cộng & Trong lát hồi hải mã, cửa sổ trong đó các đầu vào kích thích cộng lại tới ngưỡng ngắn hơn $2$~ms; nó bị rút ngắn bởi ức chế thuận chiều đến ngay sau kích thích; trong đuôi gai, nơi ức chế yếu hơn, cửa sổ rộng hơn & \cite{Pouille2001} & đã đo \\
9 & Đường dây trễ từ sợi trục (bảng~\ref{tab:eetypes}) & Ở cú lợn, đã đo độ trễ trong các sợi trục đi tới bộ phát hiện trùng khớp: chiều dài sợi trục khác nhau cho độ trễ khác nhau, và các bộ phát hiện được điều chỉnh theo hiệu thời gian âm thanh đến & \cite{Carr1990} & đã đo \\
10 & Bộ tạo nhịp khi thức và tế bào im lặng khi ngủ & Nơ-ron \nt{SERO} phát xung gần như đều đặn ban ngày, chậm lại trong giấc ngủ sóng chậm và gần như im lặng trong giấc ngủ REM (chi tiết ở chương~\ref{sec:sleep}) & \cite{JacobsAzmitia1992,McGintyHarper1976} & đã đo \\
11 & Linh kiện do các kênh ion và các kênh quảng bá điều chỉnh chúng quyết định (mệnh đề~\ref{prop:eetypes}) & Đã chứng minh trên các mạng nhỏ: các chất điều biến thay đổi tính chất riêng của nơ-ron và độ mạnh của khớp thần kinh, nên cùng một mạch cho ra các nhịp khác nhau & \cite{Marder2012} & đã đo \\
12 & Bộ não dùng sự tái cấu hình chế độ để tính toán (mệnh đề~\ref{prop:eetypes}) & Chưa có thí nghiệm trực tiếp nào trong đó việc tái cấu hình chế độ của tế bào là cần thiết để giải một nhiệm vụ; chỉ có các thí nghiệm trong đó chất điều biến thay đổi đầu ra của mạch & \cite{Marder2012} & giả thuyết \\
\bottomrule
\end{longtable}}

\section{Chương~\ref{sec:dendrites}. Số đuôi gai}

Cấu tạo của các lớp và số đầu vào đã được đo bằng giải phẫu và ghi điện; ước tính~\eqref{eq:dendriteload} là vật lý đơn giản của tụ điện, còn lời giải thích về tính toán cho số đầu vào dựa trên lý thuyết và mô hình.

{\footnotesize
\setlength{\tabcolsep}{3pt}\renewcommand{\arraystretch}{1.15}
\begin{longtable}{>{\raggedright}p{0.5cm}>{\raggedright}p{4.0cm}>{\raggedright}p{5.6cm}>{\raggedright}p{2.3cm}>{\raggedright\arraybackslash}p{1.7cm}}
\toprule
TT & Khẳng định & Thí nghiệm và điều đã đo & Nguồn & Trạng thái \\
\midrule\endfirsthead
\toprule
TT & Khẳng định & Thí nghiệm và điều đã đo & Nguồn & Trạng thái \\
\midrule\endhead
1 & Điện dung riêng của màng $c_m\approx1$~\textmu F/cm$^2$ & Màng được lấy từ thân nơ-ron vào pipet, đặt một bậc điện áp và từ sự suy giảm của dòng điện dung tìm ra điện dung toàn phần, rồi chia cho diện tích: $0.9$~\textmu F/cm$^2$ ở ba lớp nơ-ron & \cite{Gentet2000} & đã đo \\
2 & Thời gian nạp $t\approx c_mA\Delta V/I$~\eqref{eq:dendriteload}: nơ-ron lớn cần hàng chục đầu vào đồng thời, nơ-ron nhỏ chỉ cần một khớp lớn & Ước tính theo định luật nạp tụ; các con số ($20$ và $300$~pF, $0.1$ và vài nA) là bậc độ lớn & suy ra trong chương & lý thuyết \\
3 & Một khớp khổng lồ cho dòng vài nanoampe và đưa ngay nơ-ron nhỏ tới ngưỡng & Ghi đồng thời đầu tận cùng và tế bào nhận trong lát não: dòng của một khớp $2$--$13$~nA, mỗi xung đầu vào gây một đáp ứng trên ngưỡng, độ trễ khoảng $0.4$~ms ở $36^\circ$C; tế bào nhận xuất \nt{GLYC} & \cite{Borst1995,BorstSoria2012} & đã đo \\
4 & Không đuôi gai: ở nơ-ron cảm giác xúc giác và đau, xung đi từ cảm biến vào tủy sống, bỏ qua thân tế bào & Cấu tạo (sợi trục chia đôi gần thân) đã được xác lập bằng giải phẫu. Việc dẫn truyền không cần thân kích thích được đã được chỉ ra trên một mô hình đã kiểm chứng của nơ-ron này: không có thân kích thích được, xung vẫn đi qua chỗ rẽ nhánh tin cậy như cũ & \cite{AmirDevor2003} & lý thuyết \\
5 & Một đuôi gai: nơ-ron khứu giác mang một loại thụ thể và truyền một đường đầu vào & Tổng quan các thí nghiệm với gen thụ thể: mỗi nơ-ron khứu giác biểu hiện một gen thụ thể từ một họ lớn & \cite{Mombaerts2004} & đã đo \\
6 & Hai đuôi gai: ở bộ phát hiện hướng âm thanh, đầu vào từ tai trái và tai phải bám vào các đuôi gai khác nhau & Giải phẫu và các bản ghi ở động vật có vú, được tổng hợp trong một tổng quan: đầu vào từ hai tai được tách ra trên hai đuôi gai đối diện & \cite{Grothe2010} & đã đo \\
7 & Tách đầu vào ra hai đuôi gai làm cổng AND chặt hơn & Được chỉ ra trên mô hình: sự bão hòa~\eqref{eq:shunt} trên một đuôi gai không cho các đầu vào từ một phía tới được ngưỡng, và khả năng phát hiện trùng khớp được cải thiện. Chưa có thí nghiệm trực tiếp nào thay đổi cách bố trí đầu vào & \cite{AgmonSnir1998} & lý thuyết \\
8 & Khoảng $7\cdot10^{10}$ nơ-ron \nt{GLUT} nhỏ của mạch học vận động, mỗi nơ-ron $\sim4$ đầu vào & Đếm tế bào bằng phân đoạn đẳng hướng: phần não này ở người có khoảng $69$ tỉ nơ-ron. Ghi một tế bào như vậy ở chuột cống sống: khi chạm, các chùm dòng rời rạc từ vài đầu vào đến, và tế bào phát xung khi chúng cộng lại & \cite{Azevedo2009,Chadderton2004} & đã đo \\
9 & Phần tử ngưỡng có $n$ đầu vào phân tách không quá $P_{\max}\approx2n$ mẫu ngẫu nhiên~\eqref{eq:cover}; với nơ-ron đầu ra của mạch học vận động, giới hạn là $\sim2\cdot10^5$, còn ước tính thực tế là $\sim4\cdot10^4$ phép ánh xạ & Giới hạn là kết quả kinh điển của hình học tổ hợp. Ước tính thực tế là lý thuyết dung lượng với trọng số chỉ dương và có dự trữ cho nhiễu, áp dụng cho số đầu vào đã đo của nơ-ron này & \cite{Cover1965,Brunel2004} & lý thuyết \\
10 & Bộ trộn với ít đầu vào cần để mở rộng đầu vào; con số ``bốn'' gần tối ưu & Lý thuyết: số chiều của biểu diễn do lớp bộ trộn tạo ra đạt cực đại xấp xỉ tại số đầu vào quan sát được trên giải phẫu; giả thuyết mở rộng ban đầu có trong các công trình kinh điển & \cite{Marr1969,Albus1971,LitwinKumar2017} & giả thuyết \\
11 & Cây lớn: $10^4$--$10^5$ đầu vào & Đếm khớp thần kinh trên ảnh hiển vi điện tử: $\approx1.7\cdot10^5$ khớp trên một nơ-ron \nt{GABA} đầu ra của mạch học vận động ở chuột cống; khoảng $30\,000$ đầu vào kích thích và $1\,700$ đầu vào ức chế trên một nơ-ron \nt{GLUT} của hồi hải mã & \cite{NapperHarvey1988,Megias2001} & đã đo \\
12 & Các nhánh của cây lớn là những mạch con riêng với ngưỡng riêng; nơ-ron là một mạng nhiều lớp nhỏ & Kích thích điểm tại hai vị trí trên đuôi gai mảnh: các đầu vào trên cùng một nhánh cộng theo đường sigmoid, trên các nhánh khác nhau thì cộng tuyến tính. Trong đuôi gai của nơ-ron vỏ não người đã tìm thấy các xung canxi phân cấp cho phép một tế bào giải một bài toán không tách được tuyến tính. Mô hình: để tái tạo một nơ-ron cần một mạng sâu & \cite{Polsky2004,Gidon2020,Beniaguev2021} & đã đo \\
13 & Đuôi gai kiêm đầu ra: mỗi nhánh của nơ-ron \nt{GABA} võng mạc là một bộ phát hiện hướng riêng & Ghi canxi hai photon trong từng nhánh: tín hiệu trong nhánh chọn lọc với chuyển động từ thân tế bào ra đầu mút, còn điện áp của thân thì không & \cite{Euler2002} & đã đo \\
14 & Số đầu vào đi theo nhiệm vụ (mệnh đề~\ref{prop:dendrites}) & Cấu tạo các lớp đã được đo (dòng 3--13); việc số đầu vào được chọn vì tốc độ hay vì phân lớp chỉ mới được chỉ ra bằng lý thuyết và mô hình cho một số lớp & \cite{LitwinKumar2017,AgmonSnir1998} & giả thuyết \\
\bottomrule
\end{longtable}}

\section{Chương~\ref{sec:sleep}. Giấc ngủ}

Các chế độ ngủ, việc phát lại và sự thu nhỏ khớp thần kinh đã được đo bằng ghi nơ-ron, vi thẩm tách và hiển vi điện tử; lịch ngủ và dao động chậm được mô tả bằng các mô hình của sách, còn mục đích của giấc ngủ phần lớn là giả thuyết.

{\footnotesize
\setlength{\tabcolsep}{3pt}\renewcommand{\arraystretch}{1.15}
\begin{longtable}{>{\raggedright}p{0.5cm}>{\raggedright}p{4.0cm}>{\raggedright}p{5.6cm}>{\raggedright}p{2.3cm}>{\raggedright\arraybackslash}p{1.7cm}}
\toprule
TT & Khẳng định & Thí nghiệm và điều đã đo & Nguồn & Trạng thái \\
\midrule\endfirsthead
\toprule
TT & Khẳng định & Thí nghiệm và điều đã đo & Nguồn & Trạng thái \\
\midrule\endhead
1 & Khi ngủ, nơ-ron vỏ não không im lặng; cái thay đổi là chế độ làm việc & Ghi dài hạn nơ-ron vỏ não ở chuột cống: trong giấc ngủ sóng chậm, tần số vẫn khác không, các giai đoạn làm việc xen kẽ với các giai đoạn im lặng; sau thời gian thức dài, tần số cao hơn ở mọi trạng thái, sau khi ngủ thì thấp hơn & \cite{Vyazovskiy2009} & đã đo \\
2 & \nt{ACET} cao ban ngày và trong giấc ngủ REM, thấp trong giấc ngủ sóng chậm (bảng~\ref{tab:sleepmodes}) & Vi thẩm tách trong vỏ não ở mèo tự do vận động: khi thức, lượng acetylcholine giải phóng cao hơn $100$--$175\%$ so với khi ngủ sóng chậm, còn trong giấc ngủ REM thì xấp xỉ mức thức yên tĩnh & \cite{Marrosu1995} & đã đo \\
3 & \nt{NORA} và \nt{SERO} lắng xuống trong giấc ngủ sóng chậm và gần như im lặng trong giấc ngủ REM & Ghi từng nơ-ron \nt{NORA} ở chuột cống và nơ-ron \nt{SERO} ở mèo theo các giai đoạn ngủ: tần số cực đại khi thức, thấp hơn trong giấc ngủ sóng chậm và gần bằng không trong giấc ngủ REM & \cite{AstonJonesBloom1981,McGintyHarper1976} & đã đo \\
4 & Trong giấc ngủ REM, cơ bị tắt bằng ức chế qua \nt{GABA} và \nt{GLYC} & Ở chuột cống, người ta đo hoạt động của các nơ-ron \nt{ACET} của cơ nhai và đưa chất chẹn trực tiếp tới chúng: tình trạng liệt chỉ được giải trừ nếu chẹn cả hai loại thụ thể \nt{GABA} lẫn thụ thể \nt{GLYC} & \cite{BrooksPeever2012} & đã đo \\
5 & Áp lực ngủ $S$ là một tụ điện với hai ngưỡng~\eqref{eq:sleeppressure}, $\tau_r=18.2$~giờ, $\tau_d=4.2$~giờ & Mô hình hai quá trình; các hằng số thời gian thu được từ cách công suất sóng chậm của EEG tăng khi thức và giảm khi ngủ, còn dạng của các ngưỡng từ thời điểm tự thức dậy & \cite{Borbely1982,Daan1984} & lý thuyết \\
6 & Cỗ máy tự đi tới $16.1$~giờ thức và $8.0$~giờ ngủ (hình~\ref{fig:twoprocess}) & Tính theo~\eqref{eq:sleeppressure} với ngưỡng dao động, tập lệnh \texttt{models/sleep\_models.py}: $16.1$~giờ thức và $7.9$--$8.0$~giờ ngủ & --- & mô hình \\
7 & Sau $40$~giờ không ngủ $S=0.90$, nhưng giấc ngủ chỉ kéo dài $8.8$~giờ: món nợ được trả bằng độ sâu, không bằng độ dài & Mô hình: \texttt{models/sleep\_models.py} cho $S=0.90$ và $8.8$~giờ. Thí nghiệm: sau $40.5$~giờ không ngủ, ở tám người, trong đêm ngủ bù đầu tiên công suất sóng chậm của EEG cao hơn đáng kể so với bình thường & \cite{Borbely1981} & mô hình \\
8 & Về mặt vật lý, $S$ là adenosine & Vi thẩm tách ở mèo: adenosine ngoại bào ở một trong các vùng điều khiển sự thức tăng lên trong thời gian thức, tăng tiếp khi bị thiếu ngủ và giảm khi ngủ. Việc $S$ chỉ là adenosine chưa được chứng minh & \cite{PorkkaHeiskanen1997,Dunwiddie2001} & giả thuyết \\
9 & Trong giấc ngủ sóng chậm, vỏ não dao động: làm việc vài trăm ms rồi im lặng, chưa tới một lần mỗi giây ($<1$~Hz, khi gây mê thường $0.3$--$0.4$~Hz) & Ghi nội bào ở mèo được gây mê: khử cực $0.8$--$1.5$~s và tăng phân cực kéo dài, nhịp $<1$~Hz, thường gặp nhất $0.3$--$0.4$~Hz; cùng sự xen kẽ làm việc và im lặng này thấy được trong giấc ngủ tự nhiên (dòng~1) & \cite{Steriade1993,Vyazovskiy2009} & đã đo \\
10 & Dao động là một nhóm có phản hồi và mệt mỏi~\eqref{eq:updown}; với $b=5$ chu kỳ $1.1$~s (giá trị mô hình) và làm việc khoảng $35\%$ thời gian, với $b=1$ làm việc đều đặn (hình~\ref{fig:updown}) & Tính toán, tập lệnh \texttt{models/sleep\_models.py}: chu kỳ $1.11$~s, tỉ lệ thời gian làm việc $0.35$ (trung bình trên một số nguyên chu kỳ); với $b=1$ tỉ lệ làm việc là $1.0$. Chu kỳ của mô hình ngắn hơn chu kỳ đo được (dòng~9) & --- & mô hình \\
11 & \nt{ACET} tắt dao động và chuyển vỏ não sang làm việc đều đặn & Kích thích nhân nơ-ron \nt{ACET} ở mèo chặn nhịp chậm ở $79\%$ nơ-ron vỏ não và chuyển chúng sang làm việc đều đặn, chủ yếu do các đợt tăng phân cực kéo dài biến mất. Acetylcholine ức chế các dòng kali chậm gây ra mệt mỏi & \cite{SteriadeAmzica1993,McCormick1992} & đã đo \\
12 & Các đợt bùng nhịp $7$--$15$~Hz do vòng \nt{GLUT}--\nt{GABA} giữa vỏ não và nút chuyển tiếp sinh ra & Trong lát nút chuyển tiếp thị giác của chồn sương, các đợt bùng được tạo ra bởi các chùm xung đáp trả của tế bào chuyển tiếp trước sự ức chế từ một nhân \nt{GABA} lân cận, mà chính chúng kích thích & \cite{vonKrosigk1993,SteriadeMcCormickSejnowski1993} & đã đo \\
13 & Việc phát lại ngày đã qua diễn ra có tua nhanh: ở hồi hải mã nhanh hơn khoảng $20$ lần, ở vỏ não $6$--$7$ lần & Trong giấc ngủ sóng chậm, các chuỗi lặp lại của nơ-ron hồi hải mã bị nén theo thời gian. Sau khi chạy trên đường chạy, các chuỗi nơ-ron vị trí được lặp lại theo cùng thứ tự trong các đợt ngắn $\sim100$~ms, nén khoảng $20$ lần; ở vỏ não, mức nén là $6$--$7$ lần & \cite{Nadasdy1999,LeeWilson2002,Euston2007} & đã đo \\
14 & Tăng cường sự phối hợp giữa phát lại và vỏ não cải thiện trí nhớ (quan hệ nhân quả) & Ở chuột cống sau khi học, kích thích gắn với các lần phát lại làm tăng liên kết của chúng với sóng chậm và các đợt bùng nhịp ở vỏ não: hôm sau trí nhớ ở mức cao, còn ở nhóm đối chứng thì ở mức ngẫu nhiên & \cite{Maingret2016} & đã đo \\
15 & Mục đích của phát lại là học trộn lẫn cho bộ nhớ chậm & Lý thuyết hai hệ thống học và các tính toán trên mạng nhân tạo: học tuần tự làm hỏng cái cũ, học trộn lẫn thì không. Chưa có thí nghiệm trực tiếp trên não cho thấy phát lại cần chính là cho việc này & \cite{McClelland1995} & giả thuyết \\
16 & Sau khi ngủ, khớp thần kinh thu nhỏ tỉ lệ với kích thước, các khớp lớn hầu như không đổi & Hiển vi điện tử ba chiều $6920$ khớp thần kinh vỏ não chuột nhắt: diện tích tiếp xúc sau khi ngủ nhỏ hơn $\sim18\%$ so với sau khi thức, mức giảm tỉ lệ với kích thước, các khớp lớn và bền không thay đổi & \cite{deVivo2017} & đã đo \\
17 & Nhiệm vụ chính của giấc ngủ là chuẩn hóa lại trọng số & Giả thuyết cân bằng nội môi khớp thần kinh; các dòng 1 và 16 phù hợp với nó, nhưng không chứng minh đó là chức năng chính & \cite{TononiCirelli2014} & giả thuyết \\
18 & Giấc ngủ REM là việc thử nghiệm mô hình thế giới trên bàn thử & Chế độ (bảng~\ref{tab:sleepmodes}) đã được đo; việc mạng chạy thử mô hình thế giới là một giả định lý thuyết, chưa có thí nghiệm & \cite{Hobson2009} & giả thuyết \\
19 & Rửa sạch não khi ngủ: câu hỏi còn mở & Một thí nghiệm: khi ngủ, khoảng gian bào rộng hơn $60\%$ và việc dọn dẹp nhanh hơn. Một thí nghiệm khác, với cách đo khác: khi ngủ và khi gây mê, việc dọn dẹp chậm hơn rõ rệt. Các kết quả mâu thuẫn nhau & \cite{Xie2013,Miao2024} & giả thuyết \\
20 & Giấc ngủ cục bộ: các vùng vỏ não của con vật đang thức thỉnh thoảng rơi vào im lặng trong chốc lát & Ở chuột cống sau thời gian thức dài, nơ-ron của một vùng vỏ não riêng lẻ im lặng trong chốc lát như trong giấc ngủ sóng chậm, kèm một sóng chậm trong EEG cục bộ; vào những lúc đó chuột mắc lỗi trong nhiệm vụ nhiều hơn & \cite{Vyazovskiy2011} & đã đo \\
\bottomrule
\end{longtable}}

\section{Chương~\ref{sec:livingmachine}. Cỗ máy từ nơ-ron sống}

Hành vi của mẫu nuôi cấy trên mảng điện cực đã được đo trong các thí nghiệm trực tiếp có ghi và kích thích; cơ chế của các đợt bùng phát dựa trên mô hình cạn kiệt khớp thần kinh và thí nghiệm trên vi mẫu nuôi cấy, còn các khẳng định về dopamine trong đĩa nuôi dựa trên những thí nghiệm đơn lẻ, mà một phần trong số đó chính tác giả cũng chỉ coi là minh họa.

{\footnotesize
\setlength{\tabcolsep}{3pt}\renewcommand{\arraystretch}{1.15}
\begin{longtable}{>{\raggedright}p{0.5cm}>{\raggedright}p{4.0cm}>{\raggedright}p{5.6cm}>{\raggedright}p{2.3cm}>{\raggedright\arraybackslash}p{1.7cm}}
\toprule
TT & Khẳng định & Thí nghiệm và điều đã đo & Nguồn & Trạng thái \\
\midrule\endfirsthead
\toprule
TT & Khẳng định & Thí nghiệm và điều đã đo & Nguồn & Trạng thái \\
\midrule\endhead
1 & Mẫu nuôi cấy trên chip: phần lớn là \nt{GLUT}, một tỷ lệ nhỏ hơn \nt{GABA} tùy mẫu; trong mẫu nuôi cấy hồi hải mã khoảng $6\%$ & Mạng hàng nghìn nơ-ron trên mảng điện cực là thí nghiệm tiêu chuẩn (dòng~3). Tỷ lệ nơ-ron \nt{GABA} trong mẫu nuôi cấy hồi hải mã được đếm bằng nhuộm GABA: khoảng $6\%$ số nơ-ron. Với mẫu nuôi cấy vỏ não, chúng tôi không đưa ra con số đã được kiểm chứng & \cite{Benson1994} & đã đo \\
2 & Organoid: mô thần kinh ba chiều cỡ khoảng một milimét từ tế bào gốc người & Từ tế bào gốc đã nuôi được organoid ba chiều có nhiều vùng não, trong đó có vỏ não với các lớp tế bào tiền thân và nơ-ron trưởng thành & \cite{Lancaster2013} & đã đo \\
3 & Không có đầu vào, mẫu nuôi cấy bùng lên thành từng đợt, cách nhau từ một giây đến vài phút tùy mẫu & $58$ mẫu nuôi cấy vỏ não mật độ từ $3\,000$ đến $50\,000$ nơ-ron được ghi trong năm tuần ($963$ bản ghi nửa giờ): sau hai tuần, ở gần như mọi mẫu, các đợt bùng phát toàn mạng chiếm ưu thế; khoảng cách giữa các đợt từ $1$ đến $300$~s và phụ thuộc mạnh vào mẫu & \cite{Wagenaar2006} & đã đo \\
4 & Đợt bùng phát kết thúc do cạn chất dẫn truyền, khoảng cách $T_{\text{đợt}}\approx\tau_{\text{ph}}\ln\frac{1}{1-x^*}$~\eqref{eq:burstinterval}; $2$--$4$~s là ước tính mô hình với $\tau_{\text{ph}}=0.8$~s, sự hồi phục đo được thì chậm hơn & Công thức được suy ra trong chương. Mô hình: mạng có khớp thần kinh bị cạn kiệt tự tạo ra các đợt bùng phát toàn mạng. Thí nghiệm trên vi mẫu nuôi cấy $4$--$30$ nơ-ron: độ dài đợt bùng phát phụ thuộc vào thời gian từ đợt trước, làm cạn túi chất dẫn truyền thì xóa các đợt; kết luận của tác giả là đợt bùng phát bị giới hạn bởi kho chất dẫn truyền. Kho hồi phục ở đó mất vài chục giây, và với cùng công thức thì cho khoảng cách vài chục giây đến vài phút & suy ra trong chương; \cite{Tsodyks2000,CohenSegal2011,Tsodyks1997} & lý thuyết \\
5 & ``Người thầy im lặng'': kích thích bị tắt khi đáp ứng mong muốn xuất hiện, và đáp ứng được củng cố & Mẫu nuôi cấy được kích thích ở tần số $0.3$--$1$~Hz cho đến khi đáp ứng đã chọn xuất hiện sau kích thích $50\pm10$~ms, rồi tắt kích thích trong $5$~phút; sau vài chu kỳ, đáp ứng mong muốn được gây ra ngay; đã vẽ các đường cong học & \cite{Shahaf2001} & đã đo \\
6 & Mẫu nuôi cấy chơi quần vợt; mức tăng có ý nghĩa của các chuỗi trong một phiên chỉ có khi phản hồi đầy đủ & Chi tiết ở chương~\ref{sec:dishbrain} và phụ lục của nó: với im lặng thay cho kích thích, cuối phiên mẫu nuôi cấy cũng chơi tốt hơn so với không có phản hồi, nhưng hiệu ứng nhỏ hơn; việc học từ phiên này sang phiên khác không bền vững & \cite{Kagan2022} & đã đo \\
7 & Việc mẫu nuôi cấy học để dự đoán đầu vào của mình là giả thuyết; những lời lớn về sự ``có tri giác'' bị phản bác & Nguyên lý năng lượng tự do là một đề xuất lý thuyết; chưa có thí nghiệm trực tiếp nào phân biệt nó với các cách giải thích đơn giản hơn. Ba mươi nhà nghiên cứu trong một bài phản hồi chỉ ra rằng sự thay đổi hành vi trong vòng lặp không chứng tỏ trí tuệ hay cảm giác & \cite{Friston2010,Balci2023} & giả thuyết \\
8 & Mẫu nuôi cấy điều khiển cơ thể ảo tới mục tiêu khi chọn đúng mẫu kích thích & Chương trình tự chọn các mẫu kích thích huấn luyện theo kết quả hiện tại; trong vài chục phút, mạng đạt được các trạng thái mong muốn, tính dẻo duy trì $80$~phút sau $2$~giờ huấn luyện. Cùng những kích thích ấy, đưa vào không gắn với hành vi, thì không có tác dụng & \cite{Bakkum2008} & đã đo \\
9 & Hồ chứa: chỉ lớp đọc ra tuyến tính $y=W_{\text{ra}}\mathbf r$ được huấn luyện~\eqref{eq:reservoir} & Lý thuyết tính toán hồ chứa: một hệ phi tuyến bất kỳ có trí nhớ tắt dần cộng với một đầu ra tuyến tính được huấn luyện & \cite{Maass2002,JaegerHaas2004} & lý thuyết \\
10 & Organoid làm hồ chứa phân biệt người nói với độ chính xác khoảng $78\%$ (theo tác giả); lớp đọc ra học theo sai số, còn liên kết của organoid thay đổi không cần thầy & Âm thanh tiếng nói của vài người được đưa vào organoid dưới dạng các mẫu dòng điện trên mảng điện cực; lớp đọc ra đã huấn luyện phân biệt người nói với độ chính xác khoảng $78\%$, theo báo cáo của tác giả. Họ cũng báo cáo rằng trong lúc huấn luyện, độ liên kết chức năng của organoid thay đổi, và gọi đó là học không giám sát bên trong organoid & \cite{Cai2023} & đã đo \\
11 & Một triệu nơ-ron tốn cho xung khoảng $10^{-4}$~W~\eqref{eq:dishpower} & Ước tính: giá một xung $10^{-10}$~J từ ngân sách năng lượng của vỏ não~\eqref{eq:energy}, nhân với số xung; chưa có phép đo trực tiếp công suất của mẫu nuôi cấy & suy ra trong chương; \cite{Attwell2001} & lý thuyết \\
12 & Dopamine theo chớp sáng: $1$~mM dopamine trong lồng, $240$ lần lặp, số xung được gây ra tăng lên & Trên nền tảng organoid từ xa, dopamine trong lồng nhạy sáng ($1$~mM) được giải phóng bằng tia cực tím ($365$~nm, $0.8$~s) khi kích thích gây ra nhiều xung hơn; qua $240$ bước (khoảng $5$~giờ), số xung nhìn chung tăng. Tác giả viết rằng chỉ từ một thí nghiệm thì không thể xác lập mối liên hệ với dopamine; không có đối chứng & \cite{Jordan2024} & giả thuyết \\
13 & Dopamine từ tế bào sống thay đổi trọng số của transistor hữu cơ lâu dài và thuận nghịch & Tế bào tiết dopamine được nuôi phía trên một phần tử neuromorphic hữu cơ; dopamine bị oxy hóa ở điện cực làm thay đổi độ dẫn của kênh; đã cho thấy thay đổi trọng số lâu dài và sự đưa về lại & \cite{Keene2020} & đã đo \\
14 & Có những organoid với nơ-ron \nt{DOPA} hoạt động; dopamine của chúng chưa được dùng làm thầy dạy & Trong organoid từ tế bào gốc người đã tìm thấy nơ-ron \nt{DOPA} trưởng thành hoạt động điện và sự sản sinh dopamine. Chúng tôi không tìm thấy trong tài liệu thí nghiệm nào dùng dopamine này dạy một mạng khác & \cite{Jo2016} & đã đo \\
15 & Organoid giữ thăng bằng cây sào: kích thích dạy do chương trình chọn tốt hơn ngẫu nhiên, không được củng cố, qua khớp thần kinh \nt{GLUT} & Organoid vỏ não chuột trong vòng kín với bài toán ``cây sào trên xe đẩy'': với phần lớn organoid, tín hiệu do học tăng cường chọn cho kết quả tốt hơn tín hiệu ngẫu nhiên hoặc không có tín hiệu; sau $45$~phút nghỉ, sự cải thiện không còn; chặn thụ thể AMPA và NMDA làm mất nó & \cite{Robbins2026} & đã đo \\
16 & Không có thanh ghi điều khiển, mạng trên bàn thử không tự quyết định được thế nào là thành công (mệnh đề~\ref{prop:livingmachine}) & Trong mọi thí nghiệm ở các dòng 5--15, tín hiệu dạy do người làm thí nghiệm hoặc chương trình đặt ra; chưa có thí nghiệm nào mà mẫu nuôi cấy tự hình thành tiêu chí thành công --- đây là kết luận từ sự vắng mặt, không phải thí nghiệm & --- & giả thuyết \\
\bottomrule
\end{longtable}}

\section{Chương~\ref{sec:dishbrain}. DishBrain}

Cấu tạo của bàn thử và kết quả chơi lấy từ một công trình thực nghiệm và phần tiếp theo của nó; các công thức về nhiễu và về sự trôi tới cấu hình tốt được suy ra trong chương và kiểm tra bằng mô hình của sách, còn cơ chế học trong đĩa nuôi thì chưa được xác lập.

{\footnotesize
\setlength{\tabcolsep}{3pt}\renewcommand{\arraystretch}{1.15}
\begin{longtable}{>{\raggedright}p{0.5cm}>{\raggedright}p{4.0cm}>{\raggedright}p{5.6cm}>{\raggedright}p{2.3cm}>{\raggedright\arraybackslash}p{1.7cm}}
\toprule
TT & Khẳng định & Thí nghiệm và điều đã đo & Nguồn & Trạng thái \\
\midrule\endfirsthead
\toprule
TT & Khẳng định & Thí nghiệm và điều đã đo & Nguồn & Trạng thái \\
\midrule\endhead
1 & Bàn thử: khoảng $800\,000$ tế bào vỏ não chuột hoặc khoảng $10^6$ tế bào người trên chip; $26\,000$ điện cực trên $8$~mm$^2$, ghi tối đa $1024$, kích thích qua $8$ & Phương pháp của công trình: chip MaxOne, $26\,000$ điện cực bạch kim trên $8$~mm$^2$, ghi $1024$ kênh, mỗi kênh $20\,000$ mẫu mỗi giây; về kỹ thuật có thể kích thích tới $32$ điện cực, kích thích độc lập được $8$; mẫu chuột được dùng từ khoảng $10$ ngày tuổi & \cite{Kagan2022} & đã đo \\
2 & Vòng lặp với nhịp $10$~ms: xung của các vùng vận động đẩy vợt (hình~\ref{fig:dishloop}) & Xung được đếm theo cửa sổ $10$~ms; độ trễ từ xung đến kích thích đáp lại khoảng $5$~ms, do bộ đệm của thiết bị quyết định & \cite{Kagan2022} & đã đo \\
3 & Đầu vào: mã vị trí (điện cực nào trong $8$) và mã tần số $4$--$40$~xung/s & Trong thí nghiệm chính, tần số kích thích tăng tuyến tính từ $4$~Hz khi bóng ở tường xa đến $40$~Hz ở vợt; biên độ xung của bóng là $75$~mV; xung chữ nhật hai pha & \cite{Kagan2022} & đã đo \\
4 & Đầu ra $\Delta p=c\,(n_\uparrow-n_\downarrow)$~\eqref{eq:dishreadout}, nhiễu của số đếm mỗi cửa sổ $1/\sqrt{RT}$ & Quy tắc hiệu số của hai vùng được mô tả trong công trình (với trọng số các vùng theo hoạt động), công thức chính xác không được đưa ra. Ước tính nhiễu là hệ quả của phép đếm Poisson, được suy ra trong chương & \cite{Kagan2022}; suy ra trong chương & lý thuyết \\
5 & Thầy dạy: sau cú trúng là kích thích đoán trước được $75$~mV, $100$~xung/s, $0.1$~s; sau cú trượt là kích thích không đoán trước được, $150$~mV, $5$~xung/s, $4$~s vào vị trí và thời điểm ngẫu nhiên, sau đó nghỉ $4$~s & Phương pháp: sau cú trúng $75$~mV, $100$~Hz, $100$~ms trên cả $8$ điện cực cùng lúc; sau cú trượt $150$~mV, $5$~Hz vào các điện cực ngẫu nhiên ở các thời điểm ngẫu nhiên trong $4$~s, sau đó $4$~s không kích thích. Mẫu nuôi cấy không có dopamine & \cite{Kagan2022} & đã đo \\
6 & Mạng hành động sao cho đầu vào của nó trở nên dễ đoán & Nguyên lý năng lượng tự do là một đề xuất lý thuyết mà từ đó các tác giả suy ra giao thức; thí nghiệm phù hợp với nó, nhưng không phân biệt nó với các cơ chế đơn giản hơn (dòng~7--8) & \cite{Friston2010,Kagan2022} & giả thuyết \\
7 & Nếu thất bại làm mạng bị lắc còn thành công thì không, thời gian ở một cấu hình $\rho\propto1/P_{\text{trượt}}$~\eqref{eq:dishdensity} (mệnh đề~\ref{prop:dishscramble}) & Suy ra từ cân bằng dòng trong xích Markov với các cú hích đối xứng, được chứng minh trong chương. Chưa có kiểm tra trực tiếp trên mẫu nuôi cấy & suy ra trong chương & lý thuyết \\
8 & Mô hình ``xáo trộn khi trượt'' học được: tỷ lệ đỡ trúng $0.21\to0.38$; khi hích sau mỗi lượt đánh $\approx0.12$, không hích $0.19$ (hình~\ref{fig:dishmodel}) & Tính toán, script \texttt{models/dishbrain\_model.py}, $40$ mô hình mẫu nuôi cấy, mỗi mô hình $2000$ lượt đánh: $0.21\to0.38$; $0.13\to0.11$; $0.19\to0.19$ & --- & mô hình \\
9 & Chuỗi cú đỡ trúng dài ra chỉ ở các mẫu sống trong vòng lặp đầy đủ; các đối chứng không học & $399$ phiên, mỗi phiên $20$~phút: $9$ mẫu chuột ($101$ phiên), $11$ mẫu người ($138$), $6$ chip chỉ có môi trường nuôi ($80$), $20$ mẫu ở trạng thái nghỉ ($42$), $3$ mô phỏng ($38$). Độ dài trung bình của lượt đánh trong $15$~phút cuối lớn hơn $5$ phút đầu chỉ ở các mẫu sống; ở tế bào không phải nơ-ron (HEK293T) và ở chip chỉ có môi trường nuôi thì không có hiệu ứng & \cite{Kagan2022} & đã đo \\
10 & Mức tăng có ý nghĩa trong một phiên chỉ khi phản hồi đầy đủ; với im lặng thay cho kích thích, cuối phiên chơi tốt hơn so với không phản hồi, nhưng hiệu ứng nhỏ hơn & $486$ phiên trong $3$ ngày: ``kích thích'' ($n=27$), ``im lặng'' ($17$), ``không phản hồi'' ($15$). Mức tăng theo thời gian chỉ có ý nghĩa ở điều kiện ``kích thích''; cuối phiên, điều kiện ``im lặng'' hơn ``không phản hồi'' với hiệu ứng nhỏ hơn & \cite{Kagan2022} & đã đo \\
11 & Hiệu ứng nhỏ; mạng có học lâu dài không là câu hỏi mở & Trong một phiên, sự cải thiện đáng tin cậy, nhưng theo chính các tác giả, việc học giữa các phiên qua vài ngày không được quan sát ổn định & \cite{Kagan2022} & đã đo \\
12 & Khi chơi, mạng tiến gần chế độ tới hạn, và càng gần thì chơi càng tốt & Phân tích các trận thác hoạt động của cùng các mẫu nuôi cấy: đầu vào có cấu trúc đưa mạng tới gần chế độ tới hạn, chất lượng chơi tương quan với mức độ gần; nhưng chỉ riêng tính tới hạn, không có thông tin về hệ quả của hành động, thì chưa đủ cho việc học & \cite{Habibollahi2023} & đã đo \\
13 & ``Tri giác'' của mẫu nuôi cấy chưa được chứng minh & Bài phản hồi của ba mươi nhà nghiên cứu: thay đổi hành vi trong vòng kín không chứng minh trí tuệ hay cảm giác; chưa có thí nghiệm nào cho thấy điều đó & \cite{Balci2023} & giả thuyết \\
14 & Đối chứng thứ tư được đề xuất: kích thích đoán trước được sau cú trượt, cùng độ dài và năng lượng (mục~\ref{sec:dishspec}) & Thí nghiệm này chưa được tiến hành; đây là đề xuất của cuốn sách & --- & giả thuyết \\
\bottomrule
\end{longtable}}


\begin{thebibliography}{99}
\bibitem{Azevedo2009} F. A. C. Azevedo et al., \emph{Equal numbers of neuronal and nonneuronal cells make the human brain an isometrically scaled-up primate brain}, J. Comp. Neurol. \textbf{513}, 532 (2009).
\bibitem{Schultz1997} W. Schultz, P. Dayan, and P. R. Montague, \emph{A neural substrate of prediction and reward}, Science \textbf{275}, 1593 (1997).
\bibitem{SuttonBarto2018} R. S. Sutton and A. G. Barto, \emph{Reinforcement Learning: An Introduction}, 2nd ed. (MIT Press, Cambridge, MA, 2018).
\bibitem{Doya2002} K. Doya, \emph{Metalearning and neuromodulation}, Neural Networks \textbf{15}, 495 (2002).
\bibitem{Strata1999} P. Strata and R. Harvey, \emph{Dale's principle}, Brain Res. Bull. \textbf{50}, 349 (1999).
\bibitem{vanVreeswijk1996} C. van Vreeswijk and H. Sompolinsky, \emph{Chaos in neuronal networks with balanced excitatory and inhibitory activity}, Science \textbf{274}, 1724 (1996).
\bibitem{AstonJones2005} G. Aston-Jones and J. D. Cohen, \emph{An integrative theory of locus coeruleus-norepinephrine function: adaptive gain and optimal performance}, Annu. Rev. Neurosci. \textbf{28}, 403 (2005).
\bibitem{Beaulieu2011} J.-M. Beaulieu and R. R. Gainetdinov, \emph{The physiology, signaling, and pharmacology of dopamine receptors}, Pharmacol. Rev. \textbf{63}, 182 (2011).
\bibitem{Sparso2001} J. Spars{\o} and S. Furber (eds.), \emph{Principles of Asynchronous Circuit Design: A Systems Perspective} (Kluwer, Boston, 2001).
\bibitem{Carr1990} C. E. Carr and M. Konishi, \emph{A circuit for detection of interaural time differences in the brain stem of the barn owl}, J. Neurosci. \textbf{10}, 3227 (1990).
\bibitem{Wang2001} X.-J. Wang, \emph{Synaptic reverberation underlying mnemonic persistent activity}, Trends Neurosci. \textbf{24}, 455 (2001).
\bibitem{Hahnloser2002} R. H. R. Hahnloser, A. A. Kozhevnikov, and M. S. Fee, \emph{An ultra-sparse code underlies the generation of neural sequences in a songbird}, Nature \textbf{419}, 65 (2002).
\bibitem{Barlow1965} H. B. Barlow and W. R. Levick, \emph{The mechanism of directionally selective units in rabbit's retina}, J. Physiol. \textbf{178}, 477 (1965).
\bibitem{Carandini2012} M. Carandini and D. J. Heeger, \emph{Normalization as a canonical neural computation}, Nat. Rev. Neurosci. \textbf{13}, 51 (2012).
\bibitem{Pouille2001} F. Pouille and M. Scanziani, \emph{Enforcement of temporal fidelity in pyramidal cells by somatic feed-forward inhibition}, Science \textbf{293}, 1159 (2001).
\bibitem{Tsodyks1997isn} M. V. Tsodyks, W. E. Skaggs, T. J. Sejnowski, and B. L. McNaughton, \emph{Paradoxical effects of external modulation of inhibitory interneurons}, J. Neurosci. \textbf{17}, 4382 (1997).
\bibitem{Buzsaki2012} G. Buzs\'aki and X.-J. Wang, \emph{Mechanisms of gamma oscillations}, Annu. Rev. Neurosci. \textbf{35}, 203 (2012).
\bibitem{Rieke1997} F. Rieke, D. Warland, R. de Ruyter van Steveninck, and W. Bialek, \emph{Spikes: Exploring the Neural Code} (MIT Press, Cambridge, MA, 1997).
\bibitem{Mainen1995} Z. F. Mainen and T. J. Sejnowski, \emph{Reliability of spike timing in neocortical neurons}, Science \textbf{268}, 1503 (1995).
\bibitem{Softky1993} W. R. Softky and C. Koch, \emph{The highly irregular firing of cortical cells is inconsistent with temporal integration of random EPSPs}, J. Neurosci. \textbf{13}, 334 (1993).
\bibitem{Thorpe1996} S. Thorpe, D. Fize, and C. Marlot, \emph{Speed of processing in the human visual system}, Nature \textbf{381}, 520 (1996).
\bibitem{Zohary1994} E. Zohary, M. N. Shadlen, and W. T. Newsome, \emph{Correlated neuronal discharge rate and its implications for psychophysical performance}, Nature \textbf{370}, 140 (1994).
\bibitem{MorenoBote2014} R. Moreno-Bote et al., \emph{Information-limiting correlations}, Nat. Neurosci. \textbf{17}, 1410 (2014).
\bibitem{Gollisch2008} T. Gollisch and M. Meister, \emph{Rapid neural coding in the retina with relative spike latencies}, Science \textbf{319}, 1108 (2008).
\bibitem{OKeefe1993} J. O'Keefe and M. L. Recce, \emph{Phase relationship between hippocampal place units and the EEG theta rhythm}, Hippocampus \textbf{3}, 317 (1993).
\bibitem{Destexhe2003} A. Destexhe, M. Rudolph, and D. Par\'e, \emph{The high-conductance state of neocortical neurons in vivo}, Nat. Rev. Neurosci. \textbf{4}, 739 (2003).
\bibitem{Tsodyks1997} M. V. Tsodyks and H. Markram, \emph{The neural code between neocortical pyramidal neurons depends on neurotransmitter release probability}, Proc. Natl. Acad. Sci. USA \textbf{94}, 719 (1997).
\bibitem{Lisman1997} J. E. Lisman, \emph{Bursts as a unit of neural information: making unreliable synapses reliable}, Trends Neurosci. \textbf{20}, 38 (1997).
\bibitem{Attwell2001} D. Attwell and S. B. Laughlin, \emph{An energy budget for signaling in the grey matter of the brain}, J. Cereb. Blood Flow Metab. \textbf{21}, 1133 (2001).
\bibitem{Lennie2003} P. Lennie, \emph{The cost of cortical computation}, Curr. Biol. \textbf{13}, 493 (2003).
\bibitem{Yagishita2014} S. Yagishita et al., \emph{A critical time window for dopamine actions on the structural plasticity of dendritic spines}, Science \textbf{345}, 1616 (2014).
\bibitem{Werfel2005} J. Werfel, X. Xie, and H. S. Seung, \emph{Learning curves for stochastic gradient descent in linear feedforward networks}, Neural Comput. \textbf{17}, 2699 (2005).
\bibitem{Doya2000} K. Doya, \emph{Reinforcement learning in continuous time and space}, Neural Comput. \textbf{12}, 219 (2000).
\bibitem{Oja1982} E. Oja, \emph{Simplified neuron model as a principal component analyzer}, J. Math. Biol. \textbf{15}, 267 (1982).
\bibitem{BiPoo1998} G.-Q. Bi and M.-M. Poo, \emph{Synaptic modifications in cultured hippocampal neurons: dependence on spike timing, synaptic strength, and postsynaptic cell type}, J. Neurosci. \textbf{18}, 10464 (1998).
\bibitem{Williams1992} R. J. Williams, \emph{Simple statistical gradient-following algorithms for connectionist reinforcement learning}, Mach. Learn. \textbf{8}, 229 (1992).
\bibitem{Bayer2005} H. M. Bayer and P. W. Glimcher, \emph{Midbrain dopamine neurons encode a quantitative reward prediction error signal}, Neuron \textbf{47}, 129 (2005).
\bibitem{Shen2008} W. Shen, M. Flajolet, P. Greengard, and D. J. Surmeier, \emph{Dichotomous dopaminergic control of striatal synaptic plasticity}, Science \textbf{321}, 848 (2008).
\bibitem{Dabney2020} W. Dabney, Z. Kurth-Nelson, N. Uchida, C. K. Starkweather, D. Hassabis, R. Munos, and M. Botvinick, \emph{A distributional code for value in dopamine-based reinforcement learning}, Nature \textbf{577}, 671 (2020).
\bibitem{Matsumoto2009} M. Matsumoto and O. Hikosaka, \emph{Two types of dopamine neuron distinctly convey positive and negative motivational signals}, Nature \textbf{459}, 837 (2009).
\bibitem{BrombergMartin2010} E. S. Bromberg-Martin, M. Matsumoto, and O. Hikosaka, \emph{Dopamine in motivational control: rewarding, aversive, and alerting}, Neuron \textbf{68}, 815 (2010).
\bibitem{Menegas2018} W. Menegas, K. Akiti, R. Amo, N. Uchida, and M. Watabe-Uchida, \emph{Dopamine neurons projecting to the posterior striatum reinforce avoidance of threatening stimuli}, Nat. Neurosci. \textbf{21}, 1421 (2018).
\bibitem{Niv2007} Y. Niv, N. D. Daw, D. Joel, and P. Dayan, \emph{Tonic dopamine: opportunity costs and the control of response vigor}, Psychopharmacology \textbf{191}, 507 (2007).
\bibitem{Mohebi2019} A. Mohebi et al., \emph{Dissociable dopamine dynamics for learning and motivation}, Nature \textbf{570}, 65 (2019).
\bibitem{Hamid2016} A. A. Hamid et al., \emph{Mesolimbic dopamine signals the value of work}, Nat. Neurosci. \textbf{19}, 117 (2016).
\bibitem{Kim2020} H. R. Kim, A. N. Malik et al., \emph{A unified framework for dopamine signals across timescales}, Cell \textbf{183}, 1600 (2020).
\bibitem{Jeong2022} H. Jeong et al., \emph{Mesolimbic dopamine release conveys causal associations}, Science \textbf{378}, eabq6740 (2022).
\bibitem{Amo2022} R. Amo et al., \emph{A gradual temporal shift of dopamine responses mirrors the progression of temporal difference error in machine learning}, Nat. Neurosci. \textbf{25}, 1082 (2022).
\bibitem{Takahashi2017} Y. K. Takahashi et al., \emph{Dopamine neurons respond to errors in the prediction of sensory features of expected rewards}, Neuron \textbf{95}, 1395 (2017).
\bibitem{Sharpe2017} M. J. Sharpe et al., \emph{Dopamine transients are sufficient and necessary for acquisition of model-based associations}, Nat. Neurosci. \textbf{20}, 735 (2017).
\bibitem{Gardner2018} M. P. H. Gardner, G. Schoenbaum, and S. J. Gershman, \emph{Rethinking dopamine as generalized prediction error}, Proc. R. Soc. B \textbf{285}, 20181645 (2018).
\bibitem{Engelhard2019} B. Engelhard et al., \emph{Specialized coding of sensory, motor and cognitive variables in VTA dopamine neurons}, Nature \textbf{570}, 509 (2019).
\bibitem{Clements1992} J. D. Clements, R. A. J. Lester, G. Tong, C. E. Jahr, and G. L. Westbrook, \emph{The time course of glutamate in the synaptic cleft}, Science \textbf{258}, 1498 (1992).
\bibitem{Hasselmo2006} M. E. Hasselmo, \emph{The role of acetylcholine in learning and memory}, Curr. Opin. Neurobiol. \textbf{16}, 710 (2006).
\bibitem{JacobsAzmitia1992} B. L. Jacobs and E. C. Azmitia, \emph{Structure and function of the brain serotonin system}, Physiol. Rev. \textbf{72}, 165 (1992).
\bibitem{Schiller1992} P. H. Schiller, \emph{The ON and OFF channels of the visual system}, Trends Neurosci. \textbf{15}, 86 (1992).
\bibitem{Grothe2010} B. Grothe, M. Pecka, and D. McAlpine, \emph{Mechanisms of sound localization in mammals}, Physiol. Rev. \textbf{90}, 983 (2010).
\bibitem{Markram2004} H. Markram et al., \emph{Interneurons of the neocortical inhibitory system}, Nat. Rev. Neurosci. \textbf{5}, 793 (2004).
\bibitem{Joshi2016} S. Joshi, Y. Li, R. M. Kalwani, and J. I. Gold, \emph{Relationships between pupil diameter and neuronal activity in the locus coeruleus, colliculi, and cingulate cortex}, Neuron \textbf{89}, 221 (2016).
\bibitem{ServanSchreiber1990} D. Servan-Schreiber, H. Printz, and J. D. Cohen, \emph{A network model of catecholamine effects: gain, signal-to-noise ratio, and behavior}, Science \textbf{249}, 892 (1990).
\bibitem{YuDayan2005} A. J. Yu and P. Dayan, \emph{Uncertainty, neuromodulation, and attention}, Neuron \textbf{46}, 681 (2005).
\bibitem{Nassar2012} M. R. Nassar et al., \emph{Rational regulation of learning dynamics by pupil-linked arousal systems}, Nat. Neurosci. \textbf{15}, 1040 (2012).
\bibitem{BouretSara2005} S. Bouret and S. J. Sara, \emph{Network reset: a simplified overarching theory of locus coeruleus noradrenaline function}, Trends Neurosci. \textbf{28}, 574 (2005).
\bibitem{Hebb1949} D. O. Hebb, \emph{The Organization of Behavior: A Neuropsychological Theory} (Wiley, New York, 1949).
\bibitem{MacKay1952} D. M. MacKay and W. S. McCulloch, \emph{The limiting information capacity of a neuronal link}, Bull. Math. Biophys. \textbf{14}, 127 (1952).
\bibitem{WilsonCowan1972} H. R. Wilson and J. D. Cowan, \emph{Excitatory and inhibitory interactions in localized populations of model neurons}, Biophys. J. \textbf{12}, 1 (1972).
\bibitem{AllenStevens1994} C. Allen and C. F. Stevens, \emph{An evaluation of causes for unreliability of synaptic transmission}, Proc. Natl. Acad. Sci. USA \textbf{91}, 10380 (1994).
\bibitem{Barlow1961} H. B. Barlow, \emph{Possible principles underlying the transformations of sensory messages}, in \emph{Sensory Communication}, ed. W. A. Rosenblith (MIT Press, Cambridge, MA, 1961), pp.~217--234.
\bibitem{Cardin2009} J. A. Cardin et al., \emph{Driving fast-spiking cells induces gamma rhythm and controls sensory responses}, Nature \textbf{459}, 663 (2009).
\bibitem{Lillicrap2020} T. P. Lillicrap, A. Santoro, L. Marris, C. J. Akerman, and G. Hinton, \emph{Backpropagation and the brain}, Nat. Rev. Neurosci. \textbf{21}, 335 (2020).
\bibitem{Fremaux2016} N. Fr\'emaux and W. Gerstner, \emph{Neuromodulated spike-timing-dependent plasticity, and theory of three-factor learning rules}, Front. Neural Circuits \textbf{9}, 85 (2016).
\bibitem{Haider2006} B. Haider, A. Duque, A. R. Hasenstaub, and D. A. McCormick, \emph{Neocortical network activity in vivo is generated through a dynamic balance of excitation and inhibition}, J. Neurosci. \textbf{26}, 4535 (2006).
\bibitem{HoltKoch1997} G. R. Holt and C. Koch, \emph{Shunting inhibition does not have a divisive effect on firing rates}, Neural Comput. \textbf{9}, 1001 (1997).
\bibitem{Chance2002} F. S. Chance, L. F. Abbott, and A. D. Reyes, \emph{Gain modulation from background synaptic input}, Neuron \textbf{35}, 773 (2002).
\bibitem{Ozeki2009} H. Ozeki, I. M. Finn, E. S. Schaffer, K. D. Miller, and D. Ferster, \emph{Inhibitory stabilization of the cortical network underlies visual surround suppression}, Neuron \textbf{62}, 578 (2009).
\bibitem{HasselmoBower1992} M. E. Hasselmo and J. M. Bower, \emph{Cholinergic suppression specific to intrinsic not afferent fiber synapses in rat piriform (olfactory) cortex}, J. Neurophysiol. \textbf{67}, 1222 (1992).
\bibitem{HasselmoAndersonBower1992} M. E. Hasselmo, B. P. Anderson, and J. M. Bower, \emph{Cholinergic modulation of cortical associative memory function}, J. Neurophysiol. \textbf{67}, 1230 (1992).
\bibitem{Hasselmo1999} M. E. Hasselmo, \emph{Neuromodulation: acetylcholine and memory consolidation}, Trends Cogn. Sci. \textbf{3}, 351 (1999).
\bibitem{Tsodyks1988} M. V. Tsodyks and M. V. Feigel'man, \emph{The enhanced storage capacity in neural networks with low activity level}, Europhys. Lett. \textbf{6}, 101 (1988).
\bibitem{KalmanBucy1961} R. E. Kalman and R. S. Bucy, \emph{New results in linear filtering and prediction theory}, J. Basic Eng. \textbf{83}, 95 (1961).
\bibitem{WilsonMcNaughton1994} M. A. Wilson and B. L. McNaughton, \emph{Reactivation of hippocampal ensemble memories during sleep}, Science \textbf{265}, 676 (1994).
\bibitem{BarnesSharp1999} N. M. Barnes and T. Sharp, \emph{A review of central 5-HT receptors and their function}, Neuropharmacology \textbf{38}, 1083 (1999).
\bibitem{Bunin1998} M. A. Bunin and R. M. Wightman, \emph{Quantitative evaluation of 5-hydroxytryptamine (serotonin) neuronal release and uptake: an investigation of extrasynaptic transmission}, J. Neurosci. \textbf{18}, 4854 (1998).
\bibitem{Sprouse1987} J. S. Sprouse and G. K. Aghajanian, \emph{Electrophysiological responses of serotoninergic dorsal raphe neurons to 5-HT$_{1A}$ and 5-HT$_{1B}$ agonists}, Synapse \textbf{1}, 3 (1987).
\bibitem{Sykova2008} E. Syková and C. Nicholson, \emph{Diffusion in brain extracellular space}, Physiol. Rev. \textbf{88}, 1277 (2008).
\bibitem{WilsonNicoll2001} R. I. Wilson and R. A. Nicoll, \emph{Endogenous cannabinoids mediate retrograde signalling at hippocampal synapses}, Nature \textbf{410}, 588 (2001).
\bibitem{Garthwaite2008} J. Garthwaite, \emph{Concepts of neural nitric oxide-mediated transmission}, Eur. J. Neurosci. \textbf{27}, 2783 (2008).
\bibitem{vandenPol2012} A. N. van den Pol, \emph{Neuropeptide transmission in brain circuits}, Neuron \textbf{76}, 98 (2012).
\bibitem{Dunwiddie2001} T. V. Dunwiddie and S. A. Masino, \emph{The role and regulation of adenosine in the central nervous system}, Annu. Rev. Neurosci. \textbf{24}, 31 (2001).
\bibitem{Skaggs1996} W. E. Skaggs and B. L. McNaughton, \emph{Replay of neuronal firing sequences in rat hippocampus during sleep following spatial experience}, Science \textbf{271}, 1870 (1996).
\bibitem{Baylor1979} D. A. Baylor, T. D. Lamb, and K.-W. Yau, \emph{Responses of retinal rods to single photons}, J. Physiol. \textbf{288}, 613 (1979).
\bibitem{Hudspeth1989} A. J. Hudspeth, \emph{How the ear's works work}, Nature \textbf{341}, 397 (1989).
\bibitem{NakaRushton1966} K. I. Naka and W. A. H. Rushton, \emph{S-potentials from colour units in the retina of fish (Cyprinidae)}, J. Physiol. \textbf{185}, 536 (1966).
\bibitem{Lichtsteiner2008} P. Lichtsteiner, C. Posch, and T. Delbruck, \emph{A $128\times128$ 120 dB 15 $\mu$s latency asynchronous temporal contrast vision sensor}, IEEE J. Solid-State Circuits \textbf{43}, 566 (2008).
\bibitem{Koch2006} K. Koch et al., \emph{How much the eye tells the brain}, Curr. Biol. \textbf{16}, 1428 (2006).
\bibitem{Robles2001} L. Robles and M. A. Ruggero, \emph{Mechanics of the mammalian cochlea}, Physiol. Rev. \textbf{81}, 1305 (2001).
\bibitem{Mombaerts2004} P. Mombaerts, \emph{Genes and ligands for odorant, vomeronasal and taste receptors}, Nat. Rev. Neurosci. \textbf{5}, 263 (2004).
\bibitem{WoodSlater2001} S. J. Wood and C. R. Slater, \emph{Safety factor at the neuromuscular junction}, Prog. Neurobiol. \textbf{64}, 393 (2001).
\bibitem{Henneman1957} E. Henneman, \emph{Relation between size of neurons and their susceptibility to discharge}, Science \textbf{126}, 1345 (1957).
\bibitem{HarrisWolpert1998} C. M. Harris and D. M. Wolpert, \emph{Signal-dependent noise determines motor planning}, Nature \textbf{394}, 780 (1998).
\bibitem{Hogan1984} N. Hogan, \emph{Adaptive control of mechanical impedance by coactivation of antagonist muscles}, IEEE Trans. Autom. Control \textbf{29}, 681 (1984).
\bibitem{McAuleyMarsden2000} J. H. McAuley and C. D. Marsden, \emph{Physiological and pathological tremors and rhythmic central motor control}, Brain \textbf{123}, 1545 (2000).
\bibitem{WolpertGhahramani2000} D. M. Wolpert and Z. Ghahramani, \emph{Computational principles of movement neuroscience}, Nat. Neurosci. \textbf{3}, 1212 (2000).
\bibitem{MarderBucher2001} E. Marder and D. Bucher, \emph{Central pattern generators and the control of rhythmic movements}, Curr. Biol. \textbf{11}, R986 (2001).
\bibitem{Miyazaki2014} K. W. Miyazaki et al., \emph{Optogenetic activation of dorsal raphe serotonin neurons enhances patience for future rewards}, Curr. Biol. \textbf{24}, 2033 (2014).
\bibitem{Fuglevand1993} A. J. Fuglevand, D. A. Winter, and A. E. Patla, \emph{Models of recruitment and rate coding organization in motor-unit pools}, J. Neurophysiol. \textbf{70}, 2470 (1993).
\bibitem{Curcio1990} C. A. Curcio, K. R. Sloan, R. E. Kalina, and A. E. Hendrickson, \emph{Human photoreceptor topography}, J. Comp. Neurol. \textbf{292}, 497 (1990).
\bibitem{Hayes1950} N. D. Hayes, \emph{Roots of the transcendental equation associated with a certain difference-differential equation}, J. London Math. Soc. \textbf{25}, 226 (1950).
\bibitem{MelzackWall1965} R. Melzack and P. D. Wall, \emph{Pain mechanisms: a new theory}, Science \textbf{150}, 971 (1965).
\bibitem{Fields2004} H. Fields, \emph{State-dependent opioid control of pain}, Nat. Rev. Neurosci. \textbf{5}, 565 (2004).
\bibitem{Levine1978} J. D. Levine, N. C. Gordon, and H. L. Fields, \emph{The mechanism of placebo analgesia}, Lancet \textbf{312}, 654 (1978).
\bibitem{Woolf1983} C. J. Woolf, \emph{Evidence for a central component of post-injury pain hypersensitivity}, Nature \textbf{306}, 686 (1983).
\bibitem{Seymour2004} B. Seymour et al., \emph{Temporal difference models describe higher-order learning in humans}, Nature \textbf{429}, 664 (2004).
\bibitem{EcclestonCrombez1999} C. Eccleston and G. Crombez, \emph{Pain demands attention: a cognitive-affective model of the interruptive function of pain}, Psychol. Bull. \textbf{125}, 356 (1999).
\bibitem{WidrowHoff1960} B. Widrow and M. E. Hoff, \emph{Adaptive switching circuits}, IRE WESCON Convention Record, part 4, 96 (1960).
\bibitem{Marr1969} D. Marr, \emph{A theory of cerebellar cortex}, J. Physiol. \textbf{202}, 437 (1969).
\bibitem{Albus1971} J. S. Albus, \emph{A theory of cerebellar function}, Math. Biosci. \textbf{10}, 25 (1971).
\bibitem{Cover1965} T. M. Cover, \emph{Geometrical and statistical properties of systems of linear inequalities with applications in pattern recognition}, IEEE Trans. Electron. Comput. \textbf{EC-14}, 326 (1965).
\bibitem{Ito2001} M. Ito, \emph{Cerebellar long-term depression: characterization, signal transduction, and functional roles}, Physiol. Rev. \textbf{81}, 1143 (2001).
\bibitem{Suvrathan2016} A. Suvrathan, H. L. Payne, and J. L. Raymond, \emph{Timing rules for synaptic plasticity matched to behavioral function}, Neuron \textbf{92}, 959 (2016).
\bibitem{Fujita1982} M. Fujita, \emph{Adaptive filter model of the cerebellum}, Biol. Cybern. \textbf{45}, 195 (1982).
\bibitem{Blakemore1998} S.-J. Blakemore, D. M. Wolpert, and C. D. Frith, \emph{Central cancellation of self-produced tickle sensation}, Nat. Neurosci. \textbf{1}, 635 (1998).
\bibitem{Smith2006} M. A. Smith, A. Ghazizadeh, and R. Shadmehr, \emph{Interacting adaptive processes with different timescales underlie short-term motor learning}, PLoS Biol. \textbf{4}, e179 (2006).
\bibitem{Kording2007} K. P. K\"ording, J. B. Tenenbaum, and R. Shadmehr, \emph{The dynamics of memory as a consequence of optimal adaptation to a changing body}, Nat. Neurosci. \textbf{10}, 779 (2007).
\bibitem{Yao2023} Z. Yao et al., \emph{A high-resolution transcriptomic and spatial atlas of cell types in the whole mouse brain}, Nature \textbf{624}, 317 (2023).
\bibitem{Sanzeni2020} A. Sanzeni, B. Akitake, H. C. Goldbach, C. E. Leedy, N. Brunel, and M. H. Histed, \emph{Inhibition stabilization is a widespread property of cortical networks}, eLife \textbf{9}, e54875 (2020).
\bibitem{Padamsey2022} Z. Padamsey, D. Katsanevaki, N. Dupuy, and N. L. Rochefort, \emph{Neocortex saves energy by reducing coding precision during food scarcity}, Neuron \textbf{110}, 280 (2022).
\bibitem{Stringer2019} C. Stringer, M. Pachitariu, N. Steinmetz, C. B. Reddy, M. Carandini, and K. D. Harris, \emph{Spontaneous behaviors drive multidimensional, brainwide activity}, Science \textbf{364}, eaav7893 (2019).
\bibitem{Bittner2017} K. C. Bittner, A. D. Milstein, C. Grienberger, S. Romani, and J. C. Magee, \emph{Behavioral time scale synaptic plasticity underlies CA1 place fields}, Science \textbf{357}, 1033 (2017).
\bibitem{WhittingtonBogacz2019} J. C. R. Whittington and R. Bogacz, \emph{Theories of error back-propagation in the brain}, Trends Cogn. Sci. \textbf{23}, 235 (2019).
\bibitem{Reimer2016} J. Reimer et al., \emph{Pupil fluctuations track rapid changes in adrenergic and cholinergic activity in cortex}, Nat. Commun. \textbf{7}, 13289 (2016).
\bibitem{Beniaguev2021} D. Beniaguev, I. Segev, and M. London, \emph{Single cortical neurons as deep artificial neural networks}, Neuron \textbf{109}, 2727 (2021).
\bibitem{Gidon2020} A. Gidon et al., \emph{Dendritic action potentials and computation in human layer 2/3 cortical neurons}, Science \textbf{367}, 83 (2020).
\bibitem{Gallego2022} G. Gallego et al., \emph{Event-based vision: a survey}, IEEE Trans. Pattern Anal. Mach. Intell. \textbf{44}, 154 (2022).
\bibitem{Berger1989} R. D. Berger, J. P. Saul, and R. J. Cohen, \emph{Transfer function analysis of autonomic regulation. I. Canine atrial rate response}, Am. J. Physiol. \textbf{256}, H142 (1989).
\bibitem{Walker2010} J. J. Walker, J. R. Terry, and S. L. Lightman, \emph{Origin of ultradian pulsatility in the hypothalamic--pituitary--adrenal axis}, Proc. R. Soc. B \textbf{277}, 1627 (2010).
\bibitem{Belchetz1978} P. E. Belchetz, T. M. Plant, Y. Nakai, E. J. Keogh, and E. Knobil, \emph{Hypophysial responses to continuous and intermittent delivery of hypothalamic gonadotropin-releasing hormone}, Science \textbf{202}, 631 (1978).
\bibitem{Czeisler1999} C. A. Czeisler et al., \emph{Stability, precision, and near-24-hour period of the human circadian pacemaker}, Science \textbf{284}, 2177 (1999).
\bibitem{TakahashiJS2017} J. S. Takahashi, \emph{Transcriptional architecture of the mammalian circadian clock}, Nat. Rev. Genet. \textbf{18}, 164 (2017).
\bibitem{Musk2019} E. Musk and Neuralink, \emph{An integrated brain-machine interface platform with thousands of channels}, J. Med. Internet Res. \textbf{21}, e16194 (2019).
\bibitem{Hochberg2006} L. R. Hochberg et al., \emph{Neuronal ensemble control of prosthetic devices by a human with tetraplegia}, Nature \textbf{442}, 164 (2006).
\bibitem{Willett2021} F. R. Willett et al., \emph{High-performance brain-to-text communication via handwriting}, Nature \textbf{593}, 249 (2021).
\bibitem{Willett2023} F. R. Willett et al., \emph{A high-performance speech neuroprosthesis}, Nature \textbf{620}, 1031 (2023).
\bibitem{Gallego2017} J. A. Gallego, M. G. Perich, L. E. Miller, and S. A. Solla, \emph{Neural manifolds for the control of movement}, Neuron \textbf{94}, 978 (2017).
\bibitem{Fernandez2021} E. Fern\'andez et al., \emph{Visual percepts evoked with an intracortical 96-channel microelectrode array inserted in human occipital cortex}, J. Clin. Invest. \textbf{131}, e151331 (2021).
\bibitem{Tritsch2012} N. X. Tritsch, J. B. Ding, and B. L. Sabatini, \emph{Dopaminergic neurons inhibit striatal output through non-canonical release of GABA}, Nature \textbf{490}, 262 (2012).
\bibitem{Zhang2015} S. Zhang et al., \emph{Dopaminergic and glutamatergic microdomains in a subset of rodent mesoaccumbens axons}, Nat. Neurosci. \textbf{18}, 386 (2015).
\bibitem{Mongillo2008} G. Mongillo, O. Barak, and M. Tsodyks, \emph{Synaptic theory of working memory}, Science \textbf{319}, 1543 (2008).
\bibitem{Stokes2015} M. G. Stokes, \emph{`Activity-silent' working memory in prefrontal cortex: a dynamic coding framework}, Trends Cogn. Sci. \textbf{19}, 394 (2015).
\bibitem{Smith2013} S. L. Smith, I. T. Smith, T. Branco, and M. H\"ausser, \emph{Dendritic spikes enhance stimulus selectivity in cortical neurons in vivo}, Nature \textbf{503}, 115 (2013).
\bibitem{Baden2016} T. Baden et al., \emph{The functional diversity of retinal ganglion cells in the mouse}, Nature \textbf{529}, 345 (2016).
\bibitem{Lottem2018} E. Lottem et al., \emph{Activation of serotonin neurons promotes active persistence in a probabilistic foraging task}, Nat. Commun. \textbf{9}, 1000 (2018).
\bibitem{Payeur2021} A. Payeur, J. Guerguiev, F. Zenke, B. A. Richards, and R. Naud, \emph{Burst-dependent synaptic plasticity can coordinate learning in hierarchical circuits}, Nat. Neurosci. \textbf{24}, 1010 (2021).
\bibitem{MehonicKenyon2022} A. Mehonic and A. J. Kenyon, \emph{Brain-inspired computing needs a master plan}, Nature \textbf{604}, 255 (2022).
\bibitem{Poe2020} G. R. Poe et al., \emph{Locus coeruleus: a new look at the blue spot}, Nat. Rev. Neurosci. \textbf{21}, 644 (2020).
\bibitem{Hangya2015} B. Hangya, S. P. Ranade, M. Lorenc, and A. Kepecs, \emph{Central cholinergic neurons are rapidly recruited by reinforcement feedback}, Cell \textbf{162}, 1155 (2015).
\bibitem{FernandezRuiz2019} A. Fern\'andez-Ruiz et al., \emph{Long-duration hippocampal sharp wave ripples improve memory}, Science \textbf{364}, 1082 (2019).
\bibitem{Herzfeld2018} D. J. Herzfeld, Y. Kojima, R. Soetedjo, and R. Shadmehr, \emph{Encoding of error and learning to correct that error by the Purkinje cells of the cerebellum}, Nat. Neurosci. \textbf{21}, 736 (2018).
\bibitem{Duan2014} B. Duan et al., \emph{Identification of spinal circuits transmitting and gating mechanical pain}, Cell \textbf{159}, 1417 (2014).
\bibitem{Tremblay2016} R. Tremblay, S. Lee, and B. Rudy, \emph{GABAergic interneurons in the neocortex: from cellular properties to circuits}, Neuron \textbf{91}, 260 (2016).
\bibitem{Ren2018} J. Ren et al., \emph{Anatomically defined and functionally distinct dorsal raphe serotonin sub-systems}, Cell \textbf{175}, 472 (2018).
\bibitem{Pfeffer2013} C. K. Pfeffer, M. Xue, M. He, Z. J. Huang, and M. Scanziani, \emph{Inhibition of inhibition in visual cortex: the logic of connections between molecularly distinct interneurons}, Nat. Neurosci. \textbf{16}, 1068 (2013).
\bibitem{Pi2013} H.-J. Pi et al., \emph{Cortical interneurons that specialize in disinhibitory control}, Nature \textbf{503}, 521 (2013).
\bibitem{Keller2018} G. B. Keller and T. D. Mrsic-Flogel, \emph{Predictive processing: a canonical cortical computation}, Neuron \textbf{100}, 424 (2018).
\bibitem{HubelWiesel1962} D. H. Hubel and T. N. Wiesel, \emph{Receptive fields, binocular interaction and functional architecture in the cat's visual cortex}, J. Physiol. \textbf{160}, 106 (1962).
\bibitem{Riesenhuber1999} M. Riesenhuber and T. Poggio, \emph{Hierarchical models of object recognition in cortex}, Nat. Neurosci. \textbf{2}, 1019 (1999).
\bibitem{Fukushima1980} K. Fukushima, \emph{Neocognitron: a self-organizing neural network model for a mechanism of pattern recognition unaffected by shift in position}, Biol. Cybern. \textbf{36}, 193 (1980).
\bibitem{Yamins2014} D. L. K. Yamins et al., \emph{Performance-optimized hierarchical models predict neural responses in higher visual cortex}, Proc. Natl. Acad. Sci. USA \textbf{111}, 8619 (2014).
\bibitem{Hung2005} C. P. Hung, G. Kreiman, T. Poggio, and J. J. DiCarlo, \emph{Fast readout of object identity from macaque inferior temporal cortex}, Science \textbf{310}, 863 (2005).
\bibitem{WiskottSejnowski2002} L. Wiskott and T. J. Sejnowski, \emph{Slow feature analysis: unsupervised learning of invariances}, Neural Comput. \textbf{14}, 715 (2002).
\bibitem{Foldiak1991} P. F\"oldi\'ak, \emph{Learning invariance from transformation sequences}, Neural Comput. \textbf{3}, 194 (1991).
\bibitem{LiDiCarlo2008} N. Li and J. J. DiCarlo, \emph{Unsupervised natural experience rapidly alters invariant object representation in visual cortex}, Science \textbf{321}, 1502 (2008).
\bibitem{RaoBallard1999} R. P. N. Rao and D. H. Ballard, \emph{Predictive coding in the visual cortex: a functional interpretation of some extra-classical receptive-field effects}, Nat. Neurosci. \textbf{2}, 79 (1999).
\bibitem{Kar2019} K. Kar, J. Kubilius, K. Schmidt, E. B. Issa, and J. J. DiCarlo, \emph{Evidence that recurrent circuits are critical to the ventral stream's execution of core object recognition behavior}, Nat. Neurosci. \textbf{22}, 974 (2019).
\bibitem{Szegedy2014} C. Szegedy et al., \emph{Intriguing properties of neural networks}, in \emph{Proc. ICLR} (2014), arXiv:1312.6199.
\bibitem{Weiss2002} Y. Weiss, E. P. Simoncelli, and E. H. Adelson, \emph{Motion illusions as optimal percepts}, Nat. Neurosci. \textbf{5}, 598 (2002).
\bibitem{Purves2011} D. Purves, W. T. Wojtach, and R. B. Lotto, \emph{Understanding vision in wholly empirical terms}, Proc. Natl. Acad. Sci. USA \textbf{108}, Suppl.~3, 15588 (2011).
\bibitem{LaferSousa2015} R. Lafer-Sousa, K. L. Hermann, and B. R. Conway, \emph{Striking individual differences in color perception uncovered by `the dress' photograph}, Curr. Biol. \textbf{25}, R545 (2015).
\bibitem{Wallisch2017} P. Wallisch, \emph{Illumination assumptions account for individual differences in the perceptual interpretation of a profoundly ambiguous stimulus in the color domain: ``The dress''}, J. Vis. \textbf{17}(4), 5 (2017).
\bibitem{Schwartz2009} O. Schwartz, T. J. Sejnowski, and P. Dayan, \emph{Perceptual organization in the tilt illusion}, J. Vis. \textbf{9}(4), 19 (2009).
\bibitem{SchillerCarvey2005} P. H. Schiller and C. E. Carvey, \emph{The Hermann grid illusion revisited}, Perception \textbf{34}, 1375 (2005).
\bibitem{Nijhawan1994} R. Nijhawan, \emph{Motion extrapolation in catching}, Nature \textbf{370}, 256 (1994).
\bibitem{Conway2005} B. R. Conway, A. Kitaoka, A. Yazdanbakhsh, C. C. Pack, and M. S. Livingstone, \emph{Neural basis for a powerful static motion illusion}, J. Neurosci. \textbf{25}, 5651 (2005).
\bibitem{OteroMillan2012} J. Otero-Millan, S. L. Macknik, and S. Martinez-Conde, \emph{Microsaccades and blinks trigger illusory rotation in the ``rotating snakes'' illusion}, J. Neurosci. \textbf{32}, 6043 (2012).
\bibitem{MorenoBote2007} R. Moreno-Bote, J. Rinzel, and N. Rubin, \emph{Noise-induced alternations in an attractor network model of perceptual bistability}, J. Neurophysiol. \textbf{98}, 1125 (2007).
\bibitem{Watanabe2018} E. Watanabe, A. Kitaoka, K. Sakamoto, M. Yasugi, and K. Tanaka, \emph{Illusory motion reproduced by deep neural networks trained for prediction}, Front. Psychol. \textbf{9}, 345 (2018).
\bibitem{GomezVilla2020} A. G\'omez-Villa, A. Mart\'in, J. Vazquez-Corral, M. Bertalm\'io, and J. Malo, \emph{Color illusions also deceive CNNs for low-level vision tasks: analysis and implications}, Vision Res. \textbf{176}, 156 (2020).
\bibitem{delCastillo1954} J. del Castillo and B. Katz, \emph{Quantal components of the end-plate potential}, J. Physiol. \textbf{124}, 560 (1954).
\bibitem{Nowak1984} L. Nowak, P. Bregestovski, P. Ascher, A. Herbet, and A. Prochiantz, \emph{Magnesium gates glutamate-activated channels in mouse central neurones}, Nature \textbf{307}, 462 (1984).
\bibitem{Malinow2002} R. Malinow and R. C. Malenka, \emph{AMPA receptor trafficking and synaptic plasticity}, Annu. Rev. Neurosci. \textbf{25}, 103 (2002).
\bibitem{Matsuzaki2004} M. Matsuzaki, N. Honkura, G. C. R. Ellis-Davies, and H. Kasai, \emph{Structural basis of long-term potentiation in single dendritic spines}, Nature \textbf{429}, 761 (2004).
\bibitem{Rudomin1999} P. Rudomin and R. F. Schmidt, \emph{Presynaptic inhibition in the vertebrate spinal cord revisited}, Exp. Brain Res. \textbf{129}, 1 (1999).
\bibitem{HutcheonYarom2000} B. Hutcheon and Y. Yarom, \emph{Resonance, oscillation and the intrinsic frequency preferences of neurons}, Trends Neurosci. \textbf{23}, 216 (2000).
\bibitem{Seung1996} H. S. Seung, \emph{How the brain keeps the eyes still}, Proc. Natl. Acad. Sci. USA \textbf{93}, 13339 (1996).
\bibitem{Hounsgaard1988} J. Hounsgaard, H. Hultborn, B. Jespersen, and O. Kiehn, \emph{Bistability of alpha-motoneurones in the decerebrate cat and in the acute spinal cat after intravenous 5-hydroxytryptophan}, J. Physiol. \textbf{405}, 345 (1988).
\bibitem{Izhikevich2007} E. M. Izhikevich, \emph{Dynamical Systems in Neuroscience: The Geometry of Excitability and Bursting} (MIT Press, Cambridge, MA, 2007).
\bibitem{AgmonSnir1998} H. Agmon-Snir, C. E. Carr, and J. Rinzel, \emph{The role of dendrites in auditory coincidence detection}, Nature \textbf{393}, 268 (1998).
\bibitem{BorstSoria2012} J. G. G. Borst and J. Soria van Hoeve, \emph{The calyx of Held synapse: from model synapse to auditory relay}, Annu. Rev. Physiol. \textbf{74}, 199 (2012).
\bibitem{Euler2002} T. Euler, P. B. Detwiler, and W. Denk, \emph{Directionally selective calcium signals in dendrites of starburst amacrine cells}, Nature \textbf{418}, 845 (2002).
\bibitem{AstonJonesBloom1981} G. Aston-Jones and F. E. Bloom, \emph{Activity of norepinephrine-containing locus coeruleus neurons in behaving rats anticipates fluctuations in the sleep-waking cycle}, J. Neurosci. \textbf{1}, 876 (1981).
\bibitem{BrooksPeever2012} P. L. Brooks and J. H. Peever, \emph{Identification of the transmitter and receptor mechanisms responsible for REM sleep paralysis}, J. Neurosci. \textbf{32}, 9785 (2012).
\bibitem{Borbely1982} A. A. Borb\'ely, \emph{A two process model of sleep regulation}, Hum. Neurobiol. \textbf{1}, 195 (1982).
\bibitem{PorkkaHeiskanen1997} T. Porkka-Heiskanen et al., \emph{Adenosine: a mediator of the sleep-inducing effects of prolonged wakefulness}, Science \textbf{276}, 1265 (1997).
\bibitem{Steriade1993} M. Steriade, A. N\'u\~nez, and F. Amzica, \emph{A novel slow ($<1$ Hz) oscillation of neocortical neurons in vivo: depolarizing and hyperpolarizing components}, J. Neurosci. \textbf{13}, 3252 (1993).
\bibitem{McCormick1992} D. A. McCormick, \emph{Neurotransmitter actions in the thalamus and cerebral cortex and their role in neuromodulation of thalamocortical activity}, Prog. Neurobiol. \textbf{39}, 337 (1992).
\bibitem{SteriadeMcCormickSejnowski1993} M. Steriade, D. A. McCormick, and T. J. Sejnowski, \emph{Thalamocortical oscillations in the sleeping and aroused brain}, Science \textbf{262}, 679 (1993).
\bibitem{Nadasdy1999} Z. N\'adasdy et al., \emph{Replay and time compression of recurring spike sequences in the hippocampus}, J. Neurosci. \textbf{19}, 9497 (1999).
\bibitem{Maingret2016} N. Maingret et al., \emph{Hippocampo-cortical coupling mediates memory consolidation during sleep}, Nat. Neurosci. \textbf{19}, 959 (2016).
\bibitem{McClelland1995} J. L. McClelland, B. L. McNaughton, and R. C. O'Reilly, \emph{Why there are complementary learning systems in the hippocampus and neocortex}, Psychol. Rev. \textbf{102}, 419 (1995).
\bibitem{TononiCirelli2014} G. Tononi and C. Cirelli, \emph{Sleep and the price of plasticity: from synaptic and cellular homeostasis to memory consolidation and integration}, Neuron \textbf{81}, 12 (2014).
\bibitem{deVivo2017} L. de Vivo et al., \emph{Ultrastructural evidence for synaptic scaling across the wake/sleep cycle}, Science \textbf{355}, 507 (2017).
\bibitem{Xie2013} L. Xie et al., \emph{Sleep drives metabolite clearance from the adult brain}, Science \textbf{342}, 373 (2013).
\bibitem{Miao2024} A. Miao et al., \emph{Brain clearance is reduced during sleep and anesthesia}, Nat. Neurosci. \textbf{27}, 1046 (2024).
\bibitem{Vyazovskiy2011} V. V. Vyazovskiy et al., \emph{Local sleep in awake rats}, Nature \textbf{472}, 443 (2011).
\bibitem{Lancaster2013} M. A. Lancaster et al., \emph{Cerebral organoids model human brain development and microcephaly}, Nature \textbf{501}, 373 (2013).
\bibitem{Jordan2024} F. D. Jordan et al., \emph{Open and remotely accessible Neuroplatform for research in wetware computing}, Front. Artif. Intell. \textbf{7}, 1376042 (2024).
\bibitem{Smirnova2023} L. Smirnova et al., \emph{Organoid intelligence (OI): the new frontier in biocomputing and intelligence-in-a-dish}, Front. Sci. \textbf{1}, 1017235 (2023).
\bibitem{Wagenaar2006} D. A. Wagenaar, J. Pine, and S. M. Potter, \emph{An extremely rich repertoire of bursting patterns during the development of cortical cultures}, BMC Neurosci. \textbf{7}, 11 (2006).
\bibitem{Shahaf2001} G. Shahaf and S. Marom, \emph{Learning in networks of cortical neurons}, J. Neurosci. \textbf{21}, 8782 (2001).
\bibitem{Kagan2022} B. J. Kagan et al., \emph{In vitro neurons learn and exhibit sentience when embodied in a simulated game-world}, Neuron \textbf{110}, 3952 (2022).
\bibitem{Friston2010} K. Friston, \emph{The free-energy principle: a unified brain theory?}, Nat. Rev. Neurosci. \textbf{11}, 127 (2010).
\bibitem{Balci2023} F. Balci et al., \emph{A response to claims of emergent intelligence and sentience in a dish}, Neuron \textbf{111}, 604 (2023).
\bibitem{Bakkum2008} D. J. Bakkum, Z. C. Chao, and S. M. Potter, \emph{Spatio-temporal electrical stimuli shape behavior of an embodied cortical network in a goal-directed learning task}, J. Neural Eng. \textbf{5}, 310 (2008).
\bibitem{Maass2002} W. Maass, T. Natschl\"ager, and H. Markram, \emph{Real-time computing without stable states: a new framework for neural computation based on perturbations}, Neural Comput. \textbf{14}, 2531 (2002).
\bibitem{JaegerHaas2004} H. Jaeger and H. Haas, \emph{Harnessing nonlinearity: predicting chaotic systems and saving energy in wireless communication}, Science \textbf{304}, 78 (2004).
\bibitem{Cai2023} H. Cai et al., \emph{Brain organoid reservoir computing for artificial intelligence}, Nat. Electron. \textbf{6}, 1032 (2023).
\bibitem{Habibollahi2023} F. Habibollahi, B. J. Kagan, A. N. Burkitt, and C. French, \emph{Critical dynamics arise during structured information presentation within embodied in vitro neuronal networks}, Nat. Commun. \textbf{14}, 5287 (2023).
\bibitem{Robbins2026} A. Robbins et al., \emph{Goal-directed learning in cortical organoids}, Cell Rep. \textbf{45}, 116984 (2026).
\bibitem{Keene2020} S. T. Keene et al., \emph{A biohybrid synapse with neurotransmitter-mediated plasticity}, Nat. Mater. \textbf{19}, 969 (2020).
\bibitem{Jo2016} J. Jo et al., \emph{Midbrain-like organoids from human pluripotent stem cells contain functional dopaminergic and neuromelanin-producing neurons}, Cell Stem Cell \textbf{19}, 248 (2016).
\bibitem{Hendry1987} S. H. Hendry, H. D. Schwark, E. G. Jones, and J. Yan, \emph{Numbers and proportions of GABA-immunoreactive neurons in different areas of monkey cerebral cortex}, J. Neurosci. \textbf{7}, 1503 (1987).
\bibitem{Tang2001synapses} Y. Tang, J. R. Nyengaard, D. M. G. De Groot, and H. J. G. Gundersen, \emph{Total regional and global number of synapses in the human brain neocortex}, Synapse \textbf{41}, 258 (2001).
\bibitem{Pakkenberg1997} B. Pakkenberg and H. J. G. Gundersen, \emph{Neocortical neuron number in humans: effect of sex and age}, J. Comp. Neurol. \textbf{384}, 312 (1997).
\bibitem{Matsuda2009} W. Matsuda, T. Furuta, K. C. Nakamura, H. Hioki, F. Fujiyama, R. Arai, and T. Kaneko, \emph{Single nigrostriatal dopaminergic neurons form widely spread and highly dense axonal arborizations in the neostriatum}, J. Neurosci. \textbf{29}, 444 (2009).
\bibitem{Pakkenberg1991} B. Pakkenberg, A. M{\o}ller, H. J. G. Gundersen, A. Mouritzen Dam, and H. Pakkenberg, \emph{The absolute number of nerve cells in substantia nigra in normal subjects and in patients with Parkinson's disease estimated with an unbiased stereological method}, J. Neurol. Neurosurg. Psychiatry \textbf{54}, 30 (1991).
\bibitem{Baker1990} K. G. Baker, G. M. Halliday, and I. T\"ork, \emph{Cytoarchitecture of the human dorsal raphe nucleus}, J. Comp. Neurol. \textbf{301}, 147 (1990).
\bibitem{Baker1991} K. G. Baker, G. M. Halliday, P. Halasz, J.-P. Hornung, L. B. Geffen, R. G. H. Cotton, and I. T\"ork, \emph{Cytoarchitecture of serotonin-synthesizing neurons in the pontine tegmentum of the human brain}, Synapse \textbf{7}, 301 (1991).
\bibitem{Airaksinen1991} M. S. Airaksinen, A. Paetau, L. Palj\"arvi, K. Reinikainen, P. Riekkinen, R. Suomalainen, and P. Panula, \emph{Histamine neurons in human hypothalamus: anatomy in normal and Alzheimer diseased brains}, Neuroscience \textbf{44}, 465 (1991).
\bibitem{Colquhoun1992} D. Colquhoun, P. Jonas, and B. Sakmann, \emph{Action of brief pulses of glutamate on AMPA/kainate receptors in patches from different neurones of rat hippocampal slices}, J. Physiol. \textbf{458}, 261 (1992).
\bibitem{DutarNicoll1988} P. Dutar and R. A. Nicoll, \emph{A physiological role for GABA$_B$ receptors in the central nervous system}, Nature \textbf{332}, 156 (1988).

\bibitem{Rauch2003} A. Rauch, G. La Camera, H.-R. L\"uscher, W. Senn, and S. Fusi, \emph{Neocortical pyramidal cells respond as integrate-and-fire neurons to in vivo-like input currents}, J. Neurophysiol. \textbf{90}, 1598 (2003).
\bibitem{KlumppEady1956} R. G. Klumpp and H. R. Eady, \emph{Some measurements of interaural time difference thresholds}, J. Acoust. Soc. Am. \textbf{28}, 859 (1956).
\bibitem{Brand2002} A. Brand, O. Behrend, T. Marquardt, D. McAlpine, and B. Grothe, \emph{Precise inhibition is essential for microsecond interaural time difference coding}, Nature \textbf{417}, 543 (2002).
\bibitem{Funahashi1989} S. Funahashi, C. J. Bruce, and P. S. Goldman-Rakic, \emph{Mnemonic coding of visual space in the monkey's dorsolateral prefrontal cortex}, J. Neurophysiol. \textbf{61}, 331 (1989).
\bibitem{WangMarkram2006} Y. Wang, H. Markram, P. H. Goodman, T. K. Berger, J. Ma, and P. S. Goldman-Rakic, \emph{Heterogeneity in the pyramidal network of the medial prefrontal cortex}, Nat. Neurosci. \textbf{9}, 534 (2006).

\bibitem{McCullochPitts1943} W. S. McCulloch and W. Pitts, \emph{A logical calculus of the ideas immanent in nervous activity}, Bull. Math. Biophys. \textbf{5}, 115 (1943).
\bibitem{Fried2002} S. I. Fried, T. A. M\"unch, and F. S. Werblin, \emph{Mechanisms and circuitry underlying directional selectivity in the retina}, Nature \textbf{420}, 411 (2002).

\bibitem{Hromadka2008} T. Hrom\'adka, M. R. DeWeese, and A. M. Zador, \emph{Sparse representation of sounds in the unanesthetized auditory cortex}, PLoS Biol. \textbf{6}, e16 (2008).
\bibitem{Abbott1997depression} L. F. Abbott, J. A. Varela, K. Sen, and S. B. Nelson, \emph{Synaptic depression and cortical gain control}, Science \textbf{275}, 221 (1997).
\bibitem{ShadlenNewsome1998} M. N. Shadlen and W. T. Newsome, \emph{The variable discharge of cortical neurons: implications for connectivity, computation, and information coding}, J. Neurosci. \textbf{18}, 3870 (1998).

\bibitem{TrulsonJacobs1979} M. E. Trulson and B. L. Jacobs, \emph{Raphe unit activity in freely moving cats: correlation with level of behavioral arousal}, Brain Res. \textbf{163}, 135 (1979).
\bibitem{GagnonParent2014} D. Gagnon and M. Parent, \emph{Distribution of VGLUT3 in highly collateralized axons from the rat dorsal raphe nucleus as revealed by single-neuron reconstructions}, PLoS One \textbf{9}, e87709 (2014).
\bibitem{VandermaelenAghajanian1983} C. P. Vandermaelen and G. K. Aghajanian, \emph{Electrophysiological and pharmacological characterization of serotonergic dorsal raphe neurons recorded extracellularly and intracellularly in rat brain slices}, Brain Res. \textbf{289}, 109 (1983).
\bibitem{CourtneyFord2016} N. A. Courtney and C. P. Ford, \emph{Mechanisms of 5-HT$_{1A}$ receptor-mediated transmission in dorsal raphe serotonin neurons}, J. Physiol. \textbf{594}, 953 (2016).
\bibitem{Hashemi2012serotonin} P. Hashemi et al., \emph{Brain dopamine and serotonin differ in regulation and its consequences}, Proc. Natl. Acad. Sci. USA \textbf{109}, 11510 (2012).
\bibitem{Black1934feedback} H. S. Black, \emph{Stabilized feedback amplifiers}, Bell Syst. Tech. J. \textbf{13}, 1 (1934).
\bibitem{KobayashiSchultz2008} S. Kobayashi and W. Schultz, \emph{Influence of reward delays on responses of dopamine neurons}, J. Neurosci. \textbf{28}, 7837 (2008).
\bibitem{Schweighofer2008} N. Schweighofer et al., \emph{Low-serotonin levels increase delayed reward discounting in humans}, J. Neurosci. \textbf{28}, 4528 (2008).

\bibitem{Malenka1988calcium} R. C. Malenka, J. A. Kauer, R. S. Zucker, and R. A. Nicoll, \emph{Postsynaptic calcium is sufficient for potentiation of hippocampal synaptic transmission}, Science \textbf{242}, 81 (1988).
\bibitem{Bliss1973LTP} T. V. P. Bliss and T. L{\o}mo, \emph{Long-lasting potentiation of synaptic transmission in the dentate area of the anaesthetized rabbit following stimulation of the perforant path}, J. Physiol. \textbf{232}, 331 (1973).
\bibitem{Kelso1986hebb} S. R. Kelso, A. H. Ganong, and T. H. Brown, \emph{Hebbian synapses in hippocampus}, Proc. Natl. Acad. Sci. USA \textbf{83}, 5326 (1986).
\bibitem{He2015eligibility} K. He et al., \emph{Distinct eligibility traces for LTP and LTD in cortical synapses}, Neuron \textbf{88}, 528 (2015).
\bibitem{Olveczky2005} B. P. \"Olveczky, A. S. Andalman, and M. S. Fee, \emph{Vocal experimentation in the juvenile songbird requires a basal ganglia circuit}, PLoS Biol. \textbf{3}, e153 (2005).
\bibitem{TumerBrainard2007} E. C. Tumer and M. S. Brainard, \emph{Performance variability enables adaptive plasticity of `crystallized' adult birdsong}, Nature \textbf{450}, 1240 (2007).
\bibitem{Eshel2015arithmetic} N. Eshel et al., \emph{Arithmetic and local circuitry underlying dopamine prediction errors}, Nature \textbf{525}, 243 (2015).
\bibitem{GraceOnn1989} A. A. Grace and S.-P. Onn, \emph{Morphology and electrophysiological properties of immunocytochemically identified rat dopamine neurons recorded in vitro}, J. Neurosci. \textbf{9}, 3463 (1989).
\bibitem{ODoherty2004actor} J. O'Doherty et al., \emph{Dissociable roles of ventral and dorsal striatum in instrumental conditioning}, Science \textbf{304}, 452 (2004).

\bibitem{NeweyPowell1987} W. K. Newey and J. L. Powell, \emph{Asymmetric least squares estimation and testing}, Econometrica \textbf{55}, 819 (1987).
\bibitem{Beierholm2013vigor} U. Beierholm et al., \emph{Dopamine modulates reward-related vigor}, Neuropsychopharmacology \textbf{38}, 1495 (2013).
\bibitem{Howe2013ramp} M. W. Howe, P. L. Tierney, S. G. Sandberg, P. E. M. Phillips, and A. M. Graybiel, \emph{Prolonged dopamine signalling in striatum signals proximity and value of distant rewards}, Nature \textbf{500}, 575 (2013).

\bibitem{Baker1989LC} K. G. Baker, I. T\"ork, J.-P. Hornung, and P. Halasz, \emph{The human locus coeruleus complex: an immunohistochemical and three dimensional reconstruction study}, Exp. Brain Res. \textbf{77}, 257 (1989).
\bibitem{Schwarz2015LC} L. A. Schwarz et al., \emph{Viral-genetic tracing of the input--output organization of a central noradrenaline circuit}, Nature \textbf{524}, 88 (2015).
\bibitem{Uematsu2017LC} A. Uematsu et al., \emph{Modular organization of the brainstem noradrenaline system coordinates opposing learning states}, Nat. Neurosci. \textbf{20}, 1602 (2017).
\bibitem{AstonJones1994vigilance} G. Aston-Jones, J. Rajkowski, P. Kubiak, and T. Alexinsky, \emph{Locus coeruleus neurons in monkey are selectively activated by attended cues in a vigilance task}, J. Neurosci. \textbf{14}, 4467 (1994).
\bibitem{Clayton2004LC} E. C. Clayton, J. Rajkowski, J. D. Cohen, and G. Aston-Jones, \emph{Phasic activation of monkey locus ceruleus neurons by simple decisions in a forced-choice task}, J. Neurosci. \textbf{24}, 9914 (2004).
\bibitem{Foote1975NE} S. L. Foote, R. Freedman, and A. P. Oliver, \emph{Effects of putative neurotransmitters on neuronal activity in monkey auditory cortex}, Brain Res. \textbf{86}, 229 (1975).
\bibitem{JepmaNieuwenhuis2011} M. Jepma and S. Nieuwenhuis, \emph{Pupil diameter predicts changes in the exploration--exploitation trade-off: evidence for the adaptive gain theory}, J. Cogn. Neurosci. \textbf{23}, 1587 (2011).
\bibitem{Tervo2014stochastic} D. G. R. Tervo et al., \emph{Behavioral variability through stochastic choice and its gating by anterior cingulate cortex}, Cell \textbf{159}, 21 (2014).
\bibitem{Usher1999LC} M. Usher, J. D. Cohen, D. Servan-Schreiber, J. Rajkowski, and G. Aston-Jones, \emph{The role of locus coeruleus in the regulation of cognitive performance}, Science \textbf{283}, 549 (1999).

\bibitem{Mesulam1983} M.-M. Mesulam, E. J. Mufson, B. H. Wainer, and A. I. Levey, \emph{Central cholinergic pathways in the rat: an overview based on an alternative nomenclature (Ch1--Ch6)}, Neuroscience \textbf{10}, 1185 (1983).
\bibitem{Marrosu1995} F. Marrosu et al., \emph{Microdialysis measurement of cortical and hippocampal acetylcholine release during sleep-wake cycle in freely moving cats}, Brain Res. \textbf{671}, 329 (1995).
\bibitem{Picciotto2012} M. R. Picciotto, M. J. Higley, and Y. S. Mineur, \emph{Acetylcholine as a neuromodulator: cholinergic signaling shapes nervous system function and behavior}, Neuron \textbf{76}, 116 (2012).
\bibitem{HasselmoSchnell1994} M. E. Hasselmo and E. Schnell, \emph{Laminar selectivity of the cholinergic suppression of synaptic transmission in rat hippocampal region CA1: computational modeling and brain slice physiology}, J. Neurosci. \textbf{14}, 3898 (1994).
\bibitem{Gil1997} Z. Gil, B. W. Connors, and Y. Amitai, \emph{Differential regulation of neocortical synapses by neuromodulators and activity}, Neuron \textbf{19}, 679 (1997).
\bibitem{KilgardMerzenich1998} M. P. Kilgard and M. M. Merzenich, \emph{Cortical map reorganization enabled by nucleus basalis activity}, Science \textbf{279}, 1714 (1998).
\bibitem{Atri2004} A. Atri et al., \emph{Blockade of central cholinergic receptors impairs new learning and increases proactive interference in a word paired-associate memory task}, Behav. Neurosci. \textbf{118}, 223 (2004).
\bibitem{Girardeau2009} G. Girardeau, K. Benchenane, S. I. Wiener, G. Buzs\'aki, and M. B. Zugaro, \emph{Selective suppression of hippocampal ripples impairs spatial memory}, Nat. Neurosci. \textbf{12}, 1222 (2009).

\input{supplement/bib_10}
\bibitem{CopenhagenJahr1989} D. R. Copenhagen and C. E. Jahr, \emph{Release of endogenous excitatory amino acids from turtle photoreceptors}, Nature \textbf{341}, 536 (1989).
\bibitem{GlowatzkiFuchs2002} E. Glowatzki and P. A. Fuchs, \emph{Transmitter release at the hair cell ribbon synapse}, Nat. Neurosci. \textbf{5}, 147 (2002).
\bibitem{Tinsley2016} J. N. Tinsley et al., \emph{Direct detection of a single photon by humans}, Nat. Commun. \textbf{7}, 12172 (2016).
\bibitem{Sellick1982} P. M. Sellick, R. Patuzzi, and B. M. Johnstone, \emph{Measurement of basilar membrane motion in the guinea pig using the M\"ossbauer technique}, J. Acoust. Soc. Am. \textbf{72}, 131 (1982).
\bibitem{Ruggero1997} M. A. Ruggero, N. C. Rich, A. Recio, S. S. Narayan, and L. Robles, \emph{Basilar-membrane responses to tones at the base of the chinchilla cochlea}, J. Acoust. Soc. Am. \textbf{101}, 2151 (1997).
\bibitem{NormannPerlman1979} R. A. Normann and I. Perlman, \emph{The effects of background illumination on the photoresponses of red and green cones}, J. Physiol. \textbf{286}, 491 (1979).
\bibitem{Laughlin1981} S. Laughlin, \emph{A simple coding procedure enhances a neuron's information capacity}, Z. Naturforsch. C \textbf{36}, 910 (1981).
\bibitem{CurcioAllen1990} C. A. Curcio and K. A. Allen, \emph{Topography of ganglion cells in human retina}, J. Comp. Neurol. \textbf{300}, 5 (1990).
\bibitem{Greenwood1990} D. D. Greenwood, \emph{A cochlear frequency-position function for several species --- 29 years later}, J. Acoust. Soc. Am. \textbf{87}, 2592 (1990).
\bibitem{Eguiluz2000} V. M. Egu\'iluz, M. Ospeck, Y. Choe, A. J. Hudspeth, and M. O. Magnasco, \emph{Essential nonlinearities in hearing}, Phys. Rev. Lett. \textbf{84}, 5232 (2000).
\bibitem{PalmerRussell1986} A. R. Palmer and I. J. Russell, \emph{Phase-locking in the cochlear nerve of the guinea-pig and its relation to the receptor potential of inner hair-cells}, Hear. Res. \textbf{24}, 1 (1986).
\bibitem{Malnic1999} B. Malnic, J. Hirono, T. Sato, and L. B. Buck, \emph{Combinatorial receptor codes for odors}, Cell \textbf{96}, 713 (1999).
\bibitem{Finger2005} T. E. Finger et al., \emph{ATP signaling is crucial for communication from taste buds to gustatory nerves}, Science \textbf{310}, 1495 (2005).
\bibitem{Huang2005} Y.-J. Huang et al., \emph{Mouse taste buds use serotonin as a neurotransmitter}, J. Neurosci. \textbf{25}, 843 (2005).
\bibitem{JohanssonVallbo1979} R. S. Johansson and A. B. Vallbo, \emph{Tactile sensibility in the human hand: relative and absolute densities of four types of mechanoreceptive units in glabrous skin}, J. Physiol. \textbf{286}, 283 (1979).
\bibitem{McKemy2002} D. D. McKemy, W. M. Neuhausser, and D. Julius, \emph{Identification of a cold receptor reveals a general role for TRP channels in thermosensation}, Nature \textbf{416}, 52 (2002).

\bibitem{Schmolesky1998} M. T. Schmolesky et al., \emph{Signal timing across the macaque visual system}, J. Neurophysiol. \textbf{79}, 3272 (1998).
\bibitem{Lampl2004} I. Lampl, D. Ferster, T. Poggio, and M. Riesenhuber, \emph{Intracellular measurements of spatial integration and the MAX operation in complex cells of the cat primary visual cortex}, J. Neurophysiol. \textbf{92}, 2704 (2004).
\bibitem{KobatakeTanaka1994} E. Kobatake and K. Tanaka, \emph{Neuronal selectivities to complex object features in the ventral visual pathway of the macaque cerebral cortex}, J. Neurophysiol. \textbf{71}, 856 (1994).
\bibitem{Albright1984} T. D. Albright, \emph{Direction and orientation selectivity of neurons in visual area MT of the macaque}, J. Neurophysiol. \textbf{52}, 1106 (1984).
\bibitem{Goodfellow2015} I. J. Goodfellow, J. Shlens, and C. Szegedy, \emph{Explaining and harnessing adversarial examples}, in \emph{Proc. ICLR} (2015), arXiv:1412.6572.
\bibitem{Elsayed2018} G. F. Elsayed et al., \emph{Adversarial examples that fool both computer vision and time-limited humans}, in \emph{Advances in Neural Information Processing Systems 31} (2018), arXiv:1802.08195.
\bibitem{Thompson1982} P. Thompson, \emph{Perceived rate of movement depends on contrast}, Vision Res. \textbf{22}, 377 (1982).
\bibitem{BarlowHill1963} H. B. Barlow and R. M. Hill, \emph{Evidence for a physiological explanation of the waterfall phenomenon and figural after-effects}, Nature \textbf{200}, 1345 (1963).
\bibitem{EaglemanSejnowski2000} D. M. Eagleman and T. J. Sejnowski, \emph{Motion integration and postdiction in visual awareness}, Science \textbf{287}, 2036 (2000).

\bibitem{Feinstein1955mu} B. Feinstein, B. Lindeg{\aa}rd, E. Nyman, and G. Wohlfart, \emph{Morphologic studies of motor units in normal human muscles}, Acta Anat. \textbf{23}, 127 (1955).
\bibitem{MilnerBrown1973rec} H. S. Milner-Brown, R. B. Stein, and R. Yemm, \emph{The orderly recruitment of human motor units during voluntary isometric contractions}, J. Physiol. \textbf{230}, 359 (1973).
\bibitem{Jones2002sdn} K. E. Jones, A. F. de C. Hamilton, and D. M. Wolpert, \emph{Sources of signal-dependent noise during isometric force production}, J. Neurophysiol. \textbf{88}, 1533 (2002).
\bibitem{GomiOsu1998stiff} H. Gomi and R. Osu, \emph{Task-dependent viscoelasticity of human multijoint arm and its spatial characteristics for interaction with environments}, J. Neurosci. \textbf{18}, 8965 (1998).
\bibitem{Jankowska1972Ia} E. Jankowska and W. J. Roberts, \emph{Synaptic actions of single interneurones mediating reciprocal Ia inhibition of motoneurones}, J. Physiol. \textbf{222}, 623 (1972).
\bibitem{Pruszynski2008rapid} J. A. Pruszynski, I. Kurtzer, and S. H. Scott, \emph{Rapid motor responses are appropriately tuned to the metrics of a visuospatial task}, J. Neurophysiol. \textbf{100}, 224 (2008).
\bibitem{SaundersKnill2003vis} J. A. Saunders and D. C. Knill, \emph{Humans use continuous visual feedback from the hand to control fast reaching movements}, Exp. Brain Res. \textbf{152}, 341 (2003).
\bibitem{FlanaganWing1997grip} J. R. Flanagan and A. M. Wing, \emph{The role of internal models in motion planning and control: evidence from grip force adjustments during movements of hand-held loads}, J. Neurosci. \textbf{17}, 1519 (1997).
\bibitem{Matsuoka1985osc} K. Matsuoka, \emph{Sustained oscillations generated by mutually inhibiting neurons with adaptation}, Biol. Cybern. \textbf{52}, 367 (1985).

\bibitem{Treede1995heat} R.-D. Treede, R. A. Meyer, S. N. Raja, and J. N. Campbell, \emph{Evidence for two different heat transduction mechanisms in nociceptive primary afferents innervating monkey skin}, J. Physiol. \textbf{483}, 747 (1995).
\bibitem{Weidner1999cfib} C. Weidner et al., \emph{Functional attributes discriminating mechano-insensitive and mechano-responsive C nociceptors in human skin}, J. Neurosci. \textbf{19}, 10184 (1999).
\bibitem{CampbellLaMotte1983first} J. N. Campbell and R. H. LaMotte, \emph{Latency to detection of first pain}, Brain Res. \textbf{266}, 203 (1983).
\bibitem{LaMotte1982hyper} R. H. LaMotte, J. G. Thalhammer, H. E. Torebj{\"o}rk, and C. J. Robinson, \emph{Peripheral neural mechanisms of cutaneous hyperalgesia following mild injury by heat}, J. Neurosci. \textbf{2}, 765 (1982).
\bibitem{Mancini2014touch} F. Mancini, T. Nash, G. D. Iannetti, and P. Haggard, \emph{Pain relief by touch: a quantitative approach}, Pain \textbf{155}, 635 (2014).
\bibitem{Fields1983rvm} H. L. Fields, J. Bry, I. Hentall, and G. Zorman, \emph{The activity of neurons in the rostral medulla of the rat during withdrawal from noxious heat}, J. Neurosci. \textbf{3}, 2545 (1983).
\bibitem{Crombez1996disrupt} G. Crombez, C. Eccleston, F. Baeyens, and P. Eelen, \emph{The disruptive nature of pain: an experimental investigation}, Behav. Res. Ther. \textbf{34}, 911 (1996).
\bibitem{Schouenborg1995wr} J. Schouenborg, H.-R. Weng, J. Kalliom{\"a}ki, and H. Holmberg, \emph{A survey of spinal dorsal horn neurones encoding the spatial organization of withdrawal reflexes in the rat}, Exp. Brain Res. \textbf{106}, 19 (1995).

\bibitem{Andersen1992cereb} B. B. Andersen, L. Korbo, and B. Pakkenberg, \emph{A quantitative study of the human cerebellum with unbiased stereological techniques}, J. Comp. Neurol. \textbf{326}, 549 (1992).
\bibitem{ShadmehrMussaIvaldi1994} R. Shadmehr and F. A. Mussa-Ivaldi, \emph{Adaptive representation of dynamics during learning of a motor task}, J. Neurosci. \textbf{14}, 3208 (1994).

\bibitem{JoseCollison1970ihr} A. D. Jose and D. Collison, \emph{The normal range and determinants of the intrinsic heart rate in man}, Cardiovasc. Res. \textbf{4}, 160 (1970).
\bibitem{Logothetis1987gbg} D. E. Logothetis, Y. Kurachi, J. Galper, E. J. Neer, and D. E. Clapham, \emph{The $\beta\gamma$ subunits of GTP-binding proteins activate the muscarinic K$^+$ channel in heart}, Nature \textbf{325}, 321 (1987).
\bibitem{KellerWood1984fb} M. E. Keller-Wood and M. F. Dallman, \emph{Corticosteroid inhibition of ACTH secretion}, Endocr. Rev. \textbf{5}, 1 (1984).
\bibitem{Walker2012glc} J. J. Walker et al., \emph{The origin of glucocorticoid hormone oscillations}, PLoS Biol. \textbf{10}, e1001341 (2012).
\bibitem{Wildt1981gnrh} L. Wildt et al., \emph{Frequency and amplitude of gonadotropin-releasing hormone stimulation and gonadotropin secretion in the rhesus monkey}, Endocrinology \textbf{109}, 376 (1981).
\bibitem{Welsh2010scn} D. K. Welsh, J. S. Takahashi, and S. A. Kay, \emph{Suprachiasmatic nucleus: cell autonomy and network properties}, Annu. Rev. Physiol. \textbf{72}, 551 (2010).
\bibitem{Khalsa2003prc} S. B. S. Khalsa, M. E. Jewett, C. Cajochen, and C. A. Czeisler, \emph{A phase response curve to single bright light pulses in human subjects}, J. Physiol. \textbf{549}, 945 (2003).
\bibitem{Sack1992blind} R. L. Sack, A. J. Lewy, M. L. Blood, L. D. Keith, and H. Nakagawa, \emph{Circadian rhythm abnormalities in totally blind people: incidence and clinical significance}, J. Clin. Endocrinol. Metab. \textbf{75}, 127 (1992).

\bibitem{Henze2000extra} D. A. Henze et al., \emph{Intracellular features predicted by extracellular recordings in the hippocampus in vivo}, J. Neurophysiol. \textbf{84}, 390 (2000).
\bibitem{Histed2009stim} M. H. Histed, V. Bonin, and R. C. Reid, \emph{Direct activation of sparse, distributed populations of cortical neurons by electrical microstimulation}, Neuron \textbf{63}, 508 (2009).
\bibitem{Orsborn2014coad} A. L. Orsborn et al., \emph{Closed-loop decoder adaptation shapes neural plasticity for skillful neuroprosthetic control}, Neuron \textbf{82}, 1380 (2014).

\bibitem{HuVervaekeStorm2002} H. Hu, K. Vervaeke, and J. F. Storm, \emph{Two forms of electrical resonance at theta frequencies, generated by M-current, h-current and persistent Na$^+$ current in rat hippocampal pyramidal cells}, J. Physiol. \textbf{545}, 783 (2002).
\bibitem{Mauro1970} A. Mauro, F. Conti, F. Dodge, and R. Schor, \emph{Subthreshold behavior and phenomenological impedance of the squid giant axon}, J. Gen. Physiol. \textbf{55}, 497 (1970).
\bibitem{Aksay2001} E. Aksay et al., \emph{In vivo intracellular recording and perturbation of persistent activity in a neural integrator}, Nat. Neurosci. \textbf{4}, 184 (2001).
\bibitem{Major2004} G. Major et al., \emph{Plasticity and tuning of the time course of analog persistent firing in a neural integrator}, Proc. Natl. Acad. Sci. USA \textbf{101}, 7745 (2004).
\bibitem{Tateno2004} T. Tateno, A. Harsch, and H. P. C. Robinson, \emph{Threshold firing frequency-current relationships of neurons in rat somatosensory cortex: type 1 and type 2 dynamics}, J. Neurophysiol. \textbf{92}, 2283 (2004).
\bibitem{Marder2012} E. Marder, \emph{Neuromodulation of neuronal circuits: back to the future}, Neuron \textbf{76}, 1 (2012).
\bibitem{McGintyHarper1976} D. J. McGinty and R. M. Harper, \emph{Dorsal raphe neurons: depression of firing during sleep in cats}, Brain Res. \textbf{101}, 569 (1976).

\bibitem{Gentet2000} L. J. Gentet, G. J. Stuart, and J. D. Clements, \emph{Direct measurement of specific membrane capacitance in neurons}, Biophys. J. \textbf{79}, 314 (2000).
\bibitem{Borst1995} J. G. G. Borst, F. Helmchen, and B. Sakmann, \emph{Pre- and postsynaptic whole-cell recordings in the medial nucleus of the trapezoid body of the rat}, J. Physiol. \textbf{489}, 825 (1995).
\bibitem{AmirDevor2003} R. Amir and M. Devor, \emph{Electrical excitability of the soma of sensory neurons is required for spike invasion of the soma, but not for through-conduction}, Biophys. J. \textbf{84}, 2181 (2003).
\bibitem{Chadderton2004} P. Chadderton, T. W. Margrie, and M. H\"ausser, \emph{Integration of quanta in cerebellar granule cells during sensory processing}, Nature \textbf{428}, 856 (2004).
\bibitem{LitwinKumar2017} A. Litwin-Kumar et al., \emph{Optimal degrees of synaptic connectivity}, Neuron \textbf{93}, 1153 (2017).
\bibitem{NapperHarvey1988} R. M. A. Napper and R. J. Harvey, \emph{Number of parallel fiber synapses on an individual Purkinje cell in the cerebellum of the rat}, J. Comp. Neurol. \textbf{274}, 168 (1988).
\bibitem{Megias2001} M. Meg\'{\i}as, Z. Emri, T. F. Freund, and A. I. Guly\'as, \emph{Total number and distribution of inhibitory and excitatory synapses on hippocampal CA1 pyramidal cells}, Neuroscience \textbf{102}, 527 (2001).
\bibitem{Polsky2004} A. Polsky, B. W. Mel, and J. Schiller, \emph{Computational subunits in thin dendrites of pyramidal cells}, Nat. Neurosci. \textbf{7}, 621 (2004).
\bibitem{Brunel2004} N. Brunel et al., \emph{Optimal information storage and the distribution of synaptic weights: perceptron versus Purkinje cell}, Neuron \textbf{43}, 745 (2004).

\bibitem{Vyazovskiy2009} V. V. Vyazovskiy et al., \emph{Cortical firing and sleep homeostasis}, Neuron \textbf{63}, 865 (2009).
\bibitem{Daan1984} S. Daan, D. G. Beersma, and A. A. Borb\'ely, \emph{Timing of human sleep: recovery process gated by a circadian pacemaker}, Am. J. Physiol. \textbf{246}, R161 (1984).
\bibitem{Borbely1981} A. A. Borb\'ely et al., \emph{Sleep deprivation: effect on sleep stages and EEG power density in man}, Electroencephalogr. Clin. Neurophysiol. \textbf{51}, 483 (1981).
\bibitem{SteriadeAmzica1993} M. Steriade, F. Amzica, and A. Nu\~nez, \emph{Cholinergic and noradrenergic modulation of the slow ($\approx0.3$ Hz) oscillation in neocortical cells}, J. Neurophysiol. \textbf{70}, 1385 (1993).
\bibitem{vonKrosigk1993} M. von Krosigk, T. Bal, and D. A. McCormick, \emph{Cellular mechanisms of a synchronized oscillation in the thalamus}, Science \textbf{261}, 361 (1993).
\bibitem{LeeWilson2002} A. K. Lee and M. A. Wilson, \emph{Memory of sequential experience in the hippocampus during slow wave sleep}, Neuron \textbf{36}, 1183 (2002).
\bibitem{Euston2007} D. R. Euston, M. Tatsuno, and B. L. McNaughton, \emph{Fast-forward playback of recent memory sequences in prefrontal cortex during sleep}, Science \textbf{318}, 1147 (2007).
\bibitem{Hobson2009} J. A. Hobson, \emph{REM sleep and dreaming: towards a theory of protoconsciousness}, Nat. Rev. Neurosci. \textbf{10}, 803 (2009).

\bibitem{CohenSegal2011} D. Cohen and M. Segal, \emph{Network bursts in hippocampal microcultures are terminated by exhaustion of vesicle pools}, J. Neurophysiol. \textbf{106}, 2314 (2011).
\bibitem{Tsodyks2000} M. Tsodyks, A. Uziel, and H. Markram, \emph{Synchrony generation in recurrent networks with frequency-dependent synapses}, J. Neurosci. \textbf{20}, RC50 (2000).
\bibitem{Benson1994} D. L. Benson, F. H. Watkins, O. Steward, and G. Banker, \emph{Characterization of GABAergic neurons in hippocampal cell cultures}, J. Neurocytol. \textbf{23}, 279 (1994).

\input{supplement/bib_22}
\end{thebibliography}


\end{document}
